\documentclass[12pt,fontset=fandol]{ctexart}

\makeatletter
\def\input@path{{../}{../../}{../../../}{../../../../}}
\makeatother

\usepackage{researchnotes}
\usepackage{arydshln,mathdots}
\newtheorem{exercise}{训练题}[section]

\title{全国大学生数学竞赛高等代数与线性代数试题选编}
\author{历年真题与专题研习}
\date{2026年10月10日}

\begin{document}
\maketitle
\section*{编写说明与阅读路线}
\addcontentsline{toc}{section}{编写说明与阅读路线}
本书面向本科高等代数与线性代数课程，从全国大学生数学竞赛的历年试题出发，按知识依赖组织真题、解析、方法与训练。全国大学生数学竞赛的英文名称为 The Chinese Mathematics Competitions，简称 CMC；它与全国大学生数学建模竞赛是不同的比赛。本书不是整套试卷汇编，纯数学分析、概率统计、抽象群论和微分几何试题不在课程选编范围。数值线性代数以及直接使用二次型、行列式的交叉题列为进阶研习。

真题部分共收录 \textbf{100 道独立题目}，来自第1至17届的初赛与决赛，含第13届初赛补赛和第14届两次初赛补赛。初赛来源为数学A类及分设A、B类以前的数学专业卷；决赛保留高低年级组标识，共同题合并计数。本书不声称收齐数学B类或非数学类中的全部代数题。初赛括号内标明赛事年度，补赛的实际日期可能在下一自然年；决赛年份表示举办年。其中第13届决赛按中国数学会实际举办记录更正为2023年，说明及核查入口见文末。

题目编供独立练习，解析编以双向页码链接与之对应。同一题在题目编、解析编和专题索引中出现，只计一道真题。专题方法和补充训练另编，不计入上述100道真题。全书附入 398 道课程补充训练及解答，来源为已有的开放课程习题选编；训练题标题明确标注“补充训练”，用来补足真题所需的基础与证明工具。页数来自正文、推导和练习的实际内容，采用项目共享的A4版式。

推荐先做线性空间、秩与行列式，再学习特征值、多项式和循环结构，随后处理幂零、二次型、内积及矩阵方程。数值算法专题放在最后。做题时先记录数域、维数和非退化条件；计算题应代回核验，证明题应检查零矩阵、零空间、重根和参数特殊值。参考解答提供一种可核查路径，不表示该题只有这一种解法。

\subsection*{来源层级与许可}
真题题面与解析主要转录、改编自 Jacobi Lie 的\href{https://github.com/lzxcr/cmc-mathematics}{《全国大学生数学竞赛题解》}；方法编节选其\href{https://github.com/lzxcr/cmc-mathematics}{《数学竞赛的方法与结构》}。所用固定版本为 \href{https://github.com/lzxcr/cmc-mathematics/tree/64916fa1fdbc3e444a652eca204d5669058ee321}{64916fa1fdbc3e444a652eca204d5669058ee321}，检索与许可核验日期为2026年10月10日。原作者撰写的解答、评注和方法稿采用\href{https://creativecommons.org/licenses/by-sa/4.0/}{CC BY-SA 4.0}；本书对这一部分的专题重组、TeX转写和修订沿用同一许可。引用试题及第三方材料的权利仍归原权利人，此许可不改变试题的权利状态。

补充训练编是一份独立的课程材料节选，选自 yhwu-is 及项目贡献者的\href{https://github.com/yhwu-is/Linear-Algebra-Left-Undone}{《线性代数：未竟之美》}和本项目已有习题集的少量计算训练，保留\href{https://creativecommons.org/licenses/by-nc-sa/4.0/}{CC BY-NC-SA 4.0}许可、原作者署名和逐题链接。不同来源部分各自保留其许可，不能把一种许可扩张到全书全部材料。L编号指原讲义的讲号、题组与题号；S编号为项目先前编写的计算练习，不是竞赛原题。本次没有将该习题集附编中的其他教材题目收入。

本书的来源层级是“公开题解转录与课程选编”，不是竞赛组委会审定的官方版本，也未声称逐题核对所有原卷扫描件。题面中明确条件、记号与退化情形的规范化表述沿用开放题解；重要补充说明与勘误另行标明。题目保留永久来源标识、题解目录序号与逐题链接；目录序号可能将填空题拆开，不能直接等同于原卷大题号。高校原卷检索入口采用\href{https://gdjpkc.xmu.edu.cn/DocumentList.aspx?cID=22&dID=383}{厦门大学高等代数课程竞赛题选目录}。

\subsection*{记号与约定}
未另行说明时，$K$ 或 $\mathbb F$ 为数域，即复数域的子域；$M_{m\times n}(K)$ 是所有 $m$ 行 $n$ 列矩阵构成的空间，$M_n(K)=M_{n\times n}(K)$。$I$、$I_n$ 或 $E$ 表示单位矩阵，$0$ 表示相应维数的零对象。$r(A)$、$\operatorname{rank}A$ 表示秩（rank），$\operatorname{tr}A$ 表示迹（trace），$\ker T$、$\operatorname{im}T$ 表示核（kernel）与像（image），$\operatorname{span}$ 表示线性张成（span）。

$A^{\mathrm T}$ 表示转置（transpose），$\overline A$ 表示逐元共轭。伴随矩阵（adjugate）统一在说明中写作 $\operatorname{adj}A$；部分题面沿用 $A^*$，在该题中专指伴随矩阵。复内积中的伴随算子（adjoint）或共轭转置也可能记为 $A^*$，此时由内积语境明确定义，不能与代数余子式组成的伴随矩阵混淆。$\mu_A$ 与 $\chi_A$ 分别为最小多项式（minimal polynomial）与特征多项式（characteristic polynomial）。

中心化子（centralizer）$C(A)=\{X:XA=AX\}$ 是与 $A$ 可交换的矩阵组成的空间；换位子（commutator）$[A,B]=AB-BA$ 是给定两个矩阵算出的一个矩阵。部分来源把这两种语境中的对象都称作“交换子”，本书保留相应题面措辞，并通过集合定义或方括号记号区分。$X\mapsto AX-XA$ 则是矩阵空间上的线性算子，其核为 $C(A)$。

实二次型 $x^{\mathrm T}Ax$ 的对称矩阵为 $(A+A^{\mathrm T})/2$；正定矩阵指实对称正定矩阵或题目明确指定的埃尔米特（Hermitian）正定矩阵。若当标准形也称 Jordan 标准形。补充训练中的 $\mathbb F[x]_n$ 表示次数严格小于 $n$ 的多项式连同零多项式，维数为 $n$。正交变换保持内积，合同变换保持实二次型的惯性；“标准形”和“规范形”的系数约定在相应解答中说明。
\tableofcontents
\clearpage


\part{历年真题题目编}

\section{线性空间、线性相关与秩}

\label{sec:questions-1}

把矩阵看成线性映射，先确定核、像和维数。秩不等式的实质往往是某个限制映射的秩—零度公式；函数或多项式构成的向量组也应回到有限维张成空间研究。

\begin{example}[初赛第 6 届（2014 年度），来源序号3]\label{ex:pre06-math:3}

设 $V$ 为 $[0,1]$ 上全部实函数构成的实向量空间，$f_1,\ldots,f_n\in V$。证明以下两条等价：

（1）$f_1,\ldots,f_n$ 线性无关；

（2）存在 $a_1,\ldots,a_n\in[0,1]$，使 $\det[f_i(a_j)]_{i,j=1}^n\ne0$。

\par\hyperref[ans:pre06-math:3]{解析见第~\pageref{ans:pre06-math:3} 页。}\quad\href{https://github.com/lzxcr/cmc-mathematics/blob/64916fa1fdbc3e444a652eca204d5669058ee321/src/problems/preliminary/06/03.md}{来源条目 pre06-math:3}

\end{example}

\begin{example}[初赛第 8 届（2016 年度），来源序号3]\label{ex:pre08-math:3}

设 $A_1,\ldots,A_{2017}$ 为 $2016$ 阶实矩阵。证明方程
\[
\det(x_1A_1+\cdots+x_{2017}A_{2017})=0
\]
有非零实向量解 $(x_1,\ldots,x_{2017})$。

\par\hyperref[ans:pre08-math:3]{解析见第~\pageref{ans:pre08-math:3} 页。}\quad\href{https://github.com/lzxcr/cmc-mathematics/blob/64916fa1fdbc3e444a652eca204d5669058ee321/src/problems/preliminary/08/03.md}{来源条目 pre08-math:3}

\end{example}

\begin{example}[初赛第 14 届第二次补赛（2022 年度），来源序号4]\label{ex:ch06-14th-A-makeup2:4}

设 $A,B$ 都是秩为 $r$ 的 $n$ 阶不可逆实矩阵，$I$ 和 $J$ 是集合 $\{1,2,\cdots,n\}$ 的两个 $r+1$ 元子集。用 $\mathbb{R}^{n\times n}$ 表示所有 $n$ 阶实矩阵构成的集合，令
\[
V=\left\{C=\left(c_{ij}\right)_{n\times n}\in\mathbb{R}^{n\times n}\ \middle|\ c_{ij}=0,\ \text{若 } i\notin I \text{ 或 } j\notin J\right\}.
\]
证明：存在 $0\neq C\in V$ 使得 $ACB=0$。

\par\hyperref[ans:ch06-14th-A-makeup2:4]{解析见第~\pageref{ans:ch06-14th-A-makeup2:4} 页。}\quad\href{https://github.com/lzxcr/cmc-mathematics/blob/64916fa1fdbc3e444a652eca204d5669058ee321/src/problems/preliminary/14-makeup2/04.md}{来源条目 ch06-14th-A-makeup2:4}

\end{example}

\begin{example}[初赛第 16 届（2024 年度），来源序号4]\label{ex:ch02-16th-A:4}

已知实数域 $\mathbb{R}$ 上的一元多项式集合 $\mathbb{R}[x]$ 在多项式加法和数乘下构成 $\mathbb{R}$ 上的一个线性空间。设 $f_{i}(x) \in \mathbb{R}[x]$ 且次数为 $n_{i}, 1 \leq i \leq 2024$，规定零多项式的次数为 $-\infty$，已知 $\sum_{i=1}^{2024} n_{i} < 2047276$，证明：$f_{1}(x), f_{2}(x), \cdots, f_{2024}(x)$ 为空间 $\mathbb{R}[x]$ 中线性相关的向量组。

\par\hyperref[ans:ch02-16th-A:4]{解析见第~\pageref{ans:ch02-16th-A:4} 页。}\quad\href{https://github.com/lzxcr/cmc-mathematics/blob/64916fa1fdbc3e444a652eca204d5669058ee321/src/problems/preliminary/16/04.md}{来源条目 ch02-16th-A:4}

\end{example}

\begin{example}[决赛第 1 届（2010 年），来源序号11]\label{ex:fin01-all:11}

设 $V$ 为复数域上的有限维线性空间，$f_1,f_2\in V^*$ 线性无关。证明每个 $\alpha\in V$ 都可写成 $\alpha=\alpha_1+\alpha_2$，使
\[
f_1(\alpha)=f_1(\alpha_2),\qquad f_2(\alpha)=f_2(\alpha_1).
\]

\par\hyperref[ans:fin01-all:11]{解析见第~\pageref{ans:fin01-all:11} 页。}\quad\href{https://github.com/lzxcr/cmc-mathematics/blob/64916fa1fdbc3e444a652eca204d5669058ee321/src/problems/finals/01/11.md}{来源条目 fin01-all:11}

\end{example}

\begin{example}[决赛第 12 届（2021 年），来源序号6，共同题]\label{ex:fin12-high:6}

设 $n$ 阶复方阵 $A_{1},\cdots,A_{2n}$ 均相似于对角阵，$\mathbb{C}^{n}$ 表示复 $n$ 维列向量空间。 证明：

（1）$\mathbb{C}^{n}=\ker A_{k}\oplus\operatorname{Im} A_{k}$。 这里
$\ker A_{k}=\left\{\alpha\mid A_{k}\alpha=0,\alpha\in\mathbb{C}^{n}\right\}$，
$\operatorname{Im} A_{k}=\left\{A_{k}\beta\mid\beta\in\mathbb{C}^{n}\right\}(k=1,\cdots,2n)$。

（2）若对所有的 $k<j$ 皆有 $A_{k}A_{j}=0\,(k,j=1,2,\cdots,2n)$，则 $A_{1},\cdots,A_{2n}$ 中至少有 $n$ 个矩阵为零矩阵。

\par\hyperref[ans:fin12-high:6]{解析见第~\pageref{ans:fin12-high:6} 页。}\quad\href{https://github.com/lzxcr/cmc-mathematics/blob/64916fa1fdbc3e444a652eca204d5669058ee321/src/problems/finals/12/06.md}{来源条目 fin12-high:6}

\end{example}

\begin{example}[决赛第 15 届（2024 年），来源序号6，共同题]\label{ex:fin15-high:6}

设 $A \in \mathbb{R}^{n\times n}$，$E$ 为单位矩阵。求证 $A^{4} = E$ 当且仅当 $\operatorname{rank}(E - A) + \operatorname{rank}\left(E + A + A^{2} + A^{3}\right) = n$。

\par\hyperref[ans:fin15-high:6]{解析见第~\pageref{ans:fin15-high:6} 页。}\quad\href{https://github.com/lzxcr/cmc-mathematics/blob/64916fa1fdbc3e444a652eca204d5669058ee321/src/problems/finals/15/06.md}{来源条目 fin15-high:6}

\end{example}

\section{行列式、伴随矩阵与多重线性}

\label{sec:questions-2}

先观察行列式的结构，再选择初等变换、分块或低秩分解。伴随矩阵题应先按原矩阵的秩分类；对三维向量的混合积，可以用行列式的多重线性统一处理。

\begin{example}[初赛第 16 届（2024 年度），来源序号3]\label{ex:ch02-16th-A:3}

设 $A=[a_{ij}]_{3 \times 3}$ 为实数域 $\mathbb{R}$ 上的 3 阶不可逆方阵。若 $A$ 的伴随矩阵 $A^{*}=[a_{ij}^{2}]_{3 \times 3}$，求 $A$。

\par\hyperref[ans:ch02-16th-A:3]{解析见第~\pageref{ans:ch02-16th-A:3} 页。}\quad\href{https://github.com/lzxcr/cmc-mathematics/blob/64916fa1fdbc3e444a652eca204d5669058ee321/src/problems/preliminary/16/03.md}{来源条目 ch02-16th-A:3}

\end{example}

\begin{example}[初赛第 17 届（2025 年度），来源序号3]\label{ex:ch01-17th-A:3}

设 $m, n$ 为大于 2 的整数，$a_1, \dots, a_{m+1}$ 为任意 $m+1$ 个有理数，$\mathbb Q^{n\times n}$ 为有理数域上 $n$ 阶方阵全体。证明：

（1）$\mathbb Q^{n\times n}$ 中存在 $m$ 个元素 $B_1, \dots, B_m$，使得行列式 $\left|B_j\right| = j(j=1, \dots, m)$ 成立。

（2）$\mathbb Q^{n\times n}$ 中存在 $m$ 个元素 $A_1, \dots, A_m$，使得下列两条同时成立：

(i) $\left|A_j\right| = a_j(j=1, \dots, m)$；(ii) $\left|A_1 - A_2 - \dots - A_m\right| = a_{m+1}$。

\par\hyperref[ans:ch01-17th-A:3]{解析见第~\pageref{ans:ch01-17th-A:3} 页。}\quad\href{https://github.com/lzxcr/cmc-mathematics/blob/64916fa1fdbc3e444a652eca204d5669058ee321/src/problems/preliminary/17/03.md}{来源条目 ch01-17th-A:3}

\end{example}

\begin{example}[决赛第 1 届（2010 年），来源序号4]\label{ex:fin01-all:4}

设 $\boldsymbol a,\boldsymbol b,\boldsymbol c\in\mathbb R^3$ 满足
$(\boldsymbol a\times\boldsymbol b)\cdot\boldsymbol c=6$，求
\[
\bigl[(\boldsymbol a+\boldsymbol b)\times(\boldsymbol b+\boldsymbol c)\bigr]
\cdot(\boldsymbol a+\boldsymbol c).
\]

\par\hyperref[ans:fin01-all:4]{解析见第~\pageref{ans:fin01-all:4} 页。}\quad\href{https://github.com/lzxcr/cmc-mathematics/blob/64916fa1fdbc3e444a652eca204d5669058ee321/src/problems/finals/01/04.md}{来源条目 fin01-all:4}

\end{example}

\begin{example}[决赛第 4 届（2013 年），来源序号3]\label{ex:fin04-all:3}

设实矩阵 $A\in M_n(\mathbb R)$ 的每个元素的绝对值均为 $2$。证明当 $n\ge3$ 时，
\[
|\det A|\le\frac{2^{n+1}}3\,n!.
\]

\par\hyperref[ans:fin04-all:3]{解析见第~\pageref{ans:fin04-all:3} 页。}\quad\href{https://github.com/lzxcr/cmc-mathematics/blob/64916fa1fdbc3e444a652eca204d5669058ee321/src/problems/finals/04/03.md}{来源条目 fin04-all:3}

\end{example}

\begin{example}[决赛第 6 届（2015 年），来源序号4，共同题]\label{ex:fin06-lo:4}

（填空题）$A=(a_{ij})$ 为 $n$ 阶实对称矩阵（$n>1$），$\mathrm{rank}(A)=n-1$，$A$ 的每行元素之和均为 $0$。设 $2,3,\cdots,n$ 为 $A$ 的全部非零特征值。用 $A_{11}$ 表示 $A$ 的元素 $a_{11}$ 所对应的代数余子式，则有 $A_{11}=\underline{\qquad\qquad}$。

\par\hyperref[ans:fin06-lo:4]{解析见第~\pageref{ans:fin06-lo:4} 页。}\quad\href{https://github.com/lzxcr/cmc-mathematics/blob/64916fa1fdbc3e444a652eca204d5669058ee321/src/problems/finals/06/04.md}{来源条目 fin06-lo:4}

\end{example}

\begin{example}[决赛第 7 届（2016 年），来源序号1，共同题]\label{ex:fin07-hi:1}

（填空题）设 $\Gamma$ 为形如下列形式的 2016 阶矩阵全体：矩阵的每行每列只有一个非零元素，且该非零元素为 1，则 $\sum_{A\in\Gamma}\left|A\right|=$ $\underline{\qquad\qquad}$。

\par\hyperref[ans:fin07-hi:1]{解析见第~\pageref{ans:fin07-hi:1} 页。}\quad\href{https://github.com/lzxcr/cmc-mathematics/blob/64916fa1fdbc3e444a652eca204d5669058ee321/src/problems/finals/07/01.md}{来源条目 fin07-hi:1}

\end{example}

\begin{example}[决赛第 8 届（2017 年），来源序号1，共同题]\label{ex:fin08-hi:1}

（填空题）设 $x^{4}+3x^{2}+2x+1=0$ 的 4 个根为 $\alpha_{1},\alpha_{2},\alpha_{3},\alpha_{4}$，则
\[
\begin{vmatrix}
\alpha_{1} & \alpha_{2} & \alpha_{3} & \alpha_{4} \\
\alpha_{2} & \alpha_{3} & \alpha_{4} & \alpha_{1} \\
\alpha_{3} & \alpha_{4} & \alpha_{1} & \alpha_{2} \\
\alpha_{4} & \alpha_{1} & \alpha_{2} & \alpha_{3}
\end{vmatrix}=\underline{\hspace{2cm}}.
\]

\par\hyperref[ans:fin08-hi:1]{解析见第~\pageref{ans:fin08-hi:1} 页。}\quad\href{https://github.com/lzxcr/cmc-mathematics/blob/64916fa1fdbc3e444a652eca204d5669058ee321/src/problems/finals/08/01.md}{来源条目 fin08-hi:1}

\end{example}

\begin{example}[决赛第 12 届（2021 年），来源序号4，共同题]\label{ex:fin12-high:4}

（填空题）设 $A$ 为 2021 阶对称矩阵，$A$ 的每一行均为 $1,2,\cdots,2021$ 的一个排列，则 $A$ 的迹 $\operatorname{tr} A=$ $\underline{\qquad\qquad}$。

\par\hyperref[ans:fin12-high:4]{解析见第~\pageref{ans:fin12-high:4} 页。}\quad\href{https://github.com/lzxcr/cmc-mathematics/blob/64916fa1fdbc3e444a652eca204d5669058ee321/src/problems/finals/12/04.md}{来源条目 fin12-high:4}

\end{example}

\begin{example}[决赛第 16 届（2025 年），来源序号8，低年级组]\label{ex:fin16-low:8}

设矩阵 $M$ 及其子矩阵 $N,P,Q$ 分别为
$M = \begin{pmatrix} a_{11} & a_{12} & a_{13} \\ a_{21} & a_{22} & a_{23} \\ a_{31} & a_{32} & a_{33} \end{pmatrix}$，$N = \begin{pmatrix} a_{21} & a_{22} \\ a_{31} & a_{32} \end{pmatrix}$，$P = \begin{pmatrix} a_{31} & a_{32} \\ a_{11} & a_{12} \end{pmatrix}$，$Q = \begin{pmatrix} a_{11} & a_{12} \\ a_{21} & a_{22} \end{pmatrix}$。
若它们的行列式分别为 $|M| = 12$，$|N| = 1$，$|P| = 2$，$|Q| = 3$，求平面上由直线 $L_1: a_{11}x + a_{12}y + a_{13} = 0$，$L_2: a_{21}x + a_{22}y + a_{23} = 0$，$L_3: a_{31}x + a_{32}y + a_{33} = 0$ 所围成的三角形的面积。

\par\hyperref[ans:fin16-low:8]{解析见第~\pageref{ans:fin16-low:8} 页。}\quad\href{https://github.com/lzxcr/cmc-mathematics/blob/64916fa1fdbc3e444a652eca204d5669058ee321/src/problems/finals/16/08.md}{来源条目 fin16-low:8}

\end{example}

\section{特征值、迹与谱的代数约束}

\label{sec:questions-3}

迹、行列式和特征多项式系数连接特征值的幂和与对称函数。讨论谱之前先说明标量域，不能从特征值相同直接推出相似，也不能从实矩阵直接推出全部特征值为实数。

\begin{example}[初赛第 7 届（2015 年度），来源序号2]\label{ex:pre07-math:2}

设 $A$ 为 $4$ 阶复矩阵，满足 $\operatorname{tr}(A^i)=i$，$i=1,2,3,4$。求 $\det A$。

\par\hyperref[ans:pre07-math:2]{解析见第~\pageref{ans:pre07-math:2} 页。}\quad\href{https://github.com/lzxcr/cmc-mathematics/blob/64916fa1fdbc3e444a652eca204d5669058ee321/src/problems/preliminary/07/02.md}{来源条目 pre07-math:2}

\end{example}

\begin{example}[初赛第 11 届（2019 年度），来源序号3]\label{ex:ch10-11th-A:3}

设 $A$ 为 $n$ 阶复方阵，$p(x)$ 为 $I-\bar{A}A$ 的特征多项式，其中 $\bar{A}$ 表示 $A$ 的共轭矩阵。证明：$p(x)$ 必为实系数多项式。

\par\hyperref[ans:ch10-11th-A:3]{解析见第~\pageref{ans:ch10-11th-A:3} 页。}\quad\href{https://github.com/lzxcr/cmc-mathematics/blob/64916fa1fdbc3e444a652eca204d5669058ee321/src/problems/preliminary/11/03.md}{来源条目 ch10-11th-A:3}

\end{example}

\begin{example}[初赛第 13 届补赛（2021 年度），来源序号4]\label{ex:ch08-13th-A-makeup:4}

设 $a_{1},a_{2},a_{3}$ 为满足 $a_{1}^{2}+a_{2}^{2}+a_{3}^{2}=1$ 的一组实数，$b_{1},b_{2}$ 为满足 $b_{1}^{2}+b_{2}^{2}=1$ 的一组实数。又设 $M_{1}$ 为 $5\times 3$ 矩阵，其每一行都为 $a_{1},a_{2},a_{3}$ 的一个排列；$M_{2}$ 是 $5\times 2$ 矩阵，其每一行都为 $b_{1},b_{2}$ 的一个排列。令 $M=\left(M_{1},M_{2}\right)$，它为 $5\times 5$ 矩阵。证明：
（1）$(\operatorname{tr}M)^{2}\leq(5+2\sqrt{6})\operatorname{rank}M$；
（2）$M$ 必有绝对值小于或等于 $\sqrt{2}+\sqrt{3}$ 的实特征值 $\lambda$。

\par\hyperref[ans:ch08-13th-A-makeup:4]{解析见第~\pageref{ans:ch08-13th-A-makeup:4} 页。}\quad\href{https://github.com/lzxcr/cmc-mathematics/blob/64916fa1fdbc3e444a652eca204d5669058ee321/src/problems/preliminary/13-makeup/04.md}{来源条目 ch08-13th-A-makeup:4}

\end{example}

\begin{example}[初赛第 14 届补赛（2022 年度），来源序号3]\label{ex:ch05-14th-A-makeup:3}

设 $n \geq 2$，$A$ 为 $n$ 阶实对称阵，$v = \left(a_1,a_2,\cdots,a_n\right)^{\mathrm{T}}$，$w = \left(b_1,b_2,\cdots,b_n\right)^{\mathrm{T}}$，分别是 $A$ 关于其特征值 $\mu_1, \mu_2$ 的实特征向量，$\mu_1 \neq \mu_2$。令

\[
B = \begin{bmatrix}
a_1+b_1 & a_1+b_2 & \cdots & a_1+b_n \\
a_2+b_1 & a_2+b_2 & \cdots & a_2+b_n \\
\vdots & \vdots & & \vdots \\
a_n+b_1 & a_n+b_2 & \cdots & a_n+b_n
\end{bmatrix}.
\]

求矩阵 $B$ 的所有特征值。

\par\hyperref[ans:ch05-14th-A-makeup:3]{解析见第~\pageref{ans:ch05-14th-A-makeup:3} 页。}\quad\href{https://github.com/lzxcr/cmc-mathematics/blob/64916fa1fdbc3e444a652eca204d5669058ee321/src/problems/preliminary/14-makeup/03.md}{来源条目 ch05-14th-A-makeup:3}

\end{example}

\begin{example}[初赛第 15 届（2023 年度），来源序号4]\label{ex:ch03-15th-A:4}

设 $n \geq 3$ 为自然数，$\theta = \frac{2\pi}{n}$。对任意 $1 \leq s,t \leq n$，取 $a_{st} = \sin((s+t)\theta)$，令矩阵 $A = \left(a_{st}\right)_{n \times n}$，计算 $E + A^{2023}$ 的行列式，其中 $E$ 为 $n$ 阶单位矩阵。

\par\hyperref[ans:ch03-15th-A:4]{解析见第~\pageref{ans:ch03-15th-A:4} 页。}\quad\href{https://github.com/lzxcr/cmc-mathematics/blob/64916fa1fdbc3e444a652eca204d5669058ee321/src/problems/preliminary/15/04.md}{来源条目 ch03-15th-A:4}

\end{example}

\begin{example}[决赛第 2 届（2011 年），来源序号4]\label{ex:fin02-all:4}

设 $A\in M_n(\mathbb C)$，定义线性变换
\[
\sigma_A:M_n(\mathbb C)\to M_n(\mathbb C),\qquad
\sigma_A(X)=AX-XA.
\]
证明：若 $A$ 可对角化，则 $\sigma_A$ 也可对角化。

\par\hyperref[ans:fin02-all:4]{解析见第~\pageref{ans:fin02-all:4} 页。}\quad\href{https://github.com/lzxcr/cmc-mathematics/blob/64916fa1fdbc3e444a652eca204d5669058ee321/src/problems/finals/02/04.md}{来源条目 fin02-all:4}

\end{example}

\begin{example}[决赛第 5 届（2014 年），来源序号2，共同题]\label{ex:fin05-hi:2}

设实二次型 $q(x)=x^{\mathsf T}Ax$，其中
\[
A=\begin{pmatrix}
2&a_0&2&-2\\
a&0&b&c\\
d&e&0&f\\
g&h&k&4
\end{pmatrix},
\]
各参数均为实数，且特征值 $2$ 的几何重数为 $3$。

（1）判断 $A$ 是否一定可对角化，并证明。
（2）当 $a_0=2$ 时，求 $q$ 在正交变换下的标准型。

\par\hyperref[ans:fin05-hi:2]{解析见第~\pageref{ans:fin05-hi:2} 页。}\quad\href{https://github.com/lzxcr/cmc-mathematics/blob/64916fa1fdbc3e444a652eca204d5669058ee321/src/problems/finals/05/02.md}{来源条目 fin05-hi:2}

\end{example}

\begin{example}[决赛第 9 届（2018 年），来源序号1，共同题]\label{ex:fin09-hi:1}

（填空题）设实方阵 $H_{1}=\begin{pmatrix}0&1\\1&0\end{pmatrix}$，$H_{n+1}=\begin{pmatrix}H_{n}&I\\I&H_{n}\end{pmatrix}$，$n \geq 1$，其中 $I$ 是与 $H_{n}$ 同阶的单位方阵，则 $\operatorname{rank}(H_{4})=$ $\underline{\qquad\qquad}$。

\par\hyperref[ans:fin09-hi:1]{解析见第~\pageref{ans:fin09-hi:1} 页。}\quad\href{https://github.com/lzxcr/cmc-mathematics/blob/64916fa1fdbc3e444a652eca204d5669058ee321/src/problems/finals/09/01.md}{来源条目 fin09-hi:1}

\end{example}

\begin{example}[决赛第 13 届（2023 年），来源序号3，共同题]\label{ex:fin13-high:3}

设 $A$ 为 $n \times n$ 实矩阵，$m$ 为正整数，$I$ 为 $n \times n$ 单位矩阵。若 $(I+A)^m = 0$，则 $A$ 的行列式 $|A| = \underline{\qquad}$。

\par\hyperref[ans:fin13-high:3]{解析见第~\pageref{ans:fin13-high:3} 页。}\quad\href{https://github.com/lzxcr/cmc-mathematics/blob/64916fa1fdbc3e444a652eca204d5669058ee321/src/problems/finals/13/03.md}{来源条目 fin13-high:3}

\end{example}

\begin{example}[决赛第 14 届（2023 年），来源序号8，低年级组]\label{ex:fin14-low:8}

设 $A$ 为复数域上的 4 阶幂等阵（即 $A^2 = A$）。证明：存在复数 $c_0, c_1, c_2, c_3$ 使得 $A$ 相似于 $\begin{bmatrix} c_0 & c_1 & c_2 & c_3 \\ c_3 & c_0 & c_1 & c_2 \\ c_2 & c_3 & c_0 & c_1 \\ c_1 & c_2 & c_3 & c_0 \end{bmatrix}.$

\par\hyperref[ans:fin14-low:8]{解析见第~\pageref{ans:fin14-low:8} 页。}\quad\href{https://github.com/lzxcr/cmc-mathematics/blob/64916fa1fdbc3e444a652eca204d5669058ee321/src/problems/finals/14/08.md}{来源条目 fin14-low:8}

\end{example}

\begin{example}[决赛第 17 届（2026 年），来源序号8，低年级组]\label{ex:fin17-low:6}

设 $A,B$ 分别为 $n$ 阶实矩阵，$E$ 为 $n$ 阶单位矩阵，满足 $A^{2025}=E$，$B^{2026}=E$，$AB=BA$。证明：$A+B+E$ 可逆。

\par\hyperref[ans:fin17-low:6]{解析见第~\pageref{ans:fin17-low:6} 页。}\quad\href{https://github.com/lzxcr/cmc-mathematics/blob/64916fa1fdbc3e444a652eca204d5669058ee321/src/problems/finals/17/08.md}{来源条目 fin17-low:6}

\end{example}

\section{多项式、整除、插值与根的约束}

\label{sec:questions-4}

本专题研究系数、根与整除关系，包含有理数域上的不可约性。将次数界和线性代数结合，经常能把无穷多个取值条件化为有限个独立条件；整系数问题还必须使用整数性。

\begin{example}[初赛第 5 届（2013 年度），来源序号2]\label{ex:pre05-math:2}

设 $B(t)=(b_{ij}(t))$ 为 $n$ 阶矩阵，$b(t)=(b_1(t),\ldots,b_n(t))^{\mathrm T}$，其中各元素均为实系数多项式。记 $d(t)=\det B(t)$，$d_i(t)$ 为用 $b(t)$ 替换 $B(t)$ 第 $i$ 列所得行列式。若 $d(t_0)=0$，且方程 $B(t_0)X=b(t_0)$ 相容，证明 $d,d_1,\ldots,d_n$ 有次数至少为 $1$ 的公因式。

\par\hyperref[ans:pre05-math:2]{解析见第~\pageref{ans:pre05-math:2} 页。}\quad\href{https://github.com/lzxcr/cmc-mathematics/blob/64916fa1fdbc3e444a652eca204d5669058ee321/src/problems/preliminary/05/02.md}{来源条目 pre05-math:2}

\end{example}

\begin{example}[初赛第 10 届（2018 年度），来源序号3]\label{ex:pre10-math:3}

称元素全为整数的矩阵为整矩阵。设 $A,B$ 为 $n$ 阶整矩阵。

（1）证明：$A$ 可逆且 $A^{-1}$ 为整矩阵，当且仅当 $|\det A|=1$。

（2）若 $A-2kB$ 对每个 $k=0,1,\ldots,2n$ 均可逆，且它们的逆矩阵均为整矩阵，证明 $A+B$ 可逆。

\par\hyperref[ans:pre10-math:3]{解析见第~\pageref{ans:pre10-math:3} 页。}\quad\href{https://github.com/lzxcr/cmc-mathematics/blob/64916fa1fdbc3e444a652eca204d5669058ee321/src/problems/preliminary/10/03.md}{来源条目 pre10-math:3}

\end{example}

\begin{example}[初赛第 12 届（2020 年度），来源序号4]\label{ex:ch09-12th-A:4}

称非常值一元 $n$ 次多项式（合并同类项后）的 $n-1$ 次项（可能为 0）为第二项。求所有 2020 次复系数首一多项式 $f(x)$，满足对 $f(x)$ 的每个复根 $x_{k}$，都存在非常值复系数首一多项式 $g_{k}(x)$ 和 $h_{k}(x)$，使得
\[
f(x)=\left(x-x_{k}\right)g_{k}(x)h_{k}(x),
\]
且 $g_{k}(x)$ 和 $h_{k}(x)$ 的第二项系数相等。

\par\hyperref[ans:ch09-12th-A:4]{解析见第~\pageref{ans:ch09-12th-A:4} 页。}\quad\href{https://github.com/lzxcr/cmc-mathematics/blob/64916fa1fdbc3e444a652eca204d5669058ee321/src/problems/preliminary/12/04.md}{来源条目 ch09-12th-A:4}

\end{example}

\begin{example}[初赛第 13 届（2021 年度），来源序号3]\label{ex:ch07-13th-A:3}

设 $f(x)=x^{2021}+a_{2020}x^{2020}+a_{2019}x^{2019}+\cdots+a_{2}x^{2}+a_{1}x+a_{0}$ 为整系数多项式，$a_{0}\neq 0$。设对任意 $0\leq k\leq 2020$ 有 $\left|a_{k}\right|\leq 40$，证明：$f(x)=0$ 的根不可能全为实数。

\par\hyperref[ans:ch07-13th-A:3]{解析见第~\pageref{ans:ch07-13th-A:3} 页。}\quad\href{https://github.com/lzxcr/cmc-mathematics/blob/64916fa1fdbc3e444a652eca204d5669058ee321/src/problems/preliminary/13/03.md}{来源条目 ch07-13th-A:3}

\end{example}

\begin{example}[初赛第 14 届（2022 年度），来源序号4]\label{ex:ch04-14th-A:4}

设 $n\geq 2$ 为正整数，证明多项式 $f(x)=x^{n}-x-1$ 在有理数域 $\mathbb{Q}$ 上不可约。

\par\hyperref[ans:ch04-14th-A:4]{解析见第~\pageref{ans:ch04-14th-A:4} 页。}\quad\href{https://github.com/lzxcr/cmc-mathematics/blob/64916fa1fdbc3e444a652eca204d5669058ee321/src/problems/preliminary/14/04.md}{来源条目 ch04-14th-A:4}

\end{example}

\begin{example}[决赛第 5 届（2014 年），来源序号4，共同题]\label{ex:fin05-hi:4}

对所有根均为实数的多项式 $p$，以 $d(p)$ 表示其最大根与最小根之差。设 $n\ge2$，求使每个这样的 $n$ 次多项式均满足
\[
d(p')\ge C\,d(p)
\]
的最大常数 $C$。

\par\hyperref[ans:fin05-hi:4]{解析见第~\pageref{ans:fin05-hi:4} 页。}\quad\href{https://github.com/lzxcr/cmc-mathematics/blob/64916fa1fdbc3e444a652eca204d5669058ee321/src/problems/finals/05/04.md}{来源条目 fin05-hi:4}

\end{example}

\begin{example}[决赛第 8 届（2017 年），来源序号2，共同题]\label{ex:fin08-hi:2}

（填空题）设 $a$ 为实数，关于 $x$ 的方程
\[
3x^{4}-8x^{3}-30x^{2}+72x+a=0
\]
有虚根的充分必要条件是 $a$ 满足 $\underline{\qquad\qquad}$。

\par\hyperref[ans:fin08-hi:2]{解析见第~\pageref{ans:fin08-hi:2} 页。}\quad\href{https://github.com/lzxcr/cmc-mathematics/blob/64916fa1fdbc3e444a652eca204d5669058ee321/src/problems/finals/08/02.md}{来源条目 fin08-hi:2}

\end{example}

\begin{example}[决赛第 16 届（2025 年），来源序号6，共同题]\label{ex:fin16-high:6}

设实系数多项式 $p(x)$ 和 $q(x)$ 分别为：
$p(x) = x^n + a_1 x^{n-1} + \cdots + a_n$，$q(x) = x^m + b_1 x^{m-1} + \cdots + b_m$。
已知 $q(x) \mid p(x)$，且对某个 $k$，有 $|b_k| > 2025^k C_m^k$，其中 $C_m^k$ 为组合数。证明：存在 $j$，使得 $|a_j| > 2024$。

\par\hyperref[ans:fin16-high:6]{解析见第~\pageref{ans:fin16-high:6} 页。}\quad\href{https://github.com/lzxcr/cmc-mathematics/blob/64916fa1fdbc3e444a652eca204d5669058ee321/src/problems/finals/16/06.md}{来源条目 fin16-high:6}

\end{example}

\section{循环向量、最小多项式与中心化子}

\label{sec:questions-5}

有循环向量时，一个与原变换交换的算子由它在该向量上的值完全确定。最小多项式不仅控制可对角化性，还决定矩阵多项式化简、中心化子空间以及插值构造的维数。

\begin{example}[初赛第 1 届（2009 年度），来源序号2]\label{ex:pre01-math:2}

设
\[
F=\begin{pmatrix}
0&0&\cdots&0&-a_n\\
1&0&\cdots&0&-a_{n-1}\\
0&1&\cdots&0&-a_{n-2}\\
\vdots&\vdots&\ddots&\vdots&\vdots\\
0&0&\cdots&1&-a_1
\end{pmatrix}\in M_n(\mathbb C).
\]

（1） 若 $A=(a_{ij})$ 满足 $AF=FA$，证明
\[
A=a_{11}I+a_{21}F+\cdots+a_{n1}F^{n-1}.
\]

（2） 求交换子空间
\[
C(F)=\{X\in M_n(\mathbb C):XF=FX\}
\]
的维数。

\par\hyperref[ans:pre01-math:2]{解析见第~\pageref{ans:pre01-math:2} 页。}\quad\href{https://github.com/lzxcr/cmc-mathematics/blob/64916fa1fdbc3e444a652eca204d5669058ee321/src/problems/preliminary/01/02.md}{来源条目 pre01-math:2}

\end{example}

\begin{example}[初赛第 3 届（2011 年度），来源序号3]\label{ex:pre03-math:3}

设 $\mathbb F$ 是数域，线性变换 $\sigma:\mathbb F^n\to\mathbb F^n$ 满足
\[
\sigma(Av)=A\sigma(v)
\qquad(A\in M_n(\mathbb F),\ v\in\mathbb F^n).
\]
证明 $\sigma=\lambda I$，其中 $\lambda\in\mathbb F$。

\par\hyperref[ans:pre03-math:3]{解析见第~\pageref{ans:pre03-math:3} 页。}\quad\href{https://github.com/lzxcr/cmc-mathematics/blob/64916fa1fdbc3e444a652eca204d5669058ee321/src/problems/preliminary/03/03.md}{来源条目 pre03-math:3}

\end{example}

\begin{example}[初赛第 14 届（2022 年度），来源序号3]\label{ex:ch04-14th-A:3}

设 $A=\begin{pmatrix}2&1\\1&1\end{pmatrix}$，矩阵 $B$ 与 $A$ 可交换，其元素均为正整数且行列式为 1。证明存在正整数 $k$ 使得 $B=A^{k}$。

\par\hyperref[ans:ch04-14th-A:3]{解析见第~\pageref{ans:ch04-14th-A:3} 页。}\quad\href{https://github.com/lzxcr/cmc-mathematics/blob/64916fa1fdbc3e444a652eca204d5669058ee321/src/problems/preliminary/14/03.md}{来源条目 ch04-14th-A:3}

\end{example}

\begin{example}[初赛第 14 届第二次补赛（2022 年度），来源序号3]\label{ex:ch06-14th-A-makeup2:3}

设 $n\geq 2$，$A_{1},A_{2},\cdots,A_{n}$ 为数域 $K$ 上的方阵，它们的极小多项式两两互素。证明：给定数域 $K$ 上的任意多项式 $f_{1}(x),f_{2}(x),\cdots,f_{n}(x)\in K[x]$，存在多项式 $f(x)\in K[x]$ 使得对所有 $i=1,2,\cdots,n$，有
\[
f\left(A_{i}\right)=f_{i}\left(A_{i}\right).
\]

\par\hyperref[ans:ch06-14th-A-makeup2:3]{解析见第~\pageref{ans:ch06-14th-A-makeup2:3} 页。}\quad\href{https://github.com/lzxcr/cmc-mathematics/blob/64916fa1fdbc3e444a652eca204d5669058ee321/src/problems/preliminary/14-makeup2/03.md}{来源条目 ch06-14th-A-makeup2:3}

\end{example}

\begin{example}[初赛第 15 届（2023 年度），来源序号3]\label{ex:ch03-15th-A:3}

设 $V$ 是复数域 $\mathbb{C}$ 上的 $n$ 维线性空间，$A$ 是 $V$ 上的一个线性变换。证明：存在 $\alpha \in V$ 使得 $\left\{\alpha, A\alpha, \cdots, A^{n-1}\alpha\right\}$ 成为 $V$ 的一组基当且仅当对于 $A$ 的任一特征值 $\lambda$ 的几何重数为 1。

\par\hyperref[ans:ch03-15th-A:3]{解析见第~\pageref{ans:ch03-15th-A:3} 页。}\quad\href{https://github.com/lzxcr/cmc-mathematics/blob/64916fa1fdbc3e444a652eca204d5669058ee321/src/problems/preliminary/15/03.md}{来源条目 ch03-15th-A:3}

\end{example}

\begin{example}[初赛第 17 届（2025 年度），来源序号4]\label{ex:ch01-17th-A:4}

设 $K$ 为数域，$\mathcal{A}: V \rightarrow V$ 是 $n$ 维 $K$-向量空间 $V$ 上的线性算子（线性变换），$\mu_A(x), \chi_A(x) \in K[x]$ 分别是 $\mathcal{A}$ 的极小多项式和特征多项式。证明：

（1）存在 $\alpha \in V$，使得 $\mu_A(x) = \mu_{A,\alpha}(x)$，其中 $\mu_{A,\alpha}(x)$ 是集合 $\{f(x) \in K[x] \mid f(\mathcal{A})(\alpha) = 0\}$ 中首项系数为 1，次数最小的多项式；

（2）$K[\mathcal{A}] \cdot \alpha := \{f(\mathcal{A})(\alpha) \mid f(x) \in K[x]\}$ 是 $\deg \mu_{A,\alpha}(x)$ 维 $\mathcal{A}$-不变子空间；

（3）设 $\mathcal{B}: V \rightarrow V$ 是任意线性算子。若 $\mu_A(x) = \chi_A(x)$，则 $\mathcal{A} \cdot \mathcal{B} = \mathcal{B} \cdot \mathcal{A} \Leftrightarrow$ 存在 $f(x) \in K[x]$，使得 $\mathcal{B} = f(\mathcal{A})$。

\par\hyperref[ans:ch01-17th-A:4]{解析见第~\pageref{ans:ch01-17th-A:4} 页。}\quad\href{https://github.com/lzxcr/cmc-mathematics/blob/64916fa1fdbc3e444a652eca204d5669058ee321/src/problems/preliminary/17/04.md}{来源条目 ch01-17th-A:4}

\end{example}

\begin{example}[决赛第 5 届（2014 年），来源序号6，低年级组]\label{ex:fin05-lo:3}

设 $A=(a_{ij})\in M_n(\mathbb R)$ 满足 $a_{ij}=0$（$j>i+1$），且每个上次对角元 $a_{i,i+1}\ne0$。已知 $A$ 有 $n$ 个线性无关的特征向量。证明
\[
W=\{X\in M_n(\mathbb R):XA=AX\}
\]
是实向量空间，且 $I,A,\ldots,A^{n-1}$ 为其一组基。

\par\hyperref[ans:fin05-lo:3]{解析见第~\pageref{ans:fin05-lo:3} 页。}\quad\href{https://github.com/lzxcr/cmc-mathematics/blob/64916fa1fdbc3e444a652eca204d5669058ee321/src/problems/finals/05/06.md}{来源条目 fin05-lo:3}

\end{example}

\begin{example}[决赛第 13 届（2023 年），来源序号6，共同题]\label{ex:fin13-high:6}

给定 $n$ 阶复方阵 $A \neq 0, b \in \mathbb{C}^n$ 为 $n$ 维列向量。复系数多项式 $f(x)$ 被称为向量 $b$ 关于 $A$ 的零化多项式是指 $f(x)$ 满足 $f(A)b = 0$ 且 $f(x) \neq 0$。在向量 $b$ 关于 $A$ 的零化多项式中其次数最低的首 1 多项式称为向量 $b$ 关于 $A$ 的最小多项式。现设 $p(\lambda)$ 为矩阵 $A$ 的最小多项式。证明：$p(\lambda)$ 必为 $\mathbb{C}^n$ 中某向量 $b$ 关于 $A$ 的最小多项式。

\par\hyperref[ans:fin13-high:6]{解析见第~\pageref{ans:fin13-high:6} 页。}\quad\href{https://github.com/lzxcr/cmc-mathematics/blob/64916fa1fdbc3e444a652eca204d5669058ee321/src/problems/finals/13/06.md}{来源条目 fin13-high:6}

\end{example}

\section{幂零结构、不变子空间与相似}

\label{sec:questions-6}

通过核与像的稳定、幂的秩和若当链识别矩阵结构。交换关系可以传递到不变子空间，但交换子本身不必为零；证明幂零性时应交代从迹的幂和到谱结论的依据。

\begin{example}[初赛第 1 届（2009 年度），来源序号3]\label{ex:pre01-math:3}

设 $V$ 是非零有限维复线性空间，线性变换 $f,g:V\to V$ 满足
\[
fg-gf=f.
\]
证明：$f$ 的全部特征值均为零，且 $f,g$ 有公共特征向量。

\par\hyperref[ans:pre01-math:3]{解析见第~\pageref{ans:pre01-math:3} 页。}\quad\href{https://github.com/lzxcr/cmc-mathematics/blob/64916fa1fdbc3e444a652eca204d5669058ee321/src/problems/preliminary/01/03.md}{来源条目 pre01-math:3}

\end{example}

\begin{example}[初赛第 2 届（2010 年度），来源序号2]\label{ex:pre02-math:2}

设
\[
B=\begin{pmatrix}0&10&30\\0&0&2010\\0&0&0\end{pmatrix}.
\]
证明不存在 $X\in M_3(\mathbb C)$ 满足 $X^2=B$。

\par\hyperref[ans:pre02-math:2]{解析见第~\pageref{ans:pre02-math:2} 页。}\quad\href{https://github.com/lzxcr/cmc-mathematics/blob/64916fa1fdbc3e444a652eca204d5669058ee321/src/problems/preliminary/02/02.md}{来源条目 pre02-math:2}

\end{example}

\begin{example}[初赛第 3 届（2011 年度），来源序号6]\label{ex:pre03-math:6}

设 $A\in M_n(\mathbb F)$，其中 $\mathbb F$ 是数域。证明 $A$ 相似于
\[
\begin{pmatrix}B&0\\0&C\end{pmatrix},
\]
其中 $B$ 可逆、$C$ 幂零。允许其中一个块的维数为零。

\par\hyperref[ans:pre03-math:6]{解析见第~\pageref{ans:pre03-math:6} 页。}\quad\href{https://github.com/lzxcr/cmc-mathematics/blob/64916fa1fdbc3e444a652eca204d5669058ee321/src/problems/preliminary/03/06.md}{来源条目 pre03-math:6}

\end{example}

\begin{example}[初赛第 8 届（2016 年度），来源序号2]\label{ex:pre08-math:2}

设 $n$ 为奇数，实矩阵 $A,B\in M_n(\mathbb R)$ 满足 $BA=0$。分别取它们的 Jordan 标准形 $J_A,J_B$，记 $A+J_A$、$B+J_B$ 的特征值集合为 $S_1,S_2$。证明 $0\in S_1\cup S_2$。

\par\hyperref[ans:pre08-math:2]{解析见第~\pageref{ans:pre08-math:2} 页。}\quad\href{https://github.com/lzxcr/cmc-mathematics/blob/64916fa1fdbc3e444a652eca204d5669058ee321/src/problems/preliminary/08/02.md}{来源条目 pre08-math:2}

\end{example}

\begin{example}[初赛第 14 届补赛（2022 年度），来源序号4]\label{ex:ch05-14th-A-makeup:4}

设 $A,B$ 为数域 $K$ 上的两个 $n$ 阶矩阵。若存在正整数 $m$ 使得 $AB-BA=A^m$，证明对任何正整数 $k$，矩阵 $A^k+E$ 的行列式等于 1，其中 $E$ 为 $n$ 阶单位矩阵。

\par\hyperref[ans:ch05-14th-A-makeup:4]{解析见第~\pageref{ans:ch05-14th-A-makeup:4} 页。}\quad\href{https://github.com/lzxcr/cmc-mathematics/blob/64916fa1fdbc3e444a652eca204d5669058ee321/src/problems/preliminary/14-makeup/04.md}{来源条目 ch05-14th-A-makeup:4}

\end{example}

\begin{example}[决赛第 9 届（2018 年），来源序号6，共同题]\label{ex:fin09-hi:6}

设 $A, B, C$ 均为 $n$ 阶复方阵，且满足
\[
AB - BA = C, \quad AC = CA, \quad BC = CB.
\]

（1）证明：$C$ 是幂零方阵；

（2）证明：$A, B, C$ 同时相似于上三角阵；

（3）若 $C \neq 0$，求 $n$ 的最小值。

\par\hyperref[ans:fin09-hi:6]{解析见第~\pageref{ans:fin09-hi:6} 页。}\quad\href{https://github.com/lzxcr/cmc-mathematics/blob/64916fa1fdbc3e444a652eca204d5669058ee321/src/problems/finals/09/06.md}{来源条目 fin09-hi:6}

\end{example}

\begin{example}[决赛第 10 届（2019 年），来源序号6，共同题]\label{ex:fin10-hi:6}

（满分 15 分）证明：任意 $n$ 阶实方阵 $A$ 可以分解成 $A = A_{0} + A_{1} + A_{2}$，其中 $A_{0} = aI_{n}$，$a$ 是实数，$A_{1}$ 与 $A_{2}$ 都是幂零方阵。

\par\hyperref[ans:fin10-hi:6]{解析见第~\pageref{ans:fin10-hi:6} 页。}\quad\href{https://github.com/lzxcr/cmc-mathematics/blob/64916fa1fdbc3e444a652eca204d5669058ee321/src/problems/finals/10/06.md}{来源条目 fin10-hi:6}

\end{example}

\begin{example}[决赛第 12 届（2021 年），来源序号3，共同题]\label{ex:fin12-high:3}

（填空题）矩阵 $\begin{pmatrix} 1 & 0 & 1 & 0 & 0 \\ 0 & 1 & 0 & 1 & 0 \\ 0 & 0 & 1 & 0 & 1 \\ 0 & 0 & 0 & 1 & 0 \\ 0 & 0 & 0 & 0 & 1 \end{pmatrix}$ 的 Jordan 标准型为$\underline{\qquad\qquad}$。

\par\hyperref[ans:fin12-high:3]{解析见第~\pageref{ans:fin12-high:3} 页。}\quad\href{https://github.com/lzxcr/cmc-mathematics/blob/64916fa1fdbc3e444a652eca204d5669058ee321/src/problems/finals/12/03.md}{来源条目 fin12-high:3}

\end{example}

\begin{example}[决赛第 14 届（2023 年），来源序号3，共同题]\label{ex:fin14-high:3}

设 $A$ 为 $n \times n$ 实矩阵。若 $A^n=0$ 但 $A^{n-1} \neq 0$，则 $A$ 的秩为$\underline{\qquad}$。

\par\hyperref[ans:fin14-high:3]{解析见第~\pageref{ans:fin14-high:3} 页。}\quad\href{https://github.com/lzxcr/cmc-mathematics/blob/64916fa1fdbc3e444a652eca204d5669058ee321/src/problems/finals/14/03.md}{来源条目 fin14-high:3}

\end{example}

\begin{example}[决赛第 15 届（2024 年），来源序号8，低年级组]\label{ex:fin15-low:8}

设 $A$ 为 $n(n \geq 3)$ 阶实可逆矩阵，且 $A^{k}$ 相似于 $A$ 对 $k = 1,2,\cdots,n!-1,n!$ 成立，证明：对一切正整数 $k$ 皆有 $A^{k}$ 相似于 $A$。

\par\hyperref[ans:fin15-low:8]{解析见第~\pageref{ans:fin15-low:8} 页。}\quad\href{https://github.com/lzxcr/cmc-mathematics/blob/64916fa1fdbc3e444a652eca204d5669058ee321/src/problems/finals/15/08.md}{来源条目 fin15-low:8}

\end{example}

\section{二次型、合同与正定性}

\label{sec:questions-7}

二次型只由矩阵的对称部分决定。相似、合同和正交相似对应不同的问题；正定矩阵不等式通常先化为单位矩阵与一个正定矩阵的比较，再处理特征值或行列式。

\begin{example}[初赛第 2 届（2010 年度），来源序号6]\label{ex:pre02-math:6}

设 $A\in M_n(\mathbb R)$，对所有实行向量 $\alpha$ 有 $\alpha A\alpha^T\ge0$。设实行向量 $\beta$ 满足 $\beta A\beta^T=0$。另假定任意实行向量 $x,y$ 都满足
\[
xAy^T\ne0\quad\Longrightarrow\quad xAy^T+yAx^T\ne0.
\]
证明对每个实行向量 $v$，均有 $vA\beta^T=0$。

\par\hyperref[ans:pre02-math:6]{解析见第~\pageref{ans:pre02-math:6} 页。}\quad\href{https://github.com/lzxcr/cmc-mathematics/blob/64916fa1fdbc3e444a652eca204d5669058ee321/src/problems/preliminary/02/06.md}{来源条目 pre02-math:6}

\end{example}

\begin{example}[初赛第 4 届（2012 年度），来源序号4]\label{ex:pre04-math:4}

设 $A,B,C$ 是 $n$ 阶正定矩阵，令
\[
P(t)=At^2+Bt+C,\qquad f(t)=\det P(t).
\]
证明 $f$ 的每个复根 $\lambda$ 都满足 $\operatorname{Re}\lambda<0$。这里正定矩阵按通常约定指实对称正定矩阵，或更一般的 Hermitian 正定矩阵。

\par\hyperref[ans:pre04-math:4]{解析见第~\pageref{ans:pre04-math:4} 页。}\quad\href{https://github.com/lzxcr/cmc-mathematics/blob/64916fa1fdbc3e444a652eca204d5669058ee321/src/problems/preliminary/04/04.md}{来源条目 pre04-math:4}

\end{example}

\begin{example}[初赛第 9 届（2017 年度），来源序号4]\label{ex:pre09-math:4}

设 $a\ne0$ 为实数，$A$ 为 $n$ 阶实反对称矩阵，即 $A^{\mathrm T}=-A$。令
\[
\mathcal T=\{(X,Y)\in M_n(\mathbb R)^2:XY=aI+A\}.
\]
证明对任意 $(X,Y),(M,N)\in\mathcal T$，均有 $XN+Y^{\mathrm T}M^{\mathrm T}\ne0$。

\par\hyperref[ans:pre09-math:4]{解析见第~\pageref{ans:pre09-math:4} 页。}\quad\href{https://github.com/lzxcr/cmc-mathematics/blob/64916fa1fdbc3e444a652eca204d5669058ee321/src/problems/preliminary/09/04.md}{来源条目 pre09-math:4}

\end{example}

\begin{example}[初赛第 10 届（2018 年度），来源序号2]\label{ex:pre10-math:2}

设实矩阵 $A=(a_{ij})\in M_n(\mathbb R)$ 满足 $a_{ii}=a>0$，且对每个 $i$ 有
\[
\sum_{j=1}^n|a_{ij}|+\sum_{j=1}^n|a_{ji}|<4a.
\]
求实二次型 $f(x)=x^{\mathrm T}Ax$ 的规范形。

\par\hyperref[ans:pre10-math:2]{解析见第~\pageref{ans:pre10-math:2} 页。}\quad\href{https://github.com/lzxcr/cmc-mathematics/blob/64916fa1fdbc3e444a652eca204d5669058ee321/src/problems/preliminary/10/02.md}{来源条目 pre10-math:2}

\end{example}

\begin{example}[初赛第 11 届（2019 年度），来源序号4]\label{ex:ch10-11th-A:4}

已知 $f_{1}$ 为实 $n$ 元正定二次型。令
\[
V=\left\{f\mid f\text{ 为实 }n\text{ 元二次型，满足：对任何实数 }k\text{ 有 }kf+f_{1}\text{ 属于恒号二次型}\right\},
\]
这里恒号二次型为 0 二次型，正定二次型及负定二次型的总称。证明：$V$ 按照通常的二次型加法和数乘构成一个实向量空间，并求这个向量空间的维数。

\par\hyperref[ans:ch10-11th-A:4]{解析见第~\pageref{ans:ch10-11th-A:4} 页。}\quad\href{https://github.com/lzxcr/cmc-mathematics/blob/64916fa1fdbc3e444a652eca204d5669058ee321/src/problems/preliminary/11/04.md}{来源条目 ch10-11th-A:4}

\end{example}

\begin{example}[决赛第 1 届（2010 年），来源序号10]\label{ex:fin01-all:10}

设 $A,B$ 为 $n$ 阶半正定实对称矩阵，且 $n-1\le\operatorname{rank}A\le n$。证明存在实可逆矩阵 $C$，使 $C^{\mathsf T}AC$ 和 $C^{\mathsf T}BC$ 都是对角矩阵。

\par\hyperref[ans:fin01-all:10]{解析见第~\pageref{ans:fin01-all:10} 页。}\quad\href{https://github.com/lzxcr/cmc-mathematics/blob/64916fa1fdbc3e444a652eca204d5669058ee321/src/problems/finals/01/10.md}{来源条目 fin01-all:10}

\end{example}

\begin{example}[决赛第 2 届（2011 年），来源序号1]\label{ex:fin02-all:1}

求所有过原点且与椭球面 $4x^2+5y^2+6z^2=1$ 相交成圆周的平面。

\par\hyperref[ans:fin02-all:1]{解析见第~\pageref{ans:fin02-all:1} 页。}\quad\href{https://github.com/lzxcr/cmc-mathematics/blob/64916fa1fdbc3e444a652eca204d5669058ee321/src/problems/finals/02/01.md}{来源条目 fin02-all:1}

\end{example}

\begin{example}[决赛第 4 届（2013 年），来源序号5]\label{ex:fin04-all:5}

设 $A=(a_{ij})\in M_3(\mathbb R)$ 为对称矩阵，$A^*$ 表示其伴随矩阵，且
\[
\det A=-12,\qquad\operatorname{tr}A=1,\qquad
(A^*-4I_3)(1,0,-2)^{\mathsf T}=0.
\]
给出一个正交变换，将
\[
f(x_1,x_2,x_3,x_4)=
\det\begin{pmatrix}
x_1^2&x_2&x_3&x_4\\
-x_2&a_{11}&a_{12}&a_{13}\\
-x_3&a_{12}&a_{22}&a_{23}\\
-x_4&a_{13}&a_{23}&a_{33}
\end{pmatrix}
\]
化为标准型。

\par\hyperref[ans:fin04-all:5]{解析见第~\pageref{ans:fin04-all:5} 页。}\quad\href{https://github.com/lzxcr/cmc-mathematics/blob/64916fa1fdbc3e444a652eca204d5669058ee321/src/problems/finals/04/05.md}{来源条目 fin04-all:5}

\end{example}

\begin{example}[决赛第 6 届（2015 年），来源序号1，共同题]\label{ex:fin06-lo:1}

（填空题）实二次型 $2x_1x_2-x_1x_3+5x_2x_3$ 的规范型为 $\underline{\qquad\qquad}$。

\par\hyperref[ans:fin06-lo:1]{解析见第~\pageref{ans:fin06-lo:1} 页。}\quad\href{https://github.com/lzxcr/cmc-mathematics/blob/64916fa1fdbc3e444a652eca204d5669058ee321/src/problems/finals/06/01.md}{来源条目 fin06-lo:1}

\end{example}

\begin{example}[决赛第 9 届（2018 年），来源序号4，共同题]\label{ex:fin09-hi:4}

（填空题）设二次型 $f(x_{1},\dots,x_{n})=(x_{1},\dots,x_{n})A\begin{pmatrix}x_{1}\\x_{2}\\\vdots\\x_{n}\end{pmatrix}$ 的矩阵 $A$ 为
\[
A=\begin{pmatrix}1&a&a&\dots&a\\a&1&a&\dots&a\\\vdots&\vdots&\vdots&\vdots&\vdots\\a&\dots&a&1&a\\a&a&\dots&a&1\end{pmatrix},
\]
其中 $n>1$，$a \in \mathbb{R}$，则 $f$ 在正交变换下的标准形为 $\underline{\qquad\qquad}$。

\par\hyperref[ans:fin09-hi:4]{解析见第~\pageref{ans:fin09-hi:4} 页。}\quad\href{https://github.com/lzxcr/cmc-mathematics/blob/64916fa1fdbc3e444a652eca204d5669058ee321/src/problems/finals/09/04.md}{来源条目 fin09-hi:4}

\end{example}

\begin{example}[决赛第 10 届（2019 年），来源序号1，共同题]\label{ex:fin10-hi:1}

（填空题）设 $A$ 为实对称方阵，$(1,0,1)$ 和 $(1,2,0)$ 构成其行向量的一个极大无关组，则有 $A = $ $\underline{\qquad\qquad}$。

\par\hyperref[ans:fin10-hi:1]{解析见第~\pageref{ans:fin10-hi:1} 页。}\quad\href{https://github.com/lzxcr/cmc-mathematics/blob/64916fa1fdbc3e444a652eca204d5669058ee321/src/problems/finals/10/01.md}{来源条目 fin10-hi:1}

\end{example}

\begin{example}[决赛第 11 届（2021 年），来源序号3，共同题]\label{ex:fin11-high:3}

（填空题）在三维空间的直角坐标系中，方程 $2x^{2}+y^{2}+z^{2}+2xy-2xz=1$ 表示的二次曲面类型是 $\underline{\qquad\qquad}$。

\par\hyperref[ans:fin11-high:3]{解析见第~\pageref{ans:fin11-high:3} 页。}\quad\href{https://github.com/lzxcr/cmc-mathematics/blob/64916fa1fdbc3e444a652eca204d5669058ee321/src/problems/finals/11/03.md}{来源条目 fin11-high:3}

\end{example}

\begin{example}[决赛第 14 届（2023 年），来源序号5，共同题]\label{ex:fin14-high:5}

设二次曲线 $Ax^2+Bxy+Cy^2+Dx+Ey+F=0$ 经过正交变换 $\begin{pmatrix} x \\ y \end{pmatrix}=Q \begin{pmatrix} x' \\ y' \end{pmatrix}$ 得此曲线的另一表达式。

（i）求证 $B^2-4AC$ 与 $F$ 为上述正交变换的不变量；

（ii）给出在上述正交变换下，交叉项 $x'y'$ 系数为零的充要条件。

\par\hyperref[ans:fin14-high:5]{解析见第~\pageref{ans:fin14-high:5} 页。}\quad\href{https://github.com/lzxcr/cmc-mathematics/blob/64916fa1fdbc3e444a652eca204d5669058ee321/src/problems/finals/14/05.md}{来源条目 fin14-high:5}

\end{example}

\begin{example}[决赛第 15 届（2024 年），来源序号1，共同题]\label{ex:fin15-high:1}

将 24 阶实对称矩阵按合同关系进行分类，即两个 24 阶实对称矩阵属于同一类当且仅当它们相互合同，则共有 $\underline{\qquad}$ 类。

\par\hyperref[ans:fin15-high:1]{解析见第~\pageref{ans:fin15-high:1} 页。}\quad\href{https://github.com/lzxcr/cmc-mathematics/blob/64916fa1fdbc3e444a652eca204d5669058ee321/src/problems/finals/15/01.md}{来源条目 fin15-high:1}

\end{example}

\begin{example}[决赛第 16 届（2025 年），来源序号1，共同题]\label{ex:fin16-high:1}

设 2025 阶实对称矩阵 $A$ 的秩为 1012，且 $A^2 = A$，则二次型 $x^{\mathrm{T}}Ax$ 的规范型为 $\underline{\qquad}$。

\par\hyperref[ans:fin16-high:1]{解析见第~\pageref{ans:fin16-high:1} 页。}\quad\href{https://github.com/lzxcr/cmc-mathematics/blob/64916fa1fdbc3e444a652eca204d5669058ee321/src/problems/finals/16/01.md}{来源条目 fin16-high:1}

\end{example}

\begin{example}[决赛第 17 届（2026 年），来源序号6，共同题]\label{ex:fin17-high:6}

设 $A,B$ 皆为 $n(n\ge2)$ 阶实方阵，$E$ 为 $n$ 阶单位阵，行列式 $|A-B|=1$，且 $AA^T-E$ 与 $BB^T-E$ 均为正定矩阵。证明：

（1）$(E-AB^T)(E-BB^T)^{-1}(E-AB^T)^T-(E-AA^T)=(E-BB^T)^{-1}-E-AB^T(E-BB^T)^{-1}-(E-BB^T)^{-1}BA^T+A\left(E+B^T(E-BB^T)^{-1}B\right)A^T$；

（2）$|E-AA^T|\cdot|E-BB^T|<|E-AB^T|^2-1$。

\par\hyperref[ans:fin17-high:6]{解析见第~\pageref{ans:fin17-high:6} 页。}\quad\href{https://github.com/lzxcr/cmc-mathematics/blob/64916fa1fdbc3e444a652eca204d5669058ee321/src/problems/finals/17/06.md}{来源条目 fin17-high:6}

\end{example}

\section{内积、正交投影与奇异值}

\label{sec:questions-8}

内积与正交结构能把代数约束转化为长度、角度和能量的约束。非正交直和中的正交投影一般不是互补投影；奇异值应由伴随乘积确定，不能直接当作特征值的绝对值。

\begin{example}[初赛第 9 届（2017 年度），来源序号3]\label{ex:pre09-math:3}

设 $\Gamma=\{W_1,\ldots,W_r\}$ 是由 $r$ 个互不相同的可逆 $n$ 阶复矩阵组成的集合，且关于乘法封闭。证明
\[
\sum_{i=1}^rW_i=0\quad\Longleftrightarrow\quad
\sum_{i=1}^r\operatorname{tr}W_i=0.
\]

\par\hyperref[ans:pre09-math:3]{解析见第~\pageref{ans:pre09-math:3} 页。}\quad\href{https://github.com/lzxcr/cmc-mathematics/blob/64916fa1fdbc3e444a652eca204d5669058ee321/src/problems/preliminary/09/03.md}{来源条目 pre09-math:3}

\end{example}

\begin{example}[初赛第 12 届（2020 年度），来源序号3]\label{ex:ch09-12th-A:3}

设 $A,B$ 均为 2020 阶正交矩阵，齐次线性方程组 $Ax=Bx\left(x\in\mathbb{R}^{2020}\right)$ 的解空间维数为 3。问：矩阵 $A,B$ 是否可能相似？证明你的结论。

\par\hyperref[ans:ch09-12th-A:3]{解析见第~\pageref{ans:ch09-12th-A:3} 页。}\quad\href{https://github.com/lzxcr/cmc-mathematics/blob/64916fa1fdbc3e444a652eca204d5669058ee321/src/problems/preliminary/12/03.md}{来源条目 ch09-12th-A:3}

\end{example}

\begin{example}[初赛第 13 届（2021 年度），来源序号4]\label{ex:ch07-13th-A:4}

设 $P$ 为对称酉矩阵，证明：存在可逆复矩阵 $Q$ 使得 $P=\bar{Q}Q^{-1}$。

\par\hyperref[ans:ch07-13th-A:4]{解析见第~\pageref{ans:ch07-13th-A:4} 页。}\quad\href{https://github.com/lzxcr/cmc-mathematics/blob/64916fa1fdbc3e444a652eca204d5669058ee321/src/problems/preliminary/13/04.md}{来源条目 ch07-13th-A:4}

\end{example}

\begin{example}[决赛第 2 届（2011 年），来源序号6]\label{ex:fin02-all:6}

设非零线性泛函 $\varphi:M_n(\mathbb R)\to\mathbb R$ 满足
$\varphi(XY)=\varphi(YX)$。定义双线性型
$\langle X,Y\rangle=\varphi(XY)$。

（1）证明该双线性型非退化。
（2）若 $A_1,\ldots,A_{n^2}$ 与 $B_1,\ldots,B_{n^2}$ 是一对对偶基，即 $\langle A_i,B_j\rangle=\delta_{ij}$，证明 $\sum_i A_iB_i$ 为数量矩阵。

\par\hyperref[ans:fin02-all:6]{解析见第~\pageref{ans:fin02-all:6} 页。}\quad\href{https://github.com/lzxcr/cmc-mathematics/blob/64916fa1fdbc3e444a652eca204d5669058ee321/src/problems/finals/02/06.md}{来源条目 fin02-all:6}

\end{example}

\begin{example}[决赛第 5 届（2014 年），来源序号12，高年级组]\label{ex:fin05-hi:11}

定义多项式列
\[
T_0(x)=1,\qquad T_1(x)=x,\qquad
T_{n+1}(x)=2xT_n(x)-T_{n-1}(x)\quad(n\ge1).
\]

（1）证明 $T_n(x)=\cos(n\arccos x)$（$-1\le x\le1$）。
（2）证明 $(T_n)$ 在内积
\[
\langle f,g\rangle=\int_{-1}^1\frac{f(x)g(x)}{\sqrt{1-x^2}}\,dx
\]
下两两正交。
（3）设 $n\ge1$，$P$ 为 $n$ 次首一实多项式。证明
\[
\|P\|_\infty\ge2^{1-n},\qquad
\|P\|_\infty:=\max_{[-1,1]}|P|,
\]
且等号当且仅当 $P=2^{1-n}T_n$ 时成立。

\par\hyperref[ans:fin05-hi:11]{解析见第~\pageref{ans:fin05-hi:11} 页。}\quad\href{https://github.com/lzxcr/cmc-mathematics/blob/64916fa1fdbc3e444a652eca204d5669058ee321/src/problems/finals/05/12.md}{来源条目 fin05-hi:11}

\end{example}

\begin{example}[决赛第 6 届（2015 年），来源序号6，共同题]\label{ex:fin06-lo:6}

设
\[
\Gamma=\left\{\begin{pmatrix}z_1&z_2\\-\overline{z_2}&\overline{z_1}\end{pmatrix}:\ z_1,z_2\in\mathbb{C}\right\},
\]
其中 $\mathbb{C}$ 表示复数域。试证明：$\forall A\in\Gamma$，$A$ 的 Jordan 标准型 $J_A$ 仍属于 $\Gamma$；进一步还存在可逆的矩阵 $P\in\Gamma$ 使得
\[
P^{-1}AP=J_A .
\]

\par\hyperref[ans:fin06-lo:6]{解析见第~\pageref{ans:fin06-lo:6} 页。}\quad\href{https://github.com/lzxcr/cmc-mathematics/blob/64916fa1fdbc3e444a652eca204d5669058ee321/src/problems/finals/06/06.md}{来源条目 fin06-lo:6}

\end{example}

\begin{example}[决赛第 7 届（2016 年），来源序号6，共同题]\label{ex:fin07-hi:6}

设 $A,B$ 为 $n$ 阶实对称矩阵。证明
\[
\operatorname{tr}\left(\left(AB\right)^{2}\right)\leq\operatorname{tr}\left(A^{2}B^{2}\right).
\]

\par\hyperref[ans:fin07-hi:6]{解析见第~\pageref{ans:fin07-hi:6} 页。}\quad\href{https://github.com/lzxcr/cmc-mathematics/blob/64916fa1fdbc3e444a652eca204d5669058ee321/src/problems/finals/07/06.md}{来源条目 fin07-hi:6}

\end{example}

\begin{example}[决赛第 8 届（2017 年），来源序号4，共同题]\label{ex:fin08-hi:4}

（填空题）记两个特征值为 1，2 的 2 阶实对称矩阵的全体为 $\Gamma$。$\forall A\in\Gamma$，$a_{21}$ 表示 $A$ 的 $(2,1)$ 位置元素，则集合 $\bigcup_{A\in\Gamma}\left\{a_{21}\right\}$ 的最小元等于 $\underline{\qquad\qquad}$。

\par\hyperref[ans:fin08-hi:4]{解析见第~\pageref{ans:fin08-hi:4} 页。}\quad\href{https://github.com/lzxcr/cmc-mathematics/blob/64916fa1fdbc3e444a652eca204d5669058ee321/src/problems/finals/08/04.md}{来源条目 fin08-hi:4}

\end{example}

\begin{example}[决赛第 10 届（2019 年），来源序号4，共同题]\label{ex:fin10-hi:4}

（填空题）设 $U$ 为 8 阶实正交方阵，$U$ 中元素皆为 $\frac{1}{2\sqrt{2}}$ 的 $3 \times 3$ 子矩阵的个数记为 $t$，则 $t$ 最多为 $\underline{\qquad\qquad}$。

\par\hyperref[ans:fin10-hi:4]{解析见第~\pageref{ans:fin10-hi:4} 页。}\quad\href{https://github.com/lzxcr/cmc-mathematics/blob/64916fa1fdbc3e444a652eca204d5669058ee321/src/problems/finals/10/04.md}{来源条目 fin10-hi:4}

\end{example}

\begin{example}[决赛第 11 届（2021 年），来源序号4，共同题]\label{ex:fin11-high:4}

（填空题）在矩阵 $A = \begin{bmatrix} 1 & -2 & 0 \\ 1 & 0 & 1 \\ 0 & 2 & 1 \end{bmatrix}$ 的奇异值分解 $A = U\Lambda V$ 中（其中 $U,V$ 为正交方阵，$\Lambda$ 为对角阵），$\Lambda = $ $\underline{\qquad\qquad}$。

\par\hyperref[ans:fin11-high:4]{解析见第~\pageref{ans:fin11-high:4} 页。}\quad\href{https://github.com/lzxcr/cmc-mathematics/blob/64916fa1fdbc3e444a652eca204d5669058ee321/src/problems/finals/11/04.md}{来源条目 fin11-high:4}

\end{example}

\begin{example}[决赛第 11 届（2021 年），来源序号6，共同题]\label{ex:fin11-high:6}

设 $V$ 是有限维欧氏空间，$V_{1}, V_{2}$ 是 $V$ 的非平凡子空间且 $V = V_{1} \oplus V_{2}$。设 $p_{1}, p_{2}$ 分别是 $V$ 到 $V_{1}, V_{2}$ 的正交投影，$\varphi = p_{1} + p_{2}$，用 $\det \varphi$ 表示线性变换 $\varphi$ 的行列式。证明：$0 < \det\varphi \leq 1$ 且 $\det\varphi = 1$ 的充要条件是 $V_{1}$ 与 $V_{2}$ 正交。

\par\hyperref[ans:fin11-high:6]{解析见第~\pageref{ans:fin11-high:6} 页。}\quad\href{https://github.com/lzxcr/cmc-mathematics/blob/64916fa1fdbc3e444a652eca204d5669058ee321/src/problems/finals/11/06.md}{来源条目 fin11-high:6}

\end{example}

\begin{example}[决赛第 13 届（2023 年），来源序号8，低年级组]\label{ex:fin13-low:8}

设矩阵 $A = \begin{pmatrix} a & b \\ c & d \end{pmatrix} \in \mathbb{R}^{2 \times 2}$，二元函数 $p(X, Y)$ 被定义为：

\[
p(X, Y) : \mathbb{R}^{2 \times 2} \times \mathbb{R}^{2 \times 2} \rightarrow \mathbb{R}, \quad (X, Y) \mapsto \operatorname{tr}\left(X A Y^{T}\right),
\]

即 $p(X, Y) = \operatorname{tr}\left(X A Y^{T}\right)$，其中 $X, Y \in \mathbb{R}^{2 \times 2}$ 为变量，$\operatorname{tr}$ 表示矩阵的迹。

（1）证明：$p(X, Y)$ 成为实向量空间 $\mathbb{R}^{2 \times 2}$ 上的内积的充要条件是

\[
\begin{cases} a > 0, \\ b = c, \\ ad > b^{2}. \end{cases}
\]

（2）进一步，若 $A$ 还使得变换

\[
\sigma : \mathbb{R}^{2 \times 2} \mapsto \mathbb{R}^{2 \times 2}, \quad X \rightarrow \begin{pmatrix} \frac{\sqrt{3}}{2} & \frac{1}{2} \\ -\frac{1}{2} & \frac{\sqrt{3}}{2} \end{pmatrix} X \begin{pmatrix} 0 & 1 \\ -1 & 0 \end{pmatrix}
\]

成为 （1）中内积空间 $\left(\mathbb{R}^{2 \times 2}, p(X, Y)\right)$ 中的一个正交变换。求 $A$。

\par\hyperref[ans:fin13-low:8]{解析见第~\pageref{ans:fin13-low:8} 页。}\quad\href{https://github.com/lzxcr/cmc-mathematics/blob/64916fa1fdbc3e444a652eca204d5669058ee321/src/problems/finals/13/08.md}{来源条目 fin13-low:8}

\end{example}

\section{矩阵方程、分块与数域问题}

\label{sec:questions-9}

矩阵方程可以视为更大向量空间上的线性方程；分块计算须保持乘法次序。跨数域的相似问题除求解线性条件外，还须保证最终换基矩阵可逆。

\begin{example}[初赛第 4 届（2012 年度），来源序号7]\label{ex:pre04-math:7}

设
\[
A=\begin{pmatrix}2&2\\2&a\end{pmatrix},\qquad
B=\begin{pmatrix}4&b\\3&1\end{pmatrix},
\qquad a,b\in\mathbb R.
\]
证明：

（1） $AX=B$ 有解而 $BY=A$ 无解，当且仅当 $a\ne2$、$b=4/3$。

（2） $A,B$ 相似，当且仅当 $a=3$、$b=2/3$。

（3） $A,B$ 在实数域上合同，当且仅当 $a<2$、$b=3$。

\par\hyperref[ans:pre04-math:7]{解析见第~\pageref{ans:pre04-math:7} 页。}\quad\href{https://github.com/lzxcr/cmc-mathematics/blob/64916fa1fdbc3e444a652eca204d5669058ee321/src/problems/preliminary/04/07.md}{来源条目 pre04-math:7}

\end{example}

\begin{example}[初赛第 5 届（2013 年度），来源序号6]\label{ex:pre05-math:6}

记 $E_{ij}$ 为第 $(i,j)$ 个元素为 $1$、其余为 $0$ 的 $n$ 阶实矩阵，$\Gamma_r$ 为秩等于 $r$ 的全部 $n$ 阶实矩阵。设映射 $\phi:M_n(\mathbb R)\to M_n(\mathbb R)$ 满足 $\phi(AB)=\phi(A)\phi(B)$。

（1）证明同属 $\Gamma_r$ 的 $A,B$ 满足 $\operatorname{rank}\phi(A)=\operatorname{rank}\phi(B)$。

（2）若 $\phi(0)=0$，且有秩为 $1$ 的 $W$ 满足 $\phi(W)\ne0$，证明存在可逆矩阵 $R$，使所有 $i,j$ 均满足 $\phi(E_{ij})=RE_{ij}R^{-1}$。

\par\hyperref[ans:pre05-math:6]{解析见第~\pageref{ans:pre05-math:6} 页。}\quad\href{https://github.com/lzxcr/cmc-mathematics/blob/64916fa1fdbc3e444a652eca204d5669058ee321/src/problems/preliminary/05/06.md}{来源条目 pre05-math:6}

\end{example}

\begin{example}[初赛第 13 届补赛（2021 年度），来源序号3]\label{ex:ch08-13th-A-makeup:3}

设 $R=\{0,1,4,9,16\}$，$U$ 为 $R$ 上的 3 阶方阵全体：
\[
U=\left\{\left(a_{ij}\right)_{3\times 3}\mid a_{ij}\in R\right\}.
\]
例如，
\[
\begin{pmatrix}1&1&0\\ 0&1&0\\ 0&0&0\end{pmatrix},\qquad \begin{pmatrix}1&4&0\\ 0&1&0\\ 0&0&0\end{pmatrix}
\]
便为 $U$ 中的两元。现设 $S$ 为 $U$ 的一个子集。证明：若 $S$ 的元素个数多于 $5^{9}-5^{3}-18$，则必存在不同的 $A,B\in S$ 使得 $AB=BA$。

\par\hyperref[ans:ch08-13th-A-makeup:3]{解析见第~\pageref{ans:ch08-13th-A-makeup:3} 页。}\quad\href{https://github.com/lzxcr/cmc-mathematics/blob/64916fa1fdbc3e444a652eca204d5669058ee321/src/problems/preliminary/13-makeup/03.md}{来源条目 ch08-13th-A-makeup:3}

\end{example}

\begin{example}[决赛第 3 届（2012 年），来源序号5]\label{ex:fin03-all:5}

设 $A\in M_{3\times2}(\mathbb R)$、$B\in M_{2\times3}(\mathbb R)$，且
\[
AB=\begin{pmatrix}8&0&-4\\-\frac32&9&-6\\-2&0&1\end{pmatrix}.
\]
求 $BA$。

\par\hyperref[ans:fin03-all:5]{解析见第~\pageref{ans:fin03-all:5} 页。}\quad\href{https://github.com/lzxcr/cmc-mathematics/blob/64916fa1fdbc3e444a652eca204d5669058ee321/src/problems/finals/03/05.md}{来源条目 fin03-all:5}

\end{example}

\begin{example}[决赛第 3 届（2012 年），来源序号6]\label{ex:fin03-all:6}

设 $\{A_i\}_{i\in I}$、$\{B_i\}_{i\in I}$ 是同阶有理矩阵的两个族。若存在实可逆矩阵 $P$，使
\[
P^{-1}A_iP=B_i\qquad(i\in I),
\]
证明也存在满足上述等式的有理可逆矩阵。

\par\hyperref[ans:fin03-all:6]{解析见第~\pageref{ans:fin03-all:6} 页。}\quad\href{https://github.com/lzxcr/cmc-mathematics/blob/64916fa1fdbc3e444a652eca204d5669058ee321/src/problems/finals/03/06.md}{来源条目 fin03-all:6}

\end{example}

\begin{example}[决赛第 8 届（2017 年），来源序号6，共同题]\label{ex:fin08-hi:6}

证明题：设 $n$ 阶方阵 $A,B$ 满足：秩 $\left(ABA\right)=$ 秩 $\left(B\right)$。证明：$AB$ 与 $BA$ 相似。

\par\hyperref[ans:fin08-hi:6]{解析见第~\pageref{ans:fin08-hi:6} 页。}\quad\href{https://github.com/lzxcr/cmc-mathematics/blob/64916fa1fdbc3e444a652eca204d5669058ee321/src/problems/finals/08/06.md}{来源条目 fin08-hi:6}

\end{example}

\section{矩阵函数与数值线性代数}

\label{sec:questions-10}

幂零矩阵把解析函数变成有限多项式，谱半径则控制线性迭代收敛。算法题同时说明输入、初值、更新和停止条件；有限步终止或收敛速度的结论应区分精确运算与浮点运算。

\begin{example}[初赛第 6 届（2014 年度），来源序号5]\label{ex:pre06-math:5}

给定正整数 $m$。证明对任意正整数 $n,\ell$，存在 $m$ 阶实矩阵 $X$，使
\[
X^n+X^\ell=I+T,\qquad
T_{ij}=\begin{cases}i-j+1,&i\ge j,\\0,&i<j.\end{cases}
\]

\par\hyperref[ans:pre06-math:5]{解析见第~\pageref{ans:pre06-math:5} 页。}\quad\href{https://github.com/lzxcr/cmc-mathematics/blob/64916fa1fdbc3e444a652eca204d5669058ee321/src/problems/preliminary/06/05.md}{来源条目 pre06-math:5}

\end{example}

\begin{example}[初赛第 7 届（2015 年度），来源序号3]\label{ex:pre07-math:3}

设 $A$ 为 $n$ 阶矩阵，全部特征值均为偶整数。证明矩阵方程
\[
X+AX-XA^2=0
\]
只有零解。

\par\hyperref[ans:pre07-math:3]{解析见第~\pageref{ans:pre07-math:3} 页。}\quad\href{https://github.com/lzxcr/cmc-mathematics/blob/64916fa1fdbc3e444a652eca204d5669058ee321/src/problems/preliminary/07/03.md}{来源条目 pre07-math:3}

\end{example}

\begin{example}[决赛第 6 届（2015 年），来源序号13，高年级组]\label{ex:fin06-hi:12}

考虑求解线性方程组 $A\mathbf x=\mathbf b$ 的迭代格式
\[
(\alpha D-C)\mathbf x^{(k+1)}=\bigl((\alpha-1)D+C^{T}\bigr)\mathbf x^{(k)}+\mathbf b,
\]
其中 $D$ 为实对称正定方阵，$C$ 是满足 $C+C^{T}=D-A$ 的实方阵，$\alpha$ 为实数。设 $A$ 是实对称正定方阵，$\alpha D-C$ 可逆，且 $\alpha>\frac12$。证明：该迭代格式对任何初始向量 $\mathbf x^{(0)}$ 都收敛。

\par\hyperref[ans:fin06-hi:12]{解析见第~\pageref{ans:fin06-hi:12} 页。}\quad\href{https://github.com/lzxcr/cmc-mathematics/blob/64916fa1fdbc3e444a652eca204d5669058ee321/src/problems/finals/06/13.md}{来源条目 fin06-hi:12}

\end{example}

\begin{example}[决赛第 7 届（2016 年），来源序号4，共同题]\label{ex:fin07-hi:4}

（填空题）若实向量 $X=\left(a,b,c\right)$ 的三个分量 $a,b,c$ 满足
\[
\begin{pmatrix} a & b \\ 0 & c \end{pmatrix}^{2016}=I_{2},
\]
则 $X=$ $\underline{\qquad\qquad}$ 或 $\underline{\qquad\qquad}$ 或 $\underline{\qquad\qquad}$ 或 $\underline{\qquad\qquad}$。

\par\hyperref[ans:fin07-hi:4]{解析见第~\pageref{ans:fin07-hi:4} 页。}\quad\href{https://github.com/lzxcr/cmc-mathematics/blob/64916fa1fdbc3e444a652eca204d5669058ee321/src/problems/finals/07/04.md}{来源条目 fin07-hi:4}

\end{example}

\begin{example}[决赛第 9 届（2018 年），来源序号8，低年级组]\label{ex:fin09-lo:8}

设 $\alpha \in (1, 2)$，$(1 - x)^{\alpha}$ 的麦克劳林级数为 $\sum_{k=0}^{\infty} a_k x^k$，$n \times n$ 实常数矩阵 $A$ 为幂零矩阵，$I$ 为单位阵。设矩阵值函数 $G(x)$ 定义为
\[
G(x) \equiv \left( g_{ij}(x) \right) := \sum_{k=0}^{\infty} a_k (xI + A)^k, \quad 0 \leq x < 1.
\]
试证：对于 $1 \leq i, j \leq n$，积分 $\int_0^1 g_{ij}(x) \mathrm{d}x$ 均存在的充分必要条件是 $A^3 = 0$。

\par\hyperref[ans:fin09-lo:8]{解析见第~\pageref{ans:fin09-lo:8} 页。}\quad\href{https://github.com/lzxcr/cmc-mathematics/blob/64916fa1fdbc3e444a652eca204d5669058ee321/src/problems/finals/09/08.md}{来源条目 fin09-lo:8}

\end{example}

\begin{example}[决赛第 9 届（2018 年），来源序号13，高年级组]\label{ex:fin09-hi:11}

（数值分析）设线性方程组 $Ax=b$ 的系数矩阵 $A$ 实对称正定。推导共轭梯度法的计算格式，并证明该格式经有限步迭代后收敛。

\par\hyperref[ans:fin09-hi:11]{解析见第~\pageref{ans:fin09-hi:11} 页。}\quad\href{https://github.com/lzxcr/cmc-mathematics/blob/64916fa1fdbc3e444a652eca204d5669058ee321/src/problems/finals/09/13.md}{来源条目 fin09-hi:11}

\end{example}

\begin{example}[决赛第 11 届（2021 年），来源序号8，低年级组]\label{ex:fin11-low:8}

设 $n\geq 2$，对于 $A\in\mathbb{R}^{n\times n}$，规定 $A^{0}$ 为 $n$ 阶单位阵 $I$，形式定义
\[
\sin A=\sum_{k=0}^{\infty}\frac{(-1)^{k}}{(2k+1)!}A^{2k+1}, \quad \cos A=\sum_{k=0}^{\infty}\frac{(-1)^{k}}{(2k)!}A^{2k}, \quad \arctan A=\sum_{k=0}^{\infty}\frac{(-1)^{k}}{2k+1}A^{2k+1}.
\]
记 $\|A\| \equiv \max_{\substack{x\in\mathbb{R}^{n}\\ \|x\|=1}} \|Ax\|$，其中 $\|x\| \equiv \sqrt{x^{T}x}$，证明：

（1）$\forall A\in\mathbb{R}^{n\times n}, \sin A,\cos A$ 均有意义，且 $(\sin A)^{2}+(\cos A)^{2}=I$。

（2）当 $\|A\|<1$ 时，$\arctan A$ 有意义，且 $\sin \arctan A=A\cos \arctan A$。

\par\hyperref[ans:fin11-low:8]{解析见第~\pageref{ans:fin11-low:8} 页。}\quad\href{https://github.com/lzxcr/cmc-mathematics/blob/64916fa1fdbc3e444a652eca204d5669058ee321/src/problems/finals/11/08.md}{来源条目 fin11-low:8}

\end{example}

\begin{example}[决赛第 11 届（2021 年），来源序号15，高年级组]\label{ex:fin11-high:13}

设 $A$ 是具有正对角元的非奇异对称矩阵。证明：若求解线性方程组 $Ax=b$ 的 Gauss-Seidel 迭代法对任意初始值都收敛，则 $A$ 为正定矩阵。

\par\hyperref[ans:fin11-high:13]{解析见第~\pageref{ans:fin11-high:13} 页。}\quad\href{https://github.com/lzxcr/cmc-mathematics/blob/64916fa1fdbc3e444a652eca204d5669058ee321/src/problems/finals/11/15.md}{来源条目 fin11-high:13}

\end{example}

\begin{example}[决赛第 13 届（2023 年），来源序号16，高年级组]\label{ex:fin13-high:13}

给出 QR 方法求矩阵特征值的算法描述，说明收敛所需的条件，并在一组适当的充分条件下证明其收敛性。

\par\hyperref[ans:fin13-high:13]{解析见第~\pageref{ans:fin13-high:13} 页。}\quad\href{https://github.com/lzxcr/cmc-mathematics/blob/64916fa1fdbc3e444a652eca204d5669058ee321/src/problems/finals/13/16.md}{来源条目 fin13-high:13}

\end{example}

\begin{example}[决赛第 14 届（2023 年），来源序号6，共同题]\label{ex:fin14-high:6}

设实系数多项式 $f(t)=\sum_{j=0}^ma_jt^j$ 满足 $f(1)=a\ne0$、$f'(1)=b\ne0$。证明：对任意整数 $n\ge2$ 及 $\varepsilon\in\mathbb C\setminus\{0\}$，存在 $n$ 阶复矩阵 $C$，使 $f(C)=aI_n+\varepsilon N$，其中 $N_{i,i+1}=1$（$1\le i<n$），其余元素均为零。

\par\hyperref[ans:fin14-high:6]{解析见第~\pageref{ans:fin14-high:6} 页。}\quad\href{https://github.com/lzxcr/cmc-mathematics/blob/64916fa1fdbc3e444a652eca204d5669058ee321/src/problems/finals/14/06.md}{来源条目 fin14-high:6}

\end{example}

\begin{example}[决赛第 14 届（2023 年），来源序号16，高年级组]\label{ex:fin14-high:13}

设 $A$ 是 $n$ 阶实方阵。

（i）给出矩阵 $A$ 的 $LU$ 分解的算法描述；

（ii）分析 $LU$ 分解算法的运算量（复杂度）；

（iii）证明：$A$ 具有唯一的 $LU$ 分解当且仅当 $A$ 的前 $n-1$ 阶顺序主子式均不为零。

\par\hyperref[ans:fin14-high:13]{解析见第~\pageref{ans:fin14-high:13} 页。}\quad\href{https://github.com/lzxcr/cmc-mathematics/blob/64916fa1fdbc3e444a652eca204d5669058ee321/src/problems/finals/14/16.md}{来源条目 fin14-high:13}

\end{example}

\begin{example}[决赛第 17 届（2026 年），来源序号1，共同题]\label{ex:fin17-high:1}

设矩阵

\[
A=
\begin{pmatrix}
\sin\frac{\pi}{1013} & -\cos\frac{\pi}{1013}\\
\cos\frac{\pi}{1013} & \sin\frac{\pi}{1013}
\end{pmatrix}^{2026},
\]

则 $A$ 的迹 $\operatorname{tr}A=\underline{\hspace{2cm}}$。

\par\hyperref[ans:fin17-high:1]{解析见第~\pageref{ans:fin17-high:1} 页。}\quad\href{https://github.com/lzxcr/cmc-mathematics/blob/64916fa1fdbc3e444a652eca204d5669058ee321/src/problems/finals/17/01.md}{来源条目 fin17-high:1}

\end{example}

\begin{example}[决赛第 17 届（2026 年），来源序号16，高年级组]\label{ex:fin17-high:13}

设矩阵 $A_n\in\mathbb R^{n\times n}$ 为三对角矩阵

\[
A_n=
\begin{pmatrix}
2 & -1 & & & \\
-1 & 2 & -1 & & \\
& \ddots & \ddots & \ddots & \\
& & -1 & 2 & -1\\
& & & -1 & 2
\end{pmatrix}.
\]

$A_n$ 的谱条件数为 $\kappa_2(A_n)=\frac{\lambda_n}{\lambda_1}$，其中 $\lambda_1,\lambda_n$ 分别为 $A_n$ 的最小与最大特征值。

（1）证明：$A_n$ 的特征值为 $\lambda_k=4\sin^2\frac{k\pi}{2(n+1)}$（$k=1,2,\dots,n$），并证明当 $n\to+\infty$ 时 $\kappa_2(A_n)\sim\frac{4n^2}{\pi^2}$。

（2）考虑用 Richardson 迭代求解线性方程组 $A_n\mathbf x=\mathbf b$：

\[
\mathbf x^{(m+1)}=\mathbf x^{(m)}-\alpha\left(A_n\mathbf x^{(m)}-\mathbf b\right),\qquad m=0,1,\cdots,\ \alpha>0 .
\]

求使迭代收敛的 $\alpha$ 取值范围，并求最优（即使收敛最快）的参数 $\alpha$。

（3）考虑用 Jacobi 迭代求解 $A_n\mathbf x=\mathbf b$，证明迭代收敛，并分析当 $n$ 充分大时收敛的快慢。

\par\hyperref[ans:fin17-high:13]{解析见第~\pageref{ans:fin17-high:13} 页。}\quad\href{https://github.com/lzxcr/cmc-mathematics/blob/64916fa1fdbc3e444a652eca204d5669058ee321/src/problems/finals/17/16.md}{来源条目 fin17-high:13}

\end{example}

\part{历年真题解析编}

\section{线性空间、线性相关与秩}

\label{sec:solutions-1}

把矩阵看成线性映射，先确定核、像和维数。秩不等式的实质往往是某个限制映射的秩—零度公式；函数或多项式构成的向量组也应回到有限维张成空间研究。

\subsection{初赛第 6 届（2014 年度），来源序号3}\label{ans:pre06-math:3}\label{pr:pre06-math:3}

对应例~\ref{ex:pre06-math:3}（第~\pageref{ex:pre06-math:3} 页）；\href{https://github.com/lzxcr/cmc-mathematics/blob/64916fa1fdbc3e444a652eca204d5669058ee321/src/problems/preliminary/06/03.md}{查看来源}。

\begin{solve}

令 $v(x)=(f_1(x),\ldots,f_n(x))^{\mathrm T}$。若这些向量不张成 $\mathbb R^n$，则有非零向量 $c$ 与它们全部正交，即 $\sum_ic_if_i(x)=0$ 对所有 $x$ 成立。因此函数线性无关时，$\{v(x):x\in[0,1]\}$ 张成 $\mathbb R^n$，可从中选出一组基 $v(a_1),\ldots,v(a_n)$，对应行列式非零。

反之，若该矩阵可逆，将任何恒等关系 $\sum_ic_if_i=0$ 在 $a_j$ 处取值，便得 $[f_i(a_j)]^{\mathrm T}c=0$，所以 $c=0$。

\end{solve}

\begin{remark}[方法与条件]

这是有限维对偶性的一个具体形式：点值泛函能够分离不同函数，便能从全部点值中选出有限个来识别给定的有限维函数空间。无需连续性或可微性。

\end{remark}

\subsection{初赛第 8 届（2016 年度），来源序号3}\label{ans:pre08-math:3}\label{pr:pre08-math:3}

对应例~\ref{ex:pre08-math:3}（第~\pageref{ex:pre08-math:3} 页）；\href{https://github.com/lzxcr/cmc-mathematics/blob/64916fa1fdbc3e444a652eca204d5669058ee321/src/problems/preliminary/08/03.md}{查看来源}。

\begin{solve}

固定 $0\ne v\in\mathbb R^{2016}$。$2017$ 个向量 $A_1v,\ldots,A_{2017}v$ 必线性相关，故存在不全为零的实数 $x_i$，使
\[
\left(\sum_{i=1}^{2017}x_iA_i\right)v=0.
\]
此矩阵有非零核，行列式为零。

\end{solve}

\begin{remark}[方法与条件]

结论适用于任意域：$n$ 阶矩阵组成的线性空间若每个非零元都可逆，其维数至多为 $n$。固定向量上的求值映射把它单射到 $n$ 维向量空间，无须分析行列式这个高次多项式。

\end{remark}

\subsection{初赛第 14 届第二次补赛（2022 年度），来源序号4}\label{ans:ch06-14th-A-makeup2:4}\label{pr:ch06-14th-A-makeup2:4}

对应例~\ref{ex:ch06-14th-A-makeup2:4}（第~\pageref{ex:ch06-14th-A-makeup2:4} 页）；\href{https://github.com/lzxcr/cmc-mathematics/blob/64916fa1fdbc3e444a652eca204d5669058ee321/src/problems/preliminary/14-makeup2/04.md}{查看来源}。

\begin{solve}

可让 $C$ 的像直接落进 $\ker A$，这样比要求 $ACB=0$ 更强的 $AC=0$ 就自动成立。允许使用 $I$ 中 $r+1$ 个坐标，恰好保证存在这样的非零像向量。记 $E_I=\operatorname{span}\{e_i:i\in I\}$，则 $\dim E_I=r+1>\operatorname{rank}A$。因此可取非零 $u\in E_I$ 使 $Au=0$。任取 $j\in J$，令
\[
C=ue_j^{\mathrm T}.
\]
它非零、仅在 $I\times J$ 内可能有非零元，故 $C\in V$；并且 $AC=0$，从而 $ACB=0$。

\end{solve}

\begin{remark}[方法与条件]

这个构造说明题设有冗余：只需 $|I|>\operatorname{rank}A$ 且 $J$ 非空，就能选到秩为 $1$ 的 $C$ 使 $AC=0$，完全不必使用 $B$ 的秩条件。在整个矩阵空间上，映射 $C\mapsto ACB$ 的秩为 $\operatorname{rank}A\operatorname{rank}B$；若限于本题的 $V$，则应改为 $\operatorname{rank}(A_{:,I})\operatorname{rank}(B_{J,:})$，不能混同这两个定义域。

像、核及秩一算子的观点，见\hyperref[app:linear-invariant]{方法卷：不变子空间与像核}。

\end{remark}

\subsection{初赛第 16 届（2024 年度），来源序号4}\label{ans:ch02-16th-A:4}\label{pr:ch02-16th-A:4}

对应例~\ref{ex:ch02-16th-A:4}（第~\pageref{ex:ch02-16th-A:4} 页）；\href{https://github.com/lzxcr/cmc-mathematics/blob/64916fa1fdbc3e444a652eca204d5669058ee321/src/problems/preliminary/16/04.md}{查看来源}。

\begin{solve}

低次多项式所在空间的维数有限，因此无关向量不能过多地挤在低次数层中。若某个 $f_i=0$，结论显然。否则将次数重新排列为 $n_1\leq\cdots\leq n_{2024}$。若多项式组线性无关，则前 $k$ 个多项式都属于维数为 $n_k+1$ 的空间 $\mathbb R[x]_{\leq n_k}$，故
\[
k\leq n_k+1\qquad(1\leq k\leq2024).
\]
相加得到
\[
\sum_{k=1}^{2024}n_k\geq\sum_{k=1}^{2024}(k-1)
=\binom{2024}{2}=2047276,
\]
与题设矛盾。

\end{solve}

\begin{remark}[方法与条件]

一般地，$N$ 个线性无关的一元多项式，按次数递增排列后必有 $\deg f_k\geq k-1$，因此次数和至少为 $\binom N2$。这来自滤过空间 $\mathbb R[x]_{\leq d}$ 的维数 $d+1$；多元多项式也可用同一维数思想，但维数函数改为 $\binom{d+r}{r}$，不再沿用一元的次数下界。

这种按嵌套子空间容量计数的观点与\hyperref[app:linear-invariant]{方法卷：不变子空间}中的维数方法相通。

\end{remark}

\subsection{决赛第 1 届（2010 年），来源序号11}\label{ans:fin01-all:11}\label{pr:fin01-all:11}

对应例~\ref{ex:fin01-all:11}（第~\pageref{ex:fin01-all:11} 页）；\href{https://github.com/lzxcr/cmc-mathematics/blob/64916fa1fdbc3e444a652eca204d5669058ee321/src/problems/finals/01/11.md}{查看来源}。

\begin{solve}

映射 $T:V\to\mathbb C^2$，$T(v)=(f_1(v),f_2(v))$ 的秩为 $2$，故可取 $v$ 满足 $f_1(v)=0$、$f_2(v)=1$。令
\[
\alpha_1=f_2(\alpha)v,\qquad \alpha_2=\alpha-\alpha_1,
\]
便有 $\alpha_1\in\ker f_1$、$\alpha_2\in\ker f_2$，从而满足要求。

\end{solve}

\begin{remark}[方法与条件]

题目的实质是 $V=\ker f_1+\ker f_2$。这个和一般并非直和：分解的自由度正是 $\ker f_1\cap\ker f_2$。

\end{remark}

\subsection{决赛第 12 届（2021 年），来源序号6，共同题}\label{ans:fin12-high:6}\label{pr:fin12-high:6}

对应例~\ref{ex:fin12-high:6}（第~\pageref{ex:fin12-high:6} 页）；\href{https://github.com/lzxcr/cmc-mathematics/blob/64916fa1fdbc3e444a652eca204d5669058ee321/src/problems/finals/12/06.md}{查看来源}。

\begin{solve}

（1）取 $A_k$ 的特征向量基。零特征值对应的向量张成核，非零特征值对应的向量张成像，故 $\mathbb C^n=\ker A_k\oplus\operatorname{Im}A_k$。

（2）条件只有 $k<j$ 的单向乘积为零，因此应按指标顺序逐个消去像空间中的分量。具体地，各像空间之和是直和。事实上，若 $v_k\in\operatorname{Im}A_k$ 且 $\sum_{k=1}^{2n}v_k=0$，作用 $A_1$ 后，由 $A_1A_j=0$（$j>1$）得到 $A_1v_1=0$；再由 $\ker A_1\cap\operatorname{Im}A_1=0$ 得 $v_1=0$。依次作用 $A_2,A_3,\ldots$ 得所有 $v_k=0$。

因此 $\sum_{k=1}^{2n}\operatorname{rank}A_k\leq n$，最多有 $n$ 个非零矩阵，即至少有 $n$ 个零矩阵。

\end{solve}

\begin{remark}[方法与条件]

证明只用 $\ker A_k\cap\operatorname{Im}A_k=0$，等价于 $\operatorname{rank}A_k^2=\operatorname{rank}A_k$；非零特征值部分无需可对角化。因此真正排除的是零特征值上的非平凡 Jordan 块。见\hyperref[app:linear-invariant]{方法卷：不变子空间}。

\end{remark}

\subsection{决赛第 15 届（2024 年），来源序号6，共同题}\label{ans:fin15-high:6}\label{pr:fin15-high:6}

对应例~\ref{ex:fin15-high:6}（第~\pageref{ex:fin15-high:6} 页）；\href{https://github.com/lzxcr/cmc-mathematics/blob/64916fa1fdbc3e444a652eca204d5669058ee321/src/problems/finals/15/06.md}{查看来源}。

\begin{solve}

记 $p(t)=1-t$、$q(t)=1+t+t^2+t^3$。因 $p,q$ 互素，Bezout 恒等式保证
\[
\ker p(A)\cap\ker q(A)=\{0\}.
\]
若 $A^4=I$，即 $p(A)q(A)=0$，写 $up+vq=1$，则任意向量 $w=u(A)p(A)w+v(A)q(A)w$；两项分别落在 $\ker q(A)$、$\ker p(A)$。因此
$\mathbb R^n=\ker p(A)\oplus\ker q(A)$，故两核维数之和为 $n$，等价于题设秩等式。

反之，秩等式意味着两核维数之和为 $n$，结合交集为零可知它们直和为全空间。在每个核上，$p(A)q(A)$ 都为零，所以 $p(A)q(A)=I-A^4=0$。

\end{solve}

\begin{remark}[方法与条件]

互素性排除两核的交集，乘积零化保证两核覆盖空间，秩等式则只记录它们的维数。一般恒等式为 $\ker(pq)(A)=\ker p(A)\oplus\ker q(A)$；其多项式工具见\hyperref[app:algebra-crt]{方法卷：中国剩余定理}，算子观点见\hyperref[app:linear-invariant]{方法卷：不变子空间}。

\end{remark}

\section{行列式、伴随矩阵与多重线性}

\label{sec:solutions-2}

先观察行列式的结构，再选择初等变换、分块或低秩分解。伴随矩阵题应先按原矩阵的秩分类；对三维向量的混合积，可以用行列式的多重线性统一处理。

\subsection{初赛第 16 届（2024 年度），来源序号3}\label{ans:ch02-16th-A:3}\label{pr:ch02-16th-A:3}

对应例~\ref{ex:ch02-16th-A:3}（第~\pageref{ex:ch02-16th-A:3} 页）；\href{https://github.com/lzxcr/cmc-mathematics/blob/64916fa1fdbc3e444a652eca204d5669058ee321/src/problems/preliminary/16/03.md}{查看来源}。

\begin{solve}

若 $\operatorname{rank}A\leq1$，则 $A^*=0$，由 $a_{ij}^2=0$ 立即得 $A=0$。下面排除 $\operatorname{rank}A=2$。

此时 $A^*$ 秩为 $1$。把它看成 $A$ 的逐项平方矩阵，则其所有二阶子式均为零，因此
\[
(a_{ij}a_{k\ell})^2=(a_{i\ell}a_{kj})^2.
\]
先说明 $A$ 必有零元。否则考察前两列，若三对行的二阶子式全部非零，则三个等式只能分别取负号：
\[
a_{11}a_{22}=-a_{12}a_{21},\quad
 a_{21}a_{32}=-a_{22}a_{31},\quad
 a_{31}a_{12}=-a_{32}a_{11}.
\]
相乘矛盾。故至少有一个二阶子式为零，其对应的伴随矩阵元素及 $A$ 元素也为零。

秩为 $1$ 的矩阵可写成 $uv^{\mathrm T}$；若 $u_i v_j=0$，则 $u_i=0$ 或 $v_j=0$，所以该元所在的行或列必全零。于是 $A^*$、从而 $A$ 有一行或一列全零。若 $A$ 的第 $i$ 行全零，则 $A^*$ 除第 $i$ 列外各列全零；由逐项平方关系，$A$ 也仅可能第 $i$ 列非零，故秩至多为 $1$，矛盾。零列的情形对转置作同样论证。

因此唯一解为 $A=0$。

\end{solve}

\subsection{初赛第 17 届（2025 年度），来源序号3}\label{ans:ch01-17th-A:3}\label{pr:ch01-17th-A:3}

对应例~\ref{ex:ch01-17th-A:3}（第~\pageref{ex:ch01-17th-A:3} 页）；\href{https://github.com/lzxcr/cmc-mathematics/blob/64916fa1fdbc3e444a652eca204d5669058ee321/src/problems/preliminary/17/03.md}{查看来源}。

\begin{solve}

（1）取 $B_j=\operatorname{diag}(j,1,\ldots,1)$ 即可。

（2）先分别固定每个行列式，再留下一个能改变差矩阵行列式的自由参数。三角块中的非对角元不会影响自身行列式；让上下两个方向的非对角元出现在不同矩阵中，它们便会在作差后相互作用。令 $c=2-m\neq0$，并取
\[
A_1=\begin{pmatrix}a_1&0\\1&1\end{pmatrix}\oplus I_{n-2},
\qquad
A_2=\begin{pmatrix}a_2&t\\0&1\end{pmatrix}\oplus I_{n-2},
\qquad
A_j=\operatorname{diag}(a_j,1,\ldots,1)\quad(3\leq j\leq m).
\]
它们的行列式分别为 $a_j$。记 $s=a_1-\sum_{j=2}^m a_j$，则
\[
A_1-\sum_{j=2}^mA_j
=\begin{pmatrix}s&-t\\1&c\end{pmatrix}\oplus cI_{n-2},
\qquad
\det\left(A_1-\sum_{j=2}^mA_j\right)=(sc+t)c^{n-2}.
\]
取有理数 $t=a_{m+1}c^{2-n}-sc$，全部条件即同时满足。

\end{solve}

\subsection{决赛第 1 届（2010 年），来源序号4}\label{ans:fin01-all:4}\label{pr:fin01-all:4}

对应例~\ref{ex:fin01-all:4}（第~\pageref{ex:fin01-all:4} 页）；\href{https://github.com/lzxcr/cmc-mathematics/blob/64916fa1fdbc3e444a652eca204d5669058ee321/src/problems/finals/01/04.md}{查看来源}。

\begin{solve}

混合积在更换三列向量时乘以相应的变换行列式，故所求值为
\[
\det\begin{pmatrix}1&0&1\\1&1&0\\0&1&1\end{pmatrix}
(\boldsymbol a\times\boldsymbol b)\cdot\boldsymbol c
=2\cdot6=12.
\]

\end{solve}

\subsection{决赛第 4 届（2013 年），来源序号3}\label{ans:fin04-all:3}\label{pr:fin04-all:3}

对应例~\ref{ex:fin04-all:3}（第~\pageref{ex:fin04-all:3} 页）；\href{https://github.com/lzxcr/cmc-mathematics/blob/64916fa1fdbc3e444a652eca204d5669058ee321/src/problems/finals/04/03.md}{查看来源}。

\begin{solve}

令 $B=A/2$，并以 $M_n$ 表示所有 $n$ 阶 $\{\pm1\}$ 矩阵行列式绝对值的最大值。沿一行展开得 $M_n\le nM_{n-1}$，所以只须证 $M_3\le4$。

对三阶矩阵，先调整列的符号使第一行全为 $1$，再从后两行各减去第一行，可见行列式是 $4$ 的整数倍。另一方面，Hadamard 不等式给出
$|\det B|\le3^{3/2}<8$，故 $M_3\le4$。从而
\[
|\det A|=2^n|\det B|\le2^n\frac{n!}{3!}\,4
=\frac{2^{n+1}}3n!.
\]

\end{solve}

\begin{remark}[方法与条件]

整数整除性与连续的长度估计相结合，给出了三阶的精确界；此后只需用余子式展开传播到任意维数。

\end{remark}

\subsection{决赛第 6 届（2015 年），来源序号4，共同题}\label{ans:fin06-lo:4}\label{pr:fin06-lo:4}

对应例~\ref{ex:fin06-lo:4}（第~\pageref{ex:fin06-lo:4} 页）；\href{https://github.com/lzxcr/cmc-mathematics/blob/64916fa1fdbc3e444a652eca204d5669058ee321/src/problems/finals/06/04.md}{查看来源}。

\begin{solve}

答案为 $(n-1)!$。令 $e=(1,\ldots,1)^T$。由 $Ae=0$ 和秩条件，$\ker A=\operatorname{span}\{e\}$。在以 $e/\sqrt n$ 为首向量的正交特征基中，$A$ 与其伴随矩阵分别为
\[
\operatorname{diag}(0,2,\ldots,n),\qquad
\operatorname{diag}(n!,0,\ldots,0).
\]
变回原基即得
\[
\operatorname{adj}A=n!\frac{ee^T}{n}=(n-1)!ee^T.
\]
取左上角元素便有 $A_{11}=(n-1)!$。

\end{solve}

\begin{remark}[方法与条件]

一般地，秩为 $n-1$ 的实对称矩阵若有核向量 $u\ne0$，则其伴随矩阵等于“非零特征值之积”乘以核上的正交投影 $uu^T/(u^Tu)$。这把余子式计算转成了\hyperref[app:linear-spectral]{谱投影}。

\end{remark}

\subsection{决赛第 7 届（2016 年），来源序号1，共同题}\label{ans:fin07-hi:1}\label{pr:fin07-hi:1}

对应例~\ref{ex:fin07-hi:1}（第~\pageref{ex:fin07-hi:1} 页）；\href{https://github.com/lzxcr/cmc-mathematics/blob/64916fa1fdbc3e444a652eca204d5669058ee321/src/problems/finals/07/01.md}{查看来源}。

\begin{solve}

答案为 $0$。交换前两行把置换矩阵两两配对，而每对行列式互为相反数。

\end{solve}

\subsection{决赛第 8 届（2017 年），来源序号1，共同题}\label{ans:fin08-hi:1}\label{pr:fin08-hi:1}

对应例~\ref{ex:fin08-hi:1}（第~\pageref{ex:fin08-hi:1} 页）；\href{https://github.com/lzxcr/cmc-mathematics/blob/64916fa1fdbc3e444a652eca204d5669058ee321/src/problems/finals/08/01.md}{查看来源}。

\begin{solve}

答案为 $0$。由 Vieta 公式，四根之和为零，故矩阵乘全 $1$ 向量为零。

\end{solve}

\subsection{决赛第 12 届（2021 年），来源序号4，共同题}\label{ans:fin12-high:4}\label{pr:fin12-high:4}

对应例~\ref{ex:fin12-high:4}（第~\pageref{ex:fin12-high:4} 页）；\href{https://github.com/lzxcr/cmc-mathematics/blob/64916fa1fdbc3e444a652eca204d5669058ee321/src/problems/finals/12/04.md}{查看来源}。

\begin{solve}

答案为 $2021\cdot1011$。每个数在全矩阵中恰出现 $2021$ 次，对称性使非对角位置的出现次数为偶数，所以每个数至少在对角线上出现一次。对角线恰有 $2021$ 个位置，故它是 $1,\ldots,2021$ 的一个排列，迹即为这些数之和。

\end{solve}

\subsection{决赛第 16 届（2025 年），来源序号8，低年级组}\label{ans:fin16-low:8}\label{pr:fin16-low:8}

对应例~\ref{ex:fin16-low:8}（第~\pageref{ex:fin16-low:8} 页）；\href{https://github.com/lzxcr/cmc-mathematics/blob/64916fa1fdbc3e444a652eca204d5669058ee321/src/problems/finals/16/08.md}{查看来源}。

\begin{solve}

分别解三个二元方程组会重复出现同一组余子式；齐次坐标把这些余子式直接组织成伴随矩阵。将三条直线的系数行向量记为 $r_1,r_2,r_3$。三个顶点的齐次坐标分别为 $r_2\times r_3$、$r_3\times r_1$、$r_1\times r_2$，其第三坐标分别为 $\det N,\det P,\det Q$，均非零，所以三线两两相交。又 $\det M\ne0$，故不共点。

将三个齐次列向量各除以第三坐标，得到第三行全为 $1$ 的顶点矩阵 $V$。这三个叉积正是 $\operatorname{adj}M$ 的三列；由 $M\operatorname{adj}M=(\det M)I$，
\[
|\det V|=\frac{(\det M)^2}{|\det N\det P\det Q|}.
\]
所以三角形面积为
\[
\boxed{\frac12|\det V|=\frac{12^2}{2\cdot1\cdot2\cdot3}=12}.
\]

\end{solve}

\section{特征值、迹与谱的代数约束}

\label{sec:solutions-3}

迹、行列式和特征多项式系数连接特征值的幂和与对称函数。讨论谱之前先说明标量域，不能从特征值相同直接推出相似，也不能从实矩阵直接推出全部特征值为实数。

\subsection{初赛第 7 届（2015 年度），来源序号2}\label{ans:pre07-math:2}\label{pr:pre07-math:2}

对应例~\ref{ex:pre07-math:2}（第~\pageref{ex:pre07-math:2} 页）；\href{https://github.com/lzxcr/cmc-mathematics/blob/64916fa1fdbc3e444a652eca204d5669058ee321/src/problems/preliminary/07/02.md}{查看来源}。

\begin{solve}

令 $p_i=\operatorname{tr}(A^i)$，$e_j$ 为特征值的 $j$ 次基本对称多项式，$e_0=1$。Newton 恒等式为
\[
ke_k=\sum_{j=1}^k(-1)^{j-1}e_{k-j}p_j.
\]
依次代入 $p_j=j$，得到
\[
e_1=1,\qquad e_2=-\frac12,\qquad e_3=\frac16,
\qquad 4e_4=\frac16+1+3-4=\frac16.
\]
故 $\boxed{\det A=1/24}$。矩阵无须可对角化，因为三角化后幂的迹仍是特征值幂和。

\end{solve}

\begin{remark}[方法与条件]

恒等式可由 $E(t)=\prod_i(1+\lambda_it)=\sum_ke_kt^k$ 的对数导数推得：$E'(t)/E(t)=\sum_{j\ge1}(-1)^{j-1}p_jt^{j-1}$，再比较系数。这样把谱的幂和转换为特征多项式系数。

\end{remark}

\subsection{初赛第 11 届（2019 年度），来源序号3}\label{ans:ch10-11th-A:3}\label{pr:ch10-11th-A:3}

对应例~\ref{ex:ch10-11th-A:3}（第~\pageref{ex:ch10-11th-A:3} 页）；\href{https://github.com/lzxcr/cmc-mathematics/blob/64916fa1fdbc3e444a652eca204d5669058ee321/src/problems/preliminary/11/03.md}{查看来源}。

\begin{solve}

取共轭会把 $\bar A A$ 变成 $A\bar A$，因此只需说明交换乘积的次序不改变特征多项式。Sylvester 行列式恒等式给出
\[
\det(tI+BC)=\det(tI+CB)
\]
对任意同阶方阵 $B,C$ 成立。于是对实数 $x$，
\[
\overline{p(x)}
=\det\bigl((x-1)I+A\bar A\bigr)
=\det\bigl((x-1)I+\bar A A\bigr)=p(x).
\]
因此 $p$ 在整个实轴上取实值，其虚部系数构成的实多项式恒为零，故 $p$ 为实系数多项式。

所用恒等式可由分块矩阵 $\begin{pmatrix}I&B\\-C&I\end{pmatrix}$ 的两个 Schur 补得到 $\det(I+BC)=\det(I+CB)$，再缩放参数，并以多项式恒等性涵盖 $t=0$。

\end{solve}

\begin{remark}[方法与条件]

实系数特征多项式只保证非实特征值共轭成对，不意味着全部特征值为实数。这里的机制是“取共轭等于交换两个因子的次序”，而 $BC$ 与 $CB$ 的特征多项式相同。

\end{remark}

\subsection{初赛第 13 届补赛（2021 年度），来源序号4}\label{ans:ch08-13th-A-makeup:4}\label{pr:ch08-13th-A-makeup:4}

对应例~\ref{ex:ch08-13th-A-makeup:4}（第~\pageref{ex:ch08-13th-A-makeup:4} 页）；\href{https://github.com/lzxcr/cmc-mathematics/blob/64916fa1fdbc3e444a652eca204d5669058ee321/src/problems/preliminary/13-makeup/04.md}{查看来源}。

\begin{solve}

（1）先用一个一般估计：任意实矩阵 $K$ 若秩为 $r$，取到其像空间的正交投影 $P$，则 $PK=K$，从而
\[
|\operatorname{tr}K|=|\langle P,K\rangle_F|
\leq\|P\|_F\|K\|_F=\sqrt r\,\|K\|_F.
\]
这里 $\|K\|_F^2$ 是元素平方和。

直接套用该估计只给出系数 $10$；要达到题设常数，可利用相似变换不改变迹和秩，却能平衡两个非对角块的大小。按前 $3$ 行与后 $2$ 行分块 $M=\begin{pmatrix}A&B\\C&D\end{pmatrix}$，则
\[
\|A\|_F^2=\|B\|_F^2=3,\qquad
\|C\|_F^2=\|D\|_F^2=2.
\]
取 $t=(2/3)^{1/4}$，以 $\operatorname{diag}(tI_3,I_2)$ 作相似变换，得到
\[
K=\begin{pmatrix}A&tB\\t^{-1}C&D\end{pmatrix},
\qquad \|K\|_F^2=5+3t^2+2t^{-2}=5+2\sqrt6.
\]
相似变换保持迹与秩，代入一般估计即得 （1）。

（2）每行之和均为 $\lambda=a_1+a_2+a_3+b_1+b_2$，所以全 $1$ 向量是相应特征向量。由 Cauchy--Schwarz 不等式，
\[
|\lambda|\leq\sqrt{3\sum_i a_i^2}+\sqrt{2\sum_j b_j^2}
=\sqrt3+\sqrt2.
\]

\end{solve}

\begin{remark}[方法与条件]

参数 $t$ 的选择是最小化 $3t^2+2t^{-2}$，其最小值为 $2\sqrt6$。这是“先平衡再估计”的实例：迹、秩和特征值不随相似变换改变，但 Frobenius 范数会改变，适当缩放能给出更好的界。

像空间的正交投影及 Frobenius 内积解释，见\hyperref[app:linear-gram]{方法卷：Gram 矩阵与投影}。

\end{remark}

\subsection{初赛第 14 届补赛（2022 年度），来源序号3}\label{ans:ch05-14th-A-makeup:3}\label{pr:ch05-14th-A-makeup:3}

对应例~\ref{ex:ch05-14th-A-makeup:3}（第~\pageref{ex:ch05-14th-A-makeup:3} 页）；\href{https://github.com/lzxcr/cmc-mathematics/blob/64916fa1fdbc3e444a652eca204d5669058ee321/src/problems/preliminary/14-makeup/03.md}{查看来源}。

\begin{solve}

每列都属于 $v$ 与全 $1$ 向量张成的空间，所以先做秩至多为二的分解，比展开 $n$ 阶行列式更直接。令 $e=(1,\ldots,1)^{\mathrm T}$，$s=e^{\mathrm T}v$、$t=e^{\mathrm T}w$。由实对称矩阵的特征向量正交性，$w^{\mathrm T}v=0$。于是
\[
B=ve^{\mathrm T}+ew^{\mathrm T}=UV^{\mathrm T},
\quad U=(v,e),\quad V=(e,w),
\qquad
V^{\mathrm T}U=\begin{pmatrix}s&n\\0&t\end{pmatrix}.
\]
由低秩行列式恒等式，
\[
\det(\lambda I_n-B)
=\lambda^{n-2}\det(\lambda I_2-V^{\mathrm T}U)
=\lambda^{n-2}(\lambda-s)(\lambda-t).
\]
因此全部特征值按代数重数计为 $\sum_i a_i$、$\sum_i b_i$ 以及 $n-2$ 个 $0$。上述多项式恒等式也涵盖 $s$ 或 $t$ 为零、二者相等的情形。

\end{solve}

\begin{remark}[方法与条件]

即使不要求 $v\perp w$，也只有一个二维谱问题：
\[
\det(\lambda I-B)=\lambda^{n-2}
\bigl((\lambda-s)(\lambda-t)-n\,w^{\mathrm T}v\bigr).
\]
所以正交性恰好使这个二次式分裂为两个给定的线性因子。

从大矩阵到小矩阵的多项式恒等式，见\hyperref[app:linear-lowrank]{方法卷：低秩分解与谱}。

\end{remark}

\subsection{初赛第 15 届（2023 年度），来源序号4}\label{ans:ch03-15th-A:4}\label{pr:ch03-15th-A:4}

对应例~\ref{ex:ch03-15th-A:4}（第~\pageref{ex:ch03-15th-A:4} 页）；\href{https://github.com/lzxcr/cmc-mathematics/blob/64916fa1fdbc3e444a652eca204d5669058ee321/src/problems/preliminary/15/04.md}{查看来源}。

\begin{solve}

令
\[
u=(\sin s\theta)_{s=1}^n,\qquad v=(\cos s\theta)_{s=1}^n.
\]
加法公式把每个矩阵元分成仅依赖行或列的两个因子，给出 $A=uv^{\mathrm T}+vu^{\mathrm T}$，因此像空间至多二维。由等距单位根的正交性（$n\geq3$），
\[
u^{\mathrm T}u=v^{\mathrm T}v=\frac n2,
\qquad u^{\mathrm T}v=0.
\]
所以 $A(u+v)=\tfrac n2(u+v)$，$A(u-v)=-\tfrac n2(u-v)$，而 $A$ 在 $\operatorname{span}\{u,v\}^{\perp}$ 上为零。全部特征值为 $n/2,-n/2$ 和 $n-2$ 个零。故
\[
\det(E+A^{2023})
=\left[1+\left(\frac n2\right)^{2023}\right]
 \left[1-\left(\frac n2\right)^{2023}\right]
=1-\left(\frac n2\right)^{4046}.
\]

\end{solve}

\begin{remark}[方法与条件]

本题先识别二维不变子空间，再把大矩阵压缩为 $\begin{pmatrix}0&n/2\\n/2&0\end{pmatrix}$，无需枚举极小多项式。一般地，若 $A=UV^{\mathrm T}$，其非零谱可在小矩阵 $V^{\mathrm T}U$ 上计算；对应的行列式恒等式为 $\det(I+UV^{\mathrm T})=\det(I+V^{\mathrm T}U)$。

\end{remark}

\subsection{决赛第 2 届（2011 年），来源序号4}\label{ans:fin02-all:4}\label{pr:fin02-all:4}

对应例~\ref{ex:fin02-all:4}（第~\pageref{ex:fin02-all:4} 页）；\href{https://github.com/lzxcr/cmc-mathematics/blob/64916fa1fdbc3e444a652eca204d5669058ee321/src/problems/finals/02/04.md}{查看来源}。

\begin{solve}

取 $A$ 的特征向量基 $v_1,\ldots,v_n$，对应特征值为 $\lambda_1,\ldots,\lambda_n$，并记其对偶基为 $\varphi_1,\ldots,\varphi_n$。秩一算子
$X_{ij}=v_i\varphi_j$ 构成 $M_n(\mathbb C)$ 的基，且
\[
AX_{ij}=\lambda_iX_{ij},\qquad X_{ij}A=\lambda_jX_{ij}.
\]
故 $\sigma_A(X_{ij})=(\lambda_i-\lambda_j)X_{ij}$，即得到一组特征向量基。

\end{solve}

\begin{remark}[方法与条件]

交换子算子的谱由原算子的特征值之差组成。特别地，若 $A$ 的不同特征值的重数为 $m_1,\ldots,m_r$，则其交换子核（与 $A$ 交换的矩阵全体）的维数为 $\sum_i m_i^2$。

\end{remark}

\subsection{决赛第 5 届（2014 年），来源序号2，共同题}\label{ans:fin05-hi:2}\label{pr:fin05-hi:2}

对应例~\ref{ex:fin05-hi:2}（第~\pageref{ex:fin05-hi:2} 页）；\href{https://github.com/lzxcr/cmc-mathematics/blob/64916fa1fdbc3e444a652eca204d5669058ee321/src/problems/finals/05/02.md}{查看来源}。

\begin{solve}

因为 $\operatorname{tr}A=6$，特征值 $2$ 的代数重数不可能为 $4$，故恰为 $3$，剩余特征值为 $0$。四个特征向量可以取为线性无关，因而 $A$ 可对角化。

当 $a_0=2$ 时，$\operatorname{rank}(A-2I)=1$。其第一行为 $(0,2,2,-2)$；比较其余各行的对角元，可知这三行均为第一行的相反数。因此
\[
A=\begin{pmatrix}2&2&2&-2\\0&0&-2&2\\0&-2&0&2\\0&-2&-2&4\end{pmatrix}.
\]
二次型由对称部分
\[
B=\frac{A+A^{\mathsf T}}2
=\begin{pmatrix}2&1&1&-1\\1&0&-2&0\\1&-2&0&0\\-1&0&0&4\end{pmatrix}
\]
决定。计算得
$\det(tI-B)=(t-2)^2(t^2-2t-11)$，故正交标准型为
\[
2y_1^2+2y_2^2+(1+2\sqrt3)y_3^2+(1-2\sqrt3)y_4^2.
\]

\end{solve}

\begin{remark}[方法与条件]

矩阵的相似对角化与二次型的正交标准化是两件不同的事。二次型只看对称部分 $(A+A^{\mathsf T})/2$；不能把非对称矩阵 $A$ 的特征值直接用作二次型的系数。

\end{remark}

\subsection{决赛第 9 届（2018 年），来源序号1，共同题}\label{ans:fin09-hi:1}\label{pr:fin09-hi:1}

对应例~\ref{ex:fin09-hi:1}（第~\pageref{ex:fin09-hi:1} 页）；\href{https://github.com/lzxcr/cmc-mathematics/blob/64916fa1fdbc3e444a652eca204d5669058ee321/src/problems/finals/09/01.md}{查看来源}。

\begin{solve}

答案为 $10$。若 $H_nv=\lambda v$，则 $(v,v)$、$(v,-v)$ 分别是 $H_{n+1}$ 对应于 $\lambda+1,\lambda-1$ 的特征向量。$H_1$ 有一组特征向量基，上述构造每步把它扩充为新矩阵的特征向量基。经过四次符号选择，因此 $H_4$ 的特征值为 $4-2j$，重数为 $\binom4j$。零特征值重数为 $6$，秩为 $16-6=10$。

\end{solve}

\subsection{决赛第 13 届（2023 年），来源序号3，共同题}\label{ans:fin13-high:3}\label{pr:fin13-high:3}

对应例~\ref{ex:fin13-high:3}（第~\pageref{ex:fin13-high:3} 页）；\href{https://github.com/lzxcr/cmc-mathematics/blob/64916fa1fdbc3e444a652eca204d5669058ee321/src/problems/finals/13/03.md}{查看来源}。

\begin{solve}

$I+A$ 幂零，故 $A$ 的全部特征值均为 $-1$，于是 $\boxed{\det A=(-1)^n}$。

\end{solve}

\subsection{决赛第 14 届（2023 年），来源序号8，低年级组}\label{ans:fin14-low:8}\label{pr:fin14-low:8}

对应例~\ref{ex:fin14-low:8}（第~\pageref{ex:fin14-low:8} 页）；\href{https://github.com/lzxcr/cmc-mathematics/blob/64916fa1fdbc3e444a652eca204d5669058ee321/src/problems/finals/14/08.md}{查看来源}。

\begin{solve}

幂等性使 $A$ 的最小多项式整除 $t(t-1)$，故 $A$ 可对角化。记其四个特征值（含重数）为 $\lambda_0,\dots,\lambda_3\in\{0,1\}$。

令 $P$ 为循环移位矩阵，即 $P_{12}=P_{23}=P_{34}=P_{41}=1$，其余元为零。它有互异特征值 $1,i,-1,-i$。插值给出次数至多 $3$ 的多项式 $h(t)=\sum_{j=0}^3c_jt^j$，使 $h(i^k)=\lambda_k$。于是 $h(P)$ 与 $A$ 具有相同的对角标准形，故相似；而 $h(P)$ 恰是题设循环矩阵。

\end{solve}

\begin{remark}[方法与条件]

循环移位矩阵的特征值互异，因而可以用插值任意指定其各特征空间上的作用。这证明每个复可对角化矩阵都相似于同阶循环矩阵；幂等性只用于保证可对角化。插值基与构造公式见\hyperref[app:algebra-interpolation]{方法卷：多项式插值}。

\end{remark}

\subsection{决赛第 17 届（2026 年），来源序号8，低年级组}\label{ans:fin17-low:6}\label{pr:fin17-low:6}

对应例~\ref{ex:fin17-low:6}（第~\pageref{ex:fin17-low:6} 页）；\href{https://github.com/lzxcr/cmc-mathematics/blob/64916fa1fdbc3e444a652eca204d5669058ee321/src/problems/finals/17/08.md}{查看来源}。

\begin{solve}

在复数域上，$A,B$ 均可对角化，因为它们的最小多项式分别整除无重根多项式 $t^{2025}-1$ 与 $t^{2026}-1$。又二者交换，$A$ 的每个特征空间都对 $B$ 不变；$B$ 在其上的最小多项式仍无重根，故可在各子空间分别对角化，得到共同特征基。于是只需证明相应特征值满足 $1+\lambda+\mu\ne0$。

若 $1+\lambda+\mu=0$，由 $|\lambda|=|\mu|=1$ 得 $|1+\lambda|=1$，从而 $\operatorname{Re}\lambda=-1/2$。因此 $\lambda,\mu$ 恰为两个非平凡三次单位根。但 $2026\equiv1\pmod3$，导致 $\mu^{2026}=\mu\ne1$，矛盾。故 $A+B+I$ 可逆。

矩阵部分归结为共同不变子空间，而最后的矛盾来自单位根的阶；见\href{https://github.com/lzxcr/cmc-mathematics/blob/64916fa1fdbc3e444a652eca204d5669058ee321/src/methods/outer/01-algebra.md}{方法卷：阶与互素分解（原稿链接）}。

\end{solve}

\section{多项式、整除、插值与根的约束}

\label{sec:solutions-4}

本专题研究系数、根与整除关系，包含有理数域上的不可约性。将次数界和线性代数结合，经常能把无穷多个取值条件化为有限个独立条件；整系数问题还必须使用整数性。

\subsection{初赛第 5 届（2013 年度），来源序号2}\label{ans:pre05-math:2}\label{pr:pre05-math:2}

对应例~\ref{ex:pre05-math:2}（第~\pageref{ex:pre05-math:2} 页）；\href{https://github.com/lzxcr/cmc-mathematics/blob/64916fa1fdbc3e444a652eca204d5669058ee321/src/problems/preliminary/05/02.md}{查看来源}。

\begin{solve}

取一个解 $X=(x_1,\ldots,x_n)^{\mathrm T}$，将 $B(t_0)$ 的列记为 $v_j$，则 $b(t_0)=\sum_jx_jv_j$。行列式对替换列线性，其中 $j\ne i$ 的项有两列相同，故
\[
d_i(t_0)=x_i\det B(t_0)=0.
\]
所以 $t-t_0$ 同时整除 $d,d_1,\ldots,d_n$，即为所求。

\end{solve}

\subsection{初赛第 10 届（2018 年度），来源序号3}\label{ans:pre10-math:3}\label{pr:pre10-math:3}

对应例~\ref{ex:pre10-math:3}（第~\pageref{ex:pre10-math:3} 页）；\href{https://github.com/lzxcr/cmc-mathematics/blob/64916fa1fdbc3e444a652eca204d5669058ee321/src/problems/preliminary/10/03.md}{查看来源}。

\begin{solve}

（1）若 $A,A^{-1}$ 都是整矩阵，则两个行列式均为整数且乘积为 $1$，所以 $\det A=\pm1$。反之，$A^{-1}=\operatorname{adj}(A)/\det A$ 仍为整矩阵。

（2）只需使用 $k=1$ 的条件。整矩阵满足
\[
A+B\equiv A-2B\pmod3.
\]
行列式是矩阵元素的整系数多项式，故由（1）得
\[
\det(A+B)\equiv\det(A-2B)\equiv\pm1\pmod3.
\]
因此 $\det(A+B)\ne0$。

\end{solve}

\begin{remark}[方法与条件]

\textbf{条件可以减弱。} $k=0,\ldots,n$ 当然足够，甚至任意一个 $k\ge1$ 的条件就够：改取模 $2k+1$ 即可。这说明原题的可逆结论来自整数性，所列测试点远多于所需。

\textbf{更多节点给出更强结论。} 令 $p(t)=\det(A-tB)$，则 $\deg p\le n$。若在任意 $2n+1$ 个互异实节点上 $p$ 都取值于 $\{1,-1\}$，那么 $p^2-1$ 的零点数超过次数，故 $p\equiv1$ 或 $p\equiv-1$。这里与节点的具体取值、间距及正负无关。因此题设实际上保证 $A-tB$ 对每个实数 $t$ 可逆，且整数 $t$ 时逆矩阵仍为整矩阵。

对本题的等距整数节点，当 $n\ge2$ 时，$k=0,\ldots,n$ 也足以推出这个更强结论。由 $p\in\mathbb Z[t]$ 得 $p(4)\equiv p(0)\pmod4$，而两者均为 $\pm1$，故相等。于是 $p(t)-p(0)=t(t-4)q(t)$，其中 $q\in\mathbb Z[t]$，代入 $t=2$ 又得 $p(2)=p(0)$。其余偶数节点模 $4$ 与 $0$ 或 $2$ 同余，故 $n+1$ 个指定值全相等，迫使 $p$ 恒定。$n=1$ 时这个强化结论不成立：$A=B=[1]$ 满足 $k=0,1$ 的条件，但 $A+B=[2]$ 的逆不是整矩阵；可逆性仍由主解保证。

\end{remark}

\subsection{初赛第 12 届（2020 年度），来源序号4}\label{ans:ch09-12th-A:4}\label{pr:ch09-12th-A:4}

对应例~\ref{ex:ch09-12th-A:4}（第~\pageref{ex:ch09-12th-A:4} 页）；\href{https://github.com/lzxcr/cmc-mathematics/blob/64916fa1fdbc3e444a652eca204d5669058ee321/src/problems/preliminary/12/04.md}{查看来源}。

\begin{solve}

答案是 $f(x)=x^{2020}$，它可取 $g=x$、$h=x^{2018}$。

设根按重数编号为 $z_1,\ldots,z_{2020}$。对每个编号 $k$，将其余 $2019$ 个根按出现次数分配给 $g_k,h_k$。第二项系数相等，由 Viète 公式等价于两组根之和相等，所以
\[
\sum_{j\neq k}\varepsilon_{kj}z_j=0,
\qquad \varepsilon_{kj}\in\{1,-1\}.
\]
即使有重根，也将每一次出现视为一个独立编号，因此每个系数仍是 $\pm1$。

各行的正负号由不同分解决定，无法直接比较；但模 $2$ 后这些符号差异全部消失。令上述方程组的整数系数矩阵为 $A$。它对角元为零，非对角元为 $\pm1$，所以模 $2$ 有 $A\equiv J-I$，其中 $J$ 为全 $1$ 矩阵。于是
\[
\det A\equiv\det(J-I)
=2019(-1)^{2019}\equiv1\pmod2.
\]
故 $A$ 可逆，全部 $z_j=0$，得到唯一答案。

\end{solve}

\begin{remark}[方法与条件]

同样的论证对任意偶数次数 $n\geq4$ 成立，唯一解为 $x^n$。奇数次数则不同：例如 $f(x)=(x-c)^n$，可取 $g=h=(x-c)^{(n-1)/2}$，任意 $c\in\mathbb C$ 都满足条件。重根本身不是模 $2$ 论证的障碍，关键是保留每个根的重数编号。

由根的对称关系建立系数约束的通用观点，见\hyperref[app:algebra-roots]{方法卷：根与系数}。

\end{remark}

\subsection{初赛第 13 届（2021 年度），来源序号3}\label{ans:ch07-13th-A:3}\label{pr:ch07-13th-A:3}

对应例~\ref{ex:ch07-13th-A:3}（第~\pageref{ex:ch07-13th-A:3} 页）；\href{https://github.com/lzxcr/cmc-mathematics/blob/64916fa1fdbc3e444a652eca204d5669058ee321/src/problems/preliminary/13/03.md}{查看来源}。

\begin{solve}

若全根实，根的平方既非负，其乘积又受非零整数常数项控制；这会对平方和给出与次数成正比的下界。反设全部根 $r_1,\ldots,r_{2021}$ 都为实数。常数项为非零整数，故
\[
\prod_{j=1}^{2021}r_j^2=a_0^2\geq1.
\]
算术---几何平均不等式及 Newton 恒等式给出
\[
2021\leq\sum_{j=1}^{2021}r_j^2
=a_{2020}^2-2a_{2019}
\leq40^2+80=1680,
\]
矛盾。所以至少有一对非实共轭根。

\end{solve}

\begin{remark}[方法与条件]

一般地，首一整系数 $n$ 次多项式若常数项非零且全根实，则必有 $n\leq a_{n-1}^2-2a_{n-2}$。因此仅控制最高的两个非首项系数，就足以给出次数障碍；不必同时估计倒数根。

这里使用的二阶 Newton 恒等式及其与 Viète 公式的关系，见\hyperref[app:algebra-roots]{方法卷：根与系数}。

\end{remark}

\subsection{初赛第 14 届（2022 年度），来源序号4}\label{ans:ch04-14th-A:4}\label{pr:ch04-14th-A:4}

对应例~\ref{ex:ch04-14th-A:4}（第~\pageref{ex:ch04-14th-A:4} 页）；\href{https://github.com/lzxcr/cmc-mathematics/blob/64916fa1fdbc3e444a652eca204d5669058ee321/src/problems/preliminary/14/04.md}{查看来源}。

\begin{solve}

对常数项非零的多项式 $p$，记 $p^\#(x)=x^{\deg p}p(1/x)$。若 $f=gh$ 是非平凡分解，由 Gauss 引理可取 $g,h\in\mathbb Z[x]$ 首一，常数项均为 $\pm1$。翻转一个因子不会改变它与互反多项式的乘积，却可能改变系数排列；稀疏性将限制这种改变。令 $k=gh^\#$，则
\[
kk^\#=ff^\#,
\qquad
f^\#=1-x^{n-1}-x^n.
\]
若 $n=2$，判别式为 $5$，结论显然。以下设 $n\geq3$。写 $k=\sum_{j=0}^n b_jx^j$；比较 $x^n$ 系数，得到
\[
\sum_{j=0}^n b_j^2=3.
\]
因为 $b_0,b_n=\pm1$，所以 $k=b_nx^n+b_ix^i+b_0$，其中 $1\leq i<n$，三个系数均为 $\pm1$。比较常数项得 $b_0=-b_n$。再比较次数小于 $n$ 的部分：
\[
[kk^\#]_{<n}=-1+b_nb_i(x^i-x^{n-i}),
\qquad
[ff^\#]_{<n}=-1-x+x^{n-1}.
\]
故 $i=1$ 或 $i=n-1$，分别给出 $k=\pm f$ 或 $k=\pm f^\#$。约去公因子可知 $h^\#=\pm h$ 或 $g^\#=\pm g$。因此 $f$ 与 $f^\#$ 必有公共复根 $z$。

然而公共根必须满足
\[
z^n=z+1,\qquad z^{n-1}=-z,
\]
故 $z^2+z+1=0$，$z$ 是三次单位根，同时 $z^n=-z^2$。左边的立方为 $1$，右边的立方为 $-1$，矛盾。因此 $f$ 在 $\mathbb Q$ 上不可约。

\end{solve}

\begin{remark}[方法与条件]

“翻转因子”保留了 $ff^\#$，而其中央系数是原多项式各系数的平方和。这个平方和仅为 $3$，整系数就迫使新多项式也只有三个非零系数；剩余系数比较使某个真因子必须互反，最后由公共根矛盾排除。这里发挥作用的是稀疏整数系数的刚性，而非对所有首一多项式都成立的根分布规律。

\end{remark}

\subsection{决赛第 5 届（2014 年），来源序号4，共同题}\label{ans:fin05-hi:4}\label{pr:fin05-hi:4}

对应例~\ref{ex:fin05-hi:4}（第~\pageref{ex:fin05-hi:4} 页）；\href{https://github.com/lzxcr/cmc-mathematics/blob/64916fa1fdbc3e444a652eca204d5669058ee321/src/problems/finals/05/04.md}{查看来源}。

\begin{solve}

答案是 $C=\sqrt{(n-2)/n}$。$n=2$ 时结论显然，以下设 $n>2$。

设 $p$ 首一，根为 $r_1\le\cdots\le r_n$。令
$D=\operatorname{diag}(r_1,\ldots,r_n)$、$e=(1,\ldots,1)^{\mathsf T}/\sqrt n$，并把 $D$ 正交压缩到 $e^\perp$，记所得对称算子为 $T$。分块行列式公式给出
\[
\det(tI_{n-1}-T)
=\det(tI_n-D)\,e^{\mathsf T}(tI_n-D)^{-1}e
=\frac1n p'(t).
\]
这里先在 $t\ne r_i$ 时计算，再按多项式恒等式延拓。因此 $p'$ 的根就是 $T$ 的特征值。

以 $e_i$ 表示标准基。在 $e^\perp$ 中取两个正交单位向量
\[
u=\frac{e_1-e_n}{\sqrt2},\qquad
v=\sqrt{\frac{n}{2(n-2)}}
\left(e_1+e_n-\frac2n\sum_{j=1}^n e_j\right).
\]
有
\[
u^{\mathsf T}Tv=u^{\mathsf T}Dv
=\frac{r_1-r_n}{2}\sqrt{\frac{n-2}{n}}.
\]
单位向量 $(u+v)/\sqrt2$ 与 $(u-v)/\sqrt2$ 的 Rayleigh 商之差为 $2u^{\mathsf T}Tv$，而任意两个 Rayleigh 商之差的绝对值不超过谱直径。故
\[
d(p')\ge2|u^{\mathsf T}Tv|=\sqrt{\frac{n-2}{n}}\,d(p).
\]
最后取 $p(x)=x(x-1)(x-1/2)^{n-2}$。令 $z=x-1/2$，则
\[
p'(x)=z^{n-3}\left(nz^2-\frac{n-2}{4}\right),
\]
从而 $d(p')=\sqrt{(n-2)/n}$、$d(p)=1$，证明常数最优。

\end{solve}

\begin{remark}[方法与条件]

实根多项式求导可以理解为实对称算子的正交压缩。Rolle 定理的零点交错随之成为谱交错，而本题的定量界来自两个明确试验向量的 Rayleigh 商。这一联系把零点几何纳入了统一的谱理论。

\end{remark}

\subsection{决赛第 8 届（2017 年），来源序号2，共同题}\label{ans:fin08-hi:2}\label{pr:fin08-hi:2}

对应例~\ref{ex:fin08-hi:2}（第~\pageref{ex:fin08-hi:2} 页）；\href{https://github.com/lzxcr/cmc-mathematics/blob/64916fa1fdbc3e444a652eca204d5669058ee321/src/problems/finals/08/02.md}{查看来源}。

\begin{solve}

答案为 $a<-37$ 或 $a>27$。记左端多项式为 $P$，则
\[
P'(x)=12(x+2)(x-1)(x-3),\qquad(P(-2),P(1),P(3))=(a-152,a+37,a-27).
\]
由四个单调区间可知，四根均为实数（计重数）当且仅当两处极小值非正、极大值非负，即 $-37\leq a\leq27$。

\end{solve}

\subsection{决赛第 16 届（2025 年），来源序号6，共同题}\label{ans:fin16-high:6}\label{pr:fin16-high:6}

对应例~\ref{ex:fin16-high:6}（第~\pageref{ex:fin16-high:6} 页）；\href{https://github.com/lzxcr/cmc-mathematics/blob/64916fa1fdbc3e444a652eca204d5669058ee321/src/problems/finals/16/06.md}{查看来源}。

\begin{solve}

因子关系直接传递的是根，而不是系数；因此先用原多项式系数约束全部根，再用 Vieta 公式约束因子的系数。反设全部 $|a_j|\le2024$。若 $p(z)=0$ 且 $|z|\ge2025$，则
\[
|z|^n\le2024\sum_{j=0}^{n-1}|z|^j
=2024\frac{|z|^n-1}{|z|-1}<|z|^n,
\]
矛盾。所以 $p$ 的全部复根的模小于 $2025$，其因子 $q$ 的根 $\zeta_1,\dots,\zeta_m$ 亦然。由 Vieta 公式，
\[
|b_k|\le\sum_{1\le i_1<\cdots<i_k\le m}|\zeta_{i_1}\cdots\zeta_{i_k}|
<\binom mk2025^k,
\]
与题设矛盾。

\end{solve}

\begin{remark}[方法与条件]

Cauchy 根界与 Vieta 公式将“原多项式系数小”传递为“任意首一因子的系数受控”。一般界和证明见\hyperref[app:algebra-roots]{方法卷：根与系数}。

\end{remark}

\section{循环向量、最小多项式与中心化子}

\label{sec:solutions-5}

有循环向量时，一个与原变换交换的算子由它在该向量上的值完全确定。最小多项式不仅控制可对角化性，还决定矩阵多项式化简、中心化子空间以及插值构造的维数。

\subsection{初赛第 1 届（2009 年度），来源序号2}\label{ans:pre01-math:2}\label{pr:pre01-math:2}

对应例~\ref{ex:pre01-math:2}（第~\pageref{ex:pre01-math:2} 页）；\href{https://github.com/lzxcr/cmc-mathematics/blob/64916fa1fdbc3e444a652eca204d5669058ee321/src/problems/preliminary/01/02.md}{查看来源}。

\begin{solve}

关键是 $F$ 的一个向量已生成整个空间：对标准基有
\[
e_j=F^{j-1}e_1,\qquad 1\le j\le n.
\]
因此，与 $F$ 交换的变换只要在 $e_1$ 上相同，就在每个 $e_j$ 上相同。

令 $p(t)=\sum_{j=1}^n a_{j1}t^{j-1}$。由 $Ae_1=p(F)e_1$ 及交换关系，
\[
Ae_j=AF^{j-1}e_1=F^{j-1}Ae_1=p(F)e_j,
\]
故 $A=p(F)$。

于是 $C(F)$ 由 $I,F,\ldots,F^{n-1}$ 张成。这些矩阵线性无关，因为把它们的线性关系作用于 $e_1$ 就得到标准基的线性关系。因此
\[
\boxed{\dim C(F)=n.}
\]

\end{solve}

\begin{remark}[方法与条件]

$e_1$ 称为 $F$ 的循环向量。更一般地，任何具有循环向量的线性变换 $T$，其交换子代数都恰为 $\mathbb C[T]$。这里决定结论的是循环结构，特征值是否互异并不重要。见\hyperref[app:linear-cyclic]{方法卷：循环向量与多项式模型}。

\end{remark}

\subsection{初赛第 3 届（2011 年度），来源序号3}\label{ans:pre03-math:3}\label{pr:pre03-math:3}

对应例~\ref{ex:pre03-math:3}（第~\pageref{ex:pre03-math:3} 页）；\href{https://github.com/lzxcr/cmc-mathematics/blob/64916fa1fdbc3e444a652eca204d5669058ee321/src/problems/preliminary/03/03.md}{查看来源}。

\begin{solve}

令 $E_{ij}$ 表示矩阵单位。$\sigma$ 与投影 $E_{ii}$ 交换，所以
\[
\sigma(e_i)=\sigma(E_{ii}e_i)=E_{ii}\sigma(e_i)\in\mathbb F e_i.
\]
写作 $\sigma(e_i)=\lambda_i e_i$。再与把 $e_j$ 送到 $e_i$ 的 $E_{ij}$ 交换，得到
\[
\lambda_i e_i
=\sigma(E_{ij}e_j)
=E_{ij}\sigma(e_j)
=\lambda_j e_i.
\]
因此所有 $\lambda_i$ 相等，结论成立。

\end{solve}

\begin{remark}[方法与条件]

这说明完整矩阵代数的中心只有标量矩阵。证明不依赖特征值，因此对任意域都有效。

\end{remark}

\subsection{初赛第 14 届（2022 年度），来源序号3}\label{ans:ch04-14th-A:3}\label{pr:ch04-14th-A:3}

对应例~\ref{ex:ch04-14th-A:3}（第~\pageref{ex:ch04-14th-A:3} 页）；\href{https://github.com/lzxcr/cmc-mathematics/blob/64916fa1fdbc3e444a652eca204d5669058ee321/src/problems/preliminary/14/03.md}{查看来源}。

\begin{solve}

比较 $AB=BA$ 的元素，得到
\[
B=\begin{pmatrix}a&b\\b&a-b\end{pmatrix},\qquad
 a>b>0,\qquad a^2-ab-b^2=1.
\]
乘以 $A^{-1}$ 保持交换关系和行列式；需要找一个在此操作下严格下降、而正性仍保持的整数参数。这里选择非对角元 $b$。对正整数 $b$ 作递降。若 $b=1$，则 $a=2$，所以 $B=A$。

若 $b>1$，方程 $a(a-b)=b^2+1$ 蕴含 $a<2b$。令
\[
B_1=A^{-1}B
=\begin{pmatrix}a-b&2b-a\\2b-a&2a-3b\end{pmatrix}
=\begin{pmatrix}a_1&b_1\\b_1&a_1-b_1\end{pmatrix}.
\]
这里 $a_1=a-b>0$、$0<b_1=2b-a<b$，且 $\det B_1=1$ 给出 $a_1(a_1-b_1)=1+b_1^2>0$。所以 $B_1$ 仍满足全部条件，而参数 $b$ 严格下降。有限步后必到达 $A$，倒推得 $B=A^k$，其中 $k\geq1$。

\end{solve}

\begin{remark}[方法与条件]

递降给出 Fibonacci 公式
\[
A^k=\begin{pmatrix}F_{2k+1}&F_{2k}\\F_{2k}&F_{2k-1}\end{pmatrix}.
\]
方程也可写成 $(2a-b)^2-5b^2=4$，因而是 Pell 型方程。$A$ 的特征值为 $\varphi^2,\varphi^{-2}$，其中 $\varphi=(1+\sqrt5)/2$；矩阵乘以 $A^{-1}$ 对应沿单位 $\varphi^2$ 递降，而不是沿 $\varphi$ 递降。

\end{remark}

\subsection{初赛第 14 届第二次补赛（2022 年度），来源序号3}\label{ans:ch06-14th-A-makeup2:3}\label{pr:ch06-14th-A-makeup2:3}

对应例~\ref{ex:ch06-14th-A-makeup2:3}（第~\pageref{ex:ch06-14th-A-makeup2:3} 页）；\href{https://github.com/lzxcr/cmc-mathematics/blob/64916fa1fdbc3e444a652eca204d5669058ee321/src/problems/preliminary/14-makeup2/03.md}{查看来源}。

\begin{solve}

矩阵代入只识别多项式模其极小多项式的余类，因此问题等价于同时满足 $f\equiv f_i\pmod{p_i}$。记 $p_i$ 为 $A_i$ 的极小多项式，$P=\prod_i p_i$、$P_i=P/p_i$。因 $P_i$ 与 $p_i$ 互素，可取 $u_i\in K[x]$ 满足
\[
u_iP_i\equiv1\pmod{p_i}.
\]
令 $f=\sum_i f_i u_iP_i$，则对每个 $i$，$f\equiv f_i\pmod{p_i}$，故 $f(A_i)=f_i(A_i)$。这就是 $K[x]$ 上的中国剩余定理。

\end{solve}

\begin{remark}[方法与条件]

所得 $f$ 在模 $P$ 意义下唯一，因此可唯一选取 $\deg f<\deg P$ 的代表。当模多项式为互异线性因子时，这就是 Lagrange 插值；当模为 $(x-\lambda_i)^{m_i}$ 时，它变成指定各点若干阶导数的 Hermite 插值，并产生广义特征子空间的谱投影。

投影多项式为何能拼接局部数据及其唯一性，见\hyperref[app:algebra-crt]{方法卷：中国剩余定理}；导数插值见\hyperref[app:algebra-interpolation]{方法卷：多项式插值}。

\end{remark}

\subsection{初赛第 15 届（2023 年度），来源序号3}\label{ans:ch03-15th-A:3}\label{pr:ch03-15th-A:3}

对应例~\ref{ex:ch03-15th-A:3}（第~\pageref{ex:ch03-15th-A:3} 页）；\href{https://github.com/lzxcr/cmc-mathematics/blob/64916fa1fdbc3e444a652eca204d5669058ee321/src/problems/preliminary/15/03.md}{查看来源}。

\begin{solve}

一个向量能否生成整个空间，取决于它在每个特征值对应的 Jordan 链中能走多远；同一特征值若有两条链，单个向量就不能分别控制它们。由 Jordan 标准形，某特征值的几何重数就是其 Jordan 块数；而其在极小多项式中的重数是最大 Jordan 块的大小。因此
\[
\text{每个特征值的几何重数为 }1
\quad\Longleftrightarrow\quad \mu_A=\chi_A.
\]
下面说明这等价于存在循环向量。

若 $\alpha,A\alpha,\ldots,A^{n-1}\alpha$ 是一组基，则不存在次数小于 $n$ 的非零多项式零化 $A$，故 $\deg\mu_A=n$，于是 $\mu_A=\chi_A$。

反之，每个特征值 $\lambda_i$ 只有一个大小为 $d_i$ 的 Jordan 块时，取该块的循环向量 $\alpha_i$，令 $\alpha=\sum_i\alpha_i$。若 $p(A)\alpha=0$，由广义特征子空间的直和性，$(x-\lambda_i)^{d_i}\mid p$ 对所有 $i$ 成立。因这些因子两两互素，$\chi_A\mid p$。所以任何次数小于 $n$ 的零化关系只能是零关系，所列 $n$ 个向量构成一组基。

\end{solve}

\begin{remark}[方法与条件]

“循环向量”“极小多项式等于特征多项式”“每个特征值只有一个 Jordan 块”是同一结构的三种表达。前两项对任意域都等价；第三项要求特征多项式在该域分裂。循环算子的交换子代数由其多项式组成，见\hyperref[pr:ch01-17th-A:4]{循环子空间与可交换算子}。

循环向量与多项式模的完整联系，见\hyperref[app:linear-cyclic]{方法卷：循环算子}。

\end{remark}

\subsection{初赛第 17 届（2025 年度），来源序号4}\label{ans:ch01-17th-A:4}\label{pr:ch01-17th-A:4}

对应例~\ref{ex:ch01-17th-A:4}（第~\pageref{ex:ch01-17th-A:4} 页）；\href{https://github.com/lzxcr/cmc-mathematics/blob/64916fa1fdbc3e444a652eca204d5669058ee321/src/problems/preliminary/17/04.md}{查看来源}。

\begin{solve}

（1）写 $\mu_A=\prod_{i=1}^s p_i^{m_i}$，其中 $p_i$ 是互异的首一不可约多项式。对两两互素的 $p_i^{m_i}$ 应用中国剩余定理，可构造多项式投影，把空间按不同不可约因子分开，得到准素分解
\[
V=\bigoplus_{i=1}^s V_i,\qquad
V_i=\ker p_i(\mathcal A)^{m_i},\qquad
\mu_{\mathcal A|_{V_i}}=p_i^{m_i}.
\]
每个限制算子的极小多项式恰为所列幂：若其中一个幂次可以降低，整个空间的零化多项式也能降低，与 $\mu_A$ 的极小性矛盾。因此可取 $\alpha_i\in V_i$ 使 $p_i(\mathcal A)^{m_i-1}\alpha_i\neq0$，于是 $\mu_{A,\alpha_i}=p_i^{m_i}$。令 $\alpha=\sum_i\alpha_i$。由直和性，$q(\mathcal A)\alpha=0$ 当且仅当每个 $q(\mathcal A)\alpha_i=0$，因此
\[
\mu_{A,\alpha}=\operatorname{lcm}(p_1^{m_1},\ldots,p_s^{m_s})=\mu_A.
\]

（2）映射 $K[x]\to V$，$q\mapsto q(\mathcal A)\alpha$ 的像为 $K[\mathcal A]\alpha$，核为主理想 $(\mu_{A,\alpha})$。故
\[
K[\mathcal A]\alpha\cong K[x]/(\mu_{A,\alpha}),
\]
其维数为 $d=\deg\mu_{A,\alpha}$，一组基为 $\alpha,\mathcal A\alpha,\ldots,\mathcal A^{d-1}\alpha$。乘以 $x$ 对应作用 $\mathcal A$，故该子空间不变。

（3）一个与 $\mathcal A$ 可交换的算子作用在 $g(\mathcal A)\alpha$ 上，完全由它作用在 $\alpha$ 上的结果决定。由前两问，$\mu_A=\chi_A$ 时可取 $\alpha$ 使 $V=K[\mathcal A]\alpha$。若 $\mathcal B$ 与 $\mathcal A$ 可交换，先写 $\mathcal B\alpha=f(\mathcal A)\alpha$，则对任意 $g$ 有
\[
\mathcal B(g(\mathcal A)\alpha)
=g(\mathcal A)\mathcal B\alpha
=f(\mathcal A)g(\mathcal A)\alpha.
\]
所以 $\mathcal B=f(\mathcal A)$。反向结论由算子与其多项式可交换立即得到。

\end{solve}

\begin{remark}[方法与条件]

把 $V$ 看作 $K[x]$-模，其中 $x\cdot v=\mathcal Av$，便能统一三问：向量生成的子模是 $K[x]/(\mu_{A,\alpha})$；条件 $\mu_A=\chi_A$ 等价于整个模循环。循环模的自同态由生成元的像决定，于是其交换子代数正是 $K[\mathcal A]$。这些结论对任意域都成立。

准素投影、向量的零化理想与循环生成的系统说明，见\hyperref[app:linear-cyclic]{方法卷：循环算子与多项式模}。

\end{remark}

\subsection{决赛第 5 届（2014 年），来源序号6，低年级组}\label{ans:fin05-lo:3}\label{pr:fin05-lo:3}

对应例~\ref{ex:fin05-lo:3}（第~\pageref{ex:fin05-lo:3} 页）；\href{https://github.com/lzxcr/cmc-mathematics/blob/64916fa1fdbc3e444a652eca204d5669058ee321/src/problems/finals/05/06.md}{查看来源}。

\begin{solve}

事实上不需要可对角化假设。记 $e_n$ 为最后一个标准基向量。由矩阵形状，$A^ke_n$ 在第 $n-k$ 行以上均为零，而其第 $n-k$ 个分量为
\[
a_{n-k,n-k+1}\cdots a_{n-1,n}\ne0\qquad(0\le k<n),
\]
其中空积为 $1$。所以 $e_n,Ae_n,\ldots,A^{n-1}e_n$ 是一组基。

$W$ 显然是线性子空间。给定 $X\in W$，唯一选取次数小于 $n$ 的多项式 $p$，使 $Xe_n=p(A)e_n$。由交换关系，对每个 $k$ 都有
\[
XA^ke_n=A^kXe_n=p(A)A^ke_n.
\]
在一组基上相同，故 $X=p(A)$。上述唯一性也证明 $I,A,\ldots,A^{n-1}$ 线性无关。

\end{solve}

\begin{remark}[方法与条件]

$e_n$ 是一个循环向量；交换算子在循环向量上的值决定其全部作用。见\hyperref[app:linear-cyclic]{方法卷：循环向量}。

\end{remark}

\subsection{决赛第 13 届（2023 年），来源序号6，共同题}\label{ans:fin13-high:6}\label{pr:fin13-high:6}

对应例~\ref{ex:fin13-high:6}（第~\pageref{ex:fin13-high:6} 页）；\href{https://github.com/lzxcr/cmc-mathematics/blob/64916fa1fdbc3e444a652eca204d5669058ee321/src/problems/finals/13/06.md}{查看来源}。

\begin{solve}

将 $p$ 分解为不同首一不可约因子的幂。对每个不可约因子 $q\mid p$，子空间
\[
V_q=\ker\bigl((p/q)(A)\bigr)
\]
都是真子空间，否则 $p/q$ 已能零化 $A$，与 $p$ 的最小性矛盾。

每个真子空间 $V_q$ 包含于某个非零线性函数 $\ell_q$ 的核；乘积 $\prod_q\ell_q$ 是非零多项式，在无限域 $\mathbb C$ 上不可能处处为零。因此可取 $b\notin\bigcup_{q\mid p}V_q$。设 $m_b$ 是 $b$ 的最小多项式，则 $m_b\mid p$。若 $m_b\ne p$，至少一个不可约因子 $q$ 在 $m_b$ 中的指数小于在 $p$ 中的指数，于是 $m_b\mid p/q$，导致 $b\in V_q$，矛盾。因此 $m_b=p$。

\end{solve}

\begin{remark}[方法与条件]

满足 $m_b=p$ 的向量构成有限个真线性子空间之并的补集，因此是稠密开集；“一般位置”的向量已记录了矩阵最小多项式的全部信息。这不意味着它是整个空间的循环向量：只有 $\deg p=n$ 时，其 Krylov 向量 $b,Ab,\dots,A^{n-1}b$ 才构成一组基。

\end{remark}

\section{幂零结构、不变子空间与相似}

\label{sec:solutions-6}

通过核与像的稳定、幂的秩和若当链识别矩阵结构。交换关系可以传递到不变子空间，但交换子本身不必为零；证明幂零性时应交代从迹的幂和到谱结论的依据。

\subsection{初赛第 1 届（2009 年度），来源序号3}\label{ans:pre01-math:3}\label{pr:pre01-math:3}

对应例~\ref{ex:pre01-math:3}（第~\pageref{ex:pre01-math:3} 页）；\href{https://github.com/lzxcr/cmc-mathematics/blob/64916fa1fdbc3e444a652eca204d5669058ee321/src/problems/preliminary/01/03.md}{查看来源}。

\begin{solve}

交换子关系对所有幂都给出信息。由乘积法则归纳可得
\[
f^kg-gf^k=kf^k\qquad(k\ge1).
\]
取迹，得到 $\operatorname{tr}(f^k)=0$。若 $\dim V=n$，Newton 恒等式表明 $f$ 的特征多项式为 $t^n$，故其全部特征值为零。

特别地，$W=\ker f\ne\{0\}$。对 $v\in W$，
\[
f(gv)=g(fv)+fv=0,
\]
所以 $W$ 在 $g$ 下不变。复线性变换 $g|_W$ 有特征向量，它同时是 $f$ 的零特征向量。

\end{solve}

\begin{remark}[方法与条件]

也可从相似变换理解第一步：令 $H(t)=e^{-tg}fe^{tg}$，则 $H'(t)=H(t)$，故 $e^{-tg}fe^{tg}=e^tf$。相似变换不改变谱，而有限集合不可能在任意正数伸缩下保持不变，除非只含零。原来的交换子关系因此表达了一种“无穷小伸缩”。

\end{remark}

\subsection{初赛第 2 届（2010 年度），来源序号2}\label{ans:pre02-math:2}\label{pr:pre02-math:2}

对应例~\ref{ex:pre02-math:2}（第~\pageref{ex:pre02-math:2} 页）；\href{https://github.com/lzxcr/cmc-mathematics/blob/64916fa1fdbc3e444a652eca204d5669058ee321/src/problems/preliminary/02/02.md}{查看来源}。

\begin{solve}

假设 $X^2=B$。由于 $B^3=0$，有 $X^6=0$，所以 $X$ 是三阶幂零矩阵。其特征多项式为 $t^3$，由 Cayley--Hamilton 定理得 $X^3=0$，从而
\[
B^2=X^4=0.
\]
但直接相乘可见 $(B^2)_{13}=10\cdot2010\ne0$，矛盾。

\end{solve}

\begin{remark}[方法与条件]

阻碍来自幂零链的长度：三维空间中的幂零变换至多有三步，而平方会缩短这条链。本题不需要逐项求解矩阵方程，也不需要列举所有 Jordan 形。

\end{remark}

\subsection{初赛第 3 届（2011 年度），来源序号6}\label{ans:pre03-math:6}\label{pr:pre03-math:6}

对应例~\ref{ex:pre03-math:6}（第~\pageref{ex:pre03-math:6} 页）；\href{https://github.com/lzxcr/cmc-mathematics/blob/64916fa1fdbc3e444a652eca204d5669058ee321/src/problems/preliminary/03/06.md}{查看来源}。

\begin{solve}

连续迭代最终消失的向量组成幂零部分，稳定保留下来的像组成可逆部分。令
\[
K=\ker A^n,\qquad W=\operatorname{im}A^n.
\]
核的递增链在至多 $n$ 步后稳定，所以 $\ker A^{2n}=\ker A^n$。若 $v=A^nu\in K\cap W$，则 $A^{2n}u=0$，从而 $A^nu=0$，即 $v=0$。再由秩零化度定理，
\[
\mathbb F^n=W\oplus K.
\]
两个子空间均在 $A$ 下不变。$A|_K$ 的 $n$ 次幂为零；像链同样稳定，故 $A(W)=W$，于是 $A|_W$ 可逆。分别选取 $W,K$ 的基并合并，就得到所需的分块矩阵。

\end{solve}

\begin{remark}[方法与条件]

这是 Fitting 分解。它只使用有限维性，对任意域都成立，并把零特征值对应的部分与其余部分分开；无需假定特征多项式在原域内分裂。见\hyperref[app:linear-invariant]{方法卷：不变子空间与分解}。

\end{remark}

\subsection{初赛第 8 届（2016 年度），来源序号2}\label{ans:pre08-math:2}\label{pr:pre08-math:2}

对应例~\ref{ex:pre08-math:2}（第~\pageref{ex:pre08-math:2} 页）；\href{https://github.com/lzxcr/cmc-mathematics/blob/64916fa1fdbc3e444a652eca204d5669058ee321/src/problems/preliminary/08/02.md}{查看来源}。

\begin{solve}

由 $\operatorname{im}A\subseteq\ker B$，有 $\operatorname{rank}A+\operatorname{rank}B\le n$。因 $n$ 为奇数，至少一个秩不超过 $(n-1)/2$。不妨设它是 $A$，则相似保持秩，故
\[
\operatorname{rank}(A+J_A)
\le\operatorname{rank}A+\operatorname{rank}J_A
=2\operatorname{rank}A<n.
\]
所以 $A+J_A$ 奇异，$0\in S_1$；另一情形同理。

\end{solve}

\begin{remark}[方法与条件]

证明只用到标准形与原矩阵同秩；可把 $J_A,J_B$ 换成任何分别与 $A,B$ 相似的矩阵。即使 Jordan 形式需要在复数域内取，实矩阵的秩在扩域后也不变。

\end{remark}

\subsection{初赛第 14 届补赛（2022 年度），来源序号4}\label{ans:ch05-14th-A-makeup:4}\label{pr:ch05-14th-A-makeup:4}

对应例~\ref{ex:ch05-14th-A-makeup:4}（第~\pageref{ex:ch05-14th-A-makeup:4} 页）；\href{https://github.com/lzxcr/cmc-mathematics/blob/64916fa1fdbc3e444a652eca204d5669058ee321/src/problems/preliminary/14-makeup/04.md}{查看来源}。

\begin{solve}

交换子的迹为零；先乘以 $A$ 的各次幂，仍能用迹的循环性消去交换子，因而得到一整列谱的幂和条件。对任意 $j\geq0$，迹的循环性给出
\[
\operatorname{tr}A^{m+j}
=\operatorname{tr}\bigl(A^j(AB-BA)\bigr)=0.
\]
设 $A$ 的互异非零特征值为 $\lambda_1,\ldots,\lambda_s$，其代数重数为 $d_1,\ldots,d_s$。取 $j=0,\ldots,s-1$，得到
\[
\sum_{i=1}^s(d_i\lambda_i^m)\lambda_i^j=0.
\]
Vandermonde 矩阵可逆，故 $d_i\lambda_i^m=0$。数域的特征为零，与 $d_i>0$、$\lambda_i\neq0$ 矛盾。因此 $A$ 全部特征值为零，即 $A$ 幂零；$I+A^k$ 的全部特征值均为 $1$，故行列式为 $1$。

\end{solve}

\begin{remark}[方法与条件]

更一般地，若 $[A,B]=p(A)$，同样的迹与 Vandermonde 论证说明每个特征值 $\lambda$ 都满足 $p(\lambda)=0$，所以 $p(A)$ 幂零；只有当 $p$ 的唯一根是 $0$ 时，才能据此推出 $A$ 幂零。Jacobson 引理的正确表述也是：若 $[A,B]$ 与 $A$ 可交换，则交换子 $[A,B]$ 幂零，而非 $A$ 必幂零。

\end{remark}

\subsection{决赛第 9 届（2018 年），来源序号6，共同题}\label{ans:fin09-hi:6}\label{pr:fin09-hi:6}

对应例~\ref{ex:fin09-hi:6}（第~\pageref{ex:fin09-hi:6} 页）；\href{https://github.com/lzxcr/cmc-mathematics/blob/64916fa1fdbc3e444a652eca204d5669058ee321/src/problems/finals/09/06.md}{查看来源}。

\begin{solve}

（1）幂零性可由所有特征值的幂和判定；这里交换关系使每个幂的迹都能写成交换子的迹。对任意 $k\geq1$，由 $BC=CB$ 及迹的循环性，
\[
\operatorname{tr}(C^k)=\operatorname{tr}((AB-BA)C^{k-1})=0.
\]
特别地，前 $n$ 个幂和都为零，Newton 恒等式递推给出特征多项式的全部非首项系数为零，即 $\det(tI-C)=t^n$，故由 Cayley--Hamilton 定理，$C^n=0$。

（2）非零空间 $W=\ker C$ 对 $A,B$ 不变，而限制在 $W$ 上的 $A,B$ 可交换。取 $A|_W$ 的一个特征空间，它仍对 $B$ 不变，再在其中取 $B$ 的特征向量，便得到 $A,B,C$ 的共同特征直线。在商空间中三个关系仍成立，对维数归纳得到共同不变完全旗标，即同时上三角化。

（3）最小值为 $3$，由 $A=E_{12},B=E_{23},C=E_{13}$ 实现。$n=1$ 时交换子为零；$n=2$ 时非零幂零的 $C$ 相似于 $J_2(0)$，直接解 $XC=CX$ 可得其交换代数只有 $aI+bC$，是交换的，因此 $A,B$ 必可交换，矛盾。

\end{solve}

\begin{remark}[方法与条件]

上三角化不是把三个矩阵各自上三角化，而是构造它们共同保持的一列子空间。本题先在 $\ker C$ 中找到共同直线，再转到商空间，正好完成这一归纳。这个结构见\hyperref[app:linear-invariant]{方法卷：不变子空间}。

\end{remark}

\subsection{决赛第 10 届（2019 年），来源序号6，共同题}\label{ans:fin10-hi:6}\label{pr:fin10-hi:6}

对应例~\ref{ex:fin10-hi:6}（第~\pageref{ex:fin10-hi:6} 页）；\href{https://github.com/lzxcr/cmc-mathematics/blob/64916fa1fdbc3e444a652eca204d5669058ee321/src/problems/finals/10/06.md}{查看来源}。

\begin{solve}

令 $a=\operatorname{tr}A/n$，只须把迹零矩阵 $B=A-aI$ 写成两个幂零矩阵之和。

先证迹零矩阵可相似为零对角矩阵。对维数归纳：若 $B$ 为标量阵，则 $B=0$；否则存在 $v$ 使 $v,Bv$ 线性无关（若每个向量都是特征向量，则比较 $u,v,u+v$ 可知 $B$ 为标量阵）。以 $v,Bv$ 开始扩充为基，得到
\[
B\sim\begin{pmatrix}0&*\\ *&C\end{pmatrix},\qquad \operatorname{tr}C=0.
\]
对 $C$ 应用归纳假设，并把换基扩充为 $\operatorname{diag}(1,Q)$，即得到整个矩阵零对角。

写 $P^{-1}BP=L+U$，其中 $L,U$ 分别为严格下、上三角部分，故 $L^n=U^n=0$。于是
\[
A=aI+PLP^{-1}+PUP^{-1}
\]
就是所求分解。

\end{solve}

\begin{remark}[方法与条件]

迹先唯一确定标量部分，零对角化再把剩余信息分给严格上、下三角部分。两部分分别幂零，却通常不交换；若额外要求它们交换，其和也会幂零，便无法表示一般迹零矩阵。换基与分块思路见\hyperref[app:linear-invariant]{方法卷：不变子空间}。

\end{remark}

\subsection{决赛第 12 届（2021 年），来源序号3，共同题}\label{ans:fin12-high:3}\label{pr:fin12-high:3}

对应例~\ref{ex:fin12-high:3}（第~\pageref{ex:fin12-high:3} 页）；\href{https://github.com/lzxcr/cmc-mathematics/blob/64916fa1fdbc3e444a652eca204d5669058ee321/src/problems/finals/12/03.md}{查看来源}。

\begin{solve}

答案为 $J_3(1)\oplus J_2(1)$。减去单位阵后，两条 Jordan 链分别是 $e_5\mapsto e_3\mapsto e_1\mapsto0$ 和 $e_4\mapsto e_2\mapsto0$。

\end{solve}

\subsection{决赛第 14 届（2023 年），来源序号3，共同题}\label{ans:fin14-high:3}\label{pr:fin14-high:3}

对应例~\ref{ex:fin14-high:3}（第~\pageref{ex:fin14-high:3} 页）；\href{https://github.com/lzxcr/cmc-mathematics/blob/64916fa1fdbc3e444a652eca204d5669058ee321/src/problems/finals/14/03.md}{查看来源}。

\begin{solve}

取 $v$ 使 $A^{n-1}v\ne0$。向量 $v,Av,\dots,A^{n-1}v$ 线性无关：若 $\sum_{j=0}^{n-1}c_jA^jv=0$，取最小的 $j$ 使 $c_j\ne0$，再作用 $A^{n-1-j}$，即得 $c_jA^{n-1}v=0$，矛盾。因此它们是一组基，其像由后 $n-1$ 个基向量生成，故 $\boxed{\operatorname{rank}A=n-1}$。

\end{solve}

\subsection{决赛第 15 届（2024 年），来源序号8，低年级组}\label{ans:fin15-low:8}\label{pr:fin15-low:8}

对应例~\ref{ex:fin15-low:8}（第~\pageref{ex:fin15-low:8} 页）；\href{https://github.com/lzxcr/cmc-mathematics/blob/64916fa1fdbc3e444a652eca204d5669058ee321/src/problems/finals/15/08.md}{查看来源}。

\begin{solve}

设 $A$ 有 $s\le n$ 个不同特征值。对任一特征值 $\lambda\ne0$，$\lambda,\lambda^2,\dots,\lambda^{s+1}$ 都属于该集合，因为 $s+1\le n!$。由抽屉原理，$\lambda$ 是阶数不超过 $s$ 的单位根，因而 $\lambda^{n!}=1$。

于是 $A^{n!}$ 的全部特征值为 $1$，再由 $A^{n!}\sim A$ 可知 $A=I+N$，其中 $N$ 幂零。任取正整数 $k$，
\[
A^k-I=Nq_k(N),\qquad q_k(0)=k\ne0.
\]
矩阵 $q_k(N)$ 可逆且与 $N$ 交换，故对所有 $j\ge1$，
$\ker(A^k-I)^j=\ker N^j$。对幂零算子 $T$，差 $\dim\ker T^j-\dim\ker T^{j-1}$ 恰为大小至少 $j$ 的 Jordan 块个数。因此各次幂核维数唯一确定全部块大小，因此 $A^k-I\sim N$，即 $A^k\sim A$。全部特征值为实数，所用 Jordan 基也可在实数域选取。

这一方法把相似问题转成核的滤过，见\hyperref[app:linear-cyclic]{方法卷：循环结构与幂零块}。

\end{solve}

\section{二次型、合同与正定性}

\label{sec:solutions-7}

二次型只由矩阵的对称部分决定。相似、合同和正交相似对应不同的问题；正定矩阵不等式通常先化为单位矩阵与一个正定矩阵的比较，再处理特征值或行列式。

\subsection{初赛第 2 届（2010 年度），来源序号6}\label{ans:pre02-math:6}\label{pr:pre02-math:6}

对应例~\ref{ex:pre02-math:6}（第~\pageref{ex:pre02-math:6} 页）；\href{https://github.com/lzxcr/cmc-mathematics/blob/64916fa1fdbc3e444a652eca204d5669058ee321/src/problems/preliminary/02/06.md}{查看来源}。

\begin{solve}

二次型只看到矩阵的对称部分。令
\[
S=\frac{A+A^T}{2}.
\]
则 $S$ 半正定，且 $\beta S\beta^T=0$。由谱分解，或由 $\beta S\beta^T=\|S^{1/2}\beta^T\|^2$，得到 $S\beta^T=0$。因此对任意 $v$，
\[
vA\beta^T+\beta Av^T=2vS\beta^T=0.
\]
题设蕴含关系的逆否命题随即给出 $vA\beta^T=0$。

\end{solve}

\begin{remark}[方法与条件]

半正定二次型在某向量上取零，意味着该向量属于其零空间；这是半正定性比一般非负数值条件更强的地方。附加假设负责把对称部分的消失传递到原双线性型。

\end{remark}

\subsection{初赛第 4 届（2012 年度），来源序号4}\label{ans:pre04-math:4}\label{pr:pre04-math:4}

对应例~\ref{ex:pre04-math:4}（第~\pageref{ex:pre04-math:4} 页）；\href{https://github.com/lzxcr/cmc-mathematics/blob/64916fa1fdbc3e444a652eca204d5669058ee321/src/problems/preliminary/04/04.md}{查看来源}。

\begin{solve}

取 $v\ne0$ 满足 $P(\lambda)v=0$，左乘 $v^*$，得到标量关系
\[
a\lambda^2+b\lambda+c=0,
\qquad a=v^*Av>0,\quad b=v^*Bv>0,\quad c=v^*Cv>0.
\]
由于 $c>0$，必有 $\lambda\ne0$。除以 $\lambda$ 后取实部，得
\[
\left(a+\frac{c}{|\lambda|^2}\right)\operatorname{Re}\lambda+b=0.
\]
所以
\[
\operatorname{Re}\lambda
=-\frac{b}{a+c/|\lambda|^2}<0.
\]

\end{solve}

\begin{remark}[方法与条件]

正定性允许沿核向量把矩阵多项式压缩成具有正系数的标量关系。这也解释了阻尼系统 $A\ddot u+B\dot u+Cu=0$ 的稳定性：所有指数模态 $e^{\lambda t}v$ 都衰减，且完全不要求 $A,B,C$ 两两交换。

\end{remark}

\subsection{初赛第 9 届（2017 年度），来源序号4}\label{ans:pre09-math:4}\label{pr:pre09-math:4}

对应例~\ref{ex:pre09-math:4}（第~\pageref{ex:pre09-math:4} 页）；\href{https://github.com/lzxcr/cmc-mathematics/blob/64916fa1fdbc3e444a652eca204d5669058ee321/src/problems/preliminary/09/04.md}{查看来源}。

\begin{solve}

记 $C=aI+A$。若 $Cv=0$，则 $0=v^{\mathrm T}Cv=a\|v\|^2$，所以 $v=0$；因此 $C,X,Y,M,N$ 均可逆。

若 $XN+Y^{\mathrm T}M^{\mathrm T}=0$，令 $P=XM^{-1}$，利用 $N=M^{-1}C$、$Y=X^{-1}C$，得到
\[
PC+C^{\mathrm T}P^{-\mathrm T}=0,
\qquad PCP^{\mathrm T}+C^{\mathrm T}=0.
\]
将后一式与其转置相加，并用 $C+C^{\mathrm T}=2aI$，得到 $2a(PP^{\mathrm T}+I)=0$，这与 $PP^{\mathrm T}+I$ 正定且 $a\ne0$ 矛盾。

\end{solve}

\begin{remark}[方法与条件]

反对称部分在取对称部分时消失，留下的是正定矩阵的合同结构。一般地，若 $C+C^{\mathrm T}$ 正定或负定，同样的论证仍适用。

\end{remark}

\subsection{初赛第 10 届（2018 年度），来源序号2}\label{ans:pre10-math:2}\label{pr:pre10-math:2}

对应例~\ref{ex:pre10-math:2}（第~\pageref{ex:pre10-math:2} 页）；\href{https://github.com/lzxcr/cmc-mathematics/blob/64916fa1fdbc3e444a652eca204d5669058ee321/src/problems/preliminary/10/02.md}{查看来源}。

\begin{solve}

二次型只依赖对称部分 $B=(A+A^{\mathrm T})/2$。其对角元为 $a$，且
\[
\sum_{j\ne i}|b_{ij}|
\le\frac12\sum_{j\ne i}(|a_{ij}|+|a_{ji}|)<a.
\]
用 $2|x_ix_j|\le x_i^2+x_j^2$，有
\[
x^{\mathrm T}Bx\ge\sum_i\left(a-\sum_{j\ne i}|b_{ij}|\right)x_i^2>0
\quad(x\ne0).
\]
故二次型正定，其规范形为 $\boxed{y_1^2+\cdots+y_n^2}$。

\end{solve}

\begin{remark}[方法与条件]

先取对称部分，再以严格对角占优控制交叉项。这一证明直接给出正定性，无须计算特征值；也说明题设中的行和、列和为何成对出现。

\end{remark}

\subsection{初赛第 11 届（2019 年度），来源序号4}\label{ans:ch10-11th-A:4}\label{pr:ch10-11th-A:4}

对应例~\ref{ex:ch10-11th-A:4}（第~\pageref{ex:ch10-11th-A:4} 页）；\href{https://github.com/lzxcr/cmc-mathematics/blob/64916fa1fdbc3e444a652eca204d5669058ee321/src/problems/preliminary/11/04.md}{查看来源}。

\begin{solve}

正定或负定二次型在非零向量处都不能为零。因此可调节参数，让 $kf+f_1$ 出现一个非零零点，以迫使整个二次型消失。显然 $\mathbb R f_1\subset V$。反之，取非零 $f\in V$，选向量 $v$ 使 $f(v)\neq0$，并令
\[
k=-\frac{f_1(v)}{f(v)}.
\]
二次型 $kf+f_1$ 在非零向量 $v$ 处取零值，因此既不正定也不负定。按题设的“恒号”定义，它只能是零二次型。又 $f_1(v)>0$，所以 $k\neq0$，得到 $f=-f_1/k$。

因此 $V=\mathbb R f_1$，确为实向量空间，且维数为 $1$。

\end{solve}

\subsection{决赛第 1 届（2010 年），来源序号10}\label{ans:fin01-all:10}\label{pr:fin01-all:10}

对应例~\ref{ex:fin01-all:10}（第~\pageref{ex:fin01-all:10} 页）；\href{https://github.com/lzxcr/cmc-mathematics/blob/64916fa1fdbc3e444a652eca204d5669058ee321/src/problems/finals/01/10.md}{查看来源}。

\begin{solve}

事实上秩的限制可以去掉。令
\[
K=\ker(A+B)=\ker A\cap\ker B,
\]
其中等号来自半正定性。在 $W=K^\perp$ 上，$H=A+B$ 正定。用 $H^{-1/2}$ 作合同变换后，两矩阵变为
\[
\widetilde A=H^{-1/2}AH^{-1/2},\qquad
\widetilde B=I-\widetilde A.
\]
二者由同一个正交变换对角化。把该变换连同 $K$ 上的任意一组基合并，就得到所需的 $C$。

\end{solve}

\begin{remark}[方法与条件]

两个半正定二次型总能同时合同对角化。关键是把它们的和作为内积；归一化以后，一个二次型就是另一个的补。此处的 $H^{-1/2}$ 仅作用在 $W$ 上，因而不要求 $A+B$ 在整个空间可逆。

\end{remark}

\subsection{决赛第 2 届（2011 年），来源序号1}\label{ans:fin02-all:1}\label{pr:fin02-all:1}

对应例~\ref{ex:fin02-all:1}（第~\pageref{ex:fin02-all:1} 页）；\href{https://github.com/lzxcr/cmc-mathematics/blob/64916fa1fdbc3e444a652eca204d5669058ee321/src/problems/finals/02/01.md}{查看来源}。

\begin{solve}

设所求平面为二维子空间 $W$。截面关于原点中心对称，故圆心为原点；于是二次型
$q=4x^2+5y^2+6z^2$ 在 $W$ 上等于 $\lambda(x^2+y^2+z^2)$。

平面 $W$ 与 $xy$ 平面有非零交向量，故 $\lambda\le5$；与 $yz$ 平面也有非零交向量，故 $\lambda\ge5$。所以 $\lambda=5$，且 $W$ 上恒有 $z^2-x^2=0$。一个平面若包含于两个平面的并，必为其中之一。因此恰有
\[
x=z\qquad\text{或}\qquad x=-z.
\]
二者的截面半径均为 $1/\sqrt5$。

\end{solve}

\begin{remark}[方法与条件]

圆截面意味着二次型在一个二维子空间上成为标量形式。夹住这个标量的是中间特征值，这比直接设一般平面再消元更清楚。

\end{remark}

\subsection{决赛第 4 届（2013 年），来源序号5}\label{ans:fin04-all:5}\label{pr:fin04-all:5}

对应例~\ref{ex:fin04-all:5}（第~\pageref{ex:fin04-all:5} 页）；\href{https://github.com/lzxcr/cmc-mathematics/blob/64916fa1fdbc3e444a652eca204d5669058ee321/src/problems/finals/04/05.md}{查看来源}。

\begin{solve}

记 $u=(x_2,x_3,x_4)^{\mathsf T}$。分块行列式公式给出
\[
f=-12x_1^2+u^{\mathsf T}A^*u.
\]
因为 $A^*=-12A^{-1}$，向量 $v=(1,0,-2)^{\mathsf T}$ 是 $A$ 对应特征值 $-3$ 的特征向量。其余两个特征值之和为 $4$、乘积为 $4$，故都为 $2$。于是 $A^*$ 在 $v^\perp$ 上为 $-6I$，在 $\mathbb Rv$ 上为 $4I$。

选取 $v^\perp$ 的标准正交基 $(0,1,0)^{\mathsf T}$、$(2,0,1)^{\mathsf T}/\sqrt5$，得到所求正交变换
\[
x_1=y_1,\qquad x_2=\frac{2y_3+y_4}{\sqrt5},\qquad
x_3=y_2,\qquad x_4=\frac{y_3-2y_4}{\sqrt5}.
\]
标准型为
\[
f=-12y_1^2-6y_2^2-6y_3^2+4y_4^2.
\]

\end{solve}

\begin{remark}[方法与条件]

无需先恢复 $A$ 的各个元素：伴随矩阵条件给出一个特征方向，迹与行列式确定剩余谱，对称性再确定正交补上的作用。

\end{remark}

\subsection{决赛第 6 届（2015 年），来源序号1，共同题}\label{ans:fin06-lo:1}\label{pr:fin06-lo:1}

对应例~\ref{ex:fin06-lo:1}（第~\pageref{ex:fin06-lo:1} 页）；\href{https://github.com/lzxcr/cmc-mathematics/blob/64916fa1fdbc3e444a652eca204d5669058ee321/src/problems/finals/06/01.md}{查看来源}。

\begin{solve}

答案为 $z_1^2+z_2^2-z_3^2$。二次型的对称矩阵为
\[
A=\begin{pmatrix}0&1&-1/2\\1&0&5/2\\-1/2&5/2&0\end{pmatrix},
\qquad \det A=-\frac52<0,\qquad\operatorname{tr}A=0.
\]
三个特征值均为非零实数，乘积为负，故负特征值有一个或三个；迹为零排除了三个全负的可能。因此正、负惯性指数分别为 $2,1$，由惯性定理得到所求规范型。

\end{solve}

\subsection{决赛第 9 届（2018 年），来源序号4，共同题}\label{ans:fin09-hi:4}\label{pr:fin09-hi:4}

对应例~\ref{ex:fin09-hi:4}（第~\pageref{ex:fin09-hi:4} 页）；\href{https://github.com/lzxcr/cmc-mathematics/blob/64916fa1fdbc3e444a652eca204d5669058ee321/src/problems/finals/09/04.md}{查看来源}。

\begin{solve}

答案为
\[
[1+(n-1)a]u_1^2+(1-a)\sum_{j=2}^nu_j^2.
\]
因为 $A=(1-a)I+a\mathbf1\mathbf1^T$，在 $\operatorname{span}(\mathbf1)$ 及其正交补上分别作标量 $1+(n-1)a$、$1-a$ 作用。

\end{solve}

\subsection{决赛第 10 届（2019 年），来源序号1，共同题}\label{ans:fin10-hi:1}\label{pr:fin10-hi:1}

对应例~\ref{ex:fin10-hi:1}（第~\pageref{ex:fin10-hi:1} 页）；\href{https://github.com/lzxcr/cmc-mathematics/blob/64916fa1fdbc3e444a652eca204d5669058ee321/src/problems/finals/10/01.md}{查看来源}。

\begin{solve}

答案为
\[
\begin{pmatrix}1&0&1\\0&-4&2\\1&2&0\end{pmatrix}\quad\text{或}\quad
\begin{pmatrix}3/2&1&1\\1&2&0\\1&0&1\end{pmatrix}.
\]
所给两个向量是矩阵的实际行向量。若它们分别位于第 $i,j$ 行，对称性要求第一个向量的第 $j$ 个分量等于第二个向量的第 $i$ 个分量。检查 $i\ne j$，只有 $(i,j)=(1,3)$ 或 $(3,2)$；剩余对角元再由秩为 $2$（行列式为零）唯一确定。

\end{solve}

\subsection{决赛第 11 届（2021 年），来源序号3，共同题}\label{ans:fin11-high:3}\label{pr:fin11-high:3}

对应例~\ref{ex:fin11-high:3}（第~\pageref{ex:fin11-high:3} 页）；\href{https://github.com/lzxcr/cmc-mathematics/blob/64916fa1fdbc3e444a652eca204d5669058ee321/src/problems/finals/11/03.md}{查看来源}。

\begin{solve}

答案为椭圆柱面。原二次式等于 $(x+y)^2+(x-z)^2$，故非负且秩为 $2$；核为直线 $x=-y=z$。因此沿该直线方向不受限制，垂直于它的截面是椭圆。正交标准式为 $3u^2+v^2=1$。

\end{solve}

\subsection{决赛第 14 届（2023 年），来源序号5，共同题}\label{ans:fin14-high:5}\label{pr:fin14-high:5}

对应例~\ref{ex:fin14-high:5}（第~\pageref{ex:fin14-high:5} 页）；\href{https://github.com/lzxcr/cmc-mathematics/blob/64916fa1fdbc3e444a652eca204d5669058ee321/src/problems/finals/14/05.md}{查看来源}。

\begin{solve}

令 $M=\begin{pmatrix}A&B/2\\B/2&C\end{pmatrix}$。变换后二次项矩阵为 $M'=Q^{\mathrm T}MQ$，所以
\[
B'^2-4A'C'=-4\det M'=-4\det M=B^2-4AC.
\]
线性变换保持原点，故常数项仍为 $F$。

交叉项消失当且仅当 $Q^{\mathrm T}MQ$ 对角化，即 $Q$ 的两列是一组 $M$ 的正交特征向量。具体地，写
\[
Q=\begin{pmatrix}\cos\theta&-\varepsilon\sin\theta\\
\sin\theta&\varepsilon\cos\theta\end{pmatrix},\quad\varepsilon=\pm1,
\]
则充要条件为 $\boxed{(C-A)\sin2\theta+B\cos2\theta=0}$，该形式也覆盖 $A=C$ 等退化情形。

正交变换保留度量，并将对称二次型沿特征方向分解；这也是消除交叉项的结构依据，见\href{https://github.com/lzxcr/cmc-mathematics/blob/64916fa1fdbc3e444a652eca204d5669058ee321/src/methods/outer/02-geometry.md}{方法卷：二次型与曲面（原稿链接）}。

\end{solve}

\subsection{决赛第 15 届（2024 年），来源序号1，共同题}\label{ans:fin15-high:1}\label{pr:fin15-high:1}

对应例~\ref{ex:fin15-high:1}（第~\pageref{ex:fin15-high:1} 页）；\href{https://github.com/lzxcr/cmc-mathematics/blob/64916fa1fdbc3e444a652eca204d5669058ee321/src/problems/finals/15/01.md}{查看来源}。

\begin{solve}

Sylvester 惯性定律表明合同类由正、负、零惯性指数 $(p,q,r)$ 唯一确定，其中 $p+q+r=24$，三者均非负。因此类数为 $\boxed{\binom{26}{2}=325}$。

\end{solve}

\subsection{决赛第 16 届（2025 年），来源序号1，共同题}\label{ans:fin16-high:1}\label{pr:fin16-high:1}

对应例~\ref{ex:fin16-high:1}（第~\pageref{ex:fin16-high:1} 页）；\href{https://github.com/lzxcr/cmc-mathematics/blob/64916fa1fdbc3e444a652eca204d5669058ee321/src/problems/finals/16/01.md}{查看来源}。

\begin{solve}

由实对称性可正交对角化；幂等性使特征值只取 $0,1$，秩给出 $1$ 的个数。因此 $A$ 正交相似于 $\operatorname{diag}(I_{1012},0_{1013})$，故规范型为 $\boxed{z_1^2+\cdots+z_{1012}^2}$。

\end{solve}

\subsection{决赛第 17 届（2026 年），来源序号6，共同题}\label{ans:fin17-high:6}\label{pr:fin17-high:6}

对应例~\ref{ex:fin17-high:6}（第~\pageref{ex:fin17-high:6} 页）；\href{https://github.com/lzxcr/cmc-mathematics/blob/64916fa1fdbc3e444a652eca204d5669058ee321/src/problems/finals/17/06.md}{查看来源}。

\begin{solve}

（1）置 $R=(I-BB^{\mathrm T})^{-1}$。展开左端，得到
\[
R-I-AB^{\mathrm T}R-RBA^{\mathrm T}
+A(I+B^{\mathrm T}RB)A^{\mathrm T},
\]
即为所需恒等式。

（2）目标是把左端整理成两个正定矩阵之和；再取行列式，展开中的混合项会给出严格增益。利用
\[
I+B^{\mathrm T}RB=(I-B^{\mathrm T}B)^{-1},\qquad
RB=B(I-B^{\mathrm T}B)^{-1},
\]
可将（1）进一步整理为
\[
(I-AB^{\mathrm T})(BB^{\mathrm T}-I)^{-1}(I-BA^{\mathrm T})=P+H,
\]
其中
\[
P=AA^{\mathrm T}-I>0,\qquad
H=(A-B)(B^{\mathrm T}B-I)^{-1}(A-B)^{\mathrm T}>0.
\]
这里 $BB^{\mathrm T}-I>0$ 保证 $B^{\mathrm T}B-I>0$，而 $\det(A-B)=1$ 保证 $H$ 正定。

对任意 $n\ge2$ 阶正定矩阵 $P,H$，设 $P^{-1/2}HP^{-1/2}$ 的特征值为 $\lambda_i>0$，则
\[
\det(P+H)=\det P\prod_i(1+\lambda_i)
>\det P\left(1+\prod_i\lambda_i\right)=\det P+\det H.
\]
另一方面 $\det H=1/\det(BB^{\mathrm T}-I)$。对上述矩阵恒等式取行列式并应用此严格不等式，得
\[
\det(I-AB^{\mathrm T})^2>
\det(AA^{\mathrm T}-I)\det(BB^{\mathrm T}-I)+1.
\]
最后两因子同时换号产生 $(-1)^{2n}=1$，正好得到题设结论。

\end{solve}

\begin{remark}[方法与条件]

证明实际给出更一般的不等式
$\det(I-AB^{\mathrm T})^2>\det(I-AA^{\mathrm T})\det(I-BB^{\mathrm T})+\det(A-B)^2$，只需 $A-B$ 可逆及两正定性条件。严格性来自正定矩阵行列式展开中的中间项，因而 $n\ge2$ 不可省略。

\end{remark}

\section{内积、正交投影与奇异值}

\label{sec:solutions-8}

内积与正交结构能把代数约束转化为长度、角度和能量的约束。非正交直和中的正交投影一般不是互补投影；奇异值应由伴随乘积确定，不能直接当作特征值的绝对值。

\subsection{初赛第 9 届（2017 年度），来源序号3}\label{ans:pre09-math:3}\label{pr:pre09-math:3}

对应例~\ref{ex:pre09-math:3}（第~\pageref{ex:pre09-math:3} 页）；\href{https://github.com/lzxcr/cmc-mathematics/blob/64916fa1fdbc3e444a652eca204d5669058ee321/src/problems/preliminary/09/03.md}{查看来源}。

\begin{solve}

记 $S=\sum_iW_i$。对任意 $W_j$，左乘 $W_j$ 是有限集合 $\Gamma$ 上的单射，故是置换，所以 $W_jS=S$。将其对 $j$ 求和得 $S^2=rS$，即 $P=S/r$ 是幂等矩阵。

幂等矩阵的最小多项式整除 $t(t-1)$，因而特征值只可能为 $0,1$，且可对角化，所以 $\operatorname{tr}P=\operatorname{rank}P$。因此 $\operatorname{tr}S=0$ 当且仅当 $P=0$，也即 $S=0$；另一个方向亦由迹的线性性直接成立。

\end{solve}

\begin{remark}[方法与条件]

$P$ 是有限群作用的平均投影，其像恰为被全部 $W_i$ 固定的向量空间。因此 $\frac1r\sum_i\operatorname{tr}W_i$ 实际是固定子空间的维数，必为非负整数。

\end{remark}

\subsection{初赛第 12 届（2020 年度），来源序号3}\label{ans:ch09-12th-A:3}\label{pr:ch09-12th-A:3}

对应例~\ref{ex:ch09-12th-A:3}（第~\pageref{ex:ch09-12th-A:3} 页）；\href{https://github.com/lzxcr/cmc-mathematics/blob/64916fa1fdbc3e444a652eca204d5669058ee321/src/problems/preliminary/12/03.md}{查看来源}。

\begin{solve}

不可能。解空间维数控制的是相对变换 $B^{-1}A$ 的不动向量；正交性再把这条信息转成行列式的符号。令 $C=B^{-1}A$，则 $C$ 为实正交矩阵，且
\[
\dim\ker(C-I)=\dim\ker(A-B)=3.
\]
实正交矩阵可复对角化，非实特征值成共轭对，每对之积为 $1$；实特征值只有 $\pm1$。$1$ 的代数重数为 $3$，故其余 $2017$ 个特征值中，$-1$ 的重数为奇数。于是
\[
\det C=-1,\qquad \det A=-\det B.
\]
相似矩阵的行列式相同，故 $A,B$ 不相似。

\end{solve}

\begin{remark}[方法与条件]

一般实正交矩阵满足
\[
\det C=(-1)^{n-\dim\ker(C-I)}.
\]
因此当 $\operatorname{rank}(A-B)$ 为奇数时，$\det A$ 与 $\det B$ 相反。这个奇偶结论依赖“实”条件，不能照搬到一般酉矩阵，因为它们的非实特征值不必成共轭对。

正交算子的谱分解及其实结构，见\hyperref[app:linear-spectral]{方法卷：谱与函数演算}。

\end{remark}

\subsection{初赛第 13 届（2021 年度），来源序号4}\label{ans:ch07-13th-A:4}\label{pr:ch07-13th-A:4}

对应例~\ref{ex:ch07-13th-A:4}（第~\pageref{ex:ch07-13th-A:4} 页）；\href{https://github.com/lzxcr/cmc-mathematics/blob/64916fa1fdbc3e444a652eca204d5669058ee321/src/problems/preliminary/13/04.md}{查看来源}。

\begin{solve}

对称性与酉性共同给出 $P\bar P=I$。因此反线性映射
\[
Jv=P\bar v
\]
是一个等距对合：$J^2=I$，且 $\langle Ju,Jv\rangle=\overline{\langle u,v\rangle}$。我们寻找一个使 $J$ 成为普通共轭的基。

记其实不动空间为 $W=\{v:Jv=v\}$。任意 $v\in\mathbb C^n$ 都有唯一分解
\[
v=\frac{v+Jv}{2}+i\frac{v-Jv}{2i},
\]
右侧两个向量均在 $W$ 中，且 $W\cap iW=\{0\}$。所以 $\mathbb C^n=W\oplus iW$（作为实空间），$\dim_{\mathbb R}W=n$。又 $W$ 中向量的内积为实数，可在其中作实 Gram--Schmidt 正交化，得到正交归一基 $u_1,\ldots,u_n$。

令 $U=(u_1,\ldots,u_n)$，则 $U$ 酉，且 $P\bar U=U$。取 $Q=\bar U$，便有
\[
\bar Q Q^{-1}=U\bar U^{-1}=P.
\]
这同时说明 $Q$ 可以取为酉矩阵。

\end{solve}

\begin{remark}[方法与条件]

映射 $J$ 的不动空间提供了 $\mathbb C^n$ 的一套“实坐标”：每个向量唯一写成 $u+iv$，其中 $u,v\in W$。因此结论是在说，这种实结构经酉换基后就是通常的复共轭。上面的 $P=U\bar U^{-1}=UU^{\mathrm T}$ 也是对称酉矩阵的 Takagi 分解；这里已经直接构造了所需分解，不必预先引用一般定理。

\end{remark}

\subsection{决赛第 2 届（2011 年），来源序号6}\label{ans:fin02-all:6}\label{pr:fin02-all:6}

对应例~\ref{ex:fin02-all:6}（第~\pageref{ex:fin02-all:6} 页）；\href{https://github.com/lzxcr/cmc-mathematics/blob/64916fa1fdbc3e444a652eca204d5669058ee321/src/problems/finals/02/06.md}{查看来源}。

\begin{solve}

记 $E_{ij}$ 为矩阵单位。由
$E_{ii}E_{ij}=E_{ij}$、$E_{ij}E_{ii}=0$（$i\ne j$），得 $\varphi(E_{ij})=0$；由
$E_{ij}E_{ji}=E_{ii}$、$E_{ji}E_{ij}=E_{jj}$，得所有 $\varphi(E_{ii})$ 相等。因此
\[
\varphi(X)=c\operatorname{tr}X\qquad(c\ne0).
\]
若 $\langle X,Y\rangle=0$ 对所有 $Y$ 成立，取 $Y=X^{\mathsf T}$ 即得 $X=0$，证明了非退化性。

量 $\sum_iA_iB_i$ 与对偶基的选择无关：第一组基作变换 $S$ 时，第二组基作逆转置变换，两者在求和中相消。因而可以选取矩阵单位基 $E_{ij}$ 及其对偶基 $c^{-1}E_{ji}$，得到
\[
\sum_{i=1}^{n^2}A_iB_i
=\frac1c\sum_{i,j=1}^n E_{ij}E_{ji}
=\frac nc I_n.
\]

\end{solve}

\begin{remark}[方法与条件]

迹是在矩阵代数上满足循环性的线性泛函的唯一来源（相差一个常数）。对偶基求和看似依赖坐标，实际上来自双线性型本身；先证明其坐标不变性，便可自由选择最方便的基。

\end{remark}

\subsection{决赛第 5 届（2014 年），来源序号12，高年级组}\label{ans:fin05-hi:11}\label{pr:fin05-hi:11}

对应例~\ref{ex:fin05-hi:11}（第~\pageref{ex:fin05-hi:11} 页）；\href{https://github.com/lzxcr/cmc-mathematics/blob/64916fa1fdbc3e444a652eca204d5669058ee321/src/problems/finals/05/12.md}{查看来源}。

\begin{solve}

余弦恒等式
$2\cos\theta\cos(n\theta)=\cos((n+1)\theta)+\cos((n-1)\theta)$
与初值相同，故递推即得第一问。令 $x=\cos\theta$，又有
\[
\langle T_n,T_m\rangle=\int_0^\pi\cos(n\theta)\cos(m\theta)\,d\theta=0
\qquad(n\ne m).
\]

递推表明 $T_n$ 的首项系数为 $2^{n-1}$，且在节点 $x_j=\cos(j\pi/n)$（$0\le j\le n$）上满足 $T_n(x_j)=(-1)^j$。记
\[
w_j=\frac1{\left|\prod_{k\ne j}(x_j-x_k)\right|}>0.
\]
Lagrange 插值公式中比较最高次项，得
\[
1=\sum_{j=0}^n(-1)^jw_jP(x_j)
\le\|P\|_\infty\sum_{j=0}^n w_j.
\]
把 $P$ 换成首一多项式 $2^{1-n}T_n$，知 $\sum_jw_j=2^{n-1}$，故得到所求下界。

若等号成立，由所有 $w_j>0$，必有 $P(x_j)=2^{1-n}(-1)^j$ 对每个 $j$ 成立。这给出 $n+1$ 个插值条件，唯一确定 $P=2^{1-n}T_n$。反向显然成立。

\end{solve}

\begin{remark}[方法与条件]

Chebyshev 多项式一面来自圆周上的 Fourier 正交性，一面解决一致范数下的最优逼近问题。极值点的交替符号决定最优性，正的插值权重进一步给出唯一性，避免了仅靠非严格变号计数可能留下的空隙。

\end{remark}

\subsection{决赛第 6 届（2015 年），来源序号6，共同题}\label{ans:fin06-lo:6}\label{pr:fin06-lo:6}

对应例~\ref{ex:fin06-lo:6}（第~\pageref{ex:fin06-lo:6} 页）；\href{https://github.com/lzxcr/cmc-mathematics/blob/64916fa1fdbc3e444a652eca204d5669058ee321/src/problems/finals/06/06.md}{查看来源}。

\begin{solve}

设 $A\in\Gamma$。直接计算有
\[
AA^*=A^*A=(|z_1|^2+|z_2|^2)I,
\]
故 $A$ 正规。令 $u=\operatorname{Re}z_1$、$r=\sqrt{(\operatorname{Im}z_1)^2+|z_2|^2}$，其特征值为 $\lambda=u+ir$ 与 $\overline\lambda=u-ir$，所以 Jordan 标准型为
\[
J_A=\operatorname{diag}(\lambda,\overline\lambda)\in\Gamma.
\]
当 $r=0$ 时 $A=uI$，取 $P=I$ 即可。以下设 $r>0$。

取 $\lambda$ 的非零特征向量 $v=(x,y)^T$。矩阵的定义给出配对关系
\[
Av=\lambda v
\quad\Longrightarrow\quad
A\begin{pmatrix}-\overline y\\\overline x\end{pmatrix}
=\overline\lambda\begin{pmatrix}-\overline y\\\overline x\end{pmatrix}.
\]
于是
\[
P=\begin{pmatrix}x&-\overline y\\y&\overline x\end{pmatrix}\in\Gamma,
\qquad \det P=|x|^2+|y|^2>0,
\]
且 $AP=PJ_A$，即 $P^{-1}AP=J_A$。交换两个特征值的顺序时，同样从对应的特征向量开始构造即可。

\end{solve}

\begin{remark}[方法与条件]

将 $v$ 单位化后，$P^*P=I$ 且 $\det P=1$，所以本题实际上给出了 $\Gamma$ 内部的酉对角化。$\Gamma$ 是四元数代数的复矩阵模型；正规性保证可对角化，反线性配对 $v\mapsto(-\overline y,\overline x)^T$ 则保证变换矩阵保留这一结构。

\end{remark}

\subsection{决赛第 7 届（2016 年），来源序号6，共同题}\label{ans:fin07-hi:6}\label{pr:fin07-hi:6}

对应例~\ref{ex:fin07-hi:6}（第~\pageref{ex:fin07-hi:6} 页）；\href{https://github.com/lzxcr/cmc-mathematics/blob/64916fa1fdbc3e444a652eca204d5669058ee321/src/problems/finals/07/06.md}{查看来源}。

\begin{solve}

当 $A,B$ 可交换时两边相等，因此考察交换子 $AB-BA$；对称性使它的平方范数恰好给出两边之差。记 $\|M\|_F^2=\operatorname{tr}(M^TM)$，由 $A^T=A,B^T=B$ 及迹的循环性，
\[
\|AB-BA\|_F^2=2\operatorname{tr}(A^2B^2)-2\operatorname{tr}(ABAB)\geq0.
\]
这就证明了不等式，且等号成立当且仅当 $AB=BA$。

\end{solve}

\begin{remark}[方法与条件]

对 Hermite 矩阵，把转置换成共轭转置，同一恒等式仍成立。等号时二者可同时酉对角化：每个 $A$ 的特征空间都被 $B$ 保持，再在其中对角化 $B$ 即可。参见\hyperref[app:linear-spectral]{方法卷：谱分解}。

\end{remark}

\subsection{决赛第 8 届（2017 年），来源序号4，共同题}\label{ans:fin08-hi:4}\label{pr:fin08-hi:4}

对应例~\ref{ex:fin08-hi:4}（第~\pageref{ex:fin08-hi:4} 页）；\href{https://github.com/lzxcr/cmc-mathematics/blob/64916fa1fdbc3e444a652eca204d5669058ee321/src/problems/finals/08/04.md}{查看来源}。

\begin{solve}

答案为 $-1/2$。谱分解给出 $A=I+uu^T$，其中 $|u|=1$，所以 $a_{21}=u_1u_2\geq-1/2$，在 $u=(1,-1)/\sqrt2$ 时取到。

\end{solve}

\subsection{决赛第 10 届（2019 年），来源序号4，共同题}\label{ans:fin10-hi:4}\label{pr:fin10-hi:4}

对应例~\ref{ex:fin10-hi:4}（第~\pageref{ex:fin10-hi:4} 页）；\href{https://github.com/lzxcr/cmc-mathematics/blob/64916fa1fdbc3e444a652eca204d5669058ee321/src/problems/finals/10/04.md}{查看来源}。

\begin{solve}

答案为 $0$。若存在这样的子矩阵，取向量在它的三个列位置上为 $1/\sqrt3$，其他位置为零，则向量长度为 $1$。乘以 $U$ 后，仅在那三个行位置上的长度已为 $3/(2\sqrt2)>1$，与正交变换保持长度矛盾。

\end{solve}

\subsection{决赛第 11 届（2021 年），来源序号4，共同题}\label{ans:fin11-high:4}\label{pr:fin11-high:4}

对应例~\ref{ex:fin11-high:4}（第~\pageref{ex:fin11-high:4} 页）；\href{https://github.com/lzxcr/cmc-mathematics/blob/64916fa1fdbc3e444a652eca204d5669058ee321/src/problems/finals/11/04.md}{查看来源}。

\begin{solve}

答案为 $\Lambda=\operatorname{diag}(3,\sqrt3,0)$，对角元次序可同时置换。因为
\[
AA^T=\begin{pmatrix}5&1&-4\\1&2&1\\-4&1&5\end{pmatrix}
\]
的特征值为 $9,3,0$，奇异值为其非负平方根。

\end{solve}

\subsection{决赛第 11 届（2021 年），来源序号6，共同题}\label{ans:fin11-high:6}\label{pr:fin11-high:6}

对应例~\ref{ex:fin11-high:6}（第~\pageref{ex:fin11-high:6} 页）；\href{https://github.com/lzxcr/cmc-mathematics/blob/64916fa1fdbc3e444a652eca204d5669058ee321/src/problems/finals/11/06.md}{查看来源}。

\begin{solve}

投影和是正定算子，而且其特征值的算术平均恰为 $1$。对 $v\ne0$，
\[
\langle\varphi v,v\rangle=\|p_1v\|^2+\|p_2v\|^2>0,
\]
因为同时投影为零会使 $v\perp(V_1+V_2)=V$。设 $\dim V=d$，则
\[
\operatorname{tr}\varphi=\dim V_1+\dim V_2=d.
\]
对其正特征值用算术---几何平均不等式，得到 $0<\det\varphi\leq1$。等号当且仅当所有特征值为 $1$，即 $\varphi=I$。此时对 $v\in V_1$，$v=p_1v+p_2v=v+p_2v$，故 $p_2v=0$，即 $V_1\perp V_2$；反向显然成立。

\end{solve}

\begin{remark}[方法与条件]

若分别取 $V_1,V_2$ 的正交基，合并成列矩阵 $Q$，则 $\varphi=QQ^T$。因此行列式也是这些单位列向量围成的平行多面体体积的平方；体积为 $1$ 正对应两组基彼此正交。这个几何解释与谱证明分别从 Gram 矩阵和特征值读取同一量，见\hyperref[app:linear-gram]{方法卷：Gram 矩阵}。

\end{remark}

\subsection{决赛第 13 届（2023 年），来源序号8，低年级组}\label{ans:fin13-low:8}\label{pr:fin13-low:8}

对应例~\ref{ex:fin13-low:8}（第~\pageref{ex:fin13-low:8} 页）；\href{https://github.com/lzxcr/cmc-mathematics/blob/64916fa1fdbc3e444a652eca204d5669058ee321/src/problems/finals/13/08.md}{查看来源}。

\begin{solve}

（1）将 $X,Y$ 的两行分别记为 $x_i^{\mathrm T},y_i^{\mathrm T}$，则
\[
p(X,Y)=\sum_{i=1}^2x_i^{\mathrm T}Ay_i.
\]
只让一行非零就得到必要性；反之，两个正定内积之和仍为内积。因此 $p$ 为内积，当且仅当 $A$ 为实对称正定矩阵。由 Sylvester 判据，这恰为 $b=c$、$a>0$、$ad>b^2$。

（2）左乘正交矩阵只混合两行，不改变它们的内积之和；限制来自右乘矩阵。记题设左右两个正交矩阵为 $U,J$，迹的循环性给出
\[
p(UXJ,UYJ)=\operatorname{tr}(XJAJ^{\mathrm T}Y^{\mathrm T}).
\]
故保内积当且仅当 $JAJ^{\mathrm T}=A$。在 $A$ 对称的条件下，此式等价于 $a=d$、$b=0$。答案是 $\boxed{A=aI_2\ (a>0)}$。

\end{solve}

\section{矩阵方程、分块与数域问题}

\label{sec:solutions-9}

矩阵方程可以视为更大向量空间上的线性方程；分块计算须保持乘法次序。跨数域的相似问题除求解线性条件外，还须保证最终换基矩阵可逆。

\subsection{初赛第 4 届（2012 年度），来源序号7}\label{ans:pre04-math:7}\label{pr:pre04-math:7}

对应例~\ref{ex:pre04-math:7}（第~\pageref{ex:pre04-math:7} 页）；\href{https://github.com/lzxcr/cmc-mathematics/blob/64916fa1fdbc3e444a652eca204d5669058ee321/src/problems/preliminary/04/07.md}{查看来源}。

\begin{solve}

\textbf{（1）} 矩阵方程 $AX=B$ 可解，等价于 $B$ 的每一列都属于 $A$ 的列空间。若 $a=2$，该列空间是 $\operatorname{span}(1,1)^T$，不含 $B$ 的第一列 $(4,3)^T$，故必须 $a\ne2$。此时 $A$ 可逆，$AX=B$ 总可解；而 $BY=A$ 可解恰当 $B$ 可逆。因此所求条件为
\[
a\ne2,\qquad \det B=4-3b=0.
\]

\textbf{（2）} 相似矩阵有相同的迹和行列式，所以
\[
a+2=5,\qquad 2a-4=4-3b,
\]
即 $a=3$、$b=2/3$。反过来，这时两者的特征多项式均为 $t^2-5t+2$，有两个不同实根；两者都相似于由这两个根组成的同一对角矩阵，故相似。

\textbf{（3）} 因为 $A$ 对称，与之合同的 $B$ 必须对称，即 $b=3$。此时 $\det B=-5<0$，惯性为一正一负。由 Sylvester 惯性定理，$A$ 与 $B$ 合同，等价于 $A$ 也一正一负，即
\[
\det A=2a-4<0\quad\Longleftrightarrow\quad a<2.
\]

\end{solve}

\begin{remark}[方法与条件]

三个问题分别由列空间、相似不变量和二次型惯性控制。对称性本身也是合同不变量，必须在比较惯性之前先检查；对非对称矩阵直接套用惯性定理没有依据。

\end{remark}

\subsection{初赛第 5 届（2013 年度），来源序号6}\label{ans:pre05-math:6}\label{pr:pre05-math:6}

对应例~\ref{ex:pre05-math:6}（第~\pageref{ex:pre05-math:6} 页）；\href{https://github.com/lzxcr/cmc-mathematics/blob/64916fa1fdbc3e444a652eca204d5669058ee321/src/problems/preliminary/05/06.md}{查看来源}。

\begin{solve}

（1）同秩矩阵之间有 $B=PAQ$，其中 $P,Q$ 可逆。由乘法性及乘积的秩界，$\operatorname{rank}\phi(B)\le\operatorname{rank}\phi(A)$；再用 $A=P^{-1}BQ^{-1}$ 得反向不等式。这里不需要预先证明 $\phi(P)$ 可逆。

（2）记 $F_{ij}=\phi(E_{ij})$。由第一问，每个 $F_{ij}$ 都非零；由 $\phi(0)=0$，
\[
F_{ij}F_{kl}=\delta_{jk}F_{il}.
\]
取 $0\ne v_1\in\operatorname{im}F_{11}$，令 $v_i=F_{i1}v_1$，则 $F_{1i}v_i=v_1$，故 $v_i\ne0$，且
\[
F_{ij}v_k=\delta_{jk}v_i.
\]
若 $\sum_kc_kv_k=0$，作用 $F_{1j}$ 得 $c_jv_1=0$，所以这些向量构成一组基。以它们为列组成 $R$，上述作用公式就是 $F_{ij}R=RE_{ij}$，结论成立。

\end{solve}

\begin{remark}[方法与条件]

题设没有要求可加性，证明也只使用矩阵单位的乘法关系。矩阵单位本身已编码一组基及坐标之间的全部搬运规则。

\end{remark}

\subsection{初赛第 13 届补赛（2021 年度），来源序号3}\label{ans:ch08-13th-A-makeup:3}\label{pr:ch08-13th-A-makeup:3}

对应例~\ref{ex:ch08-13th-A-makeup:3}（第~\pageref{ex:ch08-13th-A-makeup:3} 页）；\href{https://github.com/lzxcr/cmc-mathematics/blob/64916fa1fdbc3e444a652eca204d5669058ee321/src/problems/preliminary/13-makeup/03.md}{查看来源}。

\begin{solve}

若没有可交换对，每个可交换族最多保留一个元素。阈值中的 $5^3$ 提示取全部对角矩阵，再补一个与它不相交的 $20$ 元族。考虑
\[
D=\{\operatorname{diag}(r_1,r_2,r_3):r_i\in R\},
\qquad
L=\{\operatorname{diag}(1,1,r)+tE_{12}:
 r\in R,\ t\in R\setminus\{0\}\}.
\]
$|D|=5^3$、$|L|=20$。$D$ 内的矩阵当然可交换；$L$ 中的对角部分彼此可交换，且都与 $E_{12}$ 可交换，所以 $L$ 也是可交换族。

若 $S$ 中不存在两个不同的可交换矩阵，则 $S$ 在每族中至多取一个元素，故
\[
|S|\leq |U\setminus(D\cup L)|+2
=5^9-5^3-20+2=5^9-5^3-18,
\]
与题设矛盾。

\end{solve}

\subsection{决赛第 3 届（2012 年），来源序号5}\label{ans:fin03-all:5}\label{pr:fin03-all:5}

对应例~\ref{ex:fin03-all:5}（第~\pageref{ex:fin03-all:5} 页）；\href{https://github.com/lzxcr/cmc-mathematics/blob/64916fa1fdbc3e444a652eca204d5669058ee321/src/problems/finals/03/05.md}{查看来源}。

\begin{solve}

记右侧矩阵为 $M$，直接计算得 $M^2=9M$、$\operatorname{rank}M=2$。因此 $A$ 为单射、$B$ 为满射，且
\[
A(BA-9I_2)B=M^2-9M=0.
\]
利用单射性和满射性消去两侧，得到 $BA=9I_2$。

\end{solve}

\begin{remark}[方法与条件]

一般地，若 $AB$ 的秩等于中间空间的维数，且 $(AB)^2=\lambda AB$，则 $BA=\lambda I$。这体现了将算子限制到其像空间的思想。

\end{remark}

\subsection{决赛第 3 届（2012 年），来源序号6}\label{ans:fin03-all:6}\label{pr:fin03-all:6}

对应例~\ref{ex:fin03-all:6}（第~\pageref{ex:fin03-all:6} 页）；\href{https://github.com/lzxcr/cmc-mathematics/blob/64916fa1fdbc3e444a652eca204d5669058ee321/src/problems/finals/03/06.md}{查看来源}。

\begin{solve}

考虑有理线性方程组
\[
A_iX=XB_i\qquad(i\in I).
\]
未知数只有有限个，即使指标集无限，也只需保留有限个线性无关的方程。因此其实解空间有一组有理矩阵基 $X_1,\ldots,X_r$。将已知解写成
$P=\sum_j t_jX_j$。

行列式连续且 $\det P\ne0$，所以在 $(t_1,\ldots,t_r)$ 的某邻域内，$\det(\sum_j s_jX_j)$ 均非零。在其中选取有理数 $s_j$，所得矩阵即为所求。

\end{solve}

\begin{remark}[方法与条件]

关键分成两层：交换关系是有理线性条件，可逆性是一个开条件。更一般地，在任意无限基域上，也可用“非零多项式不可能处处为零”替代有理数的稠密性，得到相似关系对扩域的下降性质。

\end{remark}

\subsection{决赛第 8 届（2017 年），来源序号6，共同题}\label{ans:fin08-hi:6}\label{pr:fin08-hi:6}

对应例~\ref{ex:fin08-hi:6}（第~\pageref{ex:fin08-hi:6} 页）；\href{https://github.com/lzxcr/cmc-mathematics/blob/64916fa1fdbc3e444a652eca204d5669058ee321/src/problems/finals/08/06.md}{查看来源}。

\begin{solve}

秩条件保证耦合块可消去。取可逆矩阵 $P,Q$，写成
\[
A=P\begin{pmatrix}I_r&0\\0&0\end{pmatrix}Q,\qquad
B=Q^{-1}\begin{pmatrix}C&D\\E&F\end{pmatrix}P^{-1}.
\]
由 $ABA=P\operatorname{diag}(C,0)Q$，题设给出 $\operatorname{rank}B=\operatorname{rank}C$。于是横向扩充为 $(C\ D)$、纵向扩充为 $\binom CE$ 都不能增加秩，故 $D$ 的列属于 $C$ 的列空间、$E$ 的行属于 $C$ 的行空间，即存在 $X,Y$ 使 $D=CX,E=YC$。于是
\[
AB\sim\begin{pmatrix}C&CX\\0&0\end{pmatrix}
=\begin{pmatrix}I&-X\\0&I\end{pmatrix}\begin{pmatrix}C&0\\0&0\end{pmatrix}\begin{pmatrix}I&X\\0&I\end{pmatrix},
\]
\[
BA\sim\begin{pmatrix}C&0\\YC&0\end{pmatrix}
=\begin{pmatrix}I&0\\Y&I\end{pmatrix}\begin{pmatrix}C&0\\0&0\end{pmatrix}\begin{pmatrix}I&0\\-Y&I\end{pmatrix}.
\]
两者相似于同一个矩阵，故互相相似。

\end{solve}

\begin{remark}[方法与条件]

一般的 $AB$ 与 $BA$ 可以在零特征值处有不同的 Jordan 结构。这里的秩条件使行、列耦合都通过 $C$ 分解，从而能用两次剪切换基消去；证明没有要求 $C$ 可逆。参见\hyperref[app:linear-invariant]{方法卷：不变子空间与分块}。

\end{remark}

\section{矩阵函数与数值线性代数}

\label{sec:solutions-10}

幂零矩阵把解析函数变成有限多项式，谱半径则控制线性迭代收敛。算法题同时说明输入、初值、更新和停止条件；有限步终止或收敛速度的结论应区分精确运算与浮点运算。

\subsection{初赛第 6 届（2014 年度），来源序号5}\label{ans:pre06-math:5}\label{pr:pre06-math:5}

对应例~\ref{ex:pre06-math:5}（第~\pageref{ex:pre06-math:5} 页）；\href{https://github.com/lzxcr/cmc-mathematics/blob/64916fa1fdbc3e444a652eca204d5669058ee321/src/problems/preliminary/06/05.md}{查看来源}。

\begin{solve}

写 $T=I+N$，其中 $N$ 严格下三角，故 $N^m=0$。令
\[
p(u)=(1+u)^n+(1+u)^\ell-2.
\]
其常数项为零、一次项系数为 $n+\ell\ne0$，所以可逐阶构造无常数项的多项式 $h(t)$，使
\[
p(h(t))\equiv t\pmod{t^m}.
\]
具体地，$t^j$ 项中新的未知系数只以 $(n+\ell)[t^j]h$ 出现，其余已由低阶系数确定，故每一步均可唯一解出。

取 $X=I+h(N)$。将 $t=N$ 代入同余式，利用 $N^m=0$，得到 $X^n+X^\ell=2I+N=I+T$。

\end{solve}

\begin{remark}[方法与条件]

下三角中的具体数字并不关键：对任何幂零 $N$，同一证明都能解 $X^n+X^\ell=2I+N$。这是\hyperref[app:algebra-interpolation]{截断多项式环中的复合逆}。

\end{remark}

\subsection{初赛第 7 届（2015 年度），来源序号3}\label{ans:pre07-math:3}\label{pr:pre07-math:3}

对应例~\ref{ex:pre07-math:3}（第~\pageref{ex:pre07-math:3} 页）；\href{https://github.com/lzxcr/cmc-mathematics/blob/64916fa1fdbc3e444a652eca204d5669058ee321/src/problems/preliminary/07/03.md}{查看来源}。

\begin{solve}

方程等价于 $(I+A)X=XA^2$，故对每个多项式 $p$ 都有 $p(I+A)X=Xp(A^2)$。取 $p$ 为 $I+A$ 的特征多项式，由 Cayley--Hamilton 定理得 $Xp(A^2)=0$。

$p$ 的根都是奇整数，而 $A^2$ 的特征值均为偶整数的平方，所以 $p(A^2)$ 的全部特征值非零，矩阵可逆。因此 $X=0$。

\end{solve}

\begin{remark}[方法与条件]

一般地，若两个矩阵 $S,T$ 的谱不相交，则 $SX=XT$ 只有零解。取 $S$ 的特征多项式即可证明，不要求任何矩阵可对角化；它表达了不同谱分量之间不能存在非零的交织映射。

\end{remark}

\subsection{决赛第 6 届（2015 年），来源序号13，高年级组}\label{ans:fin06-hi:12}\label{pr:fin06-hi:12}

对应例~\ref{ex:fin06-hi:12}（第~\pageref{ex:fin06-hi:12} 页）；\href{https://github.com/lzxcr/cmc-mathematics/blob/64916fa1fdbc3e444a652eca204d5669058ee321/src/problems/finals/06/13.md}{查看来源}。

\begin{solve}

令 $M=\alpha D-C$、$N=(\alpha-1)D+C^T$，则 $M-N=A$。记精确解为 $x_*=A^{-1}b$，相邻两步误差为 $u=x^{(k)}-x_*$、$v=x^{(k+1)}-x_*$。由 $Mv=Nu$，令 $d=u-v$，即得 $Md=Au$。

在能量范数 $\|w\|_A^2=w^TAw$ 下，
\[
\begin{aligned}
\|u\|_A^2-\|v\|_A^2
&=2d^TAu-d^TAd\\
&=d^T(M+M^T-A)d\\
&=(2\alpha-1)d^TDd.
\end{aligned}
\]
若 $u\ne0$，则 $d=M^{-1}Au\ne0$，所以能量严格下降。迭代映射 $G=M^{-1}N$ 是线性的；由有限维单位球面的紧性，存在统一常数 $q<1$ 使 $\|Gu\|_A\le q\|u\|_A$。因此
\[
\|x^{(k)}-x_*\|_A\le q^k\|x^{(0)}-x_*\|_A\longrightarrow0,
\]
迭代对任意初值均收敛到精确解。

\end{solve}

\begin{remark}[方法与条件]

这是\hyperref[app:linear-iteration]{线性迭代的能量方法}。一般地，若 $A$ 正定，分裂 $A=M-N$ 满足
\[
H=M+M^T-A>0,
\]
则同一恒等式给出收敛。此条件还保证 $M$ 可逆，因为 $M+M^T=A+H>0$。仅有 $M$ 的对称部分正定并不够：取 $A=D=I$、$C=0$、$\alpha=1/4$，则 $M=I/4>0$，而迭代矩阵为 $-3I$。

在本题这族分裂中，$\alpha=1/2$ 时能量保持，$\alpha<1/2$ 且 $M$ 可逆时非零误差的能量严格增加。因此 $\alpha>1/2$ 也正是从任意初值收敛所需的范围。

\end{remark}

\subsection{决赛第 7 届（2016 年），来源序号4，共同题}\label{ans:fin07-hi:4}\label{pr:fin07-hi:4}

对应例~\ref{ex:fin07-hi:4}（第~\pageref{ex:fin07-hi:4} 页）；\href{https://github.com/lzxcr/cmc-mathematics/blob/64916fa1fdbc3e444a652eca204d5669058ee321/src/problems/finals/07/04.md}{查看来源}。

\begin{solve}

答案为 $(1,0,1),(-1,0,-1),(1,t,-1),(-1,t,1)$，其中 $t\in\mathbb R$。对角元须为 $\pm1$；同号时非对角元为 $2016a^{2015}b$，故 $b=0$；异号时矩阵的平方已是单位阵，$b$ 可任取。

\end{solve}

\subsection{决赛第 9 届（2018 年），来源序号8，低年级组}\label{ans:fin09-lo:8}\label{pr:fin09-lo:8}

对应例~\ref{ex:fin09-lo:8}（第~\pageref{ex:fin09-lo:8} 页）；\href{https://github.com/lzxcr/cmc-mathematics/blob/64916fa1fdbc3e444a652eca204d5669058ee321/src/problems/finals/09/08.md}{查看来源}。

\begin{solve}

幂零矩阵的高次幂最终消失，因此可以把矩阵幂级数改写为有限个标量导数，从中直接读取端点奇性。设 $r$ 为最小的正整数使 $A^r=0$。对固定 $0\leq x<1$，按 $A$ 的幂展开并逐项求导标量幂级数，得到有限的函数演算公式
\[
G(x)=\sum_{j=0}^{r-1}(-1)^j\binom\alpha j(1-x)^{\alpha-j}A^j.
\]
这也证明原级数收敛，因为每个系数级数是收敛圆内的导数级数。

若 $A^3=0$，只含 $j=0,1,2$，而 $\alpha-j>-1$，故所有元素均可积。

若 $A^3\ne0$，则 $r-1\geq3$。取 $(A^{r-1})_{uv}\ne0$，最高次项在该元素中不能被较低奇性项抵消：
\[
g_{uv}(x)\sim(-1)^{r-1}\binom\alpha{r-1}(A^{r-1})_{uv}(1-x)^{\alpha-r+1}.
\]
系数非零且指数 $\alpha-r+1<-1$，故这一元素的广义积分发散，必要性得证。

\end{solve}

\begin{remark}[方法与条件]

若 $\alpha>-1$ 非整数，同一推导给出可积条件 $r-1<\alpha+1$。整数情形不能直接套用这个阈值，因为二项式系数会截断：若 $\alpha$ 为非负整数，$G$ 本身就是矩阵多项式。有限函数演算与 Jordan 块的联系见\hyperref[app:linear-cyclic]{方法卷：循环子空间与函数演算}。

\end{remark}

\subsection{决赛第 9 届（2018 年），来源序号13，高年级组}\label{ans:fin09-hi:11}\label{pr:fin09-hi:11}

对应例~\ref{ex:fin09-hi:11}（第~\pageref{ex:fin09-hi:11} 页）；\href{https://github.com/lzxcr/cmc-mathematics/blob/64916fa1fdbc3e444a652eca204d5669058ee321/src/problems/finals/09/13.md}{查看来源}。

\begin{solve}

将方程看作严格凸二次函数 $q(x)=\frac12x^TAx-b^Tx$ 的极小化问题。若搜索方向两两 $A$ 正交，沿新方向最小化便不会破坏旧方向上的极小性；共轭梯度法利用残差只需一步修正就能产生这样的方向。

任取 $x_0$，令 $r_0=b-Ax_0,p_0=r_0$。残差为零时停止，否则采用
\[
\begin{aligned}
\alpha_k&=\frac{r_k^Tr_k}{p_k^TAp_k},&x_{k+1}&=x_k+\alpha_kp_k,\\
r_{k+1}&=r_k-\alpha_kAp_k,&\beta_k&=\frac{r_{k+1}^Tr_{k+1}}{r_k^Tr_k},\\
p_{k+1}&=r_{k+1}+\beta_kp_k.
\end{aligned}
\]
下面同时推导系数并证明不变式：非零残差两两正交，方向两两 $A$ 正交，且两组向量的前 $k+1$ 项张成同一空间。

假设结论在第 $k$ 步前成立。由于 $p_k-r_k$ 属于旧方向空间，$r_k^Tp_k=\|r_k\|^2$，沿 $p_k$ 的精确线搜索便给出上式 $\alpha_k$，使 $r_{k+1}\perp p_k$。对 $j<k$，残差更新和 $p_j^TAp_k=0$ 又给出 $r_{k+1}\perp p_j$，故它也正交于全部旧残差。

对 $j<k$，利用
\[
Ap_j=\frac{r_j-r_{j+1}}{\alpha_j}
\]
可知 $r_{k+1}^TAp_j=0$；只须消去它与 $p_k$ 的 $A$ 内积。由 $r_{k+1}\perp r_k$，
\[
r_{k+1}^TAp_k=-\frac{\|r_{k+1}\|^2}{\alpha_k},
\]
所以令 $p_{k+1}=r_{k+1}+\beta_kp_k$，恰得到上述 $\beta_k$。张成空间的等式也随此递推保持，归纳完成。正定性保证所有未终止步骤的分母为正。

在精确算术下，若已有 $n$ 个非零方向，它们两两 $A$ 正交，必构成全空间的基；此时 $r_n$ 正交于全空间，只能为零。因此至多 $n$ 步求得唯一解。

\end{solve}

\begin{remark}[方法与条件]

有限终止来自有限维空间中的正交性，而不是收缩因子的极限。浮点舍入会破坏严格正交，故“至多 $n$ 步”并非实际误差保证。其投影及迭代观点见\hyperref[app:linear-iteration]{方法卷：线性迭代}；对于一般非对称问题，直接转为正规方程还会把二范数条件数平方。

\end{remark}

\subsection{决赛第 11 届（2021 年），来源序号8，低年级组}\label{ans:fin11-low:8}\label{pr:fin11-low:8}

对应例~\ref{ex:fin11-low:8}（第~\pageref{ex:fin11-low:8} 页）；\href{https://github.com/lzxcr/cmc-mathematics/blob/64916fa1fdbc3e444a652eca204d5669058ee321/src/problems/finals/11/08.md}{查看来源}。

\begin{solve}

矩阵幂级数在绝对收敛范围内保持标量解析恒等式，因为只涉及同一个矩阵的幂，它们彼此可交换。

（1）由 $\|A^k\|\leq\|A\|^k$，正弦、余弦级数对所有矩阵都绝对收敛。定义复矩阵指数后，
\[
\cos A=\frac{e^{iA}+e^{-iA}}2,\qquad\sin A=\frac{e^{iA}-e^{-iA}}{2i},\qquad e^{iA}e^{-iA}=I,
\]
直接平方相加便得 $\sin^2A+\cos^2A=I$。

（2）$\|A\|<1$ 保证反正切级数绝对收敛。标量恒等式
\[
\sin(\arctan z)-z\cos(\arctan z)=0\qquad(|z|<1)
\]
是该圆盘上的解析恒等式。将其幂级数复合、乘积展开；所有项由以 $\|A\|<1$ 代入的绝对收敛级数控制，故可重排。各系数均为零，代入 $A$ 即得 $\sin(\arctan A)=A\cos(\arctan A)$。

\end{solve}

\begin{remark}[方法与条件]

同一个矩阵生成的幂代数是交换的，绝对收敛又保证复合与乘积可以逐项计算，这才是标量恒等式能够移植的两项依据。对两个任意矩阵，三角加法公式一般不成立。系统的函数演算观点见\hyperref[app:linear-cyclic]{方法卷：循环子空间与函数演算}。

\end{remark}

\subsection{决赛第 11 届（2021 年），来源序号15，高年级组}\label{ans:fin11-high:13}\label{pr:fin11-high:13}

对应例~\ref{ex:fin11-high:13}（第~\pageref{ex:fin11-high:13} 页）；\href{https://github.com/lzxcr/cmc-mathematics/blob/64916fa1fdbc3e444a652eca204d5669058ee321/src/problems/finals/11/15.md}{查看来源}。

\begin{solve}

Gauss--Seidel 每次只改变一个坐标，并沿该坐标精确最小化二次能量。考察误差的齐次迭代，能量取 $q(e)=\frac12e^TAe$。若一次坐标更新将第 $i$ 个坐标增加 $\delta$，则 $\delta=-(Ae)_i/a_{ii}$，因而
\[
q(e+\delta e_i)-q(e)=-\frac12a_{ii}\delta^2\leq0.
\]
对任意初始误差 $e_0$，假设保证迭代趋于零，故 $q(e_0)\geq\lim q(e_k)=0$。因此 $A$ 半正定；再由非奇异性，$A$ 正定。

\end{solve}

\begin{remark}[方法与条件]

对角元为正只保证沿各坐标方向能量可下降；若全局二次型仍有负方向，从该方向出发能量就不可能下降到零。这把“任意初值收敛”转成了正定性的必要条件。它与共轭梯度法使用同一个二次能量，却选择不同的搜索方向，见\hyperref[app:linear-iteration]{方法卷：线性迭代}。

\end{remark}

\subsection{决赛第 13 届（2023 年），来源序号16，高年级组}\label{ans:fin13-high:13}\label{pr:fin13-high:13}

对应例~\ref{ex:fin13-high:13}（第~\pageref{ex:fin13-high:13} 页）；\href{https://github.com/lzxcr/cmc-mathematics/blob/64916fa1fdbc3e444a652eca204d5669058ee321/src/problems/finals/13/16.md}{查看来源}。

\begin{solve}

无位移 QR 迭代定义为
\[
A_0=A,\qquad A_k=Q_kR_k,\qquad A_{k+1}=R_kQ_k=Q_k^{\mathrm T}A_kQ_k,
\]
其中 $Q_k$ 正交，$R_k$ 上三角，非奇异时约定 $R_k$ 对角元为正。各步保持特征值；下三角部分足够小时，可由对角元近似读取实特征值。实矩阵出现共轭复特征值时，应允许实 Schur 形式中的 $2\times2$ 块。

\textbf{收敛需要附加条件。} 例如转角不属于 $\pi\mathbb Z$ 的平面旋转矩阵，其 QR 迭代保持不变，并不趋于上三角阵。下面给出一个标准的充分条件及准确结论。

把每步的正交基相乘，便能将 QR 迭代与 $A^k$ 的列空间联系起来；谱模分离决定这些列空间的极限。设 $A=X\Lambda X^{-1}$，实特征值满足 $|\lambda_1|>\cdots>|\lambda_n|>0$，且 $X^{-1}=LU$ 存在无选主元、$L$ 为单位下三角阵的 LU 分解。令 $\mathcal Q_k=Q_0\cdots Q_{k-1}$，则
\[
A_k=\mathcal Q_k^{\mathrm T}A\mathcal Q_k,\qquad
A^k=\mathcal Q_k(R_{k-1}\cdots R_0).
\]
又
\[
A^k=X(\Lambda^kL\Lambda^{-k})\Lambda^kU,
\qquad \Lambda^kL\Lambda^{-k}\longrightarrow I,
\]
因为严格下三角位置的因子是 $(\lambda_i/\lambda_j)^k$，$i>j$。

对 $X(\Lambda^kL\Lambda^{-k})$ 作对角元取正的 QR 分解 $\widehat Q_k\widehat R_k$。QR 分解在非奇异矩阵处连续，故 $\widehat Q_k\to Q$，其中 $X=QR$。比较 $A^k$ 的两种 QR 分解，存在对角符号矩阵 $D_k$ 使 $\mathcal Q_k=\widehat Q_kD_k$，从而
\[
D_kA_kD_k=\widehat Q_k^{\mathrm T}A\widehat Q_k
\longrightarrow Q^{\mathrm T}AQ=R\Lambda R^{-1}.
\]
右端上三角，故 $A_k$ 的严格下三角元趋于零，对角元趋于相应 $\lambda_i$。上三角元可能因符号振荡而不收敛；若 $A$ 对称，则右端必为对角阵，此时 $A_k\to\Lambda$。

\end{solve}

\subsection{决赛第 14 届（2023 年），来源序号6，共同题}\label{ans:fin14-high:6}\label{pr:fin14-high:6}

对应例~\ref{ex:fin14-high:6}（第~\pageref{ex:fin14-high:6} 页）；\href{https://github.com/lzxcr/cmc-mathematics/blob/64916fa1fdbc3e444a652eca204d5669058ee321/src/problems/finals/14/06.md}{查看来源}。

\begin{solve}

要解 $f(C)=aI+\varepsilon N$，先在标量层面把 $f$ 在 $1$ 附近求逆，只保留到 $n-1$ 阶。设
\[
f(1+w)=a+bw+c_2w^2+\cdots,
\qquad h(z)=u_1z+\cdots+u_{n-1}z^{n-1}.
\]
要求 $f(1+h(z))\equiv a+z\pmod{z^n}$。一次项给出 $u_1=b^{-1}$；对每个 $j\ge2$，第 $j$ 次项系数形如 $bu_j+$ 已由 $u_1,\ldots,u_{j-1}$ 确定的量，故因 $b\ne0$ 可逐次唯一解出 $u_j$。

取 $C=I+h(\varepsilon N)$。由于 $N^n=0$，将上述同余式代入 $z=\varepsilon N$ 时余项恰为零，从而 $f(C)=aI+\varepsilon N$。这里 $\varepsilon$ 无须足够小，若 $\varepsilon$ 为实数，构造出的 $C$ 也为实矩阵。

\end{solve}

\begin{remark}[方法与条件]

非零一次项保证有限阶复合逆存在；幂零性则把“模高次项成立”变成精确恒等式。其代数构造见\hyperref[app:algebra-interpolation]{方法卷：截断多项式与局部逆}。

\end{remark}

\subsection{决赛第 14 届（2023 年），来源序号16，高年级组}\label{ans:fin14-high:13}\label{pr:fin14-high:13}

对应例~\ref{ex:fin14-high:13}（第~\pageref{ex:fin14-high:13} 页）；\href{https://github.com/lzxcr/cmc-mathematics/blob/64916fa1fdbc3e444a652eca204d5669058ee321/src/problems/finals/14/16.md}{查看来源}。

\begin{solve}

约定 $L$ 为单位下三角阵、$U$ 为上三角阵。无选主元 Gaussian 消元在第 $k$ 步（$1\le k<n$）要求当前主元 $u_{kk}\ne0$，然后计算
\[
l_{ik}=u_{ik}/u_{kk}\quad(i>k),\qquad
u_{ij}\leftarrow u_{ij}-l_{ik}u_{kj}\quad(i,j>k),
\]
并将主元下方元素置零。记录全部乘子即得 $A=LU$；若主元为零，此无选主元算法不能继续，一般需采用带行置换的 $PA=LU$。

每步用 $n-k$ 次除法及 $(n-k)^2$ 次乘法、同样次数减法，总运算量为
\[
\sum_{k=1}^{n-1}\bigl((n-k)+2(n-k)^2\bigr)
=\frac23n^3+O(n^2).
\]
最后一个主元无须用于消元，因此即使 $A$ 奇异，仍可能存在唯一分解。记 $\Delta_k$ 为第 $k$ 阶顺序主子式，$\Delta_0=1$。若 $\Delta_1,\dots,\Delta_{n-1}$ 均非零，则消元主元 $u_{kk}=\Delta_k/\Delta_{k-1}$ 非零，每一步乘子及剩余矩阵被唯一确定，因此存在唯一分解。

反之，设分解存在且某个 $u_{kk}=0$，$k<n$。取 $E=e_{k+1}e_k^{\mathrm T}$。对任意实数 $t$，
\[
A=L(I+tE)(I-tE)U.
\]
因 $E^2=0$，首因子仍为单位下三角阵；由于 $U$ 第 $k$ 行在前 $k$ 列全为零，$(I-tE)U$ 仍上三角。这给出不同的分解，违背唯一性。因此 $u_{11},\dots,u_{n-1,n-1}$ 均非零，从而 $\Delta_k=\prod_{j\le k}u_{jj}\ne0$。

\end{solve}

\subsection{决赛第 17 届（2026 年），来源序号1，共同题}\label{ans:fin17-high:1}\label{pr:fin17-high:1}

对应例~\ref{ex:fin17-high:1}（第~\pageref{ex:fin17-high:1} 页）；\href{https://github.com/lzxcr/cmc-mathematics/blob/64916fa1fdbc3e444a652eca204d5669058ee321/src/problems/finals/17/01.md}{查看来源}。

\begin{solve}

括号内是旋转角 $\theta=\pi/2-\pi/1013$ 的平面旋转矩阵。其 $2026$ 次幂的旋转角为 $1011\pi$，所以 $A=-I_2$，$\boxed{\operatorname{tr}A=-2}$。

\end{solve}

\subsection{决赛第 17 届（2026 年），来源序号16，高年级组}\label{ans:fin17-high:13}\label{pr:fin17-high:13}

对应例~\ref{ex:fin17-high:13}（第~\pageref{ex:fin17-high:13} 页）；\href{https://github.com/lzxcr/cmc-mathematics/blob/64916fa1fdbc3e444a652eca204d5669058ee321/src/problems/finals/17/16.md}{查看来源}。

\begin{solve}

（1）离散 Dirichlet 边界条件 $v_0=v_{n+1}=0$ 提示取
\[
v_j^{(k)}=\sin\frac{jk\pi}{n+1}\quad(1\le j\le n).
\]
由正弦加法公式，$-v_{j-1}^{(k)}+2v_j^{(k)}-v_{j+1}^{(k)}=\lambda_kv_j^{(k)}$，其中
\[
\lambda_k=2-2\cos\frac{k\pi}{n+1}=4\sin^2\frac{k\pi}{2(n+1)}.
\]
这 $n$ 个特征值互异，因而已是全部谱。于是
\[
\kappa_2(A_n)=\cot^2\frac\pi{2(n+1)}\sim\frac{4n^2}{\pi^2}.
\]
（2）误差迭代矩阵为 $I-\alpha A_n$，其谱半径为
\[
\max\{|1-\alpha\lambda_1|,|1-\alpha\lambda_n|\}.
\]
对任意初值收敛当且仅当 $\boxed{0<\alpha<2/\lambda_n}$。谱区间上 $|1-\alpha\lambda|$ 是凸函数，最大值在两端；使两端取相反且等大的值，即可最小化最大误差，故
\[
\boxed{\alpha_*=\frac2{\lambda_1+\lambda_n}=\frac12},\qquad
\rho_*=\frac{\lambda_n-\lambda_1}{\lambda_n+\lambda_1}=\cos\frac\pi{n+1}.
\]
（3）Jacobi 的对角阵为 $2I$，因此它恰是最优参数的 Richardson 迭代，必收敛。其收敛因子满足
\[
\rho=1-\frac{\pi^2}{2(n+1)^2}+O(n^{-4}).
\]
故最坏情形下将误差缩小到原来的 $\varepsilon$ 倍，需要约 $2(n+1)^2\log(1/\varepsilon)/\pi^2$ 步；矩阵阶数增大时收敛显著变慢。

\end{solve}

\begin{remark}[方法与条件]

$A_n$ 是一维 Dirichlet Laplace 算子的差分离散，正弦向量对应连续问题的正弦本征函数。低频误差衰减缓慢，提示用预条件或多重网格改善网格加密后的迭代效率。

本题的慢收敛由最小特征值对应的低频模决定；一般谱迭代与 QR、Krylov 方法的联系见\hyperref[app:linear-iteration]{方法卷：谱与迭代}。

\end{remark}

\part{专题方法与证明工具}

本编主要节选并改编自 Jacobi Lie 的方法稿，作者、版本与 CC BY-SA 4.0 许可见编写说明。以下内容是方法讲解，不作为额外真题计数。环与模的语言属于进阶阅读，其作用是解释多项式整除与矩阵结构之间的联系。



\section{矩阵与算子的结构、范数和迭代}

矩阵是线性映射在一组基下的表示。结构问题先选不变空间或合适基，估计问题再选与几何相容的范数。把这两步分开，既能避免大行列式展开，也能解释矩阵公式中哪些性质与坐标无关。

\subsection{算子范数、谱半径与扰动}

\label{app:linear-norm}

给定向量范数，诱导算子范数定义为 $\|A\|=\sup_{\|x\|=1}\|Ax\|$。由定义有 $\|Ax\|\leq\|A\|\|x\|$ 和 $\|AB\|\leq\|A\|\|B\|$。欧氏范数诱导的 $2$ 范数等于最大奇异值：$\|A\|_2^2=\lambda_{\max}(A^*A)$。Frobenius 范数 $\|A\|_F^2=\operatorname{tr}(A^*A)$ 则是所有奇异值平方之和；二者不同，但都可用于乘积估计。

谱半径 $\rho(A)=\max|\lambda|$ 满足 $\rho(A)\leq\|A\|$。对正规矩阵有 $\rho(A)=\|A\|_2$，一般矩阵则可能相差很大，例如 $\begin{pmatrix}0&M\\0&0\end{pmatrix}$ 的谱半径为零而范数为 $|M|$。特征值不能单独描述短时间的放大。

若 $\|E\|<1$，几何级数在范数中绝对收敛，且
\[
(I-E)^{-1}=\sum_{k=0}^\infty E^k,
\qquad\|(I-E)^{-1}\|\leq\frac1{1-\|E\|}.
\]
因此可逆 $A$ 在 $\|A^{-1}E\|<1$ 下满足
\[
\|(A+E)^{-1}-A^{-1}\|
\leq\frac{\|A^{-1}\|^2\|E\|}{1-\|A^{-1}E\|}.
\]
证明先提出 $A$，再用 $X^{-1}-Y^{-1}=X^{-1}(Y-X)Y^{-1}$。条件数 $\kappa(A)=\|A\|\|A^{-1}\|$ 衡量求解线性方程对相对误差的敏感度；可逆只回答能否求解，条件数回答解是否稳定。

\subsection{不变子空间、商空间与滤过}

\label{app:linear-invariant}

若 $W\subseteq V$ 满足 $T(W)\subseteq W$，先取 $W$ 的基再扩充为 $V$ 的基，矩阵呈分块上三角形
\[
[T]=\begin{pmatrix}A&B\\0&C\end{pmatrix}.
\]
$A$ 表示限制 $T|_W$，$C$ 表示商空间 $V/W$ 上的诱导算子。故特征多项式分解为 $\chi_T=\chi_A\chi_C$，迹相加、行列式相乘。耦合块 $B$ 不影响这些量，却可能影响 Jordan 结构，不能随意删掉。

一列不变空间 $0=V_0\subset V_1\subset\cdots\subset V_m=V$ 称为滤过。按这列空间选基，上述过程逐层进行。在多项式空间中，次数不超过 $k$ 的子空间常天然不变；只需看算子对最高次项的作用，就能读出三角矩阵对角元。递归求系数因此具有结构来源。

\textbf{应用。} 若 $AB=BA$，则 $A$ 的每个特征空间及广义特征空间都被 $B$ 保持，因为 $(A-\lambda I)^kB=B(A-\lambda I)^k$。在复数域上，先取非零普通特征空间 $\ker(A-\lambda I)$，再在其中取 $B$ 的特征向量，就得到共同特征向量。对有限个两两交换的矩阵，逐次限制到已有的共同特征空间，可得共同不变直线；再对商空间按维数归纳，便得到同时上三角化。实数域上若特征多项式不分裂，例如平面旋转，这个结论需改成实分块形式。

\subsection{谱定理与正交几何}

\label{app:linear-spectral}

复矩阵正规指 $AA^*=A^*A$。谱定理断言它可酉对角化；实对称矩阵可正交对角化。一个简短证明从酉 Schur 分解出发：正规上三角矩阵第一行与第一列的平方模和相同，迫使第一行非对角项为零，再归纳。Hermite 情形的特征值为实数，酉矩阵的特征值在单位圆上。

Hermite 矩阵 $H$ 有正交谱分解 $H=\sum_\lambda\lambda P_\lambda$，其中 $P_\lambda$ 两两正交且和为 $I$。因此
\[
\lambda_{\min}\|x\|^2\leq x^*Hx\leq\lambda_{\max}\|x\|^2.
\]
正定等价于所有特征值为正，也等价于存在唯一正定平方根 $H^{1/2}$。这将二次型估计变成谱区间估计，并提供自然坐标 $y=H^{1/2}x$。

若两个 Hermite 矩阵交换，则它们可同时酉对角化：一个矩阵的各特征空间被另一个保持，再分别使用谱定理。一般矩阵的同时上三角化与这里的同时正交对角化所需条件不同；后者多出的正交性正是范数计算变得精确的原因。

\subsection{低秩分解与小矩阵}

\label{app:linear-lowrank}

若 $A=UV$，其中 $U$ 为 $n\times r$、$V$ 为 $r\times n$，应先计算小矩阵 $VU$。分块消元给出 Sylvester 恒等式
\[
\det(I_n+tUV)=\det(I_r+tVU).
\]
当 $n\geq r$ 时，$\det(\lambda I_n-UV)=\lambda^{n-r}\det(\lambda I_r-VU)$，先对 $\lambda\neq0$ 推导，再由多项式恒等延至所有 $\lambda$。所以非零特征值及代数重数由小矩阵完全决定。

进一步，$(UV)^k=U(VU)^{k-1}V$（$k\geq1$）。若 $VU=cI_r$，便有 $A^2=cA$；若 $VU$ 只满足其他多项式关系，则把相应关系传回即可。不能仅由“秩很低”推断特定的二次恒等式。

\textbf{应用。} 矩阵 $a_{ij}=\cos(\theta_i-\theta_j)$ 可写成 $A=cc^T+ss^T$，其中 $c_i=\cos\theta_i$、$s_i=\sin\theta_i$。秩至多为 $2$，全部非零谱信息由
\[
\begin{pmatrix}\sum c_i^2&\sum c_is_i\\\sum c_is_i&\sum s_i^2\end{pmatrix}
\]
给出。三角恒等式的真正作用是把列空间缩到两个方向。

\subsection{循环向量与有限函数演算}

\label{app:linear-cyclic}

在域 $K$ 上定义 $p(x)\cdot v=p(T)v$，便把 $V$ 看成 $K[x]$ 模。固定 $v$，映射 $p\mapsto p(T)v$ 的核由唯一首一多项式 $\mu_v$ 生成，故
\[
K[T]v\cong K[x]/(\mu_v),\qquad\dim K[T]v=\deg\mu_v.
\]
若 $v$ 是循环向量，即 $K[T]v=V$，每个与 $T$ 交换的 $S$ 都是 $T$ 的多项式：写 $Sv=p(T)v$，再利用 $S(q(T)v)=q(T)Sv$。这说明交换代数为何在循环情形中特别简单。多块情形由\hyperref[app:algebra-modules]{PID 模结构定理}统一解释。

设极小多项式为 $\mu_T$。多项式函数 $p(T)$ 只依赖 $p$ 模 $\mu_T$ 的余类。若 $\mu_T=\prod_j(x-\lambda_j)^{m_j}$ 在 $\mathbb C$ 上分裂，解析函数 $f$ 的矩阵值由各点的导数 $f^{(k)}(\lambda_j)$（$0\leq k<m_j$）确定：用 Hermite 插值找次数小于 $\deg\mu_T$ 的 $p$ 匹配这些导数，再定义 $f(T)=p(T)$。对 Jordan 块 $T=\lambda I+N$，它就是有限式
\[
f(T)=\sum_{k=0}^{m-1}\frac{f^{(k)}(\lambda)}{k!}N^k,
\qquad N^m=0.
\]
对于整函数，这与范数绝对收敛的幂级数定义一致。幂零部分使无限解析过程退化为有限代数运算。

\subsection{Gram 矩阵、对偶证书与 Schur 补}

\label{app:linear-gram}

向量 $u_1,\ldots,u_m$ 的 Gram 矩阵 $G=(\langle u_j,u_i\rangle)_{ij}$ 满足 $c^*Gc=\|\sum_jc_ju_j\|^2\geq0$，正定当且仅当这些向量线性无关。约定内积对第一变量线性。最佳逼近系数由正规方程 $Gc=(\langle v,u_i\rangle)_i$ 求出；即使生成元相关，投影向量仍唯一，只是系数可能不唯一。

对 Hermite 分块矩阵，若 $A>0$，配方给出
\[
\begin{aligned}
\begin{pmatrix}x\\y\end{pmatrix}^*
\begin{pmatrix}A&B\\B^*&C\end{pmatrix}
\begin{pmatrix}x\\y\end{pmatrix}
={}&(x+A^{-1}By)^*A(x+A^{-1}By)\\
&+y^*(C-B^*A^{-1}B)y.
\end{aligned}
\]
故整体半正定当且仅当 Schur 补 $C-B^*A^{-1}B$ 半正定。这个等价关系统一了配方、不等式最佳常数、条件协方差和消元；$A$ 不可逆时需加入像空间条件或改用广义逆，不能直接沿用公式。

\textbf{应用。} 对 $A>0$ 和非零 $c$，约束 $c^*x=1$ 下 $x^*Ax$ 的最小值为 $1/(c^*A^{-1}c)$。令 $y=A^{1/2}x$，约束成为 $\langle y,A^{-1/2}c\rangle=1$，由 Cauchy--Schwarz 即得，等号在 $x=A^{-1}c/(c^*A^{-1}c)$ 取得。变换坐标把一般二次问题恢复为正交投影。

对列向量 $u_1,\ldots,u_n$ 组成的方阵 $U$，还有 Hadamard 不等式
\[
|\det U|\leq\prod_{j=1}^n\|u_j\|_2.
\]
若列相关，两边关系直接成立；否则 Gram--Schmidt 将体积写成每列到此前列空间距离的乘积，而每个距离不超过原列长度。等号当且仅当至少一列为零，或所有列两两正交；若各列均非零而线性相关，左端为零、右端为正，不能取等。它是 Gram 行列式的几何估计，也常与整数行列式的整除性结合，得到比单纯实数上界更强的结论。

\subsection{线性迭代与能量下降}

\label{app:linear-iteration}

迭代 $x_{k+1}=Gx_k+c$ 若有不动点 $x_*$，误差满足 $e_k=G^ke_0$。所有初值都收敛到同一不动点 $x_*$，当且仅当 $\rho(G)<1$。充分性可由 Jordan 块说明：$G^k$ 是有限个多项式因子乘 $\lambda^k$；必要性可取特征向量。某个现成范数下 $\|G\|<1$ 是方便的充分条件，却不是必要条件。

\textbf{分裂迭代的统一判据。} 设 $A=A^*>0$、$M$ 可逆，$G=I-M^{-1}A$。由 $M(I-G)=A$ 展开可得
\[
A-G^*AG=(I-G)^*(M+M^*-A)(I-G).
\]
若 $M+M^*-A>0$，右端正定，因而 $\|Gx\|_A<\|x\|_A$ 对非零 $x$ 成立。在有限维单位球面上取最大值，就有 $\|G\|_A<1$，迭代收敛。这个能量恒等式把非对称迭代矩阵的谱问题转成一个对称矩阵的正定性；仅有 $M+M^*>0$ 则不足。

对实对称正定 $A$，解 $Ax=b$ 等价于最小化 $q(x)=x^TAx/2-b^Tx$。沿方向 $p$ 的精确最小化步长为 $\alpha=p^T(b-Ax)/(p^TAp)$。若方向两两 $A$ 正交，后一步不会破坏此前已经满足的正交条件；在精确算术下，共轭梯度法至多经过 $n$ 个独立方向终止。

这种有限终止性不等于浮点运算中的实际步数保证。更有解释力的是谱与多项式估计：记 $r_0=b-Ax_0$，第 $k$ 步在仿射空间 $x_0+\mathcal K_k(A,r_0)$ 中最小化误差的 $A$ 范数，其中 $\mathcal K_k(A,r_0)=\operatorname{span}\{r_0,Ar_0,\ldots,A^{k-1}r_0\}$。误差可写为 $p_k(A)e_0$，其中 $\deg p_k\leq k$、$p_k(0)=1$，故可用谱区间上的多项式逼近控制。标准 Chebyshev 估计给出
\[
\|e_k\|_A\leq2\left(\frac{\sqrt\kappa-1}{\sqrt\kappa+1}\right)^k\|e_0\|_A,
\qquad\kappa=\frac{\lambda_{\max}(A)}{\lambda_{\min}(A)}.
\]
这里 $\|x\|_A^2=x^TAx$。竞赛中通常只需证明能量下降或正交不变量；这一推广解释了为什么改善谱分布的预处理能够加速计算。



\section{环、模、域与矩阵结构}

多项式记录代数关系，商结构把关系变成运算规则，群作用把对称性变成计数。许多题的计算看似不同，实际都在寻找可保留的信息：取模保留余式，插值保留有限点处的数据，元素的幂保留阶的某个素数部分。本章从这些操作的依据出发，说明它们如何连接。

\subsection{环、理想与商：把关系纳入运算}

\label{app:algebra-quotient}

本章的环均指交换含幺环。域中的每个非零元都可除；一般环只保留加法、乘法及其分配律，例如整数环 $\mathbb Z$ 和多项式环 $K[t]$。没有非零零因子的非零环称为整环：$ab=0$ 时必有 $a=0$ 或 $b=0$。整数与域上的多项式都满足这一条件。

若想把一批元素视为零，这批元素必须对加减封闭，并且乘任意环元后仍在其中。满足这两条规则的子集 $I$ 称为理想。定义 $a\equiv b\pmod I$ 当且仅当 $a-b\in I$，则剩余类可以相加、相乘，并且结果与代表元选择无关；所得环记为 $R/I$。理想的吸收性质正是商上乘法有定义的原因。

例如，$R/(d)$ 表示强制关系 $d=0$；$(d)=\{rd:r\in R\}$ 称为主理想。每个理想均为主理想的整环称为主理想整环，简称 PID。整数的除法和域上多项式的带余除法分别证明 $\mathbb Z$、$K[t]$ 是 PID；下一节给出多项式情形的证明。对一般环不应默认每个理想都由单个元素生成。

\textbf{例：用商构造复数。} 在 $\mathbb R[t]/(t^2+1)$ 中，$t$ 的剩余类满足 $t^2=-1$，每个元素唯一写成 $a+bt$。映射 $a+bt\mapsto a+bi$ 保持加法、乘法且双射，故
\[
\mathbb R[t]/(t^2+1)\cong\mathbb C.
\]
这里商是域，与 $t^2+1$ 在 $\mathbb R[t]$ 中不可约对应。一般地，若 $p\in K[t]$ 不可约，非零剩余类 $\bar f$ 的代表元与 $p$ 互素；Bézout 恒等式 $uf+vp=1$ 给出逆元 $\bar u$。反之，若 $p=ab$ 且两因子次数都较低，$\bar a\bar b=0$ 而两者均非零，商便不是域。

商也可保留幂零关系。例如 $\mathbb R[t]/(t^2)$ 中 $t\ne0$ 而 $t^2=0$；这正适合记录一阶数据。它与 $\mathbb C$ 同为二维实向量空间，乘法结构却不同。理解商时，既要找一组代表元，也要写清它们之间的乘法关系。

\subsection{整除、理想与 Bézout 恒等式}

\label{app:algebra-bezout}

设 $K$ 为域。多项式带余除法使每个非零理想 $I\subseteq K[x]$ 都由一个首一多项式生成：取 $I$ 中次数最低的首一元 $d$，对任意 $f\in I$ 作除法，余式仍属 $I$，故只能为零。特别地，由 $p,q$ 生成的理想为 $(d)$，其中 $d=\gcd(p,q)$，所以存在 $u,v\in K[x]$ 满足
\[
up+vq=d.
\]
这比“没有公共根”更有用：当 $p,q$ 互素时，它给出显式的恒等分解 $up+vq=1$，可以代入矩阵、作用于向量，或在商环中使用。

例如，若 $p,q$ 互素且 $p(A)q(A)=0$，则
\[
w=u(A)p(A)w+v(A)q(A)w
\]
将任意向量分成分别属于 $\ker q(A)$ 和 $\ker p(A)$ 的两部分；交集为零也由同一恒等式推出。这是\hyperref[pr:fin15-high:6]{矩阵四次方与秩等式题}的核心。算子分解的系统表述见\hyperref[app:linear-invariant]{不变子空间}。

一个需要区分的前提是系数域：在 $\mathbb R[x]$ 中没有共同实根不等于互素，例如 $x^2+1$ 与 $(x^2+1)^2$。应比较不可约因子，或在代数闭包中比较全部根。

\subsection{中国剩余定理：把独立条件合在一起}

\label{app:algebra-crt}

设 $p_1,\ldots,p_s\in K[x]$ 两两互素，$P=\prod_i p_i$。令 $P_i=P/p_i$，用 Bézout 恒等式选择 $u_i$ 使 $u_iP_i\equiv1\pmod{p_i}$，记 $e_i=u_iP_i$。它们满足
\[
e_i\equiv\begin{cases}1\pmod{p_i},\\0\pmod{p_j},\quad j\ne i.\end{cases}
\]
因此指定任意余式 $r_i$ 后，$f=\sum_i r_ie_i$ 满足全部同余条件，且解在模 $P$ 意义下唯一：两个解之差被每个 $p_i$ 整除，因互素而被乘积整除。这证明
\[
K[x]/(P)\cong\prod_{i=1}^sK[x]/(p_i).
\]
在商环中，$e_i^2=e_i$、$e_ie_j=0$（$i\ne j$）、$\sum_i e_i=1$。这些不是额外巧合，而是乘积结构的坐标投影。

若若干矩阵 $A_i$ 的最小多项式 $p_i$ 两两互素，令 $f\equiv f_i\pmod{p_i}$，就有 $f(A_i)=f_i(A_i)$。同一多项式因而能在彼此独立的代数分量上实现不同操作。这正是\hyperref[pr:ch06-14th-A-makeup2:3]{多矩阵同时插值题}；对单个算子的应用见\hyperref[app:linear-cyclic]{循环结构}。

整数的中国剩余定理有完全相同的证明：把 $K[x]$ 换成 $\mathbb Z$，把互素多项式换成互素整数。无论在哪种情形，先构造“在一个分量为一、其他分量为零”的元素，比同时解全部系数方程更清楚。

\subsection{插值、重根数据与有限阶逆}

\label{app:algebra-interpolation}

在互异点 $a_1,\ldots,a_n$ 给定值 $b_1,\ldots,b_n$，次数小于 $n$ 的插值多项式唯一。令 $P(x)=\prod_j(x-a_j)$，中国剩余定理在一次因子处具体化为
\[
f(x)=\sum_{j=1}^nb_j\ell_j(x),\qquad
\ell_j(x)=\frac{P(x)}{(x-a_j)P'(a_j)}.
\]
因为 $\ell_j(a_i)=\delta_{ij}$，存在性直接可验；两个解之差有 $n$ 个不同根而次数小于 $n$，故为零。

比较最高次系数有时比求出 $f$ 更有用。对 $f(x)=x^m$（$0\le m<n$），得到
\[
\sum_{j=1}^n\frac{a_j^m}{P'(a_j)}
=\begin{cases}0,&m<n-1,\\1,&m=n-1.\end{cases}
\]
这给出\href{https://github.com/lzxcr/cmc-mathematics/blob/64916fa1fdbc3e444a652eca204d5669058ee321/src/problems/finals/15/11.md}{根与导数加权和（原稿链接）}。它也等价于部分分式
\[
\frac1{P(z)}=\sum_j\frac1{P'(a_j)(z-a_j)}
\]
在无穷远的系数比较，因此与\href{https://github.com/lzxcr/cmc-mathematics/blob/64916fa1fdbc3e444a652eca204d5669058ee321/src/methods/inner/08-complex.md}{留数方法（原稿链接）}自然相通。非首一多项式只需保留首项系数；有重根时 $P'(a_j)=0$，不能继续套同一公式。

重根数据应改为指定模 $(x-a_j)^{m_j}$ 的余式。这些幂仍两两互素，所以中国剩余定理照常适用。在特征零时，这相当于同时指定函数值及前 $m_j-1$ 阶导数，称为 Hermite 插值。在任意特征下，直接使用余式比除以阶乘更稳妥。

有限阶复合逆也在同一类商环中进行。设
\[
f(c+w)=a+bw+c_2w^2+\cdots,\qquad b\ne0.
\]
求 $h(z)=u_1z+\cdots+u_{n-1}z^{n-1}$ 使 $f(c+h(z))\equiv a+z\pmod{z^n}$。第一步给 $u_1=b^{-1}$；第 $j$ 步只有 $bu_j$ 含新未知量，其他部分由此前系数确定，因此可递归求解。若 $N^n=0$，代入 $z=N$ 后同余即成为精确矩阵恒等式。这解释了\hyperref[pr:fin14-high:6]{矩阵多项式逆像题}为何不需要限制扰动大小。

\hyperref[pr:fin14-low:8]{循环矩阵相似题}则在互异单位根处插值：循环移位矩阵的 Fourier 特征基固定，只需指定多项式在各特征值上的取值，就能控制对应循环矩阵的全部特征值。

\subsection{根界与系数：由对称函数传递估计}

\label{app:algebra-roots}

设首一多项式为 $P(z)=z^n+a_1z^{n-1}+\cdots+a_n$，记 $M=\max_j|a_j|$。若 $M>0$ 且 $|z|\ge1+M$，则
\[
M\sum_{j=0}^{n-1}|z|^j
=M\frac{|z|^n-1}{|z|-1}<|z|^n.
\]
最高次项不能被其余项抵消，所以每个根满足 $|z|<1+M$。$M=0$ 时全部根为零。这是 Cauchy 根界。

若一个首一因子 $Q$ 的次数为 $m$，其全部根 $\zeta_1,\ldots,\zeta_m$ 也在同一圆盘内。Vieta 公式把第 $k$ 个系数写成 $\binom mk$ 个 $k$ 根乘积之和，因此
\[
|[z^{m-k}]Q|\le\binom mk\bigl(\max_j|\zeta_j|\bigr)^k.
\]
这给出一个可靠的推理方向：原多项式系数界决定根界，根界再决定任意首一因子的系数界。\hyperref[pr:fin16-high:6]{因子系数异常题}正是它的逆否命题。

根数界还有不同用途：域上非零 $d$ 次多项式至多有 $d$ 个不同根，因为每出现一个根就可除去一个一次因子。这一结论在整环上也成立，但有零因子的环上根数可能超过次数。后面将用它识别有限乘法群。

\subsection{PID 上的模与矩阵的不变因子}

\label{app:algebra-modules}

模把向量空间的标量域换成环：加法仍组成交换群，标量乘法满足通常的结合律与分配律。整数模就是交换群，其中整数乘法由重复相加定义。对域 $K$ 上的有限维空间 $V$ 及算子 $A$，令
\[
f(t)\cdot v=f(A)v,
\]
便使 $V$ 成为 $K[t]$ 模。一个子空间是子模，恰好意味着它在 $A$ 下不变；两个矩阵相似，恰好意味着对应的 $K[t]$ 模同构。这里把“换基后是什么矩阵”变成了“带这套乘法的空间是什么结构”。

\textbf{结构定理（引用）。} 若 $R$ 为 PID，$M$ 为有限生成 $R$ 模，则
\[
M\cong R^r\oplus R/(d_1)\oplus\cdots\oplus R/(d_s),
\qquad d_1\mid d_2\mid\cdots\mid d_s,
\]
其中 $r\ge0$，$d_i$ 是非零非单位元；允许 $s=0$。$r$ 与 $d_i$ 在乘单位元的意义下唯一。这是 PID 上有限生成模结构定理，本章将它作为标准定理引用；下一节的 Smith 标准形解释如何从关系矩阵读出这些量。在 $\mathbb Z$ 上约定 $d_i>0$，在 $K[t]$ 上约定首一，即消除单位元的不确定性。

非零有限维空间 $V$ 的每个向量都被某个非零多项式消去，因为 $v,Av,A^2v,\ldots$ 必线性相关。因此上述 $K[t]$ 模没有自由部分，分解为 $K[t]/(d_i)$ 的直和，且 $\dim_K V=\sum_i\deg d_i$。在一个分量中，若
\[
d(t)=t^m+c_{m-1}t^{m-1}+\cdots+c_0,
\]
则 $1,t,\ldots,t^{m-1}$ 为基，乘 $t$ 的矩阵各列依次为 $e_2,\ldots,e_m,-\sum_{j=0}^{m-1}c_je_{j+1}$，称为伴随矩阵 $C(d)$。直接计算其特征多项式得 $d$；使乘法算子为零的多项式又恰为 $(d)$ 中的元素，故最小多项式也是 $d$。

于是每个 $A$ 都相似于唯一确定的分块矩阵
\[
\operatorname{diag}(C(d_1),\ldots,C(d_s)),
\qquad
\chi_A=\prod_{i=1}^s d_i,\qquad m_A=d_s.
\]
这是有理标准形，$d_i$ 称为不变因子；“有理”意指在原系数域内构造，无须先添加特征值。最后一式来自直和的消去多项式必须同时被所有 $d_i$ 整除，而它们已按整除关系排列。特别地，最小多项式等于特征多项式当且仅当只有一个分量；这一情形的直接向量证明见\hyperref[app:linear-cyclic]{循环向量}。

\textbf{例：对角矩阵也可有二阶循环块。} 对 $A=\operatorname{diag}(2,2,3)$，取 $v=e_1+e_3$。向量 $v,Av=2e_1+3e_3$ 线性无关，且 $A^2v=5Av-6v$，所以在基 $e_2,v,Av$ 下，
\[
A\sim\operatorname{diag}\!\left(2,
\begin{pmatrix}0&-6\\1&5\end{pmatrix}\right).
\]
不变因子是 $t-2$ 与 $(t-2)(t-3)$；它们的乘积给特征多项式，末项给最小多项式。原来按特征值分开的 $2,3$ 两个一维空间合成一个循环块，这正是中国剩余定理 $K[t]/((t-2)(t-3))\cong K[t]/(t-2)\oplus K[t]/(t-3)$ 的矩阵版本。

\subsection{Smith 标准形与 Jordan 形式的来源}

\label{app:algebra-smith}

有限生成模可用“生成元与关系”描述。若 $B$ 是 PID $R$ 上的 $m\times n$ 矩阵，则
\[
M=R^m/B R^n
\]
以 $m$ 个标准向量为生成元，以 $B$ 的各列为关系。对生成元换基相当于可逆行变换，对关系列表换基相当于可逆列变换，所以 $B$ 与 $UBV$ 给出同构的商模。这里 $U,V$ 可独立选择；矩阵相似要求左右变换互逆，解决的是另一种分类问题。

\textbf{Smith 标准形（引用）。} 存在可逆 $U,V$，使 $UBV$ 为矩形对角矩阵，非零对角元满足 $d_1\mid\cdots\mid d_r$。于是
\[
M\cong R^{m-r}\oplus\bigoplus_{i=1}^r R/(d_i).
\]
单位元对应的商为零，可删除。此定理与上一节的结构定理相对应。它的消元原理是用 Bézout 操作把两个元素替换为最大公因子与相应倍数，再使当前主元整除其余子矩阵；完整存在性证明不在这里展开。

计算时常只需各阶子式。令 $\Delta_k$ 为所有 $k$ 阶子式的最大公因子，$\Delta_0=1$，统一取正整数或首一多项式。则对 $1\le k\le r$，
\[
\Delta_k=d_1\cdots d_k,
\qquad d_k=\frac{\Delta_k}{\Delta_{k-1}}.
\]
理由是可逆行列变换保留各阶子式生成的理想：变换后的每个子式是旧子式的线性组合，逆变换给出反向包含；对满足整除链的对角矩阵，子式理想恰由前 $k$ 个对角元的乘积生成。

\textbf{例：有限交换群的关系矩阵。} 取
\[
B=\begin{pmatrix}2&4\\6&8\end{pmatrix}.
\]
各元素的最大公因子为 $2$，行列式绝对值为 $8$，所以 $d_1=2,d_2=4$，得到
\[
\mathbb Z^2/B\mathbb Z^2\cong\mathbb Z/2\mathbb Z\oplus\mathbb Z/4\mathbb Z.
\]
行列式只给群的阶 $8$；一阶子式进一步区分它与循环群 $\mathbb Z/8\mathbb Z$。

对算子 $A\in M_n(K)$，其 $K[t]$ 模恰有关系矩阵 $tI-A$：关系 $te_j=Ae_j$ 将每个多项式向量归约为常向量，而映射 $\sum_j f_j(t)e_j\mapsto\sum_j f_j(A)e_j$ 说明没有额外关系。因此
\[
V_A\cong K[t]^n/(tI-A)K[t]^n.
\]
所以 $tI-A$ 的 Smith 非单位因子就是 $A$ 的不变因子。以上一节的 $A=\operatorname{diag}(2,2,3)$ 为例，
\[
\Delta_1=1,\qquad\Delta_2=t-2,\qquad
\Delta_3=(t-2)^2(t-3),
\]
从而 Smith 对角元为 $1,t-2,(t-2)(t-3)$，删除单位项后恢复同一组不变因子。

若最小多项式在 $K$ 上分裂，将每个不变因子按不同根分解，再用中国剩余定理，便得到形如 $K[t]/((t-\lambda)^m)$ 的分量。在基
\[
(t-\lambda)^{m-1},\ (t-\lambda)^{m-2},\ldots,1
\]
下，乘 $t$ 的矩阵是对角元为 $\lambda$、上副对角元为 $1$ 的 Jordan 块 $J_m(\lambda)$。因此 Jordan 形式来自不变因子的进一步互素分解。例如 $J_2(1)\oplus J_1(1)$ 的不变因子为 $t-1,(t-1)^2$，关系矩阵的 Smith 形式为 $\operatorname{diag}(1,t-1,(t-1)^2)$。

若多项式不分裂，仍可使用有理标准形。例如实矩阵 $\begin{pmatrix}0&-1\\1&0\end{pmatrix}$ 对应 $\mathbb R[t]/(t^2+1)$；它在实数上保持这个二阶块，扩张到复数后才分为特征值 $i,-i$ 的两个一阶块。是否添加根，决定的是可用的分解尺度。

\subsection{域扩张、最小多项式与乘法矩阵}

\label{app:algebra-extension}

设 $K$ 是域 $L$ 的子域，$\alpha\in L$。若有非零 $f\in K[t]$ 满足 $f(\alpha)=0$，称 $\alpha$ 在 $K$ 上代数。代入映射 $K[t]\to L$ 的核是非零主理想，其中首一生成元 $m_\alpha$ 称为最小多项式。它必不可约：若分解为两个正次数因子，代入后乘积为零，域中便有某个因子先消去 $\alpha$，与次数最小矛盾。

由商结构及不可约性，
\[
K[\alpha]\cong K[t]/(m_\alpha)
\]
本身已经是域，故等于包含 $K,\alpha$ 的最小子域 $K(\alpha)$。若 $\deg m_\alpha=d$，则 $1,\alpha,\ldots,\alpha^{d-1}$ 为其 $K$ 基：带余除法给生成性，次数最小给线性无关性。所以扩张次数 $[K(\alpha):K]$ 就是 $d$。若 $\alpha$ 超越，代入映射是单射，各次幂线性无关，便没有这种有限维描述。

在这个幂基下，乘法算子 $M_\alpha:x\mapsto\alpha x$ 的矩阵为 $C(m_\alpha)$，因而最小多项式与特征多项式都等于 $m_\alpha$。这把域元素的代数关系与线性算子的消去关系直接对应起来。

\textbf{例：$\mathbb Q(\sqrt2)$ 内的乘法。} 相对于基 $1,\sqrt2$，有
\[
M_{\sqrt2}=\begin{pmatrix}0&2\\1&0\end{pmatrix},
\qquad
M_{1+\sqrt2}=\begin{pmatrix}1&2\\1&1\end{pmatrix}.
\]
第二个矩阵的特征多项式为 $t^2-2t-1$。因为 $1+\sqrt2\notin\mathbb Q$，它也就是该元素的最小多项式；令 $\beta=1+\sqrt2$，关系 $\beta^2-2\beta=1$ 立即给出
\[
\beta^{-1}=\beta-2=\sqrt2-1.
\]
对有限扩张，乘法矩阵的迹、行列式分别定义元素的域迹与域范数，它们与基无关；这里 $\operatorname{Tr}(\beta)=2$、$\operatorname{N}(\beta)=-1$。

最小多项式与特征多项式的相等依赖于所用空间。若在更大的有限扩张 $L/K$ 上考察同一 $M_\alpha$，设 $r=[L:K(\alpha)]$，选择 $L$ 的 $K(\alpha)$ 基即可把乘法算子分成 $r$ 个相同的 $C(m_\alpha)$ 块。因此
\[
m_{M_\alpha}=m_\alpha,\qquad
\chi_{M_\alpha}=m_\alpha^{\,r}.
\]
例如在 $\mathbb Q(\sqrt2,\sqrt3)$ 上乘 $\sqrt2$，最小多项式仍是 $t^2-2$，特征多项式则是 $(t^2-2)^2$。有限域中也可先用不可约多项式的商添加根，再用 Frobenius 描述根之间的对称，见\href{https://github.com/lzxcr/cmc-mathematics/blob/64916fa1fdbc3e444a652eca204d5669058ee321/src/methods/outer/01-algebra.md}{有限域的构造（原稿链接）}。

\part{课程补充训练与解答}

本编为独立的课程练习选编，所有条目均非本书新增的全国大学生数学竞赛真题。材料及修订沿用 CC BY-NC-SA 4.0，作者、许可链接与来源编号规则见编写说明。各题保留原来源链接；S类为本项目已有计算练习。题后的参考解答适合用于逐步核验，但未经过本次逐题独立审定。

\section{补充训练：线性无关、基与维数}

本节包含 26 道训练题，建议在完成相关专题真题后选择性练习。

\begin{exercise}[补充训练 L3-A-1]\label{prob:L3-A-1}
下列命题是否正确？若正确请证明，否则举出反例.
        \begin{enumerate}
            \item 若 $\alpha_1, \ldots, \alpha_m \ (m > 2)$ 线性相关，则其中每一向量都是其余向量的线性组合；

            \item 若 $\alpha_1, \ldots, \alpha_m$ 线性无关，则其中每一向量不是其余向量的线性组合，这个命题的等价命题应如何叙述？

            \item $\alpha_1, \ldots, \alpha_m \ (m > 2)$ 线性无关的充要条件是任意两个向量都线性无关；

            \item  若 $\alpha_1, \alpha_2$ 线性相关，$ \beta_1, \beta_2$ 线性相关，则 $\alpha_1 + \beta_1, \alpha_2 + \beta_2$ 也线性相关；

            \item 若 $\alpha_1, \ldots, \alpha_n$ 线性无关，则 $\alpha_1 + \alpha_2, \alpha_2 + \alpha_3, \ldots, \alpha_{n-1} + \alpha_n, \alpha_n + \alpha_1$ 也线性无关；

            \item 若 $\alpha_1, \alpha_2, \alpha_3$ 线性相关，则 $\alpha_1 + \alpha_2, \alpha_2 + \alpha_3, \alpha_3 + \alpha_1$ 也线性相关；

            \item 设 $B = \{\alpha_1, \alpha_2, \alpha_3\}$ 是 $ \mathbf{R}^3 $ 的一组基，非零向量 $\alpha_0 \in \mathbf{R}^3$，则 $\{\alpha_0 + \alpha_1, \alpha_0 + \alpha_2, \alpha_0 + \alpha_3\}$ (其中三个向量均是非零向量) 也是 $\mathbf{R}^3$ 的一组基；

            \item  设 $B = \{\alpha_1, \alpha_2\}$ 是 $\mathbf{R}^2$ 的一组基，则 $\{\alpha_1 + \alpha_2, \alpha_1 - \alpha_2\}$ 也是 $\mathbf{R}^2$ 的一组基；

            \item 一个有限维线性空间内只含有有限个子空间；

            \item  若 $W_1, W_2$ 是 $ \mathbf{R}^n $ 的两个子空间，$B_1, B_2$ 分别是 $W_1, W_2$ 的基，则存在 $ \mathbf{R}^n $ 的一组基 $B$，使得 $B \supseteq B_1 \cup B_2$。

        \end{enumerate}
\end{exercise}
\noindent 来源：\href{https://github.com/yhwu-is/Linear-Algebra-Left-Undone/blob/new/%E8%AE%B2%E4%B9%89/%E4%B8%93%E9%A2%98/3%20%E6%9C%89%E9%99%90%E7%BB%B4%E7%BA%BF%E6%80%A7%E7%A9%BA%E9%97%B4.tex}{L3-A-1}。
\begin{solve}
\begin{enumerate}
                \item 错. 反例：$\alpha_1=(1,0),\alpha_2=(2,0),\alpha_3=(0,1)$，则 $\alpha_1,\alpha_2,\alpha_3$ 线性相关而 $\alpha_3$ 不是 $\alpha_1.\alpha_2$ 的线性组合.

                \item 对. 该命题的等价命题（逆否命题）是：若存在一个向量是其余向量的线性组合，则 $\alpha_1,\ldots,\alpha_m$ 线性相关. 这正是原讲义定理 3.1 的内容，因而成立.

                \item 错. 反例：$\alpha_1=(1,0),\alpha_2=(0,1),\alpha_3=(1,1)$，则 $\alpha_1,\alpha_2,\alpha_3$ 两两无关，而三者线性相关. 可证两两无关是向量组无关的必要条件.

                \item 错. 反例：$\alpha_1=(1,0),\alpha_2=(0,0),\beta_1=(0,0),\beta_2=(0,1)$，有 $\alpha_1,\alpha_2$ 相关，$\beta_1,\beta_2$ 相关，而 $\alpha_1+\beta_1$ 与 $\alpha_2+\beta_2$ 线性无关.

                \item 错. 若 $\alpha_1,\ldots,\alpha_n$ 线性无关，有
                    \[\lambda_1\alpha_1+\cdots+\lambda_n\alpha_n=0\implies\lambda_1=\lambda_2=\cdots=\lambda_n=0,\]
                    判断 $\alpha_1+\alpha_2,\alpha_2+\alpha_3,\ldots,\alpha_n+\alpha_1$ 是否无关. 设
                    \[\lambda_1'(\alpha_1+\alpha_2)+\cdots+\lambda_n'(\alpha_n+\alpha_1)=0,\]
                    则
                    \[(\lambda_n'+\lambda_1')\alpha_1+(\lambda_1'+\lambda_2')\alpha_2+\cdots+(\lambda_{n-1}'+\lambda_{n}')\alpha_n=0,\]
                    则
                    \[\begin{cases} \begin{aligned}
                                \lambda_n'+\lambda_1'       & = 0               \\
                                                            & \vdotswithin{ = } \\
                                \lambda_{n-1}'+\lambda_{n}' & = 0               \\
                            \end{aligned} \end{cases}.\]
                    解该方程可得 $\lambda_n'=(-1)^{n-1}\lambda_1'$，因此当 $n$ 为偶数时，上述方程组有非零解，则向量组相关，而当 $n$ 为奇数时，向量组无关. 综上，该命题不成立.

                \item 对. 由原讲义定理 3.1，不妨设 $\alpha_3$ 可由 $\alpha_1,\alpha_2$ 线性表示，则 $\alpha_1+\alpha_2,\alpha_2+\alpha_3,\alpha_3+\alpha_1$ 均可由 $\alpha_1,\alpha_2$ 线性表示，再由原讲义定理 3.3 可知，$\alpha_1+\alpha_2,\alpha_2+\alpha_3,\alpha_3+\alpha_1$ 线性相关.

                \item 错. 反例：取 $\alpha_0=\alpha_1-\alpha_2-\alpha_3$，则 $\alpha_0+\alpha_1=(\alpha_0+\alpha_2)+(\alpha_0+\alpha_3)$，三者线性相关，不是 $\mathbf{R}^3$ 的基.

                \item 对. 判断 $\alpha_1+\alpha_2$ 与 $\alpha_1-\alpha_2$ 是否无关.
                    \[\lambda_1(\alpha_1+\alpha_2)+\lambda_2(\alpha_1-\alpha_2)=0,\]
                    则有$(\lambda_1+\lambda_2)\alpha_1+(\lambda_1-\lambda_2)\alpha_2=0$，则 $\lambda_1+\lambda_2=0,\lambda_1-\lambda_2=0\implies \lambda_1=\lambda_2=0$，因此线性无关且个数等于维数，是一组基.

                \item 错. 反例：$\mathbf{R}^2$ 中过原点的直线 $L_0$ 是 $\mathbf{R}^2$ 的一个子空间. 显然这样的直线有无数条.

                \item 错. 反例：$\mathbf{R}^3$ 中，子空间 $W_1=\operatorname{span}(e_1,e_2)$，$W_2=\operatorname{span}(e_1+e_2,e_3)$，则 $B_1\cup B_2=\{e_1,e_2,e_3,e_1+e_2\}$，显然 $\mathbf{R}^3$ 中的任一组基都不可能包含四个元素.
            \end{enumerate}
\end{solve}

\begin{exercise}[补充训练 L3-A-2]\label{prob:L3-A-2}
证明：如果向量组线性相关，把每个向量去掉$m$个位置一致的分量，得到的缩短组仍线性相关；如果向量组线性无关，把每个向量添加$m$个位置一致的分量，得到的加长组仍线性无关.
\end{exercise}
\noindent 来源：\href{https://github.com/yhwu-is/Linear-Algebra-Left-Undone/blob/new/%E8%AE%B2%E4%B9%89/%E4%B8%93%E9%A2%98/3%20%E6%9C%89%E9%99%90%E7%BB%B4%E7%BA%BF%E6%80%A7%E7%A9%BA%E9%97%B4.tex}{L3-A-2}。
\begin{solve}
\begin{enumerate}
                \item 若向量组线性相关，则对应该方程组有无穷多解. 去掉 $m$ 个分量，相当于删去该方程组中的任意 $m$ 行方程，依然有无穷多解. 这是因为对于原方程组的任意一个解，将其带入被削减后的方程组也依然成立. 故线性相关得证.

                \item 若向量组线性无关，对应原方程组仅有唯一解，也就是全零解. 增加 $m$ 个分量相当于增加 $m$ 个方程，依然只有唯一解，因为若出现非零解，代入原方程组对应的方程中不会成立，矛盾. 故线性无关得证.
            \end{enumerate}
\end{solve}

\begin{exercise}[补充训练 L3-A-3]\label{prob:L3-A-3}
$a$取何值时，$\beta_1=(1,3,6,2)^\mathrm{T},\beta_2=(2,1,2,-1)^\mathrm{T},\beta_3=(1,-1,a,-2)^\mathrm{T}$线性无关？
\end{exercise}
\noindent 来源：\href{https://github.com/yhwu-is/Linear-Algebra-Left-Undone/blob/new/%E8%AE%B2%E4%B9%89/%E4%B8%93%E9%A2%98/3%20%E6%9C%89%E9%99%90%E7%BB%B4%E7%BA%BF%E6%80%A7%E7%A9%BA%E9%97%B4.tex}{L3-A-3}。
\begin{solve}
方程组：$x_1\beta_1+x_2\beta_2+x_3\beta_3=0$ 系数矩阵
            \[\begin{pmatrix}
                    1 & 2  & 1  \\
                    3 & 1  & -1 \\
                    6 & 2  & a  \\
                    2 & -1 &
                    -2\end{pmatrix}\rightarrow\begin{pmatrix}
                    1 & 2   & 1   \\
                    0 & -5  & -4  \\
                    0 & -10 & a-6 \\
                    0 & -5  & -4
                \end{pmatrix}\rightarrow\begin{pmatrix}
                    1 & 2  & 1   \\
                    0 & -5 & -4  \\
                    0 & 0  & a+2 \\
                    0 & 0  & 0
                \end{pmatrix},\]
            仅全零解的条件是 $a\neq-2$，此时向量组线性无关.
\end{solve}

\begin{exercise}[补充训练 L3-A-4]\label{prob:L3-A-4}
设$\alpha_1,\alpha_2,\ldots,\alpha_n\in\mathbf{F}^n$. 证明：$\alpha_1,\alpha_2,\ldots,\alpha_n$线性无关的充要条件是$\mathbf{F}^n$中任一向量都可以由它们线性表示.
\end{exercise}
\noindent 来源：\href{https://github.com/yhwu-is/Linear-Algebra-Left-Undone/blob/new/%E8%AE%B2%E4%B9%89/%E4%B8%93%E9%A2%98/3%20%E6%9C%89%E9%99%90%E7%BB%B4%E7%BA%BF%E6%80%A7%E7%A9%BA%E9%97%B4.tex}{L3-A-4}。
\begin{solve}
\begin{enumerate}
                \item 必要性：$\alpha_1,\ldots,\alpha_n$ 线性无关，对于 $F^n$ 中的任一向量 $\beta$, $\alpha_1,\ldots,\alpha_n,\beta$ 的向量个数大于维数 $n$，则线性相关. 由原讲义定理 3.2，$\beta$ 可被 $\alpha_1,\ldots,\alpha_n$ 唯一表示.

                \item 充分性：由于 $F^n$ 中任意向量均可被 $\alpha_1,\ldots,\alpha_n$ 线性表示，并且向量个数等于维数. 则 $\alpha_1,\ldots,\alpha_n$ 是 $F^n$ 的一组基. 则 $\alpha_1,\ldots,\alpha_n$ 线性无关.

                      $^*$ 更详细的证明：对于 $F^n$ 的一组基 $e_1,\ldots,e_n$，其可被 $\alpha_1,\ldots,\alpha_n$ 表示. 若 $\alpha_1,\ldots,\alpha_n$ 线性相关，不妨设 $\alpha_n$ 可被 $\alpha_1,\ldots,\alpha_{n-1}$ 表示，则有 $e_1,\ldots,e_n$ 可被 $\alpha_1,\ldots,\alpha_{n-1}$ 表示. 由于 $e_1,\ldots,e_n$ 线性无关. 根据原讲义定理 3.3，$n\leqslant n-1$，矛盾. 因此得证.
            \end{enumerate}
\end{solve}

\begin{exercise}[补充训练 L3-A-5]\label{prob:L3-A-5}
设$S_1=\{\alpha_1,\ldots,\alpha_s\},S_2=\{\beta_1,\ldots,\beta_t\}$是向量空间$V$的两个线性无关的子集，证明：$\alpha_1,\ldots,\alpha_s,\beta_1,\ldots,\beta_t$线性无关$\iff \operatorname{span}(S_1)\cap \operatorname{span}(S_2)=\{0\}$.
\end{exercise}
\noindent 来源：\href{https://github.com/yhwu-is/Linear-Algebra-Left-Undone/blob/new/%E8%AE%B2%E4%B9%89/%E4%B8%93%E9%A2%98/3%20%E6%9C%89%E9%99%90%E7%BB%B4%E7%BA%BF%E6%80%A7%E7%A9%BA%E9%97%B4.tex}{L3-A-5}。
\begin{solve}
\begin{enumerate}
                \item 必要性：对于 $\forall v\in\operatorname{span}(S_1) \cap \operatorname{span}(S_2)$ 有 $v=a_1\alpha_1+\cdots+a_s\alpha_s=b_1\beta_1+\cdots+b_t\beta_t$. 由于 $\alpha_1,\ldots,\alpha_s,\beta_1,\ldots,\beta_t$ 线性无关 $\implies a_1=\cdots=a_s=b_1=\cdots=b_t=0$，则 $v=0$，即 $\operatorname{span}(S_1)\cap\operatorname{span}(S_2)=\{0\}$.

                \item 充分性：考虑反证法. 如果 $\alpha_1,\ldots,\alpha_s,\beta_1,\ldots,\beta_t$ 线性相关，则存在不全为零的系数使得 $a_1\alpha_1+\cdots+a_s\alpha_s+b_1\beta_1\cdots+b_t\beta_t=0$. 因此存在一个向量 $v=a_1\alpha_1+\cdots+a_s\alpha_s=-(b_1\beta_1\cdots+b_t\beta_t)\neq 0$ 且 $v\in\operatorname{span}(S_1),v\in\operatorname{span}(S_2)$. 即存在非零向量 $v$ 属于 $\operatorname{span}(S_1),v\in\operatorname{span}(S_2)$，矛盾！则充分性得证.
            \end{enumerate}
\end{solve}

\begin{exercise}[补充训练 L3-A-7]\label{prob:L3-A-7}
已知$\alpha_1=(1,2,4,3),\alpha_2=(1,-1,-6,6),\alpha_3=(-2,-1,2,-9),\alpha_4=(1,1,-2,7),\beta=(4,2,4,a)$.
        \begin{enumerate}
            \item 求子空间$\operatorname{span}(\alpha_1,\alpha_2,\alpha_3,\alpha_4)$的维数和一组基；

            \item 求$a$的值使得$\beta\in W$，并求$\beta$在 (1) 所选基下的坐标.
        \end{enumerate}
\end{exercise}
\noindent 来源：\href{https://github.com/yhwu-is/Linear-Algebra-Left-Undone/blob/new/%E8%AE%B2%E4%B9%89/%E4%B8%93%E9%A2%98/3%20%E6%9C%89%E9%99%90%E7%BB%B4%E7%BA%BF%E6%80%A7%E7%A9%BA%E9%97%B4.tex}{L3-A-7}。
\begin{solve}
\begin{enumerate}
                \item 也就是求 $\alpha_1,\alpha_2,\alpha_3,\alpha_4$ 的极大线性无关组. 利用讲义中所述求法：方程组
                      \[ a_1\alpha_1+\cdots+a_4\alpha_4=0 \]
                      对应系数矩阵 
\[\begin{pmatrix}
                              1 & 1  & -2 & 1  \\
                              2 & -1 & -1 & 1  \\
                              4 & -6 & 2  & -2 \\
                              3 & 6  & -9 & 7\end{pmatrix}\]
. 化简为行阶梯型：
\[\begin{pmatrix}
                              1 & 1  & -2 & 1  \\
                              0 & -3 & 3  & -1 \\
                              0 & 0  & 0  & 3  \\
                              0 & 0  & 0  & 0\end{pmatrix}\]
，因此 $\alpha_1,\alpha_2,\alpha_3,\alpha_4$ 有非零解，这四个向量线性相关. （其实此处已知矩阵秩为 3，即维数是 3）.

                      再选取 $\alpha_1,\alpha_2,\alpha_4$ 来求解方程 $a_1\alpha_1+a_2\alpha_2+a_4\alpha_4=0$：
                      \[\begin{pmatrix}
                              1 & 1  & 1  \\
                              2 & -1 & 1  \\
                              4 & -6 & -2 \\
                              3 & 6  & 7\end{pmatrix}\rightarrow\begin{pmatrix}
                              1 & 1  & 1  \\
                              0 & -3 & -2 \\
                              0 & 0  & 2  \\
                              0 & 0  & 0\end{pmatrix},\]
                      因此该方程组只有全零解，即 $\alpha_1,\alpha_2,\alpha_4$ 是$\alpha_1,\alpha_2,\alpha_3,\alpha_4$的极大线性无关组. 则 $\operatorname{span}(\alpha_1,\alpha_2,\alpha_4)$ 的维数是 3. 一组基是 $\alpha_1,\alpha_2,\alpha_4$ .

                \item 也就是 $a_1\alpha_1+\cdots+a_4\alpha_4=\beta$有解：增广矩阵 
\[\begin{pmatrix}
                              1 & 1  & -2 & 1  & 4 \\
                              2 & -1 & -1 & 1  & 2 \\
                              4 & -6 & 2  & -2 & 4 \\
                              3 & 6  & -9 & 7  & a\end{pmatrix}\]
化为 
\[\begin{pmatrix}
                              1 & 1  & -2 & 1         & 4   \\
                              0 & -3 & 3  & -1        & -6  \\
                              0 & 0  & 0  & -\frac 83 & 8   \\
                              0 & 0  & 0  & 0         & a-9\end{pmatrix}\]
，若方程有解，则 $a=9$. 求坐标，取 $x_3=0$ 代入得 $\beta=4\alpha_1+3\alpha_2-3\alpha_4$.
            \end{enumerate}
\end{solve}

\begin{exercise}[补充训练 L3-A-9]\label{prob:L3-A-9}
证明：$B=\{1,x-a,(x-a)^2\}\enspace(a\neq 0)$是$\mathbf{R}[x]_3$的一组基，并求$\mathbf{R}[x]_3$的自然基$\{1,x,x^2\}$中每个向量关于基$B$的坐标.
\end{exercise}
\noindent 来源：\href{https://github.com/yhwu-is/Linear-Algebra-Left-Undone/blob/new/%E8%AE%B2%E4%B9%89/%E4%B8%93%E9%A2%98/3%20%E6%9C%89%E9%99%90%E7%BB%B4%E7%BA%BF%E6%80%A7%E7%A9%BA%E9%97%B4.tex}{L3-A-9}。
\begin{solve}
把这三个多项式关于自然基的坐标作为列，得到
\[P=\begin{pmatrix}1&-a&a^2\\0&1&-2a\\0&0&1\end{pmatrix}.\]
由于 $\det P=1$，三个多项式线性无关，因而构成三维空间的一组基。实际上 $a=0$ 时结论也成立。由
\[1=1,\qquad x=a+(x-a),\qquad x^2=a^2+2a(x-a)+(x-a)^2\]
得到所求坐标依次为 $(1,0,0)^{\mathrm T}$、$(a,1,0)^{\mathrm T}$、$(a^2,2a,1)^{\mathrm T}$。
\end{solve}

\begin{exercise}[补充训练 L3-A-10]\label{prob:L3-A-10}
已知向量组$A=\{\alpha_1,\alpha_2,\alpha_3\},\enspace B=\{\alpha_1,\alpha_2,\alpha_3,\alpha_4\},\enspace C=\{\alpha_1,\alpha_2,\alpha_3,\alpha_5\}$的秩分别为$r(A)=r(B)=3,\enspace r(C)=4$. 证明：$\{\alpha_1,\alpha_2,\alpha_3,\alpha_5-\alpha_4\}$的秩为4.
\end{exercise}
\noindent 来源：\href{https://github.com/yhwu-is/Linear-Algebra-Left-Undone/blob/new/%E8%AE%B2%E4%B9%89/%E4%B8%93%E9%A2%98/3%20%E6%9C%89%E9%99%90%E7%BB%B4%E7%BA%BF%E6%80%A7%E7%A9%BA%E9%97%B4.tex}{L3-A-10}。
\begin{solve}
等价于证明 $\alpha_1,\alpha_2,\alpha_3,\alpha_5-\alpha_4$ 线性无关. 即求解
            \begin{equation}
                \lambda_1\alpha_1+\lambda_2\alpha_2+\lambda_3\alpha_3+\lambda_4(\alpha_5-\alpha_4)=0. \tag{*} \label{eq:3:A.10}
            \end{equation}

            由于 $r(A)=r(B)=3$ 可得 $A$ 线性无关. $B$ 线性相关. 由原讲义定理 3.2 得 $\alpha_4$ 可由 $\alpha_1,\alpha_2,\alpha_3$ 唯一表示：$\alpha_4=\mu_1\alpha_1+\mu_2\alpha_2+\mu_3\alpha_3$. 则代入 式~\eqref{eq:3:A.10}. 有
            \[(\lambda_1-\mu_1\lambda_4)\alpha_1+(\lambda_2-\mu_2\lambda_4)\alpha_2+(\lambda_3-\mu_3\lambda_4)\alpha_3+\lambda_4\alpha_5=0,\]
            因为 $r(C)=4$，$\alpha_1,\alpha_2,\alpha_3,\alpha_5$ 线性无关. 有 $\lambda_4=0,\lambda_1=\mu_1\lambda_4=0,\lambda_2=\mu_2\lambda_4=0,\lambda_3=\mu_3\lambda_4=0$. 故 $\alpha_1,\alpha_2,\alpha_3,\alpha_5-\alpha_4$ 线性无关. 原题得证.
\end{solve}

\begin{exercise}[补充训练 L3-A-11]\label{prob:L3-A-11}
设向量组$\alpha_1,\alpha_2,\ldots,\alpha_s$的秩为$r$. 在其中任取$m$个向量$\alpha_{i1},\alpha_{i2},\ldots,\alpha_{im}$，证明：向量组$\alpha_{i1},\alpha_{i2},\ldots,\alpha_{im}$的秩$\geqslant r+m-s$.
\end{exercise}
\noindent 来源：\href{https://github.com/yhwu-is/Linear-Algebra-Left-Undone/blob/new/%E8%AE%B2%E4%B9%89/%E4%B8%93%E9%A2%98/3%20%E6%9C%89%E9%99%90%E7%BB%B4%E7%BA%BF%E6%80%A7%E7%A9%BA%E9%97%B4.tex}{L3-A-11}。
\begin{solve}
相当于从 $\alpha_1,\ldots,\alpha_s$ 向量中选取 $s-m$ 个向量丢弃，剩余向量的秩：
            \[r(\alpha_{i1},\ldots,\alpha_{im})\geqslant r-(s-m) =r+m-s.\]
\end{solve}

\begin{exercise}[补充训练 L3-B-1]\label{prob:L3-B-1}
已知$\alpha_1\neq 0$，则$\alpha_1,\alpha_2,\ldots,\alpha_n$线性相关的充要条件是存在$i\enspace(2 \leqslant i \leqslant n)$使得$\alpha_i$可由$\alpha_1,\alpha_2,\ldots,\alpha_{i-1}$线性表示，且表示法唯一.
\end{exercise}
\noindent 来源：\href{https://github.com/yhwu-is/Linear-Algebra-Left-Undone/blob/new/%E8%AE%B2%E4%B9%89/%E4%B8%93%E9%A2%98/3%20%E6%9C%89%E9%99%90%E7%BB%B4%E7%BA%BF%E6%80%A7%E7%A9%BA%E9%97%B4.tex}{L3-B-1}。
\begin{solve}
充分性显然成立，下证必要性：由于 $\alpha_1,\ldots,\alpha_n$ 线性相关，则存在 $m$，其能使得 $\alpha_1,\ldots,\alpha_m$ 线性无关的最大下标，有 $1\leqslant m<n$. 因此 $i=m+1$，$\alpha_1,\ldots,\alpha_{i-1}$ 线性无关，$\alpha_1,\ldots,\alpha_i$ 线性相关. 可得 $\alpha_i$ 可被 $\alpha_1,\ldots,\alpha_{i-1}$ 唯一表示.
\end{solve}

\begin{exercise}[补充训练 L3-B-2]\label{prob:L3-B-2}
证明以下两个结论：
        \begin{enumerate}
            \item 设$U$和$W$都是$V$的非零子空间，如果$U\subseteq W$，那么$\dim U \leqslant \dim W$；

            \item 设$U$和$W$都是$V$的非零子空间，$U\subseteq W$，且$\dim U = \dim W$，则$U = W$.
        \end{enumerate}
\end{exercise}
\noindent 来源：\href{https://github.com/yhwu-is/Linear-Algebra-Left-Undone/blob/new/%E8%AE%B2%E4%B9%89/%E4%B8%93%E9%A2%98/3%20%E6%9C%89%E9%99%90%E7%BB%B4%E7%BA%BF%E6%80%A7%E7%A9%BA%E9%97%B4.tex}{L3-B-2}。
\begin{solve}
\begin{enumerate}
                \item 设 $U$ 的一组基为 $u_1,\ldots,u_m$，$W$ 的一组基为 $w_1,\ldots,w_n$. 由于 $U\subseteq W$，则 $u_1,\ldots,u_m$ 可由 $w_1,\ldots,w_n$ 线性表示，且 $u_1,\ldots,u_m$ 线性无关. 由原讲义定理 3.3 的等价命题可得 $m\leqslant n$，则 $m=\dim U\leqslant\dim W=n$ 的得证.

                \item 因为 $\dim U=\dim W$，则 $u_1,\ldots,u_m$ 也是 $W$ 的一组基. 则 $W$ 的任意向量均可由 $u_1,\ldots,u_m$ 表示，可得 $W\subseteq U$，而 $U\subseteq W$，故有 $U=W$ 得证.
            \end{enumerate}
\end{solve}

\begin{exercise}[补充训练 L3-B-3]\label{prob:L3-B-3}
设向量组$\alpha_1,\alpha_2,\ldots,\alpha_n$线性无关. 证明：在向量组$\beta,\alpha_1,\alpha_2,\ldots,\alpha_n$中至多有一个向量$\alpha_i\enspace(1 \leqslant i \leqslant n)$可被其前面的$i$个向量$\beta,\alpha_1,\alpha_2,\ldots,\alpha_{i-1}$线性表示.
\end{exercise}
\noindent 来源：\href{https://github.com/yhwu-is/Linear-Algebra-Left-Undone/blob/new/%E8%AE%B2%E4%B9%89/%E4%B8%93%E9%A2%98/3%20%E6%9C%89%E9%99%90%E7%BB%B4%E7%BA%BF%E6%80%A7%E7%A9%BA%E9%97%B4.tex}{L3-B-3}。
\begin{solve}
若 $\alpha_i$ 能由其前面的向量表示，则表达式中 $\beta$ 的系数不能为零，否则 $\alpha_1,\ldots,\alpha_n$ 线性相关。因此可解出
\[\beta\in\operatorname{span}(\alpha_1,\ldots,\alpha_i).\]
若还有 $j>i$ 满足相同条件，把上式代入 $\alpha_j$ 关于 $\beta,\alpha_1,\ldots,\alpha_{j-1}$ 的表示，便得到 \[\alpha_j\in\operatorname{span}(\alpha_1,\ldots,\alpha_{j-1}),\]与已知线性无关矛盾。因此这样的下标至多有一个。
\end{solve}

\begin{exercise}[补充训练 L3-B-4]\label{prob:L3-B-4}
证明：$1,e^{\lambda_1\cdot x},e^{\lambda_2\cdot x}$（$\lambda_1\neq\lambda_2$且均不为0）线性无关.
\end{exercise}
\noindent 来源：\href{https://github.com/yhwu-is/Linear-Algebra-Left-Undone/blob/new/%E8%AE%B2%E4%B9%89/%E4%B8%93%E9%A2%98/3%20%E6%9C%89%E9%99%90%E7%BB%B4%E7%BA%BF%E6%80%A7%E7%A9%BA%E9%97%B4.tex}{L3-B-4}。
\begin{solve}
设 $k_0+k_1e^{\lambda_1x}+k_2e^{\lambda_2x}=0$ 对所有实数 $x$ 成立。求导一次、两次，并在 $x=0$ 取值，得到
\[\begin{pmatrix}\lambda_1&\lambda_2\\\lambda_1^2&\lambda_2^2\end{pmatrix}\begin{pmatrix}k_1\\k_2\end{pmatrix}=0.\]
系数行列式为 $\lambda_1\lambda_2(\lambda_2-\lambda_1)\ne0$，故 $k_1=k_2=0$；代回恒等式得 $k_0=0$。于是三个函数线性无关。
\end{solve}

\begin{exercise}[补充训练 L3-B-5]\label{prob:L3-B-5}
设线性空间$V(\mathbf{F})$中，向量$\beta$是$\alpha_1,\ldots,\alpha_r$的线性组合，但不是$\alpha_1,\ldots,\alpha_{r-1}$的线性组合. 证明：$\operatorname{span}(\alpha_1,\ldots,\alpha_{r-1},\alpha_r)=\operatorname{span}(\alpha_1,\ldots,\alpha_{r-1},\beta)$.
\end{exercise}
\noindent 来源：\href{https://github.com/yhwu-is/Linear-Algebra-Left-Undone/blob/new/%E8%AE%B2%E4%B9%89/%E4%B8%93%E9%A2%98/3%20%E6%9C%89%E9%99%90%E7%BB%B4%E7%BA%BF%E6%80%A7%E7%A9%BA%E9%97%B4.tex}{L3-B-5}。
\begin{solve}
只需证明 $\alpha_r$ 可以被 $\alpha_1,\ldots,\alpha_{r-1},\beta$ 表示即可. 由于 $\beta$ 是 $\alpha_1,\ldots,\alpha_r$ 的线性组合，若 $\lambda_r=0$，则 $\beta$ 是 $\alpha_1,\ldots,\alpha_{r-1}$ 的线性组合. 这与条件矛盾. 因此 \[\alpha_r=-\dfrac 1{\lambda_r}(\lambda_1\alpha_1+\cdots+\lambda_{r-1}\alpha_{r-1}-\beta),\] 则这两组向量等价，即 $\operatorname{span}(\alpha_1,\ldots,\alpha_{r-1},\alpha_r)=\operatorname{span}(\alpha_1,\ldots,\alpha_{r-1},\beta)$，得证.
\end{solve}

\begin{exercise}[补充训练 L3-B-6]\label{prob:L3-B-6}
\label{item:3:正实数线性空间}
        设$\mathbf{R}^+$是所有正实数组成的集合，加法和数乘定义如下：
        \[ \forall a,b \in \mathbf{R}_+,\enspace k\in \mathbf{R}\colon\enspace a\oplus b = ab,\enspace k\odot a = a^k \]
        则 $\mathbf{R}^+$关于这一加法和数乘构成一个实线性空间. 求$\mathbf{R}^+$的一组基.
\end{exercise}
\noindent 来源：\href{https://github.com/yhwu-is/Linear-Algebra-Left-Undone/blob/new/%E8%AE%B2%E4%B9%89/%E4%B8%93%E9%A2%98/3%20%E6%9C%89%E9%99%90%E7%BB%B4%E7%BA%BF%E6%80%A7%E7%A9%BA%E9%97%B4.tex}{L3-B-6}。
\begin{solve}
分析该实线性空间，可以看出加法单位元为 1，数乘单位元为 1. 我们给出一组基：$e$，其中 $e$ 为自然对数的底数. 当然, 2，3 或者 10 都可以作为一组基. 接下来我们验证 $e$ 是 $\mathbf{R}^+$ 的基：$\forall a\in \mathbf{R}^+,\exists k=\ln a\in\mathbf{R}$，满足 $k\odot e=e^k=a$，则 $\operatorname{span}(e)=\mathbf{R}^+$ 成立. 由于该向量组只有一个元素，且并非设该空间的零元 1，则 $e$ 是线性无关的. 得证.
\end{solve}

\begin{exercise}[补充训练 L3-C-1]\label{prob:L3-C-1}
已知$m$个向量$\alpha_1,\alpha_2,\ldots,\alpha_m$线性相关，但其中任意$m-1$个都线性无关，证明：
        \begin{enumerate}
            \item 若$k_1\alpha_1+\cdots+k_m\alpha_m=0$，则$k_1,\ldots,k_m$全为0或全不为0；

            \item 若以下等式成立
                  \begin{align*}
                      k_1\alpha_1+\cdots+k_m\alpha_m & =0 \\
                      l_1\alpha_1+\cdots+l_m\alpha_m & =0
                  \end{align*}
                  其中$l_1\neq 0$，证明：$\dfrac{k_1}{l_1}=\cdots=\dfrac{k_m}{l_m}$.
        \end{enumerate}
\end{exercise}
\noindent 来源：\href{https://github.com/yhwu-is/Linear-Algebra-Left-Undone/blob/new/%E8%AE%B2%E4%B9%89/%E4%B8%93%E9%A2%98/3%20%E6%9C%89%E9%99%90%E7%BB%B4%E7%BA%BF%E6%80%A7%E7%A9%BA%E9%97%B4.tex}{L3-C-1}。
\begin{solve}
\begin{enumerate}
                \item 若 不全为 0. 不妨设设至少有 $k_i=0$，则有 $k_1\alpha_1+\cdots+k_{i-1}\alpha_{i-1}+k_{i+1}\alpha_{i+1}+\cdots+k_m\alpha_m=0$，并且系数不全为 0. 因此 $\alpha_1,\ldots,\alpha_{i-1},\alpha_{i+1},\ldots,\alpha_m$ 这 $m-1$ 个向量相关，与题设矛盾. 则原题得证.

                \item $l_1\neq 0$，则 $l_2,\ldots,l_m$ 均不为 0.
                      \begin{enumerate}
                          \item 若 $k_1=\cdots=k_m=0$. 原式显然成立.

                          \item 若 $k_1,\ldots,k_m$ 不全为 0. 则
                                \begin{gather*}
                                    l_1(k_1\alpha_1+\cdots+k_m\alpha_m)=0, \\
                                    k_1(l_1\alpha_1+\cdots+l_m\alpha_m)=0.
                                \end{gather*}
                                两式相减，得
                                \[(k_2l_1-k_1l_2)\alpha_2+\cdots+(k_ml_1-k_1l_m)\alpha_m=0,\]
                                因为 $\alpha_2,\ldots,\alpha_m$ 线性无关. 则以上系数均为 0. 故$\dfrac {k_2}{l_2}=\dfrac {k_1}{l_1},\ldots,\dfrac {k_m}{l_m}=\dfrac {k_1}{l_1}$. 得证.
                      \end{enumerate}
            \end{enumerate}
\end{solve}

\begin{exercise}[补充训练 L3-C-2]\label{prob:L3-C-2}
（替换定理）设$\alpha_1,\alpha_2,\ldots,\alpha_r$线性无关，且可以被$\beta_1,\beta_2,\ldots,\beta_n$线性表示，则可以将$\beta_1,\beta_2,\ldots,\beta_n$中的$r$个向量替换成$\alpha_1,\alpha_2,\ldots,\alpha_r$后得到与$\beta_1,\beta_2,\ldots,\beta_n$等价的新向量组（注：可以使用数学归纳法证明）.
\end{exercise}
\noindent 来源：\href{https://github.com/yhwu-is/Linear-Algebra-Left-Undone/blob/new/%E8%AE%B2%E4%B9%89/%E4%B8%93%E9%A2%98/3%20%E6%9C%89%E9%99%90%E7%BB%B4%E7%BA%BF%E6%80%A7%E7%A9%BA%E9%97%B4.tex}{L3-C-2}。
\begin{solve}
采用逐个替换的证明。开始时 $\beta_1,\ldots,\beta_n$ 张成一个包含全部 $\alpha_j$ 的空间。假设已经把其中 $k$ 个向量替换为 $\alpha_1,\ldots,\alpha_k$，且张成空间不变，其中 $0\le k<r$。将 $\alpha_{k+1}$ 写成当前生成组的线性组合。这个表示中至少有一个尚未替换的 $\beta_j$ 的系数非零，否则 $\alpha_{k+1}$ 在 $\alpha_1,\ldots,\alpha_k$ 的张成空间中，违背 $\alpha_1,\ldots,\alpha_r$ 线性无关。解出该 $\beta_j$，用 $\alpha_{k+1}$ 替换它，张成空间仍不变。重复 $r$ 次，就得到由全部 $\alpha_1,\ldots,\alpha_r$ 和原组中 $n-r$ 个 $\beta_j$ 组成的等价向量组。这里同时说明 $r\le n$，否则替换过程中会在放入全部 $\alpha_j$ 前耗尽可替换向量，与线性无关矛盾。
\end{solve}

\begin{exercise}[补充训练 L3-C-3]\label{prob:L3-C-3}
设线性空间$V=\mathbf{F}^n$. 证明：
        \begin{enumerate}
            \item 存在$V$的子空间$W$，使得$W$的任一非零向量的分量均不为0；

            \item 若$V$的子空间$W$的任一非零向量的分量均不为0，则$\dim W=1$；

            \item 若$V$的子空间$W$的任一非零向量的零分量个数均不超过$r$，则$\dim W \leqslant r+1$.
        \end{enumerate}
\end{exercise}
\noindent 来源：\href{https://github.com/yhwu-is/Linear-Algebra-Left-Undone/blob/new/%E8%AE%B2%E4%B9%89/%E4%B8%93%E9%A2%98/3%20%E6%9C%89%E9%99%90%E7%BB%B4%E7%BA%BF%E6%80%A7%E7%A9%BA%E9%97%B4.tex}{L3-C-3}。
\begin{solve}
\begin{enumerate}
                \item $\alpha=(1,1,\ldots,1)^{\mathrm{T}}, W=\operatorname{span}(\alpha)$,显然 $W$ 是满足条件的一维子空间.

                \item 考虑反证法：若 $\dim W>1$，则 $W$ 中存在线性无关的两向量. 由条件，
                      \[a_1,\ldots,a_n,b_1,\ldots,b_n\neq 0,\]
                      可设 $a_1=kb_1,k\in F$，因此 $\alpha-k\beta=(0,a_2-kb-2,\ldots,a_n-kb_n)^{\mathrm{T}}\in W$. 且由于 $\alpha,\beta$ 无关，$\alpha-k\beta\neq 0$ 但存在分量为 0，这与条件矛盾. 故 $\dim W=1$.

                \item 考虑反证法：若 $\dim W>r+1$，则存在 $r+2$ 个线性无关的向量，设为
                      \begin{align*}
                          \alpha_1     & = (a_{11},a_{12},\ldots,a_{1(r+1)},\ldots,a_{1n})^{\mathrm{T}},                 \\
                                       & \vdotswithin{ = }                                                              \\
                          \alpha_{r+2} & = (a_{(r+2)1},a_{(r+2)2},\ldots,a_{(r+2)(r+1)},\ldots,a_{(r+2)n})^{\mathrm{T}}.
                      \end{align*}
                      取这些向量的前 $r+1$ 个分量组成新的向量组：
                      \begin{align*}
                          \beta_1     & = (a_{11},a_{12},\ldots,a_{1(r+1)})^{\mathrm{T}},             \\
                                      & \vdotswithin{ = }                                            \\
                          \beta_{r+2} & = (a_{(r+2)1},a_{(r+2)2},\ldots,a_{(r+2)(r+1)})^{\mathrm{T}}.
                      \end{align*}
                      由于 $\beta_1,\ldots,\beta_{r+2}$ 是 $r+2$ 个 $r+1$ 维向量，其必然线性相关，则存在不全为 0 的系数 $\lambda_1,\ldots,\lambda_{r+2}$，$\lambda_1\beta_1+\cdots+\lambda_{r+2}\beta_{r+2}=0$.  由于 $\alpha_1,\ldots,\alpha_{r+2}$ 线性无关知：$\lambda_1\alpha_1+\cdots+\lambda_{r+2}\alpha_{r+2}\neq 0$，且其属于 $W$. 但其前 $r+1$ 个分量均为 0，这与条件矛盾. 故 $\dim W\leqslant r+1$ 得证.
            \end{enumerate}
\end{solve}

\begin{exercise}[补充训练 L3-C-4]\label{prob:L3-C-4}
设 $\mathbf{Q}(\sqrt[3]{2}) = \{a+b\sqrt[3]{2}+c\sqrt[3]{4}\mid a,b,c\in\mathbf{Q}\}$，证明 $\mathbf{Q}(\sqrt[3]{2})$ 是 $\mathbf{Q}$ 上的线性空间并求其维数.
\end{exercise}
\noindent 来源：\href{https://github.com/yhwu-is/Linear-Algebra-Left-Undone/blob/new/%E8%AE%B2%E4%B9%89/%E4%B8%93%E9%A2%98/3%20%E6%9C%89%E9%99%90%E7%BB%B4%E7%BA%BF%E6%80%A7%E7%A9%BA%E9%97%B4.tex}{L3-C-4}。
\begin{solve}
先证明 $\mathbf{Q}(\sqrt[3]{2})$ 是 $\mathbf{Q}$ 上的线性空间：

            $\mathbf{Q}(\sqrt[3]{2})$ 对通常的加法构成交换群以及其封闭性是显然的，在此不再赘述，下阐述数乘性质：
            \begin{enumerate}
                \item $\forall v = a+b\sqrt[3]{2}+c\sqrt[3]{4} \in \mathbf{Q}(\sqrt[3]{2}), \enspace 1 (a+b\sqrt[3]{2}+c\sqrt[3]{4}) = a+b\sqrt[3]{2}+c\sqrt[3]{4}$；

                \item $\forall v = a+b\sqrt[3]{2}+c\sqrt[3]{4} \in \mathbf{Q}(\sqrt[3]{2}), \forall \lambda, \mu \in \mathbf{Q}, \enspace \lambda(\mu v) = (\lambda \mu) v$；

                \item $\forall v = a+b\sqrt[3]{2}+c\sqrt[3]{4} \in \mathbf{Q}(\sqrt[3]{2}), \forall \lambda, \mu \in \mathbf{Q}, \enspace (\lambda + \mu) v = (\lambda + \mu)a + (\lambda + \mu) b\sqrt[3]{2} + (\lambda + \mu) c\sqrt[3]{4} = \lambda(a+b\sqrt[3]{2}+c\sqrt[3]{4}) + \mu(a+b\sqrt[3]{2}+c\sqrt[3]{4}) = \lambda v + \mu v$；

                \item $\forall v = a_1 + b_1\sqrt[3]{2} + c_1\sqrt[3]{4}, u = a_2 + b_2\sqrt[3]{2} + c_2\sqrt[3]{4} \in \mathbf{Q}(\sqrt[3]{2}), \forall \lambda \in \mathbf{Q}, \enspace \lambda(v+u) = \lambda (a_1 + a_2) + \lambda (b_1 + b_2)\sqrt[3]{2} + \lambda (c_1 + c_2)\sqrt[3]{4} = \lambda (a_1 + b_1\sqrt[3]{2} + c_1\sqrt[3]{4}) + \lambda (a_2 + b_2\sqrt[3]{2} + c_2\sqrt[3]{4}) = \lambda v + \lambda u$；

                \item （封闭性）$\forall v = a+b\sqrt[3]{2}+c\sqrt[3]{4} \in \mathbf{Q}(\sqrt[3]{2}), \forall \lambda \in \mathbf{Q}, \enspace \lambda v = \lambda a + \lambda b\sqrt[3]{2} + \lambda c\sqrt[3]{4} \in \mathbf{Q}(\sqrt[3]{2})$.
            \end{enumerate}
            故 $\mathbf{Q}(\sqrt[3]{2})$ 是 $\mathbf{Q}$ 上的线性空间.

            下求 $\mathbf{Q}(\sqrt[3]{2})$ 的维数. 考虑找 $\mathbf{Q}(\sqrt[3]{2})$ 的一组基，我们给出 $\{1, \sqrt[3]{2}, \sqrt[3]{4}\}$，下证其为 $\mathbf{Q}(\sqrt[3]{2})$ 的一组基：

            \begin{enumerate}
                \item （线性无关）令 $k_1 + k_2 \sqrt[3]{2} + k_3 \sqrt[3]{4} = 0$，其中 $k_1, k_2, k_3 \in \mathbf{Q}$，由于左边只有第一项为有理数，故有 $k_1 = 0$，进而有 $\sqrt[3]{2} (k_2 + k_3 \sqrt[3]{2}) = 0$，又可得到 $k_2 = 0$，并且 $k_3 = 0$. 故 $\{1, \sqrt[3]{2}, \sqrt[3]{4}\}$ 线性无关.

                \item （张成空间）由 $\mathbf{Q}(\sqrt[3]{2})$ 的定义易得.
            \end{enumerate}
            综上，$\mathbf{Q}(\sqrt[3]{2})$ 是 $\mathbf{Q}$ 上的线性空间，且 $\dim \mathbf{Q}(\sqrt[3]{2}) = 3$.
\end{solve}

\begin{exercise}[补充训练 L3-C-5]\label{prob:L3-C-5}
设$\mathbf{K} \subseteq \mathbf{F} \subseteq \mathbf{E}$是三个数域，已知$\mathbf{F}$作为$\mathbf{K}$上的线性空间是$n$维的，$\mathbf{E}$作为$\mathbf{F}$上的线性空间是$m$维的，证明：$\mathbf{E}$作为$\mathbf{K}$上的线性空间是$mn$维的.
\end{exercise}
\noindent 来源：\href{https://github.com/yhwu-is/Linear-Algebra-Left-Undone/blob/new/%E8%AE%B2%E4%B9%89/%E4%B8%93%E9%A2%98/3%20%E6%9C%89%E9%99%90%E7%BB%B4%E7%BA%BF%E6%80%A7%E7%A9%BA%E9%97%B4.tex}{L3-C-5}。
\begin{solve}
设 $f_1,\ldots,f_n$ 是 $\mathbf F$ 作为 $\mathbf K$ 空间的一组基，$e_1,\ldots,e_m$ 是 $\mathbf E$ 作为 $\mathbf F$ 空间的一组基。考虑 $mn$ 个向量 $f_je_i$（$1\le i\le m,1\le j\le n$）。若
\[\sum_{i=1}^m\sum_{j=1}^n k_{ij}f_je_i=0,\qquad k_{ij}\in\mathbf K,\]
则由 $e_i$ 在 $\mathbf F$ 上线性无关，得每个 $\sum_jk_{ij}f_j=0$；再由 $f_j$ 在 $\mathbf K$ 上线性无关，得到所有 $k_{ij}=0$。故这 $mn$ 个向量线性无关。任取 $e\in\mathbf E$，先写成 $e=\sum_i\mu_ie_i$，再把每个 $\mu_i\in\mathbf F$ 写成 $\sum_jk_{ij}f_j$，即得 $e$ 在这些乘积向量下的线性表示。它们既张成又线性无关，所以构成基，所求维数为 $mn$。
\end{solve}

\begin{exercise}[补充训练 L3-例-1]\label{prob:L3-例-1}
\begin{enumerate}[label=(\arabic*)]
        \item 判断$\mathbf{R}^3$中向量$(1,1,0),(0,1,1),(1,0,-1)$的线性相关性；

        \item 判断$\mathbf{R}^3$中向量$(1,-3,1),(-1,2,-2),(1,1,3)$的线性相关性；

        \item \label{item:3:线性相关性:3}
              判断$\mathbf{R}[x]_3$中$p_1(x)=1+x,\ p_2(x)=1-x,\ p_3(x)=x+x^2$的线性相关性；

        \item 判断连续函数全体构成的线性空间中$1,\ \sin^2x,\ \cos^2x$的线性相关性；

        \item \label{item:3:线性相关性:5}
              判断连续函数全体构成的线性空间中$1,\ 2^x,\ 2^{-x}$的线性相关性.
    \end{enumerate}
\end{exercise}
\noindent 来源：\href{https://github.com/yhwu-is/Linear-Algebra-Left-Undone/blob/new/%E8%AE%B2%E4%B9%89/%E4%B8%93%E9%A2%98/3%20%E6%9C%89%E9%99%90%E7%BB%B4%E7%BA%BF%E6%80%A7%E7%A9%BA%E9%97%B4.tex}{L3-例-1}。
\begin{solve}
\begin{enumerate}
        \item 根据定义，应求解方程
              \[\lambda_1(1,1,0) + \lambda_2(0,1,1) + \lambda_3(1,0,-1) = 0,\]
              即
              \[ \begin{cases}
                      \lambda_1 + \lambda_3 = 0 \\
                      \lambda_1 + \lambda_2 = 0 \\
                      \lambda_2 - \lambda_3 = 0
                  \end{cases} \]
              解得基础解系$k(1,-1,-1)^\mathrm{T}$，所以存在非零解，向量组线性相关.

        \item 求解方程
              \[\lambda_1(1,-3,1) + \lambda_2(-1,2,-2) + \lambda_3(1,1,3) = 0,\]
              即
              \[ \begin{cases}
                      \lambda_1 - \lambda_2 + \lambda_3 = 0    \\
                      -3\lambda_1 + 2\lambda_2 + \lambda_3 = 0 \\
                      \lambda_1 - 2\lambda_2 + 3\lambda_3 = 0
                  \end{cases} \]
              解得 $\lambda_1 = \lambda_2 = \lambda_3 = 0$，此向量组线性无关.

        \item 求解方程
              \begin{align*}
                  \lambda_1p_1(x) + \lambda_2p_2(x) + \lambda_3p_3(x)                           & = 0 \\
                  (\lambda_1 + \lambda_2) + (\lambda_1 - \lambda_2 + \lambda_3)x + \lambda_3x^2 & = 0
              \end{align*}
              所以需要求解方程组
              \[ \begin{cases}
                      \lambda_1 + \lambda_2 = 0             \\
                      \lambda_1 - \lambda_2 + \lambda_3 = 0 \\
                      \lambda_3 = 0
                  \end{cases} \]
              解得 $\lambda_1 = \lambda_2 = \lambda_3 = 0$，此向量组线性无关.

        \item 易知 $-1 + \sin^2x + \cos^2x = 0$，对应系数为 $-1, 1, 1$，不全为零，所以此向量组线性相关.

        \item 求解方程
              \[\lambda_1 + \lambda_2·2^{-x} + \lambda_3·2^x = 0\]
              很明显会发现仅凭此方程是难以求解的，方程数目不足. 注意到此方程应该对于任意的 $x$ 均成立，所以取 $x = 0, x = 1, x = -1$，得到方程组
              \[ \begin{cases}
                      \lambda_1 + \lambda_2 + \lambda_3 = 0              \\[1ex]
                      \lambda_1 + \dfrac{1}{2}\lambda_2 + 2\lambda_3 = 0 \\[1ex]
                      \lambda_1 + 2\lambda_2 + \dfrac{1}{2}\lambda_3 = 0
                  \end{cases} \]
              解得 $\lambda_1 = \lambda_2 = \lambda_3 = 0$，此向量组线性无关.
    \end{enumerate}
\end{solve}

\begin{exercise}[补充训练 L3-例-2]\label{prob:L3-例-2}
判断 $\mathbf{R}^3$ 中向量组 $\{\alpha_1, \alpha_2, \alpha_3\}$ 和 $\{\beta_1, \beta_2, \beta_3\}$ 的线性相关性, 其中

              \[
                  \alpha_1 = (1, 1, 0), \quad \alpha_2 = (0, 1, 1), \quad \alpha_3 = (1, 0, -1),
              \]
              \[
                  \beta_1 = (1, -3, 1), \quad \beta_2 = (-1, 2, -2), \quad \beta_3 = (1, 1, 3).
              \]
\end{exercise}
\noindent 来源：\href{https://github.com/yhwu-is/Linear-Algebra-Left-Undone/blob/new/%E8%AE%B2%E4%B9%89/%E4%B8%93%E9%A2%98/3%20%E6%9C%89%E9%99%90%E7%BB%B4%E7%BA%BF%E6%80%A7%E7%A9%BA%E9%97%B4.tex}{L3-例-2}。
\begin{solve}
容易看出, $\alpha_3 = \alpha_1 - \alpha_2$, 所以 $\alpha_1, \alpha_2, \alpha_3$ 线性相关. 但对 $\beta_1, \beta_2, \beta_3$ 不易看出是否有线性关系, 要按定义来判别. 设

              \begin{equation} \label{eq:线性相关性:1}
                  x_1 \beta_1 + x_2 \beta_2 + x_3 \beta_3 = 0
              \end{equation}

              即

              \[
                  x_1(1, -3, 1) + x_2(-1, 2, -2) + x_3(1, 1, 3) = (0, 0, 0).
              \]

              这个向量方程等价于以下的三个元线性齐次方程组:

              \[
                  \begin{cases}
                      x_1 - x_2 + x_3 = 0    \\
                      -3x_1 + 2x_2 + x_3 = 0 \\
                      x_1 - 2x_2 + 3x_3 = 0
                  \end{cases},
              \]

              容易解得这个方程组只有零解 $x_1 = x_2 = x_3 = 0$. 即只有全为零的 $x_1, x_2, x_3$ 才使得式~\eqref{eq:线性相关性:1} 成立, 故 $\beta_1, \beta_2, \beta_3$ 线性无关.

              一般若 $\beta_1 = (a_1, b_1, c_1), \beta_2 = (a_2, b_2, c_2), \beta_3 = (a_3, b_3, c_3)$, 则 $\beta_1, \beta_2, \beta_3$ 线性相关（线性无关）的充要条件是三元线性齐次方程组

              \[
                  \begin{cases}
                      a_1 x_1 + a_2 x_2 + a_3 x_3 = 0 \\
                      b_1 x_1 + b_2 x_2 + b_3 x_3 = 0 \\
                      c_1 x_1 + c_2 x_2 + c_3 x_3 = 0
                  \end{cases}
              \]

              有非零解（只有零解）. 用此法可得 $\mathbf{R}^3$ 中任何 4 个向量, $\mathbf{R}^n$ 中任何 $n + 1$ 个向量都线性相关.

              总的来说，我们有以下结论：

              列向量组$\alpha_1,\alpha_2,\ldots,\alpha_m$线性相关$\iff$齐次线性方程组$x_1\alpha_1+x_2\alpha_2+\cdots+x_m\alpha_m=0$有非零解;

              列向量组$\alpha_1,\alpha_2,\ldots,\alpha_m$线性无关$\iff$齐次线性方程组$x_1\alpha_1+x_2\alpha_2+\cdots+x_m\alpha_m=0$只有零解.
\end{solve}

\begin{exercise}[补充训练 L3-例-10]\label{prob:L3-例-10}
已知 $\mathbf{R}^4$ 的一个子集 $S = \{a_1, a_2, a_3, a_4\}$, 其中
    \[
        a_1 = (1,1,0,1), \quad a_2 = (0,1,2,4), \quad
        a_3 = (2,1,-2,-2), \quad a_4 = (0,1,1,1).
    \]
    试求 $\operatorname{span}(S)$ 的维数及其一组基$B$。
\end{exercise}
\noindent 来源：\href{https://github.com/yhwu-is/Linear-Algebra-Left-Undone/blob/new/%E8%AE%B2%E4%B9%89/%E4%B8%93%E9%A2%98/3%20%E6%9C%89%E9%99%90%E7%BB%B4%E7%BA%BF%E6%80%A7%E7%A9%BA%E9%97%B4.tex}{L3-例-10}。
\begin{solve}
显然这一问题的关键就是求出 $S$ 的一组极大线性无关组，因为 $S$ 的任一极大线性无关组 $B$ 都是 $\operatorname{span}(S)$ 的基. 为了达到这一目标，我们首先将四个向量按列排成矩阵，即

    \begin{equation} \label{eq:极大线性无关组:1}
        \begin{pmatrix}
            1 & 0 & 2  & 0 \\
            1 & 1 & 1  & 1 \\
            0 & 2 & -2 & 1 \\
            1 & 4 & -2 & 1
        \end{pmatrix}
    \end{equation}

    对上述矩阵作初等行变换，所得的阶梯矩阵为

    \begin{equation} \label{eq:极大线性无关组:2}
        \begin{pmatrix}
            1 & 0 & 2  & 0 \\
            0 & 1 & -1 & 0 \\
            0 & 0 & 0  & 1 \\
            0 & 0 & 0  & 0
        \end{pmatrix}
    \end{equation}

    根据上面的引理，我们知道矩阵~\eqref{eq:极大线性无关组:1} 中的列的线性相关性与矩阵~\eqref{eq:极大线性无关组:2} 中的列的线性相关性是一致的，而 \eqref{eq:极大线性无关组:2} 中的经过初等变换的线性相关关系非常容易看出来：我们直接找每一行第一个非零元素所在的列对应的向量即可得到极大线性无关组，即 $S$ 的一个极大线性无关组为 $\{a_1, a_2, a_4\}$. 故 $a_1, a_2, a_4$ 是 $\operatorname{span}(S)$ 的一组基，$\operatorname{span}(S)$ 的维数为 3.
\end{solve}

\begin{exercise}[补充训练 L3-例-11]\label{prob:L3-例-11}
设$V=\mathbf{R}[x]_4$，我们已有向量组$B=\{1+x,x^3+x^2+3x\}$，请将其扩充为$V$的一组基.
\end{exercise}
\noindent 来源：\href{https://github.com/yhwu-is/Linear-Algebra-Left-Undone/blob/new/%E8%AE%B2%E4%B9%89/%E4%B8%93%E9%A2%98/3%20%E6%9C%89%E9%99%90%E7%BB%B4%E7%BA%BF%E6%80%A7%E7%A9%BA%E9%97%B4.tex}{L3-例-11}。
\begin{solve}
按讲义中的方法，取定 $\mathbf{R}[x]_4$ 的自然基，给出$B$中向量的坐标
    \[\alpha_1 = (1, 1, 0, 0), \alpha_2 = (0, 3, 1, 1).\]
    再取一组基
    \[\beta_1 = (1, 0, 0, 0), \beta_2 = (0, 1, 0, 0), \beta_3 = (0, 0, 1, 0), \beta_4 = (0, 0, 0, 1).\]
    将这 6 个向量排列成矩阵，求解极大线性无关组.
    \[ \begin{pmatrix}
            1 & 0 & 1 & 0 & 0 & 0 \\
            1 & 3 & 0 & 1 & 0 & 0 \\
            0 & 1 & 0 & 0 & 1 & 0 \\
            0 & 1 & 0 & 0 & 0 & 1
        \end{pmatrix}
        \xRightarrow{\text{初等行变换}}
        \begin{pmatrix}
            1 & 0 & 1  & 0 & 0 & 0  \\
            0 & 1 & 0  & 0 & 0 & 1  \\
            0 & 0 & -1 & 1 & 0 & -3 \\
            0 & 0 & 0  & 0 & 1 & -1
        \end{pmatrix} \]
    则可取 $\alpha_1, \alpha_2, \beta_1, \beta_3$作为$V$的一组基，即$\{1 + x, 3x + x^2 + x^3, 1, x^2\}$.
\end{solve}

\begin{exercise}[补充训练 L3-例-12]\label{prob:L3-例-12}
分别求$p(x)=a_0+a_1x+a_2x^2$在基$B_1=\{1,x,x^2\}$和$B_2=\{1,x-1,(x-1)^2\}$下的坐标.
\end{exercise}
\noindent 来源：\href{https://github.com/yhwu-is/Linear-Algebra-Left-Undone/blob/new/%E8%AE%B2%E4%B9%89/%E4%B8%93%E9%A2%98/3%20%E6%9C%89%E9%99%90%E7%BB%B4%E7%BA%BF%E6%80%A7%E7%A9%BA%E9%97%B4.tex}{L3-例-12}。
\begin{solve}
通过待定系数法解方程即可.

    在 $B_1$ 下：$(a_0, a_1, a_2)$. 在 $B_2$ 下：$(a_0 + a_1 + a_2, a_1 + 2a_2, a_2)$.
\end{solve}

\begin{exercise}[补充训练 L3-例-13]\label{prob:L3-例-13}
求$\mathbf{R}[x]_4$中向量组$\{p_1=x^3+x+1,p_2=4x^3+2x^2+x,p_3=2x^3-2x^2+x+2,p_4=x^3+x^2+x\}$的极大线性无关组.
\end{exercise}
\noindent 来源：\href{https://github.com/yhwu-is/Linear-Algebra-Left-Undone/blob/new/%E8%AE%B2%E4%B9%89/%E4%B8%93%E9%A2%98/3%20%E6%9C%89%E9%99%90%E7%BB%B4%E7%BA%BF%E6%80%A7%E7%A9%BA%E9%97%B4.tex}{L3-例-13}。
\begin{solve}
相对于自然基 $(1,x,x^2,x^3)$，把四个多项式的坐标作为列，得
\[M=\begin{pmatrix}1&0&2&0\\1&1&1&1\\0&2&-2&1\\1&4&2&1\end{pmatrix}.\]
行变换 $R_2\leftarrow R_2-R_1$、$R_4\leftarrow R_4-R_1$ 后，再作 $R_3\leftarrow R_3-2R_2$、$R_4\leftarrow R_4-4R_2$，得到
\[\begin{pmatrix}1&0&2&0\\0&1&-1&1\\0&0&0&-1\\0&0&4&-3\end{pmatrix}.\]
后两行也线性无关，故秩为 $4$。因此原四个多项式全部线性无关，它们自身就是极大线性无关组，也是 $\mathbf R[x]_4$ 的一组基。
\end{solve}

\section{补充训练：子空间、和与直和}

本节包含 15 道训练题，建议在完成相关专题真题后选择性练习。

\begin{exercise}[补充训练 L4-A-1]\label{prob:L4-A-1}
设$V=\{(a_1,a_2,a_3,a_4) \mid a_1+a_2+a_3+a_4=0\}$，$W=\{(a_1,a_2,a_3,a_4) \mid a_1-a_2-a_3+a_4=0,a_1+a_2+a_3-a_4=0\}$.
        \begin{enumerate}
            \item 证明：$V$和$W$为$\mathbf{R}^4$的子空间；
            \item 分别求$V \cap W$，$V+W$以及$W$的补空间的维数与一组基.
        \end{enumerate}
\end{exercise}
\noindent 来源：\href{https://github.com/yhwu-is/Linear-Algebra-Left-Undone/blob/new/%E8%AE%B2%E4%B9%89/%E4%B8%93%E9%A2%98/4%20%E7%BA%BF%E6%80%A7%E7%A9%BA%E9%97%B4%E7%9A%84%E8%BF%90%E7%AE%97.tex}{L4-A-1}。
\begin{solve}
\begin{enumerate}
                \item 显然任意$\alpha\in v,\beta \in w$, 有 $\alpha,\beta\in \mathbf{R}^4$ 成立，因此只需证明 $v,w$ 封闭即可. 由于 $v,w$ 内向量的约束条件都是齐次线性方程，封闭性自然满足，则 $v,w$ 都是 $\mathbf{R}^4$ 的子空间.

                \item \begin{enumerate}
                          \item $V\cap W=\{(a_1,a_2,a_3,a_4) \mid a_1+a_2+a_3+a_4=0,a_1-a_2-a_3+a_4=0,a_1+a_2+a_3-a_4=0\}$ 或者是 $\{(0,a,-a,0) \mid \alpha\in \mathbf{R}\}$.

                          \item 由维数公式：$\dim V+\dim W=\dim (V\cap W)+\dim (V+W)$，易得$\dim V=3,\enspace \dim W=2,\enspace \dim (V\cap W)=1$，因此$\dim (V+W)=4$，则有$V+W=\mathbf{R}^4$.

                          \item 先得到$W$的基：列出齐次线性方程组：
                                \[\begin{cases}
                                        a_1-a_2-a_3+a_4=0 \\
                                        a_1+a_2+a_3-a_4=0\end{cases}
                                \]
                                高斯消元得到行阶梯矩阵
                                \[\begin{pmatrix}
                                        1 & 0 & 0 & 0  \\
                                        0 & 1 & 1 & -1
                                    \end{pmatrix}\]
                                秩为2，基分别为$\beta_1= (0,-1,1,0)$，$\beta_2= (0,1,0,1)$，则$W$的补空间维数是$\dim \mathbf{R}^4 - \dim W=2$.

                                利用基的扩张求其补空间的基. 也即求$\beta_1,\beta_2,e_1,e_2,e_3,e_4$的极大无关组，其中$e_1,e_2,e_3,e_4$是自然基.
                                \[\begin{pmatrix}
                                        0  & 0 & 1 & 0 & 0 & 0 \\
                                        -1 & 1 & 0 & 1 & 0 & 0 \\
                                        1  & 0 & 0 & 0 & 1 & 0 \\
                                        0  & 1 & 0 & 0 & 0 & 1
                                    \end{pmatrix}\rightarrow
                                    \begin{pmatrix}1 & 0 & 0 & 0 & 1 & 0  \\
                 0 & 1 & 0 & 1 & 1 & 0  \\
                 0 & 0 & 1 & 0 & 0 & 0  \\
                 0 & 0 & 0 & 1 & 1 & -1
                                    \end{pmatrix}\]
                                则$\beta_1,\beta_2,e_1,e_2$即是扩张后的基，因此$W$的补空间的一组基为$e_1,e_2$.
                      \end{enumerate}
            \end{enumerate}
\end{solve}

\begin{exercise}[补充训练 L4-A-2]\label{prob:L4-A-2}
设 $f_1=-1+x,\ f_2=1-x^2,\ f_3=1-x^3,\ g_1=x-x^2,\ g_2=x+x^3,\ V_1=\operatorname{span}(f_1,f_2,f_3),\ V_2=\operatorname{span}(g_1,g_2)$，求：
        \begin{enumerate}
            \item $V_1+V_2$ 的基和维数；
            \item $V_1 \cap V_2$ 的基和维数；
            \item $V_2$ 在 $\mathbf{R}[x]_4$ 空间的补.
        \end{enumerate}
\end{exercise}
\noindent 来源：\href{https://github.com/yhwu-is/Linear-Algebra-Left-Undone/blob/new/%E8%AE%B2%E4%B9%89/%E4%B8%93%E9%A2%98/4%20%E7%BA%BF%E6%80%A7%E7%A9%BA%E9%97%B4%E7%9A%84%E8%BF%90%E7%AE%97.tex}{L4-A-2}。
\begin{solve}
\begin{enumerate}
                \item 先将各向量用坐标表示：
                      \begin{gather*}
                          f_1=(-1,1,0,0),f_2=(1,0,-1,0),f_3=(1,0,0,-1) \\
                          g_1=(0,1,-1,0),g_2=(0,1,0,1) \\
                          V_1+V_2=\operatorname{span}(f_1,f_2,f_3,g_1,g_2)
                      \end{gather*}
                      只需求 $f_1,f_2,f_3,g_1,g_2$ 的极大线性无关组即可.
                      \[\begin{pmatrix}
                              -1 & 1  & 1  & 0  & 0 \\
                              1  & 0  & 0  & 1  & 1 \\
                              0  & -1 & 0  & -1 & 0 \\
                              0  & 0  & -1 & 0  & 1
                          \end{pmatrix}\rightarrow
                          \begin{pmatrix}1 & -1 & -1 & 0 & 0 \\
                 0 & 1  & 1  & 1 & 1 \\
                 0 & 0  & 1  & 0 & 1 \\
                 0 & 0  & 0  & 0 & 1
                          \end{pmatrix}\]
                      则极大线性无关组为 $f_1,f_2,f_3,g_2$，因此 $V_1+V_2$ 的基为 $f_1,f_2,f_3,g_2$，维数为4.

                \item 易得$\dim V_1=3,\enspace \dim V_2=2$. 则$\dim (V_1\cap V_2)= \dim V_1+\dim V_2-\dim (V_1+V_2)=1$. 只需找到属于$V_1\cap V_2$的一个向量，其就是$V_1\cap V_2$的基.
                      \[U=\lambda_1f_1+\lambda_2f_2+\lambda_3f_3= \mu_1g_1+\mu_2g_2,\]
                      已知$g_1$可被$f_1,f_2,f_3,g_2$表示，则只需取$\mu_1=1$，再求解其它系数即解得：
                      \[\lambda_1=1,\lambda_2=1,\lambda_3=0,\mu_2=0,\]
                      因此$u=g_1=(0,1,-1,0)$或者$x-x^2$是$V_1\cap V_2$的一组基.

                \item 只需求$g_1,g_2,e_1,e_2,e_3,e_4$的极大无关组即可.
                      \[\begin{pmatrix}
                              0  & 0 & 1 & 0 & 0 & 0 \\
                              1  & 1 & 0 & 1 & 0 & 0 \\
                              -1 & 0 & 0 & 0 & 1 & 0 \\
                              0  & 1 & 0 & 0 & 0 & 1\end{pmatrix}\rightarrow
                          \begin{pmatrix}
                              1 & 1 & 0 & 1 & 0 & 0  \\
                              0 & 1 & 0 & 1 & 1 & 0  \\
                              0 & 0 & 1 & 0 & 0 & 0  \\
                              0 & 0 & 0 & 1 & 1 & -1\end{pmatrix}\]
                      则$g_1,g_2,e_1,e_2$是一组极大无关组. $V_2$在$\mathbf{R}[x]_4$的补是$\operatorname{span}(e_1,e_2)$其中$e_1=1,e_2=x$.
            \end{enumerate}
\end{solve}

\begin{exercise}[补充训练 L4-A-5]\label{prob:L4-A-5}
设 $V_1,V_2$ 是线性空间 $V$ 的两个非平凡子空间，证明: $\exists \alpha \in V$,使得$\alpha \notin V_1$ 且 $\alpha \notin V_2$. 并在$\mathbf{R}^3$ 中举一例.
\end{exercise}
\noindent 来源：\href{https://github.com/yhwu-is/Linear-Algebra-Left-Undone/blob/new/%E8%AE%B2%E4%B9%89/%E4%B8%93%E9%A2%98/4%20%E7%BA%BF%E6%80%A7%E7%A9%BA%E9%97%B4%E7%9A%84%E8%BF%90%E7%AE%97.tex}{L4-A-5}。
\begin{solve}
\textbf{证:} 因为 $V_1, V_2$ 是非平凡子空间, 所以存在 $\alpha \notin V_1$. 若 $\alpha \notin V_2$, 则命题得证；若 $\alpha \in V_2$, 另有 $\beta \notin V_2$, 此时若 $\beta \notin V_1$, 命题也得证. 设 $\beta \in V_1$, 则 $\alpha \notin V_1, \alpha \in V_2, \beta \in V_1, \beta \notin V_2$.
            考虑 $\alpha + \beta \in V_1$, 由 $\beta \in V_1 \implies \alpha \in V_1$, 从而得出矛盾, 所以 $\alpha + \beta \notin V_1$. 类似可证 $\alpha + \beta \notin V_2$. 于是命题成立.
\end{solve}

\begin{exercise}[补充训练 L4-A-6]\label{prob:L4-A-6}
设 $V_1,\ldots,V_m (m > 2)$ 是线性空间$V$的$m$个非平凡子空间，证明：$V$ 中存在一个同时不属于任何一个$V_i(1 \leqslant i \leqslant m)$的向量. 并在$\mathbf{R}^3$ 中举一例.
\end{exercise}
\noindent 来源：\href{https://github.com/yhwu-is/Linear-Algebra-Left-Undone/blob/new/%E8%AE%B2%E4%B9%89/%E4%B8%93%E9%A2%98/4%20%E7%BA%BF%E6%80%A7%E7%A9%BA%E9%97%B4%E7%9A%84%E8%BF%90%E7%AE%97.tex}{L4-A-6}。
\begin{solve}
\label{eg:4:A:5}
          \textbf{证:} （归纳法）当 $r=2$ 时由上题知成立.
          假设 $r = k-1$ 时也成立, 即存在 $\alpha \in V$
          使得 $\alpha$ 同时不属于 $V_i (i=1,2,\ldots,k-1)$
          下面证明 $r=k$ 时也成立. 显然当此 $\alpha \notin V_k$ 则 $\alpha$ 即为所求.
           若 $\alpha \in V_k$, 由于 $V_k$ 是 $V$ 的真子空间, 故有 $\beta \notin V_k$, 易知对任意 $V_i (i=1,2,\ldots,k-1)$ 都至多有一个 $k_i$, 使得 $\beta + k_i \alpha \in V_i; i=1,2,\ldots,k-1$, 我们取异于 $k_i$ 的数 $k$, 则必有 $\beta + k \alpha \notin V_i; i=1,2,\ldots,k-1$. 又由于 $\alpha \in V_k, \beta \notin V_k$ 故有 $\beta + k \alpha \notin V_k$, 故 $\beta + k \alpha \notin V_k \enspace (i = 1,2,\ldots,k)$.
\end{solve}

\begin{exercise}[补充训练 L4-B-1]\label{prob:L4-B-1}
已知$V_1$是线性方程组\[\begin{cases}
                3x_1+4x_2-5x_3+7x_4=0 \\
                4x_1+11x_2-13x_3+16x_4=0
            \end{cases}\]
        的解空间，$V_2$是线性方程组\[\begin{cases}
                2x_1-3x_2+3x_3-2x_4=0 \\
                7x_1-2x_2+x_3+3x_4=0
            \end{cases}\]
        的解空间，分别求$V_1 \cap V_2$与$V_1+V_2$的基和维数.
\end{exercise}
\noindent 来源：\href{https://github.com/yhwu-is/Linear-Algebra-Left-Undone/blob/new/%E8%AE%B2%E4%B9%89/%E4%B8%93%E9%A2%98/4%20%E7%BA%BF%E6%80%A7%E7%A9%BA%E9%97%B4%E7%9A%84%E8%BF%90%E7%AE%97.tex}{L4-B-1}。
\begin{solve}
\begin{enumerate}
                \item $V_1\cap V_2$ 即是两组方程组合并后的解空间.
                      \[\begin{pmatrix}
                              3 & 4  & -5  & 7  \\
                              4 & 11 & -13 & 16 \\
                              2 & -3 & 3   & -2 \\
                              7 & -2 & 1   & 3
                          \end{pmatrix}\rightarrow
                          \begin{pmatrix}1 & 7   & -8 & 9   \\
                 0 & -17 & 19 & -20 \\
                 0 & 0   & 0  & 0   \\
                 0 & 0   & 0  & 0
                          \end{pmatrix}\]
                      则解向量$u_1=\left(\frac {3}{17},\frac{19}{17},1,0\right)^\mathrm{T},u_2=\left(\frac{13}{17},-\frac{20}{17},0,1\right)^\mathrm{T}$. $u_1,u_2$是$V_1\cap V_2$的基，则其维数为2.

                \item 易得$\dim V_1=2,\dim V_2=2$，则由维数公式$\dim (V_1+V_2)=\dim V_1+\dim V_2-\dim (V_1\cap V_2)= 2$又$V_1\cap V_2\subseteq V_1+V_2$，且二者维数相等. 则$V_1\cap V_2=V_1+V_2$，亦即$V_1=V_2$. 所以$V_1+V_2$的基也是$u_1,u_2$，同$V_1\cap V_2$.
            \end{enumerate}
\end{solve}

\begin{exercise}[补充训练 L4-B-2]\label{prob:L4-B-2}
设$W_1,W_2$是线性空间$V(\mathbf{F})$的两个子空间. 证明以下命题等价：
        \begin{enumerate}
            \item $W_1 \cup W_2$为$V$ 的子空间；

            \item $W_1 \subseteq W_2$或$W_2 \subseteq W_1$；

            \item $W_1 \cup W_2=W_1+W_2$.
        \end{enumerate}
\end{exercise}
\noindent 来源：\href{https://github.com/yhwu-is/Linear-Algebra-Left-Undone/blob/new/%E8%AE%B2%E4%B9%89/%E4%B8%93%E9%A2%98/4%20%E7%BA%BF%E6%80%A7%E7%A9%BA%E9%97%B4%E7%9A%84%E8%BF%90%E7%AE%97.tex}{L4-B-2}。
\begin{solve}
先证 (1) 蕴含 (2)。若两子空间互不包含，可取 $u\in W_1\setminus W_2$ 和 $v\in W_2\setminus W_1$。若 $u+v\in W_1$，则 $v=(u+v)-u\in W_1$，矛盾；同理 $u+v\notin W_2$。因此并集不对加法封闭，与 (1) 矛盾。再证 (2) 蕴含 (3)：若 $W_1\subseteq W_2$，则并与和都等于 $W_2$；另一种包含关系同理。最后，和空间 $W_1+W_2$ 总是子空间，故 (3) 蕴含 (1)。这构成了完整的循环证明。
\end{solve}

\begin{exercise}[补充训练 L4-B-3]\label{prob:L4-B-3}
证明：每个 $n$ 维空间均可以表示为 $n$ 个一维子空间的直和.
\end{exercise}
\noindent 来源：\href{https://github.com/yhwu-is/Linear-Algebra-Left-Undone/blob/new/%E8%AE%B2%E4%B9%89/%E4%B8%93%E9%A2%98/4%20%E7%BA%BF%E6%80%A7%E7%A9%BA%E9%97%B4%E7%9A%84%E8%BF%90%E7%AE%97.tex}{L4-B-3}。
\begin{solve}
我们很自然地想到，可以任取 $n$ 维空间 $V$ 的一组基 $\alpha_1,\alpha_2,\ldots,\alpha_n$，然后取 $V_i = \operatorname{span}(\alpha_i)$. 如果 $V = V_1 \oplus V_2 \oplus \cdots \oplus V_n$，我们就成功地为每个 $n$ 维空间找到了写成 $n$ 个一维子空间直和的分解形式. 下面我们来证明这一构造的正确性：
            \begin{enumerate}
                \item
                \item
            \end{enumerate}
\end{solve}

\begin{exercise}[补充训练 L4-B-4]\label{prob:L4-B-4}
设$V$是域$\mathbf{F}$上的$n$维线性空间，$\alpha_1,\alpha_2,\ldots,\alpha_n$是$V$的一组基，且
        \begin{gather*}
            V_1=\operatorname{span}(\alpha_1+2\alpha_2+\cdots+n\alpha_n) \\
            V_2=\left\{k_1\alpha_1+k_2\alpha_2+\cdots+k_n\alpha_n \;\middle|\; k_1+\dfrac{k_2}{2}+\cdots+\dfrac{k_n}{n}=0\right\}
        \end{gather*}
        证明：
        \begin{enumerate}
            \item $V_2$是$V$的子空间；

            \item $V=V_1\oplus V_2$.
        \end{enumerate}
\end{exercise}
\noindent 来源：\href{https://github.com/yhwu-is/Linear-Algebra-Left-Undone/blob/new/%E8%AE%B2%E4%B9%89/%E4%B8%93%E9%A2%98/4%20%E7%BA%BF%E6%80%A7%E7%A9%BA%E9%97%B4%E7%9A%84%E8%BF%90%E7%AE%97.tex}{L4-B-4}。
\begin{solve}
\begin{enumerate}
                \item 显然$V_2\subseteq V$，只需证明$V_2$封闭即可. 对于$v=k_1\alpha_1+\cdots+k_n\alpha_n,v'=k_1'\alpha_1+\cdots+k_n'\alpha_n$. $\lambda v+\mu v'=(\lambda k_1+\mu k_1')\alpha_1+\cdots+(\lambda k_n+\mu k'_n)\alpha_n$，因此  $\lambda v+\mu v'\in V_2\in V_2$，$V_2$ 封闭，则$V_2$是$V$的子空间.

                \item 设$W=V_1+V_2$. 对任意 $v\in V_1,v=\lambda \alpha_1+2\lambda\alpha_2+\cdots+n\lambda\alpha_n$，则 $\lambda+2\dfrac{\lambda}2+\allowbreak\cdots+n\dfrac{\lambda}n=n\lambda$，因此当$v\ne 0$时，$v\not\in v_2$，则 $V_1\cap V_2=\{0\}$，有$W=V_1\oplus V_2$.

                      下证$W=V$，易得$\dim V_1=1$. 对于$V_2$，写成坐标形式，则其是一个方程组的解空间，系数矩阵 
\[A=\begin{pmatrix}
                              1 & \frac 12 & \cdots & \frac 1n \\
                              0 &          &        &          \\
                                & \ddots   &        &          \\
                                &          & \ddots &          \\
                                &          &        & 0
                          \end{pmatrix}\]
. 方程组为 $AK=0$，则 $\dim V_2=n-r(A)=n-1$，因此$\dim W=\dim V_1+\dim V_2=n=\dim V$，又$W\subseteq V$,则$W=V$得证，综上$V=V_1\oplus V_2$.
            \end{enumerate}
\end{solve}

\begin{exercise}[补充训练 L4-B-6]\label{prob:L4-B-6}
已知$V_1,V_2$是有限维线性空间$V$的子空间，且$\dim(V_1+V_2)=\dim(V_1 \cap V_2)+1$. 证明：要么$V_1 \subseteq V_2$，要么$V_2 \subseteq V_1$.
\end{exercise}
\noindent 来源：\href{https://github.com/yhwu-is/Linear-Algebra-Left-Undone/blob/new/%E8%AE%B2%E4%B9%89/%E4%B8%93%E9%A2%98/4%20%E7%BA%BF%E6%80%A7%E7%A9%BA%E9%97%B4%E7%9A%84%E8%BF%90%E7%AE%97.tex}{L4-B-6}。
\begin{solve}
\begin{enumerate}
                \item 法一：反证法. 如果$\exists v_1\in V_1,v_1\not\in V_2$,且 $\exists v_2\in V_2,v_2\not\in V_1$，则有 $v_1,v_2\not \in V_1\cap V_2$ 并且 $v_1,v_2$ 线性无关. 设 $A=\{\alpha_1,\ldots ,\alpha_n\}$是 $V_1\cap V_2$的基. 由基扩张，设$B=\{\alpha_1,\ldots ,\alpha_n,\alpha_{n+1},\ldots ,\alpha_{n+k}\}$是$V_1+V_2$的基. 由于$v_1,v_2$不在$V_1\cap V_2$中，因此$A$无法表出$v_1,v_2$. 则$\alpha_1,\ldots ,\alpha_n,v_1,v_2$ 线性无关. 又$v_1,v_2\in V_1+V_2$，则$\alpha_1,\ldots ,\alpha_n,v_1,v_2$可由$\alpha_1,\ldots ,\alpha_{n+k}$ 线性表出. 由定理 $3.3$ 得 $n+2\leqslant n+k$，即 $\dim (V_1+V_2)\ge\dim (V_1\cap V_2)+2$，这与条件矛盾. 因此，$\forall v_1\in V_1,v_1\in V_2$ 或 $\forall v_2\in V_2,v_2\in V_1$. 即 $V_1\subseteq V_2$ 或 $V_2\subseteq V_1$，得证.

                \item 法二：维数公式. $\dim V_1+\dim V_2=\dim  (V_1+V_2)+\dim  (V_1\cap V_2)=2\dim  (V_1\cap V_2)+1$，即 $(\dim V_1-\dim  (V_1\cap V_2))+ (\dim V_2-\dim  (V_1\cap V_2))=1$. 于是，要么 $\dim V_1-\dim  (V_1\cap V_2)=0$，要么 $\dim V_2-\dim  (V_1\cap V_2)=0$，即 $V_1\subseteq V_2$ 或 $V_2\subseteq V_1$，得证.
            \end{enumerate}
\end{solve}

\begin{exercise}[补充训练 L4-B-7]\label{prob:L4-B-7}
判断下列说法是否正确：
        \begin{enumerate}
            \item 若$V \subseteq V_1 \cup V_2 \cup \cdots \cup V_s$，则$V=(V_1 \cap V)\cup(V_2 \cap V)\cup\cdots\cup(V_s \cap V)$；

            \item 若$V \subseteq V_1+V_2+\cdots +V_s$，则$V=(V_1 \cap V)+(V_2 \cap V)+\cdots+(V_s \cap V)$.
        \end{enumerate}
\end{exercise}
\noindent 来源：\href{https://github.com/yhwu-is/Linear-Algebra-Left-Undone/blob/new/%E8%AE%B2%E4%B9%89/%E4%B8%93%E9%A2%98/4%20%E7%BA%BF%E6%80%A7%E7%A9%BA%E9%97%B4%E7%9A%84%E8%BF%90%E7%AE%97.tex}{L4-B-7}。
\begin{solve}
\begin{enumerate}
              \item 正确，若$\alpha\in V$，则$\alpha \in V_1\cup V_2\cdots \cup V_s$，即$\alpha \in V\cap (V_1\cup V_2\cdots \cup V_s)$，则$\alpha \in (V\cap V_1)\cup \cdots \cup (V\cap V_s)$，则$V\subseteq(V_1\cap V)\cup \cdots \cup (V_s\cap V)$. 另一边类似的同样成立，则有 $V=(V_1\cap V)\cup \cdots \cup (V_s\cap V)$ 得证.

              \item 错误，反例：设 $V_1,V_2,V$ 是平面$K$上三条过原点$O$的不重合直线，则$V\subseteq K=V_1+V_2$，但$V\cap V_1=\{0\},V\cap V_2=\{0\}$，$V\ne (V\cap V_1)+(V\cap V_2)$.
          \end{enumerate}
\end{solve}

\begin{exercise}[补充训练 L4-B-9]\label{prob:L4-B-9}
设$A_1$和$A_2$均为$V$的仿射子集，证明：$A_1\cap A_2$是$V$的仿射子集或空集（可推广至任意交）.
\end{exercise}
\noindent 来源：\href{https://github.com/yhwu-is/Linear-Algebra-Left-Undone/blob/new/%E8%AE%B2%E4%B9%89/%E4%B8%93%E9%A2%98/4%20%E7%BA%BF%E6%80%A7%E7%A9%BA%E9%97%B4%E7%9A%84%E8%BF%90%E7%AE%97.tex}{L4-B-9}。
\begin{solve}
若交集为空，结论成立。否则取 $v\in A_1\cap A_2$。由于两个仿射子集都经过 $v$，可写成 $A_1=v+U_1$、$A_2=v+U_2$，其中 $U_1,U_2$ 是线性子空间。对任意 $x\in V$，
\[x\in A_1\cap A_2\iff x-v\in U_1\cap U_2.\]
故 $A_1\cap A_2=v+(U_1\cap U_2)$，是仿射子集。任意一族非空交集也可选取同一交点并照此证明；单点集通常只是仿射子集，仅当该点为零时才是线性子空间。
\end{solve}

\begin{exercise}[补充训练 L4-C-1]\label{prob:L4-C-1}
设$W_0,W_1,W_2,\ldots,W_s$是线性空间$V$的$s+1$个非平凡子空间，且$W_0 \subseteq W_1 \cup W_2 \cup \cdots \cup W_s$. 证明：必存在$i$使得$W_0\subseteq W_i$.
\end{exercise}
\noindent 来源：\href{https://github.com/yhwu-is/Linear-Algebra-Left-Undone/blob/new/%E8%AE%B2%E4%B9%89/%E4%B8%93%E9%A2%98/4%20%E7%BA%BF%E6%80%A7%E7%A9%BA%E9%97%B4%E7%9A%84%E8%BF%90%E7%AE%97.tex}{L4-C-1}。
\begin{solve}
数域是无限域。先证明无限域上的向量空间 $U$ 不可能是有限个真子空间 $H_1,\ldots,H_s$ 的并。对 $s$ 归纳：$s=1$ 时结论就是“真子空间不等于全空间”。设 $s\ge2$，由归纳假设可取 $u\notin H_1\cup\cdots\cup H_{s-1}$，另取 $v\notin H_s$。对 $i<s$，直线 $\{v+tu:t\in\mathbf F\}$ 与 $H_i$ 至多有一个交点，否则两个交点相减将推出 $u\in H_i$。对 $H_s$，若 $u\in H_s$ 则无交点，否则同样至多一个交点。因为域无限，可避开这有限多个参数，得到不在任何 $H_i$ 中的向量，完成归纳。
现在在 $U=W_0$ 中应用结论。由题设，$W_0=\bigcup_{i=1}^s(W_0\cap W_i)$，所以至少一个交空间不是 $W_0$ 的真子空间，即 $W_0=W_0\cap W_i$，从而 $W_0\subseteq W_i$。
\end{solve}

\begin{exercise}[补充训练 L4-C-4]\label{prob:L4-C-4}
设$v_1,\ldots,v_m\in V$. 令
        \[A=\{\lambda_1v_1+\cdots+\lambda_mv_m \mid \lambda_1,\ldots,\lambda_m\in\mathbf{F}\text{~且~}\lambda_1+\cdots+\lambda_m=1\}.\]
        证明：
        \begin{enumerate}
            \item $A$是$V$的仿射子集；

            \item $V$的每个包含$v_1,\ldots,v_m$的仿射子集均包含$A$；

            \item 存在某个$v\in V$和$V$的子空间$U$使得$A=v+U$且$\dim U\leqslant m-1$.
        \end{enumerate}
\end{exercise}
\noindent 来源：\href{https://github.com/yhwu-is/Linear-Algebra-Left-Undone/blob/new/%E8%AE%B2%E4%B9%89/%E4%B8%93%E9%A2%98/4%20%E7%BA%BF%E6%80%A7%E7%A9%BA%E9%97%B4%E7%9A%84%E8%BF%90%E7%AE%97.tex}{L4-C-4}。
\begin{solve}
\begin{enumerate}
            \item 对于 $v=\lambda_1 v_1+\cdots+\lambda_m v_m\in A$ 和 $w=\eta_1 v_1+\cdots+\eta_m v_m\in A$, 且 $\lambda_1,\cdots,\lambda_m\in  \mathbf{F}$, $\lambda_1+\cdots+\lambda_m=1$ , $\eta_1,\cdots,\eta_m\in \mathbf{F}$, $\eta_1+\cdots+\eta_m=1$. 对于任意的$\lambda\in \mathbf{F}$, 我们有
            \begin{align*}
                \lambda v+(1-\lambda) w=\sum_{i=1}^m(\lambda\lambda_i+(1-\lambda)\eta_i)v_i.
            \end{align*}
            注意到
            \[ \sum_{i=1}^m(\lambda\lambda_i+(1-\lambda)\eta_i)=\lambda\sum_{i=1}^m\lambda_i+(1-\lambda)\sum_{i=1}^n\eta_i=\lambda+(1-\lambda)=1, \]
            我们有$\lambda v+(1-\lambda) w\in A$. 故 $A$ 是 $V$ 的一个仿射子集
        \item 我们使用数学归纳法来证明，对于任意 $V$ 的仿射子集 $W$ ,其包含$v_1$, $\cdots$, $v_m$$k\le m$,
            如果 $\lambda_1+\cdots+\lambda_k=1$,
            我们有 \[ \sum_{j=1}^k\lambda_jv_j\in W. \]

            当$k=1$ 和 $k=2$ 时，结论是显然的.
            假设 $k$ 时成立,
            接下来对于 $k+1$ ($k+1\leqslant m$).
            我们假设$\lambda_1+\cdots+\lambda_{k+1}=1$.
            若 $\lambda_{k+1}=1$, 那么\[\sum_{j=1}^{k+1}\lambda_jv_j=v_{k+1}\in W.\]
            若$\lambda_{k+1}\ne 1$,那么\[ \frac{1}{1-\lambda_{k+1}}(\lambda_1+\cdots+\lambda_k)=1. \]
            由归纳假设，我们有\[ \frac{1}{1-\lambda_{k+1}}(\lambda_1v_1+\cdots+\lambda_kv_k)\in W. \]
            也有 \[ (1-a_{k+1})\left(\frac{1}{1-\lambda_{k+1}}(\lambda_1v_1+\cdots+\lambda_kv_k)\right)+a_{k+1}v_{k+1}\in W, \]
            即 \[ \lambda_1v_1+\cdots+\lambda_{k+1}v_{k+1}\in W. \]
            故而 $\lambda_1,\cdots,\lambda_m\in  \mathbf{F}$, $\lambda_1+\cdots+\lambda_m=1$,
            \[\lambda_1 v_1+\cdots+\lambda_m v_m\in W,\]有$A\subset W$.
        \item  注意到$\lambda_1+\cdots+\lambda_m=1$, 我们有\[ \lambda_1 v_1+\cdots+\lambda_m v_m=v_1+\lambda_2(v_2-v_1)+\cdots+\lambda_m(v_m-v_1). \]所以 $A\subset v_1+\operatorname{span}(v_2-v_1,\cdots,v_m-v_1)$.
            类似的，对于任意\[v\in v_1+\operatorname{span}(v_2-v_1,\cdots,v_m-v_1),\]$v$ 可以被写为 \[ v_1+\sum_{i=2}^m \lambda_i(v_i-v_1)=(1-\lambda_2-\cdots-\lambda_m)v_1+\sum_{i=2}^m \lambda_i v_i \]
            对于 $\lambda_2$, $\cdots$, $\lambda_m\in \mathbf{F}$. 注意到\[ (1-\lambda_2-\cdots-\lambda_m)+\sum_{i=2}^m \lambda_i=1, \]
            可以推出 $v_1+\operatorname{span}(v_2-v_1,\cdots,v_m-v_1)\subset A$.
            因此\[A=v_1+\operatorname{span}(v_2-v_1,\cdots,v_m-v_1).\]
            $v=v_1, U=\operatorname{span}(v_2-v_1,\cdots,v_m-v_1)$, 有 $\dim U\le m-1$
        \end{enumerate}
\end{solve}

\begin{exercise}[补充训练 L4-例-2]\label{prob:L4-例-2}
设 $\alpha_1 = (1, 0, -1, 0)$，$\alpha_2 = (0, 1, 2, 1)$，$\alpha_3 = (2, 1, 0, 1)$，是四维实行向量空间 $V$ 中的向量，他们张成的子空间为 $V_1$；又设向量 $\beta_1 = (-1, 1, 1, 1)$，$\beta_2 = (1, -1, -3, -1)$，$\beta_3 = (-1, 1, -1, 1)$ 张成的子空间为 $V_2$，求 $V_1$ 和 $V_2$ 的交与和的基.
\end{exercise}
\noindent 来源：\href{https://github.com/yhwu-is/Linear-Algebra-Left-Undone/blob/new/%E8%AE%B2%E4%B9%89/%E4%B8%93%E9%A2%98/4%20%E7%BA%BF%E6%80%A7%E7%A9%BA%E9%97%B4%E7%9A%84%E8%BF%90%E7%AE%97.tex}{L4-例-2}。
\begin{solve}
\begin{enumerate}
        \item 方法一.  $V_1 +V_2$ 是由 $\alpha_i$ 和 $\beta_i$ 生成的，因此只需要求出这 $6$ 个向量的极大线性无关组即可. 将这 $6$ 个向量按列分块方式拼成矩阵，并用初等行变换将其化为阶梯形矩阵：
              \begin{align*}
                  \begin{pmatrix}
                      1  & 0 & 2 & -1 & 1  & -1 \\
                      0  & 1 & 1 & 1  & -1 & 1  \\
                      -1 & 2 & 0 & 1  & -3 & -1 \\
                      0  & 1 & 1 & 1  & -1 & 1
                  \end{pmatrix}
                  \xrightarrow{}
                  \begin{pmatrix}
                      1 & 0 & 2 & -1 & 1  & -1 \\
                      0 & 1 & 1 & 1  & -1 & 1  \\
                      0 & 2 & 2 & 0  & -2 & -2 \\
                      0 & 0 & 0 & 0  & 0  & 0
                  \end{pmatrix} \\
                  \xrightarrow{}
                  \begin{pmatrix}
                      1 & 0 & 2 & -1 & 1  & -1 \\
                      0 & 1 & 1 & 1  & -1 & 1  \\
                      0 & 0 & 0 & -2 & 0  & -4 \\
                      0 & 0 & 0 & 0  & 0  & 0
                  \end{pmatrix}
              \end{align*}

              所以就可以取 $\alpha_1$，$\alpha_2$，$\beta_1$ 为 $V_1 + V_2$ 的基（不唯一）.

              下面再来取 $V_1\cap V_2$ 的基，首先注意到 $\alpha_1$，$\alpha_2$ 是 $V_1$ 的基（从上面的矩阵即可看出），又不难验证 $\beta_1$，$\beta_2$ 是 $V_2$ 的基，$V_2$ 中的向量可以表示为 $\beta_1$，$\beta_2$ 的线性组合. 假设 $t_1\beta_1 + t_2\beta_2$ 属于 $V_1$，则向量组 $\alpha_1, \alpha_2, t_1\beta_1 + t_2\beta_2$ 和向量组 $\alpha_1, \alpha_2$ 的秩相等（因为 $\alpha_1, \alpha_2$ 是 $V_1$ 的基）. 因此，我们可以用矩阵方法来求出参数 $t_1, t_2$. 注意到
              \[ \begin{pmatrix}
                      1  & 0 & -t + t_2   \\
                      0  & 1 & t_1 - t_2  \\
                      -1 & 2 & t_1 - 3t_2 \\
                      0  & 1 & -t_1 - t_2
                  \end{pmatrix} \xrightarrow{} \begin{pmatrix}
                      1 & 0 & -t + t_2  \\
                      0 & 1 & t_1 - t_2 \\
                      0 & 2 & -2t_2     \\
                      0 & 0 & 0
                  \end{pmatrix} \xrightarrow{} \begin{pmatrix}
                      1 & 0 & -t + t_2  \\
                      0 & 1 & t_1 - t_2 \\
                      0 & 0 & -2t_1     \\
                      0 & 0 & 0
                  \end{pmatrix} \]

              所以可以得出当且仅当 $t_1 = 0$ 时 $t_1\beta_1 + t_2\beta_2$ 属于 $V_1$，所以 $V_1 \cap V_2$ 的基可取为 $\beta_2$.

        \item 方法二. 求 $V_1 + V_2$ 的基同方法一，现用解线性方程组的方法来求 $V_1 \cap V_2$ 的基. 因为 $\alpha_1$，$\alpha_2$ 是 $V_1$ 的基，$\beta_1$，$\beta_2$ 是 $V_2$ 的基，故对任一向量 $\gamma \in V_1 \cap V_2$，$\gamma = x_1\alpha_1 + x_2\alpha_2 = -x_3\beta_1 - x_4\beta_2$. 因此，求向量 $\gamma$ 等价于求解线性方程组

              \[ x_1\alpha_1 + x_2\alpha_2 + x_3\beta_1 + x_4\beta_2 = 0. \]

              上述线性方程的通解是 $(x_1, x_2, x_3, x_4) = k(-1, 1, 0, 1)$，从而 $\gamma = -k(\alpha_1 - \alpha_2) = -k\beta_2 (k \in \mathbf{R})$，于是 $\beta_2$ 是 $V_1 \cap V_2$ 的基.
    \end{enumerate}
\end{solve}

\begin{exercise}[补充训练 L4-例-8]\label{prob:L4-例-8}
设$A$是$\mathbf{R}$上的$2\times 3$矩阵：
    \[A=\begin{pmatrix}
            1 & -1 & 2 \\ 1 & 0 & -1
        \end{pmatrix}.\]
    \begin{enumerate}
        \item 求齐次线性方程组$AX=0$的解空间$W$的一组基；

        \item 求商空间$\mathbf{R}^3/W$的维数和一组基.
    \end{enumerate}
\end{exercise}
\noindent 来源：\href{https://github.com/yhwu-is/Linear-Algebra-Left-Undone/blob/new/%E8%AE%B2%E4%B9%89/%E4%B8%93%E9%A2%98/4%20%E7%BA%BF%E6%80%A7%E7%A9%BA%E9%97%B4%E7%9A%84%E8%BF%90%E7%AE%97.tex}{L4-例-8}。
\begin{solve}
\begin{enumerate}
        \item 首先求 $AX=0$ 的解，得到其基础解系为$k(1,3,1)^{\mathrm{T}},k \in \mathbf{R}$,故其解空间的一组基为$\{(1,3,1)^{\mathrm{T}}\}$.
        \item 根据原讲义“商空间的维数公式”处 的证明过程，我们将 $W$ 的基扩张为 $\mathbf{R^3}$ 的一组基，如 $\{(1,3,1)^{\mathrm{T}},(1,0,0)^{\mathrm{T}},(0,1,0)^{\mathrm{T}}\}$，然后将扩张的基中``加上'' $W$ 变成仿射子集，即可得到商空间的一组基：
        \[
            \{(1,0,0)^{\mathrm{T}}+W,(0,1,0)^{\mathrm{T}}+W\}
        \]
        维数为 $2$. 这可以从基的数量得到，也可以由 $\mathbf{R}^3$ 的维数减去 $W$ 的维数得到.

    \end{enumerate}
\end{solve}

\section{补充训练：线性映射、核与像}

本节包含 30 道训练题，建议在完成相关专题真题后选择性练习。

\begin{exercise}[补充训练 L5-A-1]\label{prob:L5-A-1}
设$\sigma: V_1\to V_2$是线性映射. 证明：$\sigma(W_1)$和$\sigma^{-1}(W_2)$分别是$V_2$和$V_1$的子空间.
\end{exercise}
\noindent 来源：\href{https://github.com/yhwu-is/Linear-Algebra-Left-Undone/blob/new/%E8%AE%B2%E4%B9%89/%E4%B8%93%E9%A2%98/5%20%E7%BA%BF%E6%80%A7%E6%98%A0%E5%B0%84.tex}{L5-A-1}。
\begin{solve}
$\sigma(W_1)$是$V_2$的子空间. $\forall v_1,v_2\in\sigma(W_1)$，$\exists w_1,w_2\in W_1$使得$\sigma(w_1)=v_1$，$\sigma(w_2)=v_2$，因此
            \[\sigma(w_1+w_2)=\sigma(w_1)+\sigma(w_2)=v_1+v_2\in\sigma(W_1),\]
            \[\sigma(\lambda w_1)=\lambda\sigma(w_1)=\lambda v_1\in\sigma(W_1),\]
            因此$\sigma(W_1)$是$V_2$的子空间.

            $\sigma^{-1}(W_2)$是$V_1$的子空间. $\forall v_1,v_2\in\sigma^{-1}(W_2)$，$\sigma(v_1),\sigma(v_2)\in W_2$，因此
            \[\sigma(v_1+v_2)=\sigma(v_1)+\sigma(v_2)\in W_2,\]
            \[\sigma(\lambda v_1)=\lambda\sigma(v_1)\in W_2,\]
            因此$v_1+v_2\in\sigma^{-1}(W_2)$，$\lambda v_1\in\sigma^{-1}(W_2)$，因此$\sigma^{-1}(W_2)$是$V_1$的子空间.
\end{solve}

\begin{exercise}[补充训练 L5-A-2]\label{prob:L5-A-2}
设$\sigma,\tau \in \mathcal{L}(V,V)$且$\sigma^2=\sigma$，$\tau^2=\tau$. 证明：
        \begin{enumerate}
            \item $\sigma^k=\sigma$（幂等变换）；

            \item 若$(\sigma+\tau)^2=\sigma+\tau$，则$\sigma\tau=\theta$（零变换）；

            \item 设$\sigma\tau=\tau\sigma$，则$(\sigma+\tau-\sigma\tau)^2=\sigma+\tau-\sigma\tau$.
        \end{enumerate}
\end{exercise}
\noindent 来源：\href{https://github.com/yhwu-is/Linear-Algebra-Left-Undone/blob/new/%E8%AE%B2%E4%B9%89/%E4%B8%93%E9%A2%98/5%20%E7%BA%BF%E6%80%A7%E6%98%A0%E5%B0%84.tex}{L5-A-2}。
\begin{solve}
\begin{enumerate}
                \item 使用数学归纳法证明即可.

                \item 由 $ (\sigma + \tau)^2 = \sigma + \tau $，
                      \[ (\sigma + \tau)^2 = \sigma^2 + + \sigma \tau + \tau \sigma + \tau^2 = \sigma + \tau \]
                      得
                      \begin{equation} \label{eq:5:A:2:1}
                          \sigma \tau + \tau \sigma = \theta
                      \end{equation}
                      式~\eqref{eq:5:A:2:1} 两边左乘 $ \sigma $ 得
                      \begin{align}
                          \sigma(\sigma \tau + \tau \sigma) & = \sigma^2 \tau + \sigma \tau \sigma = \sigma \theta \notag    \\
                                                            & = \sigma \tau + \sigma \tau \sigma = \theta \label{eq:5:A:2:2}
                      \end{align}
                      式~\eqref{eq:5:A:2:1} 两边右乘 $ \sigma $ 得
                      \begin{align}
                          (\sigma \tau + \tau \sigma)\sigma & = \sigma \tau \sigma + \tau \sigma^2 = \theta \sigma \notag    \\
                                                            & = \sigma \tau \sigma + \tau \sigma = \theta \label{eq:5:A:2:3}
                      \end{align}
                      式~\eqref{eq:5:A:2:3} 减去式~\eqref{eq:5:A:2:2} 得
                      \begin{equation}
                          \sigma \tau - \tau \sigma = \theta \label{eq:5:A:2:4}
                      \end{equation}
                      由式~\eqref{eq:5:A:2:1} 与式~\eqref{eq:5:A:2:4} 得
                      \[ \sigma \tau = \theta \]

                \item 若 $ \sigma \tau = \tau \sigma $，则
                      \begin{align*}
                              & (\sigma + \tau - \sigma \tau)^2                                                                                             \\
                          ={} & \sigma^2 + \tau^2 + 2 \sigma \tau - \sigma \tau \sigma - \sigma \tau^2 - \sigma^2 \tau - \tau \sigma \tau + \sigma^2 \tau^2 \\
                          ={} & \sigma + \tau + 2 \sigma \tau - \sigma \tau \sigma - \sigma \tau - \sigma \tau - \tau \sigma - \sigma \tau + \sigma \tau    \\
                          ={} & \sigma + \tau - \sigma \tau
                      \end{align*}
            \end{enumerate}
\end{solve}

\begin{exercise}[补充训练 L5-A-3]\label{prob:L5-A-3}
是否存在$\mathbf{R}^3$到$\mathbf{R}^2$的线性映射$\sigma$使得$\sigma(1,-1,1)=(1,0)$，$\sigma(1,1,1)=(0,1)$？
\end{exercise}
\noindent 来源：\href{https://github.com/yhwu-is/Linear-Algebra-Left-Undone/blob/new/%E8%AE%B2%E4%B9%89/%E4%B8%93%E9%A2%98/5%20%E7%BA%BF%E6%80%A7%E6%98%A0%E5%B0%84.tex}{L5-A-3}。
\begin{solve}
存在. 设
          \begin{align}
              \sigma(1, -1, 1) & = \sigma(\vec{e}_1 - \vec{e}_2 + \vec{e}_3) \notag                 \\
                               & = \sigma(\vec{e}_1) - \sigma(\vec{e}_2) + \sigma(\vec{e}_3) \notag \\
                               & = (1, 0) = \vec{\varepsilon}_1 \label{eq:5:A:3:1}
          \end{align} \\
          \begin{align}
              \sigma(1, 1, 1) & = \sigma(\vec{e}_1 + \vec{e}_2 + \vec{e}_3) \notag                 \\
                              & = \sigma(\vec{e}_1) + \sigma(\vec{e}_2) + \sigma(\vec{e}_3) \notag \\
                              & = (0, 1) = \vec{\varepsilon}_2 \label{eq:5:A:3:2}
          \end{align}
          可取 $ \sigma(\vec{e}_3) = (0, 0) $. 联立式~\eqref{eq:5:A:3:1} 与式~\eqref{eq:5:A:3:2} 解得
          \begin{align*}
              \sigma(\vec{e}_1) & = \frac{1}{2} (\vec{\varepsilon}_1 + \vec{\varepsilon}_2) \\
              \sigma(\vec{e}_2) & = \frac{1}{2} (\vec{\varepsilon}_1 - \vec{\varepsilon}_2)
          \end{align*}
          因此存在线性映射 $ \sigma : \mathbf{R}^3 \to \mathbf{R}^2 $ 满足题设条件，其关于 $ \mathbf{R}^3 $ 基的像为
          \begin{align*}
              \sigma(\vec{e}_1) & = \frac{1}{2} (\vec{\varepsilon}_1 + \vec{\varepsilon}_2) \\
              \sigma(\vec{e}_2) & = \frac{1}{2} (\vec{\varepsilon}_1 - \vec{\varepsilon}_2) \\
              \sigma(\vec{e}_3) & = (0, 0)
          \end{align*}
\end{solve}

\begin{exercise}[补充训练 L5-A-4]\label{prob:L5-A-4}
是否存在$\mathbf{R}^2$到$\mathbf{R}^3$的线性映射$\sigma$使得$\sigma(3,2)=(1,0,0)$，$\sigma(1,5)=(1,1,0)$，$\sigma(-1,4)=(1,1,1)$？
\end{exercise}
\noindent 来源：\href{https://github.com/yhwu-is/Linear-Algebra-Left-Undone/blob/new/%E8%AE%B2%E4%B9%89/%E4%B8%93%E9%A2%98/5%20%E7%BA%BF%E6%80%A7%E6%98%A0%E5%B0%84.tex}{L5-A-4}。
\begin{solve}
不存在. 不存在从 $ \mathbf{R}^2 $ 到 $ \mathbf{R}^3 $ 的满射.
\end{solve}

\begin{exercise}[补充训练 L5-A-5]\label{prob:L5-A-5}
求$\sigma(x_1,x_2,\ldots,x_n)=(x_1,0,\ldots,0)$的像、核与秩.
\end{exercise}
\noindent 来源：\href{https://github.com/yhwu-is/Linear-Algebra-Left-Undone/blob/new/%E8%AE%B2%E4%B9%89/%E4%B8%93%E9%A2%98/5%20%E7%BA%BF%E6%80%A7%E6%98%A0%E5%B0%84.tex}{L5-A-5}。
\begin{solve}
\[ \sigma(x_1, x_2, \ldots, x_n) = x_1(1, 0, \ldots, 0) \]
          所以
          \begin{gather*}
              \sigma(V) = \operatorname{span}(\vec{e}_1) \\
              r(\sigma) = 1
          \end{gather*}
          由于 $ \sigma(\vec{e}_1) = \vec{e}_1,\enspace \sigma(\vec{e}_2) = \cdots = \sigma(\vec{e}_n) = \vec{0} $，所以
          \begin{gather*}
              \ker \sigma = \operatorname{span}(\vec{e}_2, \vec{e}_3, \ldots, \vec{e}_n) \\
              \dim \ker \sigma = n - 1
          \end{gather*}
\end{solve}

\begin{exercise}[补充训练 L5-A-6]\label{prob:L5-A-6}
证明：同构映射的逆、复合仍然是同构映射.
\end{exercise}
\noindent 来源：\href{https://github.com/yhwu-is/Linear-Algebra-Left-Undone/blob/new/%E8%AE%B2%E4%B9%89/%E4%B8%93%E9%A2%98/5%20%E7%BA%BF%E6%80%A7%E6%98%A0%E5%B0%84.tex}{L5-A-6}。
\begin{solve}
其逆映射是同构映射是显然的.我们只证明复合映射

            对于两个同构映射 $\sigma: V_1 \to V_2, \tau: V_2 \to V_3$，我们有其核空间
            $\ker \sigma \tau$,$\tau$是单的，将$0$映射到$0$，$\sigma$也是单的，故$\sigma\tau$是单的，将$0$映射到$0$；
            $\sigma$是满的，$\tau$也是满的，故$\sigma\tau$是满的，故$\sigma\tau$是双射，是同构映射.


            如果从矩阵的角度来理解：
            同构映射对应的矩阵是可逆矩阵；
            可逆矩阵的逆矩阵，可逆矩阵的乘积矩阵仍然是可逆矩阵.
\end{solve}

\begin{exercise}[补充训练 L5-B-1]\label{prob:L5-B-1}
已知
        \begin{gather*}
            \alpha_1=(1,-1),\enspace\alpha_2=(2,-1),\enspace\alpha_3=(-3,2) \\
            \beta_1=(1,0),\enspace\beta_2=(0,1),\enspace\beta_3=(1,1)
        \end{gather*}
        是否存在$\sigma\in \mathcal{L}(\mathbf{R}^2,\mathbf{R}^2)$，使得$\sigma(\alpha_i)=\beta_i,\enspace i=1,2,3$？
\end{exercise}
\noindent 来源：\href{https://github.com/yhwu-is/Linear-Algebra-Left-Undone/blob/new/%E8%AE%B2%E4%B9%89/%E4%B8%93%E9%A2%98/5%20%E7%BA%BF%E6%80%A7%E6%98%A0%E5%B0%84.tex}{L5-B-1}。
\begin{solve}
不存在. 假设存在，则由 $ \alpha_1 + \alpha_2 + \alpha_3 = \vec{0} $ 有
          \begin{align*}
              \sigma(\vec{0}) & = \sigma(\alpha_1 + \alpha_2 + \alpha_3)                 \\
                              & = \sigma(\alpha_1) + \sigma(\alpha_2) + \sigma(\alpha_3) \\
                              & = \beta_1 + \beta_2 + \beta_3 = \vec{0}
          \end{align*}
          这与 $ \beta_1 + \beta_2 + \beta_3 = (2, 2) $ 矛盾.
\end{solve}

\begin{exercise}[补充训练 L5-B-2]\label{prob:L5-B-2}
设$\alpha_1,\alpha_2$是线性空间$V(\mathbf{F})$的一组基，$x_1\alpha_1+x_2\alpha_2 \in V$. 定义$T(x_1\alpha_1+x_2\alpha_2)=r_1x_1\alpha_1+r_2x_2\alpha_2$，其中$r_1,r_2$是域$\mathbf{F}$中的两个常数. 证明：$T$是$V$上的一个线性变换. 当$V=\mathbf{R}^2$时，说明$T$的几何意义.
\end{exercise}
\noindent 来源：\href{https://github.com/yhwu-is/Linear-Algebra-Left-Undone/blob/new/%E8%AE%B2%E4%B9%89/%E4%B8%93%E9%A2%98/5%20%E7%BA%BF%E6%80%A7%E6%98%A0%E5%B0%84.tex}{L5-B-2}。
\begin{solve}
$ \forall x_1, x_2, y_1, y_2, k_1, k_2 \in \mathbf{F} $，有
          \begin{align*}
                  & T(k_1(x_1 \alpha_1 + x_2 \alpha_2) + k_2(y_1 \alpha_1 + y_2 \alpha_2))                \\
              ={} & r_1(k_1 x_1 + k_2 y_1) \alpha_1 + r_2(k_1 x_2 + k_2 y_2) \alpha_2                     \\
              ={} & k_1 (r_1 x_1 \alpha_1 + r_2 x_2 \alpha_2) + k_2 (r_1 y_1 \alpha_1 + r_2 y_2 \alpha_2) \\
              ={} & k_1 T(x_1 \alpha_1 + x_2 \alpha_2) + k_2 T(y_1 \alpha_1 + y_2 \alpha_2)
          \end{align*}
          因此 $ T $ 是线性映射. 其几何意义：将 $ \mathbf{R}^2 $ 中向量沿 $ x, y $ 轴分别变为原来的 $ r_1, r_2 $ 倍.
\end{solve}

\begin{exercise}[补充训练 L5-B-3]\label{prob:L5-B-3}
已知$\mathbf{R}$上的线性变换$\sigma(x_1,x_2)=(x_1-x_2,x_1+x_2)$，$\tau(x_1,x_2)=(x_1-x_2,x_2-x_1)$.
        \begin{enumerate}
            \item 求$\sigma^2(x_1,x_2)$；

            \item $\sigma$是否可逆？如可逆，求$\sigma^{-1}(x_1,x_2)$；

            \item 求$\xi\in \mathcal{L}(\mathbf{R}^2,\mathbf{R}^2)$，使得$\xi\tau=\theta$（零变换）.
        \end{enumerate}
\end{exercise}
\noindent 来源：\href{https://github.com/yhwu-is/Linear-Algebra-Left-Undone/blob/new/%E8%AE%B2%E4%B9%89/%E4%B8%93%E9%A2%98/5%20%E7%BA%BF%E6%80%A7%E6%98%A0%E5%B0%84.tex}{L5-B-3}。
\begin{solve}
\begin{enumerate}
                \item \begin{align*}
                          \sigma^2(x_1, x_2) & = \sigma(\sigma(x_1, x_2))                               \\
                                             & = \sigma(x_1 - x_2, x_1 + x_2)                           \\
                                             & = ((x_1 - x_2) - (x_1 + x_2), (x_1 - x_2) + (x_1 + x_2)) \\
                                             & = (-2 x_2, 2 x_1)
                      \end{align*}

                \item $ \sigma $ 可逆的充分必要条件是存在线性变换 $ \tau $ 使得 $ \sigma \tau = I $. 于是，$ \forall \alpha \in \mathbf{R}^2 $，令 $ (\tau \sigma)(\alpha) = \alpha $，即
                      \[ (\tau \sigma)(x_1, x_2) = \tau(x_1 - x_2, x_1 + x_2) = \tau(y_1, y_2) = (x_1, x_2) \]
                      解得
                      \begin{gather*}
                          x_1 = \frac{y_1 + y_2}{2} \\
                          x_2 = \frac{y_1 - y_2}{2}
                      \end{gather*}
                      所以
                      \[ \tau(x_1, x_2) = \sigma^{-1}(x_1, x_2) = \left(\frac{x_1 + x_2}{2}, \frac{x_1 - x_2}{2}\right) \]

                \item 当 $ \xi = \theta $ 时，显然满足 $ \xi \tau = \theta $. 当 $ \xi \neq \theta $ 时，
                      \[ \xi \tau(x_1, x_2) = \xi(x_1 - x_2, x_2 - x_1) = (0, 0) \]
                      记 $ y_1 = x_1 - x_2,\enspace y_2 = x_2 - x_1 $. 由于 $ y_1 + y_2 = 0 $，
                      \[ \xi(y_1, y_2) = (y_1 + y_2, y_1 + y_2) \]
                      即 $ \xi(x_1, x_2) = (x_1 + x_2, x_1 + x_2) $ 满足 $ \xi \tau = \theta $.
            \end{enumerate}
\end{solve}

\begin{exercise}[补充训练 L5-B-4]\label{prob:L5-B-4}
已知$\mathbf{R}^3$上的两个线性变换$\sigma,\tau$为：
        \begin{gather*}
            \sigma(x_1,x_2,x_3)=(x_3,0,0) \\
            \tau(x_1,x_2,x_3)=(x_1+x_2+x_3,x_1-x_2,0)
        \end{gather*}
        \begin{enumerate}
            \item 求$r(\sigma),\enspace r(\tau),\enspace \operatorname{im}\sigma,\enspace \ker\sigma$；

            \item 求$r(\tau\sigma),\enspace r(\sigma\tau),\enspace r(\sigma+\tau)$；

            \item 求$\operatorname{im}\tau+\ker\tau$.
        \end{enumerate}
\end{exercise}
\noindent 来源：\href{https://github.com/yhwu-is/Linear-Algebra-Left-Undone/blob/new/%E8%AE%B2%E4%B9%89/%E4%B8%93%E9%A2%98/5%20%E7%BA%BF%E6%80%A7%E6%98%A0%E5%B0%84.tex}{L5-B-4}。
\begin{solve}
\begin{enumerate}
                \item \begin{align*}
                          r(\sigma)   & = 1                          \\
                          r(\tau)     & = 2                          \\
                          \operatorname{im} \sigma  & = \operatorname{span}(\vec{e}_1)            \\
                          \ker \sigma & = \operatorname{span}(\vec{e}_1, \vec{e}_2)
                      \end{align*}

                \item \begin{gather*}
                          (\sigma \tau)(x_1, x_2, x_3) = \sigma(x_1 + x_2 + x_3, x_1 - x_2, 0) = (0, 0, 0) \\
                          r(\sigma \tau) = 0
                      \end{gather*}
                      \begin{gather*}
                          (\tau \sigma)(x_1, x_2, x_3) = \tau(x_3, 0, 0) = (x_3, x_3, 0) \\
                          \tau \sigma \neq \theta = \sigma \tau
                      \end{gather*}
                      求 $ r(\tau \sigma) $，方法 1：
                      \begin{gather*}
                          (\tau \sigma)(V) = \operatorname{span}((1, 1, 0)) \\
                          r(\tau \sigma) = 1
                      \end{gather*}
                      方法 2：由 $ \tau \sigma \neq \theta $ 知
                      \[ 1 \leqslant r(\tau \sigma) \leqslant \min\{r(\tau), r(\sigma)\} = 1 \]
                      因此 $ r(\tau \sigma) = 1 $.

                      \begin{gather*}
                          \begin{aligned}
                              (\sigma + \tau)(x_1, x_2, x_3) & = \sigma(x_1, x_2, x_3) + \tau(x_1, x_2, x_3) \\
                                                             & = (x_1 + x_2 + 2 x_3, x_1 - x_2, 0)           \\
                                                             & = x_1(1, 1, 0) + x_2(1, -1, 0) + x_3(2, 0, 0)
                          \end{aligned} \\
                          r(\sigma + \tau) = r\{(1, 1, 0), (1, -1, 0), (2, 0, 0)\} = 2
                      \end{gather*}

                \item \begin{gather*}
                          \begin{aligned}
                              \operatorname{im} \tau & = \operatorname{span}((1, 1, 0), (1, -1, 0), (2, 0, 0)) \\
                                       & = \operatorname{span}((1, 1, 0), (1, -1, 0))
                          \end{aligned} \\
                          \ker \tau = \operatorname{span}((1, 1, -2))
                      \end{gather*}
                      可得
                      \[ \operatorname{im} \tau + \ker \tau = \operatorname{span}((1, 1, 0), (1, -1, 0), (1, 1, -2)) = \mathbf{R}^3 \]
            \end{enumerate}
\end{solve}

\begin{exercise}[补充训练 L5-B-5]\label{prob:L5-B-5}
设 $\sigma$ 是线性空间 $V$ 上的线性变换，如果 $\sigma^{k-1}(\alpha) \neq 0$，但 $\sigma^{k}(\alpha) = 0$，证明：\\
        $\alpha,\sigma(\alpha),\ldots,\sigma^{k-1}(\alpha)\enspace(k>0)$ 线性无关（本题还有对应的矩阵版本，解法基本一致）.
\end{exercise}
\noindent 来源：\href{https://github.com/yhwu-is/Linear-Algebra-Left-Undone/blob/new/%E8%AE%B2%E4%B9%89/%E4%B8%93%E9%A2%98/5%20%E7%BA%BF%E6%80%A7%E6%98%A0%E5%B0%84.tex}{L5-B-5}。
\begin{solve}
使用反证法. 假设 $ \alpha, \sigma(\alpha), \ldots, \sigma^{k - 1}(\alpha) $ 线性相关，则存在不全为 0 的 $ l_0, l_1, \ldots, l_{k - 1} $ 使得
          \[ l_0 \alpha + l_1 \sigma(\alpha) + \cdots + l_{k - 1} \sigma^{k - 1}(\alpha) = \vec{0} \]
          记 $ i = \min\{k \mid l_k \neq 0\} $，则
          \[ l_i \sigma^{i}(\alpha) + l_{i + 1} \sigma^{i + 1}(\alpha) + \cdots + l_{k - 1} \sigma^{k - 1}(\alpha) = \vec{0} \]
          两边同时作用 $ \sigma^{k - i - 1} $. 由于 $ \sigma^{k}(\alpha) = \vec{0} $，故 $ \sigma^{k + 1}(\alpha) = \sigma^{k + 2}(\alpha) = \cdots = \vec{0} $，上式成为
          \[ l_i \sigma^{k - 1}(\alpha) = \vec{0} \]
          而 $ l_i \neq 0,\enspace \sigma^{k - 1}(\alpha) \neq \vec{0} $，矛盾！
\end{solve}

\begin{exercise}[补充训练 L5-B-6]\label{prob:L5-B-6}
设$\mathbf{R}[x]_3$是次数小于3的实系数多项式和全体零多项式一起组成的集合关于多项式加法和数乘多项式运算构成的实数域上的线性空间.
        \begin{enumerate}
            \item 证明：$W=\{f(x)\in \mathbf{R}[x]_3 \mid f(1)=0\}$是$\mathbf{R}[x]_3$的一个子空间，并求$W$的维数和一组基；

            \item 定义从$\mathbf{R}[x]_3$到$\mathbf{R}$的线性映射$\sigma(f(x))=f(1)$，证明：$\sigma$为线性映射，并求$\operatorname{im}\sigma$和$\dim\ker\sigma$；

            \item 设$f,g,h \in \mathbf{R}[x]_3$且$f(1)=g(1)=h(1)=0$，证明：$f,g,h$线性相关.
        \end{enumerate}
\end{exercise}
\noindent 来源：\href{https://github.com/yhwu-is/Linear-Algebra-Left-Undone/blob/new/%E8%AE%B2%E4%B9%89/%E4%B8%93%E9%A2%98/5%20%E7%BA%BF%E6%80%A7%E6%98%A0%E5%B0%84.tex}{L5-B-6}。
\begin{solve}
\begin{enumerate}
\item 零多项式在 $W$ 中；若 $f,g\in W$ 且 $a,b\in\mathbf R$，则 $(af+bg)(1)=af(1)+bg(1)=0$，所以 $W$ 是子空间。任意 $f=ax^2+bx+c\in W$ 满足 $c=-a-b$，从而 $f=a(x^2-1)+b(x-1)$。两个生成多项式比较最高次项即可证明线性无关，故它们是一组基且 $\dim W=2$。
\item 对任意标量 $a,b$，$\sigma(af+bg)=a\sigma(f)+b\sigma(g)$，故 $\sigma$ 线性。常值多项式可取任意实数，故 $\operatorname{im}\sigma=\mathbf R$；$\ker\sigma=W$，故 $\dim\ker\sigma=2$。
\item 三个多项式均属于二维空间 $W$，而其中任意三个向量必线性相关。\end{enumerate}
\end{solve}

\begin{exercise}[补充训练 L5-B-7]\label{prob:L5-B-7}
设$\sigma(p(x))=p'(x)$（求导），$\forall p(x) \in \mathbf{R}[x]_n$.
        \begin{enumerate}
            \item 证明：$\sigma$是$\mathbf{R}[x]_n$上的线性变换；

            \item 求$\sigma$的值域和$r(\sigma)$，说明$\sigma$是否可逆；

            \item 求$\sigma$的核及其维数；

            \item 求$r(\sigma)+\dim\ker\sigma$，问：$\mathbf{R}[x]_n=\ker\sigma+\operatorname{im} \sigma$是否成立.
        \end{enumerate}
\end{exercise}
\noindent 来源：\href{https://github.com/yhwu-is/Linear-Algebra-Left-Undone/blob/new/%E8%AE%B2%E4%B9%89/%E4%B8%93%E9%A2%98/5%20%E7%BA%BF%E6%80%A7%E6%98%A0%E5%B0%84.tex}{L5-B-7}。
\begin{solve}
\begin{enumerate}
                \item 证明：$ \forall p_1(x), p_2(x) \in \mathbf{R}[x]_n,\enspace \forall k_1, k_2 \in \mathbf{R} $，有
                      \begin{align*}
                          \sigma(k_1 p_1 + k_2 p_2) & = (k_1 p_1(x) + k_2 p_2(x))'                                    \\
                                                    & = k_1 p_1'(x) + k_2 p_2'(x) = k_1 \sigma(p_1) + k_2 \sigma(p_2)
                      \end{align*}
                      因此 $ \sigma $ 是 $ \mathbf{R}[x]_n $ 上的线性变换.

                \item \begin{gather*}
                          \begin{aligned}
                              \operatorname{im} \sigma & = \operatorname{span}(\sigma(1), \sigma(x), \sigma(x^2), \ldots, \sigma(x^{n - 1})) \\
                                         & = \operatorname{span}(1, 2x, \ldots, (n - 1) x^{n - 2})                             \\
                                         & = \operatorname{span}(1, x, \ldots, x^{n - 2})
                          \end{aligned} \\
                          r(\sigma) = n - 1
                      \end{gather*}
                      可知 $ \sigma $ 不是单射，因此不可逆.

                \item 由 $ \sigma(p(x)) = 0 $ 可知 $ p(x) = c $（常数）. 因此 $ \ker \sigma = \operatorname{span}(1),\enspace \dim \ker \sigma = 1 $.

                \item \begin{gather*}
                          r(\sigma) + \dim \ker \sigma = (n - 1) + 1 = n \\
                          \operatorname{im} \sigma + \ker \sigma = \operatorname{span}(1, x, \ldots, x^{n - 2}) = \operatorname{im} \sigma \neq \mathbf{R}[x]_n
                      \end{gather*}
            \end{enumerate}
\end{solve}

\begin{exercise}[补充训练 L5-B-8]\label{prob:L5-B-8}
设$V$为有限维线性空间，$T\in \mathcal{L}(V,V)$且$T$不是恒等变换也不是零变换，问：下列情况是否可能发生？如果可能请举例，不可能请说明理由.
        \begin{enumerate}
            \item $\operatorname{im} T \cap \ker T = \{0\}$；

            \item $\operatorname{im} T \subseteq \ker T$；

            \item $\ker T = \operatorname{im} T$；

            \item $\ker T \subseteq \operatorname{im} T$.
        \end{enumerate}
\end{exercise}
\noindent 来源：\href{https://github.com/yhwu-is/Linear-Algebra-Left-Undone/blob/new/%E8%AE%B2%E4%B9%89/%E4%B8%93%E9%A2%98/5%20%E7%BA%BF%E6%80%A7%E6%98%A0%E5%B0%84.tex}{L5-B-8}。
\begin{solve}
\begin{enumerate}
                \item 可能. 例如 $ \sigma(x, y) = (x + y, x + y) $.
                      \begin{gather*}
                          \operatorname{im} \sigma = \operatorname{span}(\vec{e}_1 + \vec{e}_2) \\
                          \ker \sigma = \operatorname{span}(\vec{e}_1 - \vec{e}_2) \\
                          \operatorname{im} \sigma \cap \ker \sigma = \{\vec{0}\}
                      \end{gather*}

                \item 可能. 例如 $ \sigma(x, y, z) = (x - y, x - y, x - y) $.
                      \begin{gather*}
                          \operatorname{im} \sigma = \operatorname{span}((1, 1, 1)) \\
                          \ker \sigma = \operatorname{span}((1, 1, 1), (1, 1, 0)) \\
                          \operatorname{im} \sigma \subseteq \ker \sigma
                      \end{gather*}

                \item 可能. 例如 $ \sigma(x, y) = (x - y, x - y) $.
                      \begin{gather*}
                          \operatorname{im} \sigma = \ker \sigma = \operatorname{span}((1, 1))
                      \end{gather*}

                \item 可能. 例如 $ \sigma(x, y) = (x, x - y) $.
                      \begin{gather*}
                          \operatorname{im} \sigma = \mathbf{R}^2 \\
                          \ker \sigma = \{\vec{0}\} \\
                          \ker \sigma \subseteq \operatorname{im} \sigma
                      \end{gather*}
            \end{enumerate}
\end{solve}

\begin{exercise}[补充训练 L5-B-9]\label{prob:L5-B-9}
若$\sigma_1,\sigma_2\in \mathcal{L}(V,V)$，判断下列说法是否正确，正确请给出证明，反之给出反例：
        \begin{enumerate}
            \item 由$r(\sigma)+\dim\ker\sigma=n$可知$V=\ker\sigma+\operatorname{im} \sigma$；

            \item 若有$\operatorname{im} \sigma \cap \ker \sigma = \{0\}$，则$V=\ker\sigma+\operatorname{im} \sigma$成立；

            \item 因为$\forall \alpha \in V$有$(\sigma_1+\sigma_2)(\alpha)=\sigma_1(\alpha)+\sigma_2(\alpha)$，所以$(\sigma_1+\sigma_2)(V)=\sigma_1(V)+\sigma_2(V)$；

            \item $(I-\sigma)(V)+\sigma(V)=V$，其中$I$为恒等映射.
        \end{enumerate}
\end{exercise}
\noindent 来源：\href{https://github.com/yhwu-is/Linear-Algebra-Left-Undone/blob/new/%E8%AE%B2%E4%B9%89/%E4%B8%93%E9%A2%98/5%20%E7%BA%BF%E6%80%A7%E6%98%A0%E5%B0%84.tex}{L5-B-9}。
\begin{solve}
\begin{enumerate}
                \item 错误. 一个反例为 $ \sigma(x_1, x_2) = (x_1 - x_2, x_1 - x_2) $，则 $ \operatorname{im} \sigma + \ker \sigma = \operatorname{span}((1, 1)) $.

                \item 正确. $ \operatorname{im} \sigma + \ker \sigma \subseteq V $，此时 $ \dim(\operatorname{im} \sigma + \ker \sigma) = \dim \operatorname{im} \sigma + \dim \ker \sigma - 0 = \dim V $，所以 $ \operatorname{im} \sigma + \ker \sigma = V $.

                \item \label{item:6:B:2:3}
                      错误. $ \forall \alpha \in V $ 有
                      \[ (\sigma_1 + \sigma_2)(\alpha) = \sigma_1(\alpha) + \sigma_2(\alpha) \in \sigma_1(V) + \sigma_2(V) \]
                      所以
                      \[ (\sigma_1 + \sigma_2)(V) \subseteq \sigma_1(V) + \sigma_2(V) \]
                      但上式中等号不一定成立. 反例：$ V = \mathbf{R}^3 $ 上的线性变换 $ \sigma_1, \sigma_2 $ 关于 $ \mathbf{R}^3 $ 的基 $ \vec{e}_1, \vec{e}_2, \vec{e}_3 $ 的像分别为
                      \begin{gather*}
                          \sigma_1(\vec{e}_1) = \vec{e}_1,\enspace \sigma_1(\vec{e}_2) = \sigma_1(\vec{e}_3) = \vec{e}_2 \\
                          \sigma_2(\vec{e}_1) = \vec{e}_1,\enspace \sigma_2(\vec{e}_2) = \sigma_2(\vec{e}_3) = \vec{e}_3 \\
                          \sigma_1(V) = \operatorname{span}(\vec{e}_1, \vec{e}_2),\enspace \sigma_2(V) = \operatorname{span}(\vec{e}_1, \vec{e}_3) \\
                          \sigma_1(V) + \sigma_2(V) = \operatorname{span}(\vec{e}_1, \vec{e}_2, \vec{e}_3) = \mathbf{R}^3
                      \end{gather*}
                      而
                      \begin{gather*}
                          (\sigma_1 + \sigma_2)(\vec{e}_1) = \vec{e}_1 + \vec{e}_1 = 2 \vec{e}_1 \\
                          (\sigma_1 + \sigma_2)(\vec{e}_2) = (\sigma_1 + \sigma_2)(\vec{e}_3) = \vec{e}_2 + \vec{e}_3 \\
                          (\sigma_1 + \sigma_2)(V) = \operatorname{span}(\vec{e}_1, \vec{e}_2 + \vec{e}_3) \neq \mathbf{R}^3
                      \end{gather*}

                \item 正确. $(I - \sigma)(V) + \sigma(V) \subseteq V$ 是显然的，另一个方向上，$ \forall \alpha \in V $ 有
                      \[ \alpha = (I - \sigma)(\alpha) + \sigma(\alpha) \in (I - \sigma)(V) + \sigma(V) \]
                      所以 $ V \subseteq (I - \sigma)(V) + \sigma(V) $. 综上有 $ V = (I - \sigma)(V) + \sigma(V) $.
            \end{enumerate}
\end{solve}

\begin{exercise}[补充训练 L5-B-10]\label{prob:L5-B-10}
设$V(\mathbf{R})$是线性空间，$\sigma$是$V(\mathbf{R})$到$\mathbf{R}^3$的同构映射，且$\sigma(\alpha_1)=(1,0,1),\enspace\allowbreak\sigma(\alpha_2)=(-2,1,0),\enspace\allowbreak\sigma(\alpha_3)=(-3,2,1),\enspace\allowbreak\sigma(\alpha_4)=(1,1,2)$.
        \begin{enumerate}
            \item $\alpha_1$在$\operatorname{span}(\alpha_2,\alpha_3)$中吗？

            \item 设$W_1=\operatorname{span}(\alpha_1,\alpha_2),\enspace W_2=\operatorname{span}(\alpha_3,\alpha_4)$，求$W_1\cap W_2$.
        \end{enumerate}
\end{exercise}
\noindent 来源：\href{https://github.com/yhwu-is/Linear-Algebra-Left-Undone/blob/new/%E8%AE%B2%E4%B9%89/%E4%B8%93%E9%A2%98/5%20%E7%BA%BF%E6%80%A7%E6%98%A0%E5%B0%84.tex}{L5-B-10}。
\begin{solve}
\begin{enumerate}
                \item 在. 因为
                      \[ \sigma(\alpha_1) = -2 \sigma(\sigma_2) + \sigma(\alpha_3) = \sigma(-2 \alpha_2 + \alpha_3) \]
                      同构映射 $ \sigma $ 可逆. 所以
                      \[ \alpha_1 = -2 \alpha_2 + \alpha_3 \in \operatorname{span}(\alpha_2, \alpha_3) \]
                \item 对于注意力惊人的同学，直接观察当然可得答案；如果没什么思路，我们可以先从定义入手；
                由于同构是将基映射到基的，我们可以将对于$\alpha$的操作看做对于其像的操作，为了方便，在这里直接用 $a$ 来代替$\alpha$的像，若 $\beta \in W_1\cap W_2$，则
                \[ \beta = x_1 a_1+x_2 a_2=x'_3 a_3+x'_4 a_4 \implies  x_1 a_1+x_2 a_2-x'_3 a_3-x'_4 a_4 = 0 \]
                改写$x'_3,x'_4$为$x_3,x_4$，则转换为
                \[ x_1 a_1+x_2 a_2+x_3 a_3+x_4 a_4 = 0 \]
                对应的系数矩阵为
                \[ \begin{pmatrix}
                    1 & -2 & -3 & 1 \\
                    0 & 1 & 2 & 1\\
                    1 & 0 & 1 & 2
                \end{pmatrix} \]
                四个未知数，三个方程，解得其基础解系为
                \[ \begin{pmatrix}
                    -1 \\
                    -2 \\
                    1 \\
                    0
                \end{pmatrix} \]

                故$\beta = k(-a_1+-2a_2=-a_3)$其中$k\in \mathbf{R}$.故所求交空间的基为$\alpha_3$(注意要对应回$V(R)$中)
            \end{enumerate}
\end{solve}

\begin{exercise}[补充训练 L5-B-11]\label{prob:L5-B-11}
设$c_1,c_2,\ldots,c_n$是$n$个互异的实常数. 证明：$\mathbf{R}[x]_n$到$\mathbf{R}^n$的一个映射$\sigma$：
        \[\sigma(p(x))=(p(c_1),p(c_2),\ldots,p(c_n))\]
        是$\mathbf{R}[x]_n$到$\mathbf{R}^n$的一个同构映射.
\end{exercise}
\noindent 来源：\href{https://github.com/yhwu-is/Linear-Algebra-Left-Undone/blob/new/%E8%AE%B2%E4%B9%89/%E4%B8%93%E9%A2%98/5%20%E7%BA%BF%E6%80%A7%E6%98%A0%E5%B0%84.tex}{L5-B-11}。
\begin{solve}
我们仅对 $ n = 3 $ 的情况给出证明. 

          先证 $ \sigma $ 是线性映射. $ \forall p(x), q(x) \in \mathbf{R}[x]_3,\enspace \forall k_1, k_2 \in \mathbf{R} $ 有
          \begin{align*}
                  & \sigma(k_1 p(x) + k_2 q(x))                                                 \\
              ={} & (k_1 p(c_1) + k_2 q(c_1), k_1 p(c_2) + k_2 q(c_2), k_1 p(c_3) + k_2 q(c_3)) \\
              ={} & k_1(p(c_1), p(c_2), p(c_3)) + k_1(q(c_1), q(c_2), q(c_3))                   \\
              ={} & k_1 \sigma(p(x)) + k_2 \sigma(q(x))
          \end{align*}

          再证 $ \sigma $ 是双射，即 $ \forall (d_1, d_2, d_3) \in \mathbf{R}^3 $，存在唯一的
          \[ p(x) = a + bx + cx^2 \in \mathbf{R}[x]_3 \]
          使 $ \sigma(p(x)) = (d_1, d_2, d_3) $. 根据
          \[ \sigma(p(x)) = (p(c_1), p(c_2), p(c_3)) \]
          以及 $ \sigma(p(x)) = (d_1, d_2, d_3) $，有
          \[ \begin{cases}
                  a + bc_1 + cc_1^2 = d_1 \\
                  a + bc_2 + cc_2^2 = d_2 \\
                  a + bc_3 + cc_3^3 = d_3
              \end{cases} \]
          方程组是关于未知元 $ a, b, c $ 的三元线性非齐次方程组，其中 $ c_1, c_2, c_3 $ 是互异的实常数. 用高斯-若当消元法，易将其增广矩阵变换为下列阶梯形矩阵，即
          \begin{align}
                             & \begin{pmatrix}[ccc|c]
                                   1 & c_1 & c_1^2 & d_1 \\
                                   1 & c_2 & c_2^2 & d_2 \\
                                   1 & c_3 & c_3^3 & d_3
                               \end{pmatrix} \notag                                                                    \\
              \implies \quad & \begin{pmatrix}[ccc|c]
                                   1 & c_1 & c_1^2     & d_1                                                         \\[1ex]
                                   0 & 1   & c_2 + c_1 & \dfrac{d_1 - d_2}{c_1 - c_2}                                \\[2ex]
                                   0 & 0   & c_3 - c_2 & \dfrac{d_3 - d_1}{c_3 - c_1} - \dfrac{d_2 - d_1}{c_2 - c_1}
                               \end{pmatrix} \label{item:6:B:6:1}
          \end{align}
          阶梯性矩阵~\eqref{item:6:B:6:1} （其中 $ c_3 - c_2, c_3 - c_1, c_2 - c_1 $ 均为非零常数）对应的方程组有唯一解 $ a, b, c $，即存在唯一的
          \[ p(x) = a + bx + cx^2 \in \mathbf{R}[x]_3 \]
          使得 $ \sigma(p(x)) = (d_1, d_2, d_3) $ 成立. 所以 $ \sigma $ 是线性双射，即 $ \mathbf{R}[x]_3 $ 到 $ \mathbf{R}^3 $ 的同构映射.
\end{solve}

\begin{exercise}[补充训练 L5-B-12]\label{prob:L5-B-12}
设$\sigma$和$\tau$分别为有限维线性空间$U\to V$和$V\to W$的线性映射，证明
        \[\dim\ker\sigma+\dim(\operatorname{im}\sigma\cap\ker\tau)=\dim\ker(\tau\sigma).\]
\end{exercise}
\noindent 来源：\href{https://github.com/yhwu-is/Linear-Algebra-Left-Undone/blob/new/%E8%AE%B2%E4%B9%89/%E4%B8%93%E9%A2%98/5%20%E7%BA%BF%E6%80%A7%E6%98%A0%E5%B0%84.tex}{L5-B-12}。
\begin{solve}
我们先观察等式的右边，在 $\ker \tau \sigma$中的元素 $u$,会满足 $\tau \sigma (u) = 0 $,仔细想想什么元素会满足这个式子
            \begin{itemize}
                \item $\sigma(u)=0$的元素，亦即$u\in \ker \sigma$；
                \item $v=\sigma(u)\neq 0$，但是$\tau(v)=0$的元素，亦即$v\in \operatorname{im} \sigma \cap \ker \tau$.
            \end{itemize}
            这恰好对应左边的两项，所以等式显然是成立的.
\end{solve}

\begin{exercise}[补充训练 L5-C-1]\label{prob:L5-C-1}
设 $V(\mathbf{F})$ 是一个 $n$ 维线性空间，$\sigma \in \mathcal{L}(V,V)$. 证明：
        \begin{enumerate}
            \item 在 $\mathbf{F}[x]$ 中有一个次数不高于 $n^2$ 的非零多项式 $p(x)$ 使 $p(\sigma) = \theta$；

            \item $\sigma$ 可逆$\iff$有一常数项不为 0 的多项式 $p(x)$ 使 $p(\sigma) = \theta$.
        \end{enumerate}
\end{exercise}
\noindent 来源：\href{https://github.com/yhwu-is/Linear-Algebra-Left-Undone/blob/new/%E8%AE%B2%E4%B9%89/%E4%B8%93%E9%A2%98/5%20%E7%BA%BF%E6%80%A7%E6%98%A0%E5%B0%84.tex}{L5-C-1}。
\begin{solve}
\begin{enumerate}
                \item \label{item:5:C:1:1}
                      $ \mathcal{L}(V, V) $ 是 $ n^2 $ 维线性空间，其中 $ n^2 + 1 $ 个元素 $ I, \sigma, \sigma^2, \ldots, \sigma^{n^2} $ 线性相关，即 $ \exists a_i\enspace (i = 1, 2, \ldots, n^2) $ 使得
                      \[ a_0 I + a_1 \sigma + \ldots + a_{n^2} \sigma^{n^2} = \theta \]
                      于是 $ \exists p(x) = a_0 + a_1 x + \ldots + a_{n^2} x^{n^2} \in \mathbf{F}[x] $ 使得 $ p(\sigma) = \theta $.

                \item 必要性：设有一常数项不为 0 的多项式
                      \[ p(x) = a_0 + a_1 x + \cdots + a_k x^k \qquad a_0 \neq 0 \]
                      满足
                      \[ p(\sigma) = a_0 I + a_1 \sigma + \cdots + a_k \sigma^k = \theta \]
                      所以
                      \[ \sigma(a_1 I + a_2 \sigma + \cdots + a_k \sigma^{k - 1}) = -a_0 I \]
                      因此
                      \[ \sigma^{-1} = -a_0^{-1} (a_1 I + a_2 \sigma + \cdots + a_k \sigma^{k - 1}) \]

                      充分性：由 \eqref{item:5:C:1:1} 可知存在次数 $ k \leqslant n^2 $ 的多项式
                      \[ p(x) = a_0 + a_1 x + a_2 x^2 + \cdots + a_k x^k \]
                      满足
                      \[ p(\sigma) = a_0 I + a_1 \sigma + \cdots + a_k \sigma^k = \theta \]
                      若 $ a_0 \neq 0 $，则 $ p(x) $ 即为所求多项式.

                      若 $ a_0 = a_1 = \cdots = a_{i - 1} = 0,\enspace a_i \neq 0 $，即
                      \[ a_i I + a_{i + 1} \sigma + \cdots + a_k \sigma^{k - i} = \theta \]
                      于是
                      \[ P(x) = a_i + a_{i + 1} x + \cdots + a_k x^{k - i} \qquad a_i \neq 0 \]
                      为所求的多项式.
            \end{enumerate}
\end{solve}

\begin{exercise}[补充训练 L5-C-2]\label{prob:L5-C-2}
已知$\sigma_1,\sigma_2,\ldots,\sigma_s$是线性空间$V$上的$s$个两两不同的线性变换，证明：在$V$中必存在向量$\alpha$使得$\sigma_1(\alpha),\sigma_2(\alpha),\ldots,\sigma_s(\alpha)$也两两不同.
\end{exercise}
\noindent 来源：\href{https://github.com/yhwu-is/Linear-Algebra-Left-Undone/blob/new/%E8%AE%B2%E4%B9%89/%E4%B8%93%E9%A2%98/5%20%E7%BA%BF%E6%80%A7%E6%98%A0%E5%B0%84.tex}{L5-C-2}。
\begin{solve}
$\alpha \in V$使得$\sigma_i(\alpha) \neq \sigma_j(\alpha)$当且仅当$\alpha \notin \ker(\sigma_i - \sigma_j)$,我们可以只考虑该核空间不是零空间的情况，因为我们所取出的$\alpha$显然不可能是零，
            又因为我们的线性映射是两两不相同的，所以该核空间一定是一个真子空间. 根据覆盖定理，对于$m$个非平凡子空间，我们总可以取到一个不属于任意一个这些子空间的向量，这样我们就找到了所需要的$\alpha$.
\end{solve}

\begin{exercise}[补充训练 L5-C-3]\label{prob:L5-C-3}
设$V$是一个$n$维线性空间，$V=W_1\oplus W_2,\enspace\sigma\in \mathcal{L}(V,V)$. 证明：$\sigma$可逆$\iff V=\sigma(W_1)+\sigma(W_2)$.
\end{exercise}
\noindent 来源：\href{https://github.com/yhwu-is/Linear-Algebra-Left-Undone/blob/new/%E8%AE%B2%E4%B9%89/%E4%B8%93%E9%A2%98/5%20%E7%BA%BF%E6%80%A7%E6%98%A0%E5%B0%84.tex}{L5-C-3}。
\begin{solve}
证明：必要性：$ \forall \alpha \in V $，由 $ \sigma $ 可逆，存在唯一的 $ \beta \in V $ 使得 $ \sigma(\beta) = \alpha $ 且 $ \beta = \beta_1 + \beta_2 $，其中 $ \beta_1 \in W_1 $，$ \beta_2 \in W_2 $. 于是
          \[ \alpha = \sigma(\beta) = \sigma(\beta_1) + \sigma(\beta_2) \in \sigma(W_1) + \sigma(W_2) \]
          所以 $ V \subseteq \sigma(W_1) + \sigma(W_2) $.

          $ \sigma(W_1), \sigma(W_2) $ 都是 $ V $ 的子空间，它们的和也是 $ V $ 的子空间. 所以 $ \sigma(W_1) + \sigma(W_2) \subseteq V $，故 $ V = \sigma(W_1) + \sigma(W_2) $.

          充分性：$ \forall \alpha \in V = \sigma(W_1) + \sigma(W_2) $，$ \exists \alpha_i \in \sigma(W_i) $ 且 $ \exists \beta_i \in W_i $ 使 $ \alpha_i = \sigma(\beta_i) $（$ i = 1, 2 $）使得
          \[ \alpha = \alpha_1 + \alpha_2 = \sigma(\beta_1) + \sigma(\beta_2) = \sigma(\beta_1 + \beta_2) = \sigma(\beta) \]
          其中 $ \beta = \beta_1 + \beta_2 \in W_1 + W_2 = V $，所以 $ \sigma $ 为满射.

          由于 $ n $ 维线性空间上的线性映射为满射时也必为单射，从而必是双射，所以 $ \sigma $ 可逆.
\end{solve}

\begin{exercise}[补充训练 L5-C-5]\label{prob:L5-C-5}
设$V_1$是有限维线性空间，$\sigma,\tau\in \mathcal{L}(V_1,V_2)$，则
        \[r(\sigma+\tau) \leqslant r(\sigma)+r(\tau).\]

        事实上前两题的结论在后续章节矩阵的秩中都会涉及，此处有兴趣的同学可以尝试从线性映射的角度理解这两个秩不等式. 如果无法找到合适方式，可以考虑化为矩阵进行证明.
\end{exercise}
\noindent 来源：\href{https://github.com/yhwu-is/Linear-Algebra-Left-Undone/blob/new/%E8%AE%B2%E4%B9%89/%E4%B8%93%E9%A2%98/5%20%E7%BA%BF%E6%80%A7%E6%98%A0%E5%B0%84.tex}{L5-C-5}。
\begin{solve}
证明：由于 $ \forall \beta \in (\sigma + \tau)(V_1),\enspace \exists \alpha \in V_1 $ 使 $ \beta = (\sigma + \tau)(\alpha) = \sigma(\alpha) + \tau(\alpha) \in \sigma(V_1) + \tau(V_1) $，所以
          \[ (\sigma + \tau)(V_1) \subseteq \sigma(V_1) + \tau(V_1) \]
          因此
          \begin{align*}
              \dim(\sigma + \tau)(V_1) & \leqslant \dim(\sigma(V_1) + \tau(V_1))     \\
                                       & \leqslant \dim \sigma(V_1) + \dim \tau(V_1)
          \end{align*}
\end{solve}

\begin{exercise}[补充训练 L5-C-6]\label{prob:L5-C-6}
设$\sigma\in \mathcal{L}(V,V)$，$\dim V=n$，且$\sigma^2=\sigma$，$I$是$V$上的恒等变换. 证明：
        \begin{enumerate}
            \item $(I-\sigma)(V) \subseteq \ker\sigma$；

            \item $r(I-\sigma)+r(\sigma)=n$.
        \end{enumerate}
\end{exercise}
\noindent 来源：\href{https://github.com/yhwu-is/Linear-Algebra-Left-Undone/blob/new/%E8%AE%B2%E4%B9%89/%E4%B8%93%E9%A2%98/5%20%E7%BA%BF%E6%80%A7%E6%98%A0%E5%B0%84.tex}{L5-C-6}。
\begin{solve}
\begin{enumerate}
                \item \label{item:5:C:6:1}
                      $ \forall \sigma \in \mathcal{L}(V, V) $，则 $ I - \sigma \in \mathcal{L}(V, V) $. $ \forall \alpha \in (I - \sigma)(V),\enspace \exists \beta \in V $，有
                      \begin{gather*}
                          \alpha = (I - \sigma)(\beta) = \beta - \sigma(\beta) \\
                          \sigma(\alpha) = \sigma(\beta - \sigma(\beta)) = \sigma(\beta) - \sigma^2(\beta)
                      \end{gather*}
                      而由于 $ \sigma^2 = \sigma $，所以 $ \sigma(\alpha) = \vec{0} $，于是 $ \alpha \in \ker \sigma $，因此 $ (I - \sigma)(V) \subseteq \ker \sigma $.

                \item 利用 $ r(\sigma + \tau) \leqslant r(\sigma) + r(\tau) $ 和 $ r(\sigma) + \dim \ker \sigma = n $，由 \eqref{item:5:C:6:1} 可得
                      \begin{equation} \label{eq:5:C:6:2:1}
                          r(I - \sigma) + r(\sigma) \leqslant n
                      \end{equation}
                      又因为
                      \begin{equation} \label{eq:5:C:6:2:2}
                          r(I - \sigma) + r(\sigma) \geqslant r(I - \sigma + \sigma) = r(I) = n
                      \end{equation}
                      于是由式~\eqref{eq:5:C:6:2:1} 和式~\eqref{eq:5:C:6:2:2} 即可得到 $ r(I - \sigma) + r(\sigma) = n $.
            \end{enumerate}
\end{solve}

\begin{exercise}[补充训练 L5-C-7]\label{prob:L5-C-7}
设$V$是一个$n$维线性空间，$\sigma\in \mathcal{L}(V,V)$，且$\sigma^2=\theta$（零映射）. 证明：
        \begin{enumerate}
            \item $ \dim \operatorname{im} \sigma \leqslant \dfrac{n}{2} $；

            \item 设$A$是$\sigma$在某组基下的矩阵，则方程组$AX=0$的基础解系至少有$\dfrac{n}{2}$个解.
        \end{enumerate}
\end{exercise}
\noindent 来源：\href{https://github.com/yhwu-is/Linear-Algebra-Left-Undone/blob/new/%E8%AE%B2%E4%B9%89/%E4%B8%93%E9%A2%98/5%20%E7%BA%BF%E6%80%A7%E6%98%A0%E5%B0%84.tex}{L5-C-7}。
\begin{solve}
证明：\begin{enumerate}
                \item \label{item:5:C:7:1}
                      $ \forall \alpha \in \operatorname{im} \sigma,\enspace \exists \beta \in V $ 使得 $ \sigma(\beta) = \alpha $. 由 $ \sigma^2 = \theta $ 可得 $ \sigma(\alpha) = \sigma^2(\beta) = \vec{0} $，因此 $ \alpha \in \ker \sigma $，从而 $ \operatorname{im} \sigma \subseteq \ker \sigma $. 于是我们得到
                      \[ n = \dim \operatorname{im} \sigma + \dim \ker \sigma \geqslant 2 \dim \operatorname{im} \sigma \]
                      即 $ \dim \operatorname{im} \sigma \leqslant \dfrac{n}{2} $.

                \item 由 \eqref{item:5:C:7:1} 可知，方程组 $ AX = \vec{0} $ 的基础解系含有 $ n - r(A) = \dim \ker \sigma \geqslant \dfrac{n}{2} $ 个解向量，所以结论成立.
            \end{enumerate}
\end{solve}

\begin{exercise}[补充训练 L5-例-4]\label{prob:L5-例-4}
写出下列映射的出发空间和到达空间，并判断其是否为线性映射：
    \begin{enumerate}
        \item $\sigma(x_1,x_2)=(x_1-x_2,x_1,x_1+x_2)$；

        \item $\sigma(x_1,x_2)=(x_1x_2,x_1+x_2)$；

        \item $\sigma(p(x))=p(x+1)-p(x),\enspace\forall p(x) \in \mathbf{R}[x]_n$；

        \item $\sigma(p(x))=p(a),\enspace\forall p(x)$，其中$a$为常数；

        \item $\sigma(\xi)=2\xi+\xi_0$，其中$\xi_0$是线性空间$V$中的一个固定向量.
    \end{enumerate}
\end{exercise}
\noindent 来源：\href{https://github.com/yhwu-is/Linear-Algebra-Left-Undone/blob/new/%E8%AE%B2%E4%B9%89/%E4%B8%93%E9%A2%98/5%20%E7%BA%BF%E6%80%A7%E6%98%A0%E5%B0%84.tex}{L5-例-4}。
\begin{solve}
\begin{enumerate}
        \item 出发空间为 $ \mathbf{R}^2 $，到达空间为 $ \mathbf{R}^3 $. $ \sigma $ 是线性映射.

              $ \forall (x_1, x_2), (y_1, y_2) \in \mathbf{R}^2,\enspace k_1, k_2 \in \mathbf{R} $, 有
              \begin{align*}
                      & \sigma(k_1(x_1, x_2) + k_2(y_1, y_2))                                                                     \\
                  ={} & ((k_1 x_1 + k_2 y_1) - (k_1 x_2 + k_2 y_2), k_1 x_1 + k_2 y_1, (k_1 x_1 + k_2 y_1) + (k_1 x_2 + k_2 y_2)) \\
                  ={} & k_1(x_1 - x_2, x_1, x_1 + x_2) + k_2(y_1 - y_2, y_1, y_1 + y_2)                                           \\
                  ={} & k_1 \sigma(x_1, x_2) + k_2 \sigma(y_1, y_2)
              \end{align*}

        \item 出发空间为 $ \mathbf{R}^2 $，到达空间为 $ \mathbf{R}^2 $. $ \sigma $ 不是线性映射.

              $ \forall (x_1, x_2), (y_1, y_2) \in \mathbf{R}^2,\enspace k_1, k_2 \in \mathbf{R} $, 有
              \begin{align*}
                  \sigma((x_1, x_2) + (y_1, y_2))
                   & = \sigma(x_1 + y_1, x_2 + y_2)                                        \\
                   & = ((x_1 + y_1)(x_2 + y_2), ((x_1 + y_1) + (x_2 + y_2))                \\
                   & = (x_1 x_2 + x_1 y_2 + y_1 x_2 + y_1 y_2, x_1 + y_1 + x_2 + y_2)      \\
                   & = (x_1 x_2 + y_1 y_2, x_1 + y_1 + x_2 + y_2) + (x_1 y_2 + y_1 x_2, 0) \\
                   & = \sigma(x_1, x_2) + \sigma(y_1, y_2) + (x_1 y_2 + y_1 x_2, 0)        \\
                   & \neq \sigma(x_1, x_2) + \sigma(y_1, y_2)
              \end{align*}

        \item 出发空间为 $ \mathbf{R}[x]_n $，到达空间为 $ \mathbf{R}[x]_{n - 1} $. $ \sigma $ 是线性映射.

              $ \forall p_1(x), p_2(x) \in \mathbf{R}[x]_n,\enspace k_1, k_2 \in \mathbf{R} $, 有
              \begin{align*}
                      & \sigma(k_1 p_1(x) + k_2 p_2(x))                     \\
                  ={} & (k_1 p_1 + k_2 p_2)(x + 1) - (k_1 p_1 + k_2 p_2)(x) \\
                  ={} & k_1(p_1(x + 1) - p_1(x)) + k_2(p_2(x + 1) - p_2(x)) \\
                  ={} & k_1 \sigma(p_1(x)) + k_2 \sigma(p_2(x))
              \end{align*}

        \item 出发空间为 $ \mathbf{R}[x] $，到达空间为 $ \mathbf{R}[x] $. $ \sigma $ 是线性映射.

              $ \forall p_1(x), p_2(x) \in \mathbf{R}[x],\enspace k_1, k_2 \in \mathbf{R} $, 有
              \begin{align*}
                      & \sigma(k_1 p_1(x) + k_2 p_2(x))         \\
                  ={} & (k_1 p_1 + k_2 p_2)(a)                  \\
                  ={} & k_1 p_1(a) + k_2 p_2(a)                 \\
                  ={} & k_1 \sigma(p_1(x)) + k_2 \sigma(p_2(x))
              \end{align*}

        \item 出发空间为 $ V $，到达空间为 $ V $. 当 $ \xi_0 = \vec{0} $ 时， $ \sigma(\xi) = 2 \xi $ 是线性映射.

              $ \forall \xi_1, \xi_2 \in V $, 有
              \begin{align*}
                  \sigma(\xi_1 + \xi_2) & = 2(\xi_1 + \xi_2)              \\
                                        & = 2 \xi_1 + 2 \xi_2             \\
                                        & = \sigma(\xi_1) + \sigma(\xi_2)
              \end{align*}

              当 $ \xi_0 \neq \vec{0} $ 时， $ \sigma $ 不是线性映射.

              $ \forall \xi_1, \xi_2 \in V $, 有
              \begin{align*}
                  \sigma(\xi_1 + \xi_2) & = 2(\xi_1 + \xi_2) + \xi_0              \\
                                        & = 2 \xi_1 + 2 \xi_2 + \xi_0             \\
                                        & = \sigma(\xi_1) + \sigma(\xi_2) - \xi_0 \\
                                        & \neq \sigma(\xi_1) + \sigma(\xi_2)
              \end{align*}
    \end{enumerate}
\end{solve}

\begin{exercise}[补充训练 L5-例-7]\label{prob:L5-例-7}
是否存在$\mathbf{R}^3$到$\mathbf{R}^2$的线性映射$\sigma$使得$\sigma(1,-1,1)^{\mathrm{T}}=(1,0)^{\mathrm{T}},\enspace\sigma(1,1,1)^{\mathrm{T}}=(0,1)^{\mathrm{T}}$？
\end{exercise}
\noindent 来源：\href{https://github.com/yhwu-is/Linear-Algebra-Left-Undone/blob/new/%E8%AE%B2%E4%B9%89/%E4%B8%93%E9%A2%98/5%20%E7%BA%BF%E6%80%A7%E6%98%A0%E5%B0%84.tex}{L5-例-7}。
\begin{solve}
事实上这里的定义完全不违反上述的必要条件，因此我们考虑证明这一线性映射存在. 根据原讲义“线性映射构造”处，我们考虑构造出$\sigma$在一组基（我们取自然基）下的像，然后就可以根据原讲义“线性映射构造”处 知道这一映射一定存在.

    实际上，根据题给条件有
    \begin{gather*}
        \sigma(1,-1,1)^{\mathrm{T}}
        =\sigma(e_1-e_2+e_3)
        =\sigma(e_1)-\sigma(e_2)+\sigma(e_3)
        =(1,0)^{\mathrm{T}} \\
        \sigma(1,1,1)^{\mathrm{T}}
        =\sigma(e_1+e_2+e_3)
        =\sigma(e_1)+\sigma(e_2)+\sigma(e_3)
        =(0,1)^{\mathrm{T}}
    \end{gather*}
    我们希望解出$\sigma(e_1),\sigma(e_2),\sigma(e_3)$，这样就可以直接根据原讲义“线性映射构造”处 构造出这一线性映射. 但这一方程组只有2个方程却有三个未知量. 事实上我们可以任意定义$\sigma(e_3)=(0,0)$，然后解方程组得到
    \[\sigma(e_1)=\dfrac{1}{2} (1,1)^{\mathrm{T}},\enspace \sigma(e_2)=\dfrac{1}{2} (-1,1)^{\mathrm{T}}\]
    又由$\sigma(e_3)=(0,0)^{\mathrm{T}}$，根据原讲义“线性映射构造”处，满足题目条件的线性映射存在.
\end{solve}

\begin{exercise}[补充训练 L5-例-10]\label{prob:L5-例-10}
指出下面各组内的两个线性空间是否同构，若同构可以进一步思考同构映射的构造：
    \begin{enumerate}
        \item 最高次不超过$n-1$的多项式构成的线性空间$\mathbf{R}[x]_n$与$\mathbf{R}^n$；

        \item 全体复数在实数域上的线性空间$\mathbf{C}(\mathbf{R})$与$\mathbf{R}^2$；

        \item 全体二元复向量$\mathbf{C}^2$在实数域上构成的线性空间$\mathbf{C}^2(\mathbf{R})$与$\mathbf{R}[x]_4$；

        \item 全体二元复向量$\mathbf{C}^2$在复数域上构成的线性空间$\mathbf{C}^2(\mathbf{C})$与$\mathcal{L}(\mathbf{R}^4,\mathbf{R})$.
    \end{enumerate}
\end{exercise}
\noindent 来源：\href{https://github.com/yhwu-is/Linear-Algebra-Left-Undone/blob/new/%E8%AE%B2%E4%B9%89/%E4%B8%93%E9%A2%98/5%20%E7%BA%BF%E6%80%A7%E6%98%A0%E5%B0%84.tex}{L5-例-10}。
\begin{solve}
\begin{enumerate}
        \item 同构，因为二者维数均为$n$. 同构映射非常简单，因为$\mathbf{R}[x]_n$在基$\{1,x,\ldots,x^{n-1}\}$下的坐标就在$\mathbf{R}^n$中，因此同构映射就是这一坐标映射：
              \[\sigma:\mathbf{R}[x]_n\to\mathbf{R}^n,\quad a_0+a_1x+\cdots+a_{n-1}x^{n-1}\mapsto(a_0,a_1,\ldots,a_{n-1}).\]
              我们很容易验证这一映射是线性双射，因此是同构映射.

        \item 同构，因为二者维数均为2（回忆原讲义“不同数域的维数”处）. 这一同构映射同上一小问理，$\mathbf{C}(\mathbf{R})$在基$\{1,i\}$下的坐标就在$\mathbf{R}^2$中，因此同构映射就是这一坐标映射：
              \[\sigma:\mathbf{C}(\mathbf{R})\to\mathbf{R}^2,\quad a+b\mathrm{i}\mapsto(a,b).\]
              事实上这就是将复数在二维平面中的表示，我们很容易验证这一映射是线性双射，因此是同构映射..

        \item 同构，因为二者维数均为4. 这一同构映射也非常简单，因为$\mathbf{C}^2(\mathbf{R})$在基
              \[\{(1,0),(\mathrm{i},0),(0,1),(0,\mathrm{i})\}\]
              下的坐标和$\mathbf{R}[x]_4$在基$\{1,x,x^2,x^3\}$下的坐标都在$\mathrm{R}^n$中，可以将它们一一对应，因此同构映射就是这一坐标映射：
              \[\sigma:\mathbf{C}^2(\mathbf{R})\to\mathbf{R}[x]_4,\quad (a+bi,c+di)\mapsto a+bx+cx^2+dx^3.\]
              我们很容易验证这一映射是线性双射，因此是同构映射.

        \item 不同构，因为$\mathbf{C}^2(\mathbf{C})$的维数为2，而$\mathcal{L}(\mathbf{R}^4,\mathbf{R})$的维数为4.
    \end{enumerate}
\end{solve}

\begin{exercise}[补充训练 L5-例-11]\label{prob:L5-例-11}
求积空间$\mathbf{R}[x]_3\times\mathbf{R}^2$的一组基.
\end{exercise}
\noindent 来源：\href{https://github.com/yhwu-is/Linear-Algebra-Left-Undone/blob/new/%E8%AE%B2%E4%B9%89/%E4%B8%93%E9%A2%98/5%20%E7%BA%BF%E6%80%A7%E6%98%A0%E5%B0%84.tex}{L5-例-11}。
\begin{solve}
我们知道$\mathbf{R}[x]_3$的一组基为$1,x,x^2$，而$\mathbf{R}^2$的一组基为$(1,0),(0,1)$. 很自然的想法是：我们可以先取$\mathbf{R}[x]_3$的一组基，$\mathbf{R}^2$的位置置零，然后反之取$\mathbf{R}^2$的一组基，$\mathbf{R}[x]_3$的位置置零，即$(1,(0,0)),\ (x,(0,0)),\ (x^2,(0,0)),\ (0,(1,0)),\ (0,(0,1))$. 我们很容易可以证明上述向量组满足基的两个条件：线性无关和张成空间.
\end{solve}

\begin{exercise}[补充训练 S-POLYMAP-1]\label{prob:S-POLYMAP-1}
在 $V=\mathbf R[x]_3$ 上定义 $T(f)(x)=f(x+1)+(-2)f'(x)$。证明 $T$ 线性，求它在基 $(1,x,x^2)$ 下的矩阵、核和像，并计算 $T(2x^{2}+1)$。
\end{exercise}
\noindent 编者补充：S-POLYMAP-1；相关主题见原讲义对应讲，题目数据与解答由本选编另行构造。
\begin{solve}
平移代入与求导均保持线性组合，故 $T(af+bg)=aT(f)+bT(g)$。分别计算基向量的像：
\[T(1)=1,\quad T(x)=x+(-1),\quad T(x^2)=x^2+(-2)x+1.\]
把这些像的坐标按列排列，得到 $[T]=\begin{pmatrix}1&-1&1\\0&1&-2\\0&0&1\end{pmatrix}$。这个上三角矩阵的对角元全为 $1$，所以可逆，$\ker T=\{0\}$、$\operatorname{im} T=V$，秩为 $3$。所求多项式的坐标为
\[[T]\begin{pmatrix}1\\0\\2\end{pmatrix}=\begin{pmatrix}3\\-4\\2\end{pmatrix},\]
因此 $T(2x^{2}+1)=2x^{2}-4x+3$。也可直接将 $x+1$ 代入原多项式后加上导数项核对。
\end{solve}

\begin{exercise}[补充训练 S-PROJ-1]\label{prob:S-PROJ-1}
设 $T(x)=Ax$，其中 $A=\begin{pmatrix}1&0&0\\0&1&0\\1&3&0\end{pmatrix}$。验证 $T^2=T$，求 $\operatorname{im} T$、$\ker T$ 的基，并把 $x=\begin{pmatrix}1\\1\\1\end{pmatrix}$ 唯一分解为像空间向量与核向量之和。判断该投影是否一定是正交投影。
\end{exercise}
\noindent 编者补充：S-PROJ-1；相关主题见原讲义对应讲，题目数据与解答由本选编另行构造。
\begin{solve}
取 $P=\begin{pmatrix}1&-3&0\\0&1&0\\1&0&1\end{pmatrix}$，有 $A=P\operatorname{diag}(1,1,0)P^{-1}$。因为对角矩阵的平方仍等于自身，$A^2=A$。像的一组基为 $P$ 的前两列，核的一组基为其第三列，故维数分别为 $2$ 与 $1$。对于题给向量，
\[u=Ax=\begin{pmatrix}1\\1\\4\end{pmatrix},\qquad v=(I-A)x=\begin{pmatrix}0\\0\\-3\end{pmatrix},\qquad x=u+v.\]
由 $A(I-A)=O$ 得 $v\in\ker T$。若一个向量同时在像与核中，则它既满足 $Az=z$ 又满足 $Az=0$，只能为零，故分解唯一。幂等性只保证沿核向像的投影；在标准实内积下，正交投影还须 $A^{\mathrm T}=A$，应另外检查这一条件。
\end{solve}

\section{补充训练：对偶空间}

本节包含 1 道训练题，建议在完成相关专题真题后选择性练习。

\begin{exercise}[补充训练 L6-例-1]\label{prob:L6-例-1}
设 $V=\mathbf R[x]_3$，定义线性泛函
\[f_1(g)=\int_0^1g(x)\,\mathrm dx,\quad f_2(g)=\int_0^2g(x)\,\mathrm dx,\quad f_3(g)=\int_0^{-1}g(x)\,\mathrm dx.\]
证明 $f_1,f_2,f_3$ 构成对偶空间 $V^*=\mathcal L(V,\mathbf R)$ 的一组基，并求 $V$ 的一组基 $g_1,g_2,g_3$，使 $f_i(g_j)=\delta_{ij}$。
\end{exercise}
\noindent 来源：\href{https://github.com/yhwu-is/Linear-Algebra-Left-Undone/blob/new/%E8%AE%B2%E4%B9%89/%E4%B8%93%E9%A2%98/6%20%E5%AF%B9%E5%81%B6%E7%A9%BA%E9%97%B4.tex}{L6-例-1}。
\begin{solve}
积分的线性性保证 $f_i\in V^*$。对 $g(x)=a+bx+cx^2$，有
\[\begin{pmatrix}f_1(g)\\f_2(g)\\f_3(g)\end{pmatrix}=M\begin{pmatrix}a\\b\\c\end{pmatrix},\qquad
M=\begin{pmatrix}1&1/2&1/3\\2&2&8/3\\-1&1/2&-1/3\end{pmatrix}.\]
直接计算 $\det M=-2\ne0$，故三个泛函线性无关；又 $\dim V^*=3$，它们构成一组基。所求 $g_j$ 的坐标是 $M^{-1}$ 的第 $j$ 列，由
\[M^{-1}=\begin{pmatrix}1&-1/6&-1/3\\1&0&1\\-3/2&1/2&-1/2\end{pmatrix}\]
得到
\[g_1=1+x-\frac32x^2,\quad g_2=-\frac16+\frac12x^2,\quad g_3=-\frac13+x-\frac12x^2.\]
等式 $MM^{-1}=I_3$ 同时验证 $f_i(g_j)=\delta_{ij}$；$M^{-1}$ 可逆也保证这三个多项式为基。
\end{solve}

\section{补充训练：矩阵运算与线性映射的矩阵表示}

本节包含 27 道训练题，建议在完成相关专题真题后选择性练习。

\begin{exercise}[补充训练 L7-A-1]\label{prob:L7-A-1}
证明：若$AB=BA$，$AC=CA$，则$A,B,C$为同阶方阵，且
        \[A(BC)=(BC)A,\enspace A(B+C)=(B+C)A.\]
\end{exercise}
\noindent 来源：\href{https://github.com/yhwu-is/Linear-Algebra-Left-Undone/blob/new/%E8%AE%B2%E4%B9%89/%E4%B8%93%E9%A2%98/7%20%E7%BA%BF%E6%80%A7%E6%98%A0%E5%B0%84%E7%9F%A9%E9%98%B5%E8%A1%A8%E7%A4%BA%E4%B8%8E%E7%9F%A9%E9%98%B5%E8%BF%90%E7%AE%97%E5%9F%BA%E7%A1%80.tex}{L7-A-1}。
\begin{solve}
由 $AB=BA$ 得 $A(BC)=(AB)C=(BA)C=B(AC)=BCA=(BC)A$，由$AC=CA$ 得 $A(B+C)=AB+AC=BA+CA=B+C$.
\end{solve}

\begin{exercise}[补充训练 L7-A-2]\label{prob:L7-A-2}
$A,B$都是$n$阶矩阵，求下列等式成立的充分条件：
        \begin{enumerate}
            \item $(A+B)^3=A^3+3A^2B+3AB^2+B^3$；

            \item $(A+B)(A-B)=A^2-B^2$.
        \end{enumerate}
\end{exercise}
\noindent 来源：\href{https://github.com/yhwu-is/Linear-Algebra-Left-Undone/blob/new/%E8%AE%B2%E4%B9%89/%E4%B8%93%E9%A2%98/7%20%E7%BA%BF%E6%80%A7%E6%98%A0%E5%B0%84%E7%9F%A9%E9%98%B5%E8%A1%A8%E7%A4%BA%E4%B8%8E%E7%9F%A9%E9%98%B5%E8%BF%90%E7%AE%97%E5%9F%BA%E7%A1%80.tex}{L7-A-2}。
\begin{solve}
\begin{enumerate}
                \item $(A+B)^3=A^3+3A^2B+3AB^2+B^3$的充分条件是$AB=BA$；

                \item $(A+B)(A-B)=A^2-B^2$的充分条件是$AB=BA$.
            \end{enumerate}
\end{solve}

\begin{exercise}[补充训练 L7-A-3]\label{prob:L7-A-3}
设$A$是$n$阶方阵且$A^n=O$，证明：
        \[(E_n-A)(E_n+A+A^2+\cdots+A^{n-1})=E_n.\]
\end{exercise}
\noindent 来源：\href{https://github.com/yhwu-is/Linear-Algebra-Left-Undone/blob/new/%E8%AE%B2%E4%B9%89/%E4%B8%93%E9%A2%98/7%20%E7%BA%BF%E6%80%A7%E6%98%A0%E5%B0%84%E7%9F%A9%E9%98%B5%E8%A1%A8%E7%A4%BA%E4%B8%8E%E7%9F%A9%E9%98%B5%E8%BF%90%E7%AE%97%E5%9F%BA%E7%A1%80.tex}{L7-A-3}。
\begin{solve}
由$A^n=O$得$A^n-E_n^n=-E_n$，即$(A-E_n)(A^{n-1}+A^{n-2}+\cdots+A+E_n)=-E_n$，即$(E_n-A)(E_n+A+A^2+\cdots+A^{n-1})=E_n$.
\end{solve}

\begin{exercise}[补充训练 L7-A-4]\label{prob:L7-A-4}
证明：若线性映射$\sigma \in \mathcal{L}(V_1,V_2)$可逆，则其逆映射唯一.
\end{exercise}
\noindent 来源：\href{https://github.com/yhwu-is/Linear-Algebra-Left-Undone/blob/new/%E8%AE%B2%E4%B9%89/%E4%B8%93%E9%A2%98/7%20%E7%BA%BF%E6%80%A7%E6%98%A0%E5%B0%84%E7%9F%A9%E9%98%B5%E8%A1%A8%E7%A4%BA%E4%B8%8E%E7%9F%A9%E9%98%B5%E8%BF%90%E7%AE%97%E5%9F%BA%E7%A1%80.tex}{L7-A-4}。
\begin{solve}
设$\sigma$的逆映射为$\tau_1,\tau_2$，则$\sigma\tau_1=\sigma\tau_2=E$，则$\sigma(\tau_1-\tau_2)=0$.

            因为 $\sigma$ 可逆，所以对于任意 $x \neq 0$，有 $\sigma(x) \neq 0$.

            而 $\forall x, \sigma((\tau_1-\tau_2)(x))=0$，故 $\tau_1-\tau_2=0, \tau_1=\tau_2$.
\end{solve}

\begin{exercise}[补充训练 L7-A-5]\label{prob:L7-A-5}
证明：有一行元素或一列元素全为0的$n$阶方阵必定不可逆.
\end{exercise}
\noindent 来源：\href{https://github.com/yhwu-is/Linear-Algebra-Left-Undone/blob/new/%E8%AE%B2%E4%B9%89/%E4%B8%93%E9%A2%98/7%20%E7%BA%BF%E6%80%A7%E6%98%A0%E5%B0%84%E7%9F%A9%E9%98%B5%E8%A1%A8%E7%A4%BA%E4%B8%8E%E7%9F%A9%E9%98%B5%E8%BF%90%E7%AE%97%E5%9F%BA%E7%A1%80.tex}{L7-A-5}。
\begin{solve}
设$A$的第$i$行全为0，则$A$的第$i$行与$A$的任意行线性相关，记 $A$ 为 $n$ 阶矩阵，则$r(A) < n$.

            故对任意 $n$ 阶矩阵 $B$，有 $r(AB) \leqslant r(A) < n = r(E_n)$，故 $A$ 不可逆.
\end{solve}

\begin{exercise}[补充训练 L7-A-6]\label{prob:L7-A-6}
设$\alpha,\beta$为三维列向量，且$\alpha\beta^\mathrm{T}=\begin{pmatrix}
                -1 & 2  & 1  \\
                1  & -2 & -1 \\
                2  & -4 & -2
            \end{pmatrix}$，求$\alpha^\mathrm{T}\beta$.
\end{exercise}
\noindent 来源：\href{https://github.com/yhwu-is/Linear-Algebra-Left-Undone/blob/new/%E8%AE%B2%E4%B9%89/%E4%B8%93%E9%A2%98/7%20%E7%BA%BF%E6%80%A7%E6%98%A0%E5%B0%84%E7%9F%A9%E9%98%B5%E8%A1%A8%E7%A4%BA%E4%B8%8E%E7%9F%A9%E9%98%B5%E8%BF%90%E7%AE%97%E5%9F%BA%E7%A1%80.tex}{L7-A-6}。
\begin{solve}
记 $\alpha = (a_1,a_2,a_3)$, $\beta = (b_1,b_2,b_3)$，则 
\[\alpha\beta^\mathrm{T} = \begin{pmatrix}
                    a_1b_1 & a_1b_2 & a_1b_3 \\
                    a_2b_1 & a_2b_2 & a_2b_3 \\
                    a_3b_1 & a_3b_2 & a_3b_3
                \end{pmatrix}\]
,
            则有：$\text{tr}(\alpha\beta^\mathrm{T}) = a_1b_1+a_2b_2+a_3b_3 = \alpha^\mathrm{T}\beta$ ，因此 $\alpha^\mathrm{T}\beta = -1-2-2 = -5$.
\end{solve}

\begin{exercise}[补充训练 L7-B-1]\label{prob:L7-B-1}
已知$f_1=1-x,f_2=1+x^2,f_3=x+2x^2$是$\mathbf{R}[x]_3$中三个元素，$\sigma$是$\mathbf{R}[x]_3$上的线性变换且满足$\sigma(f_1)=2+x^2,\sigma(f_2)=x,\sigma(f_3)=1+x+x^2$.
        \begin{enumerate}
            \item 证明：$f_1,f_2,f_3$构成$\mathbf{R}[x]_3$的一组基；

            \item 求$\sigma$在基$f_1,f_2,f_3$下的矩阵；

            \item 设$f=1+2x+3x^2$，求$\sigma(f)$.
        \end{enumerate}
\end{exercise}
\noindent 来源：\href{https://github.com/yhwu-is/Linear-Algebra-Left-Undone/blob/new/%E8%AE%B2%E4%B9%89/%E4%B8%93%E9%A2%98/7%20%E7%BA%BF%E6%80%A7%E6%98%A0%E5%B0%84%E7%9F%A9%E9%98%B5%E8%A1%A8%E7%A4%BA%E4%B8%8E%E7%9F%A9%E9%98%B5%E8%BF%90%E7%AE%97%E5%9F%BA%E7%A1%80.tex}{L7-B-1}。
\begin{solve}
\begin{enumerate}
                \item 只需要证明$f_1,f_2,f_3$是线性无关的即可，事实上，这也十分显然(考虑其在$1,x,x^2$下的坐标表示线性无关),想必读者可以自行证明

                \item 方法一：设 $ (\sigma(f_1), \sigma(f_2), \sigma(f_3)) = (f_1, f_2, f_3) A $，则由
                      \[ (\sigma(f_1), \sigma(f_2), \sigma(f_3)) = (1, x, x^2) \begin{pmatrix}
                              2 & 0 & 1 \\
                              0 & 1 & 1 \\
                              1 & 0 & 1
                          \end{pmatrix} = (1, x, x^2) \begin{pmatrix}
                              1  & 1 & 0 \\
                              -1 & 0 & 1 \\
                              0  & 1 & 2
                          \end{pmatrix} A \]
                      可得
                      \begin{align*}
                          A & = \begin{pmatrix}
                                    1  & 1 & 0 \\
                                    -1 & 0 & 1 \\
                                    0  & 1 & 2
                                \end{pmatrix}^{-1}
                          \begin{pmatrix}
                              2 & 0 & 1 \\
                              0 & 1 & 1 \\
                              1 & 0 & 1
                          \end{pmatrix}           \\
                            & = \begin{pmatrix}
                                    -1 & -2 & 1  \\
                                    2  & 2  & -1 \\
                                    -1 & -1 & 1
                                \end{pmatrix}
                          \begin{pmatrix}
                              2 & 0 & 1 \\
                              0 & 1 & 1 \\
                              1 & 0 & 1
                          \end{pmatrix}
                          = \begin{pmatrix}
                                -1 & -2 & -2 \\
                                3  & 2  & 3  \\
                                -1 & -1 & -1
                            \end{pmatrix}
                      \end{align*}

                      方法二：通过待定系数法解方程组
                      \[ \begin{cases}
                              \sigma(f_1) = -f_1 + 3 f_2 - f_3   \\
                              \sigma(f_2) = -2 f_1 + 2 f_2 - f_3 \\
                              \sigma(f_3) = -2 f_1 + 3 f_2 - f_3
                          \end{cases} \]
                      解得
                      \[ (\sigma(f_1), \sigma(f_2), \sigma(f_3)) = (f_1, f_2, f_3) \begin{pmatrix}
                              -1 & -2 & -2 \\
                              3  & 2  & 3  \\
                              -1 & -1 & -1
                          \end{pmatrix} \]

                \item \begin{align*}
                          \sigma(f) & = \sigma\left((1, x, x^2)
                          \begin{pmatrix} 1 \\ 2 \\ 3 \end{pmatrix}\right)= \sigma\left((f_1, f_2, f_3)
                          \begin{pmatrix}
                              1  & 1 & 0 \\
                              -1 & 0 & 1 \\
                              0  & 1 & 2
                          \end{pmatrix}^{-1} \begin{pmatrix} 1 \\ 2 \\ 3 \end{pmatrix}\right) \\
                                    & = (\sigma(f_1), \sigma(f_2), \sigma(f_3))
                          \begin{pmatrix}
                              1  & 1 & 0 \\
                              -1 & 0 & 1 \\
                              0  & 1 & 2
                          \end{pmatrix}^{-1} \begin{pmatrix} 1 \\ 2 \\ 3 \end{pmatrix}        \\
                                    & = (1, x, x^2)
                          \begin{pmatrix}
                              2 & 0 & 1 \\
                              0 & 1 & 1 \\
                              1 & 0 & 1
                          \end{pmatrix}
                          \begin{pmatrix}
                              1  & 1 & 0 \\
                              -1 & 0 & 1 \\
                              0  & 1 & 2
                          \end{pmatrix}^{-1}
                          \begin{pmatrix} 1 \\ 2 \\ 3 \end{pmatrix}                           \\
                                    & = (1, x, x^2)
                          \begin{pmatrix} -4 \\ 3 \\ 2 \end{pmatrix}                          \\
                                    & = 2x^2 + 3x - 4
                      \end{align*}

                      或也可采用待定系数法求出 $ f = -2 f_1 + 3 f_2 $，所以 $ \sigma(f) = -2 \sigma(f_1) + 3 \sigma(f_2) = 2x^2 + 3x - 4 $.
            \end{enumerate}
\end{solve}

\begin{exercise}[补充训练 L7-B-3]\label{prob:L7-B-3}
设 $\mathbf{R}[x]_4$ 是数域 $\mathbf{R}$ 上次数小于 4 的多项式所构成的线性空间（约定零多项式次数为 $-\infty$）. $\mathbf{M}_2(\mathbf{R})$ 是 $\mathbf{R}$ 上 2 阶方阵所构成的线性空间. 定义 $T \colon \mathbf{R}[x]_4 \to \mathbf{M}_2(\mathbf{R})$ 如下：对 $f(x) \in \mathbf{R}[x]_4$，
        \[T(f(x))=\begin{pmatrix}f(0) & f(1) \\ f(-1) & f(0)\end{pmatrix}\]
        \begin{enumerate}
            \item 求出 $T$ 的核空间 $N(T)$ 和像空间 $R(T)$；

            \item 求$T$在$\mathbf{R}[x]_4$和$\mathbf{M}_2(\mathbf{R})$的基下的矩阵表示.
        \end{enumerate}
\end{exercise}
\noindent 来源：\href{https://github.com/yhwu-is/Linear-Algebra-Left-Undone/blob/new/%E8%AE%B2%E4%B9%89/%E4%B8%93%E9%A2%98/7%20%E7%BA%BF%E6%80%A7%E6%98%A0%E5%B0%84%E7%9F%A9%E9%98%B5%E8%A1%A8%E7%A4%BA%E4%B8%8E%E7%9F%A9%E9%98%B5%E8%BF%90%E7%AE%97%E5%9F%BA%E7%A1%80.tex}{L7-B-3}。
\begin{solve}
\begin{enumerate}
                \item \label{item:7:B:5:1}
                      求核空间，即求使得 $ T(f(x)) $ 为零矩阵的 $ f(x) $ 构成的空间. 设 $ f(x) = ax^3 + bx^2 + cx + d $，则有
                      \[ \begin{cases}
                              f(0) = d = 0             \\
                              f(1) = a + b + c + d = 0 \\
                              f(-1) = -a + b -c + d = 0
                          \end{cases} \]
                      解方程，令 $ a = t $ 有
                      \[ \begin{cases}
                              a = t  \\
                              b = 0  \\
                              c = -t \\
                              d = 0
                          \end{cases} \implies f(x) = t(x^3 - x) \]
                      故 $ N(T) = \operatorname{span}(x^3 - x) $.

                      求像空间，取 $ \mathbf{R}[x]_4 $ 的自然基 $ 1, x, x^2, x^3 $.
                      \[ \begin{matrix}
                              T(1) = \begin{pmatrix} 1 & 1 \\ 1 & 1 \end{pmatrix}   & T(x) = \begin{pmatrix} 0 & 1 \\ -1 & 0 \end{pmatrix}   \\
                              T(x^2) = \begin{pmatrix} 0 & 1 \\ 1 & 0 \end{pmatrix} & T(x^3) = \begin{pmatrix} 0 & 1 \\ -1 & 0 \end{pmatrix}
                          \end{matrix} \]
                      求它们的极大线性无关组. 我们发现 $ T(x) = T(x^3) $，故先舍弃 $ T(x^3) $，然后令
                      \begin{align*}
                                     & k_1 T(1) + k_2 T(x) + k_3 T(x^2) = 0                                             \\
                          \implies{} & \begin{pmatrix} k_1 & k_1 + k_2 + k_3 \\ k_1 - k_2 + k_3 & k_1 \end{pmatrix} = 0 \\
                          \implies{} & \begin{cases}
                                           k_1 = 0             \\
                                           k_1 + k_2 + k_3 = 0 \\
                                           k_1 - k_2 + k_3 = 0 \\
                                       \end{cases}                                                              \\
                          \implies{} & k_1 = k_2 = k_3 = 0
                      \end{align*}
                      故 $ T(1), T(x), T(x^2) $ 线性无关. 故 $ R(T) = \operatorname{span}\left(\begin{pmatrix} 1 & 1 \\ 1 & 1 \end{pmatrix}, \begin{pmatrix} 0 & 1 \\ -1 & 0 \end{pmatrix}, \begin{pmatrix} 0 & 1 \\ 1 & 0 \end{pmatrix}\right) $.

                \item 我们只需要将$2 \times 2$矩阵的第二列拼到第一列下方就可以变为我们熟悉的列向量形式，所以其矩阵表示也是简单的
                      \[
                        T(1,x,x^2,x^3)=(e_1,e_2,e_3,e_4)\begin{pmatrix}
                            1 & 0 & 0 & 0 \\
                            1 & -1 & 1 & -1 \\
                            1 & 1 & 1 & 1 \\
                            1 & 0 & 0 & 0
                        \end{pmatrix}
                      \]

                      其中$e_1=\begin{pmatrix}1&0\\0&0\end{pmatrix}$,$e_2=\begin{pmatrix}0&0\\1&0\end{pmatrix}$,
                      $e_3=\begin{pmatrix}0&1\\0&0\end{pmatrix}$,$e_4=\begin{pmatrix}0&0\\0&1\end{pmatrix}$.
            \end{enumerate}
\end{solve}

\begin{exercise}[补充训练 L7-B-4]\label{prob:L7-B-4}
设$A=\begin{pmatrix}
                1 & -1 & -1 \\ -1 & 1 & 1 \\ 0 & -4 & 2
            \end{pmatrix},\enspace\xi_1=(-1,1,-2)^\mathrm{T}$.
        \begin{enumerate}
            \item 求满足$A\xi_2=\xi_1$及$A^2\xi_3=\xi_1$的所有$\xi_2,\xi_3$；

            \item 证明：$\xi_1,\xi_2,\xi_3$线性无关.
        \end{enumerate}
\end{exercise}
\noindent 来源：\href{https://github.com/yhwu-is/Linear-Algebra-Left-Undone/blob/new/%E8%AE%B2%E4%B9%89/%E4%B8%93%E9%A2%98/7%20%E7%BA%BF%E6%80%A7%E6%98%A0%E5%B0%84%E7%9F%A9%E9%98%B5%E8%A1%A8%E7%A4%BA%E4%B8%8E%E7%9F%A9%E9%98%B5%E8%BF%90%E7%AE%97%E5%9F%BA%E7%A1%80.tex}{L7-B-4}。
\begin{solve}
\begin{enumerate}
                \item 对非齐次线性方程组$AX=\xi_1$，\\
                      
\[\bar{A}=\begin{pmatrix}
                              1  & -1 & -1 & -1 \\
                              -1 & 1  & 1  & 1  \\
                              0  & -4 & -2 & -2
                          \end{pmatrix}
                          \rightarrow\begin{pmatrix}
                              1 & -1 & -1       & -1       \\
                              0 & 1  & \frac 12 & \frac 12 \\
                              0 & 0  & 0        & 0
                          \end{pmatrix}
                          \rightarrow\begin{pmatrix}
                              1 & 0 & -\frac 12 & -\frac 12 \\
                              0 & 1 & \frac 12  & \frac 12  \\
                              0 & 0 & 0         & 0
                          \end{pmatrix}\]
，则\\
                      
\[\xi_2=C_1\begin{pmatrix}
                              \frac 12  \\
                              -\frac 12 \\
                              1
                          \end{pmatrix} + \begin{pmatrix}
                              -\frac 12 \\
                              \frac 12  \\
                              0
                          \end{pmatrix}=\dfrac 12\begin{pmatrix}
                              C_1 - 1  \\
                              -C_1 + 1 \\
                              2C_1
                          \end{pmatrix}\]
（其中$C_1$为任意常数）.\\
                      $A^2=\begin{pmatrix}
                              2  & 2  & 0 \\
                              -2 & -2 & 0 \\
                              4  & 4  & 0
                          \end{pmatrix}$，对齐次线性方程组$A^2X=\xi_1$，\\
                      
\[\bar{B}=\begin{pmatrix}
                              A^2 & \xi_1
                          \end{pmatrix}=\begin{pmatrix}
                              2  & 2  & 0 & 1  \\
                              -2 & -2 & 0 & 1  \\
                              4  & 4  & 0 & -2
                          \end{pmatrix}
                          \rightarrow\begin{pmatrix}
                              1 & 1 & 0 & -\frac 12 \\
                              0 & 0 & 0 & 0         \\
                              0 & 0 & 0 & 0
                          \end{pmatrix}\]
，\\
                      则$A^2X=\xi_1$的通解\\
                      
\[\xi_3=C_2\begin{pmatrix}
                              -1 \\
                              1  \\
                              0
                          \end{pmatrix}+C_3\begin{pmatrix}
                              0 \\
                              0 \\
                              1
                          \end{pmatrix}+\begin{pmatrix}
                              -\frac 12 \\
                              0         \\
                              0
                          \end{pmatrix}=\begin{pmatrix}
                              -C_2 - \frac 12 \\
                              C_2             \\
                              C_3
                          \end{pmatrix}\]
（其中$C_2, C_3$为任意常数）.
                \item
                      因为$|\xi_1,\xi_2,\xi_3|=\dfrac 12\begin{vmatrix}
                              -1 & C_1 - 1  & -C_2 - \frac 12 \\
                              1  & -C_1 + 1 & C_2             \\
                              -2 & 2C_1     & C_3
                          \end{vmatrix}=-\dfrac 12\neq 0$，\\
                      所以$\xi_1,\xi_2,\xi_3$线性无关.
            \end{enumerate}
\end{solve}

\begin{exercise}[补充训练 L7-B-5]\label{prob:L7-B-5}
已知$V$为$n$维线性空间，$\sigma\in \mathcal{L}(V,V)$，且$\ker\sigma=\operatorname{im} \sigma$，证明：
        \begin{enumerate}
            \item $n$为偶数；

            \item 存在$V$的一组基$\alpha_1,\ldots,\alpha_n$使得
                  \[\sigma(\alpha_1,\ldots,\alpha_n)=(\alpha_1,\ldots,\alpha_n)\begin{pmatrix}
                          O & E_{\frac{n}{2}} \\ O & O
                      \end{pmatrix}.\]
        \end{enumerate}
\end{exercise}
\noindent 来源：\href{https://github.com/yhwu-is/Linear-Algebra-Left-Undone/blob/new/%E8%AE%B2%E4%B9%89/%E4%B8%93%E9%A2%98/7%20%E7%BA%BF%E6%80%A7%E6%98%A0%E5%B0%84%E7%9F%A9%E9%98%B5%E8%A1%A8%E7%A4%BA%E4%B8%8E%E7%9F%A9%E9%98%B5%E8%BF%90%E7%AE%97%E5%9F%BA%E7%A1%80.tex}{L7-B-5}。
\begin{solve}
证明：\begin{enumerate}
                \item 由于 $ \ker \sigma = \operatorname{im} \sigma $，由 $ \dim \operatorname{im} \sigma + \dim \ker \sigma = \dim V $ 可得.

                \item 设 $ \beta_1, \ldots, \beta_n $ 为 $ V $ 的一组基，则
                      \[ \operatorname{im} \sigma = \operatorname{span}(\sigma(\beta_1), \ldots, \sigma(\beta_n)) = \ker \sigma \]
                      设 $ \sigma(\beta_1), \ldots, \sigma(\beta_{\frac{n}{2}}) $ 为 $ \operatorname{im} \sigma $ 的基，则可以证明
                      \[ \sigma(\beta_1), \ldots, \sigma(\beta_{\frac{n}{2}}), \beta_1, \ldots, \beta_{\frac{n}{2}} \]
                      线性无关，且 $ \sigma $ 在此基下的矩阵即为所求的形式.
            \end{enumerate}
\end{solve}

\begin{exercise}[补充训练 L7-B-6]\label{prob:L7-B-6}
若$f(x)$是$x$的实系数$m$次多项式：
        \[f(x)=a_mx^m+a_{m-1}x^{m-1}+\cdots+a_1x+a_0\]
        则有矩阵多项式：
        \[f(A)=a_mA^m+a_{m-1}A^{m-1}+\cdots+a_1A+a_0E\]
        其中 $A^0=E$.
        \begin{enumerate}
            \item 若$A$为对角矩阵$B=\begin{pmatrix}
                          \lambda_1 & 0 \\ 0 & \lambda_2
                      \end{pmatrix}$，证明：$f(A)=\begin{pmatrix}
                          f(\lambda_1) & 0 \\ 0 & f(\lambda_2)
                      \end{pmatrix}$；

            \item 若$A=PBP^{-1}$，证明：$f(A)=Pf(B)P^{-1}$.
        \end{enumerate}
\end{exercise}
\noindent 来源：\href{https://github.com/yhwu-is/Linear-Algebra-Left-Undone/blob/new/%E8%AE%B2%E4%B9%89/%E4%B8%93%E9%A2%98/7%20%E7%BA%BF%E6%80%A7%E6%98%A0%E5%B0%84%E7%9F%A9%E9%98%B5%E8%A1%A8%E7%A4%BA%E4%B8%8E%E7%9F%A9%E9%98%B5%E8%BF%90%E7%AE%97%E5%9F%BA%E7%A1%80.tex}{L7-B-6}。
\begin{solve}
\begin{enumerate}
                \item 因为 $A^k=\begin{pmatrix}\lambda_1^k & 0 \\ 0 & \lambda_2^k\end{pmatrix}$，所以 $f(A)=a_mA^m+a_{m-1}A^{m-1}+\cdots+a_1A+a_0E = \begin{pmatrix}f(\lambda_1) & 0 \\ 0 & f(\lambda_2)\end{pmatrix}$.
                \item $A=PBP^{-1}$，则 $A^2=(PBP^{-1})(PBP^{-1})=PB^2P^{-1}$. 由归纳法得 $A^k=PB^kP^{-1}$，于是
                      \begin{align*}
                          f(A) & = a_mPB^mP^{-1}+a_{m-1}PB^{m-1}P^{-1}+\cdots+a_0                         \\
                               & =P\begin{pmatrix}f(\lambda_1) & 0 \\ 0 & f(\lambda_2)\end{pmatrix}P^{-1} \\
                               & =Pf(B)P^{-1}
                      \end{align*}
            \end{enumerate}
\end{solve}

\begin{exercise}[补充训练 L7-B-14]\label{prob:L7-B-14}
设$\mathbf{F}$为数域，$V_1=\{A\in\mathbf{F}^{n\times n} \mid A^\mathrm{T}=A\},\enspace
            V_2=\{A\in\mathbf{F}^{n\times n} \mid A^\mathrm{T}=-A\},\enspace V_3=\{A\in\mathbf{F}^{n\times n} \mid A\text{~为上三角矩阵}\}$.
        \begin{enumerate}
            \item 证明：$V_1,V_2,V_3$都是$\mathbf{F}^{n\times n}$的子空间；

            \item 证明：$\mathbf{F}^{n\times n}=V_1+V_3$但不为直和，$\mathbf{F}^{n\times n}=V_2\oplus V_3$.
        \end{enumerate}
\end{exercise}
\noindent 来源：\href{https://github.com/yhwu-is/Linear-Algebra-Left-Undone/blob/new/%E8%AE%B2%E4%B9%89/%E4%B8%93%E9%A2%98/7%20%E7%BA%BF%E6%80%A7%E6%98%A0%E5%B0%84%E7%9F%A9%E9%98%B5%E8%A1%A8%E7%A4%BA%E4%B8%8E%E7%9F%A9%E9%98%B5%E8%BF%90%E7%AE%97%E5%9F%BA%E7%A1%80.tex}{L7-B-14}。
\begin{solve}
三个集合均含零矩阵，而且它们的定义条件都对线性组合封闭：转置保持线性组合，上三角矩阵相加或数乘仍是上三角矩阵，所以它们是子空间。
对任意 $A=(a_{ij})$，定义对称矩阵 $S$，使其对角线及下三角部分与 $A$ 相同，即 $s_{ij}=a_{ij}$（$i\ge j$），$s_{ij}=a_{ji}$（$i<j$）。于是 $A-S$ 为严格上三角矩阵，说明 $V_1+V_3$ 是全空间。但 $I\in V_1\cap V_3$，故该和不是直和。
另定义反对称矩阵 $K$：$k_{ij}=a_{ij}$（$i>j$）、$k_{ij}=-a_{ji}$（$i<j$）、$k_{ii}=0$。那么 $A-K$ 是上三角矩阵，说明 $V_2+V_3$ 为全空间。一个既上三角又反对称的矩阵，其严格下三角元、严格上三角元均为零；数域上 $2\ne0$，对角元也为零。因此交为 $\{O\}$，从而 $\mathbf F^{n\times n}=V_2\oplus V_3$。
\end{solve}

\begin{exercise}[补充训练 L7-B-15]\label{prob:L7-B-15}
设\[W_1=\left\{\begin{pmatrix}
            x & -x \\ y & z
        \end{pmatrix} \;\middle|\; x,y,z\in \mathbf{F} \right\},W_2=\left\{\begin{pmatrix}
            a & b \\ -a & c
        \end{pmatrix} \;\middle|\; a,b,c\in \mathbf{F} \right\}.\]

        \begin{enumerate}
            \item 证明：$W_1,W_2$是$\mathbf{M}_2(\mathbf{F})$的子空间，并求$\dim W_1,\dim W_2,\dim(W_1+W_2),\dim(W_1\cap W_2)$;

            \item 求$W_1\cap W_2$的一组基，并求$A=\begin{pmatrix}
                        3 & -3 \\ -3 & 1
                    \end{pmatrix}$关于这组基的坐标.
        \end{enumerate}
\end{exercise}
\noindent 来源：\href{https://github.com/yhwu-is/Linear-Algebra-Left-Undone/blob/new/%E8%AE%B2%E4%B9%89/%E4%B8%93%E9%A2%98/7%20%E7%BA%BF%E6%80%A7%E6%98%A0%E5%B0%84%E7%9F%A9%E9%98%B5%E8%A1%A8%E7%A4%BA%E4%B8%8E%E7%9F%A9%E9%98%B5%E8%BF%90%E7%AE%97%E5%9F%BA%E7%A1%80.tex}{L7-B-15}。
\begin{solve}
\begin{enumerate}
                \item 只需证明 $W_1$ 封闭即可. 对于 $A=\begin{pmatrix}x & -x \\ y & z\end{pmatrix},A'=\begin{pmatrix} x' & -x' \\ y' & z' \end{pmatrix}$ 有 
\[\lambda A+\mu A'=\begin{pmatrix}
                            \lambda x+\mu x' & -(\lambda x+\mu x') \\
                            \lambda y+\mu y  & \lambda z+\mu z
                        \end{pmatrix}\]
，则$W_1$封闭，$W_1$是$\mathbf{M}_2(\mathbf{F})$的子空间.	$W_2$同理可证.

                \item $W_1$的基：$B_1=\begin{pmatrix}1&-1\\0&0\end{pmatrix},B_2=\begin{pmatrix}0&0\\1&0\end{pmatrix},B_3=\begin{pmatrix}0&0\\0&1\end{pmatrix}$.

                    $W_2$的基：$B_1=\begin{pmatrix}1&0\\-1&0\end{pmatrix},B_2=\begin{pmatrix}0&1\\0&0\end{pmatrix},B_3=\begin{pmatrix}0&0\\0&1\end{pmatrix}$.

                    则$\dim W_1=3,\dim W_2=3$. 要求$\dim  (W_1+W_2)$，只需求$B_1,B_2,\ldots ,B_6$的极大无关组即可. 可知 $B_1,B_2,B_3,B_4$是极大线性无关组.	$\dim(W_1+W_2)=4$. 根据维数公式$\dim (W_1\cap W_2)=\dim W_1+\dim W_2-\dim (W_1+W_2)=2$.

                \item $W_1\cap W_2=\left\{\begin{pmatrix}x&-x\\-x&y\end{pmatrix} \mid x,y\in \mathbf{F}\right\}$. 则一组基为$E_1=\begin{pmatrix}1&-1\\-1&0\end{pmatrix},E_2=\begin{pmatrix}0&0\\0&1\end{pmatrix}$. $A$ 的坐标即为 $(3,1)$.
            \end{enumerate}
\end{solve}

\begin{exercise}[补充训练 L7-C-2]\label{prob:L7-C-2}
设 $A,B,C$ 为二阶复方阵，且 $A,B,C$ 在 $\mathbf{M}_2(\mathbf{C})$ 中线性无关. 证明：存在$z_1,z_2,z_3 \in \mathbf{C}$使得 $z_1A+z_2B+z_3C$ 为可逆矩阵.
\end{exercise}
\noindent 来源：\href{https://github.com/yhwu-is/Linear-Algebra-Left-Undone/blob/new/%E8%AE%B2%E4%B9%89/%E4%B8%93%E9%A2%98/7%20%E7%BA%BF%E6%80%A7%E6%98%A0%E5%B0%84%E7%9F%A9%E9%98%B5%E8%A1%A8%E7%A4%BA%E4%B8%8E%E7%9F%A9%E9%98%B5%E8%BF%90%E7%AE%97%E5%9F%BA%E7%A1%80.tex}{L7-C-2}。
\begin{solve}
首先设
            \[A=\begin{pmatrix}a_1 & a_2 \\ a_3 & a_4\end{pmatrix},B=\begin{pmatrix}b_1 & b_2 \\ b_3 & b_4\end{pmatrix},C=\begin{pmatrix}c_1 & c_2 \\ c_3 & c_4\end{pmatrix}.\]
            由于 $A,B,C$ 在 $\mathbf{M}_2(\mathbf{C})$ 中线性无关，所以将 $A,B,C$ 的元素排为一列，可知矩阵
            \[\begin{pmatrix}a_1 & b_1 & c_1 \\ a_2 & b_2 & c_2 \\ a_3 & b_3 & c_3 \\ a_4 & b_4 & c_4\end{pmatrix}\]
            的秩为 3，这里不妨设前三个行向量线性无关，即有
            \[D=\begin{pmatrix}a_1 & b_1 & c_1 \\ a_2 & b_2 & c_2 \\ a_3 & b_3 & c_3\end{pmatrix}\]
            为可逆矩阵.

            另外，注意到对任意的 $x_1,x_2,x_3$，有
            \[x_1A+x_2B+x_3C=\begin{pmatrix}a_1x_1+b_1x_2+c_1x_3 & a_2x_1+b_2x_2+c_2x_3 \\ a_3x_1+b_3x_2+c_3x_3 & a_4x_1+b_4x_2+c_4x_3\end{pmatrix}.\]
            现在考虑方程组
            \[\begin{cases}a_1x_1+b_1x_2+c_1x_3 = 0 \\ a_2x_1+b_2x_2+c_2x_3 = 1 \\ a_3x_1+b_3x_2+c_3x_3 = 1\end{cases}\]
            其系数矩阵为 $D$，这是一个可逆矩阵，所以上述方程存在唯一解，不妨记为 $(x_1',x_2',x_3')$，此时就有 $\lvert x_1'A+x_2'B+x_3'C \rvert = \begin{vmatrix}0 & 1 \\ 1 & \ast\end{vmatrix} = -1$.

            所以 $x_1'A+x_2'B+x_3'C$ 为可逆矩阵.
\end{solve}

\begin{exercise}[补充训练 L7-例-1]\label{prob:L7-例-1}
已知$\sigma \in \mathcal{L}(\mathbf{R}^3,\mathbf{R}^3)$且$\sigma(x_1,x_2,x_3)^{\mathrm{T}}=(x_1+x_2,x_1-x_3, x_2)^{\mathrm{T}}$
    \begin{enumerate}
        \item 求$\sigma$的像空间和核空间；

        \item 求$\sigma$关于$\mathbf{R}^3$自然基的矩阵.
    \end{enumerate}
\end{exercise}
\noindent 来源：\href{https://github.com/yhwu-is/Linear-Algebra-Left-Undone/blob/new/%E8%AE%B2%E4%B9%89/%E4%B8%93%E9%A2%98/7%20%E7%BA%BF%E6%80%A7%E6%98%A0%E5%B0%84%E7%9F%A9%E9%98%B5%E8%A1%A8%E7%A4%BA%E4%B8%8E%E7%9F%A9%E9%98%B5%E8%BF%90%E7%AE%97%E5%9F%BA%E7%A1%80.tex}{L7-例-1}。
\begin{solve}
\begin{enumerate}
        \item 求像空间和核空间的方法我们在之前已经介绍过，为了计算方便取$\mathbf{R}^3$的自然基$e_1,e_2,e_3$计算有：
              \[
                \operatorname{im}\sigma
                =\operatorname{span}(\sigma(e_1),\sigma(e_2),\sigma(e_3))
                =\operatorname{span}(
                        (1,1,0)^{\mathrm{T}},
                        (1,0,1)^{\mathrm{T}},
                        (0,-1,0)^{\mathrm{T}}
                    )
                =\mathbf{R}^3
              \]
              对于核空间，解方程$\sigma(\alpha)=0$即可. 我们也可以用更简洁的方式书写：
              \[
                \ker\sigma
                =\{
                    (x_1,x_2,x_3)^{\mathrm{T}} \mid
                    \sigma(x_1,x_2,x_3)^{\mathrm{T}}
                    =(0,0,0)^{\mathrm{T}}
                \}
                =\{(0,0,0)^{\mathrm{T}}\}
              \]
              即方程只有零解，核空间可以记为$\ker\sigma=\{0\}$（只含零元的空间的一般记法）.

        \item 根据原讲义“线性映射矩阵表示”处，应先写出$\sigma$在出发空间一组基（按题目要求是$\mathbf{R}^3$自然基）下的像，并将像表示为到达空间基（按题目要求是$\mathbf{R}^3$自然基）的线性组合，即
            \begin{gather*}
                \sigma(e_1) = (1,1,0)^{\mathrm{T}}
                =e_1+e_2=(e_1,e_2,e_3)\begin{pmatrix}
                    1 \\ 1 \\ 0
                \end{pmatrix} \\
                \sigma(e_2) = (1,0,1)^{\mathrm{T}}
                =e_1+e_3=(e_1,e_2,e_3)\begin{pmatrix}
                    1 \\ 0 \\ 1
                \end{pmatrix} \\
                \sigma(e_3) = (0,-1,0)^{\mathrm{T}}
                =-e_2=(e_1,e_2,e_3)\begin{pmatrix}
                    0 \\ -1 \\ 0
                \end{pmatrix}
            \end{gather*}
            接下来把坐标依次按列称矩阵就得到了本题需要求解的矩阵：
            \[
                \mathbf{M}(\sigma)=\begin{pmatrix}
                    1 & 1 & 0  \\
                    1 & 0 & -1 \\
                    0 & 1 & 0
                \end{pmatrix}
            \]
    \end{enumerate}
\end{solve}

\begin{exercise}[补充训练 L7-例-2]\label{prob:L7-例-2}
设$\sigma\colon\mathbf{R}^2\to\mathbf{R}^2$是绕原点逆时针旋转$\theta$角的变换，求$\sigma$在$\mathbf{R}^2$的自然基下的矩阵表示.
\end{exercise}
\noindent 来源：\href{https://github.com/yhwu-is/Linear-Algebra-Left-Undone/blob/new/%E8%AE%B2%E4%B9%89/%E4%B8%93%E9%A2%98/7%20%E7%BA%BF%E6%80%A7%E6%98%A0%E5%B0%84%E7%9F%A9%E9%98%B5%E8%A1%A8%E7%A4%BA%E4%B8%8E%E7%9F%A9%E9%98%B5%E8%BF%90%E7%AE%97%E5%9F%BA%E7%A1%80.tex}{L7-例-2}。
\begin{solve}
求解的过程是很自然简单的，我们只需要考虑$\sigma$在常用基$e_1,e_2$下的像，即
    \[
    \begin{cases}
        \sigma(e_1)=\cos\theta\cdot e_1+\sin\theta\cdot e_2=(e_1,e_2)\begin{pmatrix}
            \cos\theta \\ \sin\theta
        \end{pmatrix} \\
        \sigma(e_2)=-\sin\theta\cdot e_1+\cos\theta\cdot e_2=(e_1,e_2)\begin{pmatrix}
            -\sin\theta \\ \cos\theta
        \end{pmatrix}
    \end{cases}
    \]
    故
    \[\mathbf{M}(\sigma)=\begin{pmatrix}
        \cos\theta & -\sin\theta \\
        \sin\theta & \cos\theta
    \end{pmatrix}.\]
\end{solve}

\begin{exercise}[补充训练 L7-例-3]\label{prob:L7-例-3}
设$A=\begin{pmatrix}
            1 & 0 & -1 \\
            1 & 1 & -3
        \end{pmatrix}, B=\begin{pmatrix}
            0 & 3 \\
            1 & 2 \\
            3 & 1
        \end{pmatrix}$，求$AB$和$BA$.
\end{exercise}
\noindent 来源：\href{https://github.com/yhwu-is/Linear-Algebra-Left-Undone/blob/new/%E8%AE%B2%E4%B9%89/%E4%B8%93%E9%A2%98/7%20%E7%BA%BF%E6%80%A7%E6%98%A0%E5%B0%84%E7%9F%A9%E9%98%B5%E8%A1%A8%E7%A4%BA%E4%B8%8E%E7%9F%A9%E9%98%B5%E8%BF%90%E7%AE%97%E5%9F%BA%E7%A1%80.tex}{L7-例-3}。
\begin{solve}
直接根据矩阵乘法定义计算即可：
    \begin{gather*}
        AB = \begin{pmatrix}
            1 \times 0 + 0 \times 1 + (-1) \times 3 & 1 \times 3 + 0 \times 2 + (-1) \times 1 \\
            1 \times 0 + 1 \times 1 + (-3) \times 3 & 1 \times 3 + 1 \times 2 + (-3) \times 1
        \end{pmatrix} = \begin{pmatrix}
                -3 & 2 \\
                -8 & 2
            \end{pmatrix} \\
        BA = \begin{pmatrix}
            0 \times 1 + 3 \times 1 & 0 \times 0 + 3 \times 1 & 0 \times (-1) + 3 \times (-3) \\
            1 \times 1 + 2 \times 1 & 1 \times 0 + 2 \times 1 & 1 \times (-1) + 2 \times (-3) \\
            3 \times 1 + 1 \times 1 & 3 \times 0 + 1 \times 1 & 3 \times (-1) + 1 \times (-3)
        \end{pmatrix} = \begin{pmatrix}
                3 & 3 & -9 \\
                3 & 2 & -7 \\
                4 & 1 & -6
            \end{pmatrix}
    \end{gather*}
\end{solve}

\begin{exercise}[补充训练 L7-例-6]\label{prob:L7-例-6}
展开矩阵多项式$(A+\lambda E)^n$.
\end{exercise}
\noindent 来源：\href{https://github.com/yhwu-is/Linear-Algebra-Left-Undone/blob/new/%E8%AE%B2%E4%B9%89/%E4%B8%93%E9%A2%98/7%20%E7%BA%BF%E6%80%A7%E6%98%A0%E5%B0%84%E7%9F%A9%E9%98%B5%E8%A1%A8%E7%A4%BA%E4%B8%8E%E7%9F%A9%E9%98%B5%E8%BF%90%E7%AE%97%E5%9F%BA%E7%A1%80.tex}{L7-例-6}。
\begin{solve}
由于$A$与$E$是可交换的，并且$A^nE^m=A^n$显然成立，因此我们结合中学学过的二项式展开，得到结果：
    \begin{align*}
        (A+\lambda E)^n & =\sum_{i=0}^nC_n^iA^i(\lambda E)^{n-i}    \\
                        & =\sum_{i=0}^nC_n^i\lambda^{n-i}A^iE^{n-i} \\
                        & =\sum_{i=0}^nC_n^i\lambda^{n-i}A^i.
    \end{align*}
\end{solve}

\begin{exercise}[补充训练 L7-例-7]\label{prob:L7-例-7}
设$f(x),g(x) \in \mathbf{F}[x],\enspace A,B \in \mathbf{M}_n(\mathbf{F})$. 证明：
    \begin{enumerate}
        \item $f(A)g(A)=g(A)f(A)$；

        \item 如果$AB=BA$，则$f(A)g(B)=g(B)f(A)$；
    \end{enumerate}
\end{exercise}
\noindent 来源：\href{https://github.com/yhwu-is/Linear-Algebra-Left-Undone/blob/new/%E8%AE%B2%E4%B9%89/%E4%B8%93%E9%A2%98/7%20%E7%BA%BF%E6%80%A7%E6%98%A0%E5%B0%84%E7%9F%A9%E9%98%B5%E8%A1%A8%E7%A4%BA%E4%B8%8E%E7%9F%A9%E9%98%B5%E8%BF%90%E7%AE%97%E5%9F%BA%E7%A1%80.tex}{L7-例-7}。
\begin{solve}
我们可以直接证明第二点，因为第一点是第二点的特例. 设$f(x)=\displaystyle\sum_{i=0}^ma_ix^i$，$g(x)=\displaystyle\sum_{j=0}^sb_jx^j$，$A^0=B^0=E$，则
    \begin{align*}
        f(A)g(B) & =\sum_{i=0}^ma_iA^i\cdot \sum_{j=0}^sb_jB^j\quad(AB=BA)                            \\
                 & =\sum_{k=0}^{m+s}\sum_{i+j=k}a_ib_jA^iB^j=\sum_{k=0}^{m+s}\sum_{i+j=k}b_ja_iB^jA^i \\
                 & =g(B)f(A).
    \end{align*}
    事实上由于$A\cdot A=A\cdot A$，因此$f(A)g(A)=g(A)f(A)$只是上面证明的结论的特例.
\end{solve}

\begin{exercise}[补充训练 L7-例-9]\label{prob:L7-例-9}
已知3阶矩阵$A=\begin{pmatrix}
            1 & 0 & 1 \\ 0 & -1 & 0 \\ -1 & 1 & -1
        \end{pmatrix}$. 定义$\mathbf{R}^{3 \times 3}$上的线性变换$\sigma(X)=AX,\enspace X \in \mathbf{R}^{3 \times 3}$. 求$\sigma$的像和核.
\end{exercise}
\noindent 来源：\href{https://github.com/yhwu-is/Linear-Algebra-Left-Undone/blob/new/%E8%AE%B2%E4%B9%89/%E4%B8%93%E9%A2%98/7%20%E7%BA%BF%E6%80%A7%E6%98%A0%E5%B0%84%E7%9F%A9%E9%98%B5%E8%A1%A8%E7%A4%BA%E4%B8%8E%E7%9F%A9%E9%98%B5%E8%BF%90%E7%AE%97%E5%9F%BA%E7%A1%80.tex}{L7-例-9}。
\begin{solve}
核空间求解较为简单，我们先求解核空间. 首先求解线性方程组$AY=0$，其中$Y$为列向量，解得其基础解系为$\eta=(1,0,-1)^\mathrm{T}$.

    记$X=(X_1,X_2,X_3)$，则$X\in\ker\sigma$即$AX=(AX_1,AX_2,AX_3)=O$，即$AX_1=AX_2=AX_3=0$，因此$X_1,X_2,X_3$都能由$\eta$线性表出，故
    \[X=(k_1\eta,k_2\eta,k_3\eta)=\begin{pmatrix}
            k_1 & k_2 & k_3 \\ 0 & 0 & 0 \\ -k_1 & -k_2 & -k_3
        \end{pmatrix},\enspace k_1,k_2,k_3\in\mathbf{R},\]
    即
\[X=k_1\begin{pmatrix}
            1 & 0 & 0 \\ 0 & 0 & 0 \\ -1 & 0 & 0
        \end{pmatrix}+k_2\begin{pmatrix}
            0 & 1 & 0 \\ 0 & 0 & 0 \\ 0 & -1 & 0
        \end{pmatrix}+k_3\begin{pmatrix}
            0 & 0 & 1 \\ 0 & 0 & 0 \\ 0 & 0 & -1
        \end{pmatrix}\]
，即$\ker\sigma$中所有元素均可由这三个矩阵线性表示，并且这三个矩阵显然线性无关，因此核空间就是这三个矩阵的线性组合，且核空间维数为3.

    注意到$\mathbf{R}^{3 \times 3}$的一组基为$E_{11},E_{12},E_{13},E_{21},E_{22},E_{23},E_{31},E_{32},E_{33}$，其中$E_{ij}$表示第$i$行第$j$列元素为1，其余元素为0的矩阵，例如$E_{23}=\begin{pmatrix}
            0 & 0 & 0 \\ 0 & 0 & 1 \\ 0 & 0 & 0
        \end{pmatrix}$.

    根据$\sigma$的定义我们可以求得
    \begin{align*}
        \sigma(E_{11})&=\sigma(E_{31})=\begin{pmatrix}
            1 & 0 & 0 \\ 0 & 0 & 0 \\ -1 & 0 & 0
        \end{pmatrix}, &
        \sigma(E_{12})&=\sigma(E_{32})=\begin{pmatrix}
            0 & 1 & 0 \\ 0 & 0 & 0 \\ 0 & -1 & 0
        \end{pmatrix}, \\
        \sigma(E_{13})&=\sigma(E_{33})=\begin{pmatrix}
            0 & 0 & 1 \\ 0 & 0 & 0 \\ 0 & 0 & -1
        \end{pmatrix}, &
        \sigma(E_{21})&=\begin{pmatrix}
            0 & 0 & 0 \\ -1 & 0 & 0 \\ 1 & 0 & 0
        \end{pmatrix}, \\
        \sigma(E_{22})&=\begin{pmatrix}
            0 & 0 & 0 \\ 0 & -1 & 0 \\ 0 & 1 & 0
        \end{pmatrix}, &
        \sigma(E_{23})&=\begin{pmatrix}
            0 & 0 & 0 \\ 0 & 0 & -1 \\ 0 & 0 & 1
        \end{pmatrix}.
    \end{align*}
    所以$\sigma$的像空间为上述六个矩阵线性扩张而成的空间，即\[\operatorname{im}\sigma=\operatorname{span}\{\sigma(E_{ij}):1\le i\le2,\ 1\le j\le3\}.\]. 又由原讲义“线性映射基本定理”处可知，$\dim\operatorname{im}\sigma=9-\dim\ker\sigma=6$，因此像空间就是这六个矩阵线性扩张而成的空间.
\end{solve}

\begin{exercise}[补充训练 L7-例-10]\label{prob:L7-例-10}
用解方程法求矩阵 $A = \begin{pmatrix}1 & -1 & 1 \\ 0 & 1 & 2 \\ 1 & 0 & 4\end{pmatrix}$ 的逆矩阵.
\end{exercise}
\noindent 来源：\href{https://github.com/yhwu-is/Linear-Algebra-Left-Undone/blob/new/%E8%AE%B2%E4%B9%89/%E4%B8%93%E9%A2%98/7%20%E7%BA%BF%E6%80%A7%E6%98%A0%E5%B0%84%E7%9F%A9%E9%98%B5%E8%A1%A8%E7%A4%BA%E4%B8%8E%E7%9F%A9%E9%98%B5%E8%BF%90%E7%AE%97%E5%9F%BA%E7%A1%80.tex}{L7-例-10}。
\begin{solve}
以 $A$ 为系数矩阵的非齐次线性方程组 $AX = b$，对于任意的 $b = (b_1, b_2, b_3)$，可以用高斯-若当消元法将其增广矩阵

    \[(A,b)=\left(\begin{array}{ccc:c}
                1 & -1 & 1 & b_1 \\
                0 & 1  & 2 & b_2 \\
                1 & 0  & 4 & b_3
            \end{array}\right)\] 化为：

    \[\left(\begin{array}{ccc:c}
                1 & 0 & 0 & 4b_1+4b_2-3b_3 \\
                0 & 1 & 0 & 2b_1+3b_2-2b_3 \\
                0 & 0 & 1 & -b_1-b_2+b_3
            \end{array}\right)\].

    因此，对任意的 $b$，方程组 $AX=b$ 有唯一解：
    \[X=\begin{pmatrix}x_1\\x_2\\x_3\end{pmatrix}=\begin{pmatrix}
            4b_1+4b_2-3b_3 \\2b_1+3b_2-2b_3\\-b_1-b_2+b_3
        \end{pmatrix}=\begin{pmatrix}
            4  & 4  & -3 \\
            2  & 3  & -2 \\
            -1 & -1 & 1
        \end{pmatrix}\begin{pmatrix}b_1\\b_2\\b_3\end{pmatrix}\] 即：
    \[A^{-1}=\begin{pmatrix}
            4  & 4  & -3 \\
            2  & 3  & -2 \\
            -1 & -1 & 1
        \end{pmatrix}\]
\end{solve}

\begin{exercise}[补充训练 L7-例-11]\label{prob:L7-例-11}
设$\alpha=(1,-1,2)^\mathrm{T},\enspace\beta=(3,1,-2)^\mathrm{T},\enspace A=\alpha\beta^\mathrm{T}$，求$A^n$.
\end{exercise}
\noindent 来源：\href{https://github.com/yhwu-is/Linear-Algebra-Left-Undone/blob/new/%E8%AE%B2%E4%B9%89/%E4%B8%93%E9%A2%98/7%20%E7%BA%BF%E6%80%A7%E6%98%A0%E5%B0%84%E7%9F%A9%E9%98%B5%E8%A1%A8%E7%A4%BA%E4%B8%8E%E7%9F%A9%E9%98%B5%E8%BF%90%E7%AE%97%E5%9F%BA%E7%A1%80.tex}{L7-例-11}。
\begin{solve}
由于$A^2=\alpha\beta^\mathrm{T}\alpha\beta^\mathrm{T}=\alpha(\beta^\mathrm{T}\alpha)\beta^\mathrm{T}=kA$，其中$k=\beta^\mathrm{T}\alpha=-2$，则$A^2=-2A$. 则可递推得到$A^n=(-2)^{n-1}A=(-2)^{n-1}\begin{pmatrix}
            3 & 1 & -2 \\ -3 & -1 & 2 \\ 6 & 2 & -4
        \end{pmatrix}$，原因在于$A^n$展开后中间会出现$n-1$个$\beta^\mathrm{T}\alpha$.
\end{solve}

\begin{exercise}[补充训练 L7-例-12]\label{prob:L7-例-12}
证明以下几点性质：
    \begin{enumerate}
        \item 反对称矩阵主对角元均为 $0$；

        \item $AA^\mathrm{T}$ 和 $A^\mathrm{T}A$ 均为对称矩阵；

        \item 设 $A, B$ 分别为 $n$ 阶对称和反对称矩阵，则 $AB + BA$ 是反对称矩阵；

        \item 对称矩阵的乘积不一定对称；

        \item 可逆的对称（反对称）矩阵的逆矩阵也是对称（反对称）矩阵.
    \end{enumerate}
\end{exercise}
\noindent 来源：\href{https://github.com/yhwu-is/Linear-Algebra-Left-Undone/blob/new/%E8%AE%B2%E4%B9%89/%E4%B8%93%E9%A2%98/7%20%E7%BA%BF%E6%80%A7%E6%98%A0%E5%B0%84%E7%9F%A9%E9%98%B5%E8%A1%A8%E7%A4%BA%E4%B8%8E%E7%9F%A9%E9%98%B5%E8%BF%90%E7%AE%97%E5%9F%BA%E7%A1%80.tex}{L7-例-12}。
\begin{solve}
\begin{enumerate}
        \item 由于$A$为反对称矩阵，因此根据定义有 $a_{ii} = -a_{ii}$，即 $a_{ii} = 0$；

        \item 由于 $(AA^\mathrm{T})^\mathrm{T} = (A^\mathrm{T})^\mathrm{T}A^\mathrm{T} = AA^\mathrm{T}$，因此 $AA^\mathrm{T}$ 为对称矩阵；同理可证 $A^\mathrm{T}A$ 为对称矩阵；

        \item 由于$A,B$分别为对称和反对称矩阵，因此$(AB+BA)^\mathrm{T}=B^\mathrm{T}A^\mathrm{T}+A^\mathrm{T}B^\mathrm{T}=-AB-BA=-(AB+BA)$，因此$AB+BA$为反对称矩阵；

        \item 注意$(AB)^\mathrm{T}=B^\mathrm{T}A^\mathrm{T}=BA$，因为矩阵乘法不一定可交换，因此$AB$不一定对称；

        \item 因为$A$可逆有$(A^{-1})^\mathrm{T}=(A^\mathrm{T})^{-1}=A^{-1}$，因此$A^{-1}$为对称矩阵；同理可证$A$反对称的情况.
    \end{enumerate}
\end{solve}

\begin{exercise}[补充训练 L7-例-14]\label{prob:L7-例-14}
$a,b,c,d$是四个实数. 证明$\begin{cases}
            a^2+b^2=1 \\
            c^2+d^2=1 \\
            ac+bd=0
        \end{cases}$成立的充分必要条件是$\begin{cases}
            a^2+c^2=1 \\
            b^2+d^2=1 \\
            ab+cd=0
        \end{cases}$.
\end{exercise}
\noindent 来源：\href{https://github.com/yhwu-is/Linear-Algebra-Left-Undone/blob/new/%E8%AE%B2%E4%B9%89/%E4%B8%93%E9%A2%98/7%20%E7%BA%BF%E6%80%A7%E6%98%A0%E5%B0%84%E7%9F%A9%E9%98%B5%E8%A1%A8%E7%A4%BA%E4%B8%8E%E7%9F%A9%E9%98%B5%E8%BF%90%E7%AE%97%E5%9F%BA%E7%A1%80.tex}{L7-例-14}。
\begin{solve}
设$A=\begin{pmatrix}
            a & b \\ c & d
        \end{pmatrix}$，则有
    \[AA^\mathrm{T}=\begin{pmatrix}
            a^2+b^2 & ac+bd \\ ac+bd & c^2+d^2
        \end{pmatrix},\enspace A^\mathrm{T}A=\begin{pmatrix}
            a^2+c^2 & ab+cd \\ ab+cd & b^2+d^2
        \end{pmatrix}.\]
    因此题中的充要条件可以转化为$AA^\mathrm{T}=E$是$A^\mathrm{T}A=E$的充要条件. 这是显然的，因为$AA^\mathrm{T}=E\iff A^{-1}=A^\mathrm{T}\iff A^\mathrm{T}A=E$成立.
\end{solve}

\begin{exercise}[补充训练 L7-例-15]\label{prob:L7-例-15}
设$n$阶矩阵$A$分块为$A=\begin{pmatrix}
                      B & O \\ C & D
                  \end{pmatrix}$，其中$B,D$分别为$k$阶、$m$阶矩阵，求当$B,D$可逆时的$A^{-1}$.
\end{exercise}
\noindent 来源：\href{https://github.com/yhwu-is/Linear-Algebra-Left-Undone/blob/new/%E8%AE%B2%E4%B9%89/%E4%B8%93%E9%A2%98/7%20%E7%BA%BF%E6%80%A7%E6%98%A0%E5%B0%84%E7%9F%A9%E9%98%B5%E8%A1%A8%E7%A4%BA%E4%B8%8E%E7%9F%A9%E9%98%B5%E8%BF%90%E7%AE%97%E5%9F%BA%E7%A1%80.tex}{L7-例-15}。
\begin{solve}
本题我们使用的方法非常直接，就是直接设出$A^{-1}$的形式，然后验证即可. 之后我们还会学习一种基于分块矩阵初等变换的进阶方法（事实上考试如果考察的话基本是本题的解法，分块矩阵初等变换是在教材中是小字部分）. 设$A^{-1}=\begin{pmatrix}
                      X & Y \\ Z & T
                  \end{pmatrix}$，其中$X,T$分别为$k,m$阶矩阵，那么我们有
              \[\begin{pmatrix}
                      B & O \\ C & D
                  \end{pmatrix}\begin{pmatrix}
                      X & Y \\ Z & T
                  \end{pmatrix}=\begin{pmatrix}
                      BX & BY \\ CX+DZ & CY+DT
                  \end{pmatrix}=\begin{pmatrix}
                      E_k & O \\ O & E_m
                  \end{pmatrix},\]
              又由题意$A$可逆，$B,D$可逆，因此$BX=E_k$可得$X=B^{-1}$，$BY=O$可得$Y=O$，$CY+DT=DT=E_m$可得$T=D^{-1}$，$CX+DZ=CB^{-1}+DZ=O$可得$Z=-D^{-1}CB^{-1}$，因此
              \[A^{-1}=\begin{pmatrix}
                      B^{-1} & O \\ -D^{-1}CB^{-1} & D^{-1}
                  \end{pmatrix}.\]
\end{solve}

\begin{exercise}[补充训练 L7-例-16]\label{prob:L7-例-16}
设\[A=\begin{pmatrix}
            1 & 2 & 0  & 0  & 0  \\
            2 & 5 & 0  & 0  & 0  \\
            0 & 0 & -2 & 1  & 0  \\
            0 & 0 & 0  & -2 & 1  \\
            0 & 0 & 0  & 0  & -2
        \end{pmatrix},\enspace B=\begin{pmatrix}
            1  & 0 & 1 & 0 \\
            -1 & 2 & 3 & 0 \\
            1  & 2 & 0 & 4 \\
            0  & 1 & 2 & 4 \\
            0  & 0 & 1 & 4
        \end{pmatrix}\]
    利用分块矩阵的方法，求$A^2,\enspace AB,\enspace A^\mathrm{T},\enspace A^{-1}$.
\end{exercise}
\noindent 来源：\href{https://github.com/yhwu-is/Linear-Algebra-Left-Undone/blob/new/%E8%AE%B2%E4%B9%89/%E4%B8%93%E9%A2%98/7%20%E7%BA%BF%E6%80%A7%E6%98%A0%E5%B0%84%E7%9F%A9%E9%98%B5%E8%A1%A8%E7%A4%BA%E4%B8%8E%E7%9F%A9%E9%98%B5%E8%BF%90%E7%AE%97%E5%9F%BA%E7%A1%80.tex}{L7-例-16}。
\begin{solve}
按照前两行和后三行分块，$A=\operatorname{diag}(A_1,A_2)$。乘积、转置与逆分别对对角块计算；矩阵 $B$ 的行按相同位置分块。直接计算得到
\[A^2=\begin{pmatrix}5&12&0&0&0\\12&29&0&0&0\\0&0&4&-4&1\\0&0&0&4&-4\\0&0&0&0&4\end{pmatrix},\qquad AB=\begin{pmatrix}-1&4&7&0\\-3&10&17&0\\-2&-3&2&-4\\0&-2&-3&-4\\0&0&-2&-8\end{pmatrix}.\]
\[A^{\mathrm T}=\begin{pmatrix}1&2&0&0&0\\2&5&0&0&0\\0&0&-2&0&0\\0&0&1&-2&0\\0&0&0&1&-2\end{pmatrix},\qquad A^{-1}=\begin{pmatrix}5&-2&0&0&0\\-2&1&0&0&0\\0&0&-\frac{1}{2}&-\frac{1}{4}&-\frac{1}{8}\\0&0&0&-\frac{1}{2}&-\frac{1}{4}\\0&0&0&0&-\frac{1}{2}\end{pmatrix}.\]
其中二阶块的行列式为 $1$，三阶上三角块的行列式为 $-8$，故两块均可逆；验证 $AA^{-1}=I_5$ 即可检验结果。
\end{solve}

\begin{exercise}[补充训练 S-BASIS-1]\label{prob:S-BASIS-1}
在 $\mathbf R^3$ 中，基 $\mathcal B$、$\mathcal C$ 的向量分别为 $P=\begin{pmatrix}1&-3&0\\0&1&0\\1&0&1\end{pmatrix}$、$Q=\begin{pmatrix}1&0&0\\0&1&0\\0&-2&1\end{pmatrix}$ 的列。若 $[v]_{\mathcal B}=\begin{pmatrix}-1\\1\\3\end{pmatrix}$，求 $v$、$[v]_{\mathcal C}$ 及从 $\mathcal B$ 坐标到 $\mathcal C$ 坐标的变换矩阵。
\end{exercise}
\noindent 编者补充：S-BASIS-1；相关主题见原讲义对应讲，题目数据与解答由本选编另行构造。
\begin{solve}
两矩阵行列式都为 $1$，因此各自的列确实构成基。由坐标定义，$v=P[v]_{\mathcal B}$，所以 $v=\begin{pmatrix}-4\\1\\2\end{pmatrix}$。又有 $v=Q[v]_{\mathcal C}$，故
\[[v]_{\mathcal C}=Q^{-1}P[v]_{\mathcal B},\qquad Q^{-1}P=\begin{pmatrix}1&-3&0\\0&1&0\\1&2&1\end{pmatrix}.\]
代入题给坐标得 $[v]_{\mathcal C}=\begin{pmatrix}-4\\1\\4\end{pmatrix}$。验证 $Q[v]_{\mathcal C}=P[v]_{\mathcal B}$ 可确认变换方向。反向变换为 $P^{-1}Q$，它是 $Q^{-1}P$ 的逆矩阵；一般不能把这两个乘积交换。这里同一个几何向量的坐标随基变化，但向量自身未改变。
\end{solve}

\section{补充训练：矩阵相抵与秩}

本节包含 22 道训练题，建议在完成相关专题真题后选择性练习。

\begin{exercise}[补充训练 L8-A-4]\label{prob:L8-A-4}
给定$\mathbf{R}^4$的两组基
        \begin{gather*}
            \alpha_1=(1,1,1,1),\ \alpha_2=(1,1,-1,-1),\ \alpha_3=(1,-1,1,-1),\ \alpha_4=(1,-1,-1,1) \\
            \beta_1=(1,1,0,1),\ \beta_2=(2,1,3,1),\ \beta_3=(1,1,0,0),\ \beta_4=(0,1,-1,-1)
        \end{gather*}
        求由基$\alpha_1,\alpha_2,\alpha_3,\alpha_4$到基$\beta_1,\beta_2,\beta_3,\beta_4$的过渡矩阵，并求向量$\xi=(1,0,0,-1)$在基$\alpha_1,\alpha_2,\alpha_3,\alpha_4$下的坐标.
\end{exercise}
\noindent 来源：\href{https://github.com/yhwu-is/Linear-Algebra-Left-Undone/blob/new/%E8%AE%B2%E4%B9%89/%E4%B8%93%E9%A2%98/8%20%E7%9B%B8%E6%8A%B5%E6%A0%87%E5%87%86%E5%BD%A2.tex}{L8-A-4}。
\begin{solve}
首先取 $ \mathbf{R}^4 $ 的自然基 $\varepsilon_1, \varepsilon_2, \varepsilon_3, \varepsilon_4$，那么
            \[ (\alpha_1, \alpha_2, \alpha_3, \alpha_4) = (\varepsilon_1, \varepsilon_2, \varepsilon_3, \varepsilon_4) A,\enspace (\beta_1, \beta_2, \beta_3, \beta_4) = (\varepsilon_1, \varepsilon_2, \varepsilon_3, \varepsilon_4) B \]
            其中
            \[ A = \begin{pmatrix}
                    1 & 1  & 1  & 1  \\
                    1 & 1  & -1 & -1 \\
                    1 & -1 & 1  & -1 \\
                    1 & -1 & -1 & 1
                \end{pmatrix},\enspace B = \begin{pmatrix}
                    1 & 2 & 1 & 0  \\
                    1 & 1 & 1 & 1  \\
                    0 & 3 & 0 & -1 \\
                    1 & 1 & 0 & -1
                \end{pmatrix} \]
            由此可知
            \[ (\beta_1, \beta_2, \beta_3, \beta_4) = (\varepsilon_1, \varepsilon_2, \varepsilon_3, \varepsilon_4) B = (\alpha_1, \alpha_2, \alpha_3, \alpha_4) A^{-1} B \]
            所以由基 $ \alpha_1, \alpha_2, \alpha_3, \alpha_4 $ 到基 $ \beta_1, \beta_2, \beta_3, \beta_4 $ 的过渡矩阵为
            \[ A^{-1} B = \dfrac{1}{4} \begin{pmatrix}
                    1 & 1  & 1  & 1  \\
                    1 & 1  & -1 & -1 \\
                    1 & -1 & 1  & -1 \\
                    1 & -1 & -1 & 1
                \end{pmatrix} \begin{pmatrix}
                    1 & 2 & 1 & 0  \\
                    1 & 1 & 1 & 1  \\
                    0 & 3 & 0 & -1 \\
                    1 & 1 & 0 & -1
                \end{pmatrix} = \dfrac{1}{4} \begin{pmatrix}
                    3  & 7  & 2 & -1 \\
                    1  & -1 & 2 & 3  \\
                    -1 & 3  & 0 & -1 \\
                    1  & -1 & 0 & -1
                \end{pmatrix} \]
            另外，由于 $ \xi $ 在基 $ \varepsilon_1, \varepsilon_2, \varepsilon_3, \varepsilon_4 $ 下的坐标为 $ X_0 = (1, 0, 0, -1)^{\mathrm{T}} $，所以在 $ \alpha_1, \alpha_2, \alpha_3, \alpha_4 $ 下的坐标为
            \[ A^{-1} X_0 = \dfrac{1}{4} \begin{pmatrix}
                    1 & 1  & 1  & 1  \\
                    1 & 1  & -1 & -1 \\
                    1 & -1 & 1  & -1 \\
                    1 & -1 & -1 & 1
                \end{pmatrix} \begin{pmatrix}
                    1 \\ 0 \\ 0 \\ -1
                \end{pmatrix} = \begin{pmatrix}
                    0 \\ \dfrac{1}{2} \\ \dfrac{1}{2} \\ 0
                \end{pmatrix} \]
\end{solve}

\begin{exercise}[补充训练 L8-A-5]\label{prob:L8-A-5}
证明：矩阵添加一列（或一行），其秩或不变，或增加1.
\end{exercise}
\noindent 来源：\href{https://github.com/yhwu-is/Linear-Algebra-Left-Undone/blob/new/%E8%AE%B2%E4%B9%89/%E4%B8%93%E9%A2%98/8%20%E7%9B%B8%E6%8A%B5%E6%A0%87%E5%87%86%E5%BD%A2.tex}{L8-A-5}。
\begin{solve}
证明：考虑矩阵 $A$ 的行向量组的极大线性无关组，若添加的一行可由其极大线性无关组线性表示，则秩不变. 否则秩增加 $1$.
\end{solve}

\begin{exercise}[补充训练 L8-A-6]\label{prob:L8-A-6}
设$A$是$s \times n$矩阵，$B$是$A$前$m$行构成的$m \times n$矩阵，证明：$r(B) \geqslant r(A) + m - s$.
\end{exercise}
\noindent 来源：\href{https://github.com/yhwu-is/Linear-Algebra-Left-Undone/blob/new/%E8%AE%B2%E4%B9%89/%E4%B8%93%E9%A2%98/8%20%E7%9B%B8%E6%8A%B5%E6%A0%87%E5%87%86%E5%BD%A2.tex}{L8-A-6}。
\begin{solve}
证明：设 $A$ 的行向量组为 $\alpha_1,\alpha_2,\ldots,\alpha_s$，$r(A)=r$; $B$ 的行向量组为 $\alpha_1,\alpha_2,\ldots,\alpha_m ,r(B)=k$.
            不妨设：$B$ 的行向量组的极大线性无关组为 $\alpha_1,\alpha_2,\ldots,\alpha_k,\alpha_{i_1},\ldots,\alpha_{i_{r-k}}$，其中 $\alpha_{i_1},\ldots,\alpha_{i_{r-k}}$（共 $r-k$ 个向量）是包含在 $\alpha_{m+1},\ldots,\alpha_s$（共 $s-m$ 个向量）之中的. 显然有 $r-k \leqslant s-m $，即
            \[r(B)=k\geqslant r+m-s=r(A)+m-s.\]
\end{solve}

\begin{exercise}[补充训练 L8-B-2]\label{prob:L8-B-2}
设$W$是$n$维线性空间$V$的一个非平凡子空间，$W$中取一组基$\delta_1,\ldots,\delta_m$，按如下两种方式将其扩充为$V$的一组基：
        \begin{align*}
            B_1 & =\{\delta_1,\ldots,\delta_m,\alpha_{m+1},\ldots,\alpha_n\} \\
            B_2 & =\{\delta_1,\ldots,\delta_m,\beta_{m+1},\ldots,\beta_n\}
        \end{align*}
        设基$B_1$到$B_2$的过渡矩阵为$P$，求商空间$V/W$的基$\alpha_{m+1}+W,\ldots,\alpha_n+W$到$\beta_{m+1}+W,\ldots,\beta_n+W$的过渡矩阵.
\end{exercise}
\noindent 来源：\href{https://github.com/yhwu-is/Linear-Algebra-Left-Undone/blob/new/%E8%AE%B2%E4%B9%89/%E4%B8%93%E9%A2%98/8%20%E7%9B%B8%E6%8A%B5%E6%A0%87%E5%87%86%E5%BD%A2.tex}{L8-B-2}。
\begin{solve}
$V$ 的基 $B_1$ 到 $B_2$ 的过渡矩阵 $P$ 具有下述形式：
          \[P=\begin{pmatrix}E_m & B_1 \\ 0 & B_2\end{pmatrix}\]
          其中 $B_1,B_2$ 分别是域 $\mathbf{F}$ 上 $m\times (n-m),(n-m)\times (n-m)$ 矩阵，
          \[\beta_j=b_{j1}\delta_1+\cdots+b_{jm}\delta_m+b_{j,m+1}\alpha_{m+1}+\cdots+b_{jn}\alpha_n\]
          其中 $j=m+1,\ldots,n$. 于是
          \[\beta_j+W=b_{j,m+1}(\alpha_{m+1}+W)+\cdots+b_{jn}(\alpha_n+W)\]
          因此商空间 $V/W$ 的基 $\alpha_{m+1}+W,\ldots,\alpha_n+W$ 到 $\beta_{m+1}+W,\ldots,\beta_n+W$ 的过渡矩阵是 $B_2$.
\end{solve}

\begin{exercise}[补充训练 L8-B-6]\label{prob:L8-B-6}
证明：当$n$为奇数时，$\alpha_1,\alpha_2,\ldots,\alpha_n$线性无关的充要条件是$\alpha_1+\alpha_2,\alpha_2+\alpha_3,\ldots,\alpha_n+\alpha_1$线性无关.
\end{exercise}
\noindent 来源：\href{https://github.com/yhwu-is/Linear-Algebra-Left-Undone/blob/new/%E8%AE%B2%E4%B9%89/%E4%B8%93%E9%A2%98/8%20%E7%9B%B8%E6%8A%B5%E6%A0%87%E5%87%86%E5%BD%A2.tex}{L8-B-6}。
\begin{solve}
证明：我们设 $ \beta_1 = \alpha_1 + \alpha_2,\enspace \ldots, \beta_{n - 1} = \alpha_{n - 1} + \alpha_n,\enspace \beta_n = \alpha_n + \alpha_1 $. 由于
            \[ (\beta_1, \beta_2, \ldots, \beta_n) = (\alpha_1, \alpha_2, \ldots, \alpha_n) A \]
            其中
            \[ A = \begin{pmatrix}
                    1 &   &        &   & 1 \\
                    1 & 1 &        &   &   \\
                        & 1 & \ddots &   &   \\
                        &   & \ddots & 1 &   \\
                        &   &        & 1 & 1
                \end{pmatrix} \]
            满足 $ |A| = 1 + (-1)^{n + 1} = 2 $，所以 $ A $ 可逆. 根据教材原讲义命题 3.10.3 可得 $ \alpha_1, \alpha_2, \ldots, \alpha_n $ 与 $ \beta_1, \beta_2, \ldots, \beta_n $ 等价. 当然，这说明 $ \alpha_1, \alpha_2, \ldots, \alpha_n $ 线性无关的充要条件是 $ \beta_1, \beta_2, \ldots, \beta_n $ 线性无关.
\end{solve}

\begin{exercise}[补充训练 L8-B-7]\label{prob:L8-B-7}
设$B=\{\beta_1,\beta_2,\ldots,\beta_n\}$是实数域$\mathbf{R}$上的线性空间$V$的一组基，$T \in \mathcal{L}(V),\enspace T(\beta_1)=\beta_2,T(\beta_2)=\beta_3,\ldots,T(\beta_{n-1}) = \beta_n,T(\beta_n)=\displaystyle\sum_{i=1}^{n}a_i\beta_i(a_i \in \mathbf{R})$，求$T$关于基$B$的表示矩阵，并求在什么条件下$T$是同构映射.
\end{exercise}
\noindent 来源：\href{https://github.com/yhwu-is/Linear-Algebra-Left-Undone/blob/new/%E8%AE%B2%E4%B9%89/%E4%B8%93%E9%A2%98/8%20%E7%9B%B8%E6%8A%B5%E6%A0%87%E5%87%86%E5%BD%A2.tex}{L8-B-7}。
\begin{solve}
$ T(\beta_1, \beta_2, \ldots, \beta_n) = (\beta_1, \beta_2, \ldots, \beta_n)A $，其中
            \[ A = \begin{pmatrix}
                      &   &        &   & a_1    \\
                    1 &   &        &   & a_2    \\
                      & 1 &        &   & a_3    \\
                      &   & \ddots &   & \vdots \\
                      &   &        & 1 & a_n
                \end{pmatrix} \]
            是 $ T $ 关于基 $ B $ 的表示矩阵.

            $ T $ 是同构 $ \iff T $ 是双射 $ \iff r(T) = n $，满秩. 所以当 $ a_1 \neq 0 $ 时，$ r(T) = n $ 满秩，此时 $ T $ 是同构映射.
\end{solve}

\begin{exercise}[补充训练 L8-B-8]\label{prob:L8-B-8}
设矩阵空间$\mathbf{R}^{2\times 2}$的子空间为
        \[V=\{X=(x_{ij})_{2\times 2} \mid x_{11}+x_{12}+x_{21}=0,\enspace x_{ij}\in \mathbf{R}\}\]
        V中的线性变换为$\sigma(X)=X+X^\mathrm{T}$，求$V$的一组基，使得$\sigma$在该基下的矩阵表示为对角矩阵.
\end{exercise}
\noindent 来源：\href{https://github.com/yhwu-is/Linear-Algebra-Left-Undone/blob/new/%E8%AE%B2%E4%B9%89/%E4%B8%93%E9%A2%98/8%20%E7%9B%B8%E6%8A%B5%E6%A0%87%E5%87%86%E5%BD%A2.tex}{L8-B-8}。
\begin{solve}
我们可以类比线性空间
          \[ W = \{ x \in \mathbf{R}^4 \mid x_1 + x_2 + x_3 = 0 \} \]
          $ W $ 的基只需求解线性方程组 $ x_1 + x_2 + x_3 = 0 $ 即可，得到基础解系为 $ (-1, 0, 1, 0)^{\mathrm{T}},\allowbreak (-1, 1, 0, 0)^{\mathrm{T}},\allowbreak (0, 0, 0, 1)^{\mathrm{T}} $.

          换回 $ V $，即有基为 $ A_1 = \begin{pmatrix} -1 & 0 \\ 1 & 0 \end{pmatrix}, A_2 = \begin{pmatrix} -1 & 1 \\ 0 & 0 \end{pmatrix}, A_3 = \begin{pmatrix} 0 & 0 \\ 0 & 1 \end{pmatrix} $. 而
          \begin{align*}
                         & \sigma(A_1) = \begin{pmatrix} -2 & 1 \\ 1 & 0 \end{pmatrix} = A_1 + A_2                      \\
                         & \sigma(A_2) = \begin{pmatrix} -2 & 1 \\ 1 & 0 \end{pmatrix} = A_1 + A_2                      \\
                         & \sigma(A_3) = \begin{pmatrix} -2 & 1 \\ 1 & 0 \end{pmatrix} = 2 A_3                          \\
              \implies{} & \sigma(A_1 + A_2) = 2(A_1 + A_2),\enspace \sigma(A_1 - A_2) = 0,\enspace \sigma(A_3) = 2 A_3
          \end{align*}
          取基 $ A_1 - A_2, A_1 + A_2, A_3 $，有
          \[ (\sigma(A_1 - A_2), \sigma(A_1 + A_2), \sigma(A_3)) = (A_1 - A_2, A_1 + A_2, A_3) = \begin{pmatrix}
                  0 & 0 & 0 \\
                  0 & 2 & 0 \\
                  0 & 0 & 2
              \end{pmatrix} \]
          为对角矩阵.
\end{solve}

\begin{exercise}[补充训练 L8-B-9]\label{prob:L8-B-9}
设$A$为$n$阶实反对称矩阵，证明：$E-A$ 可逆.
\end{exercise}
\noindent 来源：\href{https://github.com/yhwu-is/Linear-Algebra-Left-Undone/blob/new/%E8%AE%B2%E4%B9%89/%E4%B8%93%E9%A2%98/8%20%E7%9B%B8%E6%8A%B5%E6%A0%87%E5%87%86%E5%BD%A2.tex}{L8-B-9}。
\begin{solve}
反证法，假设不可逆，则存在非零实向量$x=(x_1,\cdots,x_n)^\mathrm{T}$使得$(E-A)x=0$，即$Ax=x$，又反对称矩阵有一个性质，即对任意的$x$，$x^\mathrm{T}Ax=0$，而$x^\mathrm{T}Ax=x^\mathrm{T}x=x_1^2+\cdots+x_n^2$，又由实数条件可知$x$为零向量，矛盾.
\end{solve}

\begin{exercise}[补充训练 L8-B-10]\label{prob:L8-B-10}
设
        \[B_1=\left\{\begin{pmatrix}
                1 & 0 \\ 0 & 0
            \end{pmatrix},\begin{pmatrix}
                0 & 1 \\ 0 & 0
            \end{pmatrix},\begin{pmatrix}
                0 & 0 \\ 1 & 0
            \end{pmatrix}\begin{pmatrix}
                0 & 0 \\ 0 & 1
            \end{pmatrix}\right\},\]
        \[B_2=\left\{\begin{pmatrix}
                1 & 0 \\ 0 & 0
            \end{pmatrix},\begin{pmatrix}
                1 & 1 \\ 0 & 0
            \end{pmatrix},\begin{pmatrix}
                1 & 1 \\ 1 & 0
            \end{pmatrix}\begin{pmatrix}
                1 & 1 \\ 1 & 1
            \end{pmatrix}\right\}.\]
        \begin{enumerate}
            \item 证明：$B_2$也是线性空间$\mathbf{M}_2(\mathbf{R})$的基；

            \item 求基$B_2$变为基$B_1$的变换矩阵；

            \item 求$\mathbf{M}_2(\mathbf{R})$的一组基$B_3=\{A_1,A_2,A_3,A_4\}$，使得$A_i^2=A_i,\enspace i=1,2,3,4$；

            \item 已知矩阵$A$关于基$B_2$的坐标为$(1,1,1,1)^\mathrm{T}$，求$A$关于基$B_3$的坐标.
        \end{enumerate}
\end{exercise}
\noindent 来源：\href{https://github.com/yhwu-is/Linear-Algebra-Left-Undone/blob/new/%E8%AE%B2%E4%B9%89/%E4%B8%93%E9%A2%98/8%20%E7%9B%B8%E6%8A%B5%E6%A0%87%E5%87%86%E5%BD%A2.tex}{L8-B-10}。
\begin{solve}
记 $ B_1 = \{\vec{e}_{11}, \vec{e}_{12}, \vec{e}_{21}, \vec{e}_{22}\} $，$ B_2 = \{\vec{g}_1, \vec{g}_2, \vec{g}_3, \vec{g}_4\} $.
            \begin{enumerate}
                \item 设 $ k_1 \vec{g}_1 + k_2 \vec{g}_2 + k_3 \vec{g}_3 + k_4 \vec{g}_4 = O $，即
                      \[ \begin{pmatrix} k_1 + k_2 + k_3 + k_4 & k_2 + k_3 + k_4 \\ k_3 + k_4 & k_4 \end{pmatrix} = \begin{pmatrix} 0 & 0 \\ 0 & 0 \end{pmatrix} \]
                      这个方程对应的四元齐次线性方程组的解为
                      \[ k_1 = k_2 = k_3 = k_4 = 0 \]
                      所以 $ \vec{g}_1, \vec{g}_2, \vec{g}_3, \vec{g}_4 $ 线性无关，从而是 $ \mathbf{M}_2(\mathbf{R}) $ 的一组基.

                \item \label{item:11:B:4:2}
                      $ \mathbf{M}_2(\mathbf{R}) \cong \mathbf{R}^4 $，所以 $ \vec{e}_{11}, \vec{e}_{12}, \vec{e}_{21}, \vec{e}_{22} $ 可表示为 $ \mathbf{R}^4 $ 的自然基 $ \vec{e}_1, \vec{e}_2, \vec{e}_3, \vec{e}_4 $，而 $ \vec{g}_1, \vec{g}_2, \vec{g}_3, \vec{g}_4 $ 可表示为 $ (1, 0, 0, 0)^{\mathrm{T}},\allowbreak (1, 1, 0, 0)^{\mathrm{T}},\allowbreak (1, 1, 1, 0)^{\mathrm{T}},\allowbreak (1, 1, 1, 1)^{\mathrm{T}} $. 于是由
                      \[ (\vec{g}_1, \vec{g}_2, \vec{g}_3, \vec{g}_4) = (\vec{e}_{11}, \vec{e}_{12}, \vec{e}_{21}, \vec{e}_{22}) C \]
                      可得
                      \[ (\vec{e}_{11}, \vec{e}_{12}, \vec{e}_{21}, \vec{e}_{22}) = (\vec{g}_1, \vec{g}_2, \vec{g}_3, \vec{g}_4) C^{-1} \]
                      其中
                      \[ C = \begin{pmatrix}
                              1 & 1 & 1 & 1 \\
                              0 & 1 & 1 & 1 \\
                              0 & 0 & 1 & 1 \\
                              0 & 0 & 0 & 1
                          \end{pmatrix},\enspace C^{-1} = \begin{pmatrix}
                              1 & -1 & 0  & 0  \\
                              0 & 1  & -1 & 0  \\
                              0 & 0  & 1  & -1 \\
                              0 & 0  & 0  & 1
                          \end{pmatrix} \]
                      所以基 $ B_2 $ 变换为基 $ B_1 $ 的变换矩阵为 $ C^{-1} $.

                \item 在 $ A^2 = A $ 中选取较为简单的. 例如，由
                      \[ \begin{pmatrix} a & b \\ 0 & 0 \end{pmatrix}^2 = \begin{pmatrix} a^2 & ab \\ 0 & 0 \end{pmatrix} = \begin{pmatrix} a & b \\ 0 & 0 \end{pmatrix} \]
                      取 $ a = 1 $，$ b = 0 $ 或 $ b = 1 $ 得
                      \[ A_1 = \begin{pmatrix} 1 & 0 \\ 0 & 0 \end{pmatrix},\enspace A_2 = \begin{pmatrix} 1 & 1 \\ 0 & 0 \end{pmatrix} \]
                      由
                      \[ \begin{pmatrix} 0 & 0 \\ c & d \end{pmatrix}^2 = \begin{pmatrix} 0 & 0 \\ cd & d^2 \end{pmatrix} = \begin{pmatrix} 0 & 0 \\ c & d \end{pmatrix} \]
                      取 $ d = 1 $，$ c = 0 $ 或 $ c = 1 $ 得
                      \[ A_3 = \begin{pmatrix} 0 & 0 \\ 0 & 1 \end{pmatrix},\enspace A_4 = \begin{pmatrix} 0 & 0 \\ 1 & 1 \end{pmatrix} \]
                      上面的 $ A_1, A_2, A_3, A_4 $ 均满足 $ A_i^2 = A_i \enspace(i = 1, \ldots, 4) $，而且线性无关. 所以它们是 $ \mathbf{M}_2(\mathbf{R}) $ 的一组基 $ B_3 $.

                \item 先求基 $ B_2 $ 变换为基 $ B_3 $ 的变换矩阵 $ D $，即
                      \[ (A_1, A_2, A_3, A_4) = (\vec{g}_1, \vec{g}_2, \vec{g}_3, \vec{g}_4) D \]
                      由 \eqref{item:11:B:4:2} 所述，此时有
                      \[ \begin{pmatrix}
                              1 & 1 & 0 & 0 \\
                              0 & 1 & 0 & 0 \\
                              0 & 0 & 0 & 1 \\
                              0 & 0 & 1 & 1
                          \end{pmatrix} = CD \]
                      所以
                      \[ D = C^{-1} \begin{pmatrix}
                              1 & 1 & 0 & 0 \\
                              0 & 1 & 0 & 0 \\
                              0 & 0 & 0 & 1 \\
                              0 & 0 & 1 & 1
                          \end{pmatrix} = \begin{pmatrix}
                              1 & 0 & 0  & 0  \\
                              0 & 1 & 0  & -1 \\
                              0 & 0 & -1 & 0  \\
                              0 & 0 & 1  & 1
                          \end{pmatrix} \]
                      由于矩阵 $ A $ 关于基 $ B_2 $ 的坐标为 $ X = (1, 1, 1, 1)^{\mathrm{T}} $，所以 $ A $ 关于基 $ B_3 $ 的坐标为
                      \[ Y = D^{-1} X = \begin{pmatrix}
                              1 & 0 & 0  & 0  \\
                              0 & 1 & 0  & -1 \\
                              0 & 0 & -1 & 0  \\
                              0 & 0 & 1  & 1
                          \end{pmatrix}^{-1} \begin{pmatrix} 1 \\ 1 \\ 1 \\ 1 \end{pmatrix} = \begin{pmatrix}
                              1 & 0 & 0  & 0 \\
                              0 & 1 & 1  & 1 \\
                              0 & 0 & -1 & 0 \\
                              0 & 0 & 1  & 1
                          \end{pmatrix}^{-1} \begin{pmatrix} 1 \\ 1 \\ 1 \\ 1 \end{pmatrix} = \begin{pmatrix} 1 \\ 3 \\ -1 \\ 2 \end{pmatrix} \]
            \end{enumerate}
\end{solve}

\begin{exercise}[补充训练 L8-B-11]\label{prob:L8-B-11}
利用相抵标准形证明以下结论：
        \begin{enumerate}
            \item 设$B_1,B_2$为$s \times n$列满秩矩阵，证明：存在$s$阶可逆矩阵$C$使得$B_2=CB_1$；

            \item 设$B_1,B_2$为$s \times n$行满秩矩阵，证明：存在$n$阶可逆矩阵$C$使得$B_2=B_1C$；

            \item 任意秩为$r$的矩阵都可以被分解为$r$个秩为1的矩阵之和；

            \item 已知$A$是$n$阶方阵，证明：存在$n$阶方阵$B$使得$A=ABA,\enspace B=BAB$.
        \end{enumerate}
\end{exercise}
\noindent 来源：\href{https://github.com/yhwu-is/Linear-Algebra-Left-Undone/blob/new/%E8%AE%B2%E4%B9%89/%E4%B8%93%E9%A2%98/8%20%E7%9B%B8%E6%8A%B5%E6%A0%87%E5%87%86%E5%BD%A2.tex}{L8-B-11}。
\begin{solve}
\begin{enumerate}
                \item 初等变换即可.

                \item 同上.

                \item 矩阵 $A$ 秩为 $r$ 可写作 $A=P\begin{pmatrix}E_r & 0 \\ 0 & 0\end{pmatrix}Q = P(E_{11}+E_{22}+\cdots+E_{rr})Q$（$E_r$ 是 $r\times r$ 的单位矩阵，$E_{ii}$ 是 $n\times n$ 的只有第 $i$ 行 $i$ 列的这个元素为 1，其他元素均为 0 的矩阵）. 每个 $PE_{ii}Q$ 都是秩为 1 的矩阵，故得证.

                \item 记 $r(A)=r$，把 $A$ 写成 $P\begin{pmatrix}E_r & 0 \\ 0 & 0\end{pmatrix}Q$ 的形式. 构造 $B=Q^{-1}\begin{pmatrix}E_r & 0 \\ 0 & 0\end{pmatrix}P^{-1}$ 可以发现其满足条件，故得证.
            \end{enumerate}
\end{solve}

\begin{exercise}[补充训练 L8-B-12]\label{prob:L8-B-12}
设 $A \in \mathbf{M}_{m \times n}(\mathbf{F})$，$r(A)=r$，$k$ 是满足条件 $r \leqslant k \leqslant n$ 的任意整数，证明存在 $n$ 阶方阵 $B$，使得 $AB=O$，且 $r(A)+r(B)=k$.
\end{exercise}
\noindent 来源：\href{https://github.com/yhwu-is/Linear-Algebra-Left-Undone/blob/new/%E8%AE%B2%E4%B9%89/%E4%B8%93%E9%A2%98/8%20%E7%9B%B8%E6%8A%B5%E6%A0%87%E5%87%86%E5%BD%A2.tex}{L8-B-12}。
\begin{solve}
$r(A)=r$ 则 $AX=0$ 的解空间维数 $\dim N(A) = n-r$. 由 $r(A)+r(B)=k$ 得 $r(B)=k-r \leqslant n-r=\dim N(A)$. 要求 $AB=O$，说明 $B$ 的列向量均为 $AX=0$ 的解，那么只需要选择合适的列向量组拼接成 $B$ 即可（这一定能做到，因为 $B$ 维数不会超过解空间维数）.
\end{solve}

\begin{exercise}[补充训练 L8-B-13]\label{prob:L8-B-13}
设$A$是$m \times n$矩阵（$m \leqslant n$），$r(A)=m$，证明：存在$n \times m$矩阵$B$使得$AB=E$.
\end{exercise}
\noindent 来源：\href{https://github.com/yhwu-is/Linear-Algebra-Left-Undone/blob/new/%E8%AE%B2%E4%B9%89/%E4%B8%93%E9%A2%98/8%20%E7%9B%B8%E6%8A%B5%E6%A0%87%E5%87%86%E5%BD%A2.tex}{L8-B-13}。
\begin{solve}
由于 $A$ 是 $m\times n$ 矩阵，$r(A)=m$，可知对于矩阵 $A$ 做初等列变换，可使其前 $m$ 列变为单位矩阵，后 $n-m$ 列变为全 0 列. 因此，存在 $n$ 阶可逆矩阵 $P$ 使得
          \[AP=\begin{pmatrix}E_m & O_{m\times (n-m)}\end{pmatrix}\]
          于是\[AP(AP)^{\mathrm{T}} = \begin{pmatrix}E & O\end{pmatrix} \begin{pmatrix}E \\ O\end{pmatrix}=E_m\]
          所以存在 $B=(PP^{\mathrm{T}}A^{\mathrm{T}})$ 为 $n\times m$ 矩阵，使 $AB=E$.
\end{solve}

\begin{exercise}[补充训练 L8-B-14]\label{prob:L8-B-14}
设$A,B \in \mathbf{M}_n(\mathbf{F})$，$r(A)+r(B) \leqslant n$，证明：存在可逆矩阵$M$，使得$AMB=O$.
\end{exercise}
\noindent 来源：\href{https://github.com/yhwu-is/Linear-Algebra-Left-Undone/blob/new/%E8%AE%B2%E4%B9%89/%E4%B8%93%E9%A2%98/8%20%E7%9B%B8%E6%8A%B5%E6%A0%87%E5%87%86%E5%BD%A2.tex}{L8-B-14}。
\begin{solve}
利用 $A,B$ 的相抵标准形. 存在 $n$ 阶可逆矩阵 $P_1,Q_1,P_2,Q_2$ 使得
          \[P_1AQ_1=\begin{pmatrix}E_{r_A} & O \\ O & O\end{pmatrix},P_2BQ_2=\begin{pmatrix}O & O \\ O & E_{r_B}\end{pmatrix}\]
          于是 \[AQ_1=P_1^{-1}\begin{pmatrix}E_{r_A} & O \\ O & O\end{pmatrix},P_2B=\begin{pmatrix}O & O \\ O & E_{r_B}\end{pmatrix}Q_2^{-1}\]
          所以 \[AQ_1P_2B=P_1^{-1}\begin{pmatrix}E_{r_A} & O \\ O & O\end{pmatrix}\begin{pmatrix}O & O \\ O & E_{r_B}\end{pmatrix}Q_2^{-1}=O\]
          取 $M=Q_1P_2$ 即可.
\end{solve}

\begin{exercise}[补充训练 L8-B-16]\label{prob:L8-B-16}
\begin{enumerate}
            \item 设 $A$ 和 $B$ 是两个 $m \times n$矩阵，证明：存在可逆矩阵 $M$ 使得$B = AM$ 的充要条件是 $A$ 和 $B$ 的列向量可以互相表示；
            \item 设 $A$ 和 $B$ 是两个 $m \times n$ 矩阵，且存在两个方阵 $C, D$ 使得 $A = BC, B = AD$，证明：存在可逆矩阵 $M$ 使得 $B = AM$.
        \end{enumerate}
\end{exercise}
\noindent 来源：\href{https://github.com/yhwu-is/Linear-Algebra-Left-Undone/blob/new/%E8%AE%B2%E4%B9%89/%E4%B8%93%E9%A2%98/8%20%E7%9B%B8%E6%8A%B5%E6%A0%87%E5%87%86%E5%BD%A2.tex}{L8-B-16}。
\begin{solve}
\begin{enumerate}
                \item 必要性：$B = AM$，则 $B$ 的每一列都为 $A$ 的列向量的线性组合，所以 $B$ 的列向量可以由 $A$ 的列向量表示. 反之，由于 $M$ 可逆，所以 $A = BM^{-1}$，则 $A$ 的每一列都为 $B$ 的列向量的线性组合，所以 $A$ 的列向量可以由 $B$ 的列向量表示.

                充分性：由于 $A$ 与 $B$ 的列向量组等价，则 $A$ 与 $B$ 的列空间相同，令其维数为 $r$. 对 $A, B$ 分别作初等列变换，存在可逆矩阵 $Q_1, Q_2 \in \mathbf{R}^{n\times n}$，使得
                \[ AQ_1 = (A_1, 0), \qquad BQ_2 = (B_1, 0), \]
                其中 $A_1, B_1 \in \mathbf{R}^{m\times r}$ 均列满秩。又因 $A$ 与 $B$ 的列空间相同，故 $A_1$ 与 $B_1$ 是同一列空间的两组基，因此存在可逆矩阵 $M_1 \in \mathbf{R}^{r\times r}$，使得 $B_1 = A_1M_1$。

                于是
                \[ BQ_2=(B_1,0) = (A_1M_1,0) = (A_1,0) \begin{pmatrix}
                    M_1 & 0 \\ 0 & I_{n-r}
                \end{pmatrix} = AQ_1 \begin{pmatrix}
                    M_1 & 0 \\ 0 & I_{n-r}
                \end{pmatrix}. \]
                因此
                \[ B = AQ_1 \begin{pmatrix}
                    M_1 & 0\\ 0 & I_{n-r}
                \end{pmatrix} Q_2^{-1}. \]
                令
                \[ M = Q_1 \begin{pmatrix}
                    M_1 & 0 \\ 0 & I_{n-r}
                    \end{pmatrix} Q_2^{-1}, \]
                则 $M$ 可逆，且 $B = AM$。

                \item 由于 $A = BC$，故 $A$ 的每一列都为 $B$ 的列向量的线性组合，同样由于$B = AD$，故 $B$ 的每一列都为 $A$ 的列向量的线性组合，所以 $A$ 与 $B$ 的列向量可以互相表示，因此它们的过渡矩阵可逆.
            \end{enumerate}
\end{solve}

\begin{exercise}[补充训练 L8-C-1]\label{prob:L8-C-1}
设矩阵$A \in \mathbf{F}^{m \times n}$，$A$的秩$r(A)=r$，定义$\mathbf{F}^{n \times p}$到$\mathbf{F}^{m \times p}$的线性映射$\sigma$，使得$\forall X \in \mathbf{F}^{n \times p}$，$\sigma(X)=AX$. 求$\sigma$核空间的维数.
\end{exercise}
\noindent 来源：\href{https://github.com/yhwu-is/Linear-Algebra-Left-Undone/blob/new/%E8%AE%B2%E4%B9%89/%E4%B8%93%E9%A2%98/8%20%E7%9B%B8%E6%8A%B5%E6%A0%87%E5%87%86%E5%BD%A2.tex}{L8-C-1}。
\begin{solve}
本题直接求核空间困难，但由于我们只需求维数，所以我们转而求像空间维数并通过维数公式求核空间维数.

          取 $ \mathbf{F}^{n \times p} $ 的自然基 $ \vec{e}_{ij} = (e_{kl})_{n \times p} $，其中 $ e_{kl} = \delta_{ik} \delta_{jl} $，即第 $ i $ 行第 $ j $ 列元素为 $ 1 $，其余元素为 $ 0 $. 则 $ \operatorname{im} \sigma = \operatorname{span}(\sigma(\vec{e}_{11}), \ldots, \sigma(\vec{e}_{np})) $.

          取 $ A $ 的列向量组 $ A = (\alpha_1, \alpha_2, \ldots, \alpha_n) $，则可将 $ \sigma(\vec{e}_{ij}) $ 排列如下：
          \[ \begin{matrix}
                  (a_1, 0, \ldots, 0) & (0, a_1, \ldots, 0) & \cdots & (0, 0, \ldots, a_1) \\
                  (a_2, 0, \ldots, 0) & (0, a_2, \ldots, 0) & \cdots & (0, 0, \ldots, a_2) \\
                  \vdots              & \vdots              & \ddots & \vdots              \\
                  (a_n, 0, \ldots, 0) & (0, a_n, \ldots, 0) & \cdots & (0, 0, \ldots, a_n)
              \end{matrix} \]
          记 $ r(A) = r $，则 $ \alpha_1, \alpha_2, \ldots, \alpha_n $ 的极大线性无关组有 $ r $ 个向量，不妨设为 $ \alpha_1, \alpha_2, \ldots, \alpha_r $. 那么下列向量：
          \[ \begin{matrix}
                  (a_{r+1}, 0, \ldots, 0) & (0, a_{r+1}, \ldots, 0) & \cdots & (0, 0, \ldots, a_{r+1}) \\
                  \vdots                  & \vdots                  & \ddots & \vdots                  \\
                  (a_n, 0, \ldots, 0)     & (0, a_n, \ldots, 0)     & \cdots & (0, 0, \ldots, a_n)
              \end{matrix} \]

          均可以被其他向量线性表出. 观察除了上述向量的剩下的向量，可以发现这 $ rp $ 个向量线性无关，它们是 $ \operatorname{im} \sigma $ 的一组基，因此 $ \dim \operatorname{im} \sigma = rp $. 由维数公式，$ \dim \ker \sigma = \dim \mathbf{F}^{n \times p} - \dim \operatorname{im} \sigma = (n - r)p $.
\end{solve}

\begin{exercise}[补充训练 L8-例-2]\label{prob:L8-例-2}
已知$\beta_1=\alpha_2+\alpha_3,\enspace\beta_2=\alpha_1+\alpha_3,\enspace\beta_3=\alpha_1+\alpha_2$，证明$\alpha_1,\alpha_2,\alpha_3$与$\beta_1,\beta_2,\beta_3$等价.
\end{exercise}
\noindent 来源：\href{https://github.com/yhwu-is/Linear-Algebra-Left-Undone/blob/new/%E8%AE%B2%E4%B9%89/%E4%B8%93%E9%A2%98/8%20%E7%9B%B8%E6%8A%B5%E6%A0%87%E5%87%86%E5%BD%A2.tex}{L8-例-2}。
\begin{solve}
事实上，$(\beta_1,\beta_2,\beta_3)=(\alpha_1,\alpha_2,\alpha_3)\begin{pmatrix}
            0 & 1 & 1 \\
            1 & 0 & 1 \\
            1 & 1 & 0
        \end{pmatrix}$，而$\begin{pmatrix}
            0 & 1 & 1 \\
            1 & 0 & 1 \\
            1 & 1 & 0
        \end{pmatrix}$可逆（可以由原讲义“可逆等价条件”处（3）（4）很容易验证），故$\alpha_1,\alpha_2,\alpha_3$与$\beta_1,\beta_2,\beta_3$等价.
\end{solve}

\begin{exercise}[补充训练 L8-例-3]\label{prob:L8-例-3}
已知三维线性空间 $V$ 的线性变换 $\sigma$ 关于基 $\alpha_1,\alpha_2,\alpha_3$ 所对应的矩阵为
    \[\begin{pmatrix}1 & 2 & -1 \\ 2 & 1 & 0 \\ 3 & 0 & 1\end{pmatrix}\]
    \begin{enumerate}
        \item 求 $\sigma$ 在基 $\beta_1,\beta_2,\beta_3$ 下对应的矩阵 $B$，其中：
              \[\beta_1=2\alpha_1+\alpha_2+3\alpha_3,\enspace \beta_2=\alpha_1+\alpha_2+2\alpha_3,\enspace \beta_3=-\alpha_1+\alpha_2+\alpha_3;\]

        \item 求 $\sigma$ 的值域 $\sigma(V)$ 和核 $\ker\sigma$；

        \item 把 $\sigma(V)$ 的基扩充为 $V$ 的基，并求 $\sigma$ 在这组基下对应的矩阵；

        \item 把 $\ker\sigma$ 的基扩充为 $V$ 的基，并求 $\sigma$ 在这组基下对应的矩阵.
    \end{enumerate}
\end{exercise}
\noindent 来源：\href{https://github.com/yhwu-is/Linear-Algebra-Left-Undone/blob/new/%E8%AE%B2%E4%B9%89/%E4%B8%93%E9%A2%98/8%20%E7%9B%B8%E6%8A%B5%E6%A0%87%E5%87%86%E5%BD%A2.tex}{L8-例-3}。
\begin{solve}
\begin{enumerate}
        \item 对新的一组基，使用过渡矩阵进行表达如下：
              \[(\beta_{1}, \beta_{2}, \beta_{3})=(\alpha_{1}, \alpha_{2}, \alpha_{3})\begin{pmatrix}
                      2 & 1 & -1 \\
                      1 & 1 & 1  \\
                      3 & 2 & 1
                  \end{pmatrix}=(\alpha_{1}, \alpha_{2}, \alpha_{3}) C,\]
              其中 $C$ 是可逆矩阵，且
              \[(\alpha_{1},\alpha_{2},\alpha_{3})=(\beta_{1},\beta_{2},\beta_{3}) C^{-1},\]
              使用原讲义“基的选择对变换矩阵的影响”处可以直接得到 $B=C^{-1} A C$，或者根据原讲义“基与矩阵的"结合律"”处，将上式代入已知条件得
              \[\sigma((\beta_{1}, \beta_{2}, \beta_{3}) C^{-1})=((\beta_{1}, \beta_{2}, \beta_{3}) C^{-1}) A,\]
              \[(\sigma(\beta_{1}, \beta_{2}, \beta_{3})) C^{-1}=(\beta_{1}, \beta_{2}, \beta_{3})(C^{-1} A),\]
              从而得 $\sigma(\beta_{1}, \beta_{2}, \beta_{3})=(\beta_{1}, \beta_{2}, \beta_{3})(C^{-1} A C)$，故 $\sigma$ 关于基 $\beta_{1}, \beta_{2}, \beta_{3}$ 下对应的矩阵为
              \[B=C^{-1} A C = \begin{pmatrix}
                      2 & 1 & -1 \\
                      1 & 1 & 1  \\
                      3 & 2 & 1
                  \end{pmatrix}^{-1}\begin{pmatrix}
                      1 & 2 & -1 \\
                      2 & 1 & 0  \\
                      3 & 0 & 1
                  \end{pmatrix}\begin{pmatrix}
                      2 & 1 & -1 \\
                      1 & 1 & 1  \\
                      3 & 2 & 1
                  \end{pmatrix}=\begin{pmatrix}
                      2 & 0 & 1  \\
                      0 & 2 & 1  \\
                      3 & 1 & -1
                  \end{pmatrix}.\]

        \item  $\sigma$的值域是 $A$ 列向量组的极大线性无关组，由于 $ A $ 的第 1 列可以由第 2 列和第 3 列线性表示，从而 $\sigma(V)=\operatorname{span}(2 \alpha_{1}+\alpha_{2}, -\alpha_{1}+\alpha_{3})$.

        \item 由于 $\alpha_{1}$ 不能由 $2 \alpha_{1}+\alpha_{2}$ 和 $-\alpha_{1}+\alpha_{3}$ 线性表示，可以把 $\sigma(V)$ 的基扩充为 $V$ 的基 $\{\alpha_{1}, 2 \alpha_{1}+\alpha_{2}, -\alpha_{1}+\alpha_{3}\}$，$\sigma$ 在这个基下对应的矩阵是
              \[\begin{pmatrix}
                      1 & 2 & -1 \\
                      0 & 1 & 0  \\
                      0 & 0 & 1
                  \end{pmatrix}^{-1}\begin{pmatrix}
                      1 & 2 & -1 \\
                      2 & 1 & 0  \\
                      3 & 0 & 1
                  \end{pmatrix}\begin{pmatrix}
                      1 & 2 & -1 \\
                      0 & 1 & 0  \\
                      0 & 0 & 1
                  \end{pmatrix}=\begin{pmatrix}
                      0 & 0 & 0  \\
                      2 & 5 & -2 \\
                      3 & 6 & -2
                  \end{pmatrix}.\]

        \item 由于 $\alpha_{1}, \alpha_{2}$ 不能由 $\alpha_{1}-2 \alpha_{2}-3 \alpha_{3}$ 线性表示，可以把 $\ker \sigma$ 的基扩 充为 $V$ 的基 $\{\alpha_{1}, \alpha_{2}, \alpha_{1}-2 \alpha_{2}-3 \alpha_{3}\}$，$\sigma$ 在这个基下对应的矩阵是
              \[\begin{pmatrix}
                      1 & 0 & 1  \\
                      0 & 1 & -2 \\
                      0 & 0 & -3
                  \end{pmatrix}^{-1}\begin{pmatrix}
                      1 & 2 & -1 \\
                      2 & 1 & 0  \\
                      3 & 0 & 1
                  \end{pmatrix}\begin{pmatrix}
                      1 & 0 & 1  \\
                      0 & 1 & -2 \\
                      0 & 0 & -3
                  \end{pmatrix}=\begin{pmatrix}
                      2 & 2 & 0 \\
                      0 & 1 & 0 \\
                      1 & 0 & 0
                  \end{pmatrix}.\]
    \end{enumerate}
\end{solve}

\begin{exercise}[补充训练 L8-例-4]\label{prob:L8-例-4}
已知矩阵 $A=\begin{pmatrix}a & b & c \\ d & e & f \\ h & x & y\end{pmatrix}$ 的逆是 $A^{-1}=\begin{pmatrix}-1 & -2 & -1 \\ 2 & 1 & 0 \\ 0 & -3 & -1\end{pmatrix}$，\\
    $B=\begin{pmatrix}a-2b & b-3c & -c \\ d-2e & e-3f & -f \\ h-2x & x-3y & -y\end{pmatrix}$. 求矩阵 $X$ 满足：

    \[X+\left(B(A^\mathrm{T}B^2)^{-1}A^\mathrm{T}\right)^{-1}=X\left(A^2(B^\mathrm{T}A)^{-1}B^\mathrm{T}\right)^{-1}(A+B)\]
\end{exercise}
\noindent 来源：\href{https://github.com/yhwu-is/Linear-Algebra-Left-Undone/blob/new/%E8%AE%B2%E4%B9%89/%E4%B8%93%E9%A2%98/8%20%E7%9B%B8%E6%8A%B5%E6%A0%87%E5%87%86%E5%BD%A2.tex}{L8-例-4}。
\begin{solve}
注意到
    \[B=A\begin{pmatrix}
            1 &  & \\ -2 & 1 & \\ & & 1
        \end{pmatrix} \begin{pmatrix}
            1 &  & \\ & 1 & \\ & -3 & 1
        \end{pmatrix} \begin{pmatrix}
            1 &  & \\ & 1 & \\ & & -1
        \end{pmatrix},\]
    从而 $B$ 由 $A$ 经过有限次初等变换得到，因 $A$ 可逆，故 $B$ 可逆. 一方面，等式左侧可以化简为
    \[X + \left(B\left(A^\mathrm{T} B^{2}\right)^{-1} A^\mathrm{T}\right)^{-1} = X + B,\]
    另一方面，等式右侧可以化简为
    \[X\left(A^{2}\left(B^\mathrm{T} A\right)^{-1} B^\mathrm{T}\right)^{-1}(A + B) = X A^{-1}(A + B) = X + X A^{-1} B,\]
    于是 
\[X = A = (A^{-1})^{-1} = \begin{pmatrix}
            -\dfrac{1}{3} & \dfrac{1}{3} & \dfrac{1}{3}  \\[0.5em]
            \dfrac{2}{3}  & \dfrac{1}{3} & -\dfrac{2}{3} \\
            -2           & -1          & 1
        \end{pmatrix}\]
.
\end{solve}

\begin{exercise}[补充训练 L8-例-6]\label{prob:L8-例-6}
设$A=\begin{pmatrix}
            1 & 0 & 2 & -4 \\ 2 & 1 & 3 & -6 \\ -1 & -1 & -1 & 2
        \end{pmatrix}$. 求
    \begin{enumerate}[label=(\arabic*)]
        \item $A$的秩$r$和相抵标准形；

        \item
              3 阶可逆矩阵$P$和 4 阶可逆矩阵$Q$使得$PAQ=\begin{pmatrix}
                      E_r & 0 \\ 0 & 0
                  \end{pmatrix}$.
    \end{enumerate}
\end{exercise}
\noindent 来源：\href{https://github.com/yhwu-is/Linear-Algebra-Left-Undone/blob/new/%E8%AE%B2%E4%B9%89/%E4%B8%93%E9%A2%98/8%20%E7%9B%B8%E6%8A%B5%E6%A0%87%E5%87%86%E5%BD%A2.tex}{L8-例-6}。
\begin{solve}
\begin{enumerate}
        \item 利用求解极大线性无关组的方法可以解得$r=2$，因此对应的相抵标准形为
              \[\begin{pmatrix}
                      1 & 0 & 0 & 0 \\ 0 & 1 & 0 & 0 \\ 0 & 0 & 0 & 0
                  \end{pmatrix}.\]
        \item 这里介绍两种方法（严格来说，两种方法是一样的）.

              法一：直接将所有行变换相乘得到$P$，所有列变换相乘得到$Q$，此处略去步骤，给出答案为
              \[P=\begin{pmatrix}
                      1 & 0 & 0 \\ -2 & 1 & 0 \\ -1 & 1 & 1
                  \end{pmatrix},\enspace Q=\begin{pmatrix}
                      1 & 0 & -2 & 4 \\ 0 & 1 & 1 & -2 \\ 0 & 0 & 1  & 0 \\ 0 & 0 & 0 & 1
                  \end{pmatrix}.\]

              法二：我们用原讲义“初等变换求逆矩阵”处中的思想，将原讲义“相抵标准形另证”中的过程具象化. 具体来说，对 $m\times n$ 阶矩阵 $A$：
              \begin{gather}
                  \left(\begin{array}{c:c}
                      A & E_m
                  \end{array}\right)\xrightarrow{\text{初等行变换}}\left(\begin{array}{c:c}
                      A_1 & P
                  \end{array}\right)\\
                  \left(\begin{array}{c}
                      A_1 \\ \hdashline E_n
                  \end{array}\right)\xrightarrow{\text{初等列变换}}\left(\begin{array}{c}
                      U_r \\ \hdashline Q
                  \end{array}\right)
              \end{gather}
              其中 $A_1$ 为（行）简化阶梯矩阵，$P = P_s \cdots P_1$，$Q = Q_1 \cdots Q_t$为所求矩阵.

              以本题为例：
              \begin{gather*}
                  (A, E) =
                  \left(\begin{array}{cccc:cccc}
                          1  & 0  & 2  & -4 & 1 & 0 & 0 \\
                          2  & 1  & 3  & -6 & 0 & 1 & 0 \\
                          -1 & -1 & -1 & 2  & 0 & 0 & 1
                      \end{array}\right) \\
                  \xrightarrow{[2] + [1] \times (-2)}
                  \left(\begin{array}{cccc:cccc}
                          1  & 0  & 2  & -4 & 1  & 0 & 0 \\
                          0  & 1  & -1 & 2  & -2 & 1 & 0 \\
                          -1 & -1 & -1 & 2  & 0  & 0 & 1
                      \end{array}\right) \\
                  \xrightarrow{\substack{[3] + [1] \\ [3] + [2]}}
                  \left(\begin{array}{cccc:cccc}
                          1 & 0 & 2  & -4 & 1  & 0 & 0 \\
                          0 & 1 & -1 & 2  & -2 & 1 & 0 \\
                          0 & 0 & 0  & 0  & -1 & 1 & 1
                      \end{array}\right) \\
                  \left(\begin{array}{c}
                      A_1 \\ \hdashline E_n
                  \end{array}\right) =
                  \left(\begin{array}{cccc}
                      1 & 0 & 2  & -4 \\
                      0 & 1 & -1 & 2  \\
                      0 & 0 & 0  & 0  \\ \hdashline
                      1 & 0 & 0  & 0  \\
                      0 & 1 & 0  & 0  \\
                      0 & 0 & 1  & 0  \\
                      0 & 0 & 0  & 1
                  \end{array}\right)
                  \xrightarrow{\substack{[3] + [1] \times (-2) \\ [3] + [2] \\ [4] + [1] \times 4 \\ [4] + [2] \times (-2)}}
                  \left(\begin{array}{cccc}
                      1 & 0 & 0  & 0  \\
                      0 & 1 & 0  & 0  \\
                      0 & 0 & 0  & 0  \\ \hdashline
                      1 & 0 & -2 & 4  \\
                      0 & 1 & 1  & -2 \\
                      0 & 0 & 1  & 0  \\
                      0 & 0 & 0  & 1
                  \end{array}\right)
              \end{gather*}
              从而可以得到 $P$，$Q$，结果同上.
    \end{enumerate}
\end{solve}

\begin{exercise}[补充训练 L8-例-7]\label{prob:L8-例-7}
设$\sigma:V\to W$为线性映射，取$V$和$W$的基$\alpha_1,\alpha_2,\ldots,\alpha_n$和$\beta_1,\beta_2,\ldots,\beta_m$，且$\sigma$在这两组基下的矩阵表示为$A=(a_{ij})_{m\times n}$.
    \begin{enumerate}
        \item 若将$\alpha_i$换为$c\alpha_i$（$c$为常数），求$\sigma$在新基下的矩阵表示；

        \item 若将$\beta_i$换为$c\beta_i$（$c$为常数），求$\sigma$在新基下的矩阵表示；

        \item 若将$\alpha_i$换为$\alpha_i+k\alpha_j$（$k$为常数），求$\sigma$在新基下的矩阵表示；

        \item 若将$\beta_i$换为$\beta_i+k\beta_j$（$k$为常数），求$\sigma$在新基下的矩阵表示；

        \item 若将$\alpha_i$和$\alpha_j$对换，求$\sigma$在新基下的矩阵表示；

        \item 若将$\beta_i$和$\beta_j$对换，求$\sigma$在新基下的矩阵表示.
    \end{enumerate}
\end{exercise}
\noindent 来源：\href{https://github.com/yhwu-is/Linear-Algebra-Left-Undone/blob/new/%E8%AE%B2%E4%B9%89/%E4%B8%93%E9%A2%98/8%20%E7%9B%B8%E6%8A%B5%E6%A0%87%E5%87%86%E5%BD%A2.tex}{L8-例-7}。
\begin{solve}
直接根据原讲义“换基公式”处，写出新旧基之间的过渡矩阵. 我们很容易发现这些过渡矩阵实际上就对应于六种初等矩阵，结果如下：
    \begin{enumerate}
        \item 矩阵表示$A_1$就是$A$的第$i$列数乘$c$；

        \item 矩阵表示$A_2$就是$A$的第$i$行数乘$1/c$；

        \item 矩阵表示$A_3$就是$A$的第$i$列加上第$j$列数乘$k$；

        \item 矩阵表示$A_4$就是$A$的第$j$行减去第$i$行数乘$k$；

        \item 矩阵表示$A_5$就是$A$的第$i$列和第$j$列对换；

        \item 矩阵表示$A_6$就是$A$的第$i$行和第$j$行对换.
    \end{enumerate}
\end{solve}

\begin{exercise}[补充训练 L8-例-8]\label{prob:L8-例-8}
设$A$为$n$阶方阵，证明：
    \begin{enumerate}
        \item 存在可逆矩阵$B$和幂等矩阵$C$（即满足$C^2=C$）使得$A=BC$；
        \item 存在对称矩阵$B$和可逆矩阵$C$使得$A=BC$.
    \end{enumerate}
\end{exercise}
\noindent 来源：\href{https://github.com/yhwu-is/Linear-Algebra-Left-Undone/blob/new/%E8%AE%B2%E4%B9%89/%E4%B8%93%E9%A2%98/8%20%E7%9B%B8%E6%8A%B5%E6%A0%87%E5%87%86%E5%BD%A2.tex}{L8-例-8}。
\begin{solve}
根据相抵标准形分解，设$r(A)=r$，则我们有$A=P\begin{pmatrix}
            E_r & O \\ O & O
        \end{pmatrix}Q$，其中$P$和$Q$都是可逆矩阵.
    \begin{enumerate}
        \item 第一问难点在于取出这个幂等矩阵 $C$，实际上，我们可以取$C=Q^{-1}\begin{pmatrix}
                      E_r & O \\ O & O
                  \end{pmatrix}Q$，则$C^2=C$（因为中间的$QQ^{-1}$会抵消，而相抵标准形是幂等的）. 为了使得 $A = BC$，令 $B = PQ$ 即可（因为$P$ 和 $Q$ 都可逆，其乘积也必定可逆）.

        \item 第二问难点在于取出这个对称矩阵，但相信读者在经过上一小问的历练后，应该会产生一种取到对称矩阵的直觉. 因为标准形是对称的，我们只需要再多写一个$P^\mathrm{T}$，令$B=P\begin{pmatrix}
                      E_r & O \\ O & O
                  \end{pmatrix}P^\mathrm{T}$，则$B$是对称矩阵，且$A=BC$，令$C=(P^\mathrm{T})^{-1}Q$即可（注意矩阵可逆则转置也可逆）.
    \end{enumerate}
\end{solve}

\begin{exercise}[补充训练 S-PARAM-1]\label{prob:S-PARAM-1}
设 $P=\begin{pmatrix}1&-3&0\\0&1&0\\1&0&1\end{pmatrix}$，$A(t)=P\operatorname{diag}(1,t-1,t-3)P^{-1}$，其中 $t\in\mathbf R$。求 $A(t)$ 的行列式和秩；分别求奇异参数下的核；判断方程 $A(t)x=\begin{pmatrix}-2\\1\\2\end{pmatrix}$ 何时有解。
\end{exercise}
\noindent 编者补充：S-PARAM-1；相关主题见原讲义对应讲，题目数据与解答由本选编另行构造。
\begin{solve}
相似变换保持行列式和秩，故
\[\det A(t)=(t-1)(t-3),\qquad r(A(t))=\begin{cases}3,&t\ne1,3,\\2,&t=1\text{ 或 }t=3.\end{cases}\]
当 $t=1$ 时，核由 $P$ 的第二列 $p_2=\begin{pmatrix}-3\\1\\0\end{pmatrix}$ 张成；当 $t=3$ 时，核由第三列 $p_3=\begin{pmatrix}0\\0\\1\end{pmatrix}$ 张成。令 $y=P^{-1}x$，题给右端在这个基下的坐标是 $(1,1,1)^{\mathrm T}$，方程成为
\[y_1=1,\qquad(t-1)y_2=1,\qquad(t-3)y_3=1.\]
所以奇异参数下方程无解，其余参数下唯一解为
\[x=P\begin{pmatrix}1\\1/(t-1)\\1/(t-3)\end{pmatrix}.\]
在讨论含参消元时必须把可能为零的主元单独分类；上面的基变换直接显示了所有特殊值。
\end{solve}

\section{补充训练：分块矩阵与矩阵运算进阶}

本节包含 40 道训练题，建议在完成相关专题真题后选择性练习。

\begin{exercise}[补充训练 L9-A-1]\label{prob:L9-A-1}
设方阵$A$满足$A^2-A-2E=O$，证明：
        \begin{enumerate}
            \item $A$和$E-A$都是可逆矩阵，并求它们的逆矩阵；

            \item $A+E$和$A-2E$不可能同时可逆.
        \end{enumerate}
\end{exercise}
\noindent 来源：\href{https://github.com/yhwu-is/Linear-Algebra-Left-Undone/blob/new/%E8%AE%B2%E4%B9%89/%E4%B8%93%E9%A2%98/9%20%E7%9F%A9%E9%98%B5%E8%BF%90%E7%AE%97%E8%BF%9B%E9%98%B6.tex}{L9-A-1}。
\begin{solve}
\begin{enumerate}
                \item 若 $ AB = kE \enspace(k \neq 0) $，则 $ A^{-1} = \dfrac{1}{k} B $. 由
                      \[ A^2 - A - 2E = A(A - E) - 2E = O \]
                      可得
                      \[ A(A - E) = 2E \text{~或~} A(E - A) = -2E. \]
                      所以
                      \[ A^{-1} = \frac{1}{2}(A-E),\enspace (E - A)^{-1} = -\frac{1}{2}A. \]

                \item 由
                      \[  A^2 - A - 2E = (A - 2E)(A + E) = O \]
                      可知，$ A + E $ 和 $ A - 2E $ 不能同时可逆，否则 $ A^2 - A - 2E $ 为零矩阵可逆，矛盾.
            \end{enumerate}
\end{solve}

\begin{exercise}[补充训练 L9-A-3]\label{prob:L9-A-3}
若$A,B$为两个$n$阶矩阵且满足$A+B=AB$，证明：
        \begin{enumerate}
            \item $A-E$和$B-E$均可逆；

            \item $AB=BA$；

            \item $r(A)=r(B)$.
        \end{enumerate}
\end{exercise}
\noindent 来源：\href{https://github.com/yhwu-is/Linear-Algebra-Left-Undone/blob/new/%E8%AE%B2%E4%B9%89/%E4%B8%93%E9%A2%98/9%20%E7%9F%A9%E9%98%B5%E8%BF%90%E7%AE%97%E8%BF%9B%E9%98%B6.tex}{L9-A-3}。
\begin{solve}
\begin{enumerate}
                \item \begin{align*}
                          (A - E)(B - E) & = AB - AE - EB + E^2 \\
                                         & = AB - A - B + E     \\
                                         & = AB - (A + B) + E   \\
                                         & = AB - AB + E        \\
                                         & = E.
                      \end{align*}
                      所以 $ A - E $ 与 $ B - E $ 互为逆矩阵.

                \item 由于 $ A - E $ 与 $ B - E $ 互为逆矩阵，所以
                      \begin{align*}
                          (B - E)(A - E) & = E              \\
                                         & = BA - B - A + E \\
                                         & = BA - AB + E.
                      \end{align*}
                      所以 $ BA - AB = O $，即 $ AB = BA $.

                \item 由 $ A = AB - B = B(A - E) $ 可得 $ r(A) \leqslant r(B) $，同理可得 $ r(B) \leqslant r(A) $，所以 $ r(A) = r(B) $.
            \end{enumerate}
\end{solve}

\begin{exercise}[补充训练 L9-A-4]\label{prob:L9-A-4}
求下列矩阵的可逆的条件与逆矩阵：
            \begin{enumerate}
                \item $\begin{pmatrix}
                        A & B \\ O & D
                    \end{pmatrix}$，其中 $A$ 是 $m$ 阶方阵，$D$ 是 $n$ 阶方阵，$B$ 是 $m \times n$ 矩阵；

                \item $\begin{pmatrix}
                        O & B \\ C & D
                    \end{pmatrix}$，其中 $B$ 是 $n$ 阶方阵，$C$ 是 $m$ 阶方阵，$D$ 是 $n \times m$ 矩阵.
            \end{enumerate}
\end{exercise}
\noindent 来源：\href{https://github.com/yhwu-is/Linear-Algebra-Left-Undone/blob/new/%E8%AE%B2%E4%B9%89/%E4%B8%93%E9%A2%98/9%20%E7%9F%A9%E9%98%B5%E8%BF%90%E7%AE%97%E8%BF%9B%E9%98%B6.tex}{L9-A-4}。
\begin{solve}
\begin{enumerate}
                \item 对于矩阵 $\begin{pmatrix}
                        A & B \\ O & D
                    \end{pmatrix}$，

                    若 $A$ 不可逆，则 $r\begin{pmatrix} A \\ O \end{pmatrix} = r(A) < m$，
                    故 $r\begin{pmatrix} A & B \\ O & D \end{pmatrix} < m+n$，原矩阵不可逆.

                    若 $D$ 不可逆，则 $r\begin{pmatrix} O & D \end{pmatrix} = r(D) < n$，
                    故 $r\begin{pmatrix} A & B \\ O & D \end{pmatrix} < m+n$，原矩阵不可逆.

                    若 $A$ 与 $D$ 均可逆，利用分块矩阵初等变换，我们有
                    \[
                        \begin{pmatrix}
                            A^{-1} & O \\ O & D^{-1}
                        \end{pmatrix} \begin{pmatrix}
                            E_m & -BD^{-1} \\ O & E_n
                        \end{pmatrix} \begin{pmatrix}
                            A & B \\ O & D
                        \end{pmatrix} = E_{m+n},
                    \]
                    故有
                    \[
                        \begin{pmatrix}
                            A & B \\ O & D
                        \end{pmatrix}^{-1} = \begin{pmatrix}
                            A^{-1} & O \\ O & D^{-1}
                        \end{pmatrix} \begin{pmatrix}
                            E_m & -BD^{-1} \\ O & E_n
                        \end{pmatrix} = \begin{pmatrix}
                            A^{-1} & -A^{-1}BD^{-1} \\ O & D^{-1}
                        \end{pmatrix}.
                    \]

                \item 对于矩阵 $\begin{pmatrix}
                        O & B \\ C & D
                    \end{pmatrix}$，

                    由于
                    \[
                        \begin{pmatrix}
                            O & E_m \\ E_n & O
                        \end{pmatrix} \begin{pmatrix}
                            O & B \\ C & D
                        \end{pmatrix} = \begin{pmatrix}
                            C & D \\ O & B
                        \end{pmatrix},
                    \]

                    而 $\begin{pmatrix} O & E_m \\ E_n & O \end{pmatrix}$ 是可逆矩阵，故
                    \[
                        \begin{pmatrix}
                            O & B \\ C & D
                        \end{pmatrix} \, \text{可逆}
                        \iff \begin{pmatrix}
                            C & D \\ O & B
                        \end{pmatrix} \, \text{可逆}
                        \iff B, C \, \text{可逆}.
                    \]

                    若 $B$ 与 $C$ 均可逆，由第一问的结论，有
                    \[
                        \left( \begin{pmatrix}
                            O & E_m \\ E_n & O
                        \end{pmatrix} \begin{pmatrix}
                            O & B \\ C & D
                        \end{pmatrix} \right)^{-1} = \begin{pmatrix}
                            C^{-1} & -C^{-1}DB^{-1} \\ O & B^{-1}
                        \end{pmatrix},
                    \]

                    故
                    \[
                        \begin{pmatrix}
                            O & B \\ C & D
                        \end{pmatrix}^{-1} = \begin{pmatrix}
                            C^{-1} & -C^{-1}DB^{-1} \\ O & B^{-1}
                        \end{pmatrix} \begin{pmatrix}
                            O & E_m \\ E_n & O
                        \end{pmatrix} = \begin{pmatrix}
                            -C^{-1}DB^{-1} & C^{-1} \\ B^{-1} & O
                        \end{pmatrix}.
                    \]

            \end{enumerate}
\end{solve}

\begin{exercise}[补充训练 L9-B-1]\label{prob:L9-B-1}
设$f(x)=1+x+\cdots+x^{m-1}$，$g(x)=1-x$，$A=\begin{pmatrix}
                a & b \\ 0 & a
            \end{pmatrix}$，计算$f(A)g(A)$.
\end{exercise}
\noindent 来源：\href{https://github.com/yhwu-is/Linear-Algebra-Left-Undone/blob/new/%E8%AE%B2%E4%B9%89/%E4%B8%93%E9%A2%98/9%20%E7%9F%A9%E9%98%B5%E8%BF%90%E7%AE%97%E8%BF%9B%E9%98%B6.tex}{L9-B-1}。
\begin{solve}
\begin{align*}
                f(A) g(A) & = (E + A + A^2 + \cdots + A^{m - 1})(E - A) \\
                          & = E - A^m = \begin{pmatrix}
                                            1 - a^m & -mba^{m - 1} \\
                                            0       & 1 - a^m
                                        \end{pmatrix}.
            \end{align*}
\end{solve}

\begin{exercise}[补充训练 L9-B-2]\label{prob:L9-B-2}
已知矩阵$A=\begin{pmatrix}
                1 & 0 & 4 \\ 0 & 1 & 2 \\ 0 & 1 & 2
            \end{pmatrix}$，求证：所有与$A$可交换的矩阵构成$\mathbf{M}_3(\mathbf{R})$的一个子空间，并求子空间的一组基.
\end{exercise}
\noindent 来源：\href{https://github.com/yhwu-is/Linear-Algebra-Left-Undone/blob/new/%E8%AE%B2%E4%B9%89/%E4%B8%93%E9%A2%98/9%20%E7%9F%A9%E9%98%B5%E8%BF%90%E7%AE%97%E8%BF%9B%E9%98%B6.tex}{L9-B-2}。
\begin{solve}
求与 $ A $ 可交换的矩阵等价于求与 $ A - E $ 可交换的矩阵 $ X $. 可解得
            \begin{align*}
                X & = \begin{pmatrix}
                            x_{11} & -2x_{11} - 2x_{32} + 2x_{33} & 4x_{32} \\
                            0      & -x_{32} + x_{33}             & 2x_{32} \\
                            0      & x_{32}                       & x_{33}
                        \end{pmatrix} \\
                    & = x_{11} \begin{pmatrix}
                                1 & -2 & 0 \\
                                0 & 0  & 0 \\
                                0 & 0  & 0
                            \end{pmatrix}
                + x_{32} \begin{pmatrix}
                            0 & -2 & 4 \\
                            0 & -1 & 2 \\
                            0 & 1  & 0
                        \end{pmatrix}
                + x_{33} \begin{pmatrix}
                            0 & 2 & 0 \\
                            0 & 1 & 0 \\
                            0 & 0 & 1
                        \end{pmatrix},
            \end{align*}
            其中 $ x_{11}, x_{32}, x_{33} $ 为任意实数. 由此我们也得到了 $ A $ 可交换的矩阵构成的子空间的一组基.
\end{solve}

\begin{exercise}[补充训练 L9-B-3]\label{prob:L9-B-3}
已知矩阵$A=\begin{pmatrix}
                1 & 0 & 1 \\ 0 & 2 & 0 \\ 1 & 0 & 1
            \end{pmatrix}$，
        \begin{enumerate}
            \item 求所有与$A$可交换的矩阵；

            \item 若$AB+E=A^2+B$，求$B$.
        \end{enumerate}
\end{exercise}
\noindent 来源：\href{https://github.com/yhwu-is/Linear-Algebra-Left-Undone/blob/new/%E8%AE%B2%E4%B9%89/%E4%B8%93%E9%A2%98/9%20%E7%9F%A9%E9%98%B5%E8%BF%90%E7%AE%97%E8%BF%9B%E9%98%B6.tex}{L9-B-3}。
\begin{solve}
\begin{enumerate}
                \item 不妨设 $ B = (b_{ij})_{3 \times 3} $ 与 $ A $ 可交换，即 $ AB = BA $，这等价于 $ (A - E)B = B(A - E) $，即
                      \[ \begin{pmatrix}
                                &   & 1 \\
                                & 1 &   \\
                              1 &   &
                          \end{pmatrix}
                          \begin{pmatrix}
                              b_{11} & b_{12} & b_{13} \\
                              b_{21} & b_{22} & b_{23} \\
                              b_{31} & b_{32} & b_{33}
                          \end{pmatrix} =
                          \begin{pmatrix}
                              b_{11} & b_{12} & b_{13} \\
                              b_{21} & b_{22} & b_{23} \\
                              b_{31} & b_{32} & b_{33}
                          \end{pmatrix}
                          \begin{pmatrix}
                                &   & 1 \\
                                & 1 &   \\
                              1 &   &
                          \end{pmatrix}. \]
                      于是
                      \[ \begin{pmatrix}
                              b_{31} & b_{32} & b_{33} \\
                              b_{21} & b_{22} & b_{23} \\
                              b_{11} & b_{12} & b_{13}
                          \end{pmatrix} = \begin{pmatrix}
                              b_{11} & b_{12} & b_{13} \\
                              b_{21} & b_{22} & b_{23} \\
                              b_{31} & b_{32} & b_{33}
                          \end{pmatrix} \]
                      对应元素相等，可得 $ b_{31} = b_{13}, b_{32} = b_{12}, b_{33} = b_{11}, b_{21} = b_{23} $，即所有与 $ A $ 可交换的矩阵为
                      \[ \begin{pmatrix}
                              b_{11} & b_{12} & b_{13} \\
                              b_{21} & b_{22} & b_{21} \\
                              b_{13} & b_{12} & b_{11}
                          \end{pmatrix}, \]
                      其中 $ b_{11}, b_{12}, b_{13}, b_{21}, b_{22} $ 为任意实数.

                \item 由于 $ AB + E = A^2 + B $，所以
                      \[ (A - E)B = A^2 - E = (A - E)(A + E). \]
                      又由于 $ |A - E| = -1 \neq 0 $，所以 $ A - E $ 可逆，进而
                      \[ B = A + E = \begin{pmatrix}
                              2 & 0 & 1 \\
                              0 & 3 & 0 \\
                              1 & 0 & 2
                          \end{pmatrix}. \]
            \end{enumerate}
\end{solve}

\begin{exercise}[补充训练 L9-B-4]\label{prob:L9-B-4}
设$A \in \mathbf{F}^{n \times n}$，令$C(A)=\{B \in \mathbf{F}^{n \times n} \mid AB=BA\}$.
        \begin{enumerate}
            \item 证明：$C(A)$为$\mathbf{F}^{n \times n}$的一个子空间；

            \item 求$C(E)$；

            \item 当$A$为对角线上元素互不相等的对角阵时，求$C(A)$的维数和一组基.
        \end{enumerate}
\end{exercise}
\noindent 来源：\href{https://github.com/yhwu-is/Linear-Algebra-Left-Undone/blob/new/%E8%AE%B2%E4%B9%89/%E4%B8%93%E9%A2%98/9%20%E7%9F%A9%E9%98%B5%E8%BF%90%E7%AE%97%E8%BF%9B%E9%98%B6.tex}{L9-B-4}。
\begin{solve}
\begin{enumerate}
                \item 首先有 $ E \in C(A) $，所以 $ C(A) $ 非空. $ \forall B_1, B_2 \in C(A), \lambda \in \mathbf{F} $，有
                      \begin{gather*}
                          A(B_1 + B_2) = AB_1 + AB_2 = B_1A + B_2A = (B_1 + B_2)A \\
                          A(\lambda B_1) = \lambda AB_1 = \lambda B_1A = (\lambda B_1)A.
                      \end{gather*}
                      所以 $ C(A) $ 是 $ \mathbf{F}^{n \times n} $ 的子空间.

                \item $ C(E) = \mathbf{F}^{n \times n} $.

                \item 我们有结论：$ C(A) $ 为全体对角矩阵构成的集合. 故 $ \dim C(A) = n $. 基矩阵 $ B_k = (b_{ij})_{n \times n} $，其中 $ b_{ij} = \delta_{ij} \delta_{jk} $. $ B_1, B_2, \ldots, B_n $ 为一组基.
            \end{enumerate}
\end{solve}

\begin{exercise}[补充训练 L9-B-5]\label{prob:L9-B-5}
设$V$是域$\mathbf{F}$上的$n$阶对称矩阵关于矩阵加法和数乘运算构成的线性空间，令
        \[U=\{A\in V \mid \operatorname{tr}(A)=0\},\enspace W=\{\lambda E \mid \lambda\in\mathbf{F}\}.\]
        \begin{enumerate}
            \item 证明：$U,W$为$V$的子空间；

            \item 分别求$U,W$的一组基和维数；

            \item 证明：$V=U\oplus W$.
        \end{enumerate}
\end{exercise}
\noindent 来源：\href{https://github.com/yhwu-is/Linear-Algebra-Left-Undone/blob/new/%E8%AE%B2%E4%B9%89/%E4%B8%93%E9%A2%98/9%20%E7%9F%A9%E9%98%B5%E8%BF%90%E7%AE%97%E8%BF%9B%E9%98%B6.tex}{L9-B-5}。
\begin{solve}
\begin{enumerate}
                \item 根据定义可知$\operatorname{tr}$是线性的，有 $\operatorname{tr}(\lambda A+\mu B)=\lambda\operatorname{tr}(A)+\mu\operatorname{tr}(B)=0$，则$U$是封闭的，又$W$封闭是显然的，则 $U,W$ 是 $V$ 的子空间.

                \item 对于$W$,基为$\{E\}$，$\dim W=1$. 对于$u$,设$E_{ij}$是 $a_{ij}=1$，其余元素为0的$n$阶方阵. 则由于$u$是对称矩阵，在非对角线元素上，$u$的基包含$E_{12}+E_{21},E_{13}+E_{31},\ldots ,E_{(n-1)n}+E_{n(n-1)}$. 对于对角线元素 $a_{11}+\cdots +a_{nn}=0$. 方程解为
                      \[\begin{pmatrix}a_{11}\\ \vdots\\ a_{nn}	\end{pmatrix}=k_1\begin{pmatrix}1 \\ -1 \\ 0\\ \vdots \\0	\end{pmatrix}=k_2\begin{pmatrix}1 \\ 0 \\ -1\\ \vdots \\0	\end{pmatrix}+\cdots+k_{n-1}\begin{pmatrix}1 \\ 0 \\ 0\\ \vdots \\-1	\end{pmatrix}.\]
                      则还有$n-1$个基$E_{11}-E_{22},E_{11}-E_{33},\ldots ,E_{11}-E_{nn}$. 故 $\dim U=\dfrac{n^2-n}2+n-2=\dfrac{n^2+n-2}2$.

                \item 设$V'=U+W$. 先证明直和，即$U\cap W=\{0\}$. 这是显然的，因为$\operatorname{tr}(\lambda E)=n\lambda$，仅当 $\lambda=0$ 时 $\lambda E\in U$ 即 $U\cap W=\{0\}$得证. 又$\dim U+\dim W=\dim V'=n=\dim V$，则$V=V'=U\oplus W$得证.
            \end{enumerate}
\end{solve}

\begin{exercise}[补充训练 L9-B-9]\label{prob:L9-B-9}
设$A$是$n$阶矩阵，$A^k=O$对某个正整数$k$成立，求证下列方阵可逆，并求它们的逆：
        \begin{enumerate}
            \item $E+A$；

            \item $E-A$；

            \item $E+A+\dfrac{1}{2!}A^2+\cdots+\dfrac{1}{(k-1)!}A^{k-1}$.
        \end{enumerate}
\end{exercise}
\noindent 来源：\href{https://github.com/yhwu-is/Linear-Algebra-Left-Undone/blob/new/%E8%AE%B2%E4%B9%89/%E4%B8%93%E9%A2%98/9%20%E7%9F%A9%E9%98%B5%E8%BF%90%E7%AE%97%E8%BF%9B%E9%98%B6.tex}{L9-B-9}。
\begin{solve}
\begin{enumerate}
                \item 由 $ A^k = O $ 有
                      \[ E = E - A^k = (E - A)(E + A + \cdots + A^{k - 1}). \]
                      故 $ E - A $ 可逆，且
                      \[ (E - A)^{-1} = E + A + \cdots + A^{k - 1}. \]

                \item 若 $ k $ 为偶数，设 $ k = 2m $，由 $ A^k = O $ 有
                      \[ E = E - A^{2m} = (E + A)(E - A)(E + A^2 + A^4 + \cdots + A^{2m - 2}). \]
                      故 $ E + A $ 可逆，且
                      \[ (E + A)^{-1} = (E - A)(E + A^2 + A^4 + \cdots + A^{2m - 2}). \]

                      若 $ k $ 为奇数，设 $ k = 2m + 1 $，由 $ A^k = O $ 有
                      \[ E = E - A^{2m + 1} = (E - A)(E + A + A^2 + \cdots + A^{2m - 1} + A^{2m}) \]
                      故 $ E - A $ 可逆，且
                      \[ (E - A)^{-1} = E + A + A^2 + \cdots + A^{2m - 1} + A^{2m} \]

                \item 由于 $ A^k = O $，于是
                      \[ e^A = E + A + \frac{1}{2!} A^2 + \cdots + \frac{1}{(k - 1)!} A^{k - 1}. \]
                      由 $ e^A e^{-A} = E $ 知 $ E + A + \dfrac{1}{2!} A^2 + \cdots + \dfrac{1}{(k - 1)!} A^{k - 1} $ 可逆，且
                      \[ (E + A + \frac{1}{2!} A^2 + \cdots + \frac{1}{(k - 1)!} A^{k - 1})^{-1} = e^{-A}. \]
            \end{enumerate}
\end{solve}

\begin{exercise}[补充训练 L9-B-10]\label{prob:L9-B-10}
设$A=\begin{pmatrix}
                1      & \cdots & 1      \\
                \vdots & \ddots & \vdots \\
                1      & \cdots & 1
            \end{pmatrix}_{n \times n}$，求：
        \begin{enumerate}
            \item 一个二次实系数多项式$f(x)=ax^2+bx$，使得$f(A)=O$；
            \item $A^{100}$；
            \item $(A+E)^3$；
            \item $(A+E)^{-1}$.
        \end{enumerate}
\end{exercise}
\noindent 来源：\href{https://github.com/yhwu-is/Linear-Algebra-Left-Undone/blob/new/%E8%AE%B2%E4%B9%89/%E4%B8%93%E9%A2%98/9%20%E7%9F%A9%E9%98%B5%E8%BF%90%E7%AE%97%E8%BF%9B%E9%98%B6.tex}{L9-B-10}。
\begin{solve}
\begin{enumerate}
                \item 
\[A^2 = \begin{pmatrix}
                        n      & \cdots & n      \\
                        \vdots & \ddots & \vdots \\
                        n      & \cdots & n
                    \end{pmatrix}\]
，故 $A^2 - nA = O$，令 $a=1, b=-n$ 即可.

                \item 利用 $A^2 = nA$，有
                    \[
                        A^{100} = n A^{99} = n^2 A^{98} = \cdots = n^{99} A = \begin{pmatrix}
                            n^{99} & \cdots & n^{99} \\
                            \vdots & \ddots & \vdots \\
                            n^{99} & \cdots & n^{99}
                        \end{pmatrix}.
                    \]

                \item 利用 $A^2 = nA$，有
                    \begin{align*}
                        (A+E)^3 &= A^3 + 3A^2 + 3A + E \\
                                &= (n^2 + 3n + 3)A + E \\
                                &= \begin{pmatrix}
                                    n^2 + 3n + 4 & n^2 + 3n + 3 & \cdots & n^2 + 3n + 3 \\
                                    n^2 + 3n + 3 & n^2 + 3n + 4 & \cdots & n^2 + 3n + 3 \\
                                    \vdots & \vdots & \ddots & \vdots \\
                                    n^2 + 3n + 3 & n^2 + 3n + 3 & \cdots & n^2 + 3n + 4
                                \end{pmatrix}.
                    \end{align*}

                \item 为得到 $(A+E)^{-1}$，我们需要构造另一个含 $A$ 的一次多项式，该式乘以 $A+E$ 后出现 $A^2 - nA$，以利用条件 $A^2 = nA$. 观察发现所需的多项式为 $A - (n+1)E$：
                    \[
                        (A+E)(A-(n+1)E) = A^2 - nA -(n+1)E = -(n+1)E,
                    \]

                    于是
                    \[
                        (A+E)^{-1} = -\frac{1}{n+1} A + E = \begin{pmatrix}
                            \dfrac{n}{n+1}  & -\dfrac{1}{n+1} & \cdots & -\dfrac{1}{n+1} \\[2em]
                            -\dfrac{1}{n+1} & \dfrac{n}{n+1}  & \cdots & -\dfrac{1}{n+1} \\[2em]
                            \vdots          & \vdots          & \ddots & \vdots          \\[2em]
                            -\dfrac{1}{n+1} & -\dfrac{1}{n+1} & \cdots & \dfrac{n}{n+1}
                        \end{pmatrix}.
                    \]
            \end{enumerate}
\end{solve}

\begin{exercise}[补充训练 L9-B-11]\label{prob:L9-B-11}
设$A,B,A+B$都是可逆矩阵，求证：$A^{-1}+B^{-1}$也是可逆矩阵.
\end{exercise}
\noindent 来源：\href{https://github.com/yhwu-is/Linear-Algebra-Left-Undone/blob/new/%E8%AE%B2%E4%B9%89/%E4%B8%93%E9%A2%98/9%20%E7%9F%A9%E9%98%B5%E8%BF%90%E7%AE%97%E8%BF%9B%E9%98%B6.tex}{L9-B-11}。
\begin{solve}
这里我们遇到的加法和逆在一起的情况，因此我们使用提逆法，不难得到 $A^{-1}+B^{-1}=A^{-1}(A+B)B^{-1}$，由此可知其逆为 $B(A+B)^{-1}A$.
\end{solve}

\begin{exercise}[补充训练 L9-B-12]\label{prob:L9-B-12}
已知矩阵$A=\begin{pmatrix}1 & 2 & a  \\
            1 & 3 & 0  \\
            2 & 7 & -a\end{pmatrix}$可以通过初等列变换转化为矩阵$B=\begin{pmatrix}1  & a & 2 \\
            0  & 1 & 1 \\
            -1 & 1 & 1\end{pmatrix}$.
        \begin{enumerate}
            \item 求常数$a$；
            \item 求满足$AP=B$的可逆矩阵$P$.
        \end{enumerate}
\end{exercise}
\noindent 来源：\href{https://github.com/yhwu-is/Linear-Algebra-Left-Undone/blob/new/%E8%AE%B2%E4%B9%89/%E4%B8%93%E9%A2%98/9%20%E7%9F%A9%E9%98%B5%E8%BF%90%E7%AE%97%E8%BF%9B%E9%98%B6.tex}{L9-B-12}。
\begin{solve}
\begin{enumerate}
                \item 由题意 $r(A)=r(B)$. $A$ 可由初等变换为 $\begin{pmatrix}1 & 2 & a \\ 1 & 3 & 0 \\ 0 & 0 & 0\end{pmatrix}$，$B$ 可由初等变换为 $\begin{pmatrix}1 & a & 2 \\ 0 & 1 & 1 \\ 0 & 0 & 2-a\end{pmatrix}$. 由于秩相同故 $2-a=0,a=2$.

                \item 
\[(A,B)=\begin{pmatrix}1 & 2 & a & 1 & a & 2 \\ 1 & 3 & 0 & 0 & 1 & 1 \\ 2 & 7 & -a & -1 & 1 & 1\end{pmatrix}\rightarrow\begin{pmatrix}1 & 0 & 6 & 3 & 4 & 4 \\ 0 & 1 & -2 & -1 & -1 & -1 \\ 0 & 0 & 0 & 0 & 0 & 0\end{pmatrix}\]


                        $X = \begin{pmatrix}-6k_1+3 \\ 2k_1-1 \\ k_1\end{pmatrix},Y = \begin{pmatrix}-6k_2+4 \\ 2k_2-1 \\ k_2\end{pmatrix},Z = \begin{pmatrix}-6k_3+4 \\ 2k_3-1 \\ k_3\end{pmatrix}$

                        $\begin{pmatrix}-6k_1+3 & -6k_2+4 & -6k_3+4 \\ 2k_1-1 & 2k_2-1 & 2k_3-1 \\ k_1 & k_2 & k_3\end{pmatrix}\rightarrow \begin{pmatrix}1 & 1 & 1 \\ 0 & 1 & 1 \\ 0 & 0 & k_3-k_2\end{pmatrix}$

                        因为 $P$ 可逆，所以 $k_2\neq k_3$，故 $P=\begin{pmatrix}-6k_1+3 & -6k_2+4 & -6k_3+4 \\ 2k_1-1 & 2k_2-1 & 2k_3-1 \\ k_1 & k_2 & k_3\end{pmatrix},k_2\neq k_3,\enspace k_1,k_2,k_3 \in \mathbf{R}$.
            \end{enumerate}
\end{solve}

\begin{exercise}[补充训练 L9-B-13]\label{prob:L9-B-13}
利用列向量线性相关性，证明矩阵秩不等式：\[|r(A)-r(B)|\leqslant r(A\pm B) \leqslant r(A)+r(B).\]
\end{exercise}
\noindent 来源：\href{https://github.com/yhwu-is/Linear-Algebra-Left-Undone/blob/new/%E8%AE%B2%E4%B9%89/%E4%B8%93%E9%A2%98/9%20%E7%9F%A9%E9%98%B5%E8%BF%90%E7%AE%97%E8%BF%9B%E9%98%B6.tex}{L9-B-13}。
\begin{solve}
我们先证明 $r(A \pm B) \leqslant r(A)+r(B)$：

            设 $A = (\alpha_1, \alpha_2, \ldots, \alpha_n), B = (\beta_1, \beta_2, \ldots, \beta_n)$，并设 $r(A) = r, r(B) = s$. 不妨设 $A$ 的列向量组的一个极大线性无关组为 $\alpha_1, \alpha_2, \ldots, \alpha_r$，$B$ 的列向量组的一个极大线性无关组为 $\beta_1, \beta_2, \ldots, \beta_s$.

            取 $\alpha_1, \alpha_2, \ldots, \alpha_r, \beta_1, \beta_2, \ldots, \beta_s$ 的一个极大线性无关组 $\gamma_1, \gamma_2, \ldots, \gamma_p$，则 $p \leqslant r+s$. 该极大线性无关组可以线性表示 $A$ 与 $B$ 的列向量组，因此可以线性表示 $\alpha_1 \pm \beta_1, \alpha_2 \pm \beta_2, \ldots, \alpha_n \pm \beta_n$，即 $A \pm B$ 的列向量组，故
            \[
                r(A \pm B) \leqslant p \leqslant r+s = r(A) + r(B).
            \]

            再证明 $|r(A) - r(B)| \leqslant r(A \pm B)$：

            不妨设 $r(A) \geqslant r(B)$，即证 $r(A) - r(B) \leqslant r(A \pm B)$. 由于
            \begin{align*}
                r(A) &\leqslant r(A + B) + r(-B) = r(A + B) + r(B), \\
                r(A) &\leqslant r(A - B) + r(B),
            \end{align*}

            即有
            \[
                r(A) - r(B) \leqslant r(A \pm B).
            \]

            同理可得若 $r(A) \leqslant r(B)$，则 $r(B) - r(A) \leqslant r(A \pm B)$. 综上，结论成立.
\end{solve}

\begin{exercise}[补充训练 L9-B-14]\label{prob:L9-B-14}
设$B$是$3 \times 1$矩阵，$C$是$1 \times 3$矩阵，证明：$r(BC) \leqslant 1$. 反之，若$A$是秩为1的$3 \times 3$矩阵，证明：存在$3 \times 1$矩阵$B$和$1 \times 3$矩阵$C$，使得$A = BC$.
\end{exercise}
\noindent 来源：\href{https://github.com/yhwu-is/Linear-Algebra-Left-Undone/blob/new/%E8%AE%B2%E4%B9%89/%E4%B8%93%E9%A2%98/9%20%E7%9F%A9%E9%98%B5%E8%BF%90%E7%AE%97%E8%BF%9B%E9%98%B6.tex}{L9-B-14}。
\begin{solve}
若 $B$ 是 $3 \times 1$ 矩阵，$C$ 是 $1 \times 3$ 矩阵，利用秩不等式 $r(BC) \leqslant \min\{r(B), r(C)\}$，以及 $r(B) \leqslant 1, r(C) \leqslant 1$，有
            \[
                r(BC) \leqslant 1.
            \]

            反之，若 $A$ 是秩为 $1$ 的 $3 \times 3$ 矩阵，则 $A$ 任意两个列向量线性相关. 任取其中的一个非零列向量 $\beta$（$r(A)=1$ 保证其存在），则 $A$ 可表示为
            \[
                A = (a_1 \beta, a_2 \beta, a_3 \beta) = \beta (a_1 , a_2 , a_3).
            \]

            令 $B = \beta, C = (a_1 , a_2 , a_3)$，即有 $A = BC$.
\end{solve}

\begin{exercise}[补充训练 L9-B-15]\label{prob:L9-B-15}
设$\alpha,\beta$为$n$维列向量，且$A=\alpha\alpha^\mathrm{T}+\beta\beta^\mathrm{T}$.
        \begin{enumerate}
            \item 证明：$r(A) \leqslant 2$；

            \item 若$\alpha,\beta$线性相关，证明：$r(A) \leqslant 1$.
        \end{enumerate}
\end{exercise}
\noindent 来源：\href{https://github.com/yhwu-is/Linear-Algebra-Left-Undone/blob/new/%E8%AE%B2%E4%B9%89/%E4%B8%93%E9%A2%98/9%20%E7%9F%A9%E9%98%B5%E8%BF%90%E7%AE%97%E8%BF%9B%E9%98%B6.tex}{L9-B-15}。
\begin{solve}
\begin{enumerate}
                \item 设 $\alpha = (a_1, a_2, \ldots, a_n)^\mathrm{T}$，$\beta = (b_1, b_2, \ldots, b_n)^\mathrm{T}$，则
                    \begin{align*}
                        \alpha \alpha^\mathrm{T} &= (a_1 \alpha, a_2 \alpha, \ldots, a_n \alpha), \\
                        \beta \beta^\mathrm{T} &= (b_1 \beta, b_2 \beta, \ldots, b_n \beta).
                    \end{align*}
                    故 $\alpha \alpha^\mathrm{T}$ 和 $\beta \beta^\mathrm{T}$ 的任意两个列向量线性相关，$r\left(\alpha \alpha^\mathrm{T}\right) \leqslant 1$，$r\left(\beta \beta^\mathrm{T}\right) \leqslant 1$，故
                    \[
                        r\left(\alpha \alpha^\mathrm{T} + \beta \beta^\mathrm{T}\right) \leqslant r\left(\alpha \alpha^\mathrm{T}\right) + r\left(\beta \beta^\mathrm{T}\right) \leqslant 2.
                    \]

                \item 若 $\alpha = 0$，则 $A = \beta \beta^\mathrm{T}$，故 $r(A) \leqslant 1$.

                    若 $\alpha \neq 0$，由于 $\alpha, \beta$ 线性相关，存在常数 $k$，使得 $\beta = k \alpha$，则
                    \[
                        A = \alpha \alpha^\mathrm{T} + (k \alpha)(k \alpha)^\mathrm{T} = (k^2 + 1) \alpha \alpha^\mathrm{T},
                    \]
                    故 $r(A) \leqslant 1$.
            \end{enumerate}
\end{solve}

\begin{exercise}[补充训练 L9-B-16]\label{prob:L9-B-16}
设$A$和$B$是两个$n$阶实方阵，证明：$r(A)=r(AB)$当且仅当存在$n$阶实方阵$C$使得$A=ABC$.
\end{exercise}
\noindent 来源：\href{https://github.com/yhwu-is/Linear-Algebra-Left-Undone/blob/new/%E8%AE%B2%E4%B9%89/%E4%B8%93%E9%A2%98/9%20%E7%9F%A9%E9%98%B5%E8%BF%90%E7%AE%97%E8%BF%9B%E9%98%B6.tex}{L9-B-16}。
\begin{solve}
\begin{enumerate}
                \item 存在 $n$ 阶实方阵 $C$ 使得 $A = ABC \implies r(A) = r(AB)$：
                    由秩不等式 $r(AB) \leqslant \min\{r(A), r(B)\}$，我们有
                    \[
                        r(ABC) \leqslant r(AB) \leqslant r(A).
                    \]

                    而由 $A=ABC$ 可知 $r(A) = r(ABC)$，故 $r(A) = r(AB)$.

                \item $r(A) = r(AB) \implies$ 存在 $n$ 阶实方阵 $C$ 使得 $A = ABC$：
                    设 $A = (\alpha_1, \alpha_2, \ldots, \alpha_n)$，$B = (b_{ij})_{n \times n}$，则
                    \begin{align*}
                        AB &= (\alpha_1, \alpha_2, \ldots, \alpha_n) \begin{pmatrix}
                                  b_{11} & b_{12} & \cdots & b_{1n} \\
                                  b_{21} & b_{22} & \cdots & b_{2n} \\
                                  \vdots & \vdots & \ddots & \vdots \\
                                  b_{n1} & b_{n2} & \cdots & b_{nn}
                              \end{pmatrix} \\
                           &= \left( \sum\limits_{i=1}^n b_{i1} \alpha_i, \sum\limits_{i=1}^n b_{i2} \alpha_i, \ldots, \sum\limits_{i=1}^n b_{in} \alpha_i \right).
                    \end{align*}
                    设 $\gamma_j = \sum\limits_{i=1}^n b_{ij} \alpha_i$. 由于 $r(A) = r(AB)$，$\alpha_1, \alpha_2, \ldots, \alpha_n$ 与 $\gamma_1, \gamma_2, \ldots, \gamma_n$ 等价，故存在 $C = (k_{ij})_{n \times n}$，使得
                    \begin{align*}
                        (\alpha_1, \alpha_2, \ldots, \alpha_n)
                        &= (\gamma_1, \gamma_2, \ldots, \gamma_n) \begin{pmatrix}
                            k_{11} & k_{12} & \cdots & k_{1n} \\
                            k_{21} & k_{22} & \cdots & k_{2n} \\
                            \vdots & \vdots & \ddots & \vdots \\
                            k_{n1} & k_{n2} & \cdots & k_{nn}
                        \end{pmatrix} \\
                        &= (\gamma_1, \gamma_2, \ldots, \gamma_n) C,
                    \end{align*}
                    即 $A = ABC$.
            \end{enumerate}
\end{solve}

\begin{exercise}[补充训练 L9-B-17]\label{prob:L9-B-17}
证明如下结论：
        \begin{enumerate}
            \item 若存在正整数$m$使得$r(A^m)=r(A^{m+1})$，则必有$r(A^m)=r(A^{m+1})=r(A^{m+2})=\cdots$；
            \item 设$A$为$n$阶方阵，证明：$r(A^n)=r(A^{n+1})=r(A^{n+2})=\cdots$.
        \end{enumerate}
\end{exercise}
\noindent 来源：\href{https://github.com/yhwu-is/Linear-Algebra-Left-Undone/blob/new/%E8%AE%B2%E4%B9%89/%E4%B8%93%E9%A2%98/9%20%E7%9F%A9%E9%98%B5%E8%BF%90%E7%AE%97%E8%BF%9B%E9%98%B6.tex}{L9-B-17}。
\begin{solve}
\begin{enumerate}
                \item 该问题可以转换为线性映射的像空间停止收缩的问题：
                    设 $A = M(\sigma)$，则 $r(A^m) = r(A^{m+1})$ 等价于 $\dim \operatorname{im} \sigma^m = \dim \operatorname{im} \sigma^{m+1}$.

                    任取 $k \geqslant 2$，由于 $\forall \beta \in \operatorname{im} \sigma^k$，存在 $\alpha \in \mathbf{F}^n$，使得 $\beta = \sigma^k(\alpha) = \sigma^{k-1}(\sigma(\alpha)) \in \operatorname{im} \sigma^{k-1}$，我们有 $\operatorname{im} \sigma^k \subseteq \operatorname{im} \sigma^{k-1}$，即
                    \[
                        \mathbf{F}^n \supseteq \operatorname{im} \sigma \supseteq \operatorname{im} \sigma^2 \supseteq \cdots \supseteq \operatorname{im} \sigma^m \supseteq \operatorname{im} \sigma^{m+1} \supseteq \cdots,
                    \]

                    又 $\dim \operatorname{im} \sigma^m = \dim \operatorname{im} \sigma^{m+1}$，故
                    \[
                        \operatorname{im} \sigma^{m+1} = \operatorname{im} \sigma^m.
                    \]

                    假设 $n=k-1 \enspace (k > m)$ 时，有 $\operatorname{im} \sigma^n = \operatorname{im} \sigma^{n+1}$ 成立，下证 $n=k$ 时该结论也成立：

                    $\forall \gamma \in \operatorname{im} \sigma^k$，存在 $\alpha \in \mathbf{F}^n$，使得 $\gamma = \sigma^k(\alpha) = \sigma(\sigma^{k-1}(\alpha))$. 设 $\beta = \sigma^{k-1}(\alpha) \in \operatorname{im} \sigma^{k-1} = \operatorname{im} \sigma^k$，则 $\gamma = \sigma(\beta) \in \operatorname{im} \sigma^{k+1}$，故 $\operatorname{im} \sigma^k \subseteq \operatorname{im} \sigma^{k+1}$，而另一个方向的包含关系已经证明，所以有
                    \[
                        \operatorname{im} \sigma^{k+1} = \operatorname{im} \sigma^k.
                    \]

                    因此，
                    \[
                        \operatorname{im} \sigma^m = \operatorname{im} \sigma^{m+1} = \operatorname{im} \sigma^{m+2} = \cdots.
                    \]

                    显然有
                    \[
                        r(A^m) = r(A^{m+1}) = r(A^{m+2}) = \cdots.
                    \]

                \item 反证法. 假设 $r(A^n) \neq r(A^{n+1})$，则由 $\operatorname{im} \sigma^{n+1} \subseteq \operatorname{im} \sigma^n$，有 $r(A^n) > r(A^{n+1}) \geqslant 0$.
                    若存在 $k < n$ 使得 $r(A^k) = r(A^{k+1})$，则由第一问可知 $r(A^n) = r(A^{n+1})$，矛盾，故不存在这样的 $k$，即
                    \[
                        n > r(A) > r(A^2) > \cdots > r(A^{n-1}) > r(A^n) > r(A^{n+1}) \geqslant 0,
                    \]

                    由于秩只能是非负整数，上式不可能成立. 故有 $r(A^n) = r(A^{n+1})$. 再由第一问可得
                    \[
                        r(A^n) = r(A^{n+1}) = r(A^{n+2}) = \cdots.
                    \]

            \end{enumerate}
\end{solve}

\begin{exercise}[补充训练 L9-C-3]\label{prob:L9-C-3}
已知数列$\{a_n\},\{b_n\}$满足$a_0=-1,\enspace b_0=3$，且
        \[\begin{cases}
                a_n=3a_{n-1}+b_{n-1}+2^{n-1} \\
                b_n=2a_{n-1}+4b_{n-1}+2^n
            \end{cases}\]
        求$\{a_n\},\{b_n\}$的通项公式.
\end{exercise}
\noindent 来源：\href{https://github.com/yhwu-is/Linear-Algebra-Left-Undone/blob/new/%E8%AE%B2%E4%B9%89/%E4%B8%93%E9%A2%98/9%20%E7%9F%A9%E9%98%B5%E8%BF%90%E7%AE%97%E8%BF%9B%E9%98%B6.tex}{L9-C-3}。
\begin{solve}
设 $v_n = (a_n, b_n, 2^n)^\mathrm{T}$，则递推式可化为
            \[
                v_n = \begin{pmatrix}
                    3 & 1 & 1 \\ 2 & 4 & 2 \\ 0 & 0 & 2
                \end{pmatrix} v_{n-1},
            \]
            其中 $v_0 = (-1, 3, 1)^\mathrm{T}$.

            设 $A = \begin{pmatrix}
                3 & 1 & 1 \\ 2 & 4 & 2 \\ 0 & 0 & 2
            \end{pmatrix}$，则
            \[
                v_n = A^n v_0.
            \]

            下面我们求 $A^n$. 注意到
            \[
                A = 2E + \begin{pmatrix}
                    1 & 1 & 1 \\ 2 & 2 & 2 \\ 0 & 0 & 0
                  \end{pmatrix}
                  = 2E + \begin{pmatrix}
                      1 \\ 2 \\ 0
                  \end{pmatrix} \begin{pmatrix}
                      1 & 1 & 1
                  \end{pmatrix}.
            \]

            令 $\alpha = (1, 2, 0)^\mathrm{T}, \beta = (1, 1, 1), B=\alpha \beta$，则
            \[
                B^n = \left( \alpha \beta \right)^n = \alpha \left( \beta \alpha \right)^{n-1} \beta = \operatorname{tr}(B)^{n-1} \alpha \beta = 3^{n-1} B.
            \]

            此时我们有
            \begin{align*}
                A^n &= (2E + B)^n \\
                    &= \tbinom{n}{0} 2^n E + \tbinom{n}{1} 2^{n-1} B + \tbinom{n}{2} 2^{n-2} B^2 + \cdots + \tbinom{n}{n} B^n \\
                    &= 2^n E + \left(\tbinom{n}{1} 2^{n-1} 3^0 + \tbinom{n}{2} 2^{n-2} 3^1 + \cdots + \tbinom{n}{n} 2^0 3^{n-1} \right) B \\
                    &= 2^n E - \frac{2^n}{3} B + \frac{1}{3} \left(\tbinom{n}{0} 2^n 3^0 + \tbinom{n}{1} 2^{n-1} 3^1 + \cdots + \tbinom{n}{n} 2^0 3^{n-1} \right) B \\
                    &= 2^n E + \frac{5^n - 2^n}{3} B \\
                    &= \frac{1}{3} \begin{pmatrix}
                        5^n + 2 \cdot 2^n & 5^n - 2^n         & 5^n - 2^n    \\
                        2(5^n - 2^n)      & 2 \cdot 5^n + 2^n & 2(5^n - 2^n) \\
                        0                 & 0                 & 3 \cdot 2^n
                    \end{pmatrix}.
            \end{align*}

            代入 $v_n = A^n v_0$，有
            \[
                v_n = \begin{pmatrix}
                    5^n - 2 \cdot 2^n \\
                    2 \cdot 5^n + 2^n \\
                    2^n
                \end{pmatrix}.
            \]

            所以 $a_n = 5^n - 2 \cdot 2^n, b_n = 2 \cdot 5^n + 2^n$.
\end{solve}

\begin{exercise}[补充训练 L9-C-4]\label{prob:L9-C-4}
设$n\geqslant 3$，证明下列矩阵不可逆：
        \[A=\begin{pmatrix}
                \cos(\alpha_1-\beta_1) & \cos(\alpha_1-\beta_2) & \cdots & \cos(\alpha_1-\beta_n) \\
                \cos(\alpha_2-\beta_1) & \cos(\alpha_2-\beta_2) & \cdots & \cos(\alpha_2-\beta_n) \\
                \vdots                 & \vdots                 & \ddots & \vdots                 \\
                \cos(\alpha_n-\beta_1) & \cos(\alpha_n-\beta_2) & \cdots & \cos(\alpha_n-\beta_n)
            \end{pmatrix},\]
        \[B=\begin{pmatrix}
                1+x_1y_1 & 1+x_1y_2 & \cdots & 1+x_1y_n \\
                1+x_2y_1 & 1+x_2y_2 & \cdots & 1+x_2y_n \\
                \vdots   & \vdots   & \ddots & \vdots   \\
                1+x_ny_1 & 1+x_ny_2 & \cdots & 1+x_ny_n
            \end{pmatrix}.\]
\end{exercise}
\noindent 来源：\href{https://github.com/yhwu-is/Linear-Algebra-Left-Undone/blob/new/%E8%AE%B2%E4%B9%89/%E4%B8%93%E9%A2%98/9%20%E7%9F%A9%E9%98%B5%E8%BF%90%E7%AE%97%E8%BF%9B%E9%98%B6.tex}{L9-C-4}。
\begin{solve}
\begin{enumerate}
                \item 可以将 $A$ 拆成一个 $n \times 2$ 矩阵与一个 $2 \times n$ 矩阵的乘积：
                    \[
                        A = \begin{pmatrix}
                            \cos \alpha_1 & \sin \alpha_1 \\
                            \cos \alpha_2 & \sin \alpha_2 \\
                            \vdots        & \vdots        \\
                            \cos \alpha_n & \sin \alpha_n
                        \end{pmatrix} \begin{pmatrix}
                            \cos \beta_1 & \cos \beta_2 & \cdots & \cos \beta_n \\
                            \sin \beta_1 & \sin \beta_2 & \cdots & \sin \beta_n
                        \end{pmatrix}.
                    \]
                    故由秩不等式 $r(AB) \leqslant \min\{r(A), r(B)\}$，有 $r(A) \leqslant 2 < n$，故 $A$ 不可逆.

                \item 可以将 $A$ 拆成两个秩为 $1$ 的矩阵的和：
                    \[
                        A = \begin{pmatrix}
                            1      & 1      & \cdots & 1      \\
                            1      & 1      & \cdots & 1      \\
                            \vdots & \vdots & \ddots & \vdots \\
                            1      & 1      & \cdots & 1
                        \end{pmatrix} + \begin{pmatrix}
                            x_1 y_1 & x_1 y_2 & \cdots & x_1 y_n \\
                            x_2 y_1 & x_2 y_2 & \cdots & x_2 y_n \\
                            \vdots  & \vdots  & \ddots & \vdots  \\
                            x_n y_1 & x_n y_2 & \cdots & x_n y_n
                        \end{pmatrix},
                    \]
                    故由秩不等式 $r(A+B) \leqslant r(A)+r(B)$，有 $r(A) \leqslant 1 + 1 = 2 < n$，故 $A$ 不可逆.
            \end{enumerate}
\end{solve}

\begin{exercise}[补充训练 L9-C-5]\label{prob:L9-C-5}
证明以下两个命题：
        \begin{enumerate}
            \item 与矩阵
\[I=\begin{pmatrix}
                          0 & 1 &   &        &   \\
                            &   & 1 &        &   \\
                            &   &   & \ddots &   \\
                            &   &   &        & 1 \\
                          1 &   &   &        & 0
                      \end{pmatrix}\]
可交换的矩阵$A$都可以写成$I$的一个多项式，即$A=a_{11}E+a_{12}I+a_{13}I^2+\cdots+a_{1n}I^{n-1}$；

            \item 与矩阵
\[J=\begin{pmatrix}
                          0 & 1 &   &        &   \\
                            &   & 1 &        &   \\
                            &   &   & \ddots &   \\
                            &   &   &        & 1 \\
                            &   &   &        & 0
                      \end{pmatrix}\]
可交换的矩阵$A$都可以写成$J$的一个多项式，即$A=a_{11}E+a_{12}J+a_{13}J^2+\cdots+a_{1n}J^{n-1}$.
        \end{enumerate}
\end{exercise}
\noindent 来源：\href{https://github.com/yhwu-is/Linear-Algebra-Left-Undone/blob/new/%E8%AE%B2%E4%B9%89/%E4%B8%93%E9%A2%98/9%20%E7%9F%A9%E9%98%B5%E8%BF%90%E7%AE%97%E8%BF%9B%E9%98%B6.tex}{L9-C-5}。
\begin{solve}
\begin{enumerate}
                \item 设 
\[B = \begin{pmatrix}
                        b_{11} & b_{12} & \cdots & b_{1n} \\
                        b_{21} & b_{22} & \cdots & b_{2n} \\
                        \vdots & \vdots & \ddots & \vdots \\
                        b_{n1} & b_{n2} & \cdots & b_{nn}
                    \end{pmatrix}\]
 与 $I$ 可交换，即 $BI=IB$，则
                    \[
                        \begin{pmatrix}
                            b_{1n} & b_{11} & b_{12} & \cdots & b_{1,n-1} \\
                            b_{2n} & b_{21} & b_{22} & \cdots & b_{2,n-1} \\
                            b_{3n} & b_{31} & b_{32} & \cdots & b_{3,n-1} \\
                            \vdots & \vdots & \vdots &        & \vdots    \\
                            b_{n-1,n} & b_{n-1,1} & b_{n-1,2} & \cdots & b_{n-1,n-1} \\
                            b_{nn} & b_{n1} & b_{n2} & \cdots & b_{n,n-1}
                        \end{pmatrix} = \begin{pmatrix}
                            b_{21} & b_{22} & b_{23} & \cdots & b_{2n} \\
                            b_{31} & b_{32} & b_{33} & \cdots & b_{3n} \\
                            b_{41} & b_{42} & b_{43} & \cdots & b_{4n} \\
                            \vdots & \vdots & \vdots &        & \vdots \\
                            b_{n1} & b_{n2} & b_{n3} & \cdots & b_{nn} \\
                            b_{11} & b_{12} & b_{13} & \cdots & b_{1n}
                        \end{pmatrix},
                    \]
                    逐一对照两矩阵中的对应元素，可得
                    \begin{align*}
                        b_{1n} &= b_{21} = b_{32} = \cdots = b_{n,n-1}, \\
                        b_{1,n-1} &= b_{2n} = b_{31} = \cdots = b_{n,n-2}, \\
                        &\vdotswithin{=} \\
                        b_{11} &= b_{22} = b_{33} = \cdots = b_{nn}.
                    \end{align*}
                    第 $k$ 行等式对应了 $I^{n-k}$，故与 $I$ 可交换的矩阵都可以写成
                    \[
                        a_{11}E + a_{12}I + a_{13}I^2 + \cdots + a_{1n}I^{n-1}.
                    \]

                \item 按照与第一问相同的方法进行操作，即可得到结果.
            \end{enumerate}
\end{solve}

\begin{exercise}[补充训练 L9-C-6]\label{prob:L9-C-6}
证明：若$n$阶矩阵$A$的各阶左上角子块矩阵都可逆，则存在$n$阶下三角矩阵$B$，使得$BA$为上三角矩阵.
\end{exercise}
\noindent 来源：\href{https://github.com/yhwu-is/Linear-Algebra-Left-Undone/blob/new/%E8%AE%B2%E4%B9%89/%E4%B8%93%E9%A2%98/9%20%E7%9F%A9%E9%98%B5%E8%BF%90%E7%AE%97%E8%BF%9B%E9%98%B6.tex}{L9-C-6}。
\begin{solve}
下使用数学归纳法证明该命题. 对于 $n=1$ 的情况，取 $B=E_1$ 即可保证 $BA$ 为上三角矩阵. 现在假设命题对 $n-1$ 阶满足条件的矩阵均成立，下证命题对 $n$ 阶满足条件的矩阵也成立：

            设 $A_n = \begin{pmatrix}
                A_{n-1} & \beta \\
                \alpha  & a_{nn}
            \end{pmatrix}$，其中 $A_{n-1}$ 是 $n-1$ 阶方阵，满足各阶左上角子块矩阵都可逆. 由归纳假设可知存在 $n-1$ 阶下三角矩阵 $B_{n-1}$，使得 $C_{n-1} = B_{n-1}A_{n-1}$ 为上三角矩阵. 下面我们构造 $n$ 阶下三角矩阵 $B_n$：令
            \[
                B_n = \begin{pmatrix}
                    B_{n-1} & \mathbf{0} \\
                    \gamma  & 1
                \end{pmatrix},
            \]
            我们要使
            \[
                C_n = B_n A_n = \begin{pmatrix}
                    B_{n-1} & \mathbf{0} \\
                    \gamma  & 1
                \end{pmatrix} \begin{pmatrix}
                    A_{n-1} & \beta \\
                    \alpha  & a_{nn}
                \end{pmatrix} = \begin{pmatrix}
                    C_{n-1}                 & B_{n-1} \beta         \\
                    \gamma A_{n-1} + \alpha & \gamma \beta + a_{nn}
                \end{pmatrix}
            \]
            为上三角矩阵，只需令 $\gamma A_{n-1} + \alpha = 0$ 即可. 由于 $A_{n-1}$ 可逆，故对任意 $\alpha$ 都有对应的 $\gamma$ 存在，于是 $B_n$ 存在. 因此，命题对 $n$ 阶矩阵也成立.
\end{solve}

\begin{exercise}[补充训练 L9-C-7]\label{prob:L9-C-7}
设$A$是数域$\mathbf{F}$上的$n$阶可逆矩阵，把$A$和$A^{-1}$做如下分块：
        \[A=\begin{pmatrix}
                A_{11} & A_{12} \\ A_{21} & A_{22}
            \end{pmatrix},\enspace A^{-1}=\begin{pmatrix}
                B_{11} & B_{12} \\ B_{21} & B_{22}
            \end{pmatrix}\]
        其中$A_{11}$是$l \times k$矩阵，$B_{11}$是$k \times l$矩阵，$l$，$k$是小于$n$的正整数. 用$W$表示$A_{12}X=0$的解空间，$U$表示$B_{12}Y=0$的解空间，其中$X$和$Y$为列向量，证明$W$与$U$同构.
\end{exercise}
\noindent 来源：\href{https://github.com/yhwu-is/Linear-Algebra-Left-Undone/blob/new/%E8%AE%B2%E4%B9%89/%E4%B8%93%E9%A2%98/9%20%E7%9F%A9%E9%98%B5%E8%BF%90%E7%AE%97%E8%BF%9B%E9%98%B6.tex}{L9-C-7}。
\begin{solve}
证明：$ \forall \alpha \in W,\enspace A_{12} \alpha = \vec{0} $.
            \begin{align*}
                \begin{pmatrix} O_{k \times 1} \\ \alpha \end{pmatrix}
                & = A^{-1}A \begin{pmatrix} O_{k \times 1} \\ \alpha \end{pmatrix} = A^{-1} \begin{pmatrix} A_{11} & A_{12} \\ A_{21} & A_{22} \end{pmatrix} \begin{pmatrix} O_{k \times 1} \\ \alpha \end{pmatrix} = A^{-1} \begin{pmatrix} O_{l \times 1} \\ A_{22} \alpha \end{pmatrix} \\
                & = \begin{pmatrix} B_{11} & B_{12} \\ B_{21} & B_{22} \end{pmatrix} \begin{pmatrix} O_{l \times 1} \\ A_{22} \alpha \end{pmatrix} = \begin{pmatrix} B_{12} A_{22} \alpha \\ B_{22} A_{22} \alpha \end{pmatrix}.
            \end{align*}
            故 $ B_{12} A_{22} \alpha = \vec{0} $，故我们可以推测如下定义：$ \sigma \in \mathcal{L}(W,U),\enspace \sigma(\alpha) = A_{22} \alpha $.
            只需证明 $ \sigma $ 是单射且满射即可.

            单射：$ \sigma(\alpha) = \sigma(\beta) \implies A_{22}(\alpha - \beta) = \vec{0} $. 又有 $ A_{12} \alpha = A_{12} \beta = \vec{0} \implies A_{12}( \alpha - \beta) = \vec{0} $. 故 $ \begin{pmatrix} A_{12} \\ A_{22} \end{pmatrix} (\alpha - \beta) = \vec{0} $. 由于 $ \begin{pmatrix} A_{12} \\ A_{22} \end{pmatrix} $ 列满秩（$ A $ 可逆），故 $ \alpha = \beta $.

            满射：$ \forall \gamma \in U,\enspace B_{12} \gamma = \vec{0} $.
            \begin{align*}
                \begin{pmatrix} O_{l \times 1} \\ \gamma \end{pmatrix}
                & = A A^{-1} \begin{pmatrix} O_{l \times 1} \\ \gamma \end{pmatrix} = A \begin{pmatrix} B_{11} & B_{12} \\ B_{21} & B_{22} \end{pmatrix} \begin{pmatrix} O_{l \times 1} \\ \gamma \end{pmatrix}      \\
                & = A \begin{pmatrix} O_{k \times 1} \\ B_{22} \gamma \end{pmatrix} = \begin{pmatrix} A_{11} & A_{12} \\ A_{21} & A_{22} \end{pmatrix} \begin{pmatrix} O_{k \times 1} \\ B_{22} \gamma \end{pmatrix} \\
                & = \begin{pmatrix} A_{12} B_{22} \gamma \\ A_{22} B_{22} \gamma \end{pmatrix}.
            \end{align*}
            故 $ \exists B_{22} \gamma \in W,\enspace A_{22} B_{22} \gamma = \gamma \in U $.
\end{solve}

\begin{exercise}[补充训练 L9-C-8]\label{prob:L9-C-8}
证明 Schur 补的商等式 $A/B=(A/C)/(B/C)$，其中 $A,B,C$ 均为可逆矩阵，且 $B$ 为 $A$ 的主子矩阵，$C$ 为 $B$ 的主子矩阵.
\end{exercise}
\noindent 来源：\href{https://github.com/yhwu-is/Linear-Algebra-Left-Undone/blob/new/%E8%AE%B2%E4%B9%89/%E4%B8%93%E9%A2%98/9%20%E7%9F%A9%E9%98%B5%E8%BF%90%E7%AE%97%E8%BF%9B%E9%98%B6.tex}{L9-C-8}。
\begin{solve}
设 $A = \begin{pmatrix}
                    A_{11} & A_{12} & A_{13} \\
                    A_{21} & A_{22} & A_{23} \\
                    A_{31} & A_{32} & A_{33}
                \end{pmatrix}$，$B = \begin{pmatrix}
                    A_{11} & A_{12} \\
                    A_{21} & A_{22}
                \end{pmatrix}$，$C = A_{11}$，则根据 Schur 补的定义，有
                \begin{align*}
                    A/C &= \begin{pmatrix}
                        A_{22} & A_{23} \\
                        A_{32} & A_{33}
                    \end{pmatrix} - \begin{pmatrix}
                        A_{21} \\ A_{31}
                    \end{pmatrix} A_{11}^{-1} \begin{pmatrix}
                        A_{12} & A_{13}
                    \end{pmatrix} \\
                    &= \begin{pmatrix}
                        A_{22} - A_{21} A_{11}^{-1} A_{12} & A_{23} - A_{21} A_{11}^{-1} A_{13} \\
                        A_{32} - A_{31} A_{11}^{-1} A_{12} & A_{33} - A_{31} A_{11}^{-1} A_{13}
                    \end{pmatrix},
                \end{align*}
                \[
                    A/B = A_{33} - \begin{pmatrix}
                        A_{31} & A_{32}
                    \end{pmatrix} \begin{pmatrix}
                        A_{11} & A_{12} \\
                        A_{21} & A_{22}
                    \end{pmatrix}^{-1} \begin{pmatrix}
                        A_{13} \\ A_{23}
                    \end{pmatrix}, \enspace
                    B/C = A_{22} - A_{21} A_{11}^{-1} A_{12}.
                \]

                设 $P_{ij} = A_{ij} - A_{i1} A_{11}^{-1} A_{1j} \enspace (i=2,3; j=2,3)$，则
                \[
                    A/C = \begin{pmatrix}
                        P_{22} & P_{23} \\
                        P_{32} & P_{33}
                    \end{pmatrix}, \enspace
                    B/C = P_{22}.
                \]
                故
                \[
                    (A/C)/(B/C) = P_{33} - P_{32} P_{22}^{-1} P_{23}.
                \]

                由分块矩阵初等变换可以求得
                \[
                    \begin{pmatrix}
                        A_{11} & A_{12} \\
                        A_{21} & A_{22}
                    \end{pmatrix}^{-1} = \begin{pmatrix}
                        A_{11}^{-1} + A_{11}^{-1} A_{12} P_{22}^{-1} A_{21} A_{11}^{-1} & -A_{11}^{-1} A_{12} P_{22}^{-1} \\
                        -P_{22}^{-1} A_{21} A_{11}^{-1} & P_{22}^{-1}
                    \end{pmatrix},
                \]
                由此可得
                \begin{align*}
                    A/B = {} & A_{33} - \begin{pmatrix}
                        A_{31} & A_{32}
                    \end{pmatrix} \begin{pmatrix}
                        A_{11}^{-1} + A_{11}^{-1} A_{12} P_{22}^{-1} A_{21} A_{11}^{-1} & -A_{11}^{-1} A_{12} P_{22}^{-1} \\
                        -P_{22}^{-1} A_{21} A_{11}^{-1} & P_{22}^{-1}
                    \end{pmatrix} \begin{pmatrix}
                        A_{13} \\ A_{23}
                    \end{pmatrix} \\
                    = {} & A_{33} - A_{31} \left(A_{11}^{-1} + A_{11}^{-1} A_{12} P_{22}^{-1} A_{21} A_{11}^{-1}\right) - A_{31} \left(-A_{11}^{-1} A_{12} P_{22}^{-1}\right) - \\
                      {} & A_{32} \left(-P_{22}^{-1} A_{21} A_{11}^{-1}\right) - A_{32} P_{22}^{-1} A_{23} \\
                    = {} & \left(A_{33} - A_{31} A_{11}^{-1} A_{13}\right) - \left(A_{32} - A_{31} A_{11}^{-1} A_{12}\right) P_{22}^{-1} \left(A_{23} - A_{21} A_{11}^{-1} A_{13}\right) \\
                    = {} & P_{33} - P_{32} P_{22}^{-1} P_{23} \\
                    = {} & (A/C)/(B/C).
                \end{align*}
\end{solve}

\begin{exercise}[补充训练 L9-C-10]\label{prob:L9-C-10}
已知$A$是一个$s \times n$矩阵，证明：$r(E_n-A^\mathrm{T}A)-r(E_s-AA^\mathrm{T})=n-s$.
\end{exercise}
\noindent 来源：\href{https://github.com/yhwu-is/Linear-Algebra-Left-Undone/blob/new/%E8%AE%B2%E4%B9%89/%E4%B8%93%E9%A2%98/9%20%E7%9F%A9%E9%98%B5%E8%BF%90%E7%AE%97%E8%BF%9B%E9%98%B6.tex}{L9-C-10}。
\begin{solve}
令
            \[
                B = \begin{pmatrix}
                    E_n & A^\mathrm{T} \\
                    A   & E_s
                \end{pmatrix},
            \]
            对 $B$ 做分块倍加变化，有
            \begin{gather*}
                \begin{pmatrix}
                    E_n & A^\mathrm{T} \\
                    A   & E_s
                \end{pmatrix} \to \begin{pmatrix}
                    E_n - A^\mathrm{T}A & A^\mathrm{T} \\
                    O                   & E_s
                \end{pmatrix} \to \begin{pmatrix}
                    E_n - A^\mathrm{T}A & O   \\
                    O                   & E_s
                \end{pmatrix}, \\
                \begin{pmatrix}
                    E_n & A^\mathrm{T} \\
                    A   & E_s
                \end{pmatrix} \to \begin{pmatrix}
                    E_n & A^\mathrm{T}        \\
                    O   & E_s - AA^\mathrm{T}
                \end{pmatrix} \to \begin{pmatrix}
                    E_n & O                   \\
                    O   & E_s - AA^\mathrm{T}
                \end{pmatrix}.
            \end{gather*}
            分块倍加变换不改变矩阵的秩，故有
            \begin{align*}
                r(B) &= r\begin{pmatrix}
                    E_n - A^\mathrm{T}A & O   \\
                    O                   & E_s
                \end{pmatrix} = r(E_n - A^\mathrm{T}A) + s, \\
                r(B) &= r\begin{pmatrix}
                    E_n & O                   \\
                    O   & E_s - AA^\mathrm{T}
                \end{pmatrix} = r(E_s - AA^\mathrm{T}) + n.
            \end{align*}
            于是有 $r(E_n-A^\mathrm{T}A)-r(E_s-AA^\mathrm{T})=n-s$.
\end{solve}

\begin{exercise}[补充训练 L9-C-12]\label{prob:L9-C-12}
$f(x)=f_1(x)f_2(x)$是多项式，且$f_1(x)$与$f_2(x)$互素，则$f(A)=O$的充要条件是$r(f_1(A))+r(f_2(A))=n$. （注：此题的推论非常多，如$A^2=A$，$A^n=E$等形式的结论都可以利用这个例子推导出）
\end{exercise}
\noindent 来源：\href{https://github.com/yhwu-is/Linear-Algebra-Left-Undone/blob/new/%E8%AE%B2%E4%B9%89/%E4%B8%93%E9%A2%98/9%20%E7%9F%A9%E9%98%B5%E8%BF%90%E7%AE%97%E8%BF%9B%E9%98%B6.tex}{L9-C-12}。
\begin{solve}
略有超纲，使用贝祖定理. $\exists u(x),v(x),u(x)f_1(x)+v(x)f_2(x)=1 $，于是
            \begin{align*}
                r\begin{pmatrix}
                    f_1(A) & O      \\
                    O      & f_2(A)
                \end{pmatrix}
                &= r\begin{pmatrix}
                        f_1(A) & f_1(A)u(A)+f_2(A)v(A) \\
                        O      & f_2(A)
                    \end{pmatrix} = r\begin{pmatrix}
                        f_1(A) & E_n    \\
                        O      & f_2(A)
                    \end{pmatrix} \\
                &= r\begin{pmatrix}
                        f_1(A)        & E_n \\
                        -f_2(A)f_1(A) & O
                    \end{pmatrix} = r\begin{pmatrix}
                        O    & E_n \\
                        f(A) & O
                    \end{pmatrix}.
            \end{align*}
            因此
            \[
                f(A) = O \iff
                r\begin{pmatrix}
                    f_1(A) & O      \\
                    O      & f_2(A)
                \end{pmatrix} = r\begin{pmatrix}
                    O    & E_n \\
                    f(A) & O
                \end{pmatrix} = n \iff
                r(f_1(A)) + r(f_2(A)) = n.
            \]
\end{solve}

\begin{exercise}[补充训练 L9-C-13]\label{prob:L9-C-13}
设$A,B$分别为$3 \times 2$和$2 \times 3$实矩阵. 若$AB=\begin{pmatrix}
                8             & 0 & -4 \\[1ex]
                -\dfrac{3}{2} & 9 & -6 \\[1ex]
                -2            & 0 & 1
            \end{pmatrix}$，求$BA$.
\end{exercise}
\noindent 来源：\href{https://github.com/yhwu-is/Linear-Algebra-Left-Undone/blob/new/%E8%AE%B2%E4%B9%89/%E4%B8%93%E9%A2%98/9%20%E7%9F%A9%E9%98%B5%E8%BF%90%E7%AE%97%E8%BF%9B%E9%98%B6.tex}{L9-C-13}。
\begin{solve}
由于 $A$ 是列满秩矩阵，$B$ 是行满秩矩阵，知存在可逆矩阵 $P_{3\times 3},Q_{2\times 2}$ 使得
            \[A=P\begin{pmatrix}E_2 \\ O\end{pmatrix},B=\begin{pmatrix}E_2 & O\end{pmatrix}Q,\]
            于是
            \[BA=\begin{pmatrix}E_2 & O\end{pmatrix}QP\begin{pmatrix}E_2 \\ O\end{pmatrix}.\]
            由 $(AB)^2=9AB$ 有 \[P\begin{pmatrix}E_2 \\ O\end{pmatrix}\begin{pmatrix}E_2 & O\end{pmatrix}QP\begin{pmatrix}E_2 \\ O\end{pmatrix}\begin{pmatrix}E_2 & O\end{pmatrix}Q=9P\begin{pmatrix}E_2 \\ O\end{pmatrix}\begin{pmatrix}E_2 & O\end{pmatrix}Q,\]
            即
            \[\begin{pmatrix}E_2 \\ O\end{pmatrix}BA\begin{pmatrix}E_2 & O\end{pmatrix}=9\begin{pmatrix}E_2 \\ O\end{pmatrix}\begin{pmatrix}E_2 & O\end{pmatrix},\]
            也就是 \[\begin{pmatrix}BA & O \\ O & O\end{pmatrix}=\begin{pmatrix}9E_2 & 0 \\ 0 & 0\end{pmatrix}.\]
            所以 $BA=9E_2$.
\end{solve}

\begin{exercise}[补充训练 L9-例-2]\label{prob:L9-例-2}
求所有与$A$可交换的矩阵，其中
    \[A=\begin{pmatrix}
            1 & 0 & 0  \\
            0 & 1 & 2  \\
            0 & 1 & -2
        \end{pmatrix}.\]
\end{exercise}
\noindent 来源：\href{https://github.com/yhwu-is/Linear-Algebra-Left-Undone/blob/new/%E8%AE%B2%E4%B9%89/%E4%B8%93%E9%A2%98/9%20%E7%9F%A9%E9%98%B5%E8%BF%90%E7%AE%97%E8%BF%9B%E9%98%B6.tex}{L9-例-2}。
\begin{solve}
设$B$为与$A$可交换的矩阵，设
    \[B=\begin{pmatrix}
            a & b & c \\
            d & e & f \\
            g & h & i
        \end{pmatrix},\]
    则
    \[AB=\begin{pmatrix}
            a    & b    & c    \\
            d+2g & e+2h & f+2i \\
            d-2g & e-2h & f-2i
        \end{pmatrix}=BA=\begin{pmatrix}
            a & b+c & 2b-2c \\
            d & e+f & 2e-2f \\
            g & h+i & 2h-2i
        \end{pmatrix}.\]
    由此可得
    \[\begin{cases}
            b=b+c      \\
            c=2b-2c    \\
            d+2g=d     \\
            e+2h=e+f   \\
            f+2i=2e-2f \\
            d-2g=g     \\
            e-2h=h+i   \\
            f-2i=2h-2i
        \end{cases},\]
    很容易解得$b=c=d=g=0$，且$f=2h,\enspace e=3h+i$，因此
    \[B=\begin{pmatrix}
            a & 0    & 0  \\
            0 & 3h+i & 2h \\
            0 & h    & i
        \end{pmatrix}.\]
\end{solve}

\begin{exercise}[补充训练 L9-例-3]\label{prob:L9-例-3}
求与矩阵$A=\begin{pmatrix}
            3  & 0  & 0 \\
            -1 & 3  & 0 \\
            0  & -1 & 3
        \end{pmatrix}$可交换的矩阵.
\end{exercise}
\noindent 来源：\href{https://github.com/yhwu-is/Linear-Algebra-Left-Undone/blob/new/%E8%AE%B2%E4%B9%89/%E4%B8%93%E9%A2%98/9%20%E7%9F%A9%E9%98%B5%E8%BF%90%E7%AE%97%E8%BF%9B%E9%98%B6.tex}{L9-例-3}。
\begin{solve}
由前述分析，取$t=3$可以使得$A-3E$对角线上元素全为0，便于计算，因此我们只需求与$A-3E$可交换的矩阵即可. 与原讲义“可交换矩阵1”处 同样的设未知数法，具体过程省略得到与$A$可交换的矩阵为
    \[B=\begin{pmatrix}
            a & 0 & 0 \\
            b & a & 0 \\
            c & b & a
        \end{pmatrix}.\]
\end{solve}

\begin{exercise}[补充训练 L9-例-10]\label{prob:L9-例-10}
设 $A$ 为 $n$ 阶可逆矩阵，考虑 $B = A + XRY$，其中 $X$ 是 $n \times r$ 矩阵，$Y$ 是 $r \times n$ 矩阵，而 $R$ 是 $r \times r$ 的可逆矩阵. 若 $R^{-1} + YA^{-1}X$ 可逆，则
    \[B^{-1}=A^{-1}-A^{-1}X(R^{-1}+YA^{-1}X)^{-1}YA^{-1}.\]
\end{exercise}
\noindent 来源：\href{https://github.com/yhwu-is/Linear-Algebra-Left-Undone/blob/new/%E8%AE%B2%E4%B9%89/%E4%B8%93%E9%A2%98/9%20%E7%9F%A9%E9%98%B5%E8%BF%90%E7%AE%97%E8%BF%9B%E9%98%B6.tex}{L9-例-10}。
\begin{solve}
注意到 $B = A + XRY = A(E + A^{-1}XRY)$，将 $A^{-1}X$ 和 $RY$ 分别视为整体，因为 $E + RYA^{-1}X = R(R^{-1} + YA^{-1}X)$ 可逆，根据原讲义“凑因子3”处，则有 $E + A^{-1}XRY$ 可逆，并且
    \begin{align*}
        (E + A^{-1}XRY)^{-1} & = E - A^{-1}X(E + RYA^{-1}X)^{-1}RY \\
                             & = E - A^{-1}X(R^{-1} + YA^{-1}X)^{-1}Y.
    \end{align*}
    因此 $B^{-1} = (E + A^{-1}XRY)^{-1}A^{-1} = A^{-1} - A^{-1}X(R^{-1} + YA^{-1}X)^{-1}YA^{-1}$.
\end{solve}

\begin{exercise}[补充训练 L9-例-11]\label{prob:L9-例-11}
求 $n$ 阶矩阵 
\[A = \begin{pmatrix}
        0 & 1 & 1 & \cdots & 1 \\
        1 & 0 & 1 & \cdots & 1 \\
        1 & 1 & 0 & \cdots & 1 \\
        \vdots & \vdots & \vdots & \ddots & \vdots \\
        1 & 1 & 1 & \cdots & 0
    \end{pmatrix}\]
 的逆矩阵.
\end{exercise}
\noindent 来源：\href{https://github.com/yhwu-is/Linear-Algebra-Left-Undone/blob/new/%E8%AE%B2%E4%B9%89/%E4%B8%93%E9%A2%98/9%20%E7%9F%A9%E9%98%B5%E8%BF%90%E7%AE%97%E8%BF%9B%E9%98%B6.tex}{L9-例-11}。
\begin{solve}
不难观察到这一矩阵与秩为 $1$ 的元素全为 $1$ 的矩阵之间只差一个单位矩阵，因此令 $\alpha = (1,1,\ldots,1)^\mathrm{T}$，则 $A = -E_n + \alpha\alpha^\mathrm{T}$，由 Sherman-Morrison 公式，我们有
    \begin{align*}
        A^{-1} &= -E_n - (1 + \alpha^\mathrm{T}(-E_n)\alpha)^{-1}(-E_n)\alpha\alpha^\mathrm{T}(-E_n) \\
                &= -E_n + \dfrac{1}{n-1}\alpha\alpha^\mathrm{T}.
    \end{align*}
\end{solve}

\begin{exercise}[补充训练 L9-例-12]\label{prob:L9-例-12}
已知方阵$A$满足$A^k=O$，其中$k$是一个正整数，求$E-A$的逆.
\end{exercise}
\noindent 来源：\href{https://github.com/yhwu-is/Linear-Algebra-Left-Undone/blob/new/%E8%AE%B2%E4%B9%89/%E4%B8%93%E9%A2%98/9%20%E7%9F%A9%E9%98%B5%E8%BF%90%E7%AE%97%E8%BF%9B%E9%98%B6.tex}{L9-例-12}。
\begin{solve}
根据我们前面的分析，结合$A^k=O$，我们猜测
    \[(E-A)^{-1}=E+A+A^2+\cdots+A^{k-1}.\]
    事实上我们直接验证
    \[(E-A)(E+A+A^2+\cdots+A^{k-1})=E-A^k=E.\]
    因此$(E-A)^{-1}=E+A+A^2+\cdots+A^{k-1}$.
\end{solve}

\begin{exercise}[补充训练 L9-例-16]\label{prob:L9-例-16}
设$A=\begin{pmatrix}0 & -1 & 0 \\ 1 & 0 & 0 \\ 0 & 0 & -1 \end{pmatrix},\enspace P^{-1}AP=B$，求$B^{2004}-2A^2$.
\end{exercise}
\noindent 来源：\href{https://github.com/yhwu-is/Linear-Algebra-Left-Undone/blob/new/%E8%AE%B2%E4%B9%89/%E4%B8%93%E9%A2%98/9%20%E7%9F%A9%E9%98%B5%E8%BF%90%E7%AE%97%E8%BF%9B%E9%98%B6.tex}{L9-例-16}。
\begin{solve}
事实上，$A^2=\begin{pmatrix}
                      -1 & 0 & 0 \\ 0 & -1 & 0 \\ 0 & 0 & 1
                  \end{pmatrix}$，因此$A^4=E$，于是$A^{2004}=(A^4)^{501}=E$，因此$B^{2004}-2A^2=E-2A^2=\begin{pmatrix}
                      3 & 0 & 0 \\ 0 & 3  & 0 \\ 0 & 0 & -1
                  \end{pmatrix}$.
\end{solve}

\begin{exercise}[补充训练 L9-例-21]\label{prob:L9-例-21}
已知$M$是秩为 1 的矩阵，记$\operatorname{tr}(M)=b$，讨论$(aE+M)^n$的计算结果.
\end{exercise}
\noindent 来源：\href{https://github.com/yhwu-is/Linear-Algebra-Left-Undone/blob/new/%E8%AE%B2%E4%B9%89/%E4%B8%93%E9%A2%98/9%20%E7%9F%A9%E9%98%B5%E8%BF%90%E7%AE%97%E8%BF%9B%E9%98%B6.tex}{L9-例-21}。
\begin{solve}
事实上，由原讲义“相抵分解”处 可知，$M^k=b^{k-1}M$，因此
              \begin{enumerate}
                  \item 当$b=0$时，$M^k=O(k\geqslant 2)$，因此
                        \[(aE+M)^n=a^nE+na^{n-1}M.\]

                  \item 当$b\neq 0$时有
                        \begin{align*}
                            (aE+M)^n & =\sum\limits_{k=0}^nC_n^ka^{n-k}M^k                                    \\
                                     & =\sum\limits_{k=0}^nC_n^ka^{n-k}b^{k-1}M                               \\
                                     & =a^nE+\dfrac{1}{b}\sum\limits_{k=1}^nC_n^ka^{n-k}b^kM                  \\
                                     & =a^nE+\dfrac{1}{b}\left(\sum\limits_{k=0}^nC_n^ka^{n-k}b^k-a^n\right)M \\
                                     & =a^nE+\dfrac{(a+b)^n-a^n}{b}M
                        \end{align*}
              \end{enumerate}
\end{solve}

\begin{exercise}[补充训练 L9-例-22]\label{prob:L9-例-22}
求矩阵的幂
\[A^n=\begin{pmatrix}
                      1  & -1 & -1 & -1 \\
                      -1 & 1  & -1 & -1 \\
                      -1 & -1 & 1  & -1 \\
                      -1 & -1 & -1 & 1
                  \end{pmatrix}^n\]
.
\end{exercise}
\noindent 来源：\href{https://github.com/yhwu-is/Linear-Algebra-Left-Undone/blob/new/%E8%AE%B2%E4%B9%89/%E4%B8%93%E9%A2%98/9%20%E7%9F%A9%E9%98%B5%E8%BF%90%E7%AE%97%E8%BF%9B%E9%98%B6.tex}{L9-例-22}。
\begin{solve}
很显然我们可以将$A$分为一个数量矩阵秩1矩阵的和，即\[A=2E+\begin{pmatrix}
                      -1 & -1 & -1 & -1 \\
                      -1 & -1 & -1 & -1 \\
                      -1 & -1 & -1 & -1 \\
                      -1 & -1 & -1 & -1
                  \end{pmatrix},\]
              然后我们利用上例的结论即可（或者不记得结论也可以将
\[\begin{pmatrix}
                      -1 & -1 & -1 & -1 \\
                      -1 & -1 & -1 & -1 \\
                      -1 & -1 & -1 & -1 \\
                      -1 & -1 & -1 & -1
                  \end{pmatrix}\]
分解成$\begin{pmatrix}
                      -1 \\ -1 \\ \vdots \\ -1
                  \end{pmatrix}\begin{pmatrix}
                      1 & 1 & \cdots & 1
                  \end{pmatrix}$然后现场推导）.
\end{solve}

\begin{exercise}[补充训练 L9-例-24]\label{prob:L9-例-24}
设$A$为三阶矩阵，$P$为三阶可逆矩阵，$P^{-1}AP=B$，其中
              \[P=\begin{pmatrix}
                      0 & 2 & -1 \\ 1 & 1 & 2 \\ -1 & -1 & -1
                  \end{pmatrix},\enspace B=\begin{pmatrix}
                      0 & 0 & -1 \\ 0 & -1 & 0 \\ -1 & 0 & 0
                  \end{pmatrix},\]
              求$A^{2024}$.
\end{exercise}
\noindent 来源：\href{https://github.com/yhwu-is/Linear-Algebra-Left-Undone/blob/new/%E8%AE%B2%E4%B9%89/%E4%B8%93%E9%A2%98/9%20%E7%9F%A9%E9%98%B5%E8%BF%90%E7%AE%97%E8%BF%9B%E9%98%B6.tex}{L9-例-24}。
\begin{solve}
事实上$A=PBP^{-1}$，因此$A^{2024}=PB^{2024}P^{-1}$，由于$B^2=E$，因此$B^{2024}=(B^2)^{1012}=E$，因此$A^{2024}=PEP^{-1}=E$.
\end{solve}

\begin{exercise}[补充训练 L9-例-25]\label{prob:L9-例-25}
当$D$和$A-BD^{-1}C$可逆时，通过分块矩阵初等变换求$P=\begin{pmatrix}A & B \\ C & D\end{pmatrix}$的逆矩阵.
\end{exercise}
\noindent 来源：\href{https://github.com/yhwu-is/Linear-Algebra-Left-Undone/blob/new/%E8%AE%B2%E4%B9%89/%E4%B8%93%E9%A2%98/9%20%E7%9F%A9%E9%98%B5%E8%BF%90%E7%AE%97%E8%BF%9B%E9%98%B6.tex}{L9-例-25}。
\begin{solve}
设$A$和$D$分别是$m$、$n$阶矩阵，根据前面的讨论，我们的想法就是对$P$和$E$同时做初等变换使得$P$变为单位矩阵，那么我们自然会想到打洞法，利用$D$可逆的条件消去左下角和右上角的矩阵. 我们首先对第二个分块行左乘$-BD^{-1}$加到第一个分块行，有
    \[\left(\begin{array}{cc:cc}
                A & B & E_m & O \\ C & D & O & E_n
            \end{array}\right)\rightarrow\left(\begin{array}{cc:cc}
                A-BD^{-1}C & O & E_m & -BD^{-1} \\ C & D & O & E_n
            \end{array}\right),\]
    这样消去了右上角的分块阵，接下来消去左下角的（当然也可以换个顺序消去，效果一致），我们此时利用$A-BD^{-1}C$可逆（继续利用$D$可逆得到的结果不具代表性，故不展开），将第一个分块行左乘$-C(A-BD^{-1}C)^{-1}$加到第二个分块行，得到
    \[\left(\begin{array}{cc:cc}
                A-BD^{-1}C & O & E_m & -BD^{-1} \\ O & D & -C(A-BD^{-1}C)^{-1} & E_n+C(A-BD^{-1}C)^{-1}BD^{-1}
            \end{array}\right),\]
    最后用$(A-BD^{-1}C)^{-1}$和$D^{-1}$分别左乘第一分块行和第二分块行，得到
    \[\left(\begin{array}{cc:cc}
                E_m & O & (A-BD^{-1}C)^{-1} & -(A-BD^{-1}C)^{-1}BD^{-1} \\ O & E_n & -D^{-1}C(A-BD^{-1}C)^{-1} & D^{-1}+D^{-1}C(A-BD^{-1}C)^{-1}BD^{-1}
            \end{array}\right),\]
    由此可得原矩阵的逆就是上述虚线右侧的
    \[\begin{pmatrix}
            (A-BD^{-1}C)^{-1} & -(A-BD^{-1}C)^{-1}BD^{-1} \\ -D^{-1}C(A-BD^{-1}C)^{-1} & D^{-1}+D^{-1}C(A-BD^{-1}C)^{-1}BD^{-1}
        \end{pmatrix}.\]
\end{solve}

\begin{exercise}[补充训练 L9-例-28]\label{prob:L9-例-28}
若$n$阶矩阵$A$的各阶左上角子块矩阵都可逆，则存在上三角矩阵$U$和主对角元全为1的下三角矩阵$L$，使得$A=LU$（$LU$分解）.
\end{exercise}
\noindent 来源：\href{https://github.com/yhwu-is/Linear-Algebra-Left-Undone/blob/new/%E8%AE%B2%E4%B9%89/%E4%B8%93%E9%A2%98/9%20%E7%9F%A9%E9%98%B5%E8%BF%90%E7%AE%97%E8%BF%9B%E9%98%B6.tex}{L9-例-28}。
\begin{solve}
由于主对角元全为1的下三角矩阵可逆，其逆矩阵也是主对角元全为1的下三角矩阵，因此只要证明存在主对角元全为1的下三角矩阵$S$使得$SA=U$，我们可以利用数学归纳法来证明这一结论.

    当$n=1$时，结论显然成立. 假设命题对$n-1$阶矩阵成立. 对$n$阶矩阵$A$，将$A$分块为
    \[A=\begin{pmatrix}
            A_1 & \alpha_1 \\ \alpha_2 & a_{nn}
        \end{pmatrix},\]
    其中$A_1$为满足命题条件的$n-1$阶矩阵，根据归纳假设，存在上三角矩阵$U_1$和$n-1$阶主对角元全为1的下三角矩阵$S_1$使得$S_1A_1=U_1$. 根据我们前面的讨论，我们可以对$A$做分块矩阵初等变换将其化为上三角块矩阵，即
    \[PA=\begin{pmatrix}
            E_{n-1} & O \\ -\alpha_2A_1^{-1} & 1
        \end{pmatrix}\begin{pmatrix}
            A_1 & \alpha_1 \\ \alpha_2 & a_{nn}
        \end{pmatrix}=\begin{pmatrix}
            A_1 & \alpha_1 \\ O & a_{nn}-\alpha_2A_1^{-1}\alpha_1
        \end{pmatrix},\]
    再取$Q=\begin{pmatrix}
            S_1 & O \\ O & 1
        \end{pmatrix}$，即得
    \[QPA=\begin{pmatrix}
            U_1 & S_1\alpha_1 \\ O & a_{nn}-\alpha_2A_1^{-1}\alpha_1
        \end{pmatrix}=U.\]
    其中$U$为上三角矩阵，$S=QP$是主对角元全为1的下三角矩阵，故存在$L=S^{-1}$和$U$使得$A=LU$. 从而由归纳法得证.
\end{solve}

\begin{exercise}[补充训练 L9-例-29]\label{prob:L9-例-29}
设矩阵$A=\begin{pmatrix}
            1 & 2 & -1 \\ 2 & 1 & -2 \\ -3 & 1 & 1
        \end{pmatrix}$，求$A$的$LU$分解，并利用$LU$分解求解方程组$Ax=b$，其中$b=(3,3,-6)^\mathrm{T}$.
\end{exercise}
\noindent 来源：\href{https://github.com/yhwu-is/Linear-Algebra-Left-Undone/blob/new/%E8%AE%B2%E4%B9%89/%E4%B8%93%E9%A2%98/9%20%E7%9F%A9%E9%98%B5%E8%BF%90%E7%AE%97%E8%BF%9B%E9%98%B6.tex}{L9-例-29}。
\begin{solve}
对于$A$，我们可以通过下述三个初等行变换化为上三角矩阵：
    \begin{enumerate}
        \item 第一行乘以$-2$加到第二行；
        \item 第一行乘以$3$加到第三行；
        \item 第二行乘以$\dfrac{7}{3}$加到第三行.
    \end{enumerate}
    记三个初等变换对应的矩阵依次为$P_1,P_2,P_3$，则上述初等变换过程可以写为
    \begin{equation}\label{eq:10:LU分解}
        P_3P_2P_1A=\begin{pmatrix}
            1 & 2 & -1 \\ 0 & -3 & 0 \\ 0 & 0 & -2
        \end{pmatrix}=U.
    \end{equation}
    即$A=(P_3P_2P_1)^{-1}U$，其中$U$为上三角矩阵，我们希望$L=(P_3P_2P_1)^{-1}$为下三角矩阵（事实上我们做的初等行变换保证了这一点），因此我们只需要求出$L$即可. 由于$P_1,P_2,P_3$均为初等矩阵，因此它们的逆矩阵也是初等矩阵，我们可以很容易地求出它们的逆矩阵为$L=(P_3P_2P_1)^{-1}=P_1^{-1}P_2^{-1}P_3^{-1}=\begin{pmatrix}
            1 & 0 & 0 \\ 2 & 1 & 0 \\ -3 & -\dfrac{7}{3} & 1
        \end{pmatrix}$. 因此根据式~\eqref{eq:10:LU分解} 可知我们得到了$A$的$LU$分解为
    \[A=\begin{pmatrix}
            1 & 0 & 0 \\ 2 & 1 & 0 \\ -3 & -\dfrac{7}{3} & 1
        \end{pmatrix}\begin{pmatrix}
            1 & 2 & -1 \\ 0 & -3 & 0 \\ 0 & 0 & -2
        \end{pmatrix}=LU.\]
    接下来我们利用分解求解方程组$Ax=b$. 我们首先令$c=Ux$，则方程组可以化为$Lc=b$，写成方程组的形式即为
    \[\begin{cases}
            c_1=3 \\ 2c_1+c_2=3 \\ -3c_1-\dfrac{7}{3}c_2+c_3=-6
        \end{cases},\]
    从第一个方程开始解很容易得到$c_1=3,c_2=-3,c_3=-4$，因此我们得到了$c=(3,-3,-4)^\mathrm{T}$. 最后我们利用$Ux=c$求解出$x$，即
    \[\begin{cases}
            x_1+2x_2-x_3=3 \\ -3x_2=-3 \\ -2x_3=-4
        \end{cases},\]
    从第三个方程开始解很容易得到$x_1=3,x_2=1,x_3=2$，因此我们得到了$x=(3,1,2)^\mathrm{T}$.
\end{solve}

\begin{exercise}[补充训练 L9-例-30]\label{prob:L9-例-30}
设$A=\begin{pmatrix}1 & 0 & 0 \\ 1 & 1 & 0 \\ 1 & 1 & 1\end{pmatrix},\ B=\begin{pmatrix}0 & 1 & 1 \\ 1 & 0 & 1 \\ 1 & 1 & 0\end{pmatrix}$，求矩阵$X$满足：
    \[AXA+BXB=AXB+BXA+A(A-B)\]
\end{exercise}
\noindent 来源：\href{https://github.com/yhwu-is/Linear-Algebra-Left-Undone/blob/new/%E8%AE%B2%E4%B9%89/%E4%B8%93%E9%A2%98/9%20%E7%9F%A9%E9%98%B5%E8%BF%90%E7%AE%97%E8%BF%9B%E9%98%B6.tex}{L9-例-30}。
\begin{solve}
由$AXA+BXB=AXB+BXA+A(A-B)$得$(A-B)X(A-B)=A(A-B)$，又$A-B=\begin{pmatrix}
            1 & -1 & -1 \\ 0 & 1 & -1 \\ 0 & 0 & 1
        \end{pmatrix}$，很容易判断$A-B$可逆（可以参考原讲义“可逆等价条件”处），因此$(A-B)X=A$，即$X=(A-B)^{-1}A$.

    我们直接对$(A-B,A)$做初等行变换，有
    \[\left(\begin{array}{ccc:ccc}
                1 & -1 & -1 & 1 & 0 & 0 \\
                0 & 1  & -1 & 1 & 1 & 0 \\
                0 & 0  & 1  & 1 & 1 & 1
            \end{array}\right)\xrightarrow{\text{初等行变换}}\left(\begin{array}{ccc:ccc}
                1 & 0 & 0 & 4 & 3 & 2 \\
                0 & 1 & 0 & 2 & 2 & 1 \\
                0 & 0 & 1 & 1 & 1 & 1
            \end{array}\right),\]
    因此$X=\begin{pmatrix}
            4 & 3 & 2 \\ 2 & 2 & 1 \\ 1 & 1 & 1
        \end{pmatrix}$.
\end{solve}

\begin{exercise}[补充训练 L9-例-32]\label{prob:L9-例-32}
若$A,B$为两个$n$阶矩阵，则
    \begin{enumerate}[label=\Alph*.]
        \item $r(A,B)=r(A^\mathrm{T},B^\mathrm{T})$

        \item $r(A,AB)=r(A)$

        \item $r(A,B)=\max\{r(A), r(B)\}$

        \item $r(A,BA)=r(A)$
    \end{enumerate}
\end{exercise}
\noindent 来源：\href{https://github.com/yhwu-is/Linear-Algebra-Left-Undone/blob/new/%E8%AE%B2%E4%B9%89/%E4%B8%93%E9%A2%98/9%20%E7%9F%A9%E9%98%B5%E8%BF%90%E7%AE%97%E8%BF%9B%E9%98%B6.tex}{L9-例-32}。
\begin{solve}
选择B项.
    \begin{enumerate}[label=\Alph*.]
        \item 要注意分块矩阵$(A\enspace B)$的转置是
              \[\begin{pmatrix}
                      A^\mathrm{T} \\ B^\mathrm{T}
                  \end{pmatrix}\]
              而非这里给出的$(A^\mathrm{T}, B^\mathrm{T})$. 因此不能应用``分块矩阵初等变换不改变矩阵的秩''，秩不一定相等. 事实上我们很好举出例子来说明这一点，例如
              \[A=\begin{pmatrix}
                      1 & 0 \\ 0 & 0
                  \end{pmatrix},\enspace B=\begin{pmatrix}
                      0 & 0 \\ 1 & 0
                  \end{pmatrix}.\]

        \item 我们在矩阵乘法一节中介绍过，左乘矩阵相当于做初等行变换，右乘矩阵相当于做初等列变换，因此$AB$相当于对$A$做了初等列变换，考察$(A\enspace AB)$的列向量组，显然$AB$的列向量组是$A$的列向量组的线性组合，因此$r(A,AB)=r(A)$.

              当然我们也可以有另一种证明方式. 事实上直接看列向量组，显然$(A,AB)$的列向量组包括了$A$的列向量组，因此$r(A,AB)\geqslant r(A)$. 反过来，$r(A,AB)=r(A(E,B))\leqslant r(A)$，其中$E$是$n$阶单位矩阵，综上$r(A,AB)=r(A)$.

        \item 显然$r(A,B)\geqslant \max\{r(A), r(B)\}$，这只需要看矩阵列向量组的秩即可轻松看出.

        \item 可以将$BA$理解为对$A$做初等行变换，但行变换不能像B选项那样直接得出秩相关的结论，因为$A$和$BA$的行混在一起，无法判断二者行秩的关系.

              另一方面，B选项的另一种证明也不可使用，因为下列写法根据分块矩阵乘法块的匹配规则是不合法的：$(A,BA)=(E,B)A$，因为$(E,B)$是$1\times 2$分块，后面应该乘以一个两行分块.

              我们直接举反例否定这一结论，例如
              \[A=\begin{pmatrix}
                      1 & 0 \\ 0 & 0
                  \end{pmatrix},\enspace B=\begin{pmatrix}
                      1 & 1 \\ 1 & 1
                  \end{pmatrix}.\]
    \end{enumerate}
\end{solve}

\section{补充训练：行列式}

本节包含 79 道训练题，建议在完成相关专题真题后选择性练习。

\begin{exercise}[补充训练 L10-A-2]\label{prob:L10-A-2}
设$\alpha_1,\alpha_2,\alpha_3$为三维列向量，令$A=(\alpha_1,\alpha_2,\alpha_3)$，且$|A|=2$，求$|\alpha_1+\alpha_2+\alpha_3,\alpha_1+3\alpha_2+9\alpha_3,\alpha_1+4\alpha_2+16\alpha_3|$.
\end{exercise}
\noindent 来源：\href{https://github.com/yhwu-is/Linear-Algebra-Left-Undone/blob/new/%E8%AE%B2%E4%B9%89/%E4%B8%93%E9%A2%98/10%20%E8%A1%8C%E5%88%97%E5%BC%8F.tex}{L10-A-2}。
\begin{solve}
使用倍加列变换的性质，可得$|\alpha_1+\alpha_2+\alpha_3,\alpha_1+3\alpha_2+9\alpha_3,\alpha_1+4\alpha_2+16\alpha_3|=6|\alpha_1,\alpha_2,\alpha_3|=12$.
\end{solve}

\begin{exercise}[补充训练 L10-A-3]\label{prob:L10-A-3}
求证以下命题：
        \begin{enumerate}
            \item 奇数阶反对称矩阵不可逆；

            \item 若$A$是$n$阶可逆对称矩阵，$B$是$n$阶反对称矩阵，则当$n$为奇数时，齐次线性方程组$(AB)X=O$有非零解.
        \end{enumerate}
\end{exercise}
\noindent 来源：\href{https://github.com/yhwu-is/Linear-Algebra-Left-Undone/blob/new/%E8%AE%B2%E4%B9%89/%E4%B8%93%E9%A2%98/10%20%E8%A1%8C%E5%88%97%E5%BC%8F.tex}{L10-A-3}。
\begin{solve}
\begin{enumerate}
                \item 首先由反对称矩阵的性质，有$A^T=-A$，其次由矩阵与行列式的性质，可得$|A^T|=|A|$，则可推出$|A|=|-A|=(-1)^n|A|$，又n为奇数，故$|A|=0$，矩阵$A$不可逆.

                \item 只需证明$AB$的秩小于$n$，即$|AB|=0$即可. $B$是奇数阶反对称矩阵，$|B|=0$,则$|AB|=|A||B|=0$得证.
            \end{enumerate}
\end{solve}

\begin{exercise}[补充训练 L10-A-4]\label{prob:L10-A-4}
设 $A,B$ 分别为 $m,n$ 阶可逆矩阵，$a=\det A$、$b=\det B$。求分块矩阵
\[M=\begin{pmatrix}A&O\\O&B\end{pmatrix},\qquad N=\begin{pmatrix}O&A\\B&O\end{pmatrix}\]
的伴随矩阵。$N$ 中第一块列宽为 $n$，第二块列宽为 $m$。
\end{exercise}
\noindent 来源：\href{https://github.com/yhwu-is/Linear-Algebra-Left-Undone/blob/new/%E8%AE%B2%E4%B9%89/%E4%B8%93%E9%A2%98/10%20%E8%A1%8C%E5%88%97%E5%BC%8F.tex}{L10-A-4}。
\begin{solve}
由 $\det M=ab$ 及 $M^{-1}=\operatorname{diag}(A^{-1},B^{-1})$，得到
\[\operatorname{adj}(M)=\begin{pmatrix}b\operatorname{adj}(A)&O\\O&a\operatorname{adj}(B)\end{pmatrix}.\]
交换 $M$ 的两个块列，共需 $mn$ 次普通列交换，得到 $N$，故 $\det N=(-1)^{mn}ab$。又
\[N^{-1}=\begin{pmatrix}O&B^{-1}\\A^{-1}&O\end{pmatrix},\]
所以
\[\operatorname{adj}(N)=(-1)^{mn}\begin{pmatrix}O&a\operatorname{adj}(B)\\b\operatorname{adj}(A)&O\end{pmatrix}.\]
注意逆矩阵的块行高度依次为 $n,m$，与原矩阵的块列宽度对应。
\end{solve}

\begin{exercise}[补充训练 L10-A-5]\label{prob:L10-A-5}
证明：上（下）三角矩阵的伴随矩阵是上（下）三角矩阵（对角矩阵为特例）.
\end{exercise}
\noindent 来源：\href{https://github.com/yhwu-is/Linear-Algebra-Left-Undone/blob/new/%E8%AE%B2%E4%B9%89/%E4%B8%93%E9%A2%98/10%20%E8%A1%8C%E5%88%97%E5%BC%8F.tex}{L10-A-5}。
\begin{solve}
先设 $A$ 可逆且为上三角矩阵，则 $A^{-1}$ 也是上三角矩阵，故 $\operatorname{adj}(A)=\det(A)A^{-1}$ 为上三角矩阵。若 $A$ 奇异，考察 $A+tI$：除有限多个 $t$ 外，它都可逆且上三角，所以 $\operatorname{adj}(A+tI)$ 的严格下三角元为零。每个这样的元素都是关于 $t$ 的多项式，在无限多个值上为零，因而恒为零；令 $t=0$，结论仍成立。下三角情形通过转置得到，因为 $\operatorname{adj}(A^{\mathrm T})=\operatorname{adj}(A)^{\mathrm T}$。
\end{solve}

\begin{exercise}[补充训练 L10-A-6]\label{prob:L10-A-6}
设 $n\ge2$，$A$ 为 $n$ 阶方阵且 $\det A=0$。证明：代数余子式矩阵的任意两个行向量以及任意两个列向量均线性相关。若两个向量均非零，则它们互为非零标量倍。
\end{exercise}
\noindent 来源：\href{https://github.com/yhwu-is/Linear-Algebra-Left-Undone/blob/new/%E8%AE%B2%E4%B9%89/%E4%B8%93%E9%A2%98/10%20%E8%A1%8C%E5%88%97%E5%BC%8F.tex}{L10-A-6}。
\begin{solve}
若 $r(A)\le n-2$，则所有 $n-1$ 阶子式为零，伴随矩阵为零。若 $r(A)=n-1$，则伴随矩阵至少有一个非零元素；又 $A\operatorname{adj}(A)=O$，所以其每一列都属于一维空间 $\ker A$，从而 $r(\operatorname{adj}(A))=1$。代数余子式矩阵为伴随矩阵的转置，两者秩相同，均不超过 $1$，因此任意两个行向量或列向量线性相关。使用“线性相关”还涵盖一个向量为零、另一个非零的情形，避免“成比例”的歧义。
\end{solve}

\begin{exercise}[补充训练 L10-A-7]\label{prob:L10-A-7}
设向量$\alpha_1,\alpha_2,\alpha_3$线性无关，讨论向量$\alpha_1-\alpha_2-2\alpha_3,\ 2\alpha_1+\alpha_2-\alpha_3,\ 3\alpha_1+\alpha_2+2\alpha_3$的线性相关性.
\end{exercise}
\noindent 来源：\href{https://github.com/yhwu-is/Linear-Algebra-Left-Undone/blob/new/%E8%AE%B2%E4%B9%89/%E4%B8%93%E9%A2%98/10%20%E8%A1%8C%E5%88%97%E5%BC%8F.tex}{L10-A-7}。
\begin{solve}
\[(\alpha_1-\alpha_2-2\alpha_3,2\alpha_1+\alpha_2-\alpha_3,3\alpha_1+\alpha_2+2\alpha_3)=(\alpha_1,\alpha_2,\alpha_3)\begin{pmatrix}
                1 & 2 & 3 \\ -1 & 1 & 1 \\ -2 & -1 & 2
            \end{pmatrix}\]
，又$\begin{vmatrix}
                1 & 2 & 3 \\ -1 & 1 & 1 \\ -2 & -1 & 2
            \end{vmatrix}=12\neq 0$，因此该矩阵可逆，则$(\alpha_1-\alpha_2-2\alpha_3,2\alpha_1+\alpha_2-\alpha_3,3\alpha_1+\alpha_2+2\alpha_3)$与$(\alpha_1,\alpha_2,\alpha_3)$秩相等，因此这三个向量线性无关.
\end{solve}

\begin{exercise}[补充训练 L10-A-8]\label{prob:L10-A-8}
设$W=\operatorname{span}(\alpha_1,\alpha_2)$是$\mathbf{R}^4$的一个子空间，其中$\alpha_1=(1,2,1,-1)^\mathrm{T}$，$\alpha_2=(1,4,-1,-1)^\mathrm{T}$，试将$\alpha_1,\alpha_2$扩充为$\mathbf{R}^4$的基.
\end{exercise}
\noindent 来源：\href{https://github.com/yhwu-is/Linear-Algebra-Left-Undone/blob/new/%E8%AE%B2%E4%B9%89/%E4%B8%93%E9%A2%98/10%20%E8%A1%8C%E5%88%97%E5%BC%8F.tex}{L10-A-8}。
\begin{solve}
只要取$\alpha_3,\alpha_4 \in \mathbf{R}^4$，使得$D(\alpha_1,\alpha_2,\alpha_3,\alpha_4) \neq 0$即可. 易得$\begin{vmatrix}
                1  & 1  & 0 & 0 \\
                2  & 4  & 0 & 0 \\
                1  & -1 & 1 & 0 \\
                -1 & 1  & 0 & 1
            \end{vmatrix}=\begin{vmatrix}
                1 & 1 \\
                2 & 4
            \end{vmatrix} \cdot \begin{vmatrix}
                1 & 0 \\
                0 & 1
            \end{vmatrix}=2 \neq 0$. 于是可取$\alpha_3=(0,0,1,0)^T,\alpha_4=(0,0,0,1)^T$，能使得${\alpha_1,\alpha_2,\alpha_3,\alpha_4}$为$\mathbf{R}^4$的一组基.
\end{solve}

\begin{exercise}[补充训练 L10-A-9]\label{prob:L10-A-9}
$ D_n=\begin{vmatrix}
                1      & 2      & 3      & 4      & \cdots & n-1    & n      \\
                x      & 1      & 2      & 3      & \cdots & n-2    & n-1    \\
                x      & x      & 1      & 2      & \cdots & n-3    & n-2    \\
                x      & x      & x      & 1      & \cdots & n-4    & n-3    \\
                \vdots & \vdots & \vdots & \vdots & \ddots & \vdots & \vdots \\
                x      & x      & x      & x      & \cdots & 1      & 2      \\
                x      & x      & x      & x      & \cdots & x      & 1      \\
            \end{vmatrix}$（提示：考虑滚动消去法）.
\end{exercise}
\noindent 来源：\href{https://github.com/yhwu-is/Linear-Algebra-Left-Undone/blob/new/%E8%AE%B2%E4%B9%89/%E4%B8%93%E9%A2%98/10%20%E8%A1%8C%E5%88%97%E5%BC%8F.tex}{L10-A-9}。
\begin{solve}
考虑滚动消去法，每行减下一行，得
            \[ \begin{vmatrix}
                    1-x    & 1      & 1      & 1      & \ldots & 1      & 1      \\
                    0      & 1-x    & 1      & 1      & \ldots & 1      & 1      \\
                    0      & 0      & 1-x    & 1      & \ldots & 1      & 1      \\
                    0      & 0      & 0      & 1-x    & \ldots & 1      & 1      \\
                    \vdots & \vdots & \vdots & \vdots & \ddots & \vdots & \vdots \\
                    0      & 0      & 0      & 0      & \ldots & 1-x    & 1      \\
                    x      & x      & x      & x      & \ldots & x      & 1      \\
                \end{vmatrix}. \]
            第一行乘$(-x)$到第$n$行，得
            \[\begin{vmatrix}
                    1-x    & 1      & 1      & 1      & \ldots & 1      & 1      \\
                    0      & 1-x    & 1      & 1      & \ldots & 1      & 1      \\
                    0      & 0      & 1-x    & 1      & \ldots & 1      & 1      \\
                    0      & 0      & 0      & 1-x    & \ldots & 1      & 1      \\
                    \vdots & \vdots & \vdots & \vdots & \ddots & \vdots & \vdots \\
                    0      & 0      & 0      & 0      & \ldots & 1-x    & 1      \\
                    x^2    & 0      & 0      & 0      & \ldots & 0      & 1-x    \\
                \end{vmatrix}.\]
            按第$n$行展开，得$D_n=(1-x)^n+(-1)^{n+1}x^2T_{n-1}$，其中
            \[T_{n-1}=\begin{vmatrix}
                    1      & 1      & 1      & \ldots & 1      & 1      \\
                    1-x    & 1      & 1      & \ldots & 1      & 1      \\
                    0      & 1-x    & 1      & \ldots & 1      & 1      \\
                    \vdots & \vdots & \vdots & \ddots & \vdots & \vdots \\
                    0      & 0      & 0      & \ldots & 1      & 1      \\
                    0      & 0      & 0      & \ldots & 1-x    & 1      \\
                \end{vmatrix}.\]
            按第$n-1$行展开$T_{n-1}$，有$T_{n-1}=T_{n-2}+(-1)(1-x)T_{n-2}=xT_{n-2}$. 由$T_2=\begin{vmatrix}
                    1&1\\1-x&1\end{vmatrix}=x$，递推得$T_{n-1}=x^{n-2}$. 代入可解得$D_n=(1-x)^n-(-x)^n$.
\end{solve}

\begin{exercise}[补充训练 L10-A-10]\label{prob:L10-A-10}
$ D_n=\begin{vmatrix}
                a_1 & b_1 &     &        &         \\
                    & a_2 & b_2 &        &         \\
                    &     & a_3 & \ddots &         \\
                    &     &     & \ddots & b_{n-1} \\
                b_n &     &     &        & a_n
            \end{vmatrix}$.
\end{exercise}
\noindent 来源：\href{https://github.com/yhwu-is/Linear-Algebra-Left-Undone/blob/new/%E8%AE%B2%E4%B9%89/%E4%B8%93%E9%A2%98/10%20%E8%A1%8C%E5%88%97%E5%BC%8F.tex}{L10-A-10}。
\begin{solve}
按第一列展开，得$D_n=\displaystyle\prod_{i=1}^na_i+(-1)^{n+1}\prod_{i=1}^nb_i$.
\end{solve}

\begin{exercise}[补充训练 L10-A-11]\label{prob:L10-A-11}
$D_n=\begin{vmatrix}
                a+b & ab  &        &     &     \\
                1   & a+b & \ddots &     &     \\
                    & 1   & \ddots & ab  &     \\
                    &     & \ddots & a+b & ab  \\
                    &     &        & 1   & a+b \\
            \end{vmatrix}$.
\end{exercise}
\noindent 来源：\href{https://github.com/yhwu-is/Linear-Algebra-Left-Undone/blob/new/%E8%AE%B2%E4%B9%89/%E4%B8%93%E9%A2%98/10%20%E8%A1%8C%E5%88%97%E5%BC%8F.tex}{L10-A-11}。
\begin{solve}
按第一列展开，得
            \begin{align*}
                D_{n} & =(a+b) D_{n-1}-\begin{vmatrix}
                                            a b    & 0      & 0      & 0   & \cdots & 0      & 0      \\
                                            1      & a+b    & a b    & 0   & \cdots & 0      & 0      \\
                                            0      & 1      & a+b    & a b & \cdots & 0      & 0      \\
                                            \vdots & \vdots & \vdots & a+b & \ddots & \vdots & \vdots \\
                                            0      & 0      & 0      & 0   & \cdots & a+b    & a b    \\
                                            0      & 0      & 0      & 0   & \cdots & 1      & a+b
                                        \end{vmatrix} \\
                        & =(a+b) D_{n-1}-a b D_{n-2}.
            \end{align*}
            \begin{enumerate}
                \item 方法一：特征方程法. 有特征方程$r^2-(a+b)r+ab=(r-a)(r-b)=0$，先考虑$a\neq b$，则$D_n=C_1a^n+C_2b^n$，$D_{1}=a+b$，$D_{2}=a^{2}+a b+b^{2}$，则有
                        \[\begin{cases}
                                C_1a+C_2b=a+b \\
                                C_1a^2+C_2b^2=a^2+ab+b^2
                            \end{cases},\]
                        考虑Cramer法则，首先判断系数行列式$\begin{vmatrix}
                                a   & b   \\
                                a^2 & b^2
                            \end{vmatrix}=ab(b-a)$，若$a=b=0$，原行列式为0；若$a=0,b\neq 0$，则有$C_2=1$，$D_n=b^n$；同理若$a=0,b\neq 0$，$D_n=a^n$. 考虑$a\neq 0,b\neq 0$，则系数行列式非0，根据Cramer法则可以解得
                        \[C_1=\frac{\begin{vmatrix}
                                    a+b        & b   \\
                                    a^2+ab+b^2 & b^2
                                \end{vmatrix}}{ab(b-a)}=\frac{a}{a-b},\quad
                            C_2=\frac{\begin{vmatrix}
                                    a   & a+b        \\
                                    a^2 & a^2+ab+b^2
                                \end{vmatrix}}{ab(b-a)}=-\frac{b}{a-b}.\]
                        故$D_n=\displaystyle\frac{a^{n+1}-b^{n+1}}{a-b}$. 考虑到$a$或$b$为0时也符合该公式，故合并为同一公式. 再考虑$a=b$，则$D_n=(C_1+C_2n)a^n$，有方程
                        \[\begin{cases}
                                (C_1+C_2)a=2a \\
                                (C_1+2C_2)a^2=3a^2
                            \end{cases},\]
                        $a=0$，则$D_n=0$；$a\neq 0$，可以解得$C_1=C_2=1$，有$D_n=(n+1)a^n$. 考虑到$a=0$时$D_n$也符合该式，故合并. 综上有
                        \[D_n=\begin{cases}
                                \displaystyle\frac{a^{n+1}-b^{n+1}}{a-b}, & a\neq b \\
                                (n+1)a^n,                                 & a=b
                            \end{cases}.\]

                \item 方法二：递推式变形法. 递推式变形, 得
                        \[\mathrm{D}_{n}-a D_{n-1}=b\left(D_{n-1}-a D_{n-2}\right),\]
                        由于 $D_{1}=a+b, D_{2}=a^{2}+a b+b^{2}$,从而利用上述递推公式得
                        \begin{align*}
                            \mathrm{D}_{n}-a D_{n-1} & =b\left(D_{n-1}-a D_{n-2}\right)                \\
                                                    & =b^{2}\left(D_{n-2}-a D_{n-3}\right)            \\
                                                    & =\cdots=b^{n-2}\left(D_{2}-a D_{1}\right)=b^{n}.
                        \end{align*}
                        故\begin{align*}
                            D_{n} & =a D_{n-1}+b^{n}=a\left(a D_{n-2}+b^{n-1}\right)+b^{n}     \\
                                  & =\cdots=a^{n-1} D_{1}+a^{n-2} b^{2}+\cdots+a b^{n-1}+b^{n} \\
                                  & =a^{n}+a^{n-1} b+\cdots+a b^{n-1}+b^{n}                    \\
                                  & =\begin{cases}
                                        \displaystyle\frac{a^{n+1}-b^{n+1}}{a-b}, & a\neq b \\
                                        (n+1)a^n,                                 & a=b
                                     \end{cases}.
                        \end{align*}
            \end{enumerate}
\end{solve}

\begin{exercise}[补充训练 L10-A-13]\label{prob:L10-A-13}
解行列式 $ \begin{vmatrix}
                a^{2} & (a+1)^{2} & (a+2)^{2} & (a+3)^{2} \\
                b^{2} & (b+1)^{2} & (b+2)^{2} & (b+3)^{2} \\
                c^{2} & (c+1)^{2} & (c+2)^{2} & (c+3)^{2} \\
                d^{2} & (d+1)^{2} & (d+2)^{2} & (d+3)^{2}
            \end{vmatrix} $.
\end{exercise}
\noindent 来源：\href{https://github.com/yhwu-is/Linear-Algebra-Left-Undone/blob/new/%E8%AE%B2%E4%B9%89/%E4%B8%93%E9%A2%98/10%20%E8%A1%8C%E5%88%97%E5%BC%8F.tex}{L10-A-13}。
\begin{solve}
\begin{enumerate}
                \item 方法一：慢慢消去.
                      \begin{align*}
                          \begin{vmatrix}
                              a^{2} & (a+1)^{2} & (a+2)^{2} & (a+3)^{2} \\
                              b^{2} & (b+1)^{2} & (b+2)^{2} & (b+3)^{2} \\
                              c^{2} & (c+1)^{2} & (c+2)^{2} & (c+3)^{2} \\
                              d^{2} & (d+1)^{2} & (d+2)^{2} & (d+3)^{2}
                          \end{vmatrix}
                           & = \begin{vmatrix}
                                   a^{2} & 2a+1 & 4a+4 & 6a+9 \\
                                   b^{2} & 2b+1 & 4b+4 & 6b+9 \\
                                   c^{2} & 2c+1 & 4c+4 & 6c+9 \\
                                   d^{2} & 2d+1 & 4d+4 & 6d+9
                               \end{vmatrix} \\
                           & = \begin{vmatrix}
                                   a^{2} & 2a+1 & 4a+4 & 4 \\
                                   b^{2} & 2b+1 & 4b+4 & 4 \\
                                   c^{2} & 2c+1 & 4c+4 & 4 \\
                                   d^{2} & 2d+1 & 4d+4 & 4
                               \end{vmatrix}    \\
                           & = 4\begin{vmatrix}
                                    a^{2} & 2a & 4a & 1 \\
                                    b^{2} & 2b & 4b & 1 \\
                                    c^{2} & 2c & 4c & 1 \\
                                    d^{2} & 2d & 4d & 1
                                \end{vmatrix}=0.
                      \end{align*}

                \item 关注到行列式中只有$(a^2, b^2, c^2, d^2)^\mathrm{T}, (a, b, c, d)^\mathrm{T}$和$(1, 1, 1, 1)^\mathrm{T}$三个线性无关列向量，则一眼看出线性相关. 因此设原行列式为$|\alpha_1, \alpha_2, \alpha_3, \alpha_4|$，有$\alpha_4=\alpha_1-3\alpha_2+3\alpha_3$，因此原行列式值为0.
            \end{enumerate}
\end{solve}

\begin{exercise}[补充训练 L10-A-14]\label{prob:L10-A-14}
设
        \[ D=\begin{vmatrix}
                1+a_1  & 1      & \cdots & 1      \\
                1      & 1+a_2  & \cdots & 1      \\
                \vdots & \vdots & \ddots & \vdots \\
                1      & 1      & \cdots & 1+a_n
            \end{vmatrix} \]
        \begin{enumerate}
            \item 用递推公式计算行列式$D$；

            \item 硬拆$D$为$2^n$个行列式，计算出结果.
        \end{enumerate}
\end{exercise}
\noindent 来源：\href{https://github.com/yhwu-is/Linear-Algebra-Left-Undone/blob/new/%E8%AE%B2%E4%B9%89/%E4%B8%93%E9%A2%98/10%20%E8%A1%8C%E5%88%97%E5%BC%8F.tex}{L10-A-14}。
\begin{solve}
\begin{enumerate}
                \item \begin{align*}
                          D_n & =\begin{vmatrix}
                                     D_{n-1} & a   \\
                                     b       & a_n
                                 \end{vmatrix},a=(0, \ldots, 0, -a_{n-1})^\mathrm{T}, b=(1, \ldots, 1) \\
                              & =a_nD_{n-1}+a_{n-1}\begin{vmatrix}
                                                       D_{n-2} & 1_{n-2} \\
                                                       1_{n-2} & 1
                                                   \end{vmatrix}
                          \quad \text{($1_{n-2}$表示由$(n-2)$个1构成的行/列向量)}                      \\
                              & =a_nD_{n-1}+a_{n-1}
                          \begin{vmatrix}
                              a_1    & 0      & 0      & \ldots & 0       & 1      \\
                              0      & a_2    & 0      & \ldots & 0       & 1      \\
                              0      & 0      & a_3    & \ldots & 0       & 1      \\
                              \vdots & \vdots & \vdots & \ddots & \vdots  & \vdots \\
                              0      & 0      & 0      & \ldots & a_{n-2} & 1      \\
                              0      & 0      & 0      & \ldots & 0       & 1      \\
                          \end{vmatrix} \quad \text{($1-n$列用第$n$列减)}                         \\
                              & =a_nD_{n-1}+\prod_{i=1}^{n-1}a_i.
                      \end{align*}
                      递推可得$\displaystyle D_n=\left(\prod_{k=1}^na_k\right)
                          \left(1+\sum_{i=1}^n\frac{1}{a_i}\right)$.

                \item 将 1 写成 $1+0$, 将 $D$ 硬拆成 $2^{n}$ 个行列式, 只有如下的 $n+1$ 个行列式非 0：
                      \begin{align*}
                          D_n = & \begin{vmatrix}
                                      1      &       &        &       \\
                                      1      & a_{2} &        &       \\
                                      \vdots &       & \ddots &       \\
                                      1      &       &        & a_{n}
                                  \end{vmatrix}
                          +\begin{vmatrix}
                               a_{1} & 1      &                \\
                                     & 1      &                \\
                                     & \vdots & \ddots &       \\
                                     & 1      &        & a_{n}
                           \end{vmatrix}+\cdots+                    \\
                                & \begin{vmatrix}
                                      a_{1} &       &        & 1      \\
                                            & a_{2} &        & 1      \\
                                            &       & \ddots & \vdots \\
                                            &       &        & 1
                                  \end{vmatrix}
                          +\begin{vmatrix}
                               a_{1} &                      \\
                                     & a_{2} &              \\
                                     &       & \ddots &     \\
                                     &       &        & a_n
                           \end{vmatrix}                       \\
                          =     & \left(\prod_{k=1}^na_k\right)
                          \left(1+\sum_{i=1}^n\frac{1}{a_i}\right).
                      \end{align*}
            \end{enumerate}
\end{solve}

\begin{exercise}[补充训练 L10-A-15]\label{prob:L10-A-15}
解行列式
        \begin{enumerate}
            \item $D=\begin{vmatrix}
                          1      & 2      & \cdots & 2      & 2      \\
                          2      & 2      & \cdots & 2      & 2      \\
                          \vdots & \vdots & \ddots & \vdots & \vdots \\
                          2      & 2      & \cdots & n-1    & 2      \\
                          2      & 2      & \cdots & 2      & n\end{vmatrix}$；

            \item $D=\begin{vmatrix}
                          1      & 2      & \cdots & n-1    & n      \\
                          2      & 3      & \cdots & n      & 1      \\
                          3      & 4      & \cdots & 1      & 2      \\
                          \vdots & \vdots & \ddots & \vdots & \vdots \\
                          n      & 1      & \cdots & n-2    & n-1
                      \end{vmatrix}$.
        \end{enumerate}
\end{exercise}
\noindent 来源：\href{https://github.com/yhwu-is/Linear-Algebra-Left-Undone/blob/new/%E8%AE%B2%E4%B9%89/%E4%B8%93%E9%A2%98/10%20%E8%A1%8C%E5%88%97%E5%BC%8F.tex}{L10-A-15}。
\begin{solve}
\begin{enumerate}
                \item 行列式中2比较多，用全是2的第二行去消，然后按第2行展开.
                      \[D=\begin{vmatrix}
                              -1     & 0      & \cdots & 0      & 0      \\
                              2      & 2      & \cdots & 2      & 2      \\
                              \vdots & \vdots & \ddots & \vdots & \vdots \\
                              0      & 0      & \cdots & n-3    & 0      \\
                              0      & 0      & \cdots & 0      & n-2
                          \end{vmatrix}=2\cdot(-1)\cdot(n-2)!=-2(n-2)!.\]

                \item \begin{enumerate}
                          \item 方法一：考虑滚动消去法，有
                                \begin{align*}
                                    D & =\begin{vmatrix}
                                             1      & 2      & 3      & \cdots & n-1    & n      \\
                                             1      & 1      & 1      & \cdots & 1      & 1-n    \\
                                             1      & 1      & 1      & \cdots & 1-n    & 1      \\
                                             \vdots & \vdots & \vdots & \ddots & \vdots & \vdots \\
                                             1      & 1      & 1-n    & \cdots & 1      & 1      \\
                                             1      & 1-n    & 1      & \cdots & 1      & 1
                                         \end{vmatrix}                                    \\&=\begin{vmatrix}
                                        0      & n+1    & 2      & \cdots & n-2    & n-1    \\
                                        0      & n      & 0      & \cdots & 0      & -n     \\
                                        0      & n      & 0      & \cdots & -n     & 0      \\
                                        \vdots & \vdots & \vdots & \ddots & \vdots & \vdots \\
                                        0      & n      & -n     & \cdots & 0      & 0      \\
                                        1      & 1-n    & 1      & \cdots & 1      & 1
                                    \end{vmatrix}\\
                                      & =(-1)^{n+1}\begin{vmatrix}
                                                       n+1    & 2      & \cdots & n-2    & n-1    \\
                                                       n      & 0      & \cdots & 0      & -n     \\
                                                       n      & 0      & \cdots & -n     & 0      \\
                                                       \vdots & \vdots & \ddots & \vdots & \vdots \\
                                                       n      & -n     & \cdots & 0      & 0
                                                   \end{vmatrix}                                   \\
                                      & =(-1)^{n+1}(-1)^{\frac{(n-2)(n-3)}{2}}
                                    \begin{vmatrix}
                                        n+1    & n-1    & \cdots & 3      & 2      \\
                                        n      & -n     & \cdots & 0      & 0      \\
                                        n      & 0      & \cdots & 0      & 0      \\
                                        \vdots & \vdots & \ddots & \vdots & \vdots \\
                                        n      & 0      & \cdots & 0      & -n
                                    \end{vmatrix}                                                  \\
                                      & =(-1)^{\frac{n^2-3n+8}{2}}\begin{vmatrix}
                                                                      \frac{n(n+1)}{2} & n-1    & \cdots & 3      & 2      \\
                                                                      0                & -n     & \cdots & 0      & 0      \\
                                                                      0                & 0      & \cdots & 0      & 0      \\
                                                                      \vdots           & \vdots & \ddots & \vdots & \vdots \\
                                                                      0                & 0      & \cdots & 0      & -n
                                                                  \end{vmatrix}=(-1)^{\frac{n^2-n+12}{2}}\frac{n^{n-1}(n+1)}{2}.
                                \end{align*}

                          \item 方法二：考虑连加法. 将第 $2,3, \ldots, n$ 列都加到第 1 列, 提出公因子 $\dfrac{1}{2} n(n+1)$, 再依次将第 $n-1$ 行乘 $(-1)$ 加到第 $n$ 行, $\cdots$, 第 2 行乘 $(-1)$ 加到第 3 行, 第 1 行乘 $(-1)$ 加到第 2 行, 然后对第 1 列展开, 得到一个 $n-1$ 阶行列式, 它的副对角元为 $1-$ $n$, 其余元素均为 1 . 再把它的各列加到第 1 列, 并把它的第 1 行乘 $(-1)$ 加到其余各行, 得
                                \begin{align*}
                                    D & =\frac{n(n+1)}{2}\begin{vmatrix}
                                                             1      & 2      & \cdots & n-1    & n      \\
                                                             1      & 3      & \cdots & n      & 1      \\
                                                             1      & 4      & \cdots & 1      & 2      \\
                                                             \vdots & \vdots & \ddots & \vdots & \vdots \\
                                                             1      & 1      & \cdots & n-2    & n-1
                                                         \end{vmatrix} \\
                                      & =\frac{n(n+1)}{2}\begin{vmatrix}
                                                             1      & 2      & 3      & \cdots & n      \\
                                                             0      & 1      & 1      & \cdots & 1-n    \\
                                                             \vdots & \vdots & \vdots & \ddots & \vdots \\
                                                             0      & 1      & 1-n    & \cdots & 1      \\
                                                             0      & 1-n    & 1      & \cdots & 1
                                                         \end{vmatrix}_{n} \\
                                      & =\frac{n(n+1)}{2}
                                    \begin{vmatrix}
                                        1      & 1      & \cdots & 1      & 1-n \\
                                        \vdots & \vdots & \ddots & \vdots & 1   \\
                                        1      & 1-n    & \cdots & 1      & 1   \\
                                        1-n    & 1      & \cdots & 1      & 1
                                    \end{vmatrix}                         \\
                                      & =\frac{n(n+1)}{2}\begin{vmatrix}
                                                             -1     & 1      & \cdots & 1      & 1-n    \\
                                                             0      & 0      & \cdots & -n     & n      \\
                                                             \vdots & \vdots & \ddots & \vdots & \vdots \\
                                                             0      & -n     & \cdots & 0      & n      \\
                                                             0      & 0      & \cdots & 0      & n
                                                         \end{vmatrix}_{n-1}
                                \end{align*}
                                将上式先对第 1 列展开, 得到一个 $n-2$ 阶行列式, 再将它对最后一行展开, 得
                                \begin{align*}
                                    D & =\frac{-n(n+1)}{2} n
                                    \begin{vmatrix}
                                        0      & 0      & \cdots & 0      & -n     \\
                                        0      & 0      & \cdots & -n     & 0      \\
                                        \vdots & \vdots & \ddots & \vdots & \vdots \\
                                        0      & -n     & \cdots & 0      & 0      \\
                                        -n     & 0      & \cdots & 0      & 0
                                    \end{vmatrix}_{n-3}                      \\
                                      & =-\frac{n^{2}(n+1)}{2}(-1)^{\frac{(n-3)(n-4)}{2}}(-n)^{n-3} \\
                                      & =(-1)^{\frac{n(n-1)}{2} }\frac{(n+1) n^{n-1}}{2}.
                                \end{align*}
                                事实上，两种方法得到的答案是等价的.
                      \end{enumerate}
            \end{enumerate}
\end{solve}

\begin{exercise}[补充训练 L10-A-16]\label{prob:L10-A-16}
记$a_{ij}=i-j$，且$A=(a_{ij})_{n\times n}$，计算$A$的行列式.
\end{exercise}
\noindent 来源：\href{https://github.com/yhwu-is/Linear-Algebra-Left-Undone/blob/new/%E8%AE%B2%E4%B9%89/%E4%B8%93%E9%A2%98/10%20%E8%A1%8C%E5%88%97%E5%BC%8F.tex}{L10-A-16}。
\begin{solve}
\begin{enumerate}
                \item 若 $n=1$，则 $|A| = |0| = 0$；

                \item 若 $n=2$, 则 $|A| = \begin{vmatrix} 0 & -1 \\ 1 & 0 \end{vmatrix} = 1$；

                \item 若 $n>2$，考虑使用滚动消去法：
                    \begin{align*}
                        |A| &= \begin{vmatrix}
                            0      & -1     & -2     & \cdots & -n+1   \\
                            1      & 0      & -1     & \cdots & -n+2   \\
                            2      & 1      & 0      & \cdots & -n+3   \\
                            \vdots & \vdots & \vdots & \ddots & \vdots \\
                            n-1    & n-2    & n-3    & \cdots & 0
                        \end{vmatrix} \\
                        &= \begin{vmatrix}
                            1 & 1 & \cdots & 1 & -n+1 \\
                            1 & 1 & \cdots & 1 & -n+2 \\
                            \vdots & \vdots & \ddots & \vdots & \vdots \\
                            1 & 1 & \cdots & 1 & 0
                        \end{vmatrix} \\
                        &= \begin{vmatrix}
                            1 & 0 & \cdots & 0 & -n+1 \\
                            1 & 0 & \cdots & 0 & -n+2 \\
                            \vdots & \vdots & \ddots & \vdots & \vdots \\
                            1 & 0 & \cdots & 0 & 0
                        \end{vmatrix}
                        = 0.
                    \end{align*}
            \end{enumerate}
\end{solve}

\begin{exercise}[补充训练 L10-A-17]\label{prob:L10-A-17}
计算行列式$D=\begin{vmatrix}
                1 & 0 & 2 & 0 & x   \\
                0 & 2 & 0 & 1 & 0   \\
                1 & 0 & 4 & 0 & x^2 \\
                0 & 5 & 0 & 3 & 0   \\
                1 & 0 & 8 & 0 & x^3
            \end{vmatrix}$.
\end{exercise}
\noindent 来源：\href{https://github.com/yhwu-is/Linear-Algebra-Left-Undone/blob/new/%E8%AE%B2%E4%B9%89/%E4%B8%93%E9%A2%98/10%20%E8%A1%8C%E5%88%97%E5%BC%8F.tex}{L10-A-17}。
\begin{solve}
我们可以先构造分块对角矩阵，然后分别处理两个低阶的行列式：
            \[
                D = (-1)^{3+3} \begin{vmatrix}
                    1 & 2 & x   & 0 & 0 \\
                    1 & 4 & x^2 & 0 & 0 \\
                    1 & 8 & x^3 & 0 & 0 \\
                    0 & 0 & 0   & 2 & 1 \\
                    0 & 0 & 0   & 5 & 3
                \end{vmatrix} = (-1)^{3+3} \begin{vmatrix}
                    1 & 2 & x   \\
                    1 & 4 & x^2 \\
                    1 & 8 & x^3
                \end{vmatrix} \begin{vmatrix}
                    2 & 1 \\
                    5 & 3
                \end{vmatrix} = \begin{vmatrix}
                    1 & 2 & x   \\
                    1 & 4 & x^2 \\
                    1 & 8 & x^3
                \end{vmatrix}.
            \]

            该行列式可以通过加边法构造范德蒙行列式：
            \[
                D = \begin{vmatrix}
                    1 & 2 & x   \\
                    1 & 4 & x^2 \\
                    1 & 8 & x^3
                \end{vmatrix} = \begin{vmatrix}
                    1 & 1 & 1 & 1   \\
                    0 & 1 & 2 & x   \\
                    0 & 1 & 4 & x^2 \\
                    0 & 1 & 8 & x^3
                \end{vmatrix} = 2x(x-1)(x-2).
            \]
\end{solve}

\begin{exercise}[补充训练 L10-B-1]\label{prob:L10-B-1}
设$D=\begin{vmatrix}
                3 & 0 & 4 & 1 \\ 2 & 3 & 1 & 4 \\ 0 & -7 & 8 & 3 \\ 5 & 3 & -2 & 2
            \end{vmatrix}$，求
        \begin{enumerate}
            \item $A_{21}+A_{22}+A_{23}+A_{24}$；

            \item $A_{31}+A_{33}$；

            \item $M_{41}+M_{42}+M_{43}+M_{44}$.
        \end{enumerate}
\end{exercise}
\noindent 来源：\href{https://github.com/yhwu-is/Linear-Algebra-Left-Undone/blob/new/%E8%AE%B2%E4%B9%89/%E4%B8%93%E9%A2%98/10%20%E8%A1%8C%E5%88%97%E5%BC%8F.tex}{L10-B-1}。
\begin{solve}
这三小问实际上都是对原行列式中的某一行进行了替换进行计算.
            \begin{enumerate}
                \item \begin{align*}
                            A_{21}+A_{22}+A_{23}+A_{24}
                            & = \begin{vmatrix}
                                    3 & 0  & 4  & 1 \\
                                    1 & 1  & 1  & 1 \\
                                    0 & -7 & 8  & 3 \\
                                    5 & 3  & -2 & 2
                                \end{vmatrix}
                            = \begin{vmatrix}
                                3 & -3 & 1  & -2 \\
                                1 & 0  & 0  & 0  \\
                                0 & -7 & 8  & 3  \\
                                5 & -2 & -7 & -3
                            \end{vmatrix}             \\
                            & = (-1)^{2+1} \begin{vmatrix}
                                                -3 & 1  & -1 \\
                                                -7 & 8  & 3  \\
                                                -2 & -7 & -3
                                            \end{vmatrix} \\
                            & = 148.
                        \end{align*}

                \item \begin{align*}
                            A_{31}+A_{33}
                            & = 1A_{31}+0A_{32}+1A_{33}+0A_{34}
                            = \begin{vmatrix}
                                3 & 0 & 4  & 1 \\
                                2 & 3 & 1  & 4 \\
                                1 & 0 & 1  & 0 \\
                                5 & 3 & -2 & 2
                            \end{vmatrix}
                            = 3 \begin{vmatrix}
                                    3 & 0 & 4  & 1 \\
                                    2 & 1 & 1  & 4 \\
                                    1 & 0 & 1  & 0 \\
                                    5 & 1 & -2 & 2
                                \end{vmatrix}                   \\
                            & = 3 \begin{vmatrix}
                                    3 & 0 & 1  & 1 \\
                                    2 & 1 & -1 & 4 \\
                                    1 & 0 & 0  & 0 \\
                                    5 & 3 & -7 & 2
                                \end{vmatrix}
                            = 3 \cdot (-1)^{3+1} \begin{vmatrix}
                                                    0 & 1  & 1 \\
                                                    1 & -1 & 4 \\
                                                    3 & -7 & 2
                                                \end{vmatrix}  \\
                            & = -12.
                        \end{align*}

                \item \begin{align*}
                            M_{41}+M_{42}+M_{43}+M_{44}
                            & = -A_{41}+A_{42}-A_{43}+A_{44}
                            = \begin{vmatrix}
                                3  & 0  & 4  & 1 \\
                                2  & 3  & 1  & 4 \\
                                0  & -7 & 8  & 3 \\
                                -1 & 1  & -1 & 1
                            \end{vmatrix}                     \\
                            & = \begin{vmatrix}
                                    3  & 3  & 1  & 4 \\
                                    2  & 5  & -1 & 6 \\
                                    0  & -7 & 8  & 3 \\
                                    -1 & 0  & 0  & 0
                                \end{vmatrix}
                            = {-1}^{4+1} \cdot (-1) \begin{vmatrix}
                                                        3  & 1  & 4 \\
                                                        5  & -1 & 6 \\
                                                        -7 & 8  & 3
                                                    \end{vmatrix} \\
                            & = -78.
                        \end{align*}
            \end{enumerate}
\end{solve}

\begin{exercise}[补充训练 L10-B-3]\label{prob:L10-B-3}
设$A,B$为三阶矩阵，且$|A|=3,|B|=2$，且$|A^{-1}+B|=2$，求$|A+B^{-1}|$.
\end{exercise}
\noindent 来源：\href{https://github.com/yhwu-is/Linear-Algebra-Left-Undone/blob/new/%E8%AE%B2%E4%B9%89/%E4%B8%93%E9%A2%98/10%20%E8%A1%8C%E5%88%97%E5%BC%8F.tex}{L10-B-3}。
\begin{solve}
\begin{align*}
                \lvert A+B^{-1} \rvert ={} & \lvert B^{-1}BA+B^{-1}E \rvert = \lvert B^{-1} \rvert \cdot \lvert BA+E \rvert                  \\
                ={}                        & \frac{1}{2} \lvert BA+A^{-1}A \rvert = \frac{1}{2} \lvert B+A^{-1} \rvert \cdot \lvert A \rvert \\
                ={}                        & \frac{3}{2} \lvert A^{-1}+B \rvert = 3.
            \end{align*}
\end{solve}

\begin{exercise}[补充训练 L10-B-4]\label{prob:L10-B-4}
设$A$为$n$阶正交矩阵，即$AA^\mathrm{T}=A^\mathrm{T}A=E$，且$|A|<0$，证明：$|E+A|=0$.
\end{exercise}
\noindent 来源：\href{https://github.com/yhwu-is/Linear-Algebra-Left-Undone/blob/new/%E8%AE%B2%E4%B9%89/%E4%B8%93%E9%A2%98/10%20%E8%A1%8C%E5%88%97%E5%BC%8F.tex}{L10-B-4}。
\begin{solve}
正交矩阵满足 $AA^{\mathrm{T}} = A^{\mathrm{T}}A = E$，所以 $\lvert AA^{\mathrm{T}} \rvert = \lvert A \rvert^2 = \lvert E \rvert = 1$. 而 $\lvert A \rvert < 0$，所以 $\lvert A \rvert = -1$.
            \[\lvert E+A \rvert = \lvert AA^{\mathrm{T}}+A \rvert = \lvert A \rvert \cdot \lvert A^{\mathrm{T}}+E \rvert = -\lvert (A+E)^{\mathrm{T}} \rvert = -\lvert E+A \rvert.\]
            故 $\lvert E+A \rvert = 0$.

            以上两道题都多次运用了 $E = AA^{-1} = A^{-1}A$ 的技巧，请大家留意.
\end{solve}

\begin{exercise}[补充训练 L10-B-5]\label{prob:L10-B-5}
设实方阵$A$的伴随矩阵$\operatorname{adj}(A)=\begin{pmatrix}
                1 & 0 & 0 \\ 0 & 1 & 0 \\ 1 & 0 & 1
            \end{pmatrix}$，且$|A|>0$，已知矩阵$B$满足$AB=E+3A$，求矩阵$B$.
\end{exercise}
\noindent 来源：\href{https://github.com/yhwu-is/Linear-Algebra-Left-Undone/blob/new/%E8%AE%B2%E4%B9%89/%E4%B8%93%E9%A2%98/10%20%E8%A1%8C%E5%88%97%E5%BC%8F.tex}{L10-B-5}。
\begin{solve}
由 $A$ 为三阶矩阵，可得 $|A|^2 = |\operatorname{adj}(A)| = 1$，而 $|A|>0$，所以 $|A|=1$. $AB=E+3A$ 经变形可得 $A(B-3E) = E$，故有
            \[
                B = 3E + A^{-1} = 3E + \frac{\operatorname{adj}(A)}{|A|} = \begin{pmatrix}
                    4 & 0 & 0 \\ 0 & 4 & 0 \\ 1 & 0 & 4
                \end{pmatrix}.
            \]
\end{solve}

\begin{exercise}[补充训练 L10-B-6]\label{prob:L10-B-6}
设$B=\begin{pmatrix}
                1 & 1 & 1 \\ 0 & 1 & 1 \\ 0 & 0 & 1
            \end{pmatrix}$，矩阵$A$满足方程$BA=B^\mathrm{T}B$，求$|\operatorname{adj}(A)-2E|$.
\end{exercise}
\noindent 来源：\href{https://github.com/yhwu-is/Linear-Algebra-Left-Undone/blob/new/%E8%AE%B2%E4%B9%89/%E4%B8%93%E9%A2%98/10%20%E8%A1%8C%E5%88%97%E5%BC%8F.tex}{L10-B-6}。
\begin{solve}
$|B|=1$，对 $BA=B^\mathrm{T}B$ 两边同时取行列式有
            \[
                |B||A| = |B^\mathrm{T}||B| = |B|^2,
            \]
            故 $|A| = 1$. 从而有
            \[
                \operatorname{adj}(A) = |A|A^{-1} = B^{-1} \left(B^\mathrm{T}\right)^{-1} B.
            \]
            $B^{-1} = \begin{pmatrix}
                1 & 0 & 0 \\ -1 & 1 & 0 \\ 0 & -1 & 1
            \end{pmatrix}$，于是有
            \[
                |\operatorname{adj}(A)-2E| = \lvert B^{-1} \left(B^\mathrm{T}\right)^{-1} B - 2 B^{-1} B \rvert
                         = \lvert B^{-1} \rvert \lvert \left(B^\mathrm{T}\right)^{-1} - 2E \rvert \lvert B \rvert
                         = \begin{vmatrix}
                            -1 & 0 & 0 \\ -1 & -1 & 0 \\ 0 & -1 & -1
                         \end{vmatrix}
                         = -1.
            \]
\end{solve}

\begin{exercise}[补充训练 L10-B-7]\label{prob:L10-B-7}
已知齐次线性方程组
        \[\begin{cases} \begin{aligned}
                    a_{11}x_1+a_{12}x_2+\cdots+a_{1n}x_n          & = 0             \\
                    a_{21}x_1+a_{22}x_2+\cdots+a_{2n}x_n          & = 0             \\
                                                                  & \vdotswithin{=} \\
                    a_{n-1,1}x_1+a_{n-1,2}x_2+\cdots+a_{n-1,n}x_n & = 0
                \end{aligned} \end{cases}\]
        设$M_j\enspace(j=1,2,\ldots,n)$表示$A=(a_{ij})_{n-1 \times n}$划掉第$j$列所得的$n-1$阶子式，证明：
        \begin{enumerate}
            \item $(M_1,-M_2,\ldots,(-1)^{n-1}M_n)$是方程组的一个解；

            \item 若$r(A)=n-1$，则方程组的解全是$(M_1,-M_2,\ldots,(-1)^{n-1}M_n)$的倍数.
        \end{enumerate}
\end{exercise}
\noindent 来源：\href{https://github.com/yhwu-is/Linear-Algebra-Left-Undone/blob/new/%E8%AE%B2%E4%B9%89/%E4%B8%93%E9%A2%98/10%20%E8%A1%8C%E5%88%97%E5%BC%8F.tex}{L10-B-7}。
\begin{solve}
\begin{enumerate}
                \item 由于行列式
                      \[ D = \begin{vmatrix}
                              a_{11}     & a_{12}     & a_{13}     & \cdots & a_{1n}     \\
                              a_{11}     & a_{12}     & a_{13}     & \cdots & a_{1n}     \\
                              a_{21}     & a_{22}     & a_{23}     & \cdots & a_{2n}     \\
                              \vdots     & \vdots     & \vdots     & \ddots & \vdots     \\
                              a_{n-1, 1} & a_{n-1, 2} & a_{n-1, 3} & \cdots & a_{n-1, n}
                          \end{vmatrix} = 0,\]
                      而 $M_1, -M_2, \ldots, (-1)^{n-1}M_n$ 恰是 $D$ 的第一行元素的代数余子式，所以将 $D$ 按第一行展开，可知\[a_{11}M_1+a_{12}(-M_2)+\cdots+a_{1n}(-1)^{n-1}M_n=0.\]
                      而 $D$ 的其他行元素与第一行元素的代数余子式乘积之和为 0，于是结论成立.

                \item 因为 $r(A) = n-1$，所以 $M_1, -M_2, \ldots, (-1)^{n-1}M_n$，不全为 0. 且该方程组解空间维数为 1，$M_1, -M_2, \ldots, (-1)^{n-1}M_n$ 正是该方程组的非零解，结论成立.
            \end{enumerate}
\end{solve}

\begin{exercise}[补充训练 L10-B-8]\label{prob:L10-B-8}
设$A,B$均为$n$阶矩阵，且$|A|=2,|B|=1$，求$|2\operatorname{adj}(A)\operatorname{adj}(B)-A^{-1}B^{-1}|$.
\end{exercise}
\noindent 来源：\href{https://github.com/yhwu-is/Linear-Algebra-Left-Undone/blob/new/%E8%AE%B2%E4%B9%89/%E4%B8%93%E9%A2%98/10%20%E8%A1%8C%E5%88%97%E5%BC%8F.tex}{L10-B-8}。
\begin{solve}
$\operatorname{adj}(A) = \lvert A \rvert A^{-1} = 2A^{-1}, \operatorname{adj}(B) = \lvert B \rvert B^{-1} = B^{-1}$，故 \[\lvert 2\operatorname{adj}(A)\operatorname{adj}(B)-A^{-1}B^{-1} \rvert = \lvert 4A^{-1}B^{-1}-A^{-1}B^{-1} \rvert = \lvert 3A^{-1}B^{-1} \rvert = \dfrac{3^n}{\lvert A \rvert \lvert B \rvert} = \dfrac{3^n}{2}.\]
\end{solve}

\begin{exercise}[补充训练 L10-B-9]\label{prob:L10-B-9}
设 $n\ge3$，$A$ 是非零实 $n$ 阶矩阵，且 $A^{\mathrm T}=\operatorname{adj}(A)$。证明：
\begin{enumerate}
\item $\det A>0$；
\item $\det A=1$；若第一行的全部元素相等，求其公共值；
\item $A$ 是正交矩阵，即 $AA^{\mathrm T}=A^{\mathrm T}A=I_n$；
\item 若 $n$ 为奇数，则 $\det(I_n-A)=0$。
\end{enumerate}
\end{exercise}
\noindent 来源：\href{https://github.com/yhwu-is/Linear-Algebra-Left-Undone/blob/new/%E8%AE%B2%E4%B9%89/%E4%B8%93%E9%A2%98/10%20%E8%A1%8C%E5%88%97%E5%BC%8F.tex}{L10-B-9}。
\begin{solve}
\begin{enumerate}
\item 由 $AA^{\mathrm T}=A\operatorname{adj}(A)=\det(A)I_n$，任取 $A$ 的一个非零行，其元素平方和等于 $\det A$，所以 $\det A>0$。实数域条件保证该平方和严格为正。
\item 对上述矩阵恒等式取行列式，得 $(\det A)^2=(\det A)^n$。因 $n\ge3$ 且 $\det A>0$，有 $\det A=1$。若第一行元素均为 $c$，由该行长度平方为 $1$ 得 $nc^2=1$，故 $c=\pm1/\sqrt n$。
\item 代回得 $AA^{\mathrm T}=I_n$，故 $A^{-1}=A^{\mathrm T}$，也有 $A^{\mathrm T}A=I_n$。
\item 因 $\det A=1$，有
\[\det(I_n-A)=\det\bigl(A(A^{\mathrm T}-I_n)\bigr)=\det(A-I_n)=(-1)^n\det(I_n-A).\]
奇数 $n$ 给出所求行列式为零。
\end{enumerate}
阶数条件不能删去：二阶矩阵 $A=2I_2$ 满足 $A^{\mathrm T}=\operatorname{adj}(A)$，但行列式为 $4$，也不是正交矩阵。
\end{solve}

\begin{exercise}[补充训练 L10-B-10]\label{prob:L10-B-10}
已知$A$是一个秩为$n-1$的$n\enspace(n \geqslant 2)$阶方阵，且已知某个元素$a_{ij}$的代数余子式$A_{ij} \neq 0$，求方程组$AX=0$的基础解系.
\end{exercise}
\noindent 来源：\href{https://github.com/yhwu-is/Linear-Algebra-Left-Undone/blob/new/%E8%AE%B2%E4%B9%89/%E4%B8%93%E9%A2%98/10%20%E8%A1%8C%E5%88%97%E5%BC%8F.tex}{L10-B-10}。
\begin{solve}
$r(A) = n-1$，则 $\lvert A \rvert = 0$ 且 $AX = 0$ 的解空间的维数为 1. 而考虑 $A\operatorname{adj}(A) = \lvert A \rvert E = 0$，且 $\exists A_{ij} \neq 0$，所以 $(A_{i1}, A_{i2}, \ldots , A_{in})^{\mathrm{T}}$ 是所求的基础解系.
\end{solve}

\begin{exercise}[补充训练 L10-B-11]\label{prob:L10-B-11}
设$D=|a_{ij}|_{n \times n}$，$A_{ij}$是$a_{ij}$的代数余子式. 求证：
        \[\begin{vmatrix}
                A_{11}    & A_{12}    & \cdots & A_{1,n-1}   \\
                A_{21}    & A_{22}    & \cdots & A_{2,n-1}   \\
                \vdots    & \vdots    & \ddots & \vdots      \\
                A_{n-1,1} & A_{n-1,2} & \cdots & A_{n-1,n-1}
            \end{vmatrix}=a_{nn}D^{n-2}.\]
\end{exercise}
\noindent 来源：\href{https://github.com/yhwu-is/Linear-Algebra-Left-Undone/blob/new/%E8%AE%B2%E4%B9%89/%E4%B8%93%E9%A2%98/10%20%E8%A1%8C%E5%88%97%E5%BC%8F.tex}{L10-B-11}。
\begin{solve}
设 \[A = \begin{pmatrix}
                a_{11} & a_{12} & \cdots & a_{1n} \\
                a_{21} & a_{22} & \cdots & a_{2n} \\
                \vdots & \vdots & \ddots & \vdots \\
                a_{n1} & a_{n2} & \cdots & a_{nn} \\
            \end{pmatrix},\]
            则 \[\operatorname{adj}(A) = \begin{pmatrix}
                    A_{11}     & A_{21}     & \cdots & A_{n-1, 1}   & A_{n1}     \\
                    A_{12}     & A_{22}     & \cdots & A_{n-1, 2}   & A_{n2}     \\
                    \vdots     & \vdots     & \ddots & \vdots       & \vdots     \\
                    A_{1, n-1} & A_{2, n-1} & \cdots & A_{n-1, n-1} & A_{n, n-1} \\
                    A_{1n}     & A_{2n}     & \cdots & A_{n-1, n}   & A_{nn}     \\
                \end{pmatrix}.\]
            注意到目标行列式是 $\operatorname{adj}(A)$ 中元素 $A_{nn}$ 的代数余子式，也就是 $\operatorname{adj}(\operatorname{adj}(A))$ 中 $(n, n)$ 位置的元素. 由 {例10.16(4)} $\operatorname{adj}(\operatorname{adj}(A)) = \lvert A \rvert^{n-2}A$ 可知结论成立.
\end{solve}

\begin{exercise}[补充训练 L10-B-12]\label{prob:L10-B-12}
设 $a_1,a_2,\ldots,a_n$ 是 $n$ 个两两互异的实数，$f_1(x),f_2(x),\ldots,f_n(x)$ 是 $n$ 个次数不大于 $n-2$ 的实系数多项式，证明：

        \[\begin{vmatrix}
                f_1(a_1) & f_1(a_2) & \cdots & f_1(a_n) \\
                f_2(a_1) & f_2(a_2) & \cdots & f_2(a_n) \\
                \vdots   & \vdots   & \ddots & \vdots   \\
                f_n(a_1) & f_n(a_2) & \cdots & f_n(a_n)
            \end{vmatrix}=0.\]
\end{exercise}
\noindent 来源：\href{https://github.com/yhwu-is/Linear-Algebra-Left-Undone/blob/new/%E8%AE%B2%E4%B9%89/%E4%B8%93%E9%A2%98/10%20%E8%A1%8C%E5%88%97%E5%BC%8F.tex}{L10-B-12}。
\begin{solve}
设 $f_i(x) = c_0 + c_1 x + c_2 x^2 + \cdots + c_{n-2} x^{n-2} \in R[x]_{n-1}$，由于 $n-1$ 维线性空间中的 $n$ 个向量必线性相关，我们可以得到 $\left\{ f_1(x), f_2(x), \ldots, f_n(x) \right\}$ 线性相关. 即存在不全为零的 $k_1, k_2, \ldots, k_n$，使得 $\forall x \in \mathbb{R}, k_1 f_1(x) + k_2 f_2(x) + \cdots + k_n f_n(x) = 0$. 设原行列式的行向量分别为 $\alpha_1, \alpha_2, \ldots, \alpha_n$，则有
            \[
                k_1 \alpha_1 + k_2 \alpha_2 + \cdots + k_n \alpha_n = \mathbf{0}.
            \]
            原行列式的行向量组线性相关，因此其值为 $0$.
\end{solve}

\begin{exercise}[补充训练 L10-B-13]\label{prob:L10-B-13}
证明：$n$维向量组$\alpha_1,\alpha_2,\ldots,\alpha_n$线性无关的充要条件是
        \[\begin{vmatrix}
                \alpha_1^\mathrm{T}\alpha_1 & \alpha_1^\mathrm{T}\alpha_2 & \cdots & \alpha_1^\mathrm{T}\alpha_n \\
                \alpha_2^\mathrm{T}\alpha_1 & \alpha_2^\mathrm{T}\alpha_2 & \cdots & \alpha_2^\mathrm{T}\alpha_n \\
                \vdots                      & \vdots                      & \ddots & \vdots                      \\
                \alpha_n^\mathrm{T}\alpha_1 & \alpha_n^\mathrm{T}\alpha_2 & \cdots & \alpha_n^\mathrm{T}\alpha_n
            \end{vmatrix}\neq 0.\]
\end{exercise}
\noindent 来源：\href{https://github.com/yhwu-is/Linear-Algebra-Left-Undone/blob/new/%E8%AE%B2%E4%B9%89/%E4%B8%93%E9%A2%98/10%20%E8%A1%8C%E5%88%97%E5%BC%8F.tex}{L10-B-13}。
\begin{solve}
令 
\[A = (\alpha_1, \alpha_2, \ldots, \alpha_n), B = \begin{pmatrix}
                \alpha_1^{\mathrm{T}}\alpha_1 & \alpha_1^{\mathrm{T}}\alpha_2 & \cdots & \alpha_1^{\mathrm{T}}\alpha_n \\
                \alpha_2^{\mathrm{T}}\alpha_1 & \alpha_2^{\mathrm{T}}\alpha_2 & \cdots & \alpha_2^{\mathrm{T}}\alpha_n \\
                \vdots                        & \vdots                        & \ddots & \vdots                        \\
                \alpha_n^{\mathrm{T}}\alpha_1 & \alpha_n^{\mathrm{T}}\alpha_2 & \cdots & \alpha_n^{\mathrm{T}}\alpha_n \\
            \end{pmatrix}\]
，显然 $B = A^{\mathrm{T}}A$，由 {11.3 秩不等式第 4 个} 有 $r(B) = r(A)$. 进而 $n$ 维向量组 $(\alpha_1, \alpha_2, \ldots, \alpha_n)$ 线性无关等价于 $r(A) = n$，也就等价于 $r(B) = n$，进而等价于 $\lvert B \rvert \neq 0$，命题得证.
\end{solve}

\begin{exercise}[补充训练 L10-B-14]\label{prob:L10-B-14}
设$a_1,\ldots,a_n$为$n$个$n$维向量，证明：向量组$a_1,\ldots,a_n$线性无关的充要条件是任一个$n$维向量都可以由其线性表示（不使用线性空间维数的方式完成）.
\end{exercise}
\noindent 来源：\href{https://github.com/yhwu-is/Linear-Algebra-Left-Undone/blob/new/%E8%AE%B2%E4%B9%89/%E4%B8%93%E9%A2%98/10%20%E8%A1%8C%E5%88%97%E5%BC%8F.tex}{L10-B-14}。
\begin{solve}
$(\implies)$ 记 $A = (\alpha_1, \alpha_2, \ldots, \alpha_n), \varepsilon_i = (0, \ldots, 1, 0, \ldots, 0)^{\mathrm{T}}$(第 $i$ 个为 1)，$E = (\varepsilon_1, \varepsilon_2, \ldots, \varepsilon_n)$. 因为 $\alpha_1, \alpha_2, \ldots, \alpha_n$ 是 $n$ 维线性无关的向量组，故 $\lvert A \rvert \neq 0$，即 $A$ 可逆，$A = EA, AA^{-1} = E$,\[(\alpha_1, \alpha_2, \ldots, \alpha_n) = (\varepsilon_1, \varepsilon_2, \ldots, \varepsilon_n)A, (\varepsilon_1, \varepsilon_2, \ldots, \varepsilon_n) = (\alpha_1, \alpha_2, \ldots, \alpha_n)A^{-1}.\] 即 $\alpha_1, \alpha_2, \ldots, \alpha_n$ 与 $\varepsilon_1, \varepsilon_2, \ldots, \varepsilon_n$ 等价. $\varepsilon_1, \varepsilon_2, \ldots, \varepsilon_n$ 为 $\mathbf{R}^n$ 的一个基，能表出任一 $n$ 维向量，故 $\alpha_1, \alpha_2, \ldots, \alpha_n$ 能表出 $\mathbf{R}^n$ 中任一 $n$ 维向量.

            $(\impliedby)$ 若 $\alpha_1, \alpha_2, \ldots, \alpha_n$ 能表出任一 $n$ 维向量，则 $\varepsilon_i = \lambda_{i1}\alpha_1+\cdots+\lambda_{in}\alpha_n, i = 1, 2, \ldots, n$. 即
            \[E = (\varepsilon_1, \varepsilon_2, \ldots, \varepsilon_n) = (\alpha_1, \alpha_2, \ldots, \alpha_n)\begin{pmatrix}
                    \lambda_{11} & \lambda_{21} & \cdots & \lambda_{n1} \\
                    \lambda_{11} & \lambda_{21} & \cdots & \lambda_{n1} \\
                    \vdots       & \vdots       & \ddots & \vdots       \\
                    \lambda_{11} & \lambda_{21} & \cdots & \lambda_{n1} \\
                \end{pmatrix}.\]
            进而 $\lvert \alpha_1, \alpha_2, \ldots, \alpha_n \rvert \begin{vmatrix}
                    \lambda_{11} & \lambda_{21} & \cdots & \lambda_{n1} \\
                    \lambda_{11} & \lambda_{21} & \cdots & \lambda_{n1} \\
                    \vdots       & \vdots       & \ddots & \vdots       \\
                    \lambda_{11} & \lambda_{21} & \cdots & \lambda_{n1} \\
                \end{vmatrix} = 1$，所以 $\lvert \alpha_1, \alpha_2, \ldots, \alpha_n \rvert \neq 0$，$\alpha_1, \alpha_2, \ldots, \alpha_n$ 线性无关.
\end{solve}

\begin{exercise}[补充训练 L10-B-15]\label{prob:L10-B-15}
设$s \times n\enspace(s\leqslant n)$矩阵为
        \[\begin{pmatrix}
                1      & a      & a^2    & \cdots & a^{n-1}    \\
                1      & a^2    & a^4    & \cdots & a^{2(n-1)} \\
                \vdots & \vdots & \vdots & \ddots & \vdots     \\
                1      & a^s    & a^{2s} & \cdots & a^{s(n-1)}
            \end{pmatrix}\]
        且$a^r\neq 1\enspace(0<r<n)$，求$A$的秩和它的列向量组的一个极大线性无关组.
\end{exercise}
\noindent 来源：\href{https://github.com/yhwu-is/Linear-Algebra-Left-Undone/blob/new/%E8%AE%B2%E4%B9%89/%E4%B8%93%E9%A2%98/10%20%E8%A1%8C%E5%88%97%E5%BC%8F.tex}{L10-B-15}。
\begin{solve}
$A$ 的前 $s$ 列组成的 $s$ 阶子式为Vandermonde行列式
            \[D = \begin{vmatrix}
                    1      & a      & a^2    & \cdots & a^{s-1}    \\
                    1      & a^2    & a^4    & \cdots & a^{2(s-1)} \\
                    \vdots & \vdots & \vdots & \ddots & \vdots     \\
                    1      & a^s    & a^{2s} & \cdots & a^{s(s-1)} \\
                \end{vmatrix}.\]
            由于当 $0 < r < n$ 时，$a^r \neq 1$，因此 $a, a^2, \ldots, a^s$ 两两不同，进而 $D \neq 0$，于是 $r(A) \geqslant s$. 又因为 $A$ 的行数是 $s$，所以 $r(A) \leqslant s$. 从而 $r(A) = s$.

            由于 $D \neq 0$，故 $D$ 的列向量线性无关，因此 $A$ 的前 $s$ 个列向量即为其列向量组的一个极大线性无关组.
\end{solve}

\begin{exercise}[补充训练 L10-B-16]\label{prob:L10-B-16}
设$A,B,C,D \in \mathbf{F}^{n \times n}$，定义变换 $ T : \mathbf{F}^{n \times n} \to \mathbf{F}^{n \times n}$ 为
        \[ T(X) = AXB+CX+XD \]
        证明：
        \begin{enumerate}
            \item $T$为$\mathbf{F}^{n \times n}$上的线性变换；

            \item 当$C=D=O$时，$T$可逆的充要条件是$|AB| \neq 0$.
        \end{enumerate}
\end{exercise}
\noindent 来源：\href{https://github.com/yhwu-is/Linear-Algebra-Left-Undone/blob/new/%E8%AE%B2%E4%B9%89/%E4%B8%93%E9%A2%98/10%20%E8%A1%8C%E5%88%97%E5%BC%8F.tex}{L10-B-16}。
\begin{solve}
\begin{enumerate}
\item 对 $X,Y\in\mathbf F^{n\times n}$ 及 $a,b\in\mathbf F$，逐项使用矩阵乘法分配律可得 $T(aX+bY)=aT(X)+bT(Y)$，故 $T$ 为线性变换。
\item 此时 $T(X)=AXB$。若 $A,B$ 均可逆，则逆映射为 $X\mapsto A^{-1}XB^{-1}$。反之，若 $A$ 奇异，取非零列向量 $u\in\ker A$ 和非零列向量 $v$，则 $X=uv^{\mathrm T}\ne O$ 且 $AXB=O$，故 $T$ 不单射。若 $B$ 奇异，取非零 $v\in\ker B^{\mathrm T}$ 与非零 $u$，同样有 $uv^{\mathrm T}B=O$。因此 $T$ 可逆当且仅当 $A,B$ 都可逆，也即 $|AB|\ne0$。\end{enumerate}
\end{solve}

\begin{exercise}[补充训练 L10-B-17]\label{prob:L10-B-17}
设$A$为$n$阶矩阵，且$r(A) < n$，又$A_{11} \neq 0$，证明：存在常数$k$，使得$(\operatorname{adj}(A))^2=k\operatorname{adj}(A)$.
\end{exercise}
\noindent 来源：\href{https://github.com/yhwu-is/Linear-Algebra-Left-Undone/blob/new/%E8%AE%B2%E4%B9%89/%E4%B8%93%E9%A2%98/10%20%E8%A1%8C%E5%88%97%E5%BC%8F.tex}{L10-B-17}。
\begin{solve}
因为 $r(A) < n, A_{11} \neq 0$，所以 $r(A) = n-1$，进而由 {例10.16 (6)} 可知 $r(\operatorname{adj}(A)) = 1$，所以有 
            \[\operatorname{adj}(A) = \begin{pmatrix} a_1 \\ a_2 \\ \vdots \\ a_n \end{pmatrix} (\lambda_1, \lambda_2, \ldots, \lambda_n).\]
            设 $(\lambda_1, \lambda_2, \ldots, \lambda_n) \begin{pmatrix} a_1 \\ a_2 \\ \vdots \\ a_n \end{pmatrix} = k$，则
            \begin{align*}
                (\operatorname{adj}(A))^2 & = \begin{pmatrix}a_1 \\ a_2 \\ \vdots \\ a_n \end{pmatrix}
                (\lambda_1, \lambda_2, \ldots, \lambda_n)
                \begin{pmatrix} a_1 \\ a_2 \\ \vdots \\ a_n \end{pmatrix}
                (\lambda_1, \lambda_2, \ldots, \lambda_n)                                     \\
                        & = \begin{pmatrix} a_1 \\ a_2 \\ \vdots \\ a_n \end{pmatrix} \cdot k
                \cdot (\lambda_1, \lambda_2, \ldots, \lambda_n) = k
                \begin{pmatrix} a_1 \\ a_2 \\ \vdots \\ a_n \end{pmatrix}
                (\lambda_1, \lambda_2, \ldots, \lambda_n) = k\operatorname{adj}(A).
            \end{align*}
\end{solve}

\begin{exercise}[补充训练 L10-B-18]\label{prob:L10-B-18}
设$V$是一个$n$维实线性空间，证明：存在$V$中的一个由可列无穷多个向量组成的向量组$\{\alpha_i \mid i\in\mathbf{Z}_+\}$，使得其中任意$n$个向量组成的向量组都是$V$的一组基.
\end{exercise}
\noindent 来源：\href{https://github.com/yhwu-is/Linear-Algebra-Left-Undone/blob/new/%E8%AE%B2%E4%B9%89/%E4%B8%93%E9%A2%98/10%20%E8%A1%8C%E5%88%97%E5%BC%8F.tex}{L10-B-18}。
\begin{solve}
$\forall a_i \in \mathbf{R}, i \in \mathbf{Z}_{+}$ 满足：若 $i \neq j$，则 $a_i \neq a_j$，考虑以Vandermonde行列式的形式进行排布. 设
            \[ A = \begin{pmatrix}
                    1         & 1         & \cdots & 1         & \cdots \\
                    a_1       & a_2       & \cdots & a_k       & \cdots \\
                    a_1^2     & a_2^2     & \cdots & a_k^2     & \cdots \\
                    \vdots    & \vdots    & \ddots & \vdots    &        \\
                    a_1^{n-1} & a_2^{n-1} & \cdots & a_k^{n-1} & \cdots
                \end{pmatrix}.\]
            考虑 $A$ 的任意 $n$ 阶主子式 $D_n$，其均构成 Vandermonde 行列式，又 $i \neq j$ 有 $a_i \neq a_j$，所以值均不为 0，也就是说任意 $n$ 个向量都线性无关，其个数也恰好为 $n$，构成 $V$ 的一组基，命题得证.
\end{solve}

\begin{exercise}[补充训练 L10-B-20]\label{prob:L10-B-20}
解行列式
        \begin{enumerate} 
                \item $ D_1=\begin{vmatrix}
                        1 & 2 & 3 & 4 \\2&3&4&1\\3&4&1&2\\4&1&2&3
                    \end{vmatrix} $

                \item $ D_2=\begin{vmatrix}
                        \lambda+2 & -1        & -1        & -1        \\
                        -1        & \lambda+2 & -1        & -1        \\
                        -1        & -1        & \lambda+2 & -1        \\
                        -1        & -1        & -1        & \lambda+2
                    \end{vmatrix} $

              

                \item $ D_3=\begin{vmatrix}
                        1^{2} & 2^{2} & 3^{2} & 4^{2} \\
                        2^{2} & 3^{2} & 4^{2} & 5^{2} \\
                        3^{2} & 4^{2} & 5^{2} & 6^{2} \\
                        4^{2} & 5^{2} & 6^{2} & 7^{2}
                    \end{vmatrix} $

                \item $ D_4=\begin{vmatrix}
                        3 & 2 & 0 & 0 \\
                        1 & 3 & 2 & 0 \\
                        0 & 1 & 3 & 2 \\
                        0 & 0 & 1 & 3
                    \end{vmatrix} $
             \end{enumerate}
\end{exercise}
\noindent 来源：\href{https://github.com/yhwu-is/Linear-Algebra-Left-Undone/blob/new/%E8%AE%B2%E4%B9%89/%E4%B8%93%E9%A2%98/10%20%E8%A1%8C%E5%88%97%E5%BC%8F.tex}{L10-B-20}。
\begin{solve}
\begin{enumerate}
                \item 连加法得结果为160；

                \item 连加法得结果为$(\lambda-1)(\lambda+3)^3$；

                \item 如果对数据敏感能看出线性相关，可以直接由$\alpha_4-\alpha_1=3(\alpha_3-\alpha_2)$证得其线性相关，则为0. 也可以使用滚动消去法，如\[
                          \begin{vmatrix}
                              1 & 4 & 9  & 16 \\
                              3 & 5 & 7  & 9  \\
                              5 & 7 & 9  & 11 \\
                              7 & 9 & 11 & 13
                          \end{vmatrix}=\begin{vmatrix}
                              1 & 4 & 9 & 16 \\
                              3 & 5 & 7 & 9  \\
                              2 & 2 & 2 & 2  \\
                              2 & 2 & 2 & 2
                          \end{vmatrix}=0.\]

                \item 当然是递推法的常见形式，但是只有4阶，所以可以暴算，转化为上三角行列式，如\[
                          \begin{vmatrix}
                              3 & 2           & 0            & 0             \\
                              0 & \frac{7}{3} & 2            & 0             \\
                              0 & 0           & \frac{15}{7} & 2             \\
                              0 & 0           & 0            & \frac{31}{15}
                          \end{vmatrix}=31.\]
            \end{enumerate}
\end{solve}

\begin{exercise}[补充训练 L10-B-21]\label{prob:L10-B-21}
解行列式 $\begin{vmatrix}
                a_{1}+a_{2}         & a_{2}+a_{3}         & \cdots & a_{n-1}+a_n         & a_n+a_{1}         \\
                a_{1}^{2}+a_{2}^{2} & a_{2}^{2}+a_{3}^{2} & \cdots & a_{n-1}^{2}+a_n^{2} & a_n^{2}+a_{1}^{2} \\
                \vdots              & \vdots              & \ddots & \vdots              & \vdots            \\
                a_{1}^{n}+a_{2}^{n} & a_{2}^{n}+a_{3}^{n} & \cdots & a_{n-1}^{n}+a_n^{n} & a_n^{n}+a_{1}^{n}
            \end{vmatrix}$.
\end{exercise}
\noindent 来源：\href{https://github.com/yhwu-is/Linear-Algebra-Left-Undone/blob/new/%E8%AE%B2%E4%B9%89/%E4%B8%93%E9%A2%98/10%20%E8%A1%8C%E5%88%97%E5%BC%8F.tex}{L10-B-21}。
\begin{solve}
暴力拆成$2^n$个行列式，其中只有2个每列各不相同，其他都因为有相同列而为0. 因此有
            \begin{align*}
                \text{原式} & =\begin{vmatrix}
                                    a_{1}     & a_{2}     & \cdots & a_{n-1}     & a_{n}     \\
                                    a_{1}^{2} & a_{2}^{2} & \cdots & a_{n-1}^{2} & a_{n}^{2} \\
                                    \vdots    & \vdots    & \ddots & \vdots      & \vdots    \\
                                    a_{1}^{n} & a_{2}^{n} & \cdots & a_{n-1}^{n} & a_{n}^{n}
                                \end{vmatrix}+
                \begin{vmatrix}
                    a_{2}     & a_{3}     & \cdots & a_{n}     & a_{1}     \\
                    a_{2}^{2} & a_{3}^{2} & \cdots & a_{n}^{2} & a_{1}^{2} \\
                    \vdots    & \vdots    & \ddots & \vdots    & \vdots    \\
                    a_{2}^{n} & a_{3}^{n} & \cdots & a_{n}^{n} & a_{1}^{n}
                \end{vmatrix}                                               \\
                            & =(1+(-1)^{n-1})\begin{vmatrix}
                                                a_{1}     & a_{2}     & \cdots & a_{n-1}     & a_{n}     \\
                                                a_{1}^{2} & a_{2}^{2} & \cdots & a_{n-1}^{2} & a_{n}^{2} \\
                                                \vdots    & \vdots    & \ddots & \vdots      & \vdots    \\
                                                a_{1}^{n} & a_{2}^{n} & \cdots & a_{n-1}^{n} & a_{n}^{n}
                                            \end{vmatrix}                \\
                            & =(1+(-1)^{n-1})\prod_{i=1}^na_i
                \begin{vmatrix}
                    1           & 1           & \cdots & 1             & 1           \\
                    a_{1}       & a_{2}       & \cdots & a_{n-1}       & a_{n}       \\
                    \vdots      & \vdots      & \ddots & \vdots        & \vdots      \\
                    a_{1}^{n-1} & a_{2}^{n-1} & \cdots & a_{n-1}^{n-1} & a_{n}^{n-1}
                \end{vmatrix}                                     \\
                            & =(1+(-1)^{n-1})\left(\prod_{i=1}^na_i\right)\prod_{1\leqslant i<j\leqslant n}(a_j-a_i). \\
            \end{align*}
\end{solve}

\begin{exercise}[补充训练 L10-B-22]\label{prob:L10-B-22}
解行列式
         \begin{enumerate}
                \item $D=\begin{vmatrix}
                              ax+by & ay+bz & az+bx \\
                              ay+bz & az+bx & ax+by \\
                              az+bx & ax+by & ay+bz
                          \end{vmatrix}$

                \item $D=\begin{vmatrix}
                              x^2+1 & xy    & xz    \\
                              xy    & y^2+1 & yz    \\
                              xz    & yz    & z^2+1
                          \end{vmatrix}$
            \end{enumerate}
\end{exercise}
\noindent 来源：\href{https://github.com/yhwu-is/Linear-Algebra-Left-Undone/blob/new/%E8%AE%B2%E4%B9%89/%E4%B8%93%E9%A2%98/10%20%E8%A1%8C%E5%88%97%E5%BC%8F.tex}{L10-B-22}。
\begin{solve}
\begin{enumerate}
\item 写成矩阵乘积：
\[M=\begin{pmatrix}a&b&0\\0&a&b\\b&0&a\end{pmatrix}\begin{pmatrix}x&y&z\\y&z&x\\z&x&y\end{pmatrix}.\]
第一个因子的行列式为 $a^3+b^3$，第二个为 $3xyz-x^3-y^3-z^3$，所以
\[D=(a^3+b^3)(3xyz-x^3-y^3-z^3).\]
这里不能把三阶行列式中的标量因子误写成一次项，也不能把立方和写成平方和。
\item 令 $u=(x,y,z)^{\mathrm T}$，则矩阵为 $I_3+uu^{\mathrm T}$。由恒等式 $\det(I+uv^{\mathrm T})=1+v^{\mathrm T}u$，得到 $D=1+x^2+y^2+z^2$。该恒等式也可由行列式按列的多线性直接证明：取两列以上的秩一部分时这些列成比例，相应行列式为零，剩下常数项和三个平方项。\end{enumerate}
\end{solve}

\begin{exercise}[补充训练 L10-B-23]\label{prob:L10-B-23}
计算行列式$|2E-\alpha_1^\mathrm{T}\beta_1-\alpha_2^\mathrm{T}\beta_2|$，其中$\alpha_1=(a_1,a_2,\ldots,a_n),\enspace \beta_1=(b_1,b_2,\ldots,b_n),\enspace \alpha_2=(c_1,c_2,\ldots,c_n),\enspace \beta_2 = (d_1,d_2,\ldots,d_n)$.（提示：利用$|\lambda E_m-AB|=\lambda^{m-n}|\lambda E_n-BA|$）
\end{exercise}
\noindent 来源：\href{https://github.com/yhwu-is/Linear-Algebra-Left-Undone/blob/new/%E8%AE%B2%E4%B9%89/%E4%B8%93%E9%A2%98/10%20%E8%A1%8C%E5%88%97%E5%BC%8F.tex}{L10-B-23}。
\begin{solve}
利用$|\lambda E_m-AB|=\lambda^{m-n}|\lambda E_n-BA|$
            \begin{align*}
                |2E-\alpha_1^\mathrm{T}\beta_1-\alpha_2^\mathrm{T}\beta_2|
                & =\begin{vmatrix}
                        2E-\begin{pmatrix}
                        \alpha_1^\mathrm{T} & \alpha_2^\mathrm{T}
                    \end{pmatrix}\begin{pmatrix}
                                        \beta_1 \\\beta_2
                                    \end{pmatrix}
                    \end{vmatrix}
                =2^{n-2}\begin{vmatrix}
                            2E-\begin{pmatrix}
                        \beta_1 \\\beta_2
                    \end{pmatrix}\begin{pmatrix}
                                        \alpha_1^\mathrm{T} & \alpha_2^\mathrm{T}
                                    \end{pmatrix}
                        \end{vmatrix}         \\
                & =2^{n-2}\begin{vmatrix}
                            2-\beta_1\alpha_1^\mathrm{T} & -\beta_1\alpha_2^\mathrm{T}  \\
                            -\beta_2\alpha_1^\mathrm{T}  & 2-\beta_2\alpha_2^\mathrm{T}
                        \end{vmatrix} \\
                & =2^{n-2}
                \left(\left(2-\sum_{i=1}^na_ib_i\right)
                \left(2-\sum_{i=1}^nc_id_i\right)
                -\left(\sum_{i=1}^na_id_i\right)
                \left(\sum_{i=1}^nb_ic_i\right)\right)
            \end{align*}
\end{solve}

\begin{exercise}[补充训练 L10-B-24]\label{prob:L10-B-24}
求$2n$阶行列式的值（空缺处都是0）：
        \[\begin{vmatrix}
                a &         &   &   &         & b \\
                  & \ddots  &   &   & \iddots &   \\
                  &         & a & b &         &   \\
                  &         & b & a &         &   \\
                  & \iddots &   &   & \ddots  &   \\
                b &         &   &   &         & a
            \end{vmatrix}.\]
\end{exercise}
\noindent 来源：\href{https://github.com/yhwu-is/Linear-Algebra-Left-Undone/blob/new/%E8%AE%B2%E4%B9%89/%E4%B8%93%E9%A2%98/10%20%E8%A1%8C%E5%88%97%E5%BC%8F.tex}{L10-B-24}。
\begin{solve}
设所求 $2n$ 阶行列式为 $D$.
            \begin{enumerate}
                \item 若 $a=0$，则
                    \[
                        D = \begin{vmatrix}
                            0 &         &   &   &         & b \\
                              & \ddots  &   &   &         &   \\
                              &         & 0 & b &         &   \\
                              &         & b & 0 &         &   \\
                              &         &   &   & \ddots  &   \\
                            b &         &   &   &         & 0
                        \end{vmatrix}
                        = (-1)^{\frac{2n(2n-1)}{2}} \begin{vmatrix}
                            b &         &   &   &         & 0 \\
                              & \ddots  &   &   &         &   \\
                              &         & b & 0 &         &   \\
                              &         & 0 & b &         &   \\
                              &         &   &   & \ddots  &   \\
                            0 &         &   &   &         & b
                        \end{vmatrix}
                        = b^{2n}.
                    \]
                \item 若 $a \neq 0$，我们将第 $i$ 列乘以 $-\dfrac{b}{a}$ 后加到第 $2n-i+1$ 列，化为下三角矩阵：
                    \[
                        D = \begin{vmatrix}
                            a &         &   &                   &         & 0                 \\
                              & \ddots  &   &                   &         &                   \\
                              &         & a & 0                 &         &                   \\
                              &         & b & a - \frac{b^2}{a} &         &                   \\
                              &         &   &                   & \ddots  &                   \\
                            b &         &   &                   &         & a - \frac{b^2}{a}
                        \end{vmatrix}
                        = a^n (a - \frac{b^2}{a})^n
                        = (a^2 - b^2)^n.
                    \]
            \end{enumerate}
\end{solve}

\begin{exercise}[补充训练 L10-C-3]\label{prob:L10-C-3}
设 $a_i, b_i$ 都是实数，利用 Cauchy-Binet 公式证明 Cauchy-Schwarz 不等式：
        \[
            \left( \sum_{i=1}^{n}\limits a_i^2 \right) \left( \sum_{i=1}^{n}\limits b_i^2 \right) \geqslant \left( \sum_{i=1}^{n}\limits a_i b_i \right)^2.
        \]
\end{exercise}
\noindent 来源：\href{https://github.com/yhwu-is/Linear-Algebra-Left-Undone/blob/new/%E8%AE%B2%E4%B9%89/%E4%B8%93%E9%A2%98/10%20%E8%A1%8C%E5%88%97%E5%BC%8F.tex}{L10-C-3}。
\begin{solve}
设

            \[
                A = \begin{pmatrix}
                    a_1 & a_2 & \cdots & a_n \\
                    b_1 & b_2 & \cdots & b_n
                \end{pmatrix}.
            \]

            由 Cauchy-Binet 公式的推论，$A A^{\mathrm{T}}$ 的任一主子式非负，故 $|AA^\mathrm{T}|$ 非负，即

            \[
                \begin{vmatrix}
                    \sum\limits_{i=1}^{n} a_i^2 & \sum\limits_{i=1}^{n} a_ib_i \\
                    \sum\limits_{i=1}^{n} a_ib_i & \sum\limits_{i=1}^{n} b_i^2
                \end{vmatrix} = |AA^\mathrm{T}| \geqslant 0.
            \]

            将上述行列式展开即有

            \[
                \left( \sum\limits_{i=1}^{n} a_i^2 \right) \left( \sum\limits_{i=1}^{n} b_i^2 \right) - \left( \sum\limits_{i=1}^{n} a_i b_i \right)^2 \geqslant 0.
            \]

            故 Cauchy-Schwarz 不等式成立.
\end{solve}

\begin{exercise}[补充训练 L10-C-4]\label{prob:L10-C-4}
设$A=(a_{ij})$是$n\enspace(n\geqslant 2)$阶整数方阵，满足对任意的$i,j$，$|A|$均可整除$a_{ij}$，证明：$|A|=\pm 1$.
\end{exercise}
\noindent 来源：\href{https://github.com/yhwu-is/Linear-Algebra-Left-Undone/blob/new/%E8%AE%B2%E4%B9%89/%E4%B8%93%E9%A2%98/10%20%E8%A1%8C%E5%88%97%E5%BC%8F.tex}{L10-C-4}。
\begin{solve}
设 $a_{ij} = |A| b_{ij}$，其中 $b_{ij}$ 为整数，则
            \[
                |A| = \begin{vmatrix}
                    |A| b_{11} & |A| b_{12} & \cdots & |A| b_{1n} \\
                    |A| b_{21} & |A| b_{22} & \cdots & |A| b_{2n} \\
                    \vdots     & \vdots     & \ddots & \vdots     \\
                    |A| b_{n1} & |A| b_{n2} & \cdots & |A| b_{nn}
                \end{vmatrix}
                = |A|^n \begin{vmatrix}
                    b_{11} & b_{12} & \cdots & b_{1n} \\
                    b_{21} & b_{22} & \cdots & b_{2n} \\
                    \vdots & \vdots & \ddots & \vdots \\
                    b_{n1} & b_{n2} & \cdots & b_{nn}
                \end{vmatrix}.
            \]

            由此可得
            \[
                |A|^{n-1} \begin{vmatrix}
                    b_{11} & b_{12} & \cdots & b_{1n} \\
                    b_{21} & b_{22} & \cdots & b_{2n} \\
                    \vdots & \vdots & \ddots & \vdots \\
                    b_{n1} & b_{n2} & \cdots & b_{nn}
                \end{vmatrix} = 1.
            \]

            由于 $a_{ij}$ 与 $b_{ij}$ 均为整数，所以 $|A|$ 与 $|b_{ij}|_{n \times n}$ 也为整数，且 $|A| \neq 0$，所以 $|A| = \pm 1$.
\end{solve}

\begin{exercise}[补充训练 L10-C-5]\label{prob:L10-C-5}
设$A,B,C,D$均为$n$阶方阵，$\lvert A \rvert \neq 0$且$AC=CA$. 证明：
        \[\begin{vmatrix}
                A & B \\ C & D
            \end{vmatrix} = |AD-CB|.\]
\end{exercise}
\noindent 来源：\href{https://github.com/yhwu-is/Linear-Algebra-Left-Undone/blob/new/%E8%AE%B2%E4%B9%89/%E4%B8%93%E9%A2%98/10%20%E8%A1%8C%E5%88%97%E5%BC%8F.tex}{L10-C-5}。
\begin{solve}
对原矩阵进行分块初等变换化为上三角块矩阵后进行计算. 因为 $\lvert A \rvert \neq 0$，所以 $A$ 可逆，进而有如下变换：
            \[\begin{pmatrix}
                    E        & O \\
                    -CA^{-1} & E
                \end{pmatrix} \begin{pmatrix}
                    A & B \\
                    C & D \\
                \end{pmatrix} = \begin{pmatrix}
                    A & B          \\
                    O & D-CA^{-1}B
                \end{pmatrix},\]
            所以
            \begin{align*}
                    & \begin{vmatrix}
                            A & B \\
                            C & D \\
                        \end{vmatrix}
                = \begin{vmatrix}
                        E        & O \\
                        -CA^{-1} & E \\
                    \end{vmatrix}
                \begin{vmatrix}
                    A & B \\
                    C & D \\
                \end{vmatrix}       \\
                ={} & \begin{vmatrix}
                            A & B          \\
                            O & D-CA^{-1}B
                        \end{vmatrix}
                = \lvert A \rvert \lvert D-CA^{-1}B \rvert = \lvert AD-ACA^{-1}B \rvert.
            \end{align*} 由于 $AC = CA$，所以有 $ACA^{-1} = CAA^{-1} = C$，所以
            \[\begin{vmatrix}
                    A & B \\
                    C & D \\
                \end{vmatrix} = \lvert AD-CB \rvert.\]
\end{solve}

\begin{exercise}[补充训练 L10-C-6]\label{prob:L10-C-6}
设$A$为$n$阶可逆矩阵，$\alpha,\beta$为$n$维列向量，证明：
        \[|A+\alpha\beta^{\mathrm{T}}|=|A|(1+\beta^\mathrm{T}A^{-1}\alpha).\]
\end{exercise}
\noindent 来源：\href{https://github.com/yhwu-is/Linear-Algebra-Left-Undone/blob/new/%E8%AE%B2%E4%B9%89/%E4%B8%93%E9%A2%98/10%20%E8%A1%8C%E5%88%97%E5%BC%8F.tex}{L10-C-6}。
\begin{solve}
这道题目我们利用了一个分块矩阵作为中间``桥梁''使得其通过分块初等变换之后能分别得到两个方向上的结果. 考虑矩阵 $\begin{pmatrix}
                A                   & \alpha \\
                -\beta^{\mathrm{T}} & 1
            \end{pmatrix}$，有
            \[\begin{pmatrix}
                    A                   & \alpha \\
                    -\beta^{\mathrm{T}} & 1
                \end{pmatrix} \begin{pmatrix}
                    E                  & O \\
                    \beta^{\mathrm{T}} & E
                \end{pmatrix} = \begin{pmatrix}
                    A+\alpha \beta^{\mathrm{T}} & \alpha \\
                    O                           & 1
                \end{pmatrix},\]
            所以
            \[\begin{vmatrix}
                    A                   & \alpha \\
                    -\beta^{\mathrm{T}} & 1
                \end{vmatrix} = \begin{vmatrix}
                    A                   & \alpha \\
                    -\beta^{\mathrm{T}} & 1
                \end{vmatrix} \begin{vmatrix}
                    E                  & O \\
                    \beta^{\mathrm{T}} & E
                \end{vmatrix} = \begin{vmatrix}
                    A+\alpha \beta^{\mathrm{T}} & \alpha \\
                    O                           & 1
                \end{vmatrix} = \lvert A+\alpha \beta^{\mathrm{T}} \rvert.\]
            另一方面，
            \[\begin{pmatrix}
                    E                        & O \\
                    \beta^{\mathrm{T}}A^{-1} & 1
                \end{pmatrix} \begin{pmatrix}
                    A                   & \alpha \\
                    -\beta^{\mathrm{T}} & 1
                \end{pmatrix} = \begin{pmatrix}
                    A & \alpha                           \\
                    O & 1+\beta^{\mathrm{T}}A^{-1}\alpha
                \end{pmatrix}.\]
            注意 $\beta^{\mathrm{T}}A^{-1}\alpha$ 的最终结果是一个数. 进而
            \[\begin{vmatrix}
                    A                   & \alpha \\
                    -\beta^{\mathrm{T}} & 1
                \end{vmatrix} = \begin{vmatrix}
                    E                        & O \\
                    \beta^{\mathrm{T}}A^{-1} & 1
                \end{vmatrix} \begin{vmatrix}
                    A                   & \alpha \\
                    -\beta^{\mathrm{T}} & 1
                \end{vmatrix} = \begin{vmatrix}
                    A & \alpha                           \\
                    O & 1+\beta^{\mathrm{T}}A^{-1}\alpha
                \end{vmatrix} = \lvert A \rvert(1+\beta^{\mathrm{T}}A^{-1}\alpha).\]
            所以有
            \[\lvert A+\alpha \beta^{\mathrm{T}} \rvert = \lvert A \rvert(1+\beta^{\mathrm{T}}A^{-1}\alpha).\]
\end{solve}

\begin{exercise}[补充训练 L10-C-7]\label{prob:L10-C-7}
设$A,B$均为$n$阶方阵，证明：
        \[\begin{vmatrix}
                A & B \\ B & A
            \end{vmatrix} = |A+B||A-B|.\]
\end{exercise}
\noindent 来源：\href{https://github.com/yhwu-is/Linear-Algebra-Left-Undone/blob/new/%E8%AE%B2%E4%B9%89/%E4%B8%93%E9%A2%98/10%20%E8%A1%8C%E5%88%97%E5%BC%8F.tex}{L10-C-7}。
\begin{solve}
依旧是对分块矩阵做初等分块变换. 考虑到
            \begin{align*}
                \left(\begin{pmatrix}
                            E & E \\
                            O & E
                        \end{pmatrix}
                \begin{pmatrix}
                    A & B \\
                    B & A
                \end{pmatrix} \right)
                \begin{pmatrix}
                    E & -E \\
                    O & E
                \end{pmatrix}
                & = \begin{pmatrix}
                        A+B & A+B \\
                        B   & A
                    \end{pmatrix}
                \begin{pmatrix}
                    E & -E \\
                    O & E
                \end{pmatrix}      \\
                & = \begin{pmatrix}
                        A+B & O   \\
                        B   & A-B
                    \end{pmatrix},
            \end{align*}
            所以有 \begin{align*}
                \begin{vmatrix}
                    A & B \\
                    B & A
                \end{vmatrix}
                & = \begin{vmatrix}
                        E & E \\
                        O & E
                    \end{vmatrix}
                \begin{vmatrix}
                    A & B \\
                    B & A
                \end{vmatrix}
                \begin{vmatrix}
                    E & -E \\
                    O & E
                \end{vmatrix}
                = \begin{vmatrix}
                        A+B & O   \\
                        B   & A-B
                    \end{vmatrix}    \\
                & =
                \lvert (A+B)(A-B) \rvert = \lvert A+B \rvert \lvert A-B \rvert
            \end{align*}
\end{solve}

\begin{exercise}[补充训练 L10-C-8]\label{prob:L10-C-8}
设$A,B,C,D$均为$n$阶方阵，且$r\begin{pmatrix}
                A & B \\ C & D
            \end{pmatrix}=n$，证明：
        \[\begin{vmatrix}
                |A| & |B| \\ |C| & |D|
            \end{vmatrix} = 0.\]
\end{exercise}
\noindent 来源：\href{https://github.com/yhwu-is/Linear-Algebra-Left-Undone/blob/new/%E8%AE%B2%E4%B9%89/%E4%B8%93%E9%A2%98/10%20%E8%A1%8C%E5%88%97%E5%BC%8F.tex}{L10-C-8}。
\begin{solve}
首先有 $\begin{vmatrix}
                \lvert A \rvert & \lvert B \rvert \\
                \lvert C \rvert & \lvert D \rvert
            \end{vmatrix} = \lvert A \rvert \lvert D \rvert-\lvert B \rvert \lvert C \rvert$.
            \begin{enumerate}
                \item 若 $\lvert A \rvert \neq 0$，即 $A$ 可逆. 因为矩阵的初等变换不改变矩阵的秩，所以由
                    \[\begin{pmatrix}
                            E        & O \\
                            -CA^{-1} & E
                        \end{pmatrix} \begin{pmatrix}
                            A & B \\
                            C & D
                        \end{pmatrix} \begin{pmatrix}
                            E & -A^{-1}B \\
                            O & E        \\
                        \end{pmatrix} = \begin{pmatrix}
                            A & O          \\
                            O & D-CA^{-1}B
                        \end{pmatrix},\]
                    条件 $r\left(\begin{pmatrix}
                                A & B \\
                                C & D
                            \end{pmatrix}\right) = n$ 以及 $A$ 可逆，可以得到
                    \[D-CA^{-1}B = O.\]
                    即若 $A$ 可逆，则 $D = CA^{-1}B$，并且
                    \[\begin{vmatrix}
                            \lvert A \rvert & \lvert B \rvert \\
                            \lvert C \rvert & \lvert D \rvert
                        \end{vmatrix} = \lvert A \rvert \lvert CA^{-1}B \rvert-\lvert B \rvert \lvert C \rvert = \lvert A \rvert \lvert C \rvert \lvert A^{-1} \rvert \lvert B \rvert-\lvert B \rvert \lvert C \rvert = 0.\]

                \item 若 $\lvert A \rvert = 0$，只需证 $\lvert B \rvert \lvert C \rvert = 0$. 若 $\lvert B \rvert \neq 0$，则由
                    \[\begin{pmatrix}
                            E        & O \\
                            -DB^{-1} & E
                        \end{pmatrix} \begin{pmatrix}
                            A & B \\
                            C & D
                        \end{pmatrix} \begin{pmatrix}
                            E        & O \\
                            -B^{-1}A & E
                        \end{pmatrix} = \begin{pmatrix}
                            O          & B \\
                            C-DB^{-1}A & O
                        \end{pmatrix},\]
                    有 $C-DB^{-1}A = O$. 注意到 $\lvert A \rvert = 0$，故 \[\lvert C \rvert = \lvert DB^{-1}A \rvert = \lvert D \rvert \lvert B^{-1} \rvert \lvert A \rvert = 0.\] 同理可证若 $\lvert C \rvert \neq 0$，则 $\lvert B \rvert = 0$.
            \end{enumerate}
            综上，结论成立.
\end{solve}

\begin{exercise}[补充训练 L10-C-9]\label{prob:L10-C-9}
设$A$为$n$阶实对称矩阵（$n>1$），$|A|=0$，证明：$A_{ii}A_{jj}=(A_{ij})^2(i,j=1,\ldots,n)$.
\end{exercise}
\noindent 来源：\href{https://github.com/yhwu-is/Linear-Algebra-Left-Undone/blob/new/%E8%AE%B2%E4%B9%89/%E4%B8%93%E9%A2%98/10%20%E8%A1%8C%E5%88%97%E5%BC%8F.tex}{L10-C-9}。
\begin{solve}
由 $|A| = 0$ 可知 $r(A) < n$. 利用结论
            \[
                r(\operatorname{adj}(A))=\begin{cases}
                    n & r(A)=n     \\
                    1 & r(A)=n-1   \\
                    0 & r(A) < n-1
                \end{cases}
            \]
            可知 $r(\operatorname{adj}(A)) \leqslant 1$.

            若 $r(\operatorname{adj}(A)) = 0$，则 $\operatorname{adj}(A) = O$，故 $A_{ii} = A_{jj} = A_{ij} = 0$，自然有 $A_{ii}A_{jj}=(A_{ij})^2$；

            若 $r(\operatorname{adj}(A)) = 1$，则 $\operatorname{adj}(A)$ 的任意两个列向量线性相关，可得 $A_{ii}A_{jj} = A_{ij}A_{ji}$. 又由 $A$ 是实对称矩阵，有 $A_{ij} = A_{ji}$，因此 $A_{ii}A_{jj} = (A_{ij})^2$.
\end{solve}

\begin{exercise}[补充训练 L10-C-10]\label{prob:L10-C-10}
分别求下列分块矩阵的伴随矩阵，所有块的尺寸均相容：
\begin{enumerate}
\item $M=\begin{pmatrix}A&C\\O&B\end{pmatrix}$，其中 $A,B$ 为方阵，允许奇异；
\item $N=\begin{pmatrix}A&B\\C&D\end{pmatrix}$，其中 $A,D$ 为方阵，且 $A$ 可逆。
\end{enumerate}
\end{exercise}
\noindent 来源：\href{https://github.com/yhwu-is/Linear-Algebra-Left-Undone/blob/new/%E8%AE%B2%E4%B9%89/%E4%B8%93%E9%A2%98/10%20%E8%A1%8C%E5%88%97%E5%BC%8F.tex}{L10-C-10}。
\begin{solve}
\begin{enumerate}
\item 当 $A,B$ 均可逆时，利用分块求逆及 $\det M=\det A\det B$ 得
\[\operatorname{adj}(M)=\begin{pmatrix}\det(B)\operatorname{adj}(A)&-\operatorname{adj}(A)C\operatorname{adj}(B)\\O&\det(A)\operatorname{adj}(B)\end{pmatrix}.\]
公式对奇异块仍成立：将 $A,B$ 同时换成 $A+tI,B+tI$ 后，除有限多个 $t$ 外都可逆；等式两边的每个元素都是 $t$ 的多项式，在无限多个值上相等，故恒等。取 $t=0$ 即得一般结论。
\item 记 Schur 补 $S=D-CA^{-1}B$。消元分解给出 $\det N=\det A\det S$。当 $S$ 可逆时，分块求逆后乘行列式得到
\[\operatorname{adj}(N)=\det(A)\begin{pmatrix}
\det(S)A^{-1}+A^{-1}B\operatorname{adj}(S)CA^{-1}&-A^{-1}B\operatorname{adj}(S)\\
-\operatorname{adj}(S)CA^{-1}&\operatorname{adj}(S)
\end{pmatrix}.\]
右端不含 $S^{-1}$，所以也适用于奇异的 $S$。严格地说，固定 $A,B,C$ 并将 $D$ 换成 $D+tI$，则 $S$ 换成 $S+tI$；两边再次是关于 $t$ 的多项式恒等式，令 $t=0$ 即可。题设仅要求 $A$ 可逆，无须额外假定整个 $N$ 可逆。
\end{enumerate}
\end{solve}

\begin{exercise}[补充训练 L10-C-11]\label{prob:L10-C-11}
以下矩阵均在复数域上考虑，且 $n\ge2$。下面三个小问探讨伴随矩阵的反问题，即对任意给定的$n$阶方阵$B$，是否存在$n$阶方阵$A$使得$\operatorname{adj}(A)=B$.
        \begin{enumerate}
            \item 证明：若$n=2$，则存在唯一的2阶方阵$A$使得$\operatorname{adj}(A)=B$；

            \item 证明：若$n > 2$，则存在$n$阶方阵$A$使得$\operatorname{adj}(A)=B$的充要条件为$r(B) \in \{0,1,n\}$，并且
                  \begin{enumerate}
                      \item $r(B)=n$时，$A=\sqrt[n-1]{|B|}B^{-1}$；

                      \item $r(B)=1$时，$A=Q^{-1}\begin{pmatrix}
                                    0 & O \\ O & X_{n-1}
                                \end{pmatrix}P^{-1}$，且$|X_{n-1}|=|PQ|$，$B=P\begin{pmatrix}
                                    1 & O \\ O & O
                                \end{pmatrix}Q$；
                  \end{enumerate}

            \item 设$A=\begin{pmatrix}
                          1 & 1 & 1 \\ 1 & 1 & 1 \\ 1 & 1 & 1
                      \end{pmatrix}$，求矩阵$B$使得$\operatorname{adj}(B)=A$.
        \end{enumerate}
\end{exercise}
\noindent 来源：\href{https://github.com/yhwu-is/Linear-Algebra-Left-Undone/blob/new/%E8%AE%B2%E4%B9%89/%E4%B8%93%E9%A2%98/10%20%E8%A1%8C%E5%88%97%E5%BC%8F.tex}{L10-C-11}。
\begin{solve}
\begin{enumerate}
\item 二阶矩阵满足 $\operatorname{adj}(\operatorname{adj}A)=A$，因此唯一解为 $A=\operatorname{adj}B$。
\item 对 $n>2$，若 $r(A)=n$，则其伴随矩阵满秩；若 $r(A)=n-1$，则至少一个 $n-1$ 阶子式非零，且 $A\operatorname{adj}A=O$，故伴随矩阵秩为 $1$；若 $r(A)\le n-2$，则全部代数余子式为零。这证明了秩条件的必要性。
若 $B$ 满秩，所有解恰为 $A=cB^{-1}$，其中 $c^{n-1}=\det B$。必要性由 $\operatorname{adj}A=\det(A)A^{-1}$ 得到；充分性直接代入即可。在复数域上，该方程有 $n-1$ 个不同的非零根，故存在解。数域不能省略：例如实数域上偶次根未必存在。
若 $r(B)=1$，写 $B=P\operatorname{diag}(1,O)Q$，其中 $P,Q$ 可逆。取任意可逆的 $n-1$ 阶矩阵 $X$ 使 $\det X=\det(PQ)$，令 $A=Q^{-1}\operatorname{diag}(0,X)P^{-1}$。利用 $\operatorname{adj}(UV)=\operatorname{adj}(V)\operatorname{adj}(U)$ 可得
\[\operatorname{adj}A=\frac{\det X}{\det(PQ)}P\operatorname{diag}(1,O)Q=B.\]
若 $B=O$，取任何秩不超过 $n-2$ 的 $A$ 即可。三种情形都可构造所需矩阵。
\item 令 $e=(1,1,1)^{\mathrm T}$，目标是 $\operatorname{adj}B=ee^{\mathrm T}$。取
\[Q=\begin{pmatrix}1/\sqrt3&1/\sqrt2&1/\sqrt6\\1/\sqrt3&-1/\sqrt2&1/\sqrt6\\1/\sqrt3&0&-2/\sqrt6\end{pmatrix}.\]
则全部解可写为 $B=Q\operatorname{diag}(0,C)Q^{\mathrm T}$，其中 $C$ 为任意满足 $\det C=3$ 的二阶矩阵。事实上，相应伴随矩阵为 $Q\operatorname{diag}(3,0,0)Q^{\mathrm T}=ee^{\mathrm T}$；反之，由秩为 $2$ 及 $Be=0,e^{\mathrm T}B=0$ 可知在这个正交基下第一行和第一列都为零，剩余二阶块必须满足此行列式条件。例如可取 $C=\operatorname{diag}(1,3)$。\end{enumerate}
\end{solve}

\begin{exercise}[补充训练 L10-C-12]\label{prob:L10-C-12}
证明$\begin{vmatrix}
                a      & c      & c      & \cdots & c      \\
                b      & a      & c      & \cdots & c      \\
                b      & b      & a      & \cdots & c      \\
                \vdots & \vdots & \vdots & \ddots & \vdots \\
                b      & b      & b      & \cdots & a
            \end{vmatrix}=\dfrac{b(a-c)^{n}-c(a-b)^{n}}{b-c}$，其中 $b \neq c$, 等式左端是 $n$ 阶行列式.
\end{exercise}
\noindent 来源：\href{https://github.com/yhwu-is/Linear-Algebra-Left-Undone/blob/new/%E8%AE%B2%E4%B9%89/%E4%B8%93%E9%A2%98/10%20%E8%A1%8C%E5%88%97%E5%BC%8F.tex}{L10-C-12}。
\begin{solve}
我们将第 $i$ 行乘 $(-1)$ 加到第 $i-1$ 行（$i$ 依次取 $n, n-1, \ldots, 2$），再对第 1 列展开,
            \[D_{n}=\begin{vmatrix}
                    a-b    & c-a    & 0      & \cdots & 0      \\
                    0      & a-b    & c-a    & \cdots & 0      \\
                    \vdots & \vdots & \vdots & \ddots & \vdots \\
                    0      & 0      & 0      & \cdots & c-a    \\
                    b      & b      & b      & \cdots & a
                \end{vmatrix}=(a-b)D_{n-1}+(-1)^{n+1}b(c-a)^{n-1},\]
            有了递推式，但是递推式中有$(-1)^{n+1}$，递推比较麻烦可能还需要处理. 如果用数归证明的话应该可以直接下手了，但是有更聪明的办法：取转置，则行列式的值不变，但是 $b$ 与 $c$ 的位置交换，由此可以写出 $D_n$ 的第二个递推式：$D_n=(a-c)D_{n-1}+(-1)^{n+1}c(b-a)^{n-1}$. 两个递推式分别乘以 $a-c$ 与 $a-b$ 后相减即有
            \[
                D_n=\dfrac{b(a-c)^{n}-c(a-b)^{n}}{b-c}.
            \]
\end{solve}

\begin{exercise}[补充训练 L10-C-13]\label{prob:L10-C-13}
计算$n$阶行列式
        \[|A|=\begin{vmatrix}
                a_1b_1 & a_1b_2 & a_1b_3 & \cdots & a_1b_n \\
                a_2b_1 & a_2b_2 & a_2b_3 & \cdots & a_2b_n \\
                a_3b_1 & a_3b_2 & a_3b_3 & \cdots & a_3b_n \\
                \vdots & \vdots & \vdots & \ddots & \vdots \\
                a_nb_1 & a_nb_2 & a_nb_3 & \cdots & a_nb_n
            \end{vmatrix}.\]
\end{exercise}
\noindent 来源：\href{https://github.com/yhwu-is/Linear-Algebra-Left-Undone/blob/new/%E8%AE%B2%E4%B9%89/%E4%B8%93%E9%A2%98/10%20%E8%A1%8C%E5%88%97%E5%BC%8F.tex}{L10-C-13}。
\begin{solve}
\begin{enumerate}
                \item $n=1$ 时，显然有 $|A|=a_1b_1$.

                \item $n \geqslant 2$ 时，
                    \[
                        A = \begin{pmatrix}
                            a_1 \\ a_2 \\ \vdots \\ a_n
                        \end{pmatrix} \begin{pmatrix}
                            b_1 & b_2 & \cdots & b_n
                        \end{pmatrix},
                    \]
                    故 $r(A) = 1 < n$，$A$ 不可逆，因此 $|A| = 0$.
            \end{enumerate}
\end{solve}

\begin{exercise}[补充训练 L10-C-14]\label{prob:L10-C-14}
计算$n$阶行列式
        \[|A|=\begin{vmatrix}
                1+a_1^2 & a_1a_2  & \cdots & a_1a_n  \\
                a_2a_1  & 1+a_2^2 & \cdots & a_2a_n  \\
                \vdots  & \vdots  & \ddots & \vdots  \\
                a_na_1  & a_na_2  & \cdots & 1+a_n^2
            \end{vmatrix}.\]
\end{exercise}
\noindent 来源：\href{https://github.com/yhwu-is/Linear-Algebra-Left-Undone/blob/new/%E8%AE%B2%E4%B9%89/%E4%B8%93%E9%A2%98/10%20%E8%A1%8C%E5%88%97%E5%BC%8F.tex}{L10-C-14}。
\begin{solve}
设
            \[
                D_n = \begin{vmatrix}
                    1+a_1^2 & a_1a_2  & \cdots & a_1a_n  \\
                    a_2a_1  & 1+a_2^2 & \cdots & a_2a_n  \\
                    \vdots  & \vdots  & \ddots & \vdots  \\
                    a_na_1  & a_na_2  & \cdots & 1+a_n^2
                \end{vmatrix},
            \]

            利用行列式公理化定义中的线性性，有
            \begin{align*}
                D_n &= \begin{vmatrix}
                    1+a_1^2 & a_1a_2  & \cdots & a_1a_n  \\
                    a_2a_1  & 1+a_2^2 & \cdots & a_2a_n  \\
                    \vdots  & \vdots  & \ddots & \vdots  \\
                    a_na_1  & a_na_2  & \cdots & a_n^2
                \end{vmatrix} + \begin{vmatrix}
                    1+a_1^2 & a_1a_2  & \cdots & 0       \\
                    a_2a_1  & 1+a_2^2 & \cdots & 0       \\
                    \vdots  & \vdots  & \ddots & \vdots  \\
                    a_na_1  & a_na_2  & \cdots & 1
                \end{vmatrix} \\
                &= a_n \begin{vmatrix}
                    1+a_1^2 & a_1a_2  & \cdots & a_1     \\
                    a_2a_1  & 1+a_2^2 & \cdots & a_2     \\
                    \vdots  & \vdots  & \ddots & \vdots  \\
                    a_na_1  & a_na_2  & \cdots & a_n
                \end{vmatrix} + D_{n-1} \\
                &= a_n \begin{vmatrix}
                    1      & 0      & \cdots & a_1     \\
                    0      & 1      & \cdots & a_2     \\
                    \vdots & \vdots & \ddots & \vdots  \\
                    0      & 0      & \cdots & a_n
                \end{vmatrix} \\
                &= D_{n-1} + a_n^2.
            \end{align*}
            而 $D_1 = 1 + a_1^2$, 因此
            \[
                |A| = D_n = 1 + \sum_{i=1}^n a_i^2.
            \]
\end{solve}

\begin{exercise}[补充训练 L10-C-15]\label{prob:L10-C-15}
计算$n$阶行列式
        \[|A|=\begin{vmatrix}
                1+x_1  & 1+x_1^2 & \cdots & 1+x_1^n \\
                1+x_2  & 1+x_2^2 & \cdots & 1+x_2^n \\
                \vdots & \vdots  & \ddots & \vdots  \\
                1+x_n  & 1+x_n^2 & \cdots & 1+x_n^n
            \end{vmatrix}.\]
\end{exercise}
\noindent 来源：\href{https://github.com/yhwu-is/Linear-Algebra-Left-Undone/blob/new/%E8%AE%B2%E4%B9%89/%E4%B8%93%E9%A2%98/10%20%E8%A1%8C%E5%88%97%E5%BC%8F.tex}{L10-C-15}。
\begin{solve}
使用硬拆法，将 $A$ 的第 $i$ 行分解为 $(1, 1, \ldots, 1) + (x_i, x_i^2, \ldots, x_i^n)$. 考虑到拆出的行列式中包含不少于两列为 $(1, 1, \ldots, 1)$ 的项均为 $0$，因此可将 $|A|$ 化简为
            \[
                \begin{vmatrix}
                    x_1    & x_1^2  & \cdots & x_1^n  \\
                    x_2    & x_2^2  & \cdots & x_2^n  \\
                    \vdots & \vdots & \ddots & \vdots \\
                    x_n    & x_n^2  & \cdots & x_n^n
                \end{vmatrix} + \sum_{i=1}^n \begin{vmatrix}
                    x_1     & x_1^2     & \cdots & x_1^n     \\
                    \vdots  & \vdots    &        & \vdots    \\
                    x_{i-1} & x_{i-1}^2 & \cdots & x_{i-1}^n \\
                    1       & 1         & \cdots & 1         \\
                    x_{i+1} & x_{i+1}^2 & \cdots & x_{i+1}^n \\
                    \vdots  & \vdots    &        & \vdots    \\
                    x_n     & x_n^2     & \cdots & x_n^n
                \end{vmatrix}.
            \]

            上式中第一项可利用 Vandermonde 行列式化简：
            \[
                \begin{vmatrix}
                    x_1    & x_1^2  & \cdots & x_1^n  \\
                    x_2    & x_2^2  & \cdots & x_2^n  \\
                    \vdots & \vdots & \ddots & \vdots \\
                    x_n    & x_n^2  & \cdots & x_n^n
                \end{vmatrix} = \prod_{i=1}^n x_i \begin{vmatrix}
                    1      & x_1    & x_1^2  & \cdots & x_1^n  \\
                    1      & x_2    & x_2^2  & \cdots & x_2^n  \\
                    \vdots & \vdots & \vdots & \ddots & \vdots \\
                    1      & x_n    & x_n^2  & \cdots & x_n^n
                \end{vmatrix} = \left(\prod_{i=1}^n x_i\right) \left(\prod_{1 \leqslant i < j \leqslant n} (x_j - x_i)\right).
            \]

            对于第二项，
            \begin{align*}
                \begin{vmatrix}
                    x_1     & x_1^2     & \cdots & x_1^n     \\
                    \vdots  & \vdots    &        & \vdots    \\
                    x_{i-1} & x_{i-1}^2 & \cdots & x_{i-1}^n \\
                    1       & 1         & \cdots & 1         \\
                    x_{i+1} & x_{i+1}^2 & \cdots & x_{i+1}^n \\
                    \vdots  & \vdots    &        & \vdots    \\
                    x_n     & x_n^2     & \cdots & x_n^n
                \end{vmatrix}
                &= \prod_{j \neq i} x_j \begin{vmatrix}
                    1      & x_1     & \cdots & x_1^{n-1}     \\
                    \vdots & \vdots  &        & \vdots        \\
                    1      & x_{i-1} & \cdots & x_{i-1}^{n-1} \\
                    1      & 1^1     & \cdots & 1^{n-1}       \\
                    1      & x_{i+1} & \cdots & x_{i-1}^{n+1} \\
                    \vdots & \vdots  &        & \vdots        \\
                    1      & x_n     & \cdots & x_n^{n-1}
                \end{vmatrix} \\
                &= \prod_{j \neq i} x_j \cdot \prod_{1 \leqslant k < j \leqslant n} (x_j - x_k) \cdot \prod_{1 \leqslant j < i} \frac{1 - x_j}{x_i - x_j} \cdot \prod_{i < j \leqslant n} \frac{x_j - 1}{x_j - x_i}  \\
                &= \prod_{j \neq i} x_j \cdot \prod_{j \neq i} \frac{x_j - 1}{x_j - x_i} \cdot \prod_{1 \leqslant k < j \leqslant n} (x_j - x_k).
            \end{align*}

            综上可得，
            \[
                |A| = \left( 1 + \sum_{i=1}^n \frac{1}{x_i} \prod_{j \neq i} \frac{x_j - 1}{x_j - x_i} \right) \prod_{i=1}^n x_i \cdot \prod_{1 \leqslant k < j \leqslant n} (x_j - x_k).
            \]
\end{solve}

\begin{exercise}[补充训练 L10-C-16]\label{prob:L10-C-16}
计算$n$阶行列式，其中$x,y,z$为任意实常数：
        \[|A|=\begin{vmatrix}
                x_1    & y      & y      & \cdots & y      \\
                z      & x_2    & y      & \cdots & y      \\
                z      & z      & x_3    & \cdots & y      \\
                \vdots & \vdots & \vdots & \ddots & \vdots \\
                z      & z      & z      & \cdots & x_n
            \end{vmatrix}.\]
\end{exercise}
\noindent 来源：\href{https://github.com/yhwu-is/Linear-Algebra-Left-Undone/blob/new/%E8%AE%B2%E4%B9%89/%E4%B8%93%E9%A2%98/10%20%E8%A1%8C%E5%88%97%E5%BC%8F.tex}{L10-C-16}。
\begin{solve}
\begin{enumerate}
                \item 若 $y = z$，
                    则
                    \[
                        |A| = \begin{vmatrix}
                            x_1    & y      & y      & \cdots & y      \\
                            y      & x_2    & y      & \cdots & y      \\
                            y      & y      & x_3    & \cdots & y      \\
                            \vdots & \vdots & \vdots & \ddots & \vdots \\
                            y      & y      & y      & \cdots & x_n
                        \end{vmatrix}.
                    \]

                    \begin{enumerate}
                        \item 若存在 $i \neq j$，使得 $x_i = x_j = y$，则 $A$ 中有两列元素全为 $y$，因此 $|A|$ 为 $0$.

                        \item 若存在唯一 $i$，使得 $x_i = y$，则
                            \begin{align*}
                                |A| &= (-1)^{i - 1 + i - 1} \begin{vmatrix}
                                    x_i    & y      & y      & \cdots & y      \\
                                    y      & x_1    & y      & \cdots & y      \\
                                    y      & y      & x_2    & \cdots & y      \\
                                    \vdots & \vdots & \vdots & \ddots & \vdots \\
                                    y      & y      & y      & \cdots & x_n
                                \end{vmatrix}
                                = \begin{vmatrix}
                                    y      & y       & y       & \cdots & y       \\
                                    0      & x_1 - y & 0       & \cdots & 0       \\
                                    0      & 0       & x_2 - y & \cdots & 0       \\
                                    \vdots & \vdots  & \vdots  & \ddots & \vdots  \\
                                    0      & 0       & 0       & \cdots & x_n - y
                                \end{vmatrix} \\
                                &= y \prod_{j \neq i} (x_j - y).
                            \end{align*}

                            \item 若不存在 $i$，使得 $x_i = y$，则
                                \begin{align*}
                                    |A| &= \begin{vmatrix}
                                        1      & y      & y      & y      & \cdots & y      \\
                                        0      & x_1    & y      & y      & \cdots & y      \\
                                        0      & y      & x_2    & y      & \cdots & y      \\
                                        0      & y      & y      & x_3    & \cdots & y      \\
                                        \vdots & \vdots & \vdots & \vdots & \ddots & \vdots \\
                                        0      & y      & y      & y      & \cdots & x_n
                                    \end{vmatrix}_{n+1} \\
                                    &= \begin{vmatrix}
                                        1      & y      & y      & y      & \cdots & y      \\
                                        -1     & x_1-y  & 0      & 0      & \cdots & 0      \\
                                        -1     & 0      & x_2-y  & 0      & \cdots & 0      \\
                                        -1     & 0      & 0      & x_3-y  & \cdots & 0      \\
                                        \vdots & \vdots & \vdots & \vdots & \ddots & \vdots \\
                                        -1     & 0      & 0      & 0      & \cdots & x_n-y
                                    \end{vmatrix}_{n+1} \\
                                    &= \begin{vmatrix}
                                        1 + \sum_{j=1}^n \frac{y}{x_j-y} & y      & y      & y      & \cdots & y      \\
                                        0                                & x_1-y  & 0      & 0      & \cdots & 0      \\
                                        0                                & 0      & x_2-y  & 0      & \cdots & 0      \\
                                        0                                & 0      & 0      & x_3-y  & \cdots & 0      \\
                                        \vdots                           & \vdots & \vdots & \vdots & \ddots & \vdots \\
                                        0                                & 0      & 0      & 0      & \cdots & x_n-y
                                    \end{vmatrix} \\
                                    &= \prod_{j=1}^n (x_j - y) + y\sum_{i=1}^n \prod_{j \neq i} (x_j - y).
                                \end{align*}
                    \end{enumerate}
                    综上，当 $y = z$ 时，
                    \[
                        |A| = \prod_{j=1}^n (x_j - y) + y\sum_{i=1}^n \prod_{j \neq i} (x_j - y).
                    \]

                \item 若 $y \neq z$，则可按照与第 12 题完全相同的方法进行处理，最后得到
                    \[
                        |A| = \frac{z \prod_{i=1}^n (x_i - y) - y \prod_{i=1}^n (x_i - z)}{z - y}.
                    \]
            \end{enumerate}
\end{solve}

\begin{exercise}[补充训练 L10-例-2]\label{prob:L10-例-2}
计算行列式 $D = \begin{vmatrix}
        1 & 2 \\
        2 & 3
    \end{vmatrix}$.
\end{exercise}
\noindent 来源：\href{https://github.com/yhwu-is/Linear-Algebra-Left-Undone/blob/new/%E8%AE%B2%E4%B9%89/%E4%B8%93%E9%A2%98/10%20%E8%A1%8C%E5%88%97%E5%BC%8F.tex}{L10-例-2}。
\begin{solve}
设 $\alpha_1 = (1, 2)$, $\alpha_2 = (2, 3)$，将行列式 $D$ 写成向量形式并用公理化定义展开可得：
    \begin{align*}
        D &= D(\alpha_1, \alpha_2) = D(e_1 + 2e_2, 2e_1 + 3e_2) \\
          &= D(e_1, 2e_1 + 3e_2) + 2 D(e_2, 2e_1 + 3e_2) \\
          &= 2D(e_1, e_1) + 3D(e_1, e_2) + 4D(e_2, e_1) + 6D(e_2, e_2) \\
          &= 3 - 4 = -1.
    \end{align*}
\end{solve}

\begin{exercise}[补充训练 L10-例-3]\label{prob:L10-例-3}
设向量$\alpha_1,\alpha_2,\beta_1,\beta_2$为三维列向量，又$A=(\alpha_1,\alpha_2,\beta_1),B=(\alpha_1,\alpha_2,\beta_2)$，且$|A|=3$，$|B|=2$，求$|2A+3B|$.
\end{exercise}
\noindent 来源：\href{https://github.com/yhwu-is/Linear-Algebra-Left-Undone/blob/new/%E8%AE%B2%E4%B9%89/%E4%B8%93%E9%A2%98/10%20%E8%A1%8C%E5%88%97%E5%BC%8F.tex}{L10-例-3}。
\begin{solve}
$2A+3B=(2\alpha_1+3\alpha_1,2\alpha_2+3\alpha_2,2\beta_1+3\beta_2)=(5\alpha_1,5\alpha_2,2\beta_1+3\beta_2)$，因此
    \begin{align*}
        |2A+3B| & =D(5\alpha_1,5\alpha_2,2\beta_1+3\beta_2)                        \\
                & =25D(\alpha_1,\alpha_2,2\beta_1+3\beta_2)                        \\
                & =25(2D(\alpha_1,\alpha_2,\beta_1)+3D(\alpha_1,\alpha_2,\beta_2)) \\
                & =25(2|A|+3|B|)=300.
    \end{align*}
\end{solve}

\begin{exercise}[补充训练 L10-例-5]\label{prob:L10-例-5}
求排列 $(3,1,4,2)$ 的逆序数.
\end{exercise}
\noindent 来源：\href{https://github.com/yhwu-is/Linear-Algebra-Left-Undone/blob/new/%E8%AE%B2%E4%B9%89/%E4%B8%93%E9%A2%98/10%20%E8%A1%8C%E5%88%97%E5%BC%8F.tex}{L10-例-5}。
\begin{solve}
我们可以直接计算：$3$ 后面有 $2$ 个数比 $3$ 小，$1$ 后面没有数比 $1$ 小，$4$ 后面有 $1$ 个数比 $4$ 小，$2$ 后面没有数比 $2$ 小，因此逆序数为 $2+0+1+0=3$.
\end{solve}

\begin{exercise}[补充训练 L10-例-11]\label{prob:L10-例-11}
求证以下命题：
    \begin{enumerate}
        \item 奇数阶反对称矩阵不可逆；

        \item 若$A$是$n$阶可逆对称矩阵，$B$是$n$阶反对称矩阵，则当$n$为奇数时，齐次线性方程组$(AB)X=O$有非零解.
    \end{enumerate}
\end{exercise}
\noindent 来源：\href{https://github.com/yhwu-is/Linear-Algebra-Left-Undone/blob/new/%E8%AE%B2%E4%B9%89/%E4%B8%93%E9%A2%98/10%20%E8%A1%8C%E5%88%97%E5%BC%8F.tex}{L10-例-11}。
\begin{solve}
\begin{enumerate}
        \item 设$A$为$n$阶反对称矩阵，则$A^\mathrm{T}=-A$，因此$|A|=|A^\mathrm{T}|=|-A|=(-1)^n|A|=-|A|$，从而$|A|=0$，故$A$不可逆.
        \item 反证法，则$(AB)X=O$只有零解，由原讲义“可逆等价条件”处 的 原讲义“item:11:可逆等价条件:4”处 可知$AB$可逆，故$|AB|=|A||B|\neq 0$，这说明$|B|\neq 0$，即$B$可逆. 但由于$B$为奇数阶反对称矩阵，因此$B$不可逆，矛盾. 因此原方程组有非零解.
    \end{enumerate}
\end{solve}

\begin{exercise}[补充训练 L10-例-12]\label{prob:L10-例-12}
设$A$为$n$阶正交矩阵，即$AA^\mathrm{T}=A^\mathrm{T}A=E$，且$|A|<0$，证明：$|E+A|=0$.
\end{exercise}
\noindent 来源：\href{https://github.com/yhwu-is/Linear-Algebra-Left-Undone/blob/new/%E8%AE%B2%E4%B9%89/%E4%B8%93%E9%A2%98/10%20%E8%A1%8C%E5%88%97%E5%BC%8F.tex}{L10-例-12}。
\begin{solve}
正交矩阵满足 $AA^{\mathrm{T}} = A^{\mathrm{T}}A = E$，所以 $\lvert AA^{\mathrm{T}} \rvert = \lvert A \rvert^2 = \lvert E \rvert = 1$. 而 $\lvert A \rvert < 0$，所以 $\lvert A \rvert = -1$.
    \[
        \lvert E+A \rvert = \lvert AA^{\mathrm{T}}+A \rvert = \lvert A \rvert \cdot \lvert A^{\mathrm{T}}+E \rvert = -\lvert (A+E)^{\mathrm{T}} \rvert = -\lvert E+A \rvert.
    \]

    故 $\lvert E+A \rvert = 0$.
\end{solve}

\begin{exercise}[补充训练 L10-例-17]\label{prob:L10-例-17}
判断矩阵$\begin{pmatrix}
            1 & 2 & 3 \\ 2 & 1 & 2 \\ 1 & 3 & 3
        \end{pmatrix}$是否可逆. 若可逆，利用伴随矩阵求其逆矩阵.
\end{exercise}
\noindent 来源：\href{https://github.com/yhwu-is/Linear-Algebra-Left-Undone/blob/new/%E8%AE%B2%E4%B9%89/%E4%B8%93%E9%A2%98/10%20%E8%A1%8C%E5%88%97%E5%BC%8F.tex}{L10-例-17}。
\begin{solve}
我们使用行列式来判断矩阵是否可逆：设 $A=\begin{pmatrix}
        1 & 2 & 3 \\ 2 & 1 & 2 \\ 1 & 3 & 3
    \end{pmatrix}$，则有 $|A| = 4$，因此 $A$ 可逆.

    接着我们使用伴随矩阵求 $A$ 的逆矩阵. 首先求 $|A|$ 中各元素的代数余子式：

    \[
        A_{11} = \begin{vmatrix}
            1 & 2 \\ 3 & 3
        \end{vmatrix} = -3, \quad
        A_{12} = - \begin{vmatrix}
            2 & 2 \\ 1 & 3
        \end{vmatrix} = -4, \quad
        A_{13} = \begin{vmatrix}
            2 & 1 \\ 1 & 3
        \end{vmatrix} = 5.
    \]

    类似地，我们有
    \begin{gather*}
        A_{21} = 3, \quad A_{22} = 0, \quad A_{23} = -1, \\
        A_{31} = 1, \quad A_{32} = 4, \quad A_{33} = -3,
    \end{gather*}

    由此可得 $A$ 的伴随矩阵 
\[\operatorname{adj}(A) = \begin{pmatrix}
            A_{11} & A_{21} & A_{31} \\
            A_{12} & A_{22} & A_{32} \\
            A_{13} & A_{23} & A_{33}
        \end{pmatrix} = \begin{pmatrix}
            -3 & 3 & 1 \\
            -4 & 0 & 4 \\
            5 & -1 & -3
        \end{pmatrix}\]
.

    因此，

    \[
        A^{-1} = \dfrac{\operatorname{adj}(A)}{|A|} = \dfrac{1}{4} \begin{pmatrix}
            -3 & 3 & 1 \\
            -4 & 0 & 4 \\
            5 & -1 & -3
            \end{pmatrix}.
    \]
\end{solve}

\begin{exercise}[补充训练 L10-例-18]\label{prob:L10-例-18}
已知 $\operatorname{adj}(A)=\begin{pmatrix}
            1 & -2 & 1 \\ 0 & 2 & -2 \\ -1 & 2 & 1
        \end{pmatrix}$，求 $A$.
\end{exercise}
\noindent 来源：\href{https://github.com/yhwu-is/Linear-Algebra-Left-Undone/blob/new/%E8%AE%B2%E4%B9%89/%E4%B8%93%E9%A2%98/10%20%E8%A1%8C%E5%88%97%E5%BC%8F.tex}{L10-例-18}。
\begin{solve}
不难求得 $|\operatorname{adj}(A)| = 4$，根据伴随矩阵行列式的性质有，$|\operatorname{adj}(A)| = |A|^2$，因此 $|A| = \pm 2$. 若 $|A| = 2$，则根据 $A\operatorname{adj}(A) = |A|E$，有

    \[A^{-1} = |A|^{-1}\operatorname{adj}(A) = \frac{1}{2}\begin{pmatrix}
            1 & -2 & 1 \\ 0 & 2 & -2 \\ -1 & 2 & 1
        \end{pmatrix} = \begin{pmatrix}
            \dfrac{1}{2} & -1 & \dfrac{1}{2} \\ 0 & 1 & -1 \\ -\dfrac{1}{2} & 1 & \dfrac{1}{2}
        \end{pmatrix}.\]

    因此 $A = (A^{-1})^{-1} = \begin{pmatrix}
            3 & 2 & 1 \\ 1 & 1 & 1 \\ 1 & 0 & 1
        \end{pmatrix}$.

    若 $|A| = -2$，则同理可以求得 $A = -\begin{pmatrix}
            3 & 2 & 1 \\ 1 & 1 & 1 \\ 1 & 0 & 1
        \end{pmatrix}$.
\end{solve}

\begin{exercise}[补充训练 L10-例-19]\label{prob:L10-例-19}
写出矩阵$\begin{pmatrix}
            1 & 5 & -2 \\ 2 & 3 & 4 \\ -1 & -3 & 0
        \end{pmatrix}$的所有一阶、二阶子式、主子式和顺序主子式.
\end{exercise}
\noindent 来源：\href{https://github.com/yhwu-is/Linear-Algebra-Left-Undone/blob/new/%E8%AE%B2%E4%B9%89/%E4%B8%93%E9%A2%98/10%20%E8%A1%8C%E5%88%97%E5%BC%8F.tex}{L10-例-19}。
\begin{solve}
\begin{enumerate}
        \item 一阶子式：根据子式定义，任意1行1列交点组成的1个元素就是一阶子式，即所有的元素本身都是一阶子式；

        \item 二阶子式：根据子式定义，任意2行2列交点组成的4个元素排成的行列式就是二阶子式，即
              \[\begin{vmatrix}
                      1 & 5 \\
                      2 & 3
                  \end{vmatrix},\begin{vmatrix}
                      1 & -2 \\
                      2 & 4
                  \end{vmatrix},\begin{vmatrix}
                      5 & -2 \\
                      3 & 4
                  \end{vmatrix},\begin{vmatrix}
                      1  & 5  \\
                      -1 & -3
                  \end{vmatrix},\begin{vmatrix}
                      1  & -2 \\
                      -1 & 0
                  \end{vmatrix},\]\[\begin{vmatrix}
                      5  & -2 \\
                      -3 & 0
                  \end{vmatrix},\begin{vmatrix}
                      2  & 3  \\
                      -1 & -3
                  \end{vmatrix},\begin{vmatrix}
                      2  & 4 \\
                      -1 & 0
                  \end{vmatrix},\begin{vmatrix}
                      3  & 4 \\
                      -3 & 0
                  \end{vmatrix};\]

        \item 主子式：根据主子式定义，要求取的行列号相同，故一阶主子式就是1行1列、2行2列、3行3列的元素，二阶主子式就是选取1、2/1、3/2、3行与列构成的子式，即
              \[\begin{vmatrix}
                      1 & 5 \\ 2 & 3
                  \end{vmatrix},\begin{vmatrix}
                      1 & -2 \\ -1 & 0
                  \end{vmatrix},\begin{vmatrix}
                      3 & 4 \\ -3 & 0
                  \end{vmatrix};\]
              三阶主子式就是矩阵本身对应的行列式，不再赘述.

        \item 顺序主子式根据定义，一阶就是第一行第一列的元素，二阶就是前两行前两节元素构成的子式，三阶就是本身的行列式.
    \end{enumerate}
\end{solve}

\begin{exercise}[补充训练 L10-例-20]\label{prob:L10-例-20}
利用定义求矩阵$\begin{pmatrix}
            1 & 1 & -1 & 3 \\ 1 & 2 & 1 & 1 \\ 3 & 3 & 0 & 4
        \end{pmatrix}$的行列式秩.
\end{exercise}
\noindent 来源：\href{https://github.com/yhwu-is/Linear-Algebra-Left-Undone/blob/new/%E8%AE%B2%E4%B9%89/%E4%B8%93%E9%A2%98/10%20%E8%A1%8C%E5%88%97%E5%BC%8F.tex}{L10-例-20}。
\begin{solve}
记该矩阵为$A$，由于$A$为3行4列矩阵，因此$r(A)\leqslant 3$. 又我们可以发现其三阶子式
    \[\begin{vmatrix}
            1 & -1 & 3 \\ 1 & 1 & 1 \\ 3 & 0 & 4
        \end{vmatrix}=-4\neq 0,\]
    故$r(A)\geqslant 3$，因此$r(A)=3$.
\end{solve}

\begin{exercise}[补充训练 L10-例-21]\label{prob:L10-例-21}
求多项式$f(x)=\begin{vmatrix}
            1 & a_1   & a_2     & a_3     \\
            1 & a_1+x & a_2     & a_3     \\
            1 & a_1   & a_2+x+1 & a_3     \\
            1 & a_1   & a_2     & a_3+x+2
        \end{vmatrix}$的所有零点.
\end{exercise}
\noindent 来源：\href{https://github.com/yhwu-is/Linear-Algebra-Left-Undone/blob/new/%E8%AE%B2%E4%B9%89/%E4%B8%93%E9%A2%98/10%20%E8%A1%8C%E5%88%97%E5%BC%8F.tex}{L10-例-21}。
\begin{solve}
事实上，本题可以直接展开求出四阶行列式的值然后解方程$f(x)=0$即可，但我们这里使用更为本质的方法. $f(x)=0$实际上就是行列式等于零，即此时$x$使得行列式中两列（或两行）线性相关了. 事实上，我们很容易发现$x=0$时，第一列与第二列成比例，故此时$f(x)=0$成立. 同理，$x=-1$和$x=-2$也是$f(x)=0$的解.

    事实上，我们知道这一行列式展开后是一个次数最高为3的多项式，因此$0,-1,-2$就是$f(x)$的所有零点.
\end{solve}

\begin{exercise}[补充训练 L10-例-22]\label{prob:L10-例-22}
计算行列式$D_{n+1}=\begin{vmatrix}
            1      & a_{1}       & a_{2}       & \cdots & a_n       \\
            1      & a_{1}+b_{1} & a_{2}       & \cdots & a_n       \\
            1      & a_{1}       & a_{2}+b_{2} & \cdots & a_n       \\
            \vdots & \vdots      & \vdots      & \ddots & \vdots    \\
            1      & a_{1}       & a_{2}       & \cdots & a_n+b_{n}
        \end{vmatrix}$.
\end{exercise}
\noindent 来源：\href{https://github.com/yhwu-is/Linear-Algebra-Left-Undone/blob/new/%E8%AE%B2%E4%B9%89/%E4%B8%93%E9%A2%98/10%20%E8%A1%8C%E5%88%97%E5%BC%8F.tex}{L10-例-22}。
\begin{solve}
观察这个行列式的特点，主对角线下方的元素与第 1 行元素对应相同，故用第 1 行的 $-1$ 倍加到下面各行便可使主对角线下方的元素全部变为0，即化为上三角形.
    \[ D_{n+1}=\begin{vmatrix}
            1 & a_{1} & a_{2} & \ldots & a_n   \\
              & b_{1} &       &        &       \\
              &       & b_{2} &        &       \\
              &       &       & \ddots &       \\
              &       &       &        & b_{n}
        \end{vmatrix}=\prod_{i=1}^n b_i \]
\end{solve}

\begin{exercise}[补充训练 L10-例-23]\label{prob:L10-例-23}
计算$n$阶行列式，其中$a_i\neq 0\enspace(i=2,\ldots,n)$：
    \[|A|=\begin{vmatrix}
            a_1    & b_2    & b_3    & \cdots & b_n    \\
            c_2    & a_2    & 0      & \cdots & 0      \\
            c_3    & 0      & a_3    & \cdots & 0      \\
            \vdots & \vdots & \vdots & \ddots & \vdots \\
            c_n    & 0      & 0      & \cdots & a_n
        \end{vmatrix}.\]
\end{exercise}
\noindent 来源：\href{https://github.com/yhwu-is/Linear-Algebra-Left-Undone/blob/new/%E8%AE%B2%E4%B9%89/%E4%B8%93%E9%A2%98/10%20%E8%A1%8C%E5%88%97%E5%BC%8F.tex}{L10-例-23}。
\begin{solve}
这一行列式因为形似箭头而被称为箭形行列式，将第$i\enspace(2\leqslant i\leqslant n)$列乘以$-c_i/a_i$加到第一列上可以消去第一列所有的$c_i$得到上三角矩阵，最后答案为
    \[
    (a_1-\sum\limits_{i=2}^n\dfrac{b_ic_i}{a_i})a_2a_3\cdots a_n.
    \]
\end{solve}

\begin{exercise}[补充训练 L10-例-24]\label{prob:L10-例-24}
计算$n$阶行列式，其中$a_i\neq 0,\enspace i=1,\ldots,n$：
    \[|A|=\begin{vmatrix}
            x_1-a_1 & x_2     & x_3     & \cdots & x_n     \\
            x_1     & x_2-a_2 & x_3     & \cdots & x_n     \\
            x_1     & x_2     & x_3-a_3 & \cdots & x_n     \\
            \vdots  & \vdots  & \vdots  & \ddots & \vdots  \\
            x_1     & x_2     & x_3     & \cdots & x_n-a_n
        \end{vmatrix}.\]
\end{exercise}
\noindent 来源：\href{https://github.com/yhwu-is/Linear-Algebra-Left-Undone/blob/new/%E8%AE%B2%E4%B9%89/%E4%B8%93%E9%A2%98/10%20%E8%A1%8C%E5%88%97%E5%BC%8F.tex}{L10-例-24}。
\begin{solve}
第一行乘以-1依次加到其它行上去可以得到一个箭形行列式，使用上一例题的结论可知答案为\[(-1)^{n-1}a_1a_2\cdots a_n(\sum\limits_{i=1}^n\dfrac{x_i}{a_i}-1).\]
\end{solve}

\begin{exercise}[补充训练 L10-例-26]\label{prob:L10-例-26}
计算 $n$ 阶行列式：
    \[
        |A| = \begin{vmatrix}
            \dfrac{1}{x_1 - 1} & 1      & x_1    & x_1^2  & \cdots & x_1^{n-2} \\[1.2em]
            \dfrac{1}{x_2 - 1} & 1      & x_2    & x_2^2  & \cdots & x_2^{n-2} \\[1.2em]
            \vdots             & \vdots & \vdots & \vdots & \ddots & \vdots    \\[1.2em]
            \dfrac{1}{x_n - 1} & 1      & x_n    & x_n^2  & \cdots & x_n^{n-2}
        \end{vmatrix}
    \]
\end{exercise}
\noindent 来源：\href{https://github.com/yhwu-is/Linear-Algebra-Left-Undone/blob/new/%E8%AE%B2%E4%B9%89/%E4%B8%93%E9%A2%98/10%20%E8%A1%8C%E5%88%97%E5%BC%8F.tex}{L10-例-26}。
\begin{solve}
对第 $i$ 行提出因子 $\dfrac{1}{x_i-1}$，则有

    \[
        |A| = \left(\displaystyle\prod_{i=1}^n \frac{1}{x_i-1} \right) \begin{vmatrix}
            1 & x_1 - 1 & x_1^2 - x_1 & \cdots & x_1^{n-1} - x_1^{n-2} \\
            1 & x_2 - 1 & x_2^2 - x_2 & \cdots & x_2^{n-1} - x_2^{n-2} \\
            1 & x_3 - 1 & x_3^2 - x_3 & \cdots & x_3^{n-1} - x_3^{n-2} \\
            \vdots & \vdots & \vdots & \ddots & \vdots \\
            1 & x_n - 1 & x_n^2 - x_n & \cdots & x_n^{n-1} - x_n^{n-2}
        \end{vmatrix}.
    \]

    通过初等列变换将 $A$ 化为范德蒙行列式的形式：

    \begin{align*}
        |A| &= \left(\displaystyle\prod_{i=1}^n \frac{1}{x_i-1} \right) \begin{vmatrix}
            1 & x_1 & x_1^2 & \cdots & x_1^{n-1} \\
            1 & x_2 & x_2^2 & \cdots & x_2^{n-1} \\
            1 & x_3 & x_3^2 & \cdots & x_3^{n-1} \\
            \vdots & \vdots & \vdots & \ddots & \vdots \\
            1 & x_n & x_n^2 & \cdots & x_n^{n-1}
        \end{vmatrix} \\
        &= \left(\displaystyle\prod_{i=1}^n \frac{1}{x_i-1} \right)
           \left(\displaystyle\prod_{1 \leqslant i < j \leqslant n} (x_j - x_i)\right).
    \end{align*}
\end{solve}

\begin{exercise}[补充训练 L10-例-29]\label{prob:L10-例-29}
计算行列式$D_n=\begin{vmatrix}
            1      & 2      & 3      & \cdots & n-1    & n      \\
            2      & 1      & 2      & \cdots & n-2    & n-1    \\
            3      & 2      & 1      & \cdots & n-3    & n-2    \\
            \vdots & \vdots & \vdots & \ddots & \vdots & \vdots \\
            n-1    & n-2    & n-3    & \cdots & 1      & 2      \\
            n      & n-1    & n-2    & \cdots & 2      & 1
        \end{vmatrix},\enspace n \geqslant 2$.
\end{exercise}
\noindent 来源：\href{https://github.com/yhwu-is/Linear-Algebra-Left-Undone/blob/new/%E8%AE%B2%E4%B9%89/%E4%B8%93%E9%A2%98/10%20%E8%A1%8C%E5%88%97%E5%BC%8F.tex}{L10-例-29}。
\begin{solve}
这是一个循环行列式，这种行列式的特点是相邻的行或列之间有共同的变化规律，所以我们可以使用``差分''的思想，从最后一行开始每行减去上一行
    \begin{align*}
        D_n & =\begin{vmatrix}
                   1      & 2      & 3      & \cdots & n-1    & n      \\
                   1      & -1     & -1     & \cdots & -1     & -1     \\
                   1      & 1      & -1     & \cdots & -1     & -1     \\
                   \vdots & \vdots & \vdots & \ddots & \vdots & \vdots \\
                   1      & 1      & 1      & \cdots & -1     & -1     \\
                   1      & 1      & 1      & \cdots & 1      & -1
               \end{vmatrix}=\begin{vmatrix}
                                 1      & 2      & 3      & \cdots & n-1    & n      \\
                                 2      & 0      & 0      & \cdots & 0      & -2     \\
                                 2      & 2      & 0      & \cdots & 0      & -2     \\
                                 \vdots & \vdots & \vdots & \ddots & \vdots & \vdots \\
                                 2      & 2      & 2      & \cdots & 0      & -2     \\
                                 1      & 1      & 1      & \cdots & 1      & -1
                             \end{vmatrix} \\
            & =\begin{vmatrix}
                   1      & 2      & 3      & \cdots & n-1    & n+1    \\
                   2      & 0      & 0      & \cdots & 0      & 0      \\
                   2      & 2      & 0      & \cdots & 0      & 0      \\
                   \vdots & \vdots & \vdots & \ddots & \vdots & \vdots \\
                   2      & 2      & 2      & \cdots & 0      & 0      \\
                   1      & 1      & 1      & \cdots & 1      & 0
               \end{vmatrix}=2^{n-2}
        \begin{vmatrix}
            1      & 2      & 3      & \cdots & n-1    & n+1    \\
            1      & 0      & 0      & \cdots & 0      & 0      \\
            1      & 1      & 0      & \cdots & 0      & 0      \\
            \vdots & \vdots & \vdots & \ddots & \vdots & \vdots \\
            1      & 1      & 1      & \cdots & 0      & 0      \\
            1      & 1      & 1      & \cdots & 1      & 0
        \end{vmatrix}                      \\
            & =(-1)^{n+1}(n+1) 2^{n-2}.
    \end{align*}
\end{solve}

\begin{exercise}[补充训练 L10-例-30]\label{prob:L10-例-30}
解行列式 $D=\begin{vmatrix}
            0      & 1      & 1      & \cdots & 1      & 1      \\
            1      & 0      & 1      & \cdots & 1      & 1      \\
            1      & 1      & 0      & \cdots & 1      & 1      \\
            \vdots & \vdots & \vdots & \ddots & \vdots & \vdots \\
            1      & 1      & 1      & \cdots & 0      & 1      \\
            1      & 1      & 1      & \cdots & 1      & 0
        \end{vmatrix}$.
\end{exercise}
\noindent 来源：\href{https://github.com/yhwu-is/Linear-Algebra-Left-Undone/blob/new/%E8%AE%B2%E4%B9%89/%E4%B8%93%E9%A2%98/10%20%E8%A1%8C%E5%88%97%E5%BC%8F.tex}{L10-例-30}。
\begin{solve}
使行列式 $D$ 变成 $n+1$ 阶行列式, 即
    \[ D=\begin{vmatrix}
            1      & 1      & 1      & \cdots & 1      & 1      \\
            0      & 0      & 1      & \cdots & 1      & 1      \\
            0      & 1      & 0      & \cdots & 1      & 1      \\
            \vdots & \vdots & \vdots & \ddots & \vdots & \vdots \\
            0      & 1      & 1      & \cdots & 0      & 1      \\
            0      & 1      & 1      & \cdots & 1      & 0
        \end{vmatrix}. \]

    再将第一行的 $-1$ 倍加到其他各行, 得:
    \[ D =\begin{vmatrix}
            1      & 1  & \cdots & 1  \\
            -1     & -1 &        &    \\
            \vdots &    & \ddots &    \\
            -1     &    &        & -1
        \end{vmatrix}. \]

    从第二列开始, 每列乘以 $-1$ 加到第一列, 得:
    \begin{align*}
        D & =\begin{vmatrix}
                 -(n-1) & 1  & \cdots & 1  \\
                        & -1 &        &    \\
                        &    & \ddots &    \\
                        &    &        & -1
             \end{vmatrix} \\
          & =(-1)^{n+1}(n-1).
    \end{align*}
\end{solve}

\begin{exercise}[补充训练 L10-例-31]\label{prob:L10-例-31}
求行列式 $D_n=\begin{vmatrix}
            1         & 1         & \cdots & 1         \\
            x_1       & x_2       & \cdots & x_n       \\
            x_1^{2}   & x_2^{2}   & \cdots & x_n^{2}   \\
            \vdots    & \vdots    & \ddots & \vdots    \\
            x_1^{n-2} & x_2^{n-2} & \cdots & x_n^{n-2} \\
            x_1^{n}   & x_2^{n}   & \cdots & x_n^{n}
        \end{vmatrix}$.
\end{exercise}
\noindent 来源：\href{https://github.com/yhwu-is/Linear-Algebra-Left-Undone/blob/new/%E8%AE%B2%E4%B9%89/%E4%B8%93%E9%A2%98/10%20%E8%A1%8C%E5%88%97%E5%BC%8F.tex}{L10-例-31}。
\begin{solve}
观察到这个行列式与范德蒙行列式十分类似，但是把 $n-1$ 次幂换成了 $n$ 次幂，因此我们可以考虑将其用加边法补成一个 $n+1$ 阶的范德蒙行列式. 普通的加边法套路是加一列只有一个元素非零的列，同时在非零元素所在行补上任意元素以化简. 如果我们在这道题中使用此方法，那么我们需要在倒数第二行补 $x_i^{n-1}$，这样的话就必须补上一个 $n+1$ 维列向量 $(0, 0, \ldots, 0, 1, 0)^{\mathrm{T}}$，无法构成范德蒙行列式.

    

    因此，我们必须另辟蹊径：假设我们已经构造出了一个 $n+1$ 阶的范德蒙行列式 $V_{n+1}$，加的行在倒数第二行，加的列在最后一列，那么我们其实只需要求这个行列式对加的行列的公共元素的余子式 $M_{n,n+1}$ 即可. 求余子式最直接的方法就是将 $n+1$ 阶的行列式按照加的列展开，我们已经知道了整个行列式的值，那么只需要将待求余子式 $M_{n,n+1}$ 与其他部分分离. 由于范德蒙行列式的结构特征，加的列必然是形如 $(1, x, x^2, \ldots, x^n)$ 的列向量，我们发现如果将 $x$ 视为主元，按该列展开后便形成了一个 $n$ 次的多项式，其中 $x^{n-1}$ 的系数正是我们需要的代数余子式 $(-1)^{2n+1} M_{n,n+1}$. 因此就有

    \begin{align*}
        V_{n+1} &= \begin{vmatrix}
            1      & 1      & \cdots & 1      & 1      \\
            x_1    & x_2    & \cdots & x_n    & x      \\
            x_1^2  & x_2^2  & \cdots & x_n^2  & x^2    \\
            \vdots & \vdots & \ddots & \vdots & \vdots \\
            x_1^{n-1} & x_2^{n-1} & \cdots & x_n^{n-1} & x^{n-1} \\
            x_1^n  & x_2^n  & \cdots & x_n^n  & x^n
        \end{vmatrix} = \sum\limits_{i=0}^n (-1)^{i+n+1} M_{i, n+1} x^i. \\
        V_{n+1} &= \left( \prod\limits_{1 \leqslant i < j \leqslant n} (x_j - x_i) \right) \left(\prod\limits_{i=1}^n (x - x_i) \right).
    \end{align*}

    上下两个多项式对应系数相等，故有

    \[
        - M_{n,n+1} = \left( \prod\limits_{1 \leqslant i < j \leqslant n} (x_j - x_i) \right) \sum\limits_{i=1}^n -x_i.
    \]

    因此，$D_n = M_{n, n+1} = \left( \displaystyle\prod_{1 \leqslant i < j \leqslant n} (x_j - x_i) \right) \displaystyle\sum_{i=1}^n x_i$.
\end{solve}

\begin{exercise}[补充训练 L10-例-32]\label{prob:L10-例-32}
计算行列式$D_n=\begin{vmatrix}
            1-a_{1} & a_{2}   &        &           &         \\
            -1      & 1-a_{2} & \ddots &           &         \\
                    & -1      & \ddots & a_{n-1}   &         \\
                    &         & \ddots & 1-a_{n-1} & a_{n}   \\
                    &         &        & -1        & 1-a_{n}
        \end{vmatrix}$.
\end{exercise}
\noindent 来源：\href{https://github.com/yhwu-is/Linear-Algebra-Left-Undone/blob/new/%E8%AE%B2%E4%B9%89/%E4%B8%93%E9%A2%98/10%20%E8%A1%8C%E5%88%97%E5%BC%8F.tex}{L10-例-32}。
\begin{solve}
把第一列的元素看成两项的和进行拆列, 得
    \begin{align*}
        D_n= & \begin{vmatrix}
                   1-a_{1} & a_{2}   &        &           &         \\
                   -1      & 1-a_{2} & \ddots &           &         \\
                   0 + 0   & -1      & \ddots & a_{n-1}   &         \\
                   0 + 0   &         & \ddots & 1-a_{n-1} & a_{n}   \\
                   0 + 0   &         &        & -1        & 1-a_{n}
               \end{vmatrix}                \\
        =    & \begin{vmatrix}
                   1  & a_{2}   &        &           &         \\
                   -1 & 1-a_{2} & \ddots &           &         \\
                      & -1      & \ddots & a_{n-1}   &         \\
                      &         & \ddots & 1-a_{n-1} & a_{n}   \\
                      &         &        & -1        & 1-a_{n}
               \end{vmatrix} + \begin{vmatrix}
                                   -a_{1} & a_{2}   &        &           &         \\
                                          & 1-a_{2} & \ddots &           &         \\
                                          & -1      & \ddots & a_{n-1}   &         \\
                                          &         & \ddots & 1-a_{n-1} & a_{n}   \\
                                          &         &        & -1        & 1-a_{n}
                               \end{vmatrix}.
    \end{align*}

    上面第一个行列式的值为 1（从第 1 行开始，每一行依次加到下一行），所以
    \[ \begin{aligned}
            D_n & =1-a_{1}\begin{vmatrix}
                              1-a_{2} & a_{3}   &        &           &         \\
                              -1      & 1-a_{3} & \ddots &           &         \\
                                      & -1      & \ddots & a_{n-1}   &         \\
                                      &         & \ddots & 1-a_{n-1} & a_{n}   \\
                                      &         &        & -1        & 1-a_{n}
                          \end{vmatrix} \\
                & =1-a_{1} D_{n-1} \cdot
        \end{aligned} \]

    这个式子对任何 $n \geqslant 2$ 都成立, 因此有
    \begin{align*}
        D_n & =1-a_{1} D_{n-1}                                          \\
            & =1-a_{1}(1-a_{2} D_{n-2})                                 \\
            & =\cdots                                                   \\
            & =1-a_{1}+a_{1} a_{2}+\cdots+(-1)^n a_{1} a_{2} \cdots a_n \\
            & =1+\sum_{i=1}^{n}(-1)^i \prod_{j=1}^i a_j.
    \end{align*}
\end{solve}

\begin{exercise}[补充训练 L10-例-33]\label{prob:L10-例-33}
定义
    \[
        f(x) = \begin{vmatrix}
            x - a_{11} & - a_{12} & \cdots & - a_{1n} \\
            - a_{21} & x - a_{22} & \cdots & - a_{2n} \\
            \vdots   & \vdots   & \ddots & \vdots   \\
            - a_{n1} & - a_{n2} & \cdots & x - a_{nn}
        \end{vmatrix},
    \]
    试使用硬拆法求
    \begin{enumerate}
        \item $f(x)$ 中 $n-1$ 次项系数 $c_{n-1}$ 的值.
        \item $f(x)$ 中常数项系数 $c_0$ 的值.
    \end{enumerate}
\end{exercise}
\noindent 来源：\href{https://github.com/yhwu-is/Linear-Algebra-Left-Undone/blob/new/%E8%AE%B2%E4%B9%89/%E4%B8%93%E9%A2%98/10%20%E8%A1%8C%E5%88%97%E5%BC%8F.tex}{L10-例-33}。
\begin{solve}
\begin{enumerate}
        \item 设列向量组 $\alpha_i = (-a_{1i}, -a_{2i}, \ldots, -a_{ni})^{\mathrm{T}}$，将行列式的第 $i$ 列分解为 $xe_i + \alpha_i$，即

            \[
                f(x) = D(xe_1 + \alpha_1, xe_2 + \alpha_2, \ldots, xe_n + \alpha_n).
            \]

            逐列进行分解，我们可以得到 $2^n$ 个 $n$ 阶行列式. 设某个行列式中有 $p$ 列为 $xe_i$，$n-p$ 列为 $\alpha_i$，对所有取 $xe_i$ 的列展开，就可以得到一个形如 $cx^p$ 的多项式. 因此，我们只需要考虑拆出的 $2^n$ 个行列式中有 $n-1$ 列为 $xe_i$ 的行列式即可. 因此有

            \[
                c_{n-1}x^{n-1} = D(\alpha_1, xe_2, \ldots, xe_n) + D(xe_1, \alpha_2, \ldots, xe_n) + \cdots + D(xe_1, xe_2, \ldots, \alpha_n).
            \]

            化简后便有 $c_{n-1} = -\displaystyle\sum_{i=1}^n a_{ii}$.

        \item 使用硬拆法，拆出的 $2^n$ 个行列式中只有 $D(\alpha_1, \alpha_2, \ldots, \alpha_n)$ 是常数，因此 $c_0 = D(\alpha_1, \alpha_2, \ldots, \alpha_n)$.
    \end{enumerate}
\end{solve}

\begin{exercise}[补充训练 L10-例-34]\label{prob:L10-例-34}
解行列式$D_n=\begin{vmatrix}
            x     & -1    &        &         &         \\
                  & x     & \ddots &         &         \\
                  &       & \ddots & -1      &         \\
                  &       &        & x       & -1      \\
            a_{0} & a_{1} & \cdots & a_{n-2} & a_{n-1}
        \end{vmatrix}$.
\end{exercise}
\noindent 来源：\href{https://github.com/yhwu-is/Linear-Algebra-Left-Undone/blob/new/%E8%AE%B2%E4%B9%89/%E4%B8%93%E9%A2%98/10%20%E8%A1%8C%E5%88%97%E5%BC%8F.tex}{L10-例-34}。
\begin{solve}
按最后一行展开，得
    \begin{align*}
        D_n & =\sum_{i=0}^{n-1}(-1)^{n+i+1}a_i\begin{vmatrix} A & O \\ O & B \end{vmatrix}
        =\sum_{i=0}^{n-1}(-1)^{n+i+1}a_i|A||B|                                                        \\
            & =\sum_{i=0}^{n-1}(-1)^{n+i+1}a_i(x^i)\left((-1)^{n-1-i}\right) =\sum_{i=0}^{n-1}a_i x^i,
    \end{align*}

    其中$A\in \mathbf{M}_i(\mathbf{R}),\enspace B\in \mathbf{M}_{n-1-i}(\mathbf{R})$，
    \[ A=\begin{pmatrix}
            x & -1 &        &    &    \\
              & x  & \ddots &    &    \\
              &    & \ddots & -1 &    \\
              &    &        & x  & -1 \\
              &    &        &    & x
        \end{pmatrix},\enspace
        B=\begin{pmatrix}
            -1 &    &        &    &    \\
            x  & -1 &        &    &    \\
               & x  & \ddots &    &    \\
               &    & \ddots & -1 &    \\
               &    &        & x  & -1
        \end{pmatrix}. \]
\end{solve}

\begin{exercise}[补充训练 L10-例-36]\label{prob:L10-例-36}
计算行列式$\begin{vmatrix}
            -1 & 1  & 1 & 2 & -1 \\
            0  & -1 & 0 & 1 & 2  \\
            2  & 1  & 1 & 3 & -1 \\
            1  & 2  & 2 & 1 & 0  \\
            0  & 3  & 0 & 1 & 3
        \end{vmatrix}$.
\end{exercise}
\noindent 来源：\href{https://github.com/yhwu-is/Linear-Algebra-Left-Undone/blob/new/%E8%AE%B2%E4%B9%89/%E4%B8%93%E9%A2%98/10%20%E8%A1%8C%E5%88%97%E5%BC%8F.tex}{L10-例-36}。
\begin{solve}
因为第一、第三列含有较多的 $0$，因此在这两列上做 Laplace 展开，得
    \begin{align*}
        |A| & = \begin{vmatrix}
                    -1 & 1 \\ 2 & 1
                \end{vmatrix} \cdot (-1)^{1+3+1+3}
        \begin{vmatrix}
            -1 & 1 & 2 \\ 2 & 1 & 0 \\ 3 & 1 & 3
        \end{vmatrix}                                 \\
            & + \begin{vmatrix}
                    -1 & 1 \\ 1 & 2
                \end{vmatrix} \cdot (-1)^{1+4+1+3}
        \begin{vmatrix}
            -1 & 1 & 2 \\ 1 & 3 & -1 \\ 3 & 1 & 3
        \end{vmatrix}                                \\
            & + \begin{vmatrix}
                    2 & 1 \\ 1 & 1
                \end{vmatrix} \cdot (-1)^{3+4+1+3}
        \begin{vmatrix}
            1 & 2 & -1 \\ -1 & 1 & 2 \\ 3 & 1 & 3
        \end{vmatrix}                                \\
            & = (-3) \times (-11) + (-3) \times 32 + 3 \times (-23) = -132.
    \end{align*}
\end{solve}

\begin{exercise}[补充训练 L10-例-37]\label{prob:L10-例-37}
计算行列式$D_n=\begin{vmatrix}
            \cos \beta & 1            &        &              &              \\
            1          & 2 \cos \beta & \ddots &              &              \\
                       & 1            & \ddots & 1            &              \\
                       &              & \ddots & 2 \cos \beta & 1            \\
                       &              &        & 1            & 2 \cos \beta
        \end{vmatrix}$.
\end{exercise}
\noindent 来源：\href{https://github.com/yhwu-is/Linear-Algebra-Left-Undone/blob/new/%E8%AE%B2%E4%B9%89/%E4%B8%93%E9%A2%98/10%20%E8%A1%8C%E5%88%97%E5%BC%8F.tex}{L10-例-37}。
\begin{solve}
\begin{enumerate}
        \item 法一：归纳法
            \begin{align*}
                D_1 & =\cos\beta,                               \\
                D_2 & =\begin{vmatrix}
                            \cos\beta & 1          \\
                            1         & 2\cos\beta
                        \end{vmatrix}=2\cos^2\beta-1=\cos2\beta.
            \end{align*}

            猜想$D_n=\cos n\beta$. 数学归纳证明：

            假设当 $n=k$ 时，结论成立，即 $D_{k}=\cos k \beta$. 现证当 $n=k+1$ 时，结论也成立. $ n=k+1 $ 时，
            \[ D_{k+1} = \begin{vmatrix}
                    \cos \beta & 1            &        &              &              \\
                    1          & 2 \cos \beta & \ddots &              &              \\
                                & 1            & \ddots & 1            &              \\
                                &              & \ddots & 2 \cos \beta & 1            \\
                                &              &        & 1            & 2 \cos \beta
                \end{vmatrix}. \]

            将 $D_{k+1}$ 按最后一行展开, 得
            \begin{align*}
                D_{k+1}={} & (-1)^{k+1+k+1} \cdot 2 \cos \beta
                \begin{vmatrix}
                    \cos \beta & 1            &        &              &              \\
                    1          & 2 \cos \beta & \ddots &              &              \\
                                & 1            & \ddots & 1            &              \\
                                &              & \ddots & 2 \cos \beta & 1            \\
                                &              &        & 1            & 2 \cos \beta
                \end{vmatrix} \\
                            & +(-1)^{k+1+k}
                \begin{vmatrix}
                    \cos \beta & 1            & 0            & \cdots & 0      \\
                    1          & 2 \cos \beta & 1            & \cdots & 0      \\
                    0          & 1            & 2 \cos \beta & \cdots & 0      \\
                    \vdots     & \vdots       & \vdots       & \ddots & \vdots \\
                    0          & 0            & 0            & \cdots & 1
                \end{vmatrix}       \\
                ={}        & 2\cos\beta D_k-D_{k-1}.
            \end{align*}

            而$D_{k}=\cos k \beta,\enspace
                D_{k-1}=\cos (k-1)\beta = \cos (k \beta-\beta) = \cos k \beta \cos \beta+\sin k \beta \sin \beta$. 所以有
            \begin{align*}
                D_{k+1} & = 2 \cos \beta D_{k}-D_{k-1}                                               \\
                        & =2 \cos \beta \cos k \beta-\cos k \beta \cos \beta-\sin k \beta \sin \beta \\
                        & =\cos k \beta \cos \beta-\sin k \beta \sin \beta                           \\
                        & =\cos (k+1) \beta.
            \end{align*}

            则证得$D_n=\cos n\beta,\enspace n\in \mathbf{N}$.

        \item 法二：递推法
            先得到 $D_n$ 的递推式：
            \[
                D_n - 2 \cos \beta D_{n-1} + D_{n-2} = 0,
            \]
            其中 $D_1 = \cos \beta, \enspace D_2 = \cos 2\beta$.

            现求其特征方程
            \[
                \lambda^2 - 2 \lambda \cos \beta + 1 = 0
            \]
            的根，有
            \[
                \lambda = 2 \cos \beta \pm 2 \mathrm{i} \sin \beta.
            \]
            故令
            \[
                D_n = C_1 \cos n\beta + C_2 \sin n\beta,
            \]
            代入初值，有
            \[
                \begin{cases}
                    C_1 \cos \beta + C_2 \sin \beta = \cos \beta, \\
                    C_1 \cos 2\beta + C_2 \sin 2\beta = \cos 2\beta.
                \end{cases}
            \]
            易知 $C_1 = 1, C_2 = 0$ 为其解，故 $D_n = \cos n\beta$.
    \end{enumerate}
\end{solve}

\begin{exercise}[补充训练 L10-例-38]\label{prob:L10-例-38}
计算行列式$D_n=
        \begin{vmatrix}
            9 & 5 &        &   &   \\
            4 & 9 & \ddots &   &   \\
              & 4 & \ddots & 5 &   \\
              &   & \ddots & 9 & 5 \\
              &   &        & 4 & 9
        \end{vmatrix}$.
\end{exercise}
\noindent 来源：\href{https://github.com/yhwu-is/Linear-Algebra-Left-Undone/blob/new/%E8%AE%B2%E4%B9%89/%E4%B8%93%E9%A2%98/10%20%E8%A1%8C%E5%88%97%E5%BC%8F.tex}{L10-例-38}。
\begin{solve}
按第一列展开, 得
    \[ D_n=9 D_{n-1}-20 D_{n-2}, \]
    即 $ D_n-9 D_{n-1}+20 D_{n-2}=0 $.
    作特征方程
    \[ x^{2}-9 x+20=0, \]
    解得 $ x_1=4,\enspace x_2=5 $. 则
    \[ D_n=A \cdot 4^n+B \cdot 5^n, \]

    当 $n=1$ 时, $9=4A+5B$；

    当 $n=2$ 时，$61=16A+25B$.

    解得$A=-4,\enspace B=5$，所以
    \[ D_n=5^{n+1}-4^{n+1}. \]
\end{solve}

\begin{exercise}[补充训练 L10-例-39]\label{prob:L10-例-39}
计算行列式 $D_n=\begin{vmatrix}
            (a_1+b_1)^{-1} & (a_1+b_2)^{-1} & \cdots & (a_1+b_n)^{-1} \\
            (a_2+b_1)^{-1} & (a_2+b_2)^{-1} & \cdots & (a_2+b_n)^{-1} \\
            \vdots         & \vdots         & \ddots & \vdots         \\
            (a_n+b_1)^{-1} & (a_n+b_2)^{-1} & \cdots & (a_n+b_n)^{-1}
        \end{vmatrix}$.
\end{exercise}
\noindent 来源：\href{https://github.com/yhwu-is/Linear-Algebra-Left-Undone/blob/new/%E8%AE%B2%E4%B9%89/%E4%B8%93%E9%A2%98/10%20%E8%A1%8C%E5%88%97%E5%BC%8F.tex}{L10-例-39}。
\begin{solve}
我们尝试求解 $D_n$ 和 $D_{n-1}$ 之间的递推式. 首先将行列式的前 $n-1$ 列每列都减去第 $n$ 列：

    \begin{align*}
        D_n & = \begin{vmatrix}
                    \dfrac{1}{a_1+b_1}     & \cdots & \dfrac{1}{a_1+b_{n-1}}     & \dfrac{1}{a_1+b_n}     \\[2ex]
                    \vdots                 & \ddots & \vdots                     & \vdots                 \\[2ex]
                    \dfrac{1}{a_{n-1}+b_1} & \cdots & \dfrac{1}{a_{n-1}+b_{n-1}} & \dfrac{1}{a_{n-1}+b_n} \\[2ex]
                    \dfrac{1}{a_n+b_1}     & \cdots & \dfrac{1}{a_n+b_{n-1}}     & \dfrac{1}{a_n+b_n}
                \end{vmatrix}                                                                       \\[2ex]
            & = \begin{vmatrix}
                    \dfrac{b_n-b_1}{(a_1+b_1)(a_1+b_n)}         & \cdots & \dfrac{b_n-b_{n-1}}{(a_1+b_{n-1})(a_1+b_n)}         & \dfrac{1}{a_1+b_n}     \\[2ex]
                    \vdots                                      & \ddots & \vdots                                              & \vdots                 \\[2ex]
                    \dfrac{b_n-b_1}{(a_{n-1}+b_1)(a_{n-1}+b_n)} & \cdots & \dfrac{b_n-b_{n-1}}{(a_{n-1}+b_{n-1})(a_{n-1}+b_n)} & \dfrac{1}{a_{n-1}+b_n} \\[2ex]
                    \dfrac{b_n-b_1}{(a_n+b_1)(a_n+b_n)}         & \cdots & \dfrac{b_n-b_{n-1}}{(a_n+b_{n-1})(a_n+b_n)}         & \dfrac{1}{a_n+b_n}
                \end{vmatrix}.
    \end{align*}
    第二步是提取公因式：
    \[ \text{上式} = \dfrac{\prod\limits_{i=1}^{n-1}(b_n-b_i)}{\prod\limits_{j=1}^{n}(a_j+b_n)}
        \begin{vmatrix}
            \dfrac{1}{a_1+b_1}     & \cdots & \dfrac{1}{a_1+b_{n-1}}     & 1      \\[2ex]
            \vdots                 & \ddots & \vdots                     & \vdots \\[2ex]
            \dfrac{1}{a_{n-1}+b_1} & \cdots & \dfrac{1}{a_{n-1}+b_{n-1}} & 1      \\[2ex]
            \dfrac{1}{a_n+b_1}     & \cdots & \dfrac{1}{a_n+b_{n-1}}     & 1
        \end{vmatrix}. \]
    第三步是将行列式的前 $n-1$ 行每行都减去第 $n$ 行：
    \[ \text{上式} = \dfrac{\prod\limits_{i=1}^{n-1}(b_n-b_i)}{\prod\limits_{j=1}^{n}(a_j+b_n)}
        \begin{vmatrix}
            \dfrac{a_n-a_1}{(a_1+b_1)(a_n+b_1)}         & \cdots & \dfrac{a_n-a_1}{(a_1+b_{n-1})(a_n+b_{n-1})}         & 0      \\[2ex]
            \vdots                                      & \ddots & \vdots                                              & \vdots \\[2ex]
            \dfrac{a_n-a_{n-1}}{(a_{n-1}+b_1)(a_n+b_1)} & \cdots & \dfrac{a_n-a_{n-1}}{(a_{n-1}+b_{n-1})(a_n+b_{n-1})} & 0      \\[2ex]
            \dfrac{1}{a_n+b_1}                          & \cdots & \dfrac{1}{a_n+b_{n-1}}                              & 1
        \end{vmatrix}. \]
    然后再次提取公因式：
    \begin{align*}
        \text{上式} & = \dfrac{\prod\limits_{i=1}^{n-1}(a_n-a_i)(b_n-b_i)}{\prod\limits_{j=1}^{n}(a_j+b_n)\prod\limits_{k=1}^{n-1}(a_n+b_k)}
        \begin{vmatrix}
            \dfrac{1}{a_1+b_1}     & \cdots & \dfrac{1}{a_1+b_{n-1}}     \\[2ex]
            \vdots                 & \ddots & \vdots                     \\[2ex]
            \dfrac{1}{a_{n-1}+b_1} & \cdots & \dfrac{1}{a_{n-1}+b_{n-1}}
        \end{vmatrix}                                                         \\[2ex]
                    & = \dfrac{\prod\limits_{i=1}^{n-1}(a_n-a_i)(b_n-b_i)}{\prod\limits_{j=1}^{n}(a_j+b_n)\prod\limits_{k=1}^{n-1}(a_n+b_k)} D_{n-1}.
    \end{align*}
    不断递推下去即得

    \[|A| = \dfrac{\prod\limits_{1 \leqslant i < j \leqslant n}(a_j - a_i)(b_j - b_i)}{\prod\limits_{i,j=1}^{n}(a_i + b_j)}.\]
\end{solve}

\begin{exercise}[补充训练 L10-例-40]\label{prob:L10-例-40}
求下列矩阵行列式的值：
            \[A=\begin{pmatrix}
                    a_1^2    & 1+a_1a_2   & \cdots & 1+a_1a_n \\
                    1+a_2a_1 & a_2^2    & \cdots & 1+a_2a_n \\
                    \vdots   & \vdots   & \ddots & \vdots   \\
                    1+a_na_1 & 1+a_na_2 & \cdots & a_n^2
                \end{pmatrix}.\]
\end{exercise}
\noindent 来源：\href{https://github.com/yhwu-is/Linear-Algebra-Left-Undone/blob/new/%E8%AE%B2%E4%B9%89/%E4%B8%93%E9%A2%98/10%20%E8%A1%8C%E5%88%97%E5%BC%8F.tex}{L10-例-40}。
\begin{solve}
我们可以将 $A$ 化为

            \[ A = -E_n + \begin{pmatrix}
                    a_1 & 1 \\ a_2 & 1 \\ \vdots & \vdots \\ a_n & 1
                \end{pmatrix} E_2^{-1} \begin{pmatrix}
                    a_1 & a_2 & \cdots & a_n \\ 1 & 1 & \cdots & 1
                \end{pmatrix}. \]

            故根据降阶公式有

            \begin{align*}
                |A| & = \left\vert-E_n +
                \begin{pmatrix}
                    a_1 & 1 \\ a_2 & 1 \\ \vdots & \vdots \\ a_n & 1
                \end{pmatrix} E_2^{-1}
                \begin{pmatrix}
                    a_1 & a_2 & \cdots & a_n \\ 1 & 1 & \cdots & 1
                \end{pmatrix}\right\vert                                                                       \\
                    & = (-1)^n \left|E_2 -
                \begin{pmatrix}
                    \sum\limits_{i=1}^n a_i^2 & \sum\limits_{i=1}^n a_i \\[2ex]
                    \sum\limits_{i=1}^n a_i   & n
                \end{pmatrix}\right|                                                          \\
                    & = (-1)^n \left((1-n)\left(1-\sum\limits_{i=1}^n a_i^2\right) - \left(\sum\limits_{i=1}^n a_i\right)^2\right).
            \end{align*}
\end{solve}

\begin{exercise}[补充训练 L10-例-41]\label{prob:L10-例-41}
求行列式 $\begin{vmatrix}
                    0      & 2a_1   & 3a_1   & \cdots & na_1     \\
                    a_2    & a_2    & 3a_2   & \cdots & na_2     \\
                    a_3    & 2a_3   & 2a_3   & \cdots & na_3     \\
                    \vdots & \vdots & \vdots & \ddots & \vdots   \\
                    a_n    & 2a_n   & 3a_n   & \cdots & (n-1)a_n \\
                \end{vmatrix}$.
\end{exercise}
\noindent 来源：\href{https://github.com/yhwu-is/Linear-Algebra-Left-Undone/blob/new/%E8%AE%B2%E4%B9%89/%E4%B8%93%E9%A2%98/10%20%E8%A1%8C%E5%88%97%E5%BC%8F.tex}{L10-例-41}。
\begin{solve}
\[ \text{原式}=\prod_{i=1}^na_i \begin{vmatrix}
                0      & 2      & 3      & \cdots & n      \\
                1      & 1      & 3      & \cdots & n      \\
                1      & 2      & 2      & \cdots & n      \\
                \vdots & \vdots & \vdots & \ddots & \vdots \\
                1      & 2      & 3      & \cdots & n-1    \\
            \end{vmatrix}. \]

            注意到这个行列式对应的矩阵可以加上一个单位矩阵构成一个每行都相等的矩阵，所以我们可以利用$|E_m-AB|=|E_n-BA|$来求解这个行列式.
            \begin{align*}
                \begin{vmatrix}
                    0      & 2      & 3      & \cdots & n      \\
                    1      & 1      & 3      & \cdots & n      \\
                    1      & 2      & 2      & \cdots & n      \\
                    \vdots & \vdots & \vdots & \ddots & \vdots \\
                    1      & 2      & 3      & \cdots & n-1
                \end{vmatrix}
                & = \begin{vmatrix}(-1)\left(E_n-\begin{pmatrix}
                        1      & 2      & 3      & \cdots & n      \\
                        1      & 2      & 3      & \cdots & n      \\
                        1      & 2      & 3      & \cdots & n      \\
                        \vdots & \vdots & \vdots & \ddots & \vdots \\
                        1      & 2      & 3      & \cdots & n
                    \end{pmatrix}\right)\end{vmatrix} \\
                & =(-1)^n\begin{vmatrix}E_n-
                            \begin{pmatrix}
                        1 \\1\\1\\\vdots\\1
                    \end{pmatrix}\begin{pmatrix}1 & 2 & 3 & \cdots & n\end{pmatrix}\end{vmatrix}.
            \end{align*}

            而 \[ \begin{vmatrix}E_n-\begin{pmatrix}
                        1 \\1\\1\\\vdots\\1
                    \end{pmatrix}\begin{pmatrix}1 & 2 & 3 & \cdots & n\end{pmatrix}\end{vmatrix}
                =1-\begin{pmatrix}1 & 2 & 3 & \cdots & n\end{pmatrix}
                \begin{pmatrix}1 \\ 1 \\ 1 \\ \vdots \\ 1\end{pmatrix}
                =-\frac{n^2+n-2}{2}. \]

            所以原式$\displaystyle =(-1)^{n+1} \cdot \frac{n^2+n-2}{2} \left(\prod_{i=1}^na_i \right)$.
\end{solve}

\begin{exercise}[补充训练 L10-例-42]\label{prob:L10-例-42}
解行列式$D_n=\begin{vmatrix}
                    \lambda & a      & a      & a      & \cdots & a      \\
                    b       & \gamma & \beta  & \beta  & \cdots & \beta  \\
                    b       & \beta  & \gamma & \beta  & \cdots & \beta  \\
                    \vdots  & \vdots & \vdots & \vdots & \ddots & \vdots \\
                    b       & \beta  & \beta  & \beta  & \cdots & \gamma
                \end{vmatrix}$.
\end{exercise}
\noindent 来源：\href{https://github.com/yhwu-is/Linear-Algebra-Left-Undone/blob/new/%E8%AE%B2%E4%B9%89/%E4%B8%93%E9%A2%98/10%20%E8%A1%8C%E5%88%97%E5%BC%8F.tex}{L10-例-42}。
\begin{solve}
从第$n$行到第3行，每行都减去上一行；再从第3列到第$n$列，每列都加到第2列，得

            \begin{align*}
                D_n & =
                \begin{vmatrix}
                    \lambda & a            & a            & a            & \cdots & a            \\
                    b       & \gamma       & \beta        & \beta        & \cdots & \beta        \\
                    0       & \beta-\gamma & \gamma-\beta & 0            & \cdots & 0            \\
                    0       & 0            & \beta-\gamma & \gamma-\beta & \cdots & 0            \\
                    \vdots  & \vdots       & \vdots       & \vdots       & \ddots & \vdots       \\
                    0       & 0            & 0            & 0            & \cdots & \gamma-\beta
                \end{vmatrix}             \\
                    & =\begin{vmatrix}
                           \lambda & (n-1)a            & a            & a            & \cdots & a            \\
                           b       & \gamma+(n-2)\beta & \beta        & \beta        & \cdots & \beta        \\
                           0       & 0                 & \gamma-\beta & 0            & \cdots & 0            \\
                           0       & 0                 & \beta-\gamma & \gamma-\beta & \cdots & 0            \\
                           \vdots  & \vdots            & \vdots       & \vdots       & \ddots & \vdots       \\
                           0       & 0                 & 0            & 0            & \cdots & \gamma-\beta
                       \end{vmatrix} \\
                    & =\begin{vmatrix}
                           \lambda & (n-1)a            \\
                           b       & \gamma+(n-2)\beta
                       \end{vmatrix} \cdot \begin{vmatrix}
                                               \gamma-\beta & 0            & \cdots & 0            \\
                                               \beta-\gamma & \gamma-\beta & \cdots & 0            \\
                                               \vdots       & \vdots       & \ddots & \vdots       \\
                                               0            & 0            & \cdots & \gamma-\beta
                                           \end{vmatrix}           \\
                    & =(\lambda \gamma+\lambda(n-2)\beta-(n-1)ab)(\gamma-\beta)^{n-2}.
            \end{align*}
\end{solve}

\begin{exercise}[补充训练 L10-例-43]\label{prob:L10-例-43}
证明 Lagrange 恒等式 $(n \geqslant 2)$：
    \[ \left( \sum_{i=1}^{n}\limits a_i^2 \right) \left( \sum_{i=1}^{n}\limits b_i^2 \right) - \left( \sum_{i=1}^{n}\limits a_i b_i \right)^2 = \sum_{1 \leqslant i < j \leqslant n} (a_i b_j - a_j b_i)^2. \]
\end{exercise}
\noindent 来源：\href{https://github.com/yhwu-is/Linear-Algebra-Left-Undone/blob/new/%E8%AE%B2%E4%B9%89/%E4%B8%93%E9%A2%98/10%20%E8%A1%8C%E5%88%97%E5%BC%8F.tex}{L10-例-43}。
\begin{solve}
左边的式子等于
    \[ \begin{vmatrix}
        \sum_{i=1}^{n}\limits a_i^2 & \sum_{i=1}^{n}\limits a_i b_i \\
        \sum_{i=1}^{n}\limits a_i b_i & \sum_{i=1}^{n}\limits b_i^2
    \end{vmatrix}, \]

    这个行列式对应的矩阵可化为：
    \[ \begin{pmatrix}
        a_1 & a_2 & \cdots & a_n \\
        b_1 & b_2 & \cdots & b_n
    \end{pmatrix}
    \begin{pmatrix}
        a_1 & b_1 \\
        a_2 & b_2 \\
        \vdots & \vdots \\
        a_n & b_n
    \end{pmatrix}. \]
    用 Cauchy-Binet 公式得

    \[
    \begin{vmatrix}
        \sum_{i=1}^{n}\limits a_i^2 & \sum_{i=1}^{n}\limits a_i b_i \\
        \sum_{i=1}^{n}\limits a_i b_i & \sum_{i=1}^{n}\limits b_i^2
    \end{vmatrix}
    = \sum_{1 \leqslant i < j \leqslant n}
    \begin{vmatrix}
        a_i & a_j \\
        b_i & b_j
    \end{vmatrix}^2
    = \sum_{1 \leqslant i < j \leqslant n} (a_i b_j - a_j b_i)^2.
    \]
\end{solve}

\section{补充训练：线性方程组与矩阵综合}

本节包含 46 道训练题，建议在完成相关专题真题后选择性练习。

\begin{exercise}[补充训练 L11-A-2]\label{prob:L11-A-2}
设$A$为四阶矩阵，$r(A)<4$，且$A_{21}\neq 0$，求方程组$AX=\vec{0}$的通解.
\end{exercise}
\noindent 来源：\href{https://github.com/yhwu-is/Linear-Algebra-Left-Undone/blob/new/%E8%AE%B2%E4%B9%89/%E4%B8%93%E9%A2%98/11%20%E6%9C%9D%E8%8A%B1%E5%A4%95%E6%8B%BE.tex}{L11-A-2}。
\begin{solve}
$r(\operatorname{adj}(A))=
              \begin{cases}
                  1 & r(A)=3 \\
                  0 & r(A)<3
              \end{cases}$，而已知$A_{21}\neq 0$，故$r(\operatorname{adj}(A))=1$，故$r(A)=3$，故$AX = 0$解空间维数为$4-r(A)=1$. 而$|A|=0$，由行列式展开可知
          \[ a_{i1}A_{21}+a_{i2}A_{22}+a_{i3}A_{23}+a_{i4}A_{24}=0 \qquad i=1,2,3,4 \]
          故解为$k(A_{21},A_{22},A_{23},A_{24})^\mathrm{T},\enspace k\in \mathbf{R}$.
\end{solve}

\begin{exercise}[补充训练 L11-A-3]\label{prob:L11-A-3}
设$A=(\alpha_1,\alpha_2,\alpha_3,\alpha_4)$为四阶矩阵，方程组$AX=\vec{0}$的通解为$X=k(1,0,-4,0)^\mathrm{T}$，求$\operatorname{adj}(A)X=0$的基础解系.
\end{exercise}
\noindent 来源：\href{https://github.com/yhwu-is/Linear-Algebra-Left-Undone/blob/new/%E8%AE%B2%E4%B9%89/%E4%B8%93%E9%A2%98/11%20%E6%9C%9D%E8%8A%B1%E5%A4%95%E6%8B%BE.tex}{L11-A-3}。
\begin{solve}
由题意得$r(A)=3<4$且$\alpha_1-4\alpha_3=0$，或$\alpha_1=4\alpha_3$. 由$r(A)=3$得$r(\operatorname{adj}(A))=1$，$\operatorname{adj}(A)X=0$的基础解系含3个线性无关的解向量. 由$\operatorname{adj}(A)A=|A|E=O$得$\alpha_1,\alpha_2,\alpha_3,\alpha_4$为$\operatorname{adj}(A)X=0$的解，从而$\alpha_1,\alpha_2,\alpha_4$或$\alpha_2,\alpha_3,\alpha_4$为$\operatorname{adj}(A)X=0$的基础解系.
\end{solve}

\begin{exercise}[补充训练 L11-A-4]\label{prob:L11-A-4}
设$A$为$n$阶实矩阵，$W=\{\beta\in\mathbf{R}^n \mid \alpha^\mathrm{T}A\beta=0,\enspace \forall \alpha\in\mathbf{R}^n\}$，证明：
        \begin{enumerate}
            \item $\dim W+r(A)=n$；

            \item $W$为$\mathbf{R}^n$的子空间.
        \end{enumerate}
\end{exercise}
\noindent 来源：\href{https://github.com/yhwu-is/Linear-Algebra-Left-Undone/blob/new/%E8%AE%B2%E4%B9%89/%E4%B8%93%E9%A2%98/11%20%E6%9C%9D%E8%8A%B1%E5%A4%95%E6%8B%BE.tex}{L11-A-4}。
\begin{solve}
$\forall \alpha \in \mathbf{R}^n, \alpha^\mathrm{T}A\beta =0\implies A\beta =0$.
\end{solve}

\begin{exercise}[补充训练 L11-A-5]\label{prob:L11-A-5}
已知4级方阵$A=(\alpha_1,\alpha_2,\alpha_3,\alpha_4)$的列向量$\alpha_1,\alpha_2,\alpha_4$线性无关，且$\alpha_1=2\alpha_2-\alpha_3$，若$\beta=\alpha_1-\alpha_2+3\alpha_4$，求方程组$AX=\beta$的通解.
\end{exercise}
\noindent 来源：\href{https://github.com/yhwu-is/Linear-Algebra-Left-Undone/blob/new/%E8%AE%B2%E4%B9%89/%E4%B8%93%E9%A2%98/11%20%E6%9C%9D%E8%8A%B1%E5%A4%95%E6%8B%BE.tex}{L11-A-5}。
\begin{solve}
$r(A)=3\implies \dim{N(A)}=4-r(A)=1$.

          $\beta = x_1\alpha_1+x_2\alpha_2+x_3\alpha_3+x_4\alpha_4\implies$特解为$(1,-1,0,3)^\mathrm{T}$

          $\alpha_1-2\alpha_2+\alpha_3=0\implies$导出组基础解系为$(1,-2,1,0)^\mathrm{T}$.

          故通解为$k(1,-2,1,0)^\mathrm{T}+(1,-1,0,3)^\mathrm{T},\enspace k\in\mathbf{R}$
\end{solve}

\begin{exercise}[补充训练 L11-A-6]\label{prob:L11-A-6}
设四元非齐次线性方程组的系数矩阵的秩为3，已知$\eta_1,\eta_2,\eta_3$是它的三个解向量，且$\eta_1=(2,3,4,5)^\mathrm{T},\eta_2+\eta_3=(1,2,3,4)^\mathrm{T}$，求该方程组的通解.
\end{exercise}
\noindent 来源：\href{https://github.com/yhwu-is/Linear-Algebra-Left-Undone/blob/new/%E8%AE%B2%E4%B9%89/%E4%B8%93%E9%A2%98/11%20%E6%9C%9D%E8%8A%B1%E5%A4%95%E6%8B%BE.tex}{L11-A-6}。
\begin{solve}
设方程组$Ax=b$，由$r(A)=3$知其导出组$Ax=0$的基础解系只含一个解向量.

          而$A\eta_1=b,A(\eta_2+\eta_3)=2b$，故$2\eta_1-(\eta_2+\eta_3)=(3,4,5,6)^\mathrm{T}$为$Ax=0$的基础解系，从而所求通解为
          \[ x = \eta_1+c(3,4,5,6)^\mathrm{T} = (2,3,4,5)^\mathrm{T} + c(3,4,5,6)^\mathrm{T} \]
          其中$c$为任意常数.
\end{solve}

\begin{exercise}[补充训练 L11-A-7]\label{prob:L11-A-7}
设$\beta_1,\beta_2,\beta_3$是$n$元非齐次线性方程组$AX=\vec{b}$的三个线性无关的解，且$r(A)=n-2$，求：
        \begin{enumerate}
            \item 导出组$AX=\vec{0}$的一个基础解系；

            \item $AX=2\vec{b}$的一般解.
        \end{enumerate}
\end{exercise}
\noindent 来源：\href{https://github.com/yhwu-is/Linear-Algebra-Left-Undone/blob/new/%E8%AE%B2%E4%B9%89/%E4%B8%93%E9%A2%98/11%20%E6%9C%9D%E8%8A%B1%E5%A4%95%E6%8B%BE.tex}{L11-A-7}。
\begin{solve}
\begin{enumerate}
                \item 取$\beta_1-\beta_2,\beta_1-\beta_3$（取法不唯一），有$A(\beta_1-\beta_2)=A(\beta_1-\beta_3)=0$且$\beta_1-\beta_2,\beta_1-\beta_3$线性无关，否则$\beta_1,\beta_2,\beta_3$线性相关.

                \item 取$2\beta_1+k_1(\beta_1+\beta_2)+k_2(\beta_2+\beta_3)$可行（取法不唯一）.
            \end{enumerate}
\end{solve}

\begin{exercise}[补充训练 L11-A-8]\label{prob:L11-A-8}
已知$A$是一个$s\times n$矩阵，证明：线性方程组$AX=\vec{b}$对任意列向量$\vec{b}_{s\times 1}$都有解的充要条件是$A$行满秩.
\end{exercise}
\noindent 来源：\href{https://github.com/yhwu-is/Linear-Algebra-Left-Undone/blob/new/%E8%AE%B2%E4%B9%89/%E4%B8%93%E9%A2%98/11%20%E6%9C%9D%E8%8A%B1%E5%A4%95%E6%8B%BE.tex}{L11-A-8}。
\begin{solve}
对于任意的$b_{s\times 1}$，$AX=b$都有解，说明$A$的列向量可以线性表出出任意$s$维向量$b$，从而$r(A)\geqslant s$，而$r(A)\leqslant s$，所以$r(A)=s$.

          反之，如果$r(A)=s$，则$A$有$s$个线性无关的列向量，它们就是$s$维空间$V$的一组基，对于$V$中任意向量$b$，都可以由$A$的列向量线性表出，从而$AX=b$有解.

          当然，也可以直接用秩不等式：$s=r(A)\leqslant r(A,b)\leqslant s$得到$r(A,b)=s$，所以$AX=b$有解.
\end{solve}

\begin{exercise}[补充训练 L11-A-9]\label{prob:L11-A-9}
设$A,B$分别是$m \times n$和$n \times s$矩阵，且$r(B)=n$，证明：若$AB=O$，则$A=O$.
\end{exercise}
\noindent 来源：\href{https://github.com/yhwu-is/Linear-Algebra-Left-Undone/blob/new/%E8%AE%B2%E4%B9%89/%E4%B8%93%E9%A2%98/11%20%E6%9C%9D%E8%8A%B1%E5%A4%95%E6%8B%BE.tex}{L11-A-9}。
\begin{solve}
$AB=0\implies r(A)+r(B)\leqslant n,r(B)=n\implies r(A)\leqslant 0\implies A=O$.
\end{solve}

\begin{exercise}[补充训练 L11-A-10]\label{prob:L11-A-10}
设$A \in \mathbf{F}^{m \times n},B \in \mathbf{F}^{(n-m) \times n}\enspace(m<n)$，$V_1,V_2$分别为齐次线性方程组$AX=\vec{0}$和$BX=\vec{0}$的解空间，证明：$\mathbf{F}^n=V_1\oplus V_2$的充要条件是$\begin{pmatrix} A \\ B \end{pmatrix}X=\vec{0}$只有零解.
\end{exercise}
\noindent 来源：\href{https://github.com/yhwu-is/Linear-Algebra-Left-Undone/blob/new/%E8%AE%B2%E4%B9%89/%E4%B8%93%E9%A2%98/11%20%E6%9C%9D%E8%8A%B1%E5%A4%95%E6%8B%BE.tex}{L11-A-10}。
\begin{solve}
必要性. 有条件可知$V_1\cap V_2=\{0\}$，$\begin{pmatrix}A\\B\end{pmatrix}\in\mathbf{F}^{n\times n}$，而$\begin{pmatrix}A\\B\end{pmatrix}x=0$只有零解，故$\begin{pmatrix}A\\B\end{pmatrix}$可逆，从而$r(A)=m,r(B)=n-m$，于是
          \begin{align*}
              \dim V_1 & =n-r(A)=n-m,       \\
              \dim V_2 & =n-r(B)=n-(n-m)=m.
          \end{align*}
          又$V_1+V_2$是$\mathbf{F}^n$的子空间，且$\dim(V_1+V_2)=\dim V_1+\dim V_2-\dim(V_1\cap V_2)=\dim\mathbf{F}^n$，故$\mathbf{F}^n=V_1\oplus V_2$.

          充分性. 若$\begin{pmatrix}A\\B\end{pmatrix}x=0$有非零解$x_1$，则$x_1\in V_1\cap V_2$. 这与$\mathbf{F}^n=V_1\oplus V_2$矛盾.
\end{solve}

\begin{exercise}[补充训练 L11-A-11]\label{prob:L11-A-11}
齐次线性方程组\[\begin{cases}
                x_2+ax_3+bx_4=0  \\
                -x_1+cx_3+dx_4=0 \\
                ax_1+cx_2-ex_4=0 \\
                bx_1+dx_2+ex_3=0
            \end{cases}\]的一般解以$x_3,x_4$作为自由未知量.
        \begin{enumerate}
            \item 求$a,b,c,d,e$满足的的条件；

            \item 求齐次线性方程组的基础解系.
        \end{enumerate}
\end{exercise}
\noindent 来源：\href{https://github.com/yhwu-is/Linear-Algebra-Left-Undone/blob/new/%E8%AE%B2%E4%B9%89/%E4%B8%93%E9%A2%98/11%20%E6%9C%9D%E8%8A%B1%E5%A4%95%E6%8B%BE.tex}{L11-A-11}。
\begin{solve}
\begin{enumerate}
                \item 由条件知，方程组系数矩阵的为2，系数矩阵
                      \[ A=\begin{pmatrix}
                              0  & 1 & a & b  \\
                              -1 & 0 & c & d  \\
                              a  & c & 0 & -e \\
                              b  & d & e & 0
                          \end{pmatrix}\rightarrow
                          \begin{pmatrix}
                              0  & 1 & a          & b       \\
                              -1 & 0 & c          & d       \\
                              0  & 0 & 0          & ad-e-bc \\
                              0  & 0 & -(ad-e-bc) & 0
                          \end{pmatrix} \]
                      故$ad-e-bc=0$.

                \item 易求得基础解系为$(c,-a,1,0)^\mathrm{T},(d,-b,0,1)^\mathrm{T}$.
            \end{enumerate}
\end{solve}

\begin{exercise}[补充训练 L11-B-1]\label{prob:L11-B-1}
设$A$为方阵，证明：$A^2=E \iff r(A+E)+r(A-E)=n$.
\end{exercise}
\noindent 来源：\href{https://github.com/yhwu-is/Linear-Algebra-Left-Undone/blob/new/%E8%AE%B2%E4%B9%89/%E4%B8%93%E9%A2%98/11%20%E6%9C%9D%E8%8A%B1%E5%A4%95%E6%8B%BE.tex}{L11-B-1}。
\begin{solve}
\begin{enumerate}
            \item 充分性：有 $(A+E)(A-E) = A^2 - E = 0$, 由定理 $BC=0 \implies r(B)+r(C) \leqslant n$，得 $r(A+E)+r(A-E) \leqslant n$；而 $r(A+E)+r(A-E) \geqslant r(A+E+A-E) = n$，故 $r(A+E)+r(A-E)=n$.
            \item 必要性：由$r(A+E)+r(A-E)=n$可得出 $A$ 的特征值仅可能为 1 或者 -1，并且可以对角化.

                因此，记 $A = PBP^{-1}$，其中 $B$ 为对角矩阵，则 $(A+E)(A-E)=P(B-E)(B+E)P^{-1}=0$，故命题得证.
        \end{enumerate}
\end{solve}

\begin{exercise}[补充训练 L11-B-2]\label{prob:L11-B-2}
设 $n\ge1$，$a_1,\ldots,a_n$ 为互异实数，$b_1,\ldots,b_n$ 为任意实数。证明：存在唯一的次数不超过 $n-1$ 的实多项式 $f$（包括零多项式），满足 $f(a_i)=b_i$（$1\le i\le n$）。
\end{exercise}
\noindent 来源：\href{https://github.com/yhwu-is/Linear-Algebra-Left-Undone/blob/new/%E8%AE%B2%E4%B9%89/%E4%B8%93%E9%A2%98/11%20%E6%9C%9D%E8%8A%B1%E5%A4%95%E6%8B%BE.tex}{L11-B-2}。
\begin{solve}
设 $f(x) = c_0+c_1x+c_2x^2+\cdots+c_{n-1}x^{n-1}$，则有
            \[\begin{cases} \begin{aligned}
                        c_0+c_1a_1+c_2a_1^2+\cdots+c_{n-1}a_1^{n-1} & = b_1,            \\
                        c_0+c_1a_2+c_2a_2^2+\cdots+c_{n-1}a_2^{n-1} & = b_2,            \\
                                                                    & \vdotswithin{ = } \\
                        c_0+c_1a_n+c_2a_n^2+\cdots+c_{n-1}a_n^{n-1} & = b_n.            \\
                    \end{aligned} \end{cases}\]
            注意此处我们研究的对象是 $(c_0, c_1, \ldots, c_{n-1})^{\mathrm{T}}$. 因为系数行列式
            \[D = \begin{vmatrix}
                    1      & a_1    & \cdots & a_1^{n-1} \\
                    1      & a_2    & \cdots & a_2^{n-1} \\
                    \vdots & \vdots & \ddots & \vdots    \\
                    1      & a_n    & \cdots & a_n^{n-1} \\
                \end{vmatrix} = \prod_{1 \leqslant j < i \leqslant n} (a_i-a_j) \neq 0.\]
            所以由 Cramer 法则，上述关于 $(c_0, c_1, \ldots, c_{n-1})^{\mathrm{T}}$ 的方程组有唯一解，所以满足条件的多项式函数 $f$ 是唯一存在的.
\end{solve}

\begin{exercise}[补充训练 L11-B-3]\label{prob:L11-B-3}
设$A$为$m \times n$矩阵，$r(A)=m$，$B$是$m$阶可逆矩阵，已知$A$的行空间$R(A^\mathrm{T})$是方程组$CX=\vec{0}$的解空间，证明：$BA$的行向量也是$CX=\vec{0}$的一个基础解系.
\end{exercise}
\noindent 来源：\href{https://github.com/yhwu-is/Linear-Algebra-Left-Undone/blob/new/%E8%AE%B2%E4%B9%89/%E4%B8%93%E9%A2%98/11%20%E6%9C%9D%E8%8A%B1%E5%A4%95%E6%8B%BE.tex}{L11-B-3}。
\begin{solve}
因 $r(A)=m$，$A$ 的 $m$ 个行向量线性无关。左乘可逆矩阵 $B$ 保持行空间与秩，所以 $BA$ 的 $m$ 个行向量仍线性无关，并张成同一行空间。把行向量转置为列向量后，有
\[C(BA)^{\mathrm T}=CA^{\mathrm T}B^{\mathrm T}=O.\]
故它们都是方程 $CX=0$ 的解；由已知该解空间就是 $A$ 的行空间，$BA$ 的各行转置后构成所求基础解系。
\end{solve}

\begin{exercise}[补充训练 L11-B-4]\label{prob:L11-B-4}
设$A$是数域$\mathbf{F}$上的一个$n$阶可逆方阵，$A$的前$r$个行向量组成的矩阵为$B$，后$n-r$个行向量组成的矩阵为$C$，$n$元线性方程组$BX=0$与$CX=0$的解空间分别为$V_1,V_2$. 证明：$\mathbf{F}^n=V_1\oplus V_2$.
\end{exercise}
\noindent 来源：\href{https://github.com/yhwu-is/Linear-Algebra-Left-Undone/blob/new/%E8%AE%B2%E4%B9%89/%E4%B8%93%E9%A2%98/11%20%E6%9C%9D%E8%8A%B1%E5%A4%95%E6%8B%BE.tex}{L11-B-4}。
\begin{solve}
先记$W=V_1+V_2$，若$\alpha\in V_1\cap V_2$，则$B\alpha=C\alpha=0$，所以
            \[A\alpha=\begin{pmatrix}
                    B \\
                    C
                \end{pmatrix}\alpha=0,\]
            由于$A$可逆，因此$\alpha=0$，即$V_1\cap V_2=\{0\}$，因此$V_1+V_2$是直和，因此只需证$W=\mathbf{F}^n$即可. 事实上，我们知道$r(B)=r,r(C)=n-r$，因此$\dim V_1=n-r,\enspace \dim V_2=n-(n-r)=r$，所以
            \[\dim W=\dim V_1+\dim V_2=n-r+r=n=\dim \mathbf{F}^n,\]
            又$W=V_1\oplus V_2\subseteq \mathbf{F}^n$，因此$W=\mathbf{F}^n$，故得证.
\end{solve}

\begin{exercise}[补充训练 L11-B-5]\label{prob:L11-B-5}
设$A$是$n$阶矩阵，且$A_{11}\neq 0$，证明：方程组$AX=\vec{b}$（$\vec{b}$为非零向量）有无穷多解的充要条件为$\operatorname{adj}(A)\vec{b}=\vec{0}$.
\end{exercise}
\noindent 来源：\href{https://github.com/yhwu-is/Linear-Algebra-Left-Undone/blob/new/%E8%AE%B2%E4%B9%89/%E4%B8%93%E9%A2%98/11%20%E6%9C%9D%E8%8A%B1%E5%A4%95%E6%8B%BE.tex}{L11-B-5}。
\begin{solve}
必要性：设$AX=b$有无穷多解，则$r(A)<n$，从而$|A|=0$，于是$\operatorname{adj}(A)b=\operatorname{adj}(A)AX=|A|X=0$.

          充分性：设$\operatorname{adj}(A)b=0$，即方程组$\operatorname{adj}(A)X=0$有非零解，则$r(\operatorname{adj}(A))<n$，又$A_{11}\neq 0$，所以$r(A)=n-1$.

          令$A=(\alpha_1,\alpha_2,\ldots,\alpha_n)$，因为$\operatorname{adj}(A)A=|A|E=O$，所以$\alpha_1,\alpha_2,\ldots,\alpha_n$为方程组$\operatorname{adj}(A)X=0$的解，又因为$A_{11}\neq 0$，所以$\alpha_2,\alpha_3,\ldots,\alpha_n$线性无关.

          由$r(\operatorname{adj}(A))=1$，得方程组$\operatorname{adj}(A)X=0$的基础解系含有$n-1$个线性无关的解向量，所以$\alpha_2,\alpha_3,\ldots,\alpha_n$为方程组$\operatorname{adj}(A)X=0$的一个基础解系.

          因为$\operatorname{adj}(A)b=0$，所以$b$为方程组$\operatorname{adj}(A)X=0$的一个解，从而$b$可由$\alpha_2,\alpha_3,\ldots,\alpha_n$线性表示，$b$也可由$\alpha_1,\alpha_2,\ldots,\alpha_n$线性表示，于是$r(A) = r\begin{pmatrix}A & b\end{pmatrix}=n-1<n$，故方程组$AX=b$有无穷多解.
\end{solve}

\begin{exercise}[补充训练 L11-B-6]\label{prob:L11-B-6}
若$n$阶矩阵$A$的各行、各列元素之和都为0，证明：$|A|$的所有元素的代数余子式都相等.
\end{exercise}
\noindent 来源：\href{https://github.com/yhwu-is/Linear-Algebra-Left-Undone/blob/new/%E8%AE%B2%E4%B9%89/%E4%B8%93%E9%A2%98/11%20%E6%9C%9D%E8%8A%B1%E5%A4%95%E6%8B%BE.tex}{L11-B-6}。
\begin{solve}
当 $n=1$ 时只有一个代数余子式，结论成立。设 $n\ge2$，令 $e=(1,\ldots,1)^{\mathrm T}$，则 $Ae=0$、$A^{\mathrm T}e=0$，故 $A$ 奇异。若 $r(A)\le n-2$，所有代数余子式都是零。若 $r(A)=n-1$，则 $\ker A=\ker A^{\mathrm T}=\operatorname{span}(e)$。由 $A\operatorname{adj}(A)=O$，伴随矩阵各列都是 $e$ 的倍数；再由 $\operatorname{adj}(A)A=O$，其各行也都是 $e^{\mathrm T}$ 的倍数。因此存在标量 $c$ 使 $\operatorname{adj}(A)=cee^{\mathrm T}$，所有代数余子式均等于 $c$。
\end{solve}

\begin{exercise}[补充训练 L11-B-7]\label{prob:L11-B-7}
已知 $\alpha_1,\alpha_2,\ldots,\alpha_s$ 是齐次线性方程组 $AX=\vec{0}$ 的一组基础解系，向量组
        \[\beta_1=t_1\alpha_1+t_2\alpha_2,\ \beta_2=t_1\alpha_2+t_2\alpha_3,\ \ldots,\ \beta_{s-1}=t_1\alpha_{s-1}+t_2\alpha_s\]
        试问当实数 $t_1,t_2$ 满足何条件时，$AX=\vec{0}$ 有基础解系包含向量 $\beta_1,\beta_2,\ldots,\beta_{s-1}$，并写出该基础解系中的其余向量.
\end{exercise}
\noindent 来源：\href{https://github.com/yhwu-is/Linear-Algebra-Left-Undone/blob/new/%E8%AE%B2%E4%B9%89/%E4%B8%93%E9%A2%98/11%20%E6%9C%9D%E8%8A%B1%E5%A4%95%E6%8B%BE.tex}{L11-B-7}。
\begin{solve}
设基础解系含 $s\ge2$ 个向量。所求条件为 $(t_1,t_2)\ne(0,0)$。
若 $t_1\ne0$，取补充向量 $\alpha_s$。这 $s$ 个向量相对于原基 $\alpha_1,\ldots,\alpha_s$ 的坐标矩阵为下三角矩阵，对角元依次为 $t_1,\ldots,t_1,1$，行列式为 $t_1^{s-1}\ne0$；因此 $\beta_1,\ldots,\beta_{s-1},\alpha_s$ 是一组基础解系。
若 $t_1=0$、$t_2\ne0$，则 $\beta_j=t_2\alpha_{j+1}$，补入 $\alpha_1$ 即可。若 $t_1=t_2=0$，则题给全部 $\beta_j$ 为零向量，不能包含在基中。若允许 $s=1$，题给向量组为空，任意 $t_1,t_2$ 都可以，直接取基础解系 $\{\alpha_1\}$。
\end{solve}

\begin{exercise}[补充训练 L11-B-8]\label{prob:L11-B-8}
已知线性方程组
        \[\begin{cases} \begin{aligned}
                    a_{11}x_1+a_{12}x_2+\cdots+a_{1,2n}x_{2n} & =0              \\
                    a_{21}x_1+a_{22}x_2+\cdots+a_{2,2n}x_{2n} & =0              \\
                                                              & \vdotswithin{=} \\
                    a_{n1}x_1+a_{n2}x_2+\cdots+a_{n,2n}x_{2n} & =0              \\
                \end{aligned}\end{cases}\]
        的一个基础解系为$(b_{11},b_{12},\ldots,b_{1,2n})^\mathrm{T},(b_{21},b_{22},\ldots,b_{2,2n})^\mathrm{T},(b_{n1},b_{n2},\ldots,b_{n,2n})^\mathrm{T}$，求解线性方程组
        \[\begin{cases} \begin{aligned}
                    b_{11}x_1+b_{12}x_2+\cdots+b_{1,2n}x_{2n} & =0              \\
                    b_{21}x_1+b_{22}x_2+\cdots+b_{2,2n}x_{2n} & =0              \\
                                                              & \vdotswithin{=} \\
                    b_{n1}x_1+b_{n2}x_2+\cdots+b_{n,2n}x_{2n} & =0              \\
                \end{aligned} \end{cases}.\]
        （注：本题的一般形式：设$A=(a_{ij})_{m\times n},m<n,r(A)=m$；$B=(b_{ij})_{n\times(n-m)},m<n,r(B)=n-m$.已知齐次线性方程组$A\vec{X}=\vec{0}$的解空间$N(A)=R(B)$（矩阵$B$的列空间），求齐次线性方程组$\displaystyle\sum_{i=1}^{n}b_{ij}y_i=0,j=1,2,\ldots,n-m$的一个基础解系.）
\end{exercise}
\noindent 来源：\href{https://github.com/yhwu-is/Linear-Algebra-Left-Undone/blob/new/%E8%AE%B2%E4%B9%89/%E4%B8%93%E9%A2%98/11%20%E6%9C%9D%E8%8A%B1%E5%A4%95%E6%8B%BE.tex}{L11-B-8}。
\begin{solve}
令原方程组$AX=0$，解系组成矩阵$B$. 则新方程组为$B^\mathrm{T}X=0$. 而$B^\mathrm{T}A^\mathrm{T}=(AB)^\mathrm{T}=0$，推测$A^\mathrm{T}$的列向量为基础解系. 而由维数公式，
          \[ r(A)=\dim V-\dim\ker(A) = 2n-n=n. \]
          故$A^\mathrm{T}$列空间维数为$n$，即$r_c(A^\mathrm{T})=n$. 而由题意$r(B)=n$（否则不是基础解系），故
          \[ \dim{\ker(B)}=\dim{V^{\prime}}-r(B)=2n-n=n=r_c(A^\mathrm{T}). \]
          猜想成立.
\end{solve}

\begin{exercise}[补充训练 L11-B-11]\label{prob:L11-B-11}
设$A$是元素全为1的$n$阶方阵.
        \begin{enumerate}
            \item 求行列式$|aE+bA|$，其中$a,b$为实常数；

            \item 已知$1<r(aE+bA)<n$，试确定$a,b$满足的条件，并求下列线性子空间的维数：
                  \[W=\{X \mid (aE+bA)X=0,\enspace X\in\mathbf{R}^n\}.\]
        \end{enumerate}
\end{exercise}
\noindent 来源：\href{https://github.com/yhwu-is/Linear-Algebra-Left-Undone/blob/new/%E8%AE%B2%E4%B9%89/%E4%B8%93%E9%A2%98/11%20%E6%9C%9D%E8%8A%B1%E5%A4%95%E6%8B%BE.tex}{L11-B-11}。
\begin{solve}
\begin{enumerate}
                \item 容易计算$|aE=bA|=(a+nb)a^{n-1}$.

                \item 由$1<r(aE+bA)<n$知$|aE+bA|=0$. 故$a\neq 0$，且$a+nb=0$，此时$aE+bA$左上角的$n-1$阶子式
                      \[ \begin{vmatrix}
                              a+b    & b      & \cdots & b      \\
                              b      & a+b    & \cdots & b      \\
                              \vdots & \vdots & \ddots & \vdots \\
                              b      & b      & \cdots & a+b
                          \end{vmatrix}=(a+(n-1)b)a^{n-2}=\frac{a^{n-1}}{n}\neq 0, \]
                      故$\dim{W}=n-r(aE+bA)=n-(n-1)=1$.
            \end{enumerate}
\end{solve}

\begin{exercise}[补充训练 L11-B-13]\label{prob:L11-B-13}
设$A=(a_{ij})_{m\times n}$，$\vec{b}$和$X$为$m$元列向量，$Y$为$n$元列向量，证明：
        \begin{enumerate}
            \item 若$AY=\vec{b}$有解，则$A^\mathrm{T}X=\vec{0}$的任一组解都满足$\vec{b}^\mathrm{T}X=0$；

            \item 方程组$AY=\vec{b}$有解的充要条件是方程组
\[\begin{pmatrix}
                          A^\mathrm{T} \\ \vec{b}^\mathrm{T}
                      \end{pmatrix}X=\begin{pmatrix}
                          \vec{0} \\ 1
                      \end{pmatrix}\]
无解（其中$\vec{0}$是$n$元零向量）.
        \end{enumerate}
\end{exercise}
\noindent 来源：\href{https://github.com/yhwu-is/Linear-Algebra-Left-Undone/blob/new/%E8%AE%B2%E4%B9%89/%E4%B8%93%E9%A2%98/11%20%E6%9C%9D%E8%8A%B1%E5%A4%95%E6%8B%BE.tex}{L11-B-13}。
\begin{solve}
\begin{enumerate}
                \item 由条件 $(A^\mathrm{T}X)^\mathrm{T}Y=X^\mathrm{T}(AY)=X^\mathrm{T}\vec{b}=0$，故 $\vec{b}^\mathrm{T}X=0$.
                \item 必要性由上一问显然成立.

                充分性：由条件可知，若 $X_0$ 使得 $A^\mathrm{T}X_0=\vec{0}$，则 $\vec{b}^\mathrm{T}X_0=0$，否则令 $X_1 =\dfrac{X_0}{\vec{b}^\mathrm{T}X_0}$，则 $X_1$ 是方程组 $\begin{pmatrix} A^\mathrm{T} \\ \vec{b}^\mathrm{T} \end{pmatrix}$ 的解，与条件矛盾.

                也就是说 $A^\mathrm{T}X=\vec{0}$ 与 $\begin{pmatrix} A^\mathrm{T} \\ \vec{b}^\mathrm{T} \end{pmatrix} X = \vec{0}$ 同解，从而 $r(A^\mathrm{T}) = r\begin{pmatrix} A^\mathrm{T} \\ \vec{b}^\mathrm{T} \end{pmatrix}$，这表明 $\vec{b}^\mathrm{T}$ 可由 $A^\mathrm{T}$ 的行向量（即 $A$ 的列向量）线性表示，因此 $AY=\vec{b}$ 有解.
            \end{enumerate}
\end{solve}

\begin{exercise}[补充训练 L11-B-14]\label{prob:L11-B-14}
判断：设 $A$ 是复数域上 $m \times n$ 阶矩阵，则矩阵秩 $r\left(A^T A\right)=r(A)$.
\end{exercise}
\noindent 来源：\href{https://github.com/yhwu-is/Linear-Algebra-Left-Undone/blob/new/%E8%AE%B2%E4%B9%89/%E4%B8%93%E9%A2%98/11%20%E6%9C%9D%E8%8A%B1%E5%A4%95%E6%8B%BE.tex}{L11-B-14}。
\begin{solve}
假. 取$A=\begin{pmatrix}
                1-i & 1+i  \\
                1+i & -1+i
            \end{pmatrix}$，有$A^\mathrm{T}A=0$.
\end{solve}

\begin{exercise}[补充训练 L11-B-15]\label{prob:L11-B-15}
证明：对于$m \times n$实矩阵$A$，方程$A^\mathrm{T}AX = A^\mathrm{T}\vec{b}$总是有解，且$A$为方阵时，$A^\mathrm{T}AX = \vec{0}$和$AX=\vec{0}$同解. 当$r(A)=n$时求其解，并证明$A(A^\mathrm{T}A)^{-1}A^\mathrm{T}$是幂等的对称矩阵.
\end{exercise}
\noindent 来源：\href{https://github.com/yhwu-is/Linear-Algebra-Left-Undone/blob/new/%E8%AE%B2%E4%B9%89/%E4%B8%93%E9%A2%98/11%20%E6%9C%9D%E8%8A%B1%E5%A4%95%E6%8B%BE.tex}{L11-B-15}。
\begin{solve}
由于 $\vec{b}$ 的任意性，$A^\mathrm{T}\vec{b}$ 可以取遍 $R(A^\mathrm{T})$ 中的任意元素，这里 $R(A^\mathrm{T})$ 表示 $A^\mathrm{T}$ 的列向量张成的空间. 因此我们只需证明 $R(A^\mathrm{T}) = R(A^\mathrm{T}A)$.

            $A^\mathrm{T}A$ 的每一列都是 $A^\mathrm{T}$ 的列向量的线性组合，因此 $R(A^\mathrm{T}A) \subseteq R(A^\mathrm{T})$. 而根据原讲义“线性方程组理论与秩不等式”处 的结论，有 $r(A^\mathrm{T}A) = r(A)$，因此 $A^\mathrm{T}A$ 的列向量组和 $A^\mathrm{T}$ 的列秩相等，即 $\dim R(A^\mathrm{T}) = \dim R(A^\mathrm{T}A)$. 从而 $R(A^\mathrm{T}) = R(A^\mathrm{T}A)$，命题得证.
\end{solve}

\begin{exercise}[补充训练 L11-B-16]\label{prob:L11-B-16}
设$A,B,C$为$n$阶实方阵，且$BAA^\mathrm{T}=CAA^\mathrm{T}$，证明：$BA=CA$.
\end{exercise}
\noindent 来源：\href{https://github.com/yhwu-is/Linear-Algebra-Left-Undone/blob/new/%E8%AE%B2%E4%B9%89/%E4%B8%93%E9%A2%98/11%20%E6%9C%9D%E8%8A%B1%E5%A4%95%E6%8B%BE.tex}{L11-B-16}。
\begin{solve}
记 $D=B-C$，由题设 $DAA^{\mathrm T}=O$，右乘 $D^{\mathrm T}$ 得
\[(DA)(DA)^{\mathrm T}=O.\]
左端第 $i$ 个对角元是实矩阵 $DA$ 第 $i$ 行元素的平方和，故各行都为零，即 $DA=O$。因此 $BA=CA$。
\end{solve}

\begin{exercise}[补充训练 L11-B-17]\label{prob:L11-B-17}
设$A$为数域$\mathbf{F}$上的$n$阶方阵，又设线性空间$\mathbf{F^n}$的两个子空间为$W_1=\{X\in\mathbf{F}^n \mid AX=\vec{0}\}$，$W_2=\{X\in\mathbf{F}^n \mid (A-E)X=\vec{0}\}$. 证明：$A^2=A \iff \mathbf{F}^n=W_1\oplus W_2$.
\end{exercise}
\noindent 来源：\href{https://github.com/yhwu-is/Linear-Algebra-Left-Undone/blob/new/%E8%AE%B2%E4%B9%89/%E4%B8%93%E9%A2%98/11%20%E6%9C%9D%E8%8A%B1%E5%A4%95%E6%8B%BE.tex}{L11-B-17}。
\begin{solve}
若 $A^2=A$，则任意 $x\in\mathbf F^n$ 可写成 $x=(x-Ax)+Ax$。其中 $A(x-Ax)=0$，而 $(A-I)Ax=0$，故两项分别属于 $W_1,W_2$。若 $z\in W_1\cap W_2$，则同时有 $Az=0$、$Az=z$，所以 $z=0$，从而是直和。
反过来，若 $\mathbf F^n=W_1\oplus W_2$，写 $x=u+v$，其中 $Au=0$、$Av=v$，则 $A^2x=Av=v=Ax$。因为这对任意 $x$ 成立，故 $A^2=A$。
\end{solve}

\begin{exercise}[补充训练 L11-B-18]\label{prob:L11-B-18}
$n$阶方阵$A,B$满足$AB=BA$，证明：$r(AB)+r(A+B) \leqslant r(A)+r(B)$.
\end{exercise}
\noindent 来源：\href{https://github.com/yhwu-is/Linear-Algebra-Left-Undone/blob/new/%E8%AE%B2%E4%B9%89/%E4%B8%93%E9%A2%98/11%20%E6%9C%9D%E8%8A%B1%E5%A4%95%E6%8B%BE.tex}{L11-B-18}。
\begin{solve}
方法一. 用分块矩阵的方法，我们知道
          \[ \begin{pmatrix}
                  A & O \\
                  O & B
              \end{pmatrix}
              \rightarrow
              \begin{pmatrix}
                  A & O \\
                  A & B
              \end{pmatrix}
              \rightarrow
              \begin{pmatrix}
                  A & A   \\
                  A & A+B
              \end{pmatrix}. \]
          结合$AB=BA$，我们知道
          \[ \begin{pmatrix}
                  A & A   \\
                  A & A+B
              \end{pmatrix}
              \underbrace{
                  \begin{pmatrix}
                      A+B & O \\
                      -A  & E
                  \end{pmatrix}
              }_{\text{非广义初等变换，难以想到}}
              =
              \begin{pmatrix}
                  AB & A   \\
                  O  & A+B
              \end{pmatrix}. \]
          于是
          \[ r(A)+r(B)=r
              \begin{pmatrix}
                  A & O \\
                  O & B
              \end{pmatrix}=
              \begin{pmatrix}
                  A & A   \\
                  A & A+B
              \end{pmatrix}\geqslant r
              \begin{pmatrix}
                  AB & A   \\
                  O  & A+B
              \end{pmatrix}\geqslant
              r(AB)+r(A+B). \]
          方法二. 设方程组$AX=0$与$BX=0$的解空间分别是$V_1, V_2$，方程组$ABX=BAX=0$与$(A+B)X=0$的解空间分别为$W_1, W_2$，则$V_1\subseteq W_1, V_2\subseteq W_1$，从而$V_1+V_2\subseteq W_1$，同时$V_1\cap V_2\subseteq W_1$，同时$V_1\cap V_2\subseteq W_2$，利用维数公式就有
          \[ \dim V_1+\dim V_2=\dim(V_1+V_2)+\dim(V_1\cap V_2)\leqslant \dim W_1+\dim W_2. \]
          即
          \[ (n-r(A))+(n-r(B))\leqslant (n-r(AB))+(n-r(A+B)). \]
          化简便知$r(A)+r(B)\geqslant r(AB)+r(A+B)$.
\end{solve}

\begin{exercise}[补充训练 L11-B-19]\label{prob:L11-B-19}
请按序证明以下结论：
        \begin{enumerate}
            \item $A,B$分别是$s \times m,m \times n$矩阵，则$ABX=\vec{0}$与$BX=\vec{0}$同解的充要条件是$r(AB)=r(B)$；

            \item $A,B$分别是$s \times m,m \times n$矩阵，且$r(AB)=r(B)$，则对任意的$n \times t$矩阵都有$r(ABC)=r(BC)$；

            \item 设$A$是$n$阶方阵，则存在正整数$k$使得$r(A^k)=r(A^{k+1})=r(A^{k+2})=\cdots$，且对任意正整数$m$，有$r(A^n)=r(A^{n+m})$.
        \end{enumerate}
\end{exercise}
\noindent 来源：\href{https://github.com/yhwu-is/Linear-Algebra-Left-Undone/blob/new/%E8%AE%B2%E4%B9%89/%E4%B8%93%E9%A2%98/11%20%E6%9C%9D%E8%8A%B1%E5%A4%95%E6%8B%BE.tex}{L11-B-19}。
\begin{solve}
\begin{enumerate}
                \item 必要性：$ABX=0$与$BX=0$同解可知它们基础解系所含向量个数相同，即
                      \[ n-r(AB)=n-r(B)\implies r(AB)=r(B). \]
                      充分性：由必要性，当$r(AB)=r(B)$时，$ABX=0$与$BX=0$的基础解系所含向量个数相同，而$BX=0$的解都是$ABX=0$的解，所以$ABX=0$与$BX=0$同解.

                \item $r(AB)=r(B)$说明$ABX=0$与$BX=0$同解，用$CX$代替$X$就得到$ABCX=0$与$BCX=0$同解，从而有$r(ABC)=r(BC)$.\\
                      \textbf{推论.}设$A$是一个方阵，且存在正整数$k$使得$r(A^{k+1})=r(A^k)$，递推就有
                      \[ r(A^k)=r(A^{k+1})=r(A^{k+2})=\cdots. \]

                \item 当$A$可逆时，结论是显然的. 当$A$不可逆时，有$r(A)\leqslant n-1$，现在考虑$n+1$个矩阵$A,A^2,\ldots,A^{n+1}$，有
                      \[ n-1\geqslant r(A)\geqslant r(A^2)\geqslant\cdots\geqslant r(A^n)\geqslant r(A^{n+1})\geqslant 0. \]
                      这$n+1$个矩阵的秩只能从$0,1,\ldots,n-1$这$n$个数中取，所以必有两个矩阵的秩相同，即存在$m(1\leqslant m\leqslant n)$使得$r(A^m)=r(A^{m+1})$，由上面的推论可得：
                      \[ r(A^m)=r(A^{m+1})=\cdots=r(A^{n})=r(A^{n+1})=\cdots. \]
                      特别的，对于任意的正整数$k$有$r(A^n)=r(A^{n+k})$.
            \end{enumerate}
\end{solve}

\begin{exercise}[补充训练 L11-B-20]\label{prob:L11-B-20}
如果齐次线性方程组\[\begin{cases}
                x_1+x_2+bx_3-x_4+x_5=0   \\
                2x_1+3x_2+x_3+x_4-2x_5=0 \\
                x_2+ax_3+3x_4-4x_5=0     \\
                -3x_1-3x_2-3bx_3+bx_4+(a+2)x_5=0
            \end{cases}\]的解空间是3维的，试求$a,b$的值与解空间的基. 解空间可能为2维吗？
\end{exercise}
\noindent 来源：\href{https://github.com/yhwu-is/Linear-Algebra-Left-Undone/blob/new/%E8%AE%B2%E4%B9%89/%E4%B8%93%E9%A2%98/11%20%E6%9C%9D%E8%8A%B1%E5%A4%95%E6%8B%BE.tex}{L11-B-20}。
\begin{solve}
对系数矩阵施行 $R_2\leftarrow R_2-2R_1$、$R_3\leftarrow R_3-R_2$（第二步使用更新后的 $R_2$）、$R_4\leftarrow R_4+3R_1$，得到
\[\begin{pmatrix}1&1&b&-1&1\\0&1&1-2b&3&-4\\0&0&a+2b-1&0&0\\0&0&0&b-3&a+5\end{pmatrix}.\]
前两行恒线性无关。因此解空间三维等价于后两行同时为零，即 $a=-5,b=3$。此时令自由变量为 $x_3,x_4,x_5$，得到
\[x_2=5x_3-3x_4+4x_5,\qquad x_1=-8x_3+4x_4-5x_5.\]
一组基础解系为
\[(-8,5,1,0,0)^{\mathrm T},\quad(4,-3,0,1,0)^{\mathrm T},\quad(-5,4,0,0,1)^{\mathrm T}.\]
解空间也可能为二维。最后两行若同时非零则相互独立；而第四行全零时必有 $a=-5,b=3$，第三行也随之为零。因此秩恰为 $3$ 的充要条件是 $a+2b=1$ 且 $b\ne3$。例如 $a=1,b=0$ 即满足要求。
\end{solve}

\begin{exercise}[补充训练 L11-B-21]\label{prob:L11-B-21}
设$W_1,W_2$分别为$n$元齐次线性方程组$AX=\vec{0}$和$BX=\vec{0}$的解空间，试构造两个$n$元齐次线性方程组，使它们的解空间分别为$W_1 \cap W_2$和$W_1+W_2$.
\end{exercise}
\noindent 来源：\href{https://github.com/yhwu-is/Linear-Algebra-Left-Undone/blob/new/%E8%AE%B2%E4%B9%89/%E4%B8%93%E9%A2%98/11%20%E6%9C%9D%E8%8A%B1%E5%A4%95%E6%8B%BE.tex}{L11-B-21}。
\begin{solve}
$W_1\cap W_2:\begin{pmatrix}
                  A \\ B
              \end{pmatrix}X=0$. \qquad
          $W_1+ W_2:\begin{cases}
                  AX = 0 \\
                  BX = 0
              \end{cases}$.
\end{solve}

\begin{exercise}[补充训练 L11-B-22]\label{prob:L11-B-22}
已知方程组$\begin{cases}
                x_1+x_2+ax_3+x_4=1 \\ -x_1+x_2-x_3+bx_4=2 \\ 2x_1+x_2+x_3+x_4=c
            \end{cases}$与$\begin{cases}
                x_1+x_4=-1 \\ x_2-2x_4=d \\ x_3+x_4=e
            \end{cases}$同解，求$a,b,c,d,e$.
\end{exercise}
\noindent 来源：\href{https://github.com/yhwu-is/Linear-Algebra-Left-Undone/blob/new/%E8%AE%B2%E4%B9%89/%E4%B8%93%E9%A2%98/11%20%E6%9C%9D%E8%8A%B1%E5%A4%95%E6%8B%BE.tex}{L11-B-22}。
\begin{solve}
方法一. 令$x_4=t$，则方程组$\begin{cases}
                x_1+x_4=-1 \\
                x_2-2x_4=d \\
                x_3+x_4=e
            \end{cases}$的一般解为$\begin{cases}
                x_1=-1-t \\
                x_2=d+2t \\
                x_3=e-t  \\
                x_4=t
            \end{cases}$.

        代入方程组$\begin{cases}
                x_1+x_2+ax_3+x_4=1  \\
                -x_1+x_2-x_3+bx_4=2 \\
                2x_1+x_2+x_3+x_4=c
            \end{cases}$可得$\begin{cases}
                (2-a)t=2-d-ae \\
                (b+4)t=1-d+e  \\
                0=c-d-e+2
            \end{cases}$. 由$t$的任意性，可得$a=2,b=-4$. 从而$\begin{cases}
                0=2-d-2e \\
                0=1-d+e  \\
                0=c-d-e+2
            \end{cases}$，解得$d=\dfrac{4}{3},e=\dfrac{1}{3},c=-\dfrac{1}{3}$.

        方法二. 由于
        \[\begin{pmatrix}
                A & b \\
                B & d
            \end{pmatrix}=
            \begin{pmatrix}
                1  & 1 & a  & 1  & 1  \\
                -1 & 1 & -1 & b  & 2  \\
                2  & 1 & 1  & 1  & c  \\
                1  & 0 & 0  & 1  & -1 \\
                0  & 1 & 0  & -2 & d  \\
                0  & 0 & 1  & 1  & e
            \end{pmatrix}\rightarrow
            \begin{pmatrix}
                1 & 0 & 0 & 0     & 0            \\
                0 & 0 & 0 & b-a+6 & 3-2d+(1-a)e  \\
                0 & 0 & 0 & 2a-4  & c-2+d(2a-1)e \\
                0 & 0 & 0 & a-2   & d-2+ae       \\
                0 & 1 & 0 & 0     & 0            \\
                0 & 0 & 1 & 0     & 0
            \end{pmatrix}.\]
        易知$r(B)=3$，由$r\begin{pmatrix}
                A & b \\
                B & d
            \end{pmatrix}=r(B)$可得$\begin{cases}
                a-2=0          \\
                2a-4=0         \\
                b-a+6=0        \\
                3-2d+(1-a)e=0  \\
                c-2+d(2a-1)e=0 \\
                d-2+ae=0
            \end{cases}$，解得$\begin{cases}
                a=2          \\
                b=-4         \\
                c=-\dfrac 13 \\[1ex]
                d=\dfrac 43  \\[1ex]
                e=\dfrac 13
            \end{cases}$.
\end{solve}

\begin{exercise}[补充训练 L11-B-23]\label{prob:L11-B-23}
设有两个非齐次线性方程组 (1) 和 (2)，它们的通解分别是$X=\gamma+t_1\eta_1+t_2\eta_2, X=\delta+k_1\xi_1+k_2\xi_2$. 其中$\gamma=(5,-3,0,0)^\mathrm{T},\eta_1=(-6,5,1,0)^\mathrm{T},\eta_2=(-5,4,0,1)^\mathrm{T},\delta=(-11,3,0,0)^\mathrm{T},\xi_1=(8,-1,1,0)^\mathrm{T},\xi_2=(10,-2,0,1)^\mathrm{T}$，求这两个方程组的公共解.
\end{exercise}
\noindent 来源：\href{https://github.com/yhwu-is/Linear-Algebra-Left-Undone/blob/new/%E8%AE%B2%E4%B9%89/%E4%B8%93%E9%A2%98/11%20%E6%9C%9D%E8%8A%B1%E5%A4%95%E6%8B%BE.tex}{L11-B-23}。
\begin{solve}
注意如果$X$是两个方程组的公共解，这等价于存在$t_1,t_2,k_1,k_2$使得
          \[ X=\gamma+t_1\eta_1+t_2\eta_2=\delta+k_1\xi_1+k_2\xi_2. \]
          从而对应有
          \[ t_1\eta_1+t_2\eta_2-k_1\xi_1-k_2\xi_2=\gamma-\delta. \]
          将$t_1,t_2,k_1,k_2$设为未知量，对上述方程组的增广矩阵进行初等行变换，可得
          \[ (\eta_1,\eta_2,-\xi_1,-\xi_2,\delta-\gamma)=
              \begin{pmatrix}
                  -6 & -5 & -8 & -10 & -16 \\
                  5  & 4  & 1  & 2   & 6   \\
                  1  & 0  & -1 & 0   & 0   \\
                  0  & 1  & 0  & -1  & 0
              \end{pmatrix}\rightarrow
              \begin{pmatrix}
                  1 & 0 & -1 & 0  & 0  \\
                  0 & 1 & 0  & -1 & 0  \\
                  0 & 0 & 1  & 1  & 1  \\
                  0 & 0 & 0  & 1  & 2.
              \end{pmatrix} \]
          故$(t_1,t_2,k_1,k_2)^\mathrm{T}$有唯一解$(-1,2,-1,2)^\mathrm{T}$. 因此公共解为
          \[ X=\gamma-\eta_1+2\eta_2=\delta-\xi_1+2\xi_2=(1,0,-1,2)\mathrm{T}. \]
\end{solve}

\begin{exercise}[补充训练 L11-B-24]\label{prob:L11-B-24}
若方程组$\begin{cases}
                x_1+x_2+x_3=0   \\
                x_1+2x_2+ax_3=0 \\
                x_1+4x_2+a^2x_3=0
            \end{cases}$与$x_1+2x_2+x_3=a-1$有公共解，求$a$的值及所有公共解.
\end{exercise}
\noindent 来源：\href{https://github.com/yhwu-is/Linear-Algebra-Left-Undone/blob/new/%E8%AE%B2%E4%B9%89/%E4%B8%93%E9%A2%98/11%20%E6%9C%9D%E8%8A%B1%E5%A4%95%E6%8B%BE.tex}{L11-B-24}。
\begin{solve}
方程组 (1)、(2) 有公共解，即方程组$\begin{cases}
                x_1+x_2+x_3=0,     \\
                x_1+2x_2+ax_3=0,   \\
                x_1+4x_2+a^2x_3=0, \\
                x_1+2x_2+x_3=a-1
            \end{cases}$有解，对其增广矩阵进行初等行变换，
        \begin{align*}
            (A,\vec{b})= & \begin{pmatrix}
                           1 & 1 & 1   & 0   \\
                           1 & 2 & a   & 0   \\
                           1 & 4 & a^2 & 0   \\
                           1 & 2 & 1   & a-1
                       \end{pmatrix}\rightarrow
            \begin{pmatrix}
                1 & 1 & 1     & 0   \\
                0 & 1 & a-1   & 0   \\
                0 & 3 & a^2-1 & 0   \\
                0 & 1 & 0     & a-1
            \end{pmatrix}\rightarrow            \\
                     & \begin{pmatrix}
                           1 & 1 & 1        & 0   \\
                           0 & 1 & a-1      & 0   \\
                           0 & 0 & a^2-3a+2 & 0   \\
                           0 & 0 & 1-a      & a-1
                       \end{pmatrix}\rightarrow
            \begin{pmatrix}
                1 & 1 & 1   & 0          \\
                0 & 1 & a-1 & 0          \\
                0 & 0 & 1-a & a-1        \\
                0 & 0 & 0   & (a-1)(a-2)
            \end{pmatrix}.
        \end{align*}
        当$a\neq 1$且$a\neq 2$时，方程组 (1)、(2) 没有公共解.

        当$a=1$时，
\[(A,\vec{b})\rightarrow\begin{pmatrix}
                1 & 1 & 1 & 0 \\
                0 & 1 & 0 & 0 \\
                0 & 0 & 0 & 0 \\
                0 & 0 & 0 & 0
            \end{pmatrix}\rightarrow\begin{pmatrix}
                1 & 0 & 1 & 0 \\
                0 & 1 & 0 & 0 \\
                0 & 0 & 0 & 0 \\
                0 & 0 & 0 & 0
            \end{pmatrix}\]
，因为$r(A)=r(A,\vec{b})=2$，所以方程组 (1)、(2) 有公共解，公共解为$X=k\begin{pmatrix} -1 \\ 0  \\ 1 \end{pmatrix}$（其中$k$为任意常数）.

        当$a=2$时，
\[(A,\vec{b})\rightarrow\begin{pmatrix}
                1 & 1 & 1  & 0 \\
                0 & 1 & 1  & 0 \\
                0 & 0 & -1 & 1 \\
                0 & 0 & 0  & 0
            \end{pmatrix}\rightarrow\begin{pmatrix}
                1 & 0 & 0 & 0  \\
                0 & 1 & 0 & 1  \\
                0 & 0 & 1 & -1 \\
                0 & 0 & 0 & 0
            \end{pmatrix}\]
，方程组 (1)、(2) 有唯一的公共解为$X=\begin{pmatrix}
                0 \\
                1 \\
                -1
            \end{pmatrix}$.
\end{solve}

\begin{exercise}[补充训练 L11-C-1]\label{prob:L11-C-1}
用方程组的理论证明：一个$n$次多项式不可能有多于$n$个不同的根.
\end{exercise}
\noindent 来源：\href{https://github.com/yhwu-is/Linear-Algebra-Left-Undone/blob/new/%E8%AE%B2%E4%B9%89/%E4%B8%93%E9%A2%98/11%20%E6%9C%9D%E8%8A%B1%E5%A4%95%E6%8B%BE.tex}{L11-C-1}。
\begin{solve}
令$f(x)=a_nx^n+\cdots+a_1x+a_0$为一个$n$次多项式，设$\lambda_1,\ldots,\lambda_{n+1}$为$f(x)$的$n+1$个不同的根，则有齐次线性方程组$\begin{cases} \begin{aligned}
                \lambda_1^n a_n+\cdots+\lambda_1a_1+a_0         & =0,               \\
                \lambda_2^n a_n+\cdots+\lambda_2a_1+a_0         & =0,               \\
                                                                & \vdotswithin{ = } \\
                \lambda_{n+1}^n a_n+\cdots+\lambda_{n+1}a_1+a_0 & =0.
            \end{aligned} \end{cases}$. 其系数矩阵为
\[A=\begin{pmatrix}
            \lambda_1^n     & \lambda_1^{n-1}     & \cdots & 1      \\
            \lambda_2^n     & \lambda_2^{n-1}     & \cdots & 1      \\
            \vdots          & \vdots              & \ddots & \vdots \\
            \lambda_{n+1}^n & \lambda_{n+1}^{n-1} & \cdots & 1
        \end{pmatrix}\]
. 显然$|A|\neq 0$，从而$a_n=a_{n-1}=\cdots=a_1=a_0=0$，这与$f(x)$是$n$次多项式矛盾.
\end{solve}

\begin{exercise}[补充训练 L11-C-3]\label{prob:L11-C-3}
已知$A$是$n$阶对称矩阵，$\beta$为$n$元非零列向量，$B=\begin{pmatrix}
                A & \beta \\ \beta^\mathrm{T} & 0
            \end{pmatrix}$，证明：
        \begin{enumerate}
            \item 若$r(A)=n$，则$B$可逆的充要条件是$\beta^\mathrm{T}A^{-1}\beta \neq \vec{0}$；

            \item 若$r(A)=r$，则$r(B)=r$的充要条件是方程组$\begin{cases}
                          AX=\beta \\ \beta^\mathrm{T}X=\vec{0}
                      \end{cases}$有解；

            \item 若$r(A)=n-1$，则$B$可逆的充要条件是$AX=\beta$无解.
        \end{enumerate}
\end{exercise}
\noindent 来源：\href{https://github.com/yhwu-is/Linear-Algebra-Left-Undone/blob/new/%E8%AE%B2%E4%B9%89/%E4%B8%93%E9%A2%98/11%20%E6%9C%9D%E8%8A%B1%E5%A4%95%E6%8B%BE.tex}{L11-C-3}。
\begin{solve}
\begin{enumerate}
                \item 由于$r(A)=n$，所以$A$可逆，由打洞原理可知$|B|=|A|(0-\beta' A^{-1}\beta)=-|A|\beta' A^{-1}\beta$，从而$B$可逆的充要条件为$\beta' A^{-1}\beta\neq 0$.

                \item 必要性. 由于
                      \[ r=r(A)\leqslant r(A,\beta)\leqslant r\begin{pmatrix} A & \beta \\ \beta' & 0 \end{pmatrix}=r(B)=r \]
                      再结合$A'=A$，可知
                      \[ r\begin{pmatrix} A & \beta \\ \beta' & 0 \end{pmatrix} =r(A,\beta)=r\begin{pmatrix} A \\ \beta' \end{pmatrix} \]
                      于是由原讲义定理 15.1可知方程组$\begin{pmatrix} A \\ \beta' \end{pmatrix}X=\begin{pmatrix} \beta \\ 0 \end{pmatrix}$有解，即$\begin{cases} AX=\beta, \\ \beta' X=0. \end{cases}$有解.

                      充分性. 由于$\begin{cases} AX=\beta \\ \beta' X=0 \end{cases}$有解，从而$AX=\beta$也有解，即有$r(A)=r(A,\beta)$. 另外$\begin{cases} AX=\beta \\ \beta' X=0 \end{cases}$有解也说明$\begin{pmatrix} A \\ \beta' \end{pmatrix}X=\begin{pmatrix} \beta \\ 0 \end{pmatrix}$有解，于是结合原讲义定理 15.1可知$r\begin{pmatrix} A & \beta \\ \beta' & 0 \end{pmatrix}=r\begin{pmatrix} A \\ \beta' \end{pmatrix}$. 而显然$r\begin{pmatrix} A \\ \beta' \end{pmatrix}=r(A,\beta)=r(A)$，于是$r(B)=r\begin{pmatrix} A & \beta \\ \beta' & 0 \end{pmatrix}=r$.

                \item 必要性. 若$B$可逆，则$B$的行向量组线性无关，从而$r(A,\beta)=n$，又由于$r(A)=n-1$，所以$r(A,\beta)=r(A)+1$，从而由原讲义定理 15.1知$AX=\beta$无解.

                      充分性. 由于$r(A)=n-1$，若$AX=\beta$无解，则由原讲义定理 15.1可知$r(A,\beta)=r(A)+1=n$，于是
                      \[ r(B)=r\begin{pmatrix}
                              A      & \beta \\
                              \beta' & 0
                          \end{pmatrix}\geqslant r(A,\beta)=n \]
                      若$r(B)=n$，则$r(B)=r\begin{pmatrix}
                              A      & \beta \\
                              \beta' & 0
                          \end{pmatrix}=r(A,\beta)$，取转置结合$A'=A$有
\[r\begin{pmatrix}
                              A      & \beta \\
                              \beta' & 0
                          \end{pmatrix}=r\begin{pmatrix}
                              A \\
                              \beta'
                          \end{pmatrix}\]
，再次结合原讲义定理 15.1可知方程组$\begin{cases}
                              AX=\beta, \\
                              \beta' X=0.
                          \end{cases}$有解，特别地，$AX=\beta$也有解，这就与已知产生了矛盾. 所以$r(B)=n+1$，即$B$可逆.
            \end{enumerate}
\end{solve}

\begin{exercise}[补充训练 L11-C-4]\label{prob:L11-C-4}
讨论$b_1,b_2,\ldots,b_n\enspace(n \geqslant 2)$满足什么条件时，下列方程组
        \[\begin{cases} \begin{aligned}
                    x_1+x_2     & =b_1            \\
                    x_2+x_3     & =b_2            \\
                                & \vdotswithin{=} \\
                    x_{n-1}+x_n & =b_{n-1}        \\
                    x_n+x_1     & =b_n
                \end{aligned} \end{cases}\]有解，并求解.
\end{exercise}
\noindent 来源：\href{https://github.com/yhwu-is/Linear-Algebra-Left-Undone/blob/new/%E8%AE%B2%E4%B9%89/%E4%B8%93%E9%A2%98/11%20%E6%9C%9D%E8%8A%B1%E5%A4%95%E6%8B%BE.tex}{L11-C-4}。
\begin{solve}
\begin{enumerate}
                \item 当$n$为偶数时，将增广矩阵$\bar{A}$的第$i\enspace(i=1,2,\ldots,n-1)$行乘以$(-1)^i$加到最后一行，得
                      \[ \bar{A}\rightarrow\begin{pmatrix}
                              1      & 1      & \cdots & 0      & 0      & b_1                                  \\
                              0      & 1      & \cdots & 0      & 0      & b_2                                  \\
                              0      & 0      & \cdots & 0      & 0      & b_3                                  \\
                              \vdots & \vdots & \ddots & \vdots & \vdots & \vdots                               \\
                              0      & 0      & \cdots & 1      & 1      & b_{n-1}                              \\
                              0      & 0      & \cdots & 0      & 0      & \displaystyle\sum_{i=1}^{n}(-1)^ib_i
                          \end{pmatrix}. \]
                      故当$\displaystyle\sum_{i=1}^{n}(-1)^ib_i=0$时，方程组有无穷多解，一般解为
                      \[\begin{cases} \begin{aligned}
                                  x_1     & =\displaystyle\sum_{i=1}^{n-1}(-1)^{i-1}b_i+(-1)^1x_n, \\
                                  x_2     & =\displaystyle\sum_{i=2}^{n-1}(-1)^{i-2}b_i+(-1)^2x_n, \\
                                          & \vdotswithin{ = }                                      \\
                                  x_{n-2} & =b_{n-2}-b_{n-1}+(-1)^{n-2}x_n,                        \\
                                  x_{n-1} & =b_{n-1}+(-1)^{n-1}x_n,                                \\
                              \end{aligned} \end{cases}\]
                      其中$x_n$为自由未知量.

                \item 当$n$为奇数时，有
                      \[ \bar{A}\rightarrow\begin{pmatrix}
                              1      & 1      & \cdots & 0      & 0      & b_1                                        \\
                              0      & 1      & \cdots & 0      & 0      & b_2                                        \\
                              0      & 0      & \cdots & 0      & 0      & b_3                                        \\
                              \vdots & \vdots & \ddots & \vdots & \vdots & \vdots                                     \\
                              0      & 0      & \cdots & 1      & 1      & b_{n-1}                                    \\
                              0      & 0      & \cdots & 0      & 2      & b_n+\displaystyle\sum_{i=1}^{n-1}(-1)^ib_i
                          \end{pmatrix}. \]
                      此时无论$b_1,b_2,\ldots,b_n\enspace(n\geqslant 2)$取何值，方程组都有唯一解为
                      \[\begin{cases} \begin{aligned}
                                  x_1     & =\displaystyle\sum_{i=1}^{n-1}(-1)^{i-1}b_i+(-1)^{n-1}x_n, \\
                                  x_2     & =\displaystyle\sum_{i=2}^{n-1}(-1)^{i-2}b_i+(-1)^{n-2}x_n, \\
                                          & \vdotswithin{ = }                                          \\
                                  x_{n-2} & =b_{n-2}-b_{n-1}+(-1)^2x_n,                                \\
                                  x_{n-1} & =b_{n-1}+(-1)^1x_n,                                        \\
                                  x_n     & =\frac 12b_n+\displaystyle\sum_{i=1}^{n-1}(-1)^ib_i.
                              \end{aligned} \end{cases}\]
            \end{enumerate}
\end{solve}

\begin{exercise}[补充训练 L11-C-5]\label{prob:L11-C-5}
已知$A$是$s \times n$矩阵，$B$是$m \times n$矩阵，$X,\vec{a},\vec{b}$分别是$n,s,m$元列向量，证明：
        \begin{enumerate}
            \item 齐次线性方程组$AX=\vec{0}$和$BX=\vec{0}$同解的充要条件是$A$与$B$行向量等价（列向量不一定）；

            \item 齐次线性方程组$AX=\vec{0}$和$BX=\vec{0}$解空间分别为$V_1,V_2$，证明：$V_1 \subseteq V_2$的充要条件是存在$m \times s$矩阵$C$使得$B=CA$；

            \item 线性方程组$AX=\vec{a}$的解都是$BX=\vec{b}$的解的充要条件是增广矩阵$(B,\vec{b})$的每个行向量都可以被$(A,\vec{a})$的行向量线性表示；

            \item 线性方程组$AX=\vec{a}$与$BX=\vec{b}$同解的充要条件是$(A,\vec{a})$与$(B,\vec{b})$行向量等价.
        \end{enumerate}
\end{exercise}
\noindent 来源：\href{https://github.com/yhwu-is/Linear-Algebra-Left-Undone/blob/new/%E8%AE%B2%E4%B9%89/%E4%B8%93%E9%A2%98/11%20%E6%9C%9D%E8%8A%B1%E5%A4%95%E6%8B%BE.tex}{L11-C-5}。
\begin{solve}
\begin{enumerate}
                \item 必要性. 设$A$的行向量为$\alpha_1,\alpha_2,\ldots,\alpha_s$，$B$的行向量为$\beta_1,\beta_2,\ldots,\beta_m$，由$AX=0$与$BX=0$同解得$AX=0$与$\begin{pmatrix}A \\ B\end{pmatrix}=0$同解，从而系数矩阵的秩相同，即行向量$\alpha_1,\alpha_2,\ldots,\alpha_s$与$\alpha_1,\alpha_2,\ldots,\alpha_s,\beta_1,\beta_2,\ldots,\beta_m$秩相同. 

                      充分性可由必要性的证明得到. 但是由于$A,B$列向量的维数都可能不同，所以不存在等价关系.

                \item 提示. $V_1\subseteq V_2$等价于$AX=0$与$\begin{pmatrix}
                              A \\
                              B
                          \end{pmatrix}X=0$同解.

                \item 必要性. 由题意可知$AX=a$与$\begin{pmatrix} A \\ B \end{pmatrix}X=\begin{pmatrix} a \\ b \end{pmatrix}$同解，从而它们的增广矩阵秩相同，即$r(A,a)=r\begin{pmatrix}
                              A & a \\
                              B & b
                          \end{pmatrix}$，可知$(A,a)$与$\begin{pmatrix}
                              A & a \\
                              B & b
                          \end{pmatrix}$的行向量组等价，从而$(B,b)$的每一个行向量都可以由$(A,a)$的行向量线性表出.\\
                      充分性可由必要性得到.

                \item 证明留给读者.
            \end{enumerate}
\end{solve}

\begin{exercise}[补充训练 L11-C-6]\label{prob:L11-C-6}
设$A=(a_{ij})_{n \times n}$是一个$n$级矩阵，证明：
        \begin{enumerate}
            \item 若$A$为复矩阵，且$|a_{ii}|>\displaystyle\sum_{j \neq i}|a_{ij}|$，那么$|A|\neq 0$；

            \item 若$A$为实矩阵，且$a_{ii}>\displaystyle\sum_{j \neq i}|a_{ij}|$，那么$|A|>0$；

            \item （推论）存在充分大的实数$M$，使得$t>M$时，$tE+A$可逆.
        \end{enumerate}
\end{exercise}
\noindent 来源：\href{https://github.com/yhwu-is/Linear-Algebra-Left-Undone/blob/new/%E8%AE%B2%E4%B9%89/%E4%B8%93%E9%A2%98/11%20%E6%9C%9D%E8%8A%B1%E5%A4%95%E6%8B%BE.tex}{L11-C-6}。
\begin{solve}
\begin{enumerate}
                \item \label{item:15:B:5:1}
                      采用反证法. 设 $\lvert A \rvert = 0$，则线性方程组 $AX = 0$ 有非零解，设为 $X_0 = (x_1, x_2, \ldots, x_n)^{\mathrm{T}}$，记
                      \[\lvert x_k \rvert = \max \{\lvert x_1 \rvert, \lvert x_2 \rvert, \ldots, \lvert x_n \rvert\}.\]
                      由 $X_0 \neq 0$ 可知 $\lvert x_k \rvert > 0$，考虑 $AX = 0$ 的第 $k$ 个方程，有 $\displaystyle\sum_{j=1}^n a_{kj}x_j = 0$，于是
                      \[\lvert a_{kk} \rvert \lvert x_k \rvert = \lvert -\displaystyle\sum_{j \neq k}a_{kj}x_j \rvert \leqslant \sum_{j \neq k}\lvert a_{kj} \rvert \lvert x_k \rvert.\]
                      约去 $\lvert x_k \rvert$ 后可得 $\lvert a_{kk} \rvert \leqslant \displaystyle\sum_{j \neq k} \lvert a_{kj} \rvert$，这与条件矛盾. 所以 $\lvert A \rvert \neq 0$.

                \item 构造实函数
                      \[f(t) = \begin{vmatrix}
                              a_{11}  & ta_{12} & ta_{13} & \cdots & ta_{1n} \\
                              ta_{21} & a_{22}  & ta_{23} & \cdots & ta_{2n} \\
                              ta_{31} & ta_{32} & a_{33}  & \cdots & ta_{3n} \\
                              \vdots  & \vdots  & \vdots  & \ddots & \vdots  \\
                              ta_{n1} & ta_{n2} & ta_{n3} & \cdots & a_{nn}  \\
                          \end{vmatrix}\]
                      由于 $A$ 是实矩阵，所以 $f(t)$ 是关于 $t$ 的一个实系数多项式（连续）函数，同时
                      \[f(0) = a_{11}a_{22}\cdots a_{nn} > 0.\]
                      当$t \in [0, 1]$ 时，还有
                      \[a_{ii} > \sum_{j \neq i} \lvert a_{ij} \rvert \geqslant \sum_{j \neq i} \lvert ta_{ij} \rvert.\]
                      由 \eqref{item:15:B:5:1} 可知 $f(t)$ 在 $[0, 1]$ 上非零，由连续函数的介值定理可知 $f(1) > 0$，即 $\lvert A \rvert > 0$.

                \item 此为直接推论不再赘述.
            \end{enumerate}
\end{solve}

\begin{exercise}[补充训练 L11-例-1]\label{prob:L11-例-1}
设$n$阶矩阵$A$的行列式$|A|\neq 0$，记$A$的前$n-1$列形成的矩阵为$A_1$，$A$的第$n$列为$\vec{b}$，问：线性方程组$A_1X=\vec{b}$是否有解？
\end{exercise}
\noindent 来源：\href{https://github.com/yhwu-is/Linear-Algebra-Left-Undone/blob/new/%E8%AE%B2%E4%B9%89/%E4%B8%93%E9%A2%98/11%20%E6%9C%9D%E8%8A%B1%E5%A4%95%E6%8B%BE.tex}{L11-例-1}。
\begin{solve}
无解；$|A|\neq 0$可知$r(A)=n$，$r(A_1)=n-1$，因此系数矩阵$A_1$的秩小于增广矩阵$A$的秩$n$，故无解.
\end{solve}

\begin{exercise}[补充训练 L11-例-2]\label{prob:L11-例-2}
求解方程组$\begin{cases}
            x_1+x_2+x_3=1          \\
            a_1x_1+a_2x_2+a_3x_3=0 \\
            a_1^2x_1+a_2^2x_2+a_3^2x_3=0
        \end{cases}$，其中$a_1,a_2,a_3$两两不等.
\end{exercise}
\noindent 来源：\href{https://github.com/yhwu-is/Linear-Algebra-Left-Undone/blob/new/%E8%AE%B2%E4%B9%89/%E4%B8%93%E9%A2%98/11%20%E6%9C%9D%E8%8A%B1%E5%A4%95%E6%8B%BE.tex}{L11-例-2}。
\begin{solve}
$a_1,a_2,a_3$两两不等时，我们有
    \[D=\begin{vmatrix}
            1     & 1     & 1     \\
            a_1   & a_2   & a_3   \\
            a_1^2 & a_2^2 & a_3^2
        \end{vmatrix}=(a_2-a_1)(a_3-a_1)(a_3-a_2)\neq 0,\]
    根据Cramer法则，方程组有唯一解，且
    \[D_1=\begin{vmatrix}
            1 & 1     & 1     \\
            0 & a_2   & a_3   \\
            0 & a_2^2 & a_3^2
        \end{vmatrix}=(a_3-a_2)a_2a_3,\]
    \[D_2=\begin{vmatrix}
            1     & 1 & 1     \\
            a_1   & 0 & a_3   \\
            a_1^2 & 0 & a_3^2
        \end{vmatrix}=(a_1-a_3)a_1a_3,\]
    \[D_3=\begin{vmatrix}
            1     & 1     & 1 \\
            a_1   & a_2   & 0 \\
            a_1^2 & a_2^2 & 0
        \end{vmatrix}=(a_2-a_1)a_1a_2,\]
    因此
    \[x_1=\dfrac{D_1}{D}=\dfrac{a_2a_3}{(a_2-a_1)(a_3-a_1)},\]\[x_2=\dfrac{D_2}{D}=\dfrac{a_1a_3}{(a_1-a_2)(a_3-a_2)},\]\[x_3=\dfrac{D_3}{D}=\dfrac{a_1a_2}{(a_1-a_3)(a_2-a_3)}.\]
\end{solve}

\begin{exercise}[补充训练 L11-例-9]\label{prob:L11-例-9}
当$k$取何值时，方程组：
    \[\begin{cases}
            x_1+x_2+kx_3=4 \\ -x_1+kx_2+x_3=k^2 \\ x_1-x_2+2x_3=-4
        \end{cases}\]
    有唯一解、无解、有无穷多解？在有解的情况下，求出方程组的全部解.
\end{exercise}
\noindent 来源：\href{https://github.com/yhwu-is/Linear-Algebra-Left-Undone/blob/new/%E8%AE%B2%E4%B9%89/%E4%B8%93%E9%A2%98/11%20%E6%9C%9D%E8%8A%B1%E5%A4%95%E6%8B%BE.tex}{L11-例-9}。
\begin{solve}
对于有无解的区别，我们一般都考虑直接使用高斯-若当消元法. 由高斯-若当消元法有（省略中间步骤直接得到阶梯矩阵）：
    \[\begin{pmatrix}
            1  & 1  & k & 4   \\
            -1 & k  & 1 & k^2 \\
            1  & -1 & 2 & -4
        \end{pmatrix}\to\begin{pmatrix}
            1 & 1 & k                     & 4      \\
            0 & 2 & k-2                   & 8      \\
            0 & 0 & \dfrac{(k+1)(4-k)}{2} & k(k-4)
        \end{pmatrix}.\]
    \begin{enumerate}
        \item 当$k=-1$时，增广矩阵秩为3，系数矩阵秩为2（或者说最后一行出现矛盾方程），无解；

        \item 当$k=4$时，增广矩阵和系数矩阵秩均为2，方程有无穷多解，解得同解为$k_1(-3,-1,1)^{\mathrm{T}}+(0,4,0)^{\mathrm{T}}(k_1\in\mathbf{R})$；

        \item 当$k\neq-1,4$时，增广矩阵和系数矩阵秩均为3，方程有唯一解，解为
              \[(\dfrac{k^2+2k}{k+1},\dfrac{k^2+2k+4}{k+1},-\dfrac{2k}{k+1})^{\mathrm{T}}.\]
    \end{enumerate}

    事实上，本题方程个数与未知数个数相等，因此可以运用 原讲义“Cramer”处解决. 首先求解系数矩阵$A$的行列式为
    \[|A|=\begin{vmatrix}
            1  & 1  & k \\
            -1 & k  & 1 \\
            1  & -1 & 2
        \end{vmatrix}=(k+1)(4-k),\]
    由Cramer法则，当$|A|\neq 0$时，方程组有唯一解，当$|A|=0$时，方程组无解或有无穷多解，其余关于有无解的讨论与上面一致.
\end{solve}

\begin{exercise}[补充训练 L11-例-10]\label{prob:L11-例-10}
已知方程组\begin{gather*}
        \begin{cases}
            x_1+2x_2+3x_3=0  \\
            2x_1+3x_2+5x_3=0 \\
            x_1+x_2+ax_3=0
        \end{cases} \\
        \begin{cases}
            x_1+bx_2+cx_3=0 \\
            2x_1+b^2x_2+(c+1)x_3=0
        \end{cases}
    \end{gather*}
    同解，求$a,b,c$的值.
\end{exercise}
\noindent 来源：\href{https://github.com/yhwu-is/Linear-Algebra-Left-Undone/blob/new/%E8%AE%B2%E4%B9%89/%E4%B8%93%E9%A2%98/11%20%E6%9C%9D%E8%8A%B1%E5%A4%95%E6%8B%BE.tex}{L11-例-10}。
\begin{solve}
设第一个方程系数矩阵为$A$，第二个方程系数矩阵为$B$，则
    \[\begin{pmatrix}
            A \\ B
        \end{pmatrix}=\begin{pmatrix}
            1 & 2 & 3 \\ 2 & 3 & 5 \\ 1 & 1 & a \\ 1 & b & c \\ 2 & b^2 & c+1
        \end{pmatrix}\to\begin{pmatrix}
            1 & 0 & 1 \\ 0 & 1 & 1 \\ 0 & 0 & a-2 \\ 0 & 0 & c-b-1 \\ 0 & 0 & c-b^2-1
        \end{pmatrix},\]
    因此由同解条件$r\begin{pmatrix}
            A \\ B
        \end{pmatrix}=r(A)=r(B)$可知必有$r\begin{pmatrix}
            A \\ B
        \end{pmatrix}=r(A)=r(B)=2$（因为$r(B)\leqslant 2$，$r\begin{pmatrix}
            A \\ B
        \end{pmatrix}\geqslant 2$）.

    从而有$a-2=c-b-1=c-b^2-1=0$. 因此$a=2,b=0,c=1$或$a=2,b=1,c=2$. 当$a=2,b=0,c=1$时，$r(B)=1\neq r(A)=2$，舍去. 故$a=2,b=1,c=2$.
\end{solve}

\begin{exercise}[补充训练 L11-例-11]\label{prob:L11-例-11}
设四元齐次线性方程组(1) 为\[\begin{cases}
            2x_1+3x_2-x_3=0 \\ x_1+2x_2+x_3-x_4=0
        \end{cases}\]已知另一个四元齐次线性方程组 (2) 的基础解系为
    \[\alpha_1=(2,-1,a+2,1)^\mathrm{T},\enspace\alpha_2=(-1,2,4,a+8)^\mathrm{T}\]
    \begin{enumerate}
        \item 求方程组 (1) 的一个基础解系；

        \item 当$a$为何值时，方程组 (1) 和 (2) 有非零公共解，并求出非零公共解.
    \end{enumerate}
\end{exercise}
\noindent 来源：\href{https://github.com/yhwu-is/Linear-Algebra-Left-Undone/blob/new/%E8%AE%B2%E4%B9%89/%E4%B8%93%E9%A2%98/11%20%E6%9C%9D%E8%8A%B1%E5%A4%95%E6%8B%BE.tex}{L11-例-11}。
\begin{solve}
\begin{enumerate}
        \item 直接给出结论：方程组 (1) 的一个基础解系为
              \[\beta_1=(5,-3,1,0)^{\mathrm{T}},\enspace\beta_2=(-3,2,0,1)^{\mathrm{T}}.\]

        \item 根据前述定理，两方程有非零公共解当且仅当$\alpha_1,\alpha_2,\beta_1,\beta_2$线性相关，可以将这$\beta_1,\beta_2,\alpha_1,\alpha_2$按列排成矩阵
              \[A = \begin{pmatrix}
                      5 & -3 & 2 & -1 \\ -3 & 2 & -1 & 2 \\ 1 & 0 & a+2 & 4 \\ 0 & 1 & 1 & a+8
                  \end{pmatrix},\]
              利用求解极大线性无关组的思想，对 $A$ 进行初等行变换得到阶梯矩阵
              \[\begin{pmatrix}
                      1 & 0 & a+2 & 4 \\ 0 & 1 & 1 & a+8 \\ 0 & 0 & 3a+3 & -2a-2 \\ 0 & 0 & 0 & -\frac{1}{3}a-\frac{1}{3}
                  \end{pmatrix},\]
              使得 $\alpha_1,\alpha_2,\beta_1,\beta_2$ 线性相关，则必有 $r(A) < 4$，即 $|A| = 0$，不难解得 $a = -1$. 此时阶梯矩阵为
              \[\begin{pmatrix}
                      1 & 0 & 1 & 4 \\ 0 & 1 & 1 & 7 \\ 0 & 0 & 0 & 0 \\ 0 & 0 & 0 & 0
                  \end{pmatrix},\]
              不难看出 $\beta_1,\beta_2$ 是向量组 $\beta_1,\beta_2,\alpha_1,\alpha_2$ 的极大线性无关组，因此 $\alpha_1,\alpha_2$ 可以完全由 $\beta_1,\beta_2$ 线性表示，故 $W_1 \cap W_2 = \operatorname{span}(\beta_1,\beta_2)$. 而求两方程组公共解就是求两子空间 $W_1$ 和 $W_2$ 的交集，故公共解为 $\beta_1,\beta_2$ 的线性组合，即
              \[k_1\beta_1 + k_2\beta_2 = k_1(5,-3,1,0)^{\mathrm{T}} + k_2(-3,2,0,1)^{\mathrm{T}}.\]
              其中 $k_1,k_2$ 为任意常数.
    \end{enumerate}
\end{solve}

\begin{exercise}[补充训练 L11-例-12]\label{prob:L11-例-12}
已知$n$维列向量组$\alpha_1,\ldots,\alpha_s$线性无关，求一齐次线性方程组以$\alpha_1,\ldots,\alpha_s$为基础解系.
\end{exercise}
\noindent 来源：\href{https://github.com/yhwu-is/Linear-Algebra-Left-Undone/blob/new/%E8%AE%B2%E4%B9%89/%E4%B8%93%E9%A2%98/11%20%E6%9C%9D%E8%8A%B1%E5%A4%95%E6%8B%BE.tex}{L11-例-12}。
\begin{solve}
设所求的齐次线性方程组为$AX=0$，令$B=(\alpha_1,\ldots,\alpha_s)$，则$B$为$n\times s$矩阵且$AB=O$，于是$B^\mathrm{T}A^\mathrm{T}=O$，解线性方程组$B^\mathrm{T}X=0$，得到其基础解系为
    \[\beta_1,\beta_2,\ldots,\beta_{n-s},\]
    令$A^\mathrm{T}=(\beta_1,\beta_2,\ldots,\beta_{n-s})$即可，因为此时$B^\mathrm{T}A^\mathrm{T}=O$，故$AB=O$.
\end{solve}

\begin{exercise}[补充训练 L11-例-13]\label{prob:L11-例-13}
设向量组$\alpha_1,\ldots,\alpha_s$线性无关，求一非齐次线性方程组$AX=\vec{b}$，使其解集以$\alpha_1,\ldots,\alpha_s$为极大线性无关组.
\end{exercise}
\noindent 来源：\href{https://github.com/yhwu-is/Linear-Algebra-Left-Undone/blob/new/%E8%AE%B2%E4%B9%89/%E4%B8%93%E9%A2%98/11%20%E6%9C%9D%E8%8A%B1%E5%A4%95%E6%8B%BE.tex}{L11-例-13}。
\begin{solve}
根据原讲义“非齐次线性方程组公共解”处，$\alpha_2-\alpha_1,\ldots,\alpha_s-\alpha_1$是$AX=\vec{0}$的基础解系，因此根据齐次线性方程组反问题的解法可以得到$A$，然后令$b=A\alpha_1$即可符合题意.
\end{solve}

\begin{exercise}[补充训练 L11-例-14]\label{prob:L11-例-14}
已知$\alpha_1=(1,2,-1,0,4)^\mathrm{T},\enspace\alpha_2=(-1,3,2,4,1)^\mathrm{T},\enspace\alpha_3=(2,9,-1,4,13)^\mathrm{T}$，且有$W=\operatorname{span}(\alpha_1,\alpha_2,\alpha_3)$.
    \begin{enumerate}
        \item 求以$W$为解空间的一个齐次线性方程组；

        \item 求以$W'=\{\eta+\alpha \mid \alpha\in W\}$为解集的一个非齐次线性方程组，其中$\eta=(1,2,1,2,1)^\mathrm{T}$.
    \end{enumerate}
\end{exercise}
\noindent 来源：\href{https://github.com/yhwu-is/Linear-Algebra-Left-Undone/blob/new/%E8%AE%B2%E4%B9%89/%E4%B8%93%E9%A2%98/11%20%E6%9C%9D%E8%8A%B1%E5%A4%95%E6%8B%BE.tex}{L11-例-14}。
\begin{solve}
首先我们通过求解极大线性无关组的方法可以得到，$\alpha_1,\alpha_2,\alpha_3$的极大线性无关组是$\alpha_1,\alpha_2$. 然后我们基于此根据前面两个例题介绍的方法，我们有如下求解过程：
    \begin{enumerate}
        \item 设所求的齐次线性方程组为$AX=0$，令$B=(\alpha_1,\alpha_2)$，解线性方程组$B^\mathrm{T}X=0$，得到其基础解系为
              \[\beta_1=(7,-1,5,0,0)^\mathrm{T},\enspace\beta_2=(8,-4,0,5,0)^\mathrm{T},\enspace\beta_3=(-2,-1,0,0,1)^\mathrm{T},\]
              令$A^\mathrm{T}=(\beta_1,\beta_2,\beta_3)$即可，即
              \[A=\begin{pmatrix}
                      7 & -1 & 5 & 0 & 0 \\ 8 & -4 & 0 & 5 & 0 \\ -2 & -1 & 0 & 0 & 1
                  \end{pmatrix}.\]
              故所求的线性方程组为$\begin{cases}
                      7x_1-x_2+5x_3=0 \\ 8x_1-4x_2+5x_4=0 \\ -2x_1-x_2+x_5=0
                  \end{cases}$.

        \item 这一问可以利用前面例子中给出的方法，具体步骤不在此赘述（注意因为$\eta$不在$W$中，因此$W'$中可以有三个线性无关向量，不可想当然认为只有$\alpha_2-\alpha_1$是$AX=0$（设$A$为本题要求的方程组的系数矩阵）的解）. 我们这里给出更简单的方法，读者只需回顾根据原讲义“通解加特解”处 即可发现奥秘. 事实上，将$x_1=1,x_2=2,x_3=1,x_4=2,x_5=1$代入有$7x_1-x_2+5x_3=10$，$8x_1-4x_2+5x_4=10$，$-2x_1-x_2+x_5=-3$，因此$\eta$就是方程组$\begin{cases}
                      7x_1-x_2+5x_3=10 \\ 8x_1-4x_2+5x_4=10 \\ -2x_1-x_2+x_5=-3
                  \end{cases}$的一个特解，从而根据原讲义“通解加特解”处 可以知道上述方程组符合题意.
    \end{enumerate}
\end{solve}

\begin{exercise}[补充训练 S-SYS-1]\label{prob:S-SYS-1}
设 $A=\begin{pmatrix}1&-6&0\\0&2&0\\1&0&3\end{pmatrix}$，$b=\begin{pmatrix}-7\\2\\8\end{pmatrix}$。求 $\det A$、$A^{-1}$，解 $Ax=b$，并利用消元结果说明解的唯一性。
\end{exercise}
\noindent 编者补充：S-SYS-1；相关主题见原讲义对应讲，题目数据与解答由本选编另行构造。
\begin{solve}
矩阵的行列式为 $\det A=6\ne0$。对增广矩阵 $(A\mid I_3)$ 作行初等变换，得到
\[(A\mid I_3)\sim(I_3\mid A^{-1}),\qquad A^{-1}=\begin{pmatrix}1&3&0\\0&\frac{1}{2}&0\\-\frac{1}{3}&-1&\frac{1}{3}\end{pmatrix}.\]
因此 $x=A^{-1}b=\begin{pmatrix}-1\\1\\3\end{pmatrix}$。也可以直接化简含常数列的增广矩阵：
\[(A\mid b)\sim \begin{pmatrix}1&0&0&-1\\0&1&0&1\\0&0&1&3\end{pmatrix}.\]
三个未知量所在的列都有主元，没有自由变量，所以此解唯一。检验时将所求向量代回 $Ax$，得到的三维列向量恰为题给 $b$；求逆过程中的等式可通过 $AA^{-1}=I_3$ 复核。
\end{solve}

\begin{exercise}[补充训练 S-NULL-1]\label{prob:S-NULL-1}
给定 $A=\begin{pmatrix}1&-3&-2&-4&-1\\0&1&1&1&1\\1&0&1&-1&2\end{pmatrix}$ 和 $b=\begin{pmatrix}-2\\1\\1\end{pmatrix}$。求 $r(A)$、$\ker A$ 的一组基及 $Ax=b$ 的全部解；再说明其解集为什么通常不是线性子空间。
\end{exercise}
\noindent 编者补充：S-NULL-1；相关主题见原讲义对应讲，题目数据与解答由本选编另行构造。
\begin{solve}
行化简得到
\[(A\mid b)\sim \begin{pmatrix}1&0&1&-1&2&1\\0&1&1&1&1&1\\0&0&0&0&0&0\end{pmatrix}.\]
前两列是主元列，所以 $r(A)=2$，核的维数为 $5-2=3$。把 $x_3,x_4,x_5$ 分别作为自由变量，可得核的一组基为下列矩阵 $N$ 的三列：
\[N=\begin{pmatrix}-1&1&-2\\-1&-1&-1\\1&0&0\\0&1&0\\0&0&1\end{pmatrix}.\]
取 $x_0=\begin{pmatrix}1\\1\\0\\0\\0\end{pmatrix}$，所有解可统一写成 $x=x_0+Nt$，其中 $t\in\mathbf R^3$ 任意。直接相乘得到 $AN=O$ 和 $Ax_0=b$，验证了这一表示；反过来，两特解之差在 $\ker A$ 中，说明没有遗漏任何解。由于此处 $b\ne0$，零向量不是解，故解集是三维仿射子空间而不是线性子空间。
\end{solve}

\section{补充训练：一元多项式}

本节包含 13 道训练题，建议在完成相关专题真题后选择性练习。

\begin{exercise}[补充训练 L13-A-1]\label{prob:L13-A-1}
设 $\mathbf F$ 是域，证明：$\mathbf F[x]$ 在通常的多项式加法和乘法下构成含幺交换环。
\end{exercise}
\noindent 来源：\href{https://github.com/yhwu-is/Linear-Algebra-Left-Undone/blob/new/%E8%AE%B2%E4%B9%89/%E4%B8%93%E9%A2%98/13%20%E5%A4%9A%E9%A1%B9%E5%BC%8F.tex}{L13-A-1}。
\begin{solve}
将多项式看成仅有有限个非零项的系数序列，规定
\[(f+g)_k=a_k+b_k,\qquad(fg)_k=\sum_{i+j=k}a_ib_j,\]
其中 $f=\sum_ia_ix^i$、$g=\sum_jb_jx^j$，未出现的系数取零。加法逐项继承域中的交换律、结合律；零多项式为加法单位元，$-f=\sum_i(-a_i)x^i$ 为加法逆元。有限支撑保证两种运算的结果仍是多项式。
乘法交换律由交换 $i,j$ 得到；若 $h=\sum_kc_kx^k$，则 $(fg)h$ 和 $f(gh)$ 的第 $r$ 个系数都为 $\sum_{i+j+k=r}a_ib_jc_k$，所以乘法结合。各系数的分配律给出 $f(g+h)=fg+fh$，另一侧同理。常数多项式 $1$ 为乘法单位元，故全部环公理成立。
\end{solve}

\begin{exercise}[补充训练 L13-A-2]\label{prob:L13-A-2}
设$p(x),q(x),r(x)\in\mathbf{F}[x]$，则
        \begin{enumerate}
            \item 若$p(x),q(x)\neq 0$，则$p(x)q(x)\neq 0$；
            \item 若$p(x)\neq 0$，且$p(x)q(x)=p(x)r(x)$，则$q(x)=r(x)$.
        \end{enumerate}
\end{exercise}
\noindent 来源：\href{https://github.com/yhwu-is/Linear-Algebra-Left-Undone/blob/new/%E8%AE%B2%E4%B9%89/%E4%B8%93%E9%A2%98/13%20%E5%A4%9A%E9%A1%B9%E5%BC%8F.tex}{L13-A-2}。
\begin{solve}
\begin{enumerate}
\item 若 $p,q$ 非零，设其次数为 $m,n$，首项系数为 $a_m,b_n\ne0$。乘积中 $x^{m+n}$ 的系数为 $a_mb_n\ne0$，故 $pq\ne0$，并有 $\deg(pq)=m+n$。
\item 由 $pq=pr$ 得 $p(q-r)=0$。若 $q-r\ne0$，第 (1) 项将给出乘积非零，矛盾。因此 $q=r$。
\end{enumerate}
此证明对任意系数域成立，无需通过选取多项式的非零函数值来论证。
\end{solve}

\begin{exercise}[补充训练 L13-A-4]\label{prob:L13-A-4}
证明：每个奇数次的实系数多项式都有实的零点.
\end{exercise}
\noindent 来源：\href{https://github.com/yhwu-is/Linear-Algebra-Left-Undone/blob/new/%E8%AE%B2%E4%B9%89/%E4%B8%93%E9%A2%98/13%20%E5%A4%9A%E9%A1%B9%E5%BC%8F.tex}{L13-A-4}。
\begin{solve}
设 $f(x)=a_nx^n+\cdots+a_0$，其中 $n$ 为奇数且 $a_n\ne0$。必要时用 $-f$ 代替 $f$，可设 $a_n>0$。由
\[\frac{f(x)}{x^n}=a_n+\frac{a_{n-1}}x+\cdots+\frac{a_0}{x^n}\longrightarrow a_n\quad(|x|\to\infty),\]
可取充分大的 $R>0$，使 $f(R)>0$ 而 $f(-R)<0$。多项式在 $[-R,R]$ 上连续，由介值定理存在 $c\in(-R,R)$ 使 $f(c)=0$。
\end{solve}

\begin{exercise}[补充训练 L13-B-1]\label{prob:L13-B-1}
设多项式$f(x)$被$(x-1),(x-2),(x-3)$除后，余式分别为$4,8,16$. 求$f(x)$被$(x-1)(x-2)(x-3)$除后的余式.
\end{exercise}
\noindent 来源：\href{https://github.com/yhwu-is/Linear-Algebra-Left-Undone/blob/new/%E8%AE%B2%E4%B9%89/%E4%B8%93%E9%A2%98/13%20%E5%A4%9A%E9%A1%B9%E5%BC%8F.tex}{L13-B-1}。
\begin{solve}
$1=\dfrac{1}{2}(x-2)(x-3)-\dfrac{1}{2}(x-4)(x-1)$，故$\dfrac{1}{2}(x-2)(x-3) \equiv1\bmod{(x-1)}$.

            同理有$-(x-1)(x-3)\equiv1\bmod{(x-2)}$，$\dfrac{1}{2}(x-1)(x-2)\equiv1\bmod{(x-3)}$.

            从而$f(x)\equiv\dfrac{1}{2}\times 4 \times(x-2)(x-3)-1\times 8\times(x-1)(x-3)+\dfrac{1}{2}\times 16\times(x-1)(x-2)\equiv 2x^2-2x+4\bmod{(x-1)(x-2)(x-3)}$.
\end{solve}

\begin{exercise}[补充训练 L13-B-2]\label{prob:L13-B-2}
设$p,q\in\mathbf{F}[x]$，证明：$p^2 \mid q^2\iff p \mid q$.
\end{exercise}
\noindent 来源：\href{https://github.com/yhwu-is/Linear-Algebra-Left-Undone/blob/new/%E8%AE%B2%E4%B9%89/%E4%B8%93%E9%A2%98/13%20%E5%A4%9A%E9%A1%B9%E5%BC%8F.tex}{L13-B-2}。
\begin{solve}
若 $p\mid q$，写 $q=ph$，则 $q^2=p^2h^2$，故 $p^2\mid q^2$。反过来，若 $q=0$，结论直接成立；若 $q\ne0$ 且 $p^2\mid q^2$，则 $p\ne0$。利用域上多项式的唯一因式分解，对每个首一不可约多项式 $\pi$，记 $\nu_\pi(p)$、$\nu_\pi(q)$ 为其在分解中的指数。整除关系给出 $2\nu_\pi(p)\le2\nu_\pi(q)$，所以 $\nu_\pi(p)\le\nu_\pi(q)$ 对所有 $\pi$ 成立，即 $p\mid q$。若 $p=0$，按 $p\mid q\iff q=ph$ 的定义，两侧都等价于 $q=0$。
\end{solve}

\begin{exercise}[补充训练 L13-B-3]\label{prob:L13-B-3}
证明：$\mathbf{F}[x]$中两个次数大于0的多项式没有公共复根的充要条件是它们互素.
\end{exercise}
\noindent 来源：\href{https://github.com/yhwu-is/Linear-Algebra-Left-Undone/blob/new/%E8%AE%B2%E4%B9%89/%E4%B8%93%E9%A2%98/13%20%E5%A4%9A%E9%A1%B9%E5%BC%8F.tex}{L13-B-3}。
\begin{solve}
\begin{enumerate}
                \item 充分性：假设多项式$p,q$有公共复根$\lambda$，则$(x-\lambda)$为$p$和$q$的公因式，这与$p$和$q$互素矛盾.
                \item 必要性：设$d(x)$是$p(x)$和$q(x)$的公因式，则$\deg d>1$，由代数基本定理知$d(x)$在复数域上有根，这表明$p$和$q$有公共复根，矛盾.从而$p$和$q$互素.
            \end{enumerate}
\end{solve}

\begin{exercise}[补充训练 L13-B-4]\label{prob:L13-B-4}
设 $\mathbf F$ 为数域，$0\ne f\in\mathbf F[x]$，$c\in\mathbf F$，$k\ge1$。证明：$c$ 恰为 $f$ 的 $k$ 重根，当且仅当
\[f(c)=f'(c)=\cdots=f^{(k-1)}(c)=0,\qquad f^{(k)}(c)\ne0.\]
\end{exercise}
\noindent 来源：\href{https://github.com/yhwu-is/Linear-Algebra-Left-Undone/blob/new/%E8%AE%B2%E4%B9%89/%E4%B8%93%E9%A2%98/13%20%E5%A4%9A%E9%A1%B9%E5%BC%8F.tex}{L13-B-4}。
\begin{solve}
在点 $c$ 处作多项式的有限 Taylor 展开：
\[f(x)=\sum_{j=0}^{\deg f}\frac{f^{(j)}(c)}{j!}(x-c)^j.\]
数域的特征为零，所以每个 $j!$ 均可逆。根的重数正是展开中第一个非零项的次数；因此重数恰为 $k$ 与题给导数条件等价。也可将 $f=(x-c)^kg$ 且 $g(c)\ne0$ 代入逐次求导，得到 $f^{(k)}(c)=k!g(c)\ne0$。
\end{solve}

\begin{exercise}[补充训练 L13-B-5]\label{prob:L13-B-5}
设$p(x)$和$q(x)$是$\mathbf{F}[x]$中的两个次数不超过$n$的多项式，若存在$\mathbf{F}$上$n+1$个不同的数$\lambda_0,\lambda_1,\ldots,\lambda_n$使得$p(\lambda_i)=q(\lambda_i),\enspace i=0,1,\ldots,n$，则$p(x)=q(x)$.
\end{exercise}
\noindent 来源：\href{https://github.com/yhwu-is/Linear-Algebra-Left-Undone/blob/new/%E8%AE%B2%E4%B9%89/%E4%B8%93%E9%A2%98/13%20%E5%A4%9A%E9%A1%B9%E5%BC%8F.tex}{L13-B-5}。
\begin{solve}
考虑多项式$f(x)=p(x)-q(x)$，$\deg f\leqslant \max\{\deg p,\deg q\}=n$，而$m\enspace(m\geqslant 0)$次多项式至多有$m$个根，这表明$\deg f=-\infty$，只能有$f(x)\equiv 0$，从而$p(x)=q(x)$.
\end{solve}

\begin{exercise}[补充训练 L13-C-1]\label{prob:L13-C-1}
设 $p,s\in\mathbf F[x]$ 且 $s\ne0$，其中 $\mathbf F$ 为域。证明存在唯一的 $q,r\in\mathbf F[x]$，使 $p=sq+r$ 且 $r=0$ 或 $\deg r<\deg s$。请用余式次数最小的方法证明存在性。
\end{exercise}
\noindent 来源：\href{https://github.com/yhwu-is/Linear-Algebra-Left-Undone/blob/new/%E8%AE%B2%E4%B9%89/%E4%B8%93%E9%A2%98/13%20%E5%A4%9A%E9%A1%B9%E5%BC%8F.tex}{L13-C-1}。
\begin{solve}
集合 $\{p-sq:q\in\mathbf F[x]\}$ 非空，因为取 $q=0$ 即含 $p$。若它含零，存在性成立；否则在所有元素中选次数最小者 $r=p-sq$。设 $\deg r=m\ge d=\deg s$，两多项式的首项系数分别为 $a,b$。则
\[r_1=r-\frac abx^{m-d}s=p-s\left(q+\frac abx^{m-d}\right)\]
仍属于该集合，而首项抵消，故 $r_1=0$ 或 $\deg r_1<m$。前一种情形与集合不含零矛盾，后一种与最小性矛盾。因此 $\deg r<d$。若另有 $p=sq_1+r_1=sq_2+r_2$ 满足要求，则 $s(q_1-q_2)=r_2-r_1$；若 $q_1\ne q_2$，左端次数至少为 $\deg s$，右端次数却更小，矛盾。故 $q_1=q_2,r_1=r_2$。
\end{solve}

\begin{exercise}[补充训练 L13-C-2]\label{prob:L13-C-2}
设 $p,q\in\mathbf F[x]$，$\deg p=n\ge1$、$\deg q=m\ge1$。用有限维线性映射证明：$p,q$ 互素当且仅当存在多项式 $u,v$ 使 $up+vq=1$。证明互素时可以选择 $\deg u<m$、$\deg v<n$。
\end{exercise}
\noindent 来源：\href{https://github.com/yhwu-is/Linear-Algebra-Left-Undone/blob/new/%E8%AE%B2%E4%B9%89/%E4%B8%93%E9%A2%98/13%20%E5%A4%9A%E9%A1%B9%E5%BC%8F.tex}{L13-C-2}。
\begin{solve}
若互素，考察线性映射
\[T:\mathbf F[x]_m\times\mathbf F[x]_n\longrightarrow\mathbf F[x]_{m+n},\qquad T(u,v)=up+vq.\]
设 $up+vq=0$，则 $p\mid vq$；由互素性和唯一因子分解可知 $p\mid v$。但 $\deg v<n=\deg p$，故 $v=0$，进而 $u=0$。所以 $T$ 单射。定义域与陪域维数都为 $m+n$，故它也是满射；特别地常值多项式 $1$ 有原像，从而得到所需次数界。反之，若 $up+vq=1$，任何公共因子都整除 $1$，所以 $p,q$ 的首一最大公因式为 $1$。
\end{solve}

\begin{exercise}[补充训练 L13-例-2]\label{prob:L13-例-2}
设$p(x)=x^4-x^3-x^2+2x-1$，$q(x)=x^3-2x+1$，求$(p(x),q(x))$，以及$u(x),v(x)$使得$(p(x),q(x))=u(x)p(x)+v(x)q(x)$.
\end{exercise}
\noindent 来源：\href{https://github.com/yhwu-is/Linear-Algebra-Left-Undone/blob/new/%E8%AE%B2%E4%B9%89/%E4%B8%93%E9%A2%98/13%20%E5%A4%9A%E9%A1%B9%E5%BC%8F.tex}{L13-例-2}。
\begin{solve}
我们可以利用带余除法进行计算：
    \begin{align*}
        p(x)   & =q(x)(x-1)+x^2-x=q(x)s_1(x)+r_1(x),   \\
        q(x)   & =r_1(x)(x+1)-x+1=r_1(x)s_2(x)+r_2(x), \\
        r_1(x) & =r_2(x)(-x)=r_2(x)s_3(x),
    \end{align*}
    因此最大公因式$(p(x),q(x))=r_2(x)=-x+1$. 接下来开始回推，我们有
    \begin{align*}
        r_2(x) & =q(x)-r_1(x)s_2(x)                   \\
               & =q(x)-(p(x)-q(x)s_1(x))s_2(x)        \\
               & =p(x)(-s_2(x))+q(x)(1+s_1(x)s_2(x)),
    \end{align*}
    所以$u(x)=-s_2(x)=-x-1$，$v(x)=1+s_1(x)s_2(x)=x^2$.
\end{solve}

\begin{exercise}[补充训练 S-GCD-1]\label{prob:S-GCD-1}
设 $f(x)=x^{2}-3x+2$，$g(x)=x^{2}-5x+4$。求首一最大公因式、首一最小公倍式，并找出多项式 $u,v$ 使 $uf+vg=\gcd(f,g)$。进一步求 $f+g$ 与 $f-g$ 的首一最大公因式。
\end{exercise}
\noindent 编者补充：S-GCD-1；相关主题见原讲义对应讲，题目数据与解答由本选编另行构造。
\begin{solve}
分解得到 $f=(x-1)(x-2)$、$g=(x-1)(x-4)$。三个根两两不同，所以
\[\gcd(f,g)=x-1,\qquad\operatorname{lcm}(f,g)=x^{3}-7x^{2}+14x-8.\]
相减得 $f-g=(2)(x-1)$，从而可以取常数多项式
\[u=\frac{1}{2},\qquad v=-\frac{1}{2}.\]
这既给出了裴蜀等式，也证明任何公共因子都整除 $x-1$。最后，由 $2f=(f+g)+(f-g)$、$2g=(f+g)-(f-g)$，且标量域为实数域，$2$ 可逆，故两对多项式的公因子相同，最后一问的答案仍为 $x-1$。
\end{solve}

\begin{exercise}[补充训练 S-INTERP-1]\label{prob:S-INTERP-1}
求次数不超过 $2$ 的实多项式 $p$，使 $p(-1)=5$、$p(0)=1$、$p(2)=5$。用待定系数法求解并说明唯一性；若任意多项式 $f$ 具有相同的三个函数值，求它除以 $(x+1)x(x-2)$ 的余式。
\end{exercise}
\noindent 编者补充：S-INTERP-1；相关主题见原讲义对应讲，题目数据与解答由本选编另行构造。
\begin{solve}
写 $p=c_0+c_1x+c_2x^2$，得到范德蒙方程组
\[\begin{pmatrix}1&-1&1\\1&0&0\\1&2&4\end{pmatrix}\begin{pmatrix}c_0\\c_1\\c_2\end{pmatrix}=\begin{pmatrix}5\\1\\5\end{pmatrix}.\]
系数行列式为 $6\ne0$，解为 $(c_0,c_1,c_2)^{\mathrm T}=\begin{pmatrix}1\\-2\\2\end{pmatrix}$，即 $p(x)=2x^{2}-2x+1$。若还有另一多项式满足条件，两者之差次数至多 $2$，却有三个不同零点，只能为零，因此解唯一。作带余除法 $f=q(x)(x+1)x(x-2)+r$，则 $\deg r<3$ 且 $r$ 与 $f$ 在三个插值点的值相同，故余式就是所求 $p$。
\end{solve}

\section{补充训练：特征值与特征向量}

本节包含 24 道训练题，建议在完成相关专题真题后选择性练习。

\begin{exercise}[补充训练 L14-A-1]\label{prob:L14-A-1}
设$S,T\in \mathcal{L}(V)$满足$ST=TS$，证明：$\ker S$和$\operatorname{im} S$都在$T$下不变.
\end{exercise}
\noindent 来源：\href{https://github.com/yhwu-is/Linear-Algebra-Left-Undone/blob/new/%E8%AE%B2%E4%B9%89/%E4%B8%93%E9%A2%98/14%20%E7%9B%B8%E4%BC%BC%E6%A0%87%E5%87%86%E5%BD%A2%EF%BC%9A%E5%8A%A8%E6%9C%BA%E4%B8%8E%E5%9F%BA%E7%A1%80.tex}{L14-A-1}。
\begin{solve}
对于$\forall v\in\ker S,Sv=0\implies STv=TSv=0$，这表明$Tv\in \ker S$，$\ker S$在$T$下不变.
            对于$\forall v \in \operatorname{im} S,v=Su\implies Tv=TSu=S(Tu)$，这表明$Tv \in \operatorname{im} S$，$\operatorname{im} S$在$T$下不变.
\end{solve}

\begin{exercise}[补充训练 L14-A-2]\label{prob:L14-A-2}
已知$\mathbf{R}^2$上线性变换$T$在基$e_1=(1,0),e_2=(0,1)$下的矩阵为$\begin{pmatrix}2 & 1 \\ 0 & 2\end{pmatrix}$. 证明：
        \begin{enumerate}
            \item 设$W_1$为由$e_1$张成的子空间，则$W_1$为$T$的不变子空间；

            \item $\mathbf{R}^2$不能表示为$T$的任何不变子空间$W_2$与$W_1$的直和.
        \end{enumerate}
\end{exercise}
\noindent 来源：\href{https://github.com/yhwu-is/Linear-Algebra-Left-Undone/blob/new/%E8%AE%B2%E4%B9%89/%E4%B8%93%E9%A2%98/14%20%E7%9B%B8%E4%BC%BC%E6%A0%87%E5%87%86%E5%BD%A2%EF%BC%9A%E5%8A%A8%E6%9C%BA%E4%B8%8E%E5%9F%BA%E7%A1%80.tex}{L14-A-2}。
\begin{solve}
\begin{enumerate}
                \item 对于$\forall v\in W_1$，有$v=ke_1,Tv=T(ke_1)=kT(e_1)=2ke_1\in W_1$，从而$W_1$为$T$的不变子空间.
                \item 容易看出$T$在自然基下的矩阵不可对角化，从而$\mathbf{R}^2$不能表示为$T$的任何不变子空间$W_2$与$W_1$的直和.
            \end{enumerate}
\end{solve}

\begin{exercise}[补充训练 L14-A-3]\label{prob:L14-A-3}
定义线性变换$T\in \mathcal{L}(\mathbf{F}^2)$为$T(x,y)=(y,0)$. 令$U=\{(x,0) \mid x\in\mathbf{F}\}$. 证明：
        \begin{enumerate}
            \item $U$在$T$下不变，且$T|_{U}$是$U$上的零线性变换；

            \item 不存在$\mathbf{F}^2$的在$T$下不变的子空间$W$使得$\mathbf{F}^2=U\oplus W$；

            \item $T/U$是$\mathbf{F}^2/U$上的0线性变换.
        \end{enumerate}
\end{exercise}
\noindent 来源：\href{https://github.com/yhwu-is/Linear-Algebra-Left-Undone/blob/new/%E8%AE%B2%E4%B9%89/%E4%B8%93%E9%A2%98/14%20%E7%9B%B8%E4%BC%BC%E6%A0%87%E5%87%86%E5%BD%A2%EF%BC%9A%E5%8A%A8%E6%9C%BA%E4%B8%8E%E5%9F%BA%E7%A1%80.tex}{L14-A-3}。
\begin{solve}
\begin{enumerate}
                \item 对于$\forall u=(x,0)\in U$，$Tu=T(x,0)=(0,0)\in U$，因此$U$在$T$下不变，且$T|_{U}$是$U$上的零线性变换.
                \item $T$在自然基下的矩阵为$\begin{pmatrix}
                    0 & 1 \\ 0 & 0
                \end{pmatrix}$，它不可对角化.
                \item 对于$\forall v+U=(x,y)+U\in \mathbf{F}^2/U$，$(T/U)(v+U)=Tv+U=(y,0)+U=0+U$，从而$T/U$是$\mathbf{F}^2/U$上的$0$线性变换.
            \end{enumerate}
\end{solve}

\begin{exercise}[补充训练 L14-A-4]\label{prob:L14-A-4}
设$V$是有限维的且$T\in \mathcal{L}(V)$，设$\lambda_1,\ldots,\lambda_m$是非零互异特征值，证明：
        \[ \dim E(\lambda_1,T)+\cdots+\dim E(\lambda_m,T)\leqslant\dim\operatorname{im} T. \]
\end{exercise}
\noindent 来源：\href{https://github.com/yhwu-is/Linear-Algebra-Left-Undone/blob/new/%E8%AE%B2%E4%B9%89/%E4%B8%93%E9%A2%98/14%20%E7%9B%B8%E4%BC%BC%E6%A0%87%E5%87%86%E5%BD%A2%EF%BC%9A%E5%8A%A8%E6%9C%BA%E4%B8%8E%E5%9F%BA%E7%A1%80.tex}{L14-A-4}。
\begin{solve}
考虑$U=E(\lambda_1,T)\oplus E(\lambda_2,T) \oplus \cdots \oplus E(\lambda_m,T)$，对于$\forall u\in U$，设$u=v_1+v_2+\cdots+v_m$，其中$v_i\in E(\lambda_i,T)$.则$v_i=T(\dfrac{v_i}{\lambda_i})$，从而$u=T(\sum\limits_{i=1}^{m}\dfrac{v_i}{\lambda_i}) \in \operatorname{im} T$.因此$U$是$\operatorname{im} T$的子空间.进而有$\dim U=\dim E(\lambda_1,T)+\cdots+\dim E(\lambda_m,T) \leqslant \dim \operatorname{im} T$.
\end{solve}

\begin{exercise}[补充训练 L14-A-5]\label{prob:L14-A-5}
设$T\in \mathcal{L}(V)$且$\dim\operatorname{im} T=k$. 证明$T$至多有$k+1$个特征值.
\end{exercise}
\noindent 来源：\href{https://github.com/yhwu-is/Linear-Algebra-Left-Undone/blob/new/%E8%AE%B2%E4%B9%89/%E4%B8%93%E9%A2%98/14%20%E7%9B%B8%E4%BC%BC%E6%A0%87%E5%87%86%E5%BD%A2%EF%BC%9A%E5%8A%A8%E6%9C%BA%E4%B8%8E%E5%9F%BA%E7%A1%80.tex}{L14-A-5}。
\begin{solve}
设$T$的非零特征值为$\lambda_1,\ldots,\lambda_m$，则$m\leqslant \dim E(\lambda_1,T)+\cdots+\dim E(\lambda_m,T)\leqslant \dim\operatorname{im} T=k$.若$0$是$T$的非零特征值，则$T$至多有$k+1$个特征值.否则$T$至多有$k$个特征值.综上可知$T$至多有$k+1$个特征值.
\end{solve}

\begin{exercise}[补充训练 L14-A-6]\label{prob:L14-A-6}
设$\sigma$是线性空间$\mathbf{R}[x]_3$上的线性变换，它在基$1,x,x^2$下的矩阵为
        \[ A=\begin{pmatrix}
                1 & 2 & 2 \\ 2 & 1 & 2 \\ 2 & 2 & 1
            \end{pmatrix}\]
        求$\sigma$的特征值与特征子空间.
\end{exercise}
\noindent 来源：\href{https://github.com/yhwu-is/Linear-Algebra-Left-Undone/blob/new/%E8%AE%B2%E4%B9%89/%E4%B8%93%E9%A2%98/14%20%E7%9B%B8%E4%BC%BC%E6%A0%87%E5%87%86%E5%BD%A2%EF%BC%9A%E5%8A%A8%E6%9C%BA%E4%B8%8E%E5%9F%BA%E7%A1%80.tex}{L14-A-6}。
\begin{solve}
令$|\lambda E-A|=0$，解得$\lambda_1=5,\lambda_2=-1$.然后解$(5E-A)X=0$，解得$X=k(1,1,1)^\mathrm{T}$.解$(-E-A)X=0$，解得$X=k_1(1,-1,0)^\mathrm{T}+k_2(1,0,-1)^\mathrm{T}$.因此$\lambda_1=5$对应的特征子空间为$\operatorname{span}(1+x+x^2)$，$\lambda_2=-1$对应的特征子空间为$\operatorname{span}(1-x,1-x^2)$.
\end{solve}

\begin{exercise}[补充训练 L14-A-7]\label{prob:L14-A-7}
设$A,P$都是3阶方阵，$P$可逆，已知$A$的特征值$\lambda_1=1,\lambda_2=-1,\lambda_3=2$，$B=A^3-5A^2$，求$|B|$，$|A+5E|$，$|5E+P^{-1}AP|$.
\end{exercise}
\noindent 来源：\href{https://github.com/yhwu-is/Linear-Algebra-Left-Undone/blob/new/%E8%AE%B2%E4%B9%89/%E4%B8%93%E9%A2%98/14%20%E7%9B%B8%E4%BC%BC%E6%A0%87%E5%87%86%E5%BD%A2%EF%BC%9A%E5%8A%A8%E6%9C%BA%E4%B8%8E%E5%9F%BA%E7%A1%80.tex}{L14-A-7}。
\begin{solve}
矩阵的行列式等于其特征值之积，知$|A|=\lambda_1\lambda_2\lambda_3=-2$.
            进一步地，$A-5E$的特征值为$-4,-6,-3$，$A+5E$和$5E+P^{-1}AP$的特征值均为$6,4,7$.从而$|A-5E|=-72,|A+5E|=|5E+P^{-1}AP|=168$.
            $|B|=|A^2(A-5E)|=|A||A||A-5E|=-288$.
\end{solve}

\begin{exercise}[补充训练 L14-A-8]\label{prob:L14-A-8}
设$A=\begin{pmatrix}
                a & -1 & c \\ 5 & b & 3 \\ 1-c & 0 & -a
            \end{pmatrix}$，$|A|=-1$，$\alpha=(-1,-1,1)^\mathrm{T}$为$\operatorname{adj}(A)$的特征向量，求$\operatorname{adj}(A)$对应于$\alpha$的特征值及$a,b,c$和$A$对应于$\alpha$的特征值$\mu$.
\end{exercise}
\noindent 来源：\href{https://github.com/yhwu-is/Linear-Algebra-Left-Undone/blob/new/%E8%AE%B2%E4%B9%89/%E4%B8%93%E9%A2%98/14%20%E7%9B%B8%E4%BC%BC%E6%A0%87%E5%87%86%E5%BD%A2%EF%BC%9A%E5%8A%A8%E6%9C%BA%E4%B8%8E%E5%9F%BA%E7%A1%80.tex}{L14-A-8}。
\begin{solve}
设$\operatorname{adj}(A)\alpha=\lambda \alpha$，则$A\operatorname{adj}(A)\alpha=\lambda A\alpha$.从而$\lambda A\alpha=|A|\alpha=-\alpha\implies A\alpha=-\dfrac{1}{\lambda}\alpha$，故$\alpha$为$A$的特征向量.

            由
\[A\alpha=\begin{pmatrix}
                a & -1 & c \\ 5 & b & 3 \\ 1-c & 0 & -a
            \end{pmatrix}\cdot\begin{pmatrix}
                -1 \\ -1 \\ 1
            \end{pmatrix}=\begin{pmatrix}
                -a+1+c \\ -2-b \\ c-1-a
            \end{pmatrix}=-\dfrac{1}{\lambda}\alpha\]
可知：
            $(-a+1+c)+(c-1-a)=0$，从而$a=c$，代入可知$b=-3,-\dfrac{1}{\lambda}=-1,\lambda=1$.
            此外有$A=\begin{pmatrix}
                a & -1 & a \\ 5 & -3 & 3 \\ 1-a & 0 & -a
            \end{pmatrix}$，由$|A|=-1$解得$a=2$，因此$c=a=2$.
\end{solve}

\begin{exercise}[补充训练 L14-A-9]\label{prob:L14-A-9}
设$A,B\in \mathbf{M}_n(\mathbf{F})$，$AB=BA$，证明：若$X$是矩阵$A$属于特征值$\lambda_0$的特征向量，则$BX\in V_{\lambda_0}$（注：本题是解决很多$AB=BA$类问题的基础）.
\end{exercise}
\noindent 来源：\href{https://github.com/yhwu-is/Linear-Algebra-Left-Undone/blob/new/%E8%AE%B2%E4%B9%89/%E4%B8%93%E9%A2%98/14%20%E7%9B%B8%E4%BC%BC%E6%A0%87%E5%87%86%E5%BD%A2%EF%BC%9A%E5%8A%A8%E6%9C%BA%E4%B8%8E%E5%9F%BA%E7%A1%80.tex}{L14-A-9}。
\begin{solve}
由于$AX=\lambda_0 X$，故$A(BX)=(AB)X=(BA)X=B(AX)=B(\lambda_0X)=\lambda_0(BX)$，从而$BX\in V_{\lambda_0}$.
\end{solve}

\begin{exercise}[补充训练 L14-B-1]\label{prob:L14-B-1}
设$\sigma\in \mathcal{L}(V,W)$，定义$\widetilde{\sigma}:(V/(\ker \sigma))\to W$如下：
        \[\widetilde{\sigma}(v+\ker\sigma)=\sigma(v).\]
        \begin{enumerate}
            \item $\widetilde{\sigma}$是良定义的，且是$(V/(\ker \sigma))$到$W$上的线性映射；

            \item $\widetilde{\sigma}$是单射；

            \item $\operatorname{im} \widetilde{\sigma}=\operatorname{im} \sigma$；

            \item $V/(\ker \sigma)$同构于$\operatorname{im} \sigma$.
        \end{enumerate}
\end{exercise}
\noindent 来源：\href{https://github.com/yhwu-is/Linear-Algebra-Left-Undone/blob/new/%E8%AE%B2%E4%B9%89/%E4%B8%93%E9%A2%98/14%20%E7%9B%B8%E4%BC%BC%E6%A0%87%E5%87%86%E5%BD%A2%EF%BC%9A%E5%8A%A8%E6%9C%BA%E4%B8%8E%E5%9F%BA%E7%A1%80.tex}{L14-B-1}。
\begin{solve}
\begin{enumerate}
                \item 对于$v_1+\ker\sigma=v_2+\ker\sigma$，有$v_1=v_2+u$，其中$u\in\ker\sigma$.因此$\widetilde{\sigma}(v_1+\ker\sigma)=\sigma(v_1)=\sigma(v_1+u)=\sigma(v_2)=\widetilde{\sigma}(v_2+\ker\sigma)$，从而$\widetilde{\sigma}$是良定义的.

                对于$\forall v_1+\ker\sigma,v_2+\ker\sigma \in V/\ker\sigma$和$l,h\in \mathbf{F}$，都有$\widetilde{\sigma}(l(v_1+\ker\sigma)+h(v_2+\ker\sigma))=\widetilde{\sigma}(lv_1+hv_2+\ker\sigma)=\sigma(lv_1+hv_2)=l\sigma(v_1)+h\sigma(v_2)=l\widetilde{\sigma}(v_1+\ker\sigma)+h\widetilde{\sigma}(v_2+\ker\sigma)$，这表明$\widetilde{\sigma}$是$V/(\ker\sigma)$到$W$上的线性映射.
                \item 对于$v_1+\ker\sigma,v_2+\ker\sigma$，若有$\sigma(v_1)=\sigma(v_2)$，则$v_1-v_2\in \ker\sigma$，从而$v_1+\ker\sigma=v_2+\ker\sigma$.因此$\widetilde{\sigma}$为单射.
                \item 对于$\forall u\in \operatorname{im}\widetilde{\sigma}$，有$u \in \operatorname{im}\sigma$，故$\operatorname{im} \widetilde{\sigma} \subseteq \operatorname{im}\sigma$.

                对于$\forall v=\operatorname{im}(w) \in\operatorname{im}\sigma$，$v=\sigma(w)=\widetilde{\sigma}(w+\ker\sigma)\in\operatorname{im}\widetilde{\sigma}$，从而$\operatorname{im}\sigma\subseteq \operatorname{im} \widetilde{\sigma}$.综上有$\operatorname{im} \widetilde{\sigma}=\operatorname{im}\sigma$.
                \item 由$(2)(3)$可知$\widetilde{\sigma}$实际上给出了一个从$V/(\ker \sigma)$到$\operatorname{im}\sigma$的一一映射，从而$V/(\ker\sigma)\cong \operatorname{im}\sigma$.
            \end{enumerate}
\end{solve}

\begin{exercise}[补充训练 L14-B-2]\label{prob:L14-B-2}
设$T\in \mathcal{L}(V)$，证明：$T$是数乘变换（即$T=cI_V$，其中$c$为非零常数，$I_V$为$V$上的恒等映射）的充要条件是$V$的每一个一维子空间都是$T$的不变子空间.
\end{exercise}
\noindent 来源：\href{https://github.com/yhwu-is/Linear-Algebra-Left-Undone/blob/new/%E8%AE%B2%E4%B9%89/%E4%B8%93%E9%A2%98/14%20%E7%9B%B8%E4%BC%BC%E6%A0%87%E5%87%86%E5%BD%A2%EF%BC%9A%E5%8A%A8%E6%9C%BA%E4%B8%8E%E5%9F%BA%E7%A1%80.tex}{L14-B-2}。
\begin{solve}
必要性是显然的.下证充分性.
            设$v_1,v_2,\ldots,v_n$为$V$的一组基，则对于每个$1\leqslant i \leqslant n$，存在$c_i\in \mathbf{F}$，使得$T(v_i)=c_iv_i$.对于$\forall i\neq j$，$T(v_i+v_j)=c_iv_i+c_jv_j\in \operatorname{span}(v_i+v_j)$，这表明$c_i=c_j$.从而所有的$c_i$均相等，$T$为数乘变换.
\end{solve}

\begin{exercise}[补充训练 L14-B-3]\label{prob:L14-B-3}
设$\sigma\in \mathcal{L}(V)$，证明：
        \begin{enumerate}
            \item $\sigma/(\operatorname{im} \sigma)=0$；

            \item $\sigma/(\ker \sigma)$是单射$\iff \ker \sigma\cap\operatorname{im} \sigma=\{0\}$.
        \end{enumerate}
\end{exercise}
\noindent 来源：\href{https://github.com/yhwu-is/Linear-Algebra-Left-Undone/blob/new/%E8%AE%B2%E4%B9%89/%E4%B8%93%E9%A2%98/14%20%E7%9B%B8%E4%BC%BC%E6%A0%87%E5%87%86%E5%BD%A2%EF%BC%9A%E5%8A%A8%E6%9C%BA%E4%B8%8E%E5%9F%BA%E7%A1%80.tex}{L14-B-3}。
\begin{solve}
\begin{enumerate}
                \item 对于$\forall v+\operatorname{im}\sigma$，$(\sigma/(\operatorname{im}\sigma))(v+\operatorname{im}\sigma)=\sigma(v)+\operatorname{im}\sigma=0+\operatorname{im}\sigma$，因此$\sigma/(\operatorname{im}\sigma)=0$.
                \item $\sigma/(\ker\sigma)$为单射$\iff \ker(\sigma/(\ker\sigma))=0+\ker\sigma\iff \ker\sigma\cap\operatorname{im}\sigma=\{0\}$.
            \end{enumerate}
\end{solve}

\begin{exercise}[补充训练 L14-B-4]\label{prob:L14-B-4}
设$V$是有限维的，$T\in \mathcal{L}(V)$且$U$在$T$下不变. 证明：$T/U$的每个特征值均为$T$的特征值.
\end{exercise}
\noindent 来源：\href{https://github.com/yhwu-is/Linear-Algebra-Left-Undone/blob/new/%E8%AE%B2%E4%B9%89/%E4%B8%93%E9%A2%98/14%20%E7%9B%B8%E4%BC%BC%E6%A0%87%E5%87%86%E5%BD%A2%EF%BC%9A%E5%8A%A8%E6%9C%BA%E4%B8%8E%E5%9F%BA%E7%A1%80.tex}{L14-B-4}。
\begin{solve}
对于$T/U$的特征值$\lambda$，设$(T/U)(v+U)=\lambda(v+U)=\lambda v+U$.则$(T/U)(v+U)=Tv+U=\lambda v+U$.
\end{solve}

\begin{exercise}[补充训练 L14-B-5]\label{prob:L14-B-5}
设$V$为$n$维复向量空间，$T\in \mathcal{L}(V)$，$T$在$V$的一组基$e_1,e_2,\ldots,e_n$下的矩阵为对角矩阵$\operatorname{diag}\{d_1,\ldots,d_n\}$，且$d_i\neq d_j\enspace(i\neq j)$.
        \begin{enumerate}
            \item 求$T$的所有一维不变子空间；

            \item 求$T$的所有不变子空间与个数.
        \end{enumerate}
\end{exercise}
\noindent 来源：\href{https://github.com/yhwu-is/Linear-Algebra-Left-Undone/blob/new/%E8%AE%B2%E4%B9%89/%E4%B8%93%E9%A2%98/14%20%E7%9B%B8%E4%BC%BC%E6%A0%87%E5%87%86%E5%BD%A2%EF%BC%9A%E5%8A%A8%E6%9C%BA%E4%B8%8E%E5%9F%BA%E7%A1%80.tex}{L14-B-5}。
\begin{solve}
\begin{enumerate}
                \item $T$的所有一维不变子空间为$\operatorname{span}(e_i),\enspace i=1,2,\ldots,n$.
                \item $T$的所有不变子空间为所有$\operatorname{span}(e_{i_1},e_{i_2},\ldots,e_{i_k})$，其中$(i_1,i_2,\ldots,i_k)\subset\{1,2,\ldots,n\}$.

                $T$的不变子空间个数为$2^n$.
            \end{enumerate}
\end{solve}

\begin{exercise}[补充训练 L14-B-6]\label{prob:L14-B-6}
给定$\mathbf{R}$上的2维线性空间$V$上的算子$T$，其在一组基$\alpha_1,\alpha_2$下的矩阵为
        \[\begin{pmatrix}
                0 & 1 \\ 1-a & 0
            \end{pmatrix}.\]
        求$T$的所有不变子空间.
\end{exercise}
\noindent 来源：\href{https://github.com/yhwu-is/Linear-Algebra-Left-Undone/blob/new/%E8%AE%B2%E4%B9%89/%E4%B8%93%E9%A2%98/14%20%E7%9B%B8%E4%BC%BC%E6%A0%87%E5%87%86%E5%BD%A2%EF%BC%9A%E5%8A%A8%E6%9C%BA%E4%B8%8E%E5%9F%BA%E7%A1%80.tex}{L14-B-6}。
\begin{solve}
令$|\lambda I-T|=0$，若$a<1$，则$\lambda_1=\sqrt{1-a},\lambda_2=-\sqrt{1-a}$，相应的特征向量为$(\alpha_1+\sqrt{1-a}\alpha_2),(\alpha_1-\sqrt{1-a}\alpha_2)$.从而$T$的所有不变子空间为$\{0\},\operatorname{span}(\alpha_1+\sqrt{1-a}\alpha_2),\operatorname{span}(\alpha_1-\sqrt{1-a}\alpha_2),V$.

            若$a=1$，则$\lambda=0$，相应的特征向量为$\alpha_1$，从而$T$的所有不变子空间为$\{0\},\operatorname{span}(\alpha_1),V$.

            若$a>1$，则$T$在$\mathbf{R}$上无特征值，$T$的不变子空间只有$\{0\}$和$V$.
\end{solve}

\begin{exercise}[补充训练 L14-B-7]\label{prob:L14-B-7}
设$A,B$都是$n$阶矩阵，且$r(A)+r(B)<n$，证明：$A,B$有相同的特征值和特征向量.
\end{exercise}
\noindent 来源：\href{https://github.com/yhwu-is/Linear-Algebra-Left-Undone/blob/new/%E8%AE%B2%E4%B9%89/%E4%B8%93%E9%A2%98/14%20%E7%9B%B8%E4%BC%BC%E6%A0%87%E5%87%86%E5%BD%A2%EF%BC%9A%E5%8A%A8%E6%9C%BA%E4%B8%8E%E5%9F%BA%E7%A1%80.tex}{L14-B-7}。
\begin{solve}
由于$r(A)+r(B)<n$，故$A$和$B$都有特征值$0$.

            由于$\dim(N(A)\cap N(B))=\dim(N(A))+\dim(N(B))-\dim(N(A)+N(B))\geqslant\dim(N(A))+\dim(N(B))-n=n-r(A)+n-r(B)-n=n-(r(A)+r(B))\ge 0$，这表明$N(A)\cap N(B)\neq \{0\}$，从而存在$A$和$B$有相同的特征向量（对应于特征值$0$）.
\end{solve}

\begin{exercise}[补充训练 L14-B-8]\label{prob:L14-B-8}
设$\lambda_1,\lambda_2,\ldots,\lambda_n$是矩阵$A=(a_{ij})_{n\times n}$的$n$个特征值，证明：$\lambda_1^2,\lambda_2^2,\ldots,\lambda_n^2$是$A^2$的$n$个特征值，且$\displaystyle\sum_{i=1}^{n}\lambda_i^2=\sum_{j=1}^{n}\displaystyle\sum_{k=1}^{n}a_{jk}a_{kj}$.
\end{exercise}
\noindent 来源：\href{https://github.com/yhwu-is/Linear-Algebra-Left-Undone/blob/new/%E8%AE%B2%E4%B9%89/%E4%B8%93%E9%A2%98/14%20%E7%9B%B8%E4%BC%BC%E6%A0%87%E5%87%86%E5%BD%A2%EF%BC%9A%E5%8A%A8%E6%9C%BA%E4%B8%8E%E5%9F%BA%E7%A1%80.tex}{L14-B-8}。
\begin{solve}
对于$\lambda_i,i=1,2,\ldots,n$，存在$u_i\in V,u_i\neq 0$，使得$Au_i=\lambda_iu_i$，从而$A^2u_i=A(\lambda_iu_i)=\lambda_iA(u_i)=\lambda_i^2u_i$，这表明$\lambda_i^2$为$A^2$的特征值，从而$\lambda_1^2,\lambda_2^2,\ldots,\lambda_n^2$为$A^2$的$n$个特征值.进而知$\displaystyle \sum_{i=1}^{n}\lambda_i^2=tr(A^2)=\sum_{j=1}^{n}\displaystyle\sum_{k=1}^{n}a_{jk}a_{kj}$.
\end{solve}

\begin{exercise}[补充训练 L14-B-9]\label{prob:L14-B-9}
设$A$为$n$阶矩阵，$X_1,X_2,X_3$为$n$元列向量，且$AX_1=kX_1\enspace(X_1\neq 0),AX_2=lX_1+kX_2,AX_3=lX_2+kX_3\enspace(l\neq 0)$. 证明：$X_1,X_2,X_3$线性无关.
\end{exercise}
\noindent 来源：\href{https://github.com/yhwu-is/Linear-Algebra-Left-Undone/blob/new/%E8%AE%B2%E4%B9%89/%E4%B8%93%E9%A2%98/14%20%E7%9B%B8%E4%BC%BC%E6%A0%87%E5%87%86%E5%BD%A2%EF%BC%9A%E5%8A%A8%E6%9C%BA%E4%B8%8E%E5%9F%BA%E7%A1%80.tex}{L14-B-9}。
\begin{solve}
设$a_1X_1+a_2X_2+a_3X_3=0$.在等式两边同时左乘$(A-kE)$可知：$a_2lX_1+a_3lX_2=0$.再左乘$(A-kE)$知$a_3l^2=0$.结合$l\neq 0$，有$a_3=0$.回代知$a_1=a_2=0$，从而$X_1,X_2,X_3$线性无关.
\end{solve}

\begin{exercise}[补充训练 L14-C-1]\label{prob:L14-C-1}
设$A,B\in \mathbf{M}_n(\mathbf{C})$，$B$的特征多项式$f(\lambda)=|\lambda E-B|$. 证明：$f(A)$可逆的充要条件是$B$的任一特征值都不是$A$的特征值.
\end{exercise}
\noindent 来源：\href{https://github.com/yhwu-is/Linear-Algebra-Left-Undone/blob/new/%E8%AE%B2%E4%B9%89/%E4%B8%93%E9%A2%98/14%20%E7%9B%B8%E4%BC%BC%E6%A0%87%E5%87%86%E5%BD%A2%EF%BC%9A%E5%8A%A8%E6%9C%BA%E4%B8%8E%E5%9F%BA%E7%A1%80.tex}{L14-C-1}。
\begin{solve}
$f(\lambda)=|\lambda E-B|=(\lambda-\lambda_1)^{\dim E(\lambda_1,B)}\cdots(\lambda-\lambda_m)^{\dim E(\lambda_m,B)}$，故$f(A)$可逆$\iff (A-\lambda_1E)^{\dim E(\lambda_1,B)}\cdots(A-\lambda_mE)^{\dim E(\lambda_m,B)}$可逆$\iff \mid(A-\lambda_1E)^{\dim E(\lambda_1,B)}\cdots(A-\lambda_mE)^{\dim E(\lambda_m,B)}\mid \neq 0 \iff \mid(A-\lambda_1E)\mid^{\dim E(\lambda_1,B)}\cdots\mid(A-\lambda_mE)\mid^{\dim E(\lambda_m,B)}\neq 0 \iff \forall i=1,2,\ldots,m,|A-\lambda_mE|\neq 0 \iff B$的任一特征值都不是$A$的特征值.
\end{solve}

\begin{exercise}[补充训练 L14-C-2]\label{prob:L14-C-2}
证明：若$AB=BA$，则$A$和$B$至少有一个共同的特征向量.
\end{exercise}
\noindent 来源：\href{https://github.com/yhwu-is/Linear-Algebra-Left-Undone/blob/new/%E8%AE%B2%E4%B9%89/%E4%B8%93%E9%A2%98/14%20%E7%9B%B8%E4%BC%BC%E6%A0%87%E5%87%86%E5%BD%A2%EF%BC%9A%E5%8A%A8%E6%9C%BA%E4%B8%8E%E5%9F%BA%E7%A1%80.tex}{L14-C-2}。
\begin{solve}
命题等价于$\sigma,\tau \in \mathcal{L}(V),\sigma\tau=\tau\sigma$，则$\sigma$和$\tau$至少有一个公共的特征向量.

            考虑$\sigma$的不变子空间$E(\lambda,\sigma)$，对于$v\in E(\lambda,\sigma)$，有$\sigma v=\lambda v \implies \sigma(\tau v)=\tau\sigma v=\lambda (\tau v)$，这表明$E(\lambda,\sigma)$是$\tau$的不变子空间.考虑限制映射$\tau|_E(\lambda,\sigma)$，它一定有特征值和相应的特征向量$u\in E(\lambda,\sigma),u \neq 0$，从而$u$为$\sigma$和$\tau$的公共特征向量.
\end{solve}

\begin{exercise}[补充训练 L14-例-4]\label{prob:L14-例-4}
设$A=\begin{pmatrix}
            1 & -1 & 0 \\ 2 & 0 & 1 \\ 1 & a & 0
        \end{pmatrix}$，且存在非零向量$\alpha$使得$A\alpha=2\alpha$，求$a$.
\end{exercise}
\noindent 来源：\href{https://github.com/yhwu-is/Linear-Algebra-Left-Undone/blob/new/%E8%AE%B2%E4%B9%89/%E4%B8%93%E9%A2%98/14%20%E7%9B%B8%E4%BC%BC%E6%A0%87%E5%87%86%E5%BD%A2%EF%BC%9A%E5%8A%A8%E6%9C%BA%E4%B8%8E%E5%9F%BA%E7%A1%80.tex}{L14-例-4}。
\begin{solve}
由题意知2是矩阵$A$的特征值，因此我们有
    \[|2E-A|=\begin{vmatrix}
            1 & 1 & 0 \\ -2 & 2 & -1 \\ -1 & -a & 2
        \end{vmatrix}=9-a=0,\]
    因此$a=9$.
\end{solve}

\begin{exercise}[补充训练 L14-例-6]\label{prob:L14-例-6}
对下列矩阵$A$的特征值，能做出怎样的断言？
    \begin{enumerate}
        \item $A$可逆/$A$不可逆/$E+A$可逆/$4E+A$不可逆；

        \item $|E - A^2| = 0$；

        \item $A^2 = E$（对合）/$A^2 = A$（幂等）/$A^k = 0$（幂零）；

        \item $A=\lambda_0E+B$（$\lambda_0$为常数，且已知$B$的$n$个特征值为$\lambda_1,\lambda_2,\ldots,\lambda_n$）；

        \item $A$为对角块矩阵，即$A=\operatorname{diag}(A_1,A_2,\ldots,A_m)$.
    \end{enumerate}
\end{exercise}
\noindent 来源：\href{https://github.com/yhwu-is/Linear-Algebra-Left-Undone/blob/new/%E8%AE%B2%E4%B9%89/%E4%B8%93%E9%A2%98/14%20%E7%9B%B8%E4%BC%BC%E6%A0%87%E5%87%86%E5%BD%A2%EF%BC%9A%E5%8A%A8%E6%9C%BA%E4%B8%8E%E5%9F%BA%E7%A1%80.tex}{L14-例-6}。
\begin{solve}
\begin{enumerate}
        \item $A$可逆时有$|A|=\lambda_1\cdots\lambda_n\neq 0$，因此$A$的特征值都不为0. 同理，$A$不可逆表明存在特征值等于0，$E+A$可逆表明$-1$不是$A$的特征值，$4E+A$不可逆表明$-4$是$A$的特征值.

        \item $|E - A^2| = |E - A||E + A| = 0$，因此 $A$ 的特征值可能为 $1$ 或 $-1$ 或二者兼有.

        \item 首先考虑对合矩阵，接下来的同理可以得到类似结论. 由于$A^2=E$，设$AX=\lambda X$，则$A^2X=\lambda^2X=X$，因此$\lambda^2=1$，即$\lambda=\pm 1$，因此$1$或$-1$是$A$的特征值.

              但这里需要强调的是，不同于前两问，前两问中我们都是说某些值是$A$的特征值，但无法保证$A$的特征值只能是某些值，但在本题这样给出矩阵方程的情况下，我们可以得到$A$的特征值只能是$\pm 1$，没有其他值. 我们用反证法，假设存在$\lambda_0\neq\pm 1$是$A$的特征值，即$AX=\lambda_0X$，则$A^2X=\lambda_0^2X\neq X$（因为$X$不是零向量），导出矛盾. 当然有同学可能会思考，$A$的特征值一定兼有$\pm 1$吗，事实上并非如此，例如$E$满足$E^2=E$，但其特征值只有1，$-E$满足$(-E)^2=E$，但其特征值只有$-1$. 并且利用下一讲对角化的结论可以知道（我们放在下一讲习题中供读者练习），满足$A^2=E$且特征值只有1的矩阵只能是$E$，特征值只有$-1$的矩阵只能是$-E$.

              注：本题解决过程中告诉我们一个解题技巧，如果看到$A$的多项式$f(A)=O$这种形式的表达式，事实上$A$的特征值只能是$f(\lambda)=0$的根，如上题中$f(A)=A^2-E$，则$f(\lambda)=\lambda^2-1$，因此$A$的特征值只能是$\pm 1$. 同理，幂等矩阵的特征值只能是0和1，幂零矩阵的特征值只能是0（这是一个重要的幂零矩阵等价条件，未来我们会再次遇到）. 这个技巧的证明放在接下来的原讲义“多项式零点为特征值”处 中.

        \item 设$BX=\lambda_iX_i(X_i\neq 0,\enspace i=1,\ldots,n)$，则$AX_i=\lambda_0X_i+BX_i=\lambda_0X_i+\lambda_iX_i=(\lambda_0+\lambda_i)X_i$，因此$\lambda_0+\lambda_i\enspace(i=1,\ldots,n)$都是$A$的特征值.

        \item \begin{align*}
                  |\lambda E-A| & =\begin{vmatrix}
                                       \lambda E_1-A_1 & 0               & \cdots & 0               \\
                                       0               & \lambda E_2-A_2 & \cdots & 0               \\
                                       \vdots          & \vdots          & \ddots & \vdots          \\
                                       0               & 0               & \cdots & \lambda E_m-A_m
                                   \end{vmatrix}
                                & =\prod_{i=1}^{m}|\lambda E_i-A_i|=0
              \end{align*}
              因此，$A_i\enspace(i=1,\ldots,m)$的特征值都是$A$的特征值.
    \end{enumerate}
\end{solve}

\begin{exercise}[补充训练 L14-例-8]\label{prob:L14-例-8}
回答以下问题：
    \begin{enumerate}
        \item 设$\alpha=(1,0,-1)^\mathrm{T}$，且$A=\alpha\alpha^\mathrm{T}$，求$|6E-A^n|$；

        \item 设$A$为三阶矩阵，其特征值为$1,-2,-1$，求$|A|$，$\operatorname{adj}(A)+3E$的特征值，$(A^{-1})^2+2E$的特征值以及$|A^2-A+E|$；

        \item 设$A$为三阶矩阵，$A^2-A-2E=O$，$|A|=2$，求$|\operatorname{adj}(A)+3E|$；

        \item 设$A$为三阶矩阵，其特征值为$-1,-1,5$，求$A_{11}+A_{22}+A_{33}$；
    \end{enumerate}
\end{exercise}
\noindent 来源：\href{https://github.com/yhwu-is/Linear-Algebra-Left-Undone/blob/new/%E8%AE%B2%E4%B9%89/%E4%B8%93%E9%A2%98/14%20%E7%9B%B8%E4%BC%BC%E6%A0%87%E5%87%86%E5%BD%A2%EF%BC%9A%E5%8A%A8%E6%9C%BA%E4%B8%8E%E5%9F%BA%E7%A1%80.tex}{L14-例-8}。
\begin{solve}
\begin{enumerate}
        \item 事实上$A=\alpha\alpha^\mathrm{T}=\begin{pmatrix}
                      1 & 0 & -1 \\ 0 & 0 & 0 \\ -1 & 0 & 1
                  \end{pmatrix}$，由$|\lambda E-A|=0$解得$A$的特征值为$\lambda_1=\lambda_2=0,\lambda_3=2$，而根据$A^n$的特征值性质和原讲义“特征值的性质”处 可知，$6E-A^n$的特征值为$6-\lambda_1^n,6-\lambda_2^n,6-\lambda_3^n$，即$6,6,6-2^n$，因此$|6E-A^n|=6^2(6-2^n)=36(6-2^n)$.

        \item 由于$A$的特征值为$1,-2,-1$，因此$|A|=1\times(-2)\times(-1)=2$，而$\operatorname{adj}(A)$的特征值为$|A|\lambda^{-1}$，因此$\operatorname{adj}(A)$的特征值为$2,-1,-2$，故$\operatorname{adj}(A)+3E$的特征值为$\operatorname{adj}(A)$的特征值加3（根据原讲义“特征值的性质”处），即为$5,2,1$，又根据$A^{-1}$和$A^2$特征值的性质可知，$(A^{-1})^2+2E$的特征值为$1^2+2,(-1/2)^2+2,(-1)^2+2$，即为$3,9/4,3$，而$A^2-A+E$的特征值根据$f(\sigma)$特征值性质的讨论可知为$1^2-1+1,(-2)^2-(-2)+1,(-1)^2-(-1)+1$，即为$1,7,3$，因此$|A^2-A+E|=1\times 7\times 3=21$.

        \item 根据原讲义“多项式零点为特征值”处，$A$的特征值只能出自$-1$和$2$. 又$|A|=2$，且$A$为3阶矩阵，因此$A$的3个特征值必为$-1,-1,2$.

              又$\operatorname{adj}(A)$的特征值为$|A|\lambda^{-1}$，因此$\operatorname{adj}(A)$的特征值为$1,-2,-2$，又根据原讲义“特征值的性质”处 的结论，$\operatorname{adj}(A)+3E$的特征值为$\operatorname{adj}(A)$的特征值加3，即$\lambda_1=\lambda_2=1,\lambda_3=4$，故$|\operatorname{adj}(A)+3E|=\lambda_1\lambda_2\lambda_3=4$.

        \item 由题意知$|A|=5$，故$\operatorname{adj}(A)$的特征值为$|A|\lambda^{-1}$即为$\mu_1=\mu_2=-5,\mu_3=1$，而$A_{11}+A_{22}+A_{33}$就是$\operatorname{adj}(A)$的迹（即矩阵对角线元素之和），因此$A_{11}+A_{22}+A_{33}=\mu_1+\mu_2+\mu_3=-9$.
    \end{enumerate}
\end{solve}

\begin{exercise}[补充训练 L14-例-13]\label{prob:L14-例-13}
设$\sigma\in\mathcal{L}(\mathbf{F}^2)$定义为$\sigma(w,z)=(-z,w)$.
    \begin{enumerate}
        \item 当$\mathbf{F}=\mathbf{R}$时，求$\sigma$的特征值和特征向量；

        \item 当$\mathbf{F}=\mathbf{C}$时，求$\sigma$的特征值和特征向量.
    \end{enumerate}
\end{exercise}
\noindent 来源：\href{https://github.com/yhwu-is/Linear-Algebra-Left-Undone/blob/new/%E8%AE%B2%E4%B9%89/%E4%B8%93%E9%A2%98/14%20%E7%9B%B8%E4%BC%BC%E6%A0%87%E5%87%86%E5%BD%A2%EF%BC%9A%E5%8A%A8%E6%9C%BA%E4%B8%8E%E5%9F%BA%E7%A1%80.tex}{L14-例-13}。
\begin{solve}
我们首先写出$\sigma$在任意一组基下的矩阵表示，为了方便，我们选取标准基$e_1=(1,0),e_2=(0,1)$，则矩阵表示为
    \[ A=\begin{pmatrix}
            0 & -1 \\ 1 & 0
        \end{pmatrix}. \]
    则其特征多项式$f(\lambda)=|\lambda E-A|=\lambda^2+1$，因此
    \begin{enumerate}
        \item 在实数域上无特征值和特征向量；

        \item 复数域上特征值为$\pm\mathrm{i}$，其中$\mathrm{i}$对应的特征向量为$(w,-w\mathrm{i})$，$-\mathrm{i}$对应的特征向量为$(w,w\mathrm{i})$，其中$w\in\mathbf{R}$且$w\neq 0$.
    \end{enumerate}
\end{solve}

\section{补充训练：不变子空间与对角化}

本节包含 41 道训练题，建议在完成相关专题真题后选择性练习。

\begin{exercise}[补充训练 L15-A-1]\label{prob:L15-A-1}
请举例：存在两个矩阵相抵但不相似.
\end{exercise}
\noindent 来源：\href{https://github.com/yhwu-is/Linear-Algebra-Left-Undone/blob/new/%E8%AE%B2%E4%B9%89/%E4%B8%93%E9%A2%98/15%20%E7%9B%B8%E4%BC%BC%E6%A0%87%E5%87%86%E5%BD%A2%EF%BC%9A%E5%A4%8D%E6%95%B0%E5%9F%9F%E4%B8%8A%E7%9A%84%E5%B0%9D%E8%AF%95%E4%B8%8E%E7%90%86%E8%AE%BA.tex}{L15-A-1}。
\begin{solve}
非常简单，可以举$\begin{pmatrix}
                1 & 0 \\ 0 & 1
            \end{pmatrix}$和$\begin{pmatrix}
                1 & 0 \\ 0 & 2
            \end{pmatrix}$的例子，也可以举一个可对角化另一个不可对角化的例子，因为上述情况矩阵对应的相似标准形不一样，故不相似.
\end{solve}

\begin{exercise}[补充训练 L15-A-2]\label{prob:L15-A-2}
求矩阵
        \[A = \begin{pmatrix}
                0  & -1 & 1 \\
                -1 & 0  & 1 \\
                1  & 1  & 0
            \end{pmatrix}\]
        的所有特征值，对应的特征子空间，以及与 $A$ 相似的一个对角矩阵.
\end{exercise}
\noindent 来源：\href{https://github.com/yhwu-is/Linear-Algebra-Left-Undone/blob/new/%E8%AE%B2%E4%B9%89/%E4%B8%93%E9%A2%98/15%20%E7%9B%B8%E4%BC%BC%E6%A0%87%E5%87%86%E5%BD%A2%EF%BC%9A%E5%A4%8D%E6%95%B0%E5%9F%9F%E4%B8%8A%E7%9A%84%E5%B0%9D%E8%AF%95%E4%B8%8E%E7%90%86%E8%AE%BA.tex}{L15-A-2}。
\begin{solve}
解$|\lambda E-A|=0$可得特征值为$1,1,-2$. 然后解$(E-A)X=0$得$X=t_1(-1,1,0)^\mathrm{T}+t_2(1,0,1)^\mathrm{T}$；解$(-2E-A)X=0$，得$X=t(-1,-1,1)^\mathrm{T}$. 可见特征值1对应的特征子空间为$\operatorname{span}((-1,1,0)^\mathrm{T},(1,0,1)^\mathrm{T})$，特征值$-2$对应的特征子空间为 $\operatorname{span}((-1,-1,1)^\mathrm{T})$. 可知与$A$相似的对角矩阵为 $\operatorname{diag}(1,1,-2)$.
\end{solve}

\begin{exercise}[补充训练 L15-A-3]\label{prob:L15-A-3}
设$A=\begin{pmatrix}
                a & b \\ c & d
            \end{pmatrix}$为二阶实矩阵.
        \begin{enumerate}
            \item 若$|A|<0$，问：$A$与对角矩阵是否相似；

            \item 若$ad-bc=1$，$|a+d|>2$，问：$A$是否可对角化.
        \end{enumerate}
\end{exercise}
\noindent 来源：\href{https://github.com/yhwu-is/Linear-Algebra-Left-Undone/blob/new/%E8%AE%B2%E4%B9%89/%E4%B8%93%E9%A2%98/15%20%E7%9B%B8%E4%BC%BC%E6%A0%87%E5%87%86%E5%BD%A2%EF%BC%9A%E5%A4%8D%E6%95%B0%E5%9F%9F%E4%B8%8A%E7%9A%84%E5%B0%9D%E8%AF%95%E4%B8%8E%E7%90%86%E8%AE%BA.tex}{L15-A-3}。
\begin{solve}
\begin{enumerate}
                \item $f(\lambda)=|\lambda E-A|=\begin{vmatrix}
                              \lambda-a & -b \\ -c & \lambda-d
                          \end{vmatrix}=\lambda^2-(a+d)\lambda+(ad-bc)=\lambda^2-(a+d)\lambda+|A|$. 由于$|A|<0$，因此判别式$\Delta=(a+d)^2-4|A|>0$，因此二阶矩阵$A$有两个互异特征值，故可对角化，因此与对角矩阵相似；

                \item 同 (1)，判别式$\Delta=(a+d)^2-4|A|>0$，因此二阶矩阵$A$有两个不同的特征值，故可对角化.
            \end{enumerate}
\end{solve}

\begin{exercise}[补充训练 L15-A-4]\label{prob:L15-A-4}
设$A=\begin{pmatrix}
                1 & 1 & a \\
                1 & a & 1 \\
                a & 1 & 1
            \end{pmatrix}$，$\beta=\begin{pmatrix}
                1 \\ 1 \\ -2
            \end{pmatrix}$，方程组$AX=\beta$有解但不唯一.
        \begin{enumerate}
            \item 求$a$的值；

            \item 求可逆矩阵$P$使得$P^{-1}AP$为对角矩阵.
        \end{enumerate}
\end{exercise}
\noindent 来源：\href{https://github.com/yhwu-is/Linear-Algebra-Left-Undone/blob/new/%E8%AE%B2%E4%B9%89/%E4%B8%93%E9%A2%98/15%20%E7%9B%B8%E4%BC%BC%E6%A0%87%E5%87%86%E5%BD%A2%EF%BC%9A%E5%A4%8D%E6%95%B0%E5%9F%9F%E4%B8%8A%E7%9A%84%E5%B0%9D%E8%AF%95%E4%B8%8E%E7%90%86%E8%AE%BA.tex}{L15-A-4}。
\begin{solve}
\begin{enumerate}
              \item 由Cramer法则，$A$为方阵且$AX=\beta$有解但不唯一，即$|A|\neq 0$，解得$a=-2$或$a=1$，分别代入可知$a=1$时增广矩阵的秩为2，而$r(A)=1$，故无解；$a=-2$时增广矩阵的秩为2，而$r(A)=2$，故有解，故$a=-2$.

              \item 容易求得特征值为$0,-3,3$，求特征向量可知$P=\begin{pmatrix}
                            1 & -1 & 1 \\ 1 & 0 & -2 \\ 1 & 1 & 1
                        \end{pmatrix}$.
          \end{enumerate}
\end{solve}

\begin{exercise}[补充训练 L15-A-5]\label{prob:L15-A-5}
设$A$为三阶矩阵，$\alpha_1,\alpha_2,\alpha_3$线性无关，且$A\alpha_1=\alpha_1,A\alpha_2=\alpha_1+\alpha_2-2\alpha_3,A\alpha_3=\alpha_1-2\alpha_2+\alpha_3$，求$A$的特征值.
\end{exercise}
\noindent 来源：\href{https://github.com/yhwu-is/Linear-Algebra-Left-Undone/blob/new/%E8%AE%B2%E4%B9%89/%E4%B8%93%E9%A2%98/15%20%E7%9B%B8%E4%BC%BC%E6%A0%87%E5%87%86%E5%BD%A2%EF%BC%9A%E5%A4%8D%E6%95%B0%E5%9F%9F%E4%B8%8A%E7%9A%84%E5%B0%9D%E8%AF%95%E4%B8%8E%E7%90%86%E8%AE%BA.tex}{L15-A-5}。
\begin{solve}
由分块矩阵乘法，
\[A(\alpha_1,\alpha_2,\alpha_3)=(A\alpha_1,A\alpha_2,A\alpha_3)=(\alpha_1,\alpha_1+\alpha_2-2\alpha_3,\alpha_1-2\alpha_2+\alpha_3)=(\alpha_1,\alpha_2,\alpha_3)\begin{pmatrix}
                1 & 1 & 1 \\ 0 & 1 & -2 \\ 0 & -2 & 1
            \end{pmatrix}\]
. 令$P=(\alpha_1,\alpha_2,\alpha_3)$，则$AP=P\begin{pmatrix}
                1 & 1 & 1 \\ 0 & 1 & -2 \\ 0 & -2 & 1
            \end{pmatrix}=PB$，即$P^{-1}AP=B$，故$A$与$B$相似，故两矩阵的特征值相同，简单求解$B$的特征值即得$A$的特征值为$1,-1,3$.
\end{solve}

\begin{exercise}[补充训练 L15-A-6]\label{prob:L15-A-6}
设三阶实对称矩阵$A$的各行元素之和为3，向量$\alpha_1=(-1,2,-1)^\mathrm{T}$，$\alpha_2=(0,-1,1)^\mathrm{T}$是方程组$AX=0$的两个解，求矩阵$P$使得$P^{-1}AP$为对角矩阵.
\end{exercise}
\noindent 来源：\href{https://github.com/yhwu-is/Linear-Algebra-Left-Undone/blob/new/%E8%AE%B2%E4%B9%89/%E4%B8%93%E9%A2%98/15%20%E7%9B%B8%E4%BC%BC%E6%A0%87%E5%87%86%E5%BD%A2%EF%BC%9A%E5%A4%8D%E6%95%B0%E5%9F%9F%E4%B8%8A%E7%9A%84%E5%B0%9D%E8%AF%95%E4%B8%8E%E7%90%86%E8%AE%BA.tex}{L15-A-6}。
\begin{solve}
每行元素之和为3，则我们知道$\alpha=(1,1,1)^\mathrm{T}$是3对应的特征向量（具体理由只需要验证$A\alpha=3\alpha$即可，这是很常用的性质），而$AX=0$的解是对应特征值0的特征向量，故由题意可知取$P=\begin{pmatrix}
                1 & -1 & 0 \\ 1 & 2 & -1 \\ 1 & -1 & 1
            \end{pmatrix}$，则$P^{-1}AP=\begin{pmatrix}
                3 & 0 & 0 \\ 0 & 0 & 0 \\ 0 & 0 & 0
            \end{pmatrix}$.
\end{solve}

\begin{exercise}[补充训练 L15-B-1]\label{prob:L15-B-1}
设$a\neq b$，且$(aE-A)(bE-A)=O$. 证明：$A$可对角化（特例：对合矩阵）；
\end{exercise}
\noindent 来源：\href{https://github.com/yhwu-is/Linear-Algebra-Left-Undone/blob/new/%E8%AE%B2%E4%B9%89/%E4%B8%93%E9%A2%98/15%20%E7%9B%B8%E4%BC%BC%E6%A0%87%E5%87%86%E5%BD%A2%EF%BC%9A%E5%A4%8D%E6%95%B0%E5%9F%9F%E4%B8%8A%E7%9A%84%E5%B0%9D%E8%AF%95%E4%B8%8E%E7%90%86%E8%AE%BA.tex}{L15-B-1}。
\begin{solve}
回忆例题中$A^2=2A$，即$A(A-2E)=O$的证明，我们知道$A$的特征值只能为$a$或$b$，然后我们可以证明$r(aE-A)+r(bE-A)\leqslant n$和$r(aE-A)+r(bE-A)\geqslant n$，因此$r(aE-A)+r(bE-A)=n$，即$A$可对角化（由于和例题完全类似，这里不再展开具体做法）.
\end{solve}

\begin{exercise}[补充训练 L15-B-2]\label{prob:L15-B-2}
证明：满足$A^2=E$且特征值只有1的矩阵只能是$E$，特征值只有$-1$的矩阵只能是$-E$.
\end{exercise}
\noindent 来源：\href{https://github.com/yhwu-is/Linear-Algebra-Left-Undone/blob/new/%E8%AE%B2%E4%B9%89/%E4%B8%93%E9%A2%98/15%20%E7%9B%B8%E4%BC%BC%E6%A0%87%E5%87%86%E5%BD%A2%EF%BC%9A%E5%A4%8D%E6%95%B0%E5%9F%9F%E4%B8%8A%E7%9A%84%E5%B0%9D%E8%AF%95%E4%B8%8E%E7%90%86%E8%AE%BA.tex}{L15-B-2}。
\begin{solve}
若$A^2=E$，则$A$可对角化.设$A=P^{-1}BP$，$B$为对角矩阵，若$A$的特征值只有$1$，则$B=E$，从而$A=P^{-1}EP=E$.若$A$的特征值只有$-1$，则$B=-E$，从而$A=-E$.
\end{solve}

\begin{exercise}[补充训练 L15-B-3]\label{prob:L15-B-3}
设$A$为三阶矩阵，$A^2=A$且$r(A)=r$，求$|A-2E|$.
\end{exercise}
\noindent 来源：\href{https://github.com/yhwu-is/Linear-Algebra-Left-Undone/blob/new/%E8%AE%B2%E4%B9%89/%E4%B8%93%E9%A2%98/15%20%E7%9B%B8%E4%BC%BC%E6%A0%87%E5%87%86%E5%BD%A2%EF%BC%9A%E5%A4%8D%E6%95%B0%E5%9F%9F%E4%B8%8A%E7%9A%84%E5%B0%9D%E8%AF%95%E4%B8%8E%E7%90%86%E8%AE%BA.tex}{L15-B-3}。
\begin{solve}
根据例题中关于$A^2=A$的讨论可知，满足$A^2=A$的矩阵$A$可对角化（设过渡矩阵为$P$），特征值0和1的重数分别为$n-r$和$r$，因此$|A-2E|=|P^{-1}||A-2E||P|=|\operatorname{diag}(-1,\ldots,-1,-2,\ldots,-2)|=(-1)^r(-2)^{n-r}$.
\end{solve}

\begin{exercise}[补充训练 L15-B-4]\label{prob:L15-B-4}
设$T\in \mathcal{L}(\mathbf{C}^3)$使得6和7是$T$的特征值，且$T$不可对角化. 证明：存在$(x,y,z)\in\mathbf{C}^3$使得$T(x,y,z)=(17+8x,\sqrt{5}+8y,2\pi+8z)$.
\end{exercise}
\noindent 来源：\href{https://github.com/yhwu-is/Linear-Algebra-Left-Undone/blob/new/%E8%AE%B2%E4%B9%89/%E4%B8%93%E9%A2%98/15%20%E7%9B%B8%E4%BC%BC%E6%A0%87%E5%87%86%E5%BD%A2%EF%BC%9A%E5%A4%8D%E6%95%B0%E5%9F%9F%E4%B8%8A%E7%9A%84%E5%B0%9D%E8%AF%95%E4%B8%8E%E7%90%86%E8%AE%BA.tex}{L15-B-4}。
\begin{solve}
三阶矩阵如果有三个特征值则一定可对角化，故由题意$T$只能有6和7两个特征值. 故$\lambda\neq 6,7$时$T-\lambda I$都是可逆的（回顾正文中的定理），故$T-8I$可逆（故是满射），因此一定存在$(x,y,z)\in\mathbf{C}^3$使得$(T-8I)(x,y,z)=(17,\sqrt{5},2\pi)$，因此存在$(x,y,z)\in\mathbf{C}^3$使得$T(x,y,z)=(17+8x,\sqrt{5}+8y,2\pi+8z)$.
\end{solve}

\begin{exercise}[补充训练 L15-B-5]\label{prob:L15-B-5}
证明：两个可对角化的同阶矩阵特征值相同（包括重数）等价于它们相似. 对于不可对角化的矩阵来说，这一结论还成立吗？
\end{exercise}
\noindent 来源：\href{https://github.com/yhwu-is/Linear-Algebra-Left-Undone/blob/new/%E8%AE%B2%E4%B9%89/%E4%B8%93%E9%A2%98/15%20%E7%9B%B8%E4%BC%BC%E6%A0%87%E5%87%86%E5%BD%A2%EF%BC%9A%E5%A4%8D%E6%95%B0%E5%9F%9F%E4%B8%8A%E7%9A%84%E5%B0%9D%E8%AF%95%E4%B8%8E%E7%90%86%E8%AE%BA.tex}{L15-B-5}。
\begin{solve}
可对角化矩阵特征值（包括重数）相等则相似标准形（对角矩阵）相等，故一定相似. 不可对角化的举例说明不成立：$\begin{pmatrix}
                0 & 1 & 0 \\ 0 & 0 & 1 \\ 0 & 0 & 0
            \end{pmatrix}$和$\begin{pmatrix}
                0 & 1 & 0 \\ 0 & 0 & 0 \\ 0 & 0 & 0
            \end{pmatrix}$的特征值全为0（重数均为3），但它们对应不同的若当标准形（之后章节会介绍），因此不相似.
\end{solve}

\begin{exercise}[补充训练 L15-B-6]\label{prob:L15-B-6}
设$A=\begin{pmatrix}
                2 & 2 & 0 \\ 8 & 2 & a \\ 0 & 0 & 6
            \end{pmatrix}$相似于对角矩阵，求常数$a$，并求可逆矩阵$P$使得$P^{-1}AP$为对角矩阵.
\end{exercise}
\noindent 来源：\href{https://github.com/yhwu-is/Linear-Algebra-Left-Undone/blob/new/%E8%AE%B2%E4%B9%89/%E4%B8%93%E9%A2%98/15%20%E7%9B%B8%E4%BC%BC%E6%A0%87%E5%87%86%E5%BD%A2%EF%BC%9A%E5%A4%8D%E6%95%B0%E5%9F%9F%E4%B8%8A%E7%9A%84%E5%B0%9D%E8%AF%95%E4%B8%8E%E7%90%86%E8%AE%BA.tex}{L15-B-6}。
\begin{solve}
直接计算$|\lambda E-A|=0$，得到$A$的特征值为$-2,6,6$，因为$A$可对角化，则特征值6对应的特征子空间为2维，即$(6E-A)X=0$的解空间为2维，根据齐次线性方程组一般理论，$r(6E-A)=3-2=1$，即$r\begin{pmatrix}
                4 & -2 & 0 \\ -8 & 4 & -a \\ 0 & 0 & 0
            \end{pmatrix}=1$，显然$a=0$. 接下来求解-2和6对应的特征向量即可，得到可行的解$P=\begin{pmatrix}
                -1 & 1 & 0 \\ 2 & 2 & 0 \\ 0 & 0 & 1
            \end{pmatrix}$.
\end{solve}

\begin{exercise}[补充训练 L15-B-7]\label{prob:L15-B-7}
设$A=(a_{ij})_{n\times n}$是上三角矩阵.
        \begin{enumerate}
            \item 求$A$的全部特征值；

            \item 若$A$主对角元互不相等，证明：$A$与对角阵相似；

            \item 若$n$个主对角元相等且$A$不为对角矩阵，证明：$A$不与对角阵相似.
        \end{enumerate}
\end{exercise}
\noindent 来源：\href{https://github.com/yhwu-is/Linear-Algebra-Left-Undone/blob/new/%E8%AE%B2%E4%B9%89/%E4%B8%93%E9%A2%98/15%20%E7%9B%B8%E4%BC%BC%E6%A0%87%E5%87%86%E5%BD%A2%EF%BC%9A%E5%A4%8D%E6%95%B0%E5%9F%9F%E4%B8%8A%E7%9A%84%E5%B0%9D%E8%AF%95%E4%B8%8E%E7%90%86%E8%AE%BA.tex}{L15-B-7}。
\begin{solve}
\begin{enumerate}
                \item 上三角矩阵特征值就是所有主对角元$a_{ii},\enspace i=1,2,\ldots,n$；

                \item 由（1）知主对角元互不相等表示$A$有$n$个互不相等的特征值，故可对角化；

                \item 由（1）知此时$A$有一个$n$重特征值，记为$c$，则所有特征向量都在$(cE-A)X=0$的解空间中，而$r(cE-A)=1$，故解空间维数为$n-1<n$，故$A$不可对角化.
            \end{enumerate}
\end{solve}

\begin{exercise}[补充训练 L15-B-8]\label{prob:L15-B-8}
已知$\mathbf{R}^3$的一个线性变换
        \[\sigma(x_1,x_2,x_3)=(2x_1-2x_2,-2x_1+x_2-2x_3,-2x_2).\]
        \begin{enumerate}
            \item 求$\sigma$关于自然基$\{e_1,e_2,e_3\}$所对应的矩阵$A$；

            \item 求$\sigma$关于基$\{(1,1,1),(0,1,1),(0,0,1)\}$所对应的矩阵$B$；

            \item 求矩阵$C_1$，使$C_1^{-1}BC_1=A$.
        \end{enumerate}
\end{exercise}
\noindent 来源：\href{https://github.com/yhwu-is/Linear-Algebra-Left-Undone/blob/new/%E8%AE%B2%E4%B9%89/%E4%B8%93%E9%A2%98/15%20%E7%9B%B8%E4%BC%BC%E6%A0%87%E5%87%86%E5%BD%A2%EF%BC%9A%E5%A4%8D%E6%95%B0%E5%9F%9F%E4%B8%8A%E7%9A%84%E5%B0%9D%E8%AF%95%E4%B8%8E%E7%90%86%E8%AE%BA.tex}{L15-B-8}。
\begin{solve}
\begin{enumerate}
                \item $A=\begin{pmatrix}
                              2 & -2 & 0 \\ -2 & 1 & -2 \\ 0 & -2 & 0
                          \end{pmatrix}$；

                \item $B=\begin{pmatrix}
                              0 & -2 & 0 \\ -3 & 1 & -2 \\ 1 & -1 & 2
                          \end{pmatrix}$；

                \item 根据讲义本讲开头的定理（线性变换在不同基下的表示），$C_1$实际上就是基$\alpha_1,\alpha_2,\alpha_3$变为自然基$e_1,e_2,e_3$的过渡矩阵，简单计算可得$\begin{pmatrix}
                              1 & 0 & 0 \\ -1 & 1 & 0 \\ 0 & -1 & 1
                          \end{pmatrix}$.
            \end{enumerate}
\end{solve}

\begin{exercise}[补充训练 L15-B-9]\label{prob:L15-B-9}
设$A=\begin{pmatrix}
                0 & 0 & 1 \\ 0 & 0 & 0 \\ 1 & 0 & 0
            \end{pmatrix},B=\begin{pmatrix}
                1 & 0 & 0 \\ 0 & 1 & 2 \\ 0 & -1 & -2
            \end{pmatrix}$，证明：$A\sim B$，并求可逆矩阵$P$使得$P^{-1}AP=B$.
\end{exercise}
\noindent 来源：\href{https://github.com/yhwu-is/Linear-Algebra-Left-Undone/blob/new/%E8%AE%B2%E4%B9%89/%E4%B8%93%E9%A2%98/15%20%E7%9B%B8%E4%BC%BC%E6%A0%87%E5%87%86%E5%BD%A2%EF%BC%9A%E5%A4%8D%E6%95%B0%E5%9F%9F%E4%B8%8A%E7%9A%84%E5%B0%9D%E8%AF%95%E4%B8%8E%E7%90%86%E8%AE%BA.tex}{L15-B-9}。
\begin{solve}
不难解得$A$和$B$的特征值均为$-1,0,1$，因此它们都与$\begin{pmatrix}
                -1 & 0 & 0 \\ 0 & 0 & 0 \\ 0 & 0 & 1
            \end{pmatrix}$相似，即$A$与$B$相似. 我们也可以解得$P_1=\begin{pmatrix}
                -1 & 0 & 1 \\ 0 & 1 & 0 \\ 1 & 0 & 1
            \end{pmatrix}$和$P_2=\begin{pmatrix}
                0 & 0 & 1 \\ -1 & -2 & 0 \\ 1 & 1 & 0
            \end{pmatrix}$使得$P_1^{-1}AP_1=P_2^{-1}BP_2=\begin{pmatrix}
                -1 & 0 & 0 \\ 0 & 0 & 0 \\ 0 & 0 & 1
            \end{pmatrix}$，故$(P_1P_2^{-1})^{-1}A(P_1P_2^{-1})=B$，即题目要求的$P=P_1P_2^{-1}=\begin{pmatrix}
                1 & -1 & -2 \\ 0 & -1 & -1 \\ 1 & 1 & 2
            \end{pmatrix}$.
\end{solve}

\begin{exercise}[补充训练 L15-B-10]\label{prob:L15-B-10}
已知$A=\begin{pmatrix}
                2 & 0 & 0 \\ 0 & 0 & 1 \\ 0 & 1 & x
            \end{pmatrix}$与$B=\begin{pmatrix}
                2 & 0 & 0 \\ 0 & y & 0 \\ 0 & 0 & -1
            \end{pmatrix}$相似.
        \begin{enumerate}
            \item 求$x$和$y$；

            \item 求一个可逆矩阵$P$，使$P^{-1}AP$为对角矩阵.
        \end{enumerate}
\end{exercise}
\noindent 来源：\href{https://github.com/yhwu-is/Linear-Algebra-Left-Undone/blob/new/%E8%AE%B2%E4%B9%89/%E4%B8%93%E9%A2%98/15%20%E7%9B%B8%E4%BC%BC%E6%A0%87%E5%87%86%E5%BD%A2%EF%BC%9A%E5%A4%8D%E6%95%B0%E5%9F%9F%E4%B8%8A%E7%9A%84%E5%B0%9D%E8%AF%95%E4%B8%8E%E7%90%86%E8%AE%BA.tex}{L15-B-10}。
\begin{solve}
\begin{enumerate}
                \item $A$与$B$相似则特征值相等，我们知道矩阵的对角线元素之和等于特征值之和，行列式等于特征值之积，则它们也一定相等. 根据这一原理非常容易解得$x=0,y=1$.

                \item 事实上$B$就是对角矩阵，实际上这里就是求过渡矩阵$P$使$A$对角化，具体步骤不再赘述，得到$P=\begin{pmatrix}
                              1 & 0 & 0 \\ 0 & 1 & 1 \\ 0 & 1 & -1
                          \end{pmatrix}$.
            \end{enumerate}
\end{solve}

\begin{exercise}[补充训练 L15-B-11]\label{prob:L15-B-11}
设$A=\begin{pmatrix}
                1 & 2 & 0 & 0  & 0 \\ 4 & 3 & 0 & 0 & 0 \\ 0 & 0 & 1 & -3 & 3 \\ 0 & 0 & 3 & -5 & 3 \\
                0 & 0 & 6 & -6 & 4
            \end{pmatrix}$，求$A$的特征值. 若$A$可对角化，求可逆矩阵$P$，使$P^{-1}AP$为对角矩阵.
\end{exercise}
\noindent 来源：\href{https://github.com/yhwu-is/Linear-Algebra-Left-Undone/blob/new/%E8%AE%B2%E4%B9%89/%E4%B8%93%E9%A2%98/15%20%E7%9B%B8%E4%BC%BC%E6%A0%87%E5%87%86%E5%BD%A2%EF%BC%9A%E5%A4%8D%E6%95%B0%E5%9F%9F%E4%B8%8A%E7%9A%84%E5%B0%9D%E8%AF%95%E4%B8%8E%E7%90%86%E8%AE%BA.tex}{L15-B-11}。
\begin{solve}
令$A_1=\begin{pmatrix}
                1 & 2 \\ 4 & 3
            \end{pmatrix}$，$A_2=\begin{pmatrix}
                1 & -3 & 3 \\ 3 & -5 & 3 \\ 6 & -6 & 4
            \end{pmatrix}$，则$A=\begin{pmatrix}
                A_1 & O \\ O & A_2
            \end{pmatrix}$. 首先得到$A_1$和$A_2$的特征值分别为$-1,5$和$-2,-2,4$. 分别求过渡矩阵$P_1$和$P_2$使得$A_1$，$A_2$对角化，容易解得$P_1=\begin{pmatrix}
                -1 & 1 \\ 1 & 2
            \end{pmatrix}$，$P_2=\begin{pmatrix}
                1 & 1 & 1 \\ 1 & 0 & 1 \\ 0 & -1 & 2
            \end{pmatrix}$. 令$P=\begin{pmatrix}
                P_1 & O \\ O & P_2
            \end{pmatrix}$，则
        \[P^{-1}AP=\begin{pmatrix}
                P_1^{-1}A_1P_1 & O \\ O & P_2^{-1}A_2P_2
            \end{pmatrix}=\begin{pmatrix}
                -1 & 0 & 0 & 0 & 0 \\ 0 & 5 & 0 & 0 & 0 \\ 0 & 0 & -2 & 0 & 0 \\ 0 & 0 & 0 & -2 & 0 \\ 0 & 0 & 0 & 0 & 4
            \end{pmatrix},\]
        即$P^{-1}AP=\operatorname{diag}(-1,5,-2,-2,4)$，故$A$与$\operatorname{diag}(-1,5,-2,-2,4)$相似. 总之，$A$的特征值为$-1,5,-2,-2,4$，过渡矩阵
\[P=\begin{pmatrix}
                P_1 & O \\ O & P_2
            \end{pmatrix}=\begin{pmatrix}
                -1 & 1 & 0 & 0 & 0 \\ 1 & 2 & 0 & 0 & 0 \\ 0 & 0 & 1 & 1 & 1 \\ 0 & 0 & 1 & 0 & 1 \\ 0 & 0 & 0 & -1 & 2
            \end{pmatrix}\]
.

        （注：通过本题也可以看出分块对角矩阵的特征值是两个对角块矩阵的并集，过渡矩阵就是两个过渡矩阵按同样分块方式排列得到的矩阵.）
\end{solve}

\begin{exercise}[补充训练 L15-B-12]\label{prob:L15-B-12}
设三阶矩阵$A$的特征值为$\lambda_1=1,\lambda_2=2,\lambda_3=3$，它们对应的特征向量为$\xi_1=(1,1,1)^\mathrm{T}, \xi_2=(1,2,4)^\mathrm{T},\xi_3=(1,3,9)^\mathrm{T}$，又$\beta=(1,1,3)^\mathrm{T}$，计算$A^n\beta$.
\end{exercise}
\noindent 来源：\href{https://github.com/yhwu-is/Linear-Algebra-Left-Undone/blob/new/%E8%AE%B2%E4%B9%89/%E4%B8%93%E9%A2%98/15%20%E7%9B%B8%E4%BC%BC%E6%A0%87%E5%87%86%E5%BD%A2%EF%BC%9A%E5%A4%8D%E6%95%B0%E5%9F%9F%E4%B8%8A%E7%9A%84%E5%B0%9D%E8%AF%95%E4%B8%8E%E7%90%86%E8%AE%BA.tex}{L15-B-12}。
\begin{solve}
回顾对角化过程可知，令$P=(\xi_1,\xi_2,\xi_3)=\begin{pmatrix}
                1 & 1 & 1 \\ 1 & 2 & 3 \\ 1 & 4 & 9
            \end{pmatrix}$，则$P^{-1}AP=\begin{pmatrix}
                1 & 0 & 0 \\ 0 & 2 & 0 \\ 0 & 0 & 3
            \end{pmatrix}$，即$P^{-1}A^nP=\begin{pmatrix}
                1 & 0 & 0 \\ 0 & 2^n & 0 \\ 0 & 0 & 3^n
            \end{pmatrix}$，故$A^n=P\begin{pmatrix}
                1 & 0 & 0 \\ 0 & 2^n & 0 \\ 0 & 0 & 3^n
            \end{pmatrix}P^{-1}$，最后可以计算得到$A^n\beta=\begin{pmatrix}
                2-2^{n+1}+3^n \\ 2-2^{n+2}+3^{n+1} \\ 2-2^{n+3}+3^{n+2}
            \end{pmatrix}$.
\end{solve}

\begin{exercise}[补充训练 L15-B-13]\label{prob:L15-B-13}
设$A=\begin{pmatrix}
                3 & 4 & 0 & 0 \\ 4 & -3 & 0 & 0 \\ 0 & 0 & 2 & 4 \\ 0 & 0 & 0 & 2
            \end{pmatrix}$，求$A^n(n\in\mathbf{N}_+)$.
\end{exercise}
\noindent 来源：\href{https://github.com/yhwu-is/Linear-Algebra-Left-Undone/blob/new/%E8%AE%B2%E4%B9%89/%E4%B8%93%E9%A2%98/15%20%E7%9B%B8%E4%BC%BC%E6%A0%87%E5%87%86%E5%BD%A2%EF%BC%9A%E5%A4%8D%E6%95%B0%E5%9F%9F%E4%B8%8A%E7%9A%84%E5%B0%9D%E8%AF%95%E4%B8%8E%E7%90%86%E8%AE%BA.tex}{L15-B-13}。
\begin{solve}
设$A=\begin{pmatrix}
                A_1 & O \\ O & A_2
            \end{pmatrix}$，其中$A_1=\begin{pmatrix}
                3 & 4 \\ 4 & -3
            \end{pmatrix}$，$A_2=\begin{pmatrix}
                2 & 4 \\ 0 & 2
            \end{pmatrix}$. 回顾矩阵运算进阶中介绍的求幂方法，我们发现$A_2=2E+B$，其中$B=\begin{pmatrix}
                0 & 4 \\ 0 & 0
            \end{pmatrix}$，且$B^2=O$，因此$A_2^n=(2E+B)^n=2^nE+n2^{n-1}B=\begin{pmatrix}
                2^n & n2^{n+1} \\ 0 & 2^n
            \end{pmatrix}$. 再看$A_1$，我们用对角化方法求幂，容易求得$P=\begin{pmatrix}
                2 & -1 \\ 1 & 2
            \end{pmatrix}$使得$P^{-1}A_1P=\begin{pmatrix}
                5 & 0 \\ 0 & -5
            \end{pmatrix}$，故
\[A_1^n=P\begin{pmatrix}
                5^n & 0 \\ 0 & (-5)^n
            \end{pmatrix}P^{-1}=5^{n-1}\begin{pmatrix}
                4+(-1)^n & 2+2(-1)^{n+1} \\ 2+2(-1)^{n+1} & 1+4(-1)^n
            \end{pmatrix}\]
，最后结合上面的讨论得到$A^n=\begin{pmatrix}
                A_1^n & O \\ O & A_2^n
            \end{pmatrix}$.
\end{solve}

\begin{exercise}[补充训练 L15-B-14]\label{prob:L15-B-14}
已知三阶矩阵$A$和三元列向量$X$，使得向量组$X,AX,A^2X$线性无关，且满足
        \[A^3X=3AX-2A^2X.\]
        \begin{enumerate}
            \item 记$P=(X,AX,A^2X)$，求三阶矩阵$B$使得$A=PBP^{-1}$；

            \item 计算行列式$|A+E|$.
        \end{enumerate}
\end{exercise}
\noindent 来源：\href{https://github.com/yhwu-is/Linear-Algebra-Left-Undone/blob/new/%E8%AE%B2%E4%B9%89/%E4%B8%93%E9%A2%98/15%20%E7%9B%B8%E4%BC%BC%E6%A0%87%E5%87%86%E5%BD%A2%EF%BC%9A%E5%A4%8D%E6%95%B0%E5%9F%9F%E4%B8%8A%E7%9A%84%E5%B0%9D%E8%AF%95%E4%B8%8E%E7%90%86%E8%AE%BA.tex}{L15-B-14}。
\begin{solve}
\begin{enumerate}
              \item 由题意，$PB=AP$，利用分块矩阵乘法可知$PB=AP=(AX,A^2X,A^3X)=(AX,A^2X,3AX-2A^2X)$（注意$(AX,A^2X,A^3X)B=(AXB,A^2XB,A^3XB)$是错误的，因为这是$1\times 3$分块和$1\times 1$分块相乘，显然不能这么乘）.

                    则$PB=(AX,A^2X,3AX-2A^2X)=(X,AX,A^X)\begin{pmatrix}
                            0 & 0 & 0 \\ 1 & 0 & 3 \\ 0 & 1 & -2
                        \end{pmatrix}$. 等号两边左乘$P^{-1}$有$B=\begin{pmatrix}
                            0 & 0 & 0 \\ 1 & 0 & 3 \\ 0 & 1 & -2
                        \end{pmatrix}$.

              \item 由于$A$与$B$相似，故有相同的特征值，易求得$B$的特征值为$0,-3,1$，则$A$的特征值也为$0,-3,1$，则$A+E$的特征值根据正文例题可知为$1,-2,2$，则$A+E$的行列式为其特征值之积，即为$-4$.
          \end{enumerate}
\end{solve}

\begin{exercise}[补充训练 L15-C-2]\label{prob:L15-C-2}
设$B=\alpha\alpha^\mathrm{T}$，其中$\alpha=(a_1,\ldots,a_n)^\mathrm{T}\neq 0\enspace(a_i\in\mathbf{R},\enspace i=1,2,\ldots,n)$.
        \begin{enumerate}
            \item 证明：$B^k=tB$，其中$k$为正整数，$t$为常数，并求$t$；

            \item 求可逆阵$P$使得$P^{-1}BP$为对角矩阵，并写出该对角矩阵.
        \end{enumerate}
\end{exercise}
\noindent 来源：\href{https://github.com/yhwu-is/Linear-Algebra-Left-Undone/blob/new/%E8%AE%B2%E4%B9%89/%E4%B8%93%E9%A2%98/15%20%E7%9B%B8%E4%BC%BC%E6%A0%87%E5%87%86%E5%BD%A2%EF%BC%9A%E5%A4%8D%E6%95%B0%E5%9F%9F%E4%B8%8A%E7%9A%84%E5%B0%9D%E8%AF%95%E4%B8%8E%E7%90%86%E8%AE%BA.tex}{L15-C-2}。
\begin{solve}
\begin{enumerate}
                \item $B^2=\alpha(\alpha^\mathrm{T}\alpha)\alpha^\mathrm{T}=m\alpha\alpha^\mathrm{T}=mB$，其中$m=\alpha^\mathrm{T}\alpha=a_1^2+\cdot+a_n^2$，且易知这就是$\operatorname{tr}(B)$. 用归纳法可知$B^k=m^{k-1}B$，令$t=m^{k-1}$则得证，其中$t=(\operatorname{tr}(B))^{k-1}$.

                \item 这是经典的秩1矩阵，由于$B$的秩为1，则$BX=0$的解空间维数为$n-1$，因此0是$B$的$n-1$重特征值. 而特征值之和等于矩阵的迹. 则剩下一个一重特征值为$\operatorname{tr}(B)$（由题意迹不为0）.

                      然后求特征向量，这里需要一些观察. 特征值0对应的特征向量即为$BX=0$的解，设解为$(x_1,\ldots,x_n)$，不失一般性地，设$a_1\neq 0$，由于$B$的秩为1，因此将$BX=0$展开为线性方程组后，每一行（除去全0的）代表的方程是完全一致的（因为它们成比例），因此我们只需考虑第一行方程
                      \[a_1^2x_1+a_1a_2x_2+\cdots+a_1a_nx_n=0,\]
                      解得$n-1$个线性无关特征向量为
                      \[X_1=(-a_2,a_1,0,\ldots,0),X_2=(-a_3,0,a_1,0,\ldots,0),\ldots,X_{n-1}=(-a_n,0,\ldots,0,a_1).\]
                      对于特征值$\operatorname{tr}(B)=\alpha^\mathrm{T}\alpha$，我们可以设特征向量为$X$，则$(\operatorname{tr}(B)E-B)X=(\alpha^\mathrm{T}\alpha E-\alpha\alpha^\mathrm{T})X=0$，即$\alpha^\mathrm{T}X=\alpha\alpha^\mathrm{T}X$，我们可以观察发现$\alpha$就是一个解，则$X_n=\alpha$是其特征向量. 则
\[P=\begin{pmatrix}
                              -a_2 & -a_3 & \cdots & -a_n & a_1 \\ a_1 & 0 & \cdots & 0 & a_2 \\ 0 & a_1 & \cdots & 0 & a_3 \\ \vdots & \vdots & \ddots & \vdots & \vdots \\ 0 & 0 & \cdots & a_1 & a_n
                          \end{pmatrix}\]
，且对角矩阵为$\operatorname{diag}(0,\ldots,0,\alpha^\mathrm{T}\alpha)$.
            \end{enumerate}
\end{solve}

\begin{exercise}[补充训练 L15-C-3]\label{prob:L15-C-3}
（秩为1的矩阵）设$n$阶矩阵$A$的元素均为1.
        \begin{enumerate}
            \item 求$A$的特征值，并求矩阵$P$使得$P^{-1}AP$为对角矩阵；

            \item 若$f(x)$是$x$的$m$次多项式，且常数项为0，证明：存在$k\in\mathbf{R}$使得$f(A)=kA$，并求出$k$；

            \item 设$B$是$n$阶实对称矩阵，每行元素之和都为$b$，若$b$是$f(\lambda)=|\lambda E-B|$的单根，求$B$属于$b$的特征向量；当$f(\lambda)=(\lambda-b)g(\lambda)$时（其中$f(B)=0$），证明：$g(B)=kA$，其中$k$为常数，$A$为元素全部为1的$n$阶矩阵.
        \end{enumerate}
\end{exercise}
\noindent 来源：\href{https://github.com/yhwu-is/Linear-Algebra-Left-Undone/blob/new/%E8%AE%B2%E4%B9%89/%E4%B8%93%E9%A2%98/15%20%E7%9B%B8%E4%BC%BC%E6%A0%87%E5%87%86%E5%BD%A2%EF%BC%9A%E5%A4%8D%E6%95%B0%E5%9F%9F%E4%B8%8A%E7%9A%84%E5%B0%9D%E8%AF%95%E4%B8%8E%E7%90%86%E8%AE%BA.tex}{L15-C-3}。
\begin{solve}
\begin{enumerate}
                \item 由题意知$A$的秩为1，因此特征值0对应的特征子空间为$AX=0$的解空间，维数为$n-1$，故特征值0的重数至少为$n-1$. 剩下可能还有一个特征值，我们知道特征值之和等于对角线元素之和，因此最后一个特征值等于$n$，$n$对应的特征子空间维数至少为1，结合0对应特征子空间维数为$n-1$可知$n$对应的特征子空间维数就是1，二者可以张成整个$\mathbf{R}^n$，因此可以对角化.

                      下面求过渡矩阵，$n$对应的特征向量只需解方程$(nE-A)X=0$，即得$X_1=(1,1,\ldots,1)^\mathrm{T}$. 0对应的$n-1$个线性无关特征向量就是$AX=0$的基础解系，即$X_2=(-1,1,0,\ldots,0)^\mathrm{T},X_3=(-1,0,1,0,\ldots,0)^\mathrm{T},\ldots,X_n=(-1,0,\ldots,0,1)^\mathrm{T}$，令
\[P=(X_1,X_2,\ldots,X_n)=\begin{pmatrix}
                              1 & -1 & -1 & \cdots & -1 \\ 1 & 1 & 0 & \cdots & 0 \\ 1 & 0 & 1 & \cdots & 0 \\ \vdots & \vdots & \vdots & \ddots & \vdots \\ 1 & 0 & 0 & \cdots & 1
                          \end{pmatrix}\]
，则有$P^{-1}AP=\Lambda=\operatorname{diag}(n,0,\ldots,0)$.

                \item 由上一小问可知$A=P\Lambda P^{-1}$，则$A^k=P\Lambda^kP^{-1}$，故
                      \[A^k=P\Lambda^kP^{-1}=P\operatorname{diag}(n^k,0,\ldots,0)P^{-1},\]
                      故
                      \begin{align*}
                          f(A) & =a_mA^m+a_{m-1}A^{m-1}+\cdots+a_1A                                                                                                   \\
                               & =P(a_m\Lambda^m+a_{m-1}\Lambda^{m-1}+\cdots+a_1\Lambda)P^{-1}                                                               \\
                               & =P\begin{pmatrix}
                                       a_mn^m+a_{m-1}n^{m-1}+\cdots+a_1n & 0 & \cdots & 0 \\ 0 & 0 & \cdots & 0 \\ \vdots & \vdots & \ddots & \vdots \\ 0 & 0 & \cdots & 0
                                   \end{pmatrix}P^{-1} \\
                               & =\begin{pmatrix}
                                      f(n) & 0 & \cdots & 0 \\ 0 & 0 & \cdots & 0 \\ \vdots & \vdots & \ddots & \vdots \\ 0 & 0 & \cdots & 0
                                  \end{pmatrix}=\dfrac{f(n)}{n}P\begin{pmatrix}
                                                                    n & 0 & \cdots & 0 \\ 0 & 0 & \cdots & 0 \\ \vdots & \vdots & \ddots & \vdots \\ 0 & 0 & \cdots & 0
                                                                \end{pmatrix}P^{-1}    \\
                               & =\dfrac{f(n)}{n}P\Lambda P^{-1}=kA.
                      \end{align*}
                      其中$k=\dfrac{f(n)}{n}$. 故$f(A)=kA$.

                \item 显然$X=(1,1,x\cdots,1)^\mathrm{T}$是$B$关于特征值$b$的特征向量，由于$b$是特征多项式的单根，则特征子空间维数也为1，故特征向量就是$kX$（其中$k$为非零常数）.

                      又$f(B)=(B-bE)g(B)=O$，则$Bg(B)=bg(B)$. 令$g(B)$的列向量组为$\alpha_1,\ldots,\alpha_n$，则上式可以写为$B(\alpha_1,\ldots,\alpha_n)=(b\alpha_1,\ldots,b\alpha_n)$，即$B\alpha_i=b\alpha_i$，即$\alpha_i$是$B$关于特征值$b$的特征向量，故$\alpha_i=k_iX=(k_i,k_i,\ldots,k_i)^{\mathrm{T}}$，即
\[g(B)=\begin{pmatrix}
                              k_1 & k_2 & \cdots & k_n \\ k_1 & k_2 & \cdots & k_n \\ \vdots & \vdots & \ddots & \vdots \\ k_1 & k_2 & \cdots & k_n
                          \end{pmatrix}\]
. 由于$B$是实对称矩阵，故$g(B)$也是实对称矩阵（很容易验证实对称矩阵经过幂次、加法运算后仍是实对称矩阵），因此$g(B)=(g(B))^\mathrm{T}$，故$k_1=k_2=\cdots=k_n=k$，即$g(B)=kA$.
            \end{enumerate}
\end{solve}

\begin{exercise}[补充训练 L15-C-4]\label{prob:L15-C-4}
设$A,B$为$n$阶矩阵，且$A+B=AB$，求证：
        \begin{enumerate}
            \item $A,B$的特征向量是公共的；

            \item $A$相似于对角矩阵当且仅当$B$相似于对角矩阵；

            \item $r(A)=r(B)$.
        \end{enumerate}
\end{exercise}
\noindent 来源：\href{https://github.com/yhwu-is/Linear-Algebra-Left-Undone/blob/new/%E8%AE%B2%E4%B9%89/%E4%B8%93%E9%A2%98/15%20%E7%9B%B8%E4%BC%BC%E6%A0%87%E5%87%86%E5%BD%A2%EF%BC%9A%E5%A4%8D%E6%95%B0%E5%9F%9F%E4%B8%8A%E7%9A%84%E5%B0%9D%E8%AF%95%E4%B8%8E%E7%90%86%E8%AE%BA.tex}{L15-C-4}。
\begin{solve}
\begin{enumerate}
                \item 设$\alpha$是$B$的特征向量，对应的特征值为$\lambda$，即$B\alpha=\lambda\alpha$，则$A\alpha+B\alpha-AB\alpha=0$，即$A\alpha+\lambda\alpha-A\lambda\alpha=0$，即$(1-\lambda)A\alpha=-\lambda\alpha$. 若$\lambda=1$，则$\alpha=0$，与它是特征向量矛盾，故$\lambda\neq 1$，从而$A\alpha=\dfrac{-\lambda}{1-\lambda}\alpha$，这就说明了$B$的特征向量就是$A$的特征向量.

                      由$A+B=AB$可知，$(A-E)(B-E)=E$，故$A-E$可逆，且$(A-E)^{-1}=B-E$，从而$E=(B-E)(A-E)=BA-A-B+E$，故$BA=A+B$，同上理可得$A$的特征向量也是$B$的，得证.

                \item \begin{enumerate}
                          \item 必要性：若$A$可对角化，则$\mathbf{R}^n$有一组由$A$的特征向量张成的基，由（1）知$B$与$A$特征向量一致，故$\mathbf{R}^n$也有一组由$B$的特征向量张成的基，即$B$可对角化.

                          \item 充分性与必要性类似.
                      \end{enumerate}

                \item 由$A+B=AB$可知$A=(A-E)B$且$B=A(B-E)$，由此可得$r(A)\leqslant r(B)\leqslant r(A)$，故$r(A)=r(B)$.
            \end{enumerate}
\end{solve}

\begin{exercise}[补充训练 L15-C-5]\label{prob:L15-C-5}
设$A,B\in \mathbf{M}_n(\mathbf{R})$，证明$A$与$B$在$\mathbf{R}$上相似当且仅当在$\mathbf{C}$上相似.

        {\kaishu 注：实际上相似这一性质与数域无关，本题是这一结论的特例.}
\end{exercise}
\noindent 来源：\href{https://github.com/yhwu-is/Linear-Algebra-Left-Undone/blob/new/%E8%AE%B2%E4%B9%89/%E4%B8%93%E9%A2%98/15%20%E7%9B%B8%E4%BC%BC%E6%A0%87%E5%87%86%E5%BD%A2%EF%BC%9A%E5%A4%8D%E6%95%B0%E5%9F%9F%E4%B8%8A%E7%9A%84%E5%B0%9D%E8%AF%95%E4%B8%8E%E7%90%86%E8%AE%BA.tex}{L15-C-5}。
\begin{solve}
\begin{enumerate}
                \item 必要性：$A$和$B$在实数域上相似，则存在可逆实矩阵$C$使得$B=C^{-1}AC$，将$A,B,C$视为复数域上的矩阵，则可知$A$和$B$在复数域上也相似.

                \item 充分性：若$A$和$B$在复数域上相似，则存在可逆复矩阵$P=P_1+\mathrm{i} P_2$使得$B=P^{-1}AP$，其中$P_1,P_2$是实矩阵，即$AP=PB$，从而$A(P_1+\mathrm{i} P_2)=(P_1+\mathrm{i} P_2)B$，比较实部虚部可知$AP_1=P_1B$且$AP_2=P_2B$. 构造实系数多项式$f(t)=|P_1+tP_2|$，显然$f(\mathrm{i})=|P|\neq 0$，则$f(t)$是非零多项式，所以$f(t)$的根只有有限多个，从而存在$t_0\in\mathbf{R}$使得$f(t_0)\neq 0$，则$|P_1+t_0P_2|\neq 0$，从而$P_1+t_0P_2$可逆，从而$A(P_1+t_0P_2)=AP_1+t_0AP_2=P_1B+t_0P_2B=(P_1+t_0P_2)B$，即$(P_1+t_0P_2)^{-1}A(P_1+t_0P_2)=B$，即$A$和$B$在实数域上相似.
            \end{enumerate}
\end{solve}

\begin{exercise}[补充训练 L15-C-6]\label{prob:L15-C-6}
设 $A,B\in\mathbf F^{n\times n}$。
\begin{enumerate}
\item 若 $A$ 在 $\mathbf F$ 中有 $n$ 个不同特征值，证明：$AB=BA$ 当且仅当 $A$ 的每个特征向量也是 $B$ 的特征向量。
\item 若 $A,B$ 都在 $\mathbf F$ 上可对角化且 $AB=BA$，证明它们可以用同一个过渡矩阵对角化。
\end{enumerate}
\end{exercise}
\noindent 来源：\href{https://github.com/yhwu-is/Linear-Algebra-Left-Undone/blob/new/%E8%AE%B2%E4%B9%89/%E4%B8%93%E9%A2%98/15%20%E7%9B%B8%E4%BC%BC%E6%A0%87%E5%87%86%E5%BD%A2%EF%BC%9A%E5%A4%8D%E6%95%B0%E5%9F%9F%E4%B8%8A%E7%9A%84%E5%B0%9D%E8%AF%95%E4%B8%8E%E7%90%86%E8%AE%BA.tex}{L15-C-6}。
\begin{solve}
\begin{enumerate}
\item 若 $AB=BA$，对 $Ax=\lambda x$ 有 $A(Bx)=\lambda Bx$，故 $B$ 保持 $A$ 的每个特征子空间。由于各特征子空间为一维，$Bx=\mu x$，所以每个非零特征向量 $x$ 也是 $B$ 的特征向量，其中允许 $\mu=0$。反过来，取 $A$ 的特征向量基，它也由 $B$ 的特征向量组成，两矩阵在此基下都是对角矩阵，因而交换。
\item 记 $A$ 的互异特征值为 $\lambda_1,\ldots,\lambda_s$，则 $\mathbf F^n=\bigoplus_i E_{\lambda_i}(A)$，每个直和分量均为 $B$ 的不变子空间。因 $B$ 可对角化，其最小多项式分裂为互异一次因子的乘积；每个限制 $B|_{E_{\lambda_i}(A)}$ 的最小多项式整除它，故该限制也可对角化。在各分量内选择 $B$ 的特征向量基，再将这些基合并，就得到 $A,B$ 的共同特征向量基。
\end{enumerate}
\end{solve}

\begin{exercise}[补充训练 L15-C-7]\label{prob:L15-C-7}
设 $A,B\in\mathbf F^{n\times n}$，$A$ 在 $\mathbf F$ 中有 $n$ 个不同的特征值，且 $AB=BA$。证明：存在唯一的次数不超过 $n-1$ 的多项式 $f\in\mathbf F[x]$，使 $B=f(A)$。
\end{exercise}
\noindent 来源：\href{https://github.com/yhwu-is/Linear-Algebra-Left-Undone/blob/new/%E8%AE%B2%E4%B9%89/%E4%B8%93%E9%A2%98/15%20%E7%9B%B8%E4%BC%BC%E6%A0%87%E5%87%86%E5%BD%A2%EF%BC%9A%E5%A4%8D%E6%95%B0%E5%9F%9F%E4%B8%8A%E7%9A%84%E5%B0%9D%E8%AF%95%E4%B8%8E%E7%90%86%E8%AE%BA.tex}{L15-C-7}。
\begin{solve}
由习题~\ref{prob:L15-C-6}，存在可逆矩阵 $P$ 使 $P^{-1}AP=\operatorname{diag}(\lambda_1,\ldots,\lambda_n)$、$P^{-1}BP=\operatorname{diag}(\mu_1,\ldots,\mu_n)$。取 Lagrange 插值多项式
\[f(x)=\sum_{i=1}^n\mu_i\prod_{j\ne i}\frac{x-\lambda_j}{\lambda_i-\lambda_j},\]
则次数不超过 $n-1$，且 $f(\lambda_i)=\mu_i$，故 $P^{-1}f(A)P=\operatorname{diag}(\mu_1,\ldots,\mu_n)$，即 $f(A)=B$。若还有 $g$ 满足同样的次数界和矩阵等式，则 $f-g$ 在 $n$ 个不同的 $\lambda_i$ 上为零，而次数小于 $n$，所以 $f=g$。
\end{solve}

\begin{exercise}[补充训练 L15-C-8]\label{prob:L15-C-8}
设$T\in \mathcal{L}(V)$，$\lambda\in\mathbf{F}$. 证明：对$V$的每个使得$T$有上三角矩阵的基，$\lambda$出现在$T$的矩阵的对角线上的次数等于$\lambda$作为$T$的特征值的重数.
\end{exercise}
\noindent 来源：\href{https://github.com/yhwu-is/Linear-Algebra-Left-Undone/blob/new/%E8%AE%B2%E4%B9%89/%E4%B8%93%E9%A2%98/15%20%E7%9B%B8%E4%BC%BC%E6%A0%87%E5%87%86%E5%BD%A2%EF%BC%9A%E5%A4%8D%E6%95%B0%E5%9F%9F%E4%B8%8A%E7%9A%84%E5%B0%9D%E8%AF%95%E4%B8%8E%E7%90%86%E8%AE%BA.tex}{L15-C-8}。
\begin{solve}
设 $T$ 在题设基下的上三角矩阵为 $A=(a_{ij})$。用独立变量 $t$ 表示特征多项式，则
\[\chi_T(t)=\det(tI-A)=\prod_{j=1}^n(t-a_{jj}).\]
因此固定的 $\lambda$ 在对角线上出现几次，因子 $t-\lambda$ 就在该乘积中出现几次，这恰是 $\lambda$ 的代数重数。这里不能把代数重数换成特征子空间的维数（几何重数）；例如 $\begin{pmatrix}\lambda&1\\0&\lambda\end{pmatrix}$ 的代数重数为 $2$，几何重数只有 $1$。
\end{solve}

\begin{exercise}[补充训练 L15-C-9]\label{prob:L15-C-9}
证明：$A$为幂零矩阵$\iff \forall k \in \mathbf{N}_+,\enspace\operatorname{tr}(A^k)=0$.
\end{exercise}
\noindent 来源：\href{https://github.com/yhwu-is/Linear-Algebra-Left-Undone/blob/new/%E8%AE%B2%E4%B9%89/%E4%B8%93%E9%A2%98/15%20%E7%9B%B8%E4%BC%BC%E6%A0%87%E5%87%86%E5%BD%A2%EF%BC%9A%E5%A4%8D%E6%95%B0%E5%9F%9F%E4%B8%8A%E7%9A%84%E5%B0%9D%E8%AF%95%E4%B8%8E%E7%90%86%E8%AE%BA.tex}{L15-C-9}。
\begin{solve}
在复数域上考察 $A$ 的特征值，不改变矩阵幂或迹。若 $A$ 幂零，则其所有特征值为零，因而每个 $A^k$ 的迹为零。
反过来，设所有正整数 $k$ 都满足 $\operatorname{tr}(A^k)=0$。若存在非零特征值，设互异非零特征值为 $\lambda_1,\ldots,\lambda_r$，代数重数为正整数 $m_1,\ldots,m_r$。由复矩阵可上三角化，有
\[\sum_{j=1}^r m_j\lambda_j^k=0,\qquad k=1,\ldots,r.\]
这一关于 $m_j$ 的线性方程组的系数行列式为
\[\left(\prod_{j=1}^r\lambda_j\right)\prod_{1\le i<j\le r}(\lambda_j-\lambda_i)\ne0,\]
故所有 $m_j=0$，矛盾。因此全部特征值为零，特征多项式为 $t^n$，由 Cayley--Hamilton 定理得 $A^n=O$，即 $A$ 幂零。
\end{solve}

\begin{exercise}[补充训练 L15-例-1]\label{prob:L15-例-1}
证明$r$阶上三角矩阵$(r>1)$
    \[J_0=\begin{pmatrix}
            \lambda_0 & 1         &        &           \\
                      & \lambda_0 & \ddots &           \\
                      &           & \ddots & 1         \\
                      &           &        & \lambda_0
        \end{pmatrix}\]
    不与对角阵相似.
\end{exercise}
\noindent 来源：\href{https://github.com/yhwu-is/Linear-Algebra-Left-Undone/blob/new/%E8%AE%B2%E4%B9%89/%E4%B8%93%E9%A2%98/15%20%E7%9B%B8%E4%BC%BC%E6%A0%87%E5%87%86%E5%BD%A2%EF%BC%9A%E5%A4%8D%E6%95%B0%E5%9F%9F%E4%B8%8A%E7%9A%84%E5%B0%9D%E8%AF%95%E4%B8%8E%E7%90%86%E8%AE%BA.tex}{L15-例-1}。
\begin{solve}
首先求出特征多项式为 $f(\lambda) = |\lambda E - J_0| = (\lambda - \lambda_0)^r$，因此 $J_0$ 只有一个特征值 $\lambda_0$，且代数重数为 $r$.

    接下来求几何重数，即 $J_0X = \lambda_0X$ 的解空间维数，即 $(\lambda_0 E - J_0)X = 0$的解空间维数，事实上由于 $r(\lambda_0 E - J_0) = r - 1$，因此解空间维数为 $r - (r - 1) = 1$，即几何重数为 $1 < r$，因此不可对角化.
\end{solve}

\begin{exercise}[补充训练 L15-例-2]\label{prob:L15-例-2}
求矩阵
    \[A=\begin{pmatrix}
            0  & -1 & 1 \\
            -1 & 0  & 1 \\
            1  & 1  & 0
        \end{pmatrix}\]
    的所有特征值，对应的特征子空间，以及与 $A$ 相似的一个对角矩阵.
\end{exercise}
\noindent 来源：\href{https://github.com/yhwu-is/Linear-Algebra-Left-Undone/blob/new/%E8%AE%B2%E4%B9%89/%E4%B8%93%E9%A2%98/15%20%E7%9B%B8%E4%BC%BC%E6%A0%87%E5%87%86%E5%BD%A2%EF%BC%9A%E5%A4%8D%E6%95%B0%E5%9F%9F%E4%B8%8A%E7%9A%84%E5%B0%9D%E8%AF%95%E4%B8%8E%E7%90%86%E8%AE%BA.tex}{L15-例-2}。
\begin{solve}
对于求解矩阵的对角化问题，首先求出其特征多项式（具体步骤不展开，实际上就是三阶行列式的计算，可以使用按行（列）展开、公式法或者初等变换化为三角矩阵等方法）$f(\lambda)=|\lambda E-A|=(\lambda-1)^2(\lambda+2)$，令$f(\lambda)=0$，解得特征值为 $\lambda_1=\lambda_2=1,\lambda_3=-2$.

    接下来求解特征向量和特征子空间，即求解$(E-A)x=0$和解$(-2E-A)x=0$，得到特征值1对应的特征子空间为$\operatorname{span}((-1,1,0)^{\mathrm{T}},(1,0,1)^{\mathrm{T}})$，特征值-2对应的特征子空间为$\operatorname{span}((-1,-1,1)^{\mathrm{T}})$.

    与$A$相似的对角矩阵实际上就是特征值排列在对角线上的结果，即 $\operatorname{diag}(1,1,-2)$.
\end{solve}

\begin{exercise}[补充训练 L15-例-4]\label{prob:L15-例-4}
已知$A=\begin{pmatrix}
            0 & \dfrac{1}{2}  & \dfrac{1}{2} \\[2ex]
            1 & -\dfrac{1}{2} & \dfrac{1}{2} \\[2ex]
            1 & -\dfrac{1}{2} & \dfrac{1}{2}
        \end{pmatrix}$，求$A^n$.
\end{exercise}
\noindent 来源：\href{https://github.com/yhwu-is/Linear-Algebra-Left-Undone/blob/new/%E8%AE%B2%E4%B9%89/%E4%B8%93%E9%A2%98/15%20%E7%9B%B8%E4%BC%BC%E6%A0%87%E5%87%86%E5%BD%A2%EF%BC%9A%E5%A4%8D%E6%95%B0%E5%9F%9F%E4%B8%8A%E7%9A%84%E5%B0%9D%E8%AF%95%E4%B8%8E%E7%90%86%E8%AE%BA.tex}{L15-例-4}。
\begin{solve}
首先求出$A$的特征多项式$f(\lambda)=|\lambda E-A|=\lambda(\lambda-1)(\lambda+1)$，令$f(\lambda)=0$，解得特征值为 $\lambda_1=0,\lambda_2=1,\lambda_3=-1$.

    接下来求解特征向量和特征子空间，实际上就是求解$(0E-A)x=0,(-E-A)x=0,(E-A)x=0$，得到特征向量为
    \[\eta_1=\begin{pmatrix}
            1 \\ 1 \\ -1
        \end{pmatrix},\enspace \eta_2=\begin{pmatrix}
            1 \\ 1 \\ 1
        \end{pmatrix},\enspace \eta_3=\begin{pmatrix}
            1 \\ -1 \\ -1
        \end{pmatrix}.\]
    所以记$P=(\eta_1,\eta_2,\eta_3)$，则$A=P\operatorname{diag}(0,1,-1)P^{-1}$，因此
    \[A^n=P\operatorname{diag}(0^n,1^n,(-1)^n)P^{-1},\]
    进一步计算得到
    \[A^n=\frac{1}{2}\begin{pmatrix}
            1+(-1)^n     & (-1)^{n+1} & 1 \\
            1+(-1)^{n+1} & (-1)^n     & 1 \\
            1+(-1)^{n+1} & (-1)^n     & 1
        \end{pmatrix}.\]
\end{solve}

\begin{exercise}[补充训练 L15-例-5]\label{prob:L15-例-5}
线性变换 $T : \mathbf{R}^3 \to \mathbf{R}^3$ 的定义是：
    \[T(x_1,x_2,x_3)=(4x_1+x_3,2x_1+3x_2+2x_3,x_1+4x_3)\]
    \begin{enumerate}
        \item 求出$T$的特征多项式及特征值；

        \item 判断$T$是否可对角化，并给出理由.
    \end{enumerate}
\end{exercise}
\noindent 来源：\href{https://github.com/yhwu-is/Linear-Algebra-Left-Undone/blob/new/%E8%AE%B2%E4%B9%89/%E4%B8%93%E9%A2%98/15%20%E7%9B%B8%E4%BC%BC%E6%A0%87%E5%87%86%E5%BD%A2%EF%BC%9A%E5%A4%8D%E6%95%B0%E5%9F%9F%E4%B8%8A%E7%9A%84%E5%B0%9D%E8%AF%95%E4%B8%8E%E7%90%86%E8%AE%BA.tex}{L15-例-5}。
\begin{solve}
\begin{enumerate}
        \item 我们知道，线性变换与其在任意一组基下的矩阵有相同的特征值和特征多项式，因此我们可以任意选取一组基（为方便我们取$\mathbf{R}^3$的自然基，然后求出$T$在这组基下的矩阵，即
              \[A=\begin{pmatrix}
                      4 & 0 & 1 \\
                      2 & 3 & 2 \\
                      1 & 0 & 4
                  \end{pmatrix}\]
              故$T$的特征多项式为$f(\lambda)=|A-\lambda E|=\begin{vmatrix}
                      4-\lambda & 0         & 1         \\
                      2         & 3-\lambda & 2         \\
                      1         & 0         & 4-\lambda
                  \end{vmatrix}=-(\lambda-3)^2(\lambda-5)$，$T$的特征值为$f(\lambda)=0$的解，即为$\lambda_1=\lambda_2=3,\lambda_3=5$.

        \item 对于$\lambda=3$，我们求解$(T-3I)\alpha$得到特征向量（其中令$\alpha=(x_1,x_2,x_3)$），解得$\alpha_1=(1,0,-1)^\mathrm{T},\alpha_2=(0,1,0)^\mathrm{T}$，因此$T$在$\lambda=3$时的特征子空间为$V_3=\operatorname{span}(\alpha_1,\alpha_2)$，维数为2，因此$\lambda=3$时代数重数与几何重数相等. 而$\lambda=5$是一重特征值，因此也必有代数重数与几何重数相等（因为一定有特征向量，那么特征子空间维数至少为1，再结合代数重数大于等于几何重数可知，特征子空间维数恰为1）. 因此$T$可对角化.
    \end{enumerate}
\end{solve}

\begin{exercise}[补充训练 L15-例-6]\label{prob:L15-例-6}
设$\alpha$和$\beta$是$\mathbf{R}^n\enspace (n>1)$中两个列向量，$A=\alpha\beta^\mathrm{T}\neq O$.
    \begin{enumerate}
        \item 求$A$的特征值；

        \item 证明：$\alpha^\mathrm{T}\beta=0\iff A$不可对角化.
    \end{enumerate}
\end{exercise}
\noindent 来源：\href{https://github.com/yhwu-is/Linear-Algebra-Left-Undone/blob/new/%E8%AE%B2%E4%B9%89/%E4%B8%93%E9%A2%98/15%20%E7%9B%B8%E4%BC%BC%E6%A0%87%E5%87%86%E5%BD%A2%EF%BC%9A%E5%A4%8D%E6%95%B0%E5%9F%9F%E4%B8%8A%E7%9A%84%E5%B0%9D%E8%AF%95%E4%B8%8E%E7%90%86%E8%AE%BA.tex}{L15-例-6}。
\begin{solve}
\begin{enumerate}
        \item 我们知道，$r(A)\leqslant\min{\{r(\alpha),r(\beta)\}}=1$，并且$A\neq O$因此$r(A)>0$，故$A$的秩为1. 而$n>1$，因此$A$一定不可逆，故0一定是$A$的特征值，且对应的特征子空间维数为$AX=0$的解空间维数，即为$n-1$.

              由此我们知道$A$最多有两个特征值，因为0的代数重数（即作为$n$次特征多项式的零点次数）必然大于等于其几何重数$n-1$，当其代数重数为$n-1$时可能还有一个一重特征值. 我们利用特征值之和等于$A$的迹来找出可能的第二个特征值. 我们设$\alpha=(a_1,a_2,\ldots,a_n)^\mathrm{T},\beta=(b_1,b_2,\ldots,b_n)^\mathrm{T}$，则
\[A=\alpha\beta^\mathrm{T}=\begin{pmatrix}
                      a_1b_1 & a_1b_2 & \cdots & a_1b_n \\
                      a_2b_1 & a_2b_2 & \cdots & a_2b_n \\
                      \vdots & \vdots & \ddots & \vdots \\
                      a_nb_1 & a_nb_2 & \cdots & a_nb_n
                  \end{pmatrix}\]
，因此$A$的迹为$\sum\limits_{i=1}^na_ib_i=\alpha^\mathrm{T}\beta=\sum\limits_{i=1}^n\lambda_i$，其中$\lambda_i$为$A$的特征值. 若$\alpha^\mathrm{T}\beta\neq 0$，则$\lambda_i=0,\enspace i=1,\ldots,n-1$，$\lambda_n=\alpha^\mathrm{T}\beta$. 若$\alpha^\mathrm{T}\beta=0$，则$A$的所有特征值均为0.

        \item 由上一小问可知，若$\alpha^\mathrm{T}\beta=0$即$A$的全部特征值为0，因此只有一个$n-1$维的特征子空间，故特征子空间直和不等于$V$，故不可对角化.

              反之，若$A$不可对角化，我们用反证法. 假设$\alpha^\mathrm{T}\beta\neq 0$，则$A$有两个特征值，一个为0，一个为$\alpha^\mathrm{T}\beta$，因此$A$有两个特征子空间，一个是0对应的$n-1$维特征子空间，一个是$\alpha^\mathrm{T}\beta$对应的一维特征子空间，因此$V$可分解为两个特征子空间的直和，与$A$不可对角化矛盾，因此$\alpha^\mathrm{T}\beta=0$.
    \end{enumerate}
\end{solve}

\begin{exercise}[补充训练 L15-例-8]\label{prob:L15-例-8}
设三阶矩阵$A$的特征值为$\lambda_1=-2,\lambda_2=1,\lambda_3=2$，对应的特征向量分别为$\alpha_1=(1,1,0)^\mathrm{T},\alpha_2=(1,0,1)^\mathrm{T},\alpha_3=(1,1,1)^\mathrm{T}$，求矩阵$A$.
\end{exercise}
\noindent 来源：\href{https://github.com/yhwu-is/Linear-Algebra-Left-Undone/blob/new/%E8%AE%B2%E4%B9%89/%E4%B8%93%E9%A2%98/15%20%E7%9B%B8%E4%BC%BC%E6%A0%87%E5%87%86%E5%BD%A2%EF%BC%9A%E5%A4%8D%E6%95%B0%E5%9F%9F%E4%B8%8A%E7%9A%84%E5%B0%9D%E8%AF%95%E4%B8%8E%E7%90%86%E8%AE%BA.tex}{L15-例-8}。
\begin{solve}
根据特征值分解可知，$A=P\Lambda P^{-1}$，其中$P=(\alpha_1,\alpha_2,\alpha_3),\Lambda=\begin{pmatrix}
            -2 & 0 & 0 \\
            0  & 1 & 0 \\
            0  & 0 & 2
        \end{pmatrix}$，因此
    \[A=P\Lambda P^{-1}=\begin{pmatrix}
            1 & 1 & 1 \\
            1 & 0 & 1 \\
            0 & 1 & 1
        \end{pmatrix}\begin{pmatrix}
            -2 & 0 & 0 \\
            0  & 1 & 0 \\
            0  & 0 & 2
        \end{pmatrix}\begin{pmatrix}
            1  & 0  & -1 \\
            1  & -1 & 0  \\
            -1 & 1  & 1
        \end{pmatrix}=\begin{pmatrix}
            -3 & 1 & 4 \\
            -4 & 2 & 4 \\
            -1 & 1 & 2
        \end{pmatrix}.\]
\end{solve}

\begin{exercise}[补充训练 L15-例-10]\label{prob:L15-例-10}
设$\sigma\in \mathcal{L}(\mathbf{C}^3)$定义为
    \[\sigma(z_1,z_2,z_3)=(6z_1+3z_2+4z_3,6z_2+2z_3,7z_3),\]求一组基使其有分块对角矩阵并写出对应的分块对角矩阵.
\end{exercise}
\noindent 来源：\href{https://github.com/yhwu-is/Linear-Algebra-Left-Undone/blob/new/%E8%AE%B2%E4%B9%89/%E4%B8%93%E9%A2%98/15%20%E7%9B%B8%E4%BC%BC%E6%A0%87%E5%87%86%E5%BD%A2%EF%BC%9A%E5%A4%8D%E6%95%B0%E5%9F%9F%E4%B8%8A%E7%9A%84%E5%B0%9D%E8%AF%95%E4%B8%8E%E7%90%86%E8%AE%BA.tex}{L15-例-10}。
\begin{solve}
不难写出 $\sigma$ 在$\mathcal{L}(\mathbf{C}^3)$的自然基下的矩阵是
    \[ \begin{pmatrix}
        6 & 3 & 4 \\
        0 & 6 & 2 \\
        0 & 0 & 7
    \end{pmatrix}, \]
    则 $\sigma$的特征值为 $\lambda_1 = 6,\lambda_2 = 7$. 解方程可以得到两个广义特征子空间 \[G(\lambda_1,\sigma) = \operatorname{span}\left((1,0,0),\left(0,\dfrac{1}{3},0\right)\right)\]，$G(\lambda_2,\sigma)=\operatorname{span}((10,2,1))$.

    取$\alpha_1=(1,0,0),\alpha_2=\left(0,\dfrac{1}{3},0\right),\alpha_3=(10,2,1)$，则$\sigma$在基$\{\alpha_1,\alpha_2,\alpha_3\}$下有分块对角矩阵
    \[ \begin{pmatrix}
        6 & 1 & 0 \\
        0 & 6 & 0 \\
        0 & 0 & 7
    \end{pmatrix}. \]
\end{solve}

\begin{exercise}[补充训练 L15-例-11]\label{prob:L15-例-11}
设$\sigma\in \mathcal{L}(\mathbf{C}^4)$定义为
    \[\sigma(z_1,z_2,z_3,z_4)=(-z_1+z_2,-z_1+z_2+z_3,z_1-z_2+2z_3+z_4,-z_1+z_2-z_3+3z_4),\]
    求$\mathbf{C}^4$的一组基，使得$\sigma$在这组基下的矩阵为上三角矩阵.
\end{exercise}
\noindent 来源：\href{https://github.com/yhwu-is/Linear-Algebra-Left-Undone/blob/new/%E8%AE%B2%E4%B9%89/%E4%B8%93%E9%A2%98/15%20%E7%9B%B8%E4%BC%BC%E6%A0%87%E5%87%86%E5%BD%A2%EF%BC%9A%E5%A4%8D%E6%95%B0%E5%9F%9F%E4%B8%8A%E7%9A%84%E5%B0%9D%E8%AF%95%E4%B8%8E%E7%90%86%E8%AE%BA.tex}{L15-例-11}。
\begin{solve}
记$\mathbf{C}^4$的自然基为$e_1,e_2,e_3,e_4$，则$\sigma$在自然基下的矩阵为
    \[
        A=\begin{pmatrix}
            -1 & 1  & 0  & 0 \\
            -1 & 1  & 1  & 0 \\
            1  & -1 & 2  & 1 \\
            -1 & 1  & -1 & 3
        \end{pmatrix}.
    \]
    下面按照原讲义“上三角矩阵存在”处 证明的法二逐步构造相似变换矩阵. 注意到 $A(1,1,0,0)^\mathrm{T} = 0$，因此 $0$ 是 $A$ 的一个特征值，对应的特征向量可取 $\alpha_1=(1,1,0,0)^\mathrm{T}$. 令 $P_1 = (\alpha_1,e_2,e_3,e_4)$，则
    \[ P_1^{-1}AP_1 = \begin{pmatrix}
            0 & 1  & 0  & 0 \\
            0 & 0  & 1  & 0 \\
            0 & -1 & 2  & 1 \\
            0 & 1  & -1 & 3
        \end{pmatrix}. \]
    其右下角三阶矩阵记为 $B_1$. 由于$B_1(1,1,0)^\mathrm{T} = (1,1,0)^\mathrm{T}$，故 $1$ 是 $B_1$ 的一个特征值. 取 $Q_2 = \begin{pmatrix}
            1 & 0 & 0 \\
            1 & 1 & 0 \\
            0 & 0 & 1
        \end{pmatrix}$，则
    \[
        Q_2^{-1}B_1Q_2
        =
        \begin{pmatrix}
            1 & 1  & 0 \\
            0 & 1  & 1 \\
            0 & -1 & 3
        \end{pmatrix}.
    \]
    再看 $Q_2^{-1}B_1Q_2$ 的右下角二阶矩阵，记为 $B_2$，有 $B_2(1,1)^\mathrm{T} = 2(1,1)^\mathrm{T}$，故 $2$ 是 $B_2$ 的一个特征值. 取 $Q_3= \begin{pmatrix}
            1 & 0 \\
            1 & 1
        \end{pmatrix}$，则
    \[
        Q_3^{-1}B_2Q_3
        =
        \begin{pmatrix}
            2 & 1 \\
            0 & 2
        \end{pmatrix}.
    \]
    至此，可以逐步构造出相似变换矩阵. 首先得到
    \[
        P_2 = Q_2
        \begin{pmatrix}
            1 & O \\ O & Q_3
        \end{pmatrix}
        =
        \begin{pmatrix}
            1 & 0 & 0 \\
            1 & 1 & 0 \\
            0 & 1 & 1
        \end{pmatrix}.
    \]
    则 $P_2$ 使得 $P_2^{-1}B_1P_2$ 为上三角矩阵. 进而得到
    \[
        P = P_1 \begin{pmatrix}
            1 & O \\ O & P_2
        \end{pmatrix} = \begin{pmatrix}
            1 & 0 & 0 & 0 \\
            1 & 1 & 0 & 0 \\
            0 & 1 & 1 & 0 \\
            0 & 0 & 1 & 1
        \end{pmatrix}. \]
    也就是说，取 $\alpha_1=(1,1,0,0)^\mathrm{T}, \alpha_2=(0,1,1,0)^\mathrm{T}, \alpha_3=(0,0,1,1)^\mathrm{T}, \alpha_4=(0,0,0,1)^\mathrm{T}$，则$\sigma$在基$\{\alpha_1,\alpha_2,\alpha_3,\alpha_4\}$下的矩阵为上三角矩阵
    \[
        P^{-1}AP = \begin{pmatrix}
            0 & 1 & 0 & 0 \\
            0 & 1 & 1 & 0 \\
            0 & 0 & 2 & 1 \\
            0 & 0 & 0 & 2
        \end{pmatrix}. \]
\end{solve}

\begin{exercise}[补充训练 L15-例-15]\label{prob:L15-例-15}
定义$N\in \mathcal{L}(\mathbf{F}^5)$为
    \[N(x_1,x_2,x_3,x_4,x_5)=(2x_2,3x_3,-x_4,4x_5,0)\]
    求$(I+N)$的一个平方根.
\end{exercise}
\noindent 来源：\href{https://github.com/yhwu-is/Linear-Algebra-Left-Undone/blob/new/%E8%AE%B2%E4%B9%89/%E4%B8%93%E9%A2%98/15%20%E7%9B%B8%E4%BC%BC%E6%A0%87%E5%87%86%E5%BD%A2%EF%BC%9A%E5%A4%8D%E6%95%B0%E5%9F%9F%E4%B8%8A%E7%9A%84%E5%B0%9D%E8%AF%95%E4%B8%8E%E7%90%86%E8%AE%BA.tex}{L15-例-15}。
\begin{solve}
由于$N^4=0$，我们设$\sqrt{I+N}=I+a_1N+a_2N^2+a_3N^3$，则
    \begin{align*}
        I+N&=(I+a_1N+a_2N^2+a_3N^3)^2 \\
        &=I+2a_1N+(a_1^2+2a_2)N^2+(2a_3+2a_1a_2)N^3,
    \end{align*}
    可解得$a_1=\dfrac{1}{2},a_2=-\dfrac{1}{8},a_3=\dfrac{1}{16}$. 代入即可算出$(I+N)$的一个平方根.
\end{solve}

\begin{exercise}[补充训练 L15-例-18]\label{prob:L15-例-18}
找出有限维实向量空间的一个线性变换$\sigma$，使得0是$\sigma$仅有的特征值但$\sigma$不是幂零线性变换.
\end{exercise}
\noindent 来源：\href{https://github.com/yhwu-is/Linear-Algebra-Left-Undone/blob/new/%E8%AE%B2%E4%B9%89/%E4%B8%93%E9%A2%98/15%20%E7%9B%B8%E4%BC%BC%E6%A0%87%E5%87%86%E5%BD%A2%EF%BC%9A%E5%A4%8D%E6%95%B0%E5%9F%9F%E4%B8%8A%E7%9A%84%E5%B0%9D%E8%AF%95%E4%B8%8E%E7%90%86%E8%AE%BA.tex}{L15-例-18}。
\begin{solve}
我们定义$\sigma\in\mathcal{L}({\mathbf{R}^3})$，$\sigma$在自然基下的矩阵为
            \[\begin{pmatrix}
                0 & 1 & 0 \\
                -1 & 0 & 0 \\
                0 & 0 & 0
            \end{pmatrix},\]
            则$0$是$\sigma$仅有的特征值，但$\sigma^3=-\sigma\neq 0$，$\sigma$不为幂零线性变换.
\end{solve}

\begin{exercise}[补充训练 L15-例-19]\label{prob:L15-例-19}
找出一个$\sigma\in \mathcal{L}(\mathbf{R}^2)$使得$\sigma^4=-I$.
\end{exercise}
\noindent 来源：\href{https://github.com/yhwu-is/Linear-Algebra-Left-Undone/blob/new/%E8%AE%B2%E4%B9%89/%E4%B8%93%E9%A2%98/15%20%E7%9B%B8%E4%BC%BC%E6%A0%87%E5%87%86%E5%BD%A2%EF%BC%9A%E5%A4%8D%E6%95%B0%E5%9F%9F%E4%B8%8A%E7%9A%84%E5%B0%9D%E8%AF%95%E4%B8%8E%E7%90%86%E8%AE%BA.tex}{L15-例-19}。
\begin{solve}
从旋转的角度考虑可以得到
            \[
            \sigma(x,y)=(x,y)\begin{pmatrix}
                \dfrac{\sqrt{2}}{2} & -\dfrac{\sqrt{2}}{2} \\[1ex]
                \dfrac{\sqrt{2}}{2} & \dfrac{\sqrt{2}}{2}
            \end{pmatrix}
            \]
            满足要求.
\end{solve}

\begin{exercise}[补充训练 L15-例-20]\label{prob:L15-例-20}
证明或给出反例：$V$上的幂零线性变换的集合是$\mathcal{L}(V)$的子空间.
\end{exercise}
\noindent 来源：\href{https://github.com/yhwu-is/Linear-Algebra-Left-Undone/blob/new/%E8%AE%B2%E4%B9%89/%E4%B8%93%E9%A2%98/15%20%E7%9B%B8%E4%BC%BC%E6%A0%87%E5%87%86%E5%BD%A2%EF%BC%9A%E5%A4%8D%E6%95%B0%E5%9F%9F%E4%B8%8A%E7%9A%84%E5%B0%9D%E8%AF%95%E4%B8%8E%E7%90%86%E8%AE%BA.tex}{L15-例-20}。
\begin{solve}
反例：$V=\mathbf{R}^2,\sigma,\tau$在自然基下的矩阵分别为
            \[ \begin{pmatrix}
                0 & 1 \\
                0 & 0
            \end{pmatrix},
            \begin{pmatrix}
                0 & 0 \\
                1 & 0
            \end{pmatrix}\]
            则$\sigma+\tau$在自然基下的矩阵为
            \[ \begin{pmatrix}
                0 & 1 \\
                1 & 0
            \end{pmatrix} \]
            表明$\sigma+\tau$不是幂零线性变换. 从而命题不成立.
\end{solve}

\begin{exercise}[补充训练 S-EIG-1]\label{prob:S-EIG-1}
设 $A=\begin{pmatrix}-1&-6&0\\0&1&0\\-4&-12&3\end{pmatrix}$。求特征值与一组特征向量基，写出对角化分解，并计算 $A^4$ 与 $\det A$。
\end{exercise}
\noindent 编者补充：S-EIG-1；相关主题见原讲义对应讲，题目数据与解答由本选编另行构造。
\begin{solve}
计算特征多项式，或逐一核对下面的列向量，可得三个不同特征值 $-1,1,3$。对应特征向量按列排列为
\[P=\begin{pmatrix}1&-3&0\\0&1&0\\1&0&1\end{pmatrix},\qquad D=\operatorname{diag}(-1,1,3).\]
直接计算 $AP=PD$ 且 $\det P=1$，从而 $A=PDP^{-1}$。因此
\[A^4=PD^4P^{-1}=\begin{pmatrix}1&0&0\\0&1&0\\-80&-240&81\end{pmatrix},\qquad \det A=-3.\]
不同特征值对应的特征向量线性无关，所以这三个向量构成一组基。该方法把高次矩阵乘法转化为标量幂；一般地 $A^m=P\operatorname{diag}((-1)^m,1,3^m)P^{-1}$，对所有非负整数 $m$ 成立。
\end{solve}

\section{补充训练：若当标准形}

本节包含 9 道训练题，建议在完成相关专题真题后选择性练习。

\begin{exercise}[补充训练 L16-例-3]\label{prob:L16-例-3}
设$V$是由两个变量$x,y$生成的次数不大于2的多项式函数（例如$xy^2$不是，因为次数等于1+2=3，$x^2+y$是）关于一般的多项式加法和数乘构成的线性空间，于是我们可以取一组基为$\{1,x,y,x^2,y^2,xy\}$，定义$\sigma\in \mathcal{L}(V)$为
    \[\sigma(f(x,y))=\dfrac{\partial}{\partial x}f(x,y).\]
    求$V$的一组基使得$\sigma$在这组基下的矩阵是若当标准形.
\end{exercise}
\noindent 来源：\href{https://github.com/yhwu-is/Linear-Algebra-Left-Undone/blob/new/%E8%AE%B2%E4%B9%89/%E4%B8%93%E9%A2%98/16%20%E8%8B%A5%E5%BD%93%E6%A0%87%E5%87%86%E5%BD%A2.tex}{L16-例-3}。
\begin{solve}
由于对多项式求导会降低次数，因此$\sigma$显然是幂零的（也可以直接写出$\sigma$在题目中给出的这组基的矩阵表示，不难得知其为对角线全为$0$的上三角矩阵），因此唯一特征值为$0$. 为了表示简便，这里用 $G_j$ 代替之前的记号 $G_j(\lambda,\sigma)=\ker (\sigma-\lambda I)^j$.

    不难得到：
    \begin{itemize}
        \item $G_1 = \ker \sigma = \operatorname{span}\{1, y, y^2\}$，$\dim G_1 = 3$；
        \item $G_2 = \ker \sigma^2 = \operatorname{span}\{1, y, y^2, x, xy\}$，$\dim G_2 = 5$；
        \item $G_3 = \ker \sigma^3 = \operatorname{span}\{1, y, y^2, x, xy, x^2\} = V$，$\dim G_3 = 6$.
    \end{itemize}

    从而根据原讲义“Young图行长度”处，Young图的行长度分别为 $r_1 = 3, r_2 = 2, r_3 = 1$，我们按照刚刚的讨论填充Young图.

    \begin{enumerate}
        \item 倒数第一层中，我们需要找到一个向量 $v_1 \in G_3 \backslash G_2$，因为这里$G_1, G_2, G_3$的基满足了包含关系，所以作差变得非常简单：直接挑出 $v_1=x^2$ 即可.
        \item 倒数第二层中，我们需要找到一个向量 $v_2 \in G_2 \backslash G_1$，且 $v_2$ 与该层已有向量 $\sigma(v_1)=2x$ 线性无关. 注意到 $G_2 = \operatorname{span}\{G_1, x, xy\}$，在 $G_2$ 中，$\{1, y, y^2, 2x, xy\}$ 线性无关且张成 $G_2$，故可取 $v_2 = xy$.
        \item 第一层中，我们需要找到一个向量 $v_3 \in G_1$，且 $v_3$ 与该层已有向量 $\sigma^2(v_1)=2, \sigma(v_2)=y$ 线性无关. 注意到 $G_1 = \operatorname{span}\{1, y, y^2\}$. 显然 $\{2, y, y^2\}$ 是 $G_1$ 的一组基，故可取 $v_3 = y^2$.
    \end{enumerate}

    从而我们得到了若当基 $\{2, 2x, x^2, y, xy, y^2\}$.
\end{solve}

\begin{exercise}[补充训练 L16-例-4]\label{prob:L16-例-4}
设$V$为$n$维复向量空间，$\sigma\in \mathcal{L}(V)$，$\sigma$在基$v_1,\ldots,v_n$下的矩阵是一个若当块
    \[\begin{pmatrix}
            \lambda & 1       & 0      & \cdots & 0       \\
            0       & \lambda & 1      & \cdots & 0       \\
            \vdots  & \vdots  & \vdots & \ddots & \vdots  \\
            0       & 0       & 0      & \cdots & \lambda
        \end{pmatrix},\]
    证明：
    \begin{enumerate}
        \item $V$中包含$v_n$的不变子空间只有$V$自身；

        \item $V$中任意非零不变子空间都包含$v_1$；

        \item $V$不能分解为两个非平凡的不变子空间的直和；

        \item $V$中有且仅有$n+1$个不变子空间，它们分别是
              \[\{0\},\operatorname{span}(v_1),\operatorname{span}(v_1,v_2),\ldots,\operatorname{span}(v_1,\ldots,v_{n-1},v_n)\]
    \end{enumerate}
\end{exercise}
\noindent 来源：\href{https://github.com/yhwu-is/Linear-Algebra-Left-Undone/blob/new/%E8%AE%B2%E4%B9%89/%E4%B8%93%E9%A2%98/16%20%E8%8B%A5%E5%BD%93%E6%A0%87%E5%87%86%E5%BD%A2.tex}{L16-例-4}。
\begin{solve}
\begin{enumerate}
        \item 设$W$是包含$v_n$的不变子空间，根据线性映射矩阵表示的定义，$\sigma(v_n)=\lambda v_n+v_{n-1}\in W$，故$v_{n-1}\in W$，同理可得$v_{n-2},\ldots,v_1\in W$，即$W=V$.

        \item 设$U$为任一非零不变子空间，那么$U$中存在一个非零向量$v$，则$v$一定可以表达为
              \[v=k_1v_1+\cdots+k_sv_s(s\leqslant n),\]
              其中$k_s\neq 0$. 由$\sigma(v)\in U$以及线性映射矩阵表示有
              \begin{align*}
                  \sigma(v) & = \lambda k_1v_1+k_2(\lambda v_2+v_1)+\cdots+k_s(\lambda v_s+v_{s-1}) \\
                            & = \lambda v+k_2v_1+\cdots+k_sv_{s-1}\in U,
              \end{align*}
              因此$k_2v_1+\cdots+k_sv_{s-1}\in U$，同理进一步作用$\sigma$有$k_3v_1+\cdots+k_sv_{s-2}\in U$，以此类推有$k_sv_1\in U$. 又$k_s\neq 0$，故$v_1\in U$，命题得证.

        \item 利用上一点显然.

        \item 显然$\{0\},\operatorname{span}(v_1),\operatorname{span}(v_1,v_2),\ldots,\operatorname{span}(v_1,\ldots,v_{n-1},v_n)$都是不变子空间，下面证明不变子空间一定在它们当中.

              设$U$为任一不变子空间，若$U=\{0\}$，则也符合结论. 我们考虑非零的情况，设$\beta_1,\ldots,\beta_s$是$U$的一组基，并且假设
              \[\beta_i=a_{i1}v_1+a_{i2}v_2+\cdots+a_{it}v_t,\enspace i=1,2,\ldots,s\]
              其中满足至少有一个$i$使得$a_{it}\neq 0$（把所有基下表示末尾的0去除即可）. 对于满足$a_{it}\neq 0$的$i$，由于$U$是不变子空间，故$\sigma(\beta_i)\in U$，与第二点证明完全相同，我们得到
              \begin{align*}
                  a_{i2}v_1+a_{i3}v_2+\cdots+a_{it}v_{t-1}\in U, \\
                  a_{i3}v_1+a_{i4}v_2+\cdots+a_{it}v_{t-2}\in U, \\
                  \cdots                                         \\
                  a_{i,t-1}v_1+a_{it}v_2\in U,                   \\
                  a_{it}v_1\in U.
              \end{align*}
              由于$a_{it}\neq 0$，故由最后的式子知道$v_1\in U$. 代入倒数第二个式子知道$v_2\in U$，以此类推，最终得到$\operatorname{span}(v_1,\ldots,v_t)\subset U$，而$U\subset\operatorname{span}(v_1,\ldots,v_t)$显然（因为每个$\beta_i$都在其中），故命题得证.
    \end{enumerate}
\end{solve}

\begin{exercise}[补充训练 L16-例-5]\label{prob:L16-例-5}
设$V$是复数域上的$n$维线性空间，$\sigma\in\mathcal{L}(V)$，若$\sigma$有$n$个不同的特征值$\lambda_1,\ldots,\lambda_n$，求$\sigma$的不变子空间的个数.
\end{exercise}
\noindent 来源：\href{https://github.com/yhwu-is/Linear-Algebra-Left-Undone/blob/new/%E8%AE%B2%E4%B9%89/%E4%B8%93%E9%A2%98/16%20%E8%8B%A5%E5%BD%93%E6%A0%87%E5%87%86%E5%BD%A2.tex}{L16-例-5}。
\begin{solve}
显然$\sigma$可对角化，每个对角元素就是一阶若当块，因此我们需要考虑组合这些若当块对应的基构成新的不变子空间. 由可对角化我们有
    \[V=V_{\lambda_1}\oplus\cdots\oplus V_{\lambda_n},\]
    结合$\dim V=n$可知每个特征子空间都是1维的. 根据我们组合的思想，我们知道对于任意的$1\leqslant j_1<j_2<\cdots<j_r\leqslant n$，$r=1,2,\ldots,n$，$V_{\lambda_{j_1}}\oplus\cdots\oplus V_{\lambda_{j_r}}$都是不变子空间. 这样的不变子空间加上零空间的个数总和为$\mathrm{C}_n^0+\mathrm{C}_n^1+\cdots+\mathrm{C}_n^n=2^n$.

    下面我们需要说明这些不变子空间是全部的不变子空间. 设$U$为任一不变子空间，若为零空间则符合要求，若不是零空间，设$\dim U=m>0$，$U$的一组基为$\beta_1,\ldots,\beta_m$，将其扩充为$V$的一组基$\beta_1,\ldots,\beta_m,\beta_{m+1},\ldots,\beta_n$，则在这组基下$\sigma$的矩阵为
    \[B=\begin{pmatrix}
            A & C \\ O & D
        \end{pmatrix},\]
    于是$\sigma\vert_U$的矩阵为$A$，对应特征多项式$|\lambda E-A|$，而$\sigma$的特征多项式$|\lambda E-B|=|\lambda E-A||\lambda E-D|$. 由于$\sigma$有$n$个不同的特征值，因此$\sigma\vert_U$必定有$m$个不同的特征值$\lambda_{l_1},\ldots,\lambda_{l_m}$，且$1\leqslant l_1<\cdots<l_m\leqslant n$. 而对每个特征值$\lambda_{l_i}$都有特征向量$\alpha_i\in U$，满足$\sigma\vert_U(\alpha_i)=\lambda_{l_i}\alpha_i$，因此$\sigma(\alpha_i)=\lambda_{l_i}\alpha_i$，即$\alpha_i\in V_{\lambda_{l_i}}$，故$V_{\lambda_{l_1}}\oplus\cdots\oplus V_{\lambda_{l_m}}\subset U$. 再结合
    \[\dim(V_{\lambda_{l_1}}\oplus\cdots\oplus V_{\lambda_{l_m}})=m=\dim U,\]
    故$U=V_{\lambda_{l_1}}\oplus\cdots\oplus V_{\lambda_{l_m}}$，故所有不变子空间都是之前讨论的$2^n$种之一，命题得证.
\end{solve}

\begin{exercise}[补充训练 L16-例-6]\label{prob:L16-例-6}
设$A=\begin{pmatrix}
            2 & 6 & -15 \\ 1 & 1 & -5 \\ 1 & 2 & -6
        \end{pmatrix}$，求$A^{m}(m\geqslant 1)$.
\end{exercise}
\noindent 来源：\href{https://github.com/yhwu-is/Linear-Algebra-Left-Undone/blob/new/%E8%AE%B2%E4%B9%89/%E4%B8%93%E9%A2%98/16%20%E8%8B%A5%E5%BD%93%E6%A0%87%E5%87%86%E5%BD%A2.tex}{L16-例-6}。
\begin{solve}
首先求 $A$ 的若当标准形，$|\lambda E - A|=(\lambda + 1)^3$，$-1$为唯一特征值，几何重数为$3$，而$A+E=\begin{pmatrix}
            3 & 6 & -15 \\ 1 & 2 & -5 \\ 1 & 2 & -5
        \end{pmatrix}$为秩为$1$的矩阵，因此由原讲义“矩阵若当标准形唯一”处，若当块有$3-1=2$个，只可能是$J_2(-1)$和$J_1(-1)$，因此$A=PJP^{-1}$，其中$J=\begin{pmatrix}
            -1 & 1 & 0 \\ 0 & -1 & 0 \\ 0 & 0 & -1
        \end{pmatrix}$，用原讲义“求若当标准形过渡矩阵”处 中的方法，可得$P=\begin{pmatrix}
            3 & 1 & 2 \\ 1 & 0 & -1 \\ 1 & 0 & 0
        \end{pmatrix}$，$P^{-1}=\begin{pmatrix}
            0 & 0 & 1 \\ 1 & 2 & -5 \\ 0 & -1 & 1
        \end{pmatrix}$，故$A^m=PJ^mP^{-1}$，其中
        \[
            J^m=\begin{pmatrix}
                (-1)^m E + m(-1)^{m-1}J_2(0) & O \\ O & (-1)^m
            \end{pmatrix}=(-1)^m\begin{pmatrix}
                1 & -m & 0 \\ 0 & 1 & 0 \\ 0 & 0 & 1
            \end{pmatrix},
        \]
        整理可得 $A^m=(-1)^m\begin{pmatrix}
            1-3m & -6m & 15m \\ -m & 1-2m & 5m \\ -m & -2m & 1+5m
        \end{pmatrix}$.
\end{solve}

\begin{exercise}[补充训练 L16-例-7]\label{prob:L16-例-7}
设$J=J_n(a)$是特征值为$a$的$n$阶若当块，求$J^m$的若当标准形，其中$m$为非零整数.当 $a=0$ 时，仅考虑正整数 $m$，因为负整数次幂无定义。
\end{exercise}
\noindent 来源：\href{https://github.com/yhwu-is/Linear-Algebra-Left-Undone/blob/new/%E8%AE%B2%E4%B9%89/%E4%B8%93%E9%A2%98/16%20%E8%8B%A5%E5%BD%93%E6%A0%87%E5%87%86%E5%BD%A2.tex}{L16-例-7}。
\begin{solve}
我们需要分别考虑$a\neq 0$和$a=0$的情况：
    \begin{enumerate}
        \item 当$a\neq 0$时，首先显然$J^m$的特征值都为$a^m$，我们先处理$m\geqslant 1$的情况，我们有
              \[J^m=(aE+J(0))^m=a^mE+\mathrm{C}_m^1a^{m-1}J_n(0)+\cdots+\mathrm{C}_m^mJ_n(0)^m.\]
              故$r(J^m-a^mE_n)=r(\mathrm{C}_m^1a^{m-1}J_n(0)+\cdots+\mathrm{C}_m^mJ_n(0)^m)=n-1$（很容易通过$J_n(0)$的幂次的特点写出这一矩阵，然后通过行列式秩判断），因此根据原讲义“若当标准形唯一”处 可知$J^m$的若当块个数为$n-(n-1)=1$，因此$J^m$的若当标准形就是$J_n(a^m)$.

              然后我们考虑$m=-1$的情况，显然$J^{-1}$的特征值为$a^{-1}$，这里我们借用原讲义“分式求逆1”处 的分式思想，得到
              \[J^{-1}=\dfrac{1}{a}E_n-\dfrac{1}{a^2}J_n(0)+\cdots+(-1)^{n-1}\dfrac{1}{a^n}J_n(0)^{n-1}.\]
              故$r(J^{-1}-\dfrac{1}{a}E_n)=r(-\dfrac{1}{a^2}J_n(0)+\cdots+(-1)^{n-1}\dfrac{1}{a^n}J_n(0)^{n-1})=n-1$，因此$J^{-1}$的若当块个数为$n-(n-1)=1$，因此$J^{-1}$的若当标准形就是$J_n(a^{-1})$.

              最后处理$m\leqslant -2$的情况，注意到$J^m=(J^{-1})^{-m}$，故根据前面$m\geqslant 1$的讨论，$J^m$的若当标准形就是$J_n((a^{-1})^{-m})=J_n(a^m)$.

              综上所述，无论$m$取什么整数值，$J^m$的若当标准形就是$J_n(a^m)$.

        \item 当$a=0$时，$J$不可逆，因此只考虑$m>0$的情形. 我们有$m\geqslant n$时$J^m=O$，此时若当标准形也就是零矩阵. 我们考虑$m<n$的情况，我们需要使用原讲义“若当标准形唯一”处 的结论来辅助我们求解$J^m$的若当标准形，因此我们需要研究$(J^m)^k$的秩. 我们做带余除法，$n=mq+r(0\leqslant r<m)$，则根据$J_n(0)$的幂次的特点，我们有
              \[r((J^m)^k)=n-mk,\enspace 0\leqslant k\leqslant q,\enspace r((J^m)^k)=0,\enspace k\geqslant q+1.\]
              因此
              \begin{enumerate}
                  \item $1\leqslant k<q$时，$J_k(0)$的个数为$r((J^m)^{k-1})+r((J^m)^{k+1})-2r((J^m)^k)=(n-m(k-1))+(n-m(k+1))-2(n-mk)=0$；
                  \item $k=q$时，$J_k(0)$的个数为$r((J^m)^{q-1})+r((J^m)^{q+1})-2r((J^m)^q)=(n-m(q-1))+0-2(n-mq)=m-r$；
                  \item $k=q+1$时，$J_k(0)$的个数为$r((J^m)^q)+r((J^m)^{q+2})-2r((J^m)^{q+1})=(n-mq)+0-0=r$.
                  \item $k>q+1$时，$J_k(0)$的个数为0.
              \end{enumerate}
              综上所述，$J^m$的若当标准形是$\operatorname{diag}(J_q(0),\ldots,J_q(0),J_{q+1}(0),\ldots,J_{q+1}(0))$，其中$J_q(0)$的个数为$m-r$，$J_{q+1}(0)$的个数为$r$.
    \end{enumerate}
\end{solve}

\begin{exercise}[补充训练 L16-例-9]\label{prob:L16-例-9}
利用若当标准形求解初值问题
    \[
        X'(t)=\begin{pmatrix}
            2 & 3 & 2 \\ 1 & 8 & 2 \\ -2 & -14 & -3
        \end{pmatrix}X(t),\qquad
        X(0)=\begin{pmatrix}
            -1 \\ 0 \\ 0
        \end{pmatrix}.
    \]
\end{exercise}
\noindent 来源：\href{https://github.com/yhwu-is/Linear-Algebra-Left-Undone/blob/new/%E8%AE%B2%E4%B9%89/%E4%B8%93%E9%A2%98/16%20%E8%8B%A5%E5%BD%93%E6%A0%87%E5%87%86%E5%BD%A2.tex}{L16-例-9}。
\begin{solve}
记
    \[
        A=\begin{pmatrix}
            2 & 3 & 2 \\ 1 & 8 & 2 \\ -2 & -14 & -3
        \end{pmatrix}.
    \]
    由原讲义“求若当标准形过渡矩阵”处 可知，$A$的若当标准形为
    \[
        J=\begin{pmatrix}
            3 & 1 & 0 \\ 0 & 3 & 0 \\ 0 & 0 & 1
        \end{pmatrix},
        \qquad
        P=\begin{pmatrix}
            1 & -1 & 2 \\ -1 & 0 & 0 \\ 2 & 0 & -1
        \end{pmatrix},
        \qquad
        P^{-1}AP=J.
    \]

    令$X(t)=PY(t)$，则
    \[
        PY'(t)=X'(t)=AX(t)=APY(t)=PJY(t).
    \]
    由于$P$可逆，故$Y'(t)=JY(t)$. 又注意到初值向量$X(0)=(-1,0,0)^\mathrm{T}$恰为$P$的第二列，因此
    \[
        Y(0)=P^{-1}X(0)=\begin{pmatrix}
            0 \\ 1 \\ 0
        \end{pmatrix}.
    \]

    现在只需求解若当标准形对应的方程组. 由于
    \[
        J=\operatorname{diag}(J_2(3),1),
    \]
    因而
    \[
        \mathrm{e}^{Jt}=\begin{pmatrix}
            \mathrm{e}^{3t} & t\mathrm{e}^{3t} & 0 \\
            0 & \mathrm{e}^{3t} & 0 \\
            0 & 0 & \mathrm{e}^{t}
        \end{pmatrix}.
    \]
    故
    \[
        Y(t)=\mathrm{e}^{Jt}Y(0)=\begin{pmatrix}
            t\mathrm{e}^{3t} \\ \mathrm{e}^{3t} \\ 0
        \end{pmatrix}.
    \]

    最后代回$X(t)=PY(t)$，得到
    \[
        X(t)=\begin{pmatrix}
            1 & -1 & 2 \\ -1 & 0 & 0 \\ 2 & 0 & -1
        \end{pmatrix}\begin{pmatrix}
            t\mathrm{e}^{3t} \\ \mathrm{e}^{3t} \\ 0
        \end{pmatrix}
        =\mathrm{e}^{3t}\begin{pmatrix}
            t-1 \\ -t \\ 2t
        \end{pmatrix}.
    \]
    这就是所求的解.
\end{solve}

\begin{exercise}[补充训练 L16-例-10]\label{prob:L16-例-10}
求解线性递推方程$a_n = 5a_{n - 1} - 8a_{n - 2} + 4a_{n - 3}$，其中$a_0 = 2, a_1 = 4, a_2 = 9$.
\end{exercise}
\noindent 来源：\href{https://github.com/yhwu-is/Linear-Algebra-Left-Undone/blob/new/%E8%AE%B2%E4%B9%89/%E4%B8%93%E9%A2%98/16%20%E8%8B%A5%E5%BD%93%E6%A0%87%E5%87%86%E5%BD%A2.tex}{L16-例-10}。
\begin{solve}
设这一递推方程的特征多项式为
    \[
        p(\lambda)=\lambda^3-5\lambda^2+8\lambda-4=(\lambda-1)(\lambda-2)^2.
    \]
    因此特征值为$1$与$2$，其中$2$是二重根. 根据前面对重根情形的讨论，数列$\{a_n\}$的通项应具有形式
    \[
        a_n=c_1\cdot 1^n+c_2\cdot 2^n+c_3\binom{n}{1}2^{n-1}
        =c_1+c_2\cdot 2^n+c_3n2^{n-1},
    \]
    其中$c_1,c_2,c_3$为待定常数.

    将初值代入上式，得到方程组
    \[
        \begin{cases}
            c_1+c_2=2, \\
            c_1+2c_2+c_3=4, \\
            c_1+4c_2+4c_3=9.
        \end{cases}
    \]
    解得
    \[
        c_1=1,\qquad c_2=1,\qquad c_3=1.
    \]
    从而
    \[
        a_n=1+2^n+n2^{n-1}.
    \]
    这就是所求的通项公式.
\end{solve}

\begin{exercise}[补充训练 L16-例-12]\label{prob:L16-例-12}
求矩阵
    \[
        A=\begin{pmatrix}
            0 & 0 & 4 \\
            1 & 0 & -6 \\
            0 & 1 & 4
        \end{pmatrix}
    \]
    在实数域上的若当标准形.
\end{exercise}
\noindent 来源：\href{https://github.com/yhwu-is/Linear-Algebra-Left-Undone/blob/new/%E8%AE%B2%E4%B9%89/%E4%B8%93%E9%A2%98/16%20%E8%8B%A5%E5%BD%93%E6%A0%87%E5%87%86%E5%BD%A2.tex}{L16-例-12}。
\begin{solve}
按照前面的总结，我们先将$A$视为复矩阵来求其复数域上的若当标准形. 计算特征多项式：
    \[
        |\lambda E-A|
        =\begin{vmatrix}
            \lambda & 0 & -4 \\
            -1 & \lambda & 6 \\
            0 & -1 & \lambda-4
        \end{vmatrix}
        =\lambda(\lambda^2-4\lambda+6)-4
        =\lambda^3-4\lambda^2+6\lambda-4.
    \]
    不难得到 $A$ 的特征值为 $2,\qquad 1+\mathrm{i},\qquad 1-\mathrm{i}$. 这三个特征值两两不同，所以$A$在复数域上的若当标准形就是对角矩阵 $J_{\mathbf{C}}=\operatorname{diag}(2,1+\mathrm{i},1-\mathrm{i})$. 现在回到实数域上. 其中实特征值$2$对应一个通常的一阶若当块$J_1(2)=(2)$；一对共轭复特征值$1\pm\mathrm{i}$对应一个二维实块 $R(1,1)=\begin{pmatrix}
            1 & 1 \\
            -1 & 1
        \end{pmatrix}$. 因而$A$在实数域上的若当标准形可取为
    \[
        J_{\mathbf{R}}=\operatorname{diag}(2,R(1,1))
        =\begin{pmatrix}
            2 & 0 & 0 \\
            0 & 1 & 1 \\
            0 & -1 & 1
        \end{pmatrix}.
    \]
\end{solve}

\begin{exercise}[补充训练 S-JORDAN-1]\label{prob:S-JORDAN-1}
设 $A=\begin{pmatrix}5&10&-3\\-1&-1&1\\0&1&2\end{pmatrix}$。求若当标准形和最小多项式，计算 $\dim\ker(A-2I)^r$（$r=1,2,3$），并给出 $A^m$ 的通式。
\end{exercise}
\noindent 编者补充：S-JORDAN-1；相关主题见原讲义对应讲，题目数据与解答由本选编另行构造。
\begin{solve}
取 $P=\begin{pmatrix}1&-3&0\\0&1&0\\1&0&1\end{pmatrix}$，逐列相乘可验证
\[AP=PJ,\qquad J=\begin{pmatrix}2&1&0\\0&2&1\\0&0&2\end{pmatrix}.\]
所以若当标准形是一个三阶块 $J_3(2)$。记 $N=A-2I$，有 $N^3=O$、$N^2\ne O$，最小多项式为 $(x-2)^3$。在若当基下，$N$ 依次把第三个基向量送到第二个、第二个送到第一个、第一个送到零，故三个核维数依次为 $1,2,3$。因为 $N$ 与 $I$ 可交换，二项式公式给出
\[A^m=2^mI+m2^{m-1}N+\binom m22^{m-2}N^2\quad(m\ge2).\]
$m=0,1$ 时分别为 $I,A$。最大若当块为三阶，说明该矩阵不能对角化。
\end{solve}

\section{补充训练：最小多项式}

本节包含 2 道训练题，建议在完成相关专题真题后选择性练习。

\begin{exercise}[补充训练 L17-例-3]\label{prob:L17-例-3}
求矩阵
    \[
        A=\begin{pmatrix}
                2 & 3 & 2 \\
                1 & 8 & 2 \\
                -2 & -14 & -3
            \end{pmatrix}
    \]
    的初等因子组.
\end{exercise}
\noindent 来源：\href{https://github.com/yhwu-is/Linear-Algebra-Left-Undone/blob/new/%E8%AE%B2%E4%B9%89/%E4%B8%93%E9%A2%98/17%20%E5%A4%9A%E9%A1%B9%E5%BC%8F%E7%9A%84%E8%BF%9B%E4%B8%80%E6%AD%A5%E8%AE%A8%E8%AE%BA.tex}{L17-例-3}。
\begin{solve}
这一矩阵正是原讲义“求若当标准形过渡矩阵”处中讨论过的矩阵. 我们已经求得它的若当标准形为
    \[
        J=\begin{pmatrix}
                3 & 1 & 0 \\
                0 & 3 & 0 \\
                0 & 0 & 1
            \end{pmatrix},
    \]
    即它有一个关于特征值$3$的二阶若当块$J_2(3)$，以及一个关于特征值$1$的一阶若当块$J_1(1)$. 根据前面的讨论可知，每一个若当块$J_r(\lambda_i)$都对应一个初等因子$(\lambda-\lambda_i)^r$，因此矩阵$A$的初等因子组为
    \[
        (\lambda-3)^2,\ \lambda-1.
    \]
\end{solve}

\begin{exercise}[补充训练 L17-例-4]\label{prob:L17-例-4}
求矩阵 $A = \begin{pmatrix}
            0 & 0 & 0 \\ 1 & 0 & 2 \\ 2 & 1 & -1
        \end{pmatrix}$ 的极小多项式.
\end{exercise}
\noindent 来源：\href{https://github.com/yhwu-is/Linear-Algebra-Left-Undone/blob/new/%E8%AE%B2%E4%B9%89/%E4%B8%93%E9%A2%98/17%20%E5%A4%9A%E9%A1%B9%E5%BC%8F%E7%9A%84%E8%BF%9B%E4%B8%80%E6%AD%A5%E8%AE%A8%E8%AE%BA.tex}{L17-例-4}。
\begin{solve}
对于矩阵$A$，我们依次考察方程组
    \[
        a_0I+A=O,\qquad a_0I+a_1A+A^2=O,\qquad a_0I+a_1A+a_2A^2+A^3=O.
    \]
    先计算
    \[
        A^2=\begin{pmatrix}
                0 & 0  & 0  \\
                4 & 2  & -2 \\
                -1 & -1 & 3
            \end{pmatrix},
        \qquad
        A^3=\begin{pmatrix}
                0 & 0  & 0  \\
                -2 & -2 & 6  \\
                5 & 3  & -5
            \end{pmatrix}.
    \]
    读者不难验证，前两个方程组均无解. 根据三个方程组可以解得 $a_0 = 0, a_1 = -2, a_2 = 1$. 于是 $A^3 + A^2 - 2A=O$. 因为次数小于 $3$ 的首一多项式都不能化零 $A$，所以 $A$ 的极小多项式为 $\lambda^3+\lambda^2-2\lambda$.
\end{solve}

\section{补充训练：内积空间与正交化}

本节包含 8 道训练题，建议在完成相关专题真题后选择性练习。

\begin{exercise}[补充训练 L19-例-1]\label{prob:L19-例-1}
验证下面的例子确实满足内积的定义：
    \begin{enumerate}
        \item $\mathbf{F}^{n}$ 上的欧几里得内积定义为：
              \[\left\langle (w_1, \ldots, w_n), (z_1, \ldots, z_n)\right\rangle = w_1\overline{z_1} + \cdots + w_n\overline{z_n} = \vec{w}\overline{\vec{z}}^{\mathrm{T}}.\]
              对应范数：
              \[\left\lVert (z_1, \ldots, z_n) \right\rVert  = \sqrt{\lvert z_1 \rvert^2 + \cdots + \lvert z_n \rvert^2}.\]

        \item \label{item:23:内积和范数:2}
              定义在 $ \left[-1, 1\right] $ 上的连续实值函数构成的向量空间可定义内积如下：
              \[\left\langle f, g\right\rangle = \int_{-1}^1f(x)g(x)\,\mathrm{d}x.\]
              对应范数：
              \[\left\lVert f \right\rVert = \sqrt{\int_{-1}^1(f(x))^2\,\mathrm{d}x}.\]
    \end{enumerate}
\end{exercise}
\noindent 来源：\href{https://github.com/yhwu-is/Linear-Algebra-Left-Undone/blob/new/%E8%AE%B2%E4%B9%89/%E4%B8%93%E9%A2%98/19%20%E5%86%85%E7%A7%AF%E7%A9%BA%E9%97%B4.tex}{L19-例-1}。
\begin{solve}
\begin{enumerate}
        \item 对于$u=(w_1,\ldots,w_n),v=(z_1,\ldots,z_n),w=(t_1,\ldots,t_n)\in\mathbf{F}^n$，$\lambda\in\mathbf{F}$，有
            \begin{enumerate}
                \item 正定性：\[\langle v,v\rangle = \sum_{i=1}^n z_i\overline{z_i} = \sum_{i=1}^n \lvert z_i\rvert^2\geqslant 0,\]
                且
                \[\langle v,v\rangle = 0 \iff z_i=0\enspace(i=1,2,\ldots,n) \iff v=\vec{0}.\]
                故满足正定性.
                \item 第一个位置的加性：\[\langle u+v,w\rangle = \sum_{i=1}^n(w_i+z_i)\overline{t_i}=\sum_{i=1}^n w_i\overline{t_i}+\sum_{i=1}^n z_i\overline{t_i}=\langle u,w\rangle+\langle v,w\rangle.\]
                故满足第一个位置的加性；
                \item 第一个位置的齐性：\[\langle \lambda u,v\rangle = \sum_{i=1}^n(\lambda w_i\overline{z_i})=\lambda\sum_{i=1}^nw_i\overline{z_i}=\lambda\langle u,v\rangle.\]
                故满足第一个位置的齐性；
                \item 共轭对称性：\[\langle u,v\rangle = \sum_{i=1}^nw_i\overline{z_i}=\overline{\sum_{i=1}^nz_i\overline{w_i}}=\overline{\langle v,u\rangle}.\]
                故满足共轭对称性.
            \end{enumerate}
            从而确实满足内积定义.
        \item 对于 $f, g, h \in C[-1,1]$，$\lambda \in \mathbf{R}$，有
            \begin{enumerate}
                \item 正定性：\[\langle f,f\rangle = \int_{-1}^1(f(x))^2\,\mathrm{d}x\geqslant 0,\]
                且
                \[\langle f,f\rangle = 0 \iff f(x)=0\enspace(\forall x\in[-1,1]) \iff f=\vec{0}.\]
                故满足正定性.
                \item 第一个位置的加性：
                \begin{align*}
                    \langle f+g,h\rangle &= \int_{-1}^1(f(x)+g(x))h(x)\,\mathrm{d}x=\int_{-1}^1f(x)h(x)\,\mathrm{d}x+\int_{-1}^1g(x)h(x)\,\mathrm{d}x \\
                    &=\langle f,h\rangle+\langle g,h\rangle.
                \end{align*}
                故满足第一个位置的加性；
                \item 第一个位置的齐性：\[\langle \lambda f,g\rangle = \int_{-1}^1(\lambda f(x))g(x)\,\mathrm{d}x=\lambda\int_{-1}^1f(x)g(x)\,\mathrm{d}x=\lambda\langle f,g\rangle.\]
                故满足第一个位置的齐性；
                \item 共轭对称性：\[\langle f,g\rangle = \int_{-1}^1f(x)g(x)\,\mathrm{d}x=\int_{-1}^1g(x)f(x)\,\mathrm{d}x=\langle g,f\rangle.\]
                故满足共轭对称性（注意到是实值函数，因此共轭对称性退化为对称性）.
            \end{enumerate}
            从而确实满足内积定义.
    \end{enumerate}
\end{solve}

\begin{exercise}[补充训练 L19-例-2]\label{prob:L19-例-2}
证明下列式子成立：
    \begin{enumerate}
        \item $\mathbf{F} = \mathbf{R}$ 时：
              \begin{gather}
                  \label{eq:23:内积和范数的性质1}
                  \langle u, v\rangle = \frac{1}{4}\left( \lVert u + v \rVert^2 - \lVert u - v \rVert^2\right) \\
                  \label{eq:23:内积和范数的性质2}
                  \langle Tu, v\rangle + \langle Tv, u\rangle = \frac{1}{2}\left(\langle T(u + v), u + v\rangle - \langle T(u - v), u - v\rangle\right)
              \end{gather}

        \item $\mathbf{F} = \mathbf{C}$ 时：
              \begin{gather}
                  \label{eq:23:内积和范数的性质3}
                  \langle u, v\rangle = \frac{1}{4}\left((\lVert u + v \rVert^2 - \lVert u - v \rVert^2) + \mathrm{i}(\lVert u + \mathrm{i} v \rVert^2 - \lVert u - \mathrm{i} v \rVert^2)\right) \\
                  \label{eq:23:内积和范数的性质4}
                  \begin{aligned}
                      \langle Tu, v\rangle ={} & \frac{1}{4}  ((\langle T(u + v), u + v\rangle - \langle T(u - v), u - v\rangle)     \\
                                               & + \mathrm{i}(\langle T(u + \mathrm{i} v), u + \mathrm{i} v\rangle + \langle T(u - \mathrm{i} v), u - \mathrm{i} v\rangle)).
                  \end{aligned}
              \end{gather}
    \end{enumerate}
\end{exercise}
\noindent 来源：\href{https://github.com/yhwu-is/Linear-Algebra-Left-Undone/blob/new/%E8%AE%B2%E4%B9%89/%E4%B8%93%E9%A2%98/19%20%E5%86%85%E7%A7%AF%E7%A9%BA%E9%97%B4.tex}{L19-例-2}。
\begin{solve}
\begin{enumerate}
        \item 利用原讲义“勾股定理展开”处 将右侧展开得证（由于是实内积空间，因此 $\operatorname{Re}(\langle u, v \rangle) = \langle u, v \rangle$）.

        \item 我们展开等式右侧的两项，利用线性算子 $T$ 的线性性：
        \begin{align*}
            & \quad \langle T(u + v), u + v\rangle - \langle T(u - v), u - v\rangle \\
            & = \left(\langle Tu, u \rangle + \langle Tu, v \rangle + \langle Tv, u \rangle + \langle Tv, v \rangle\right) - \left(\langle Tu, u \rangle - \langle Tu, v \rangle - \langle Tv, u \rangle + \langle Tv, v \rangle\right) \\
            & = 2 (\langle Tu, v \rangle + \langle Tv, u \rangle),
        \end{align*}
        两边同除以 $2$ 即得.

        \item 利用原讲义“勾股定理展开”处 将右侧展开：
        \begin{gather*}
            \lVert u + v \rVert^2 - \lVert u - v \rVert^2 = 4 \operatorname{Re}(\langle u, v \rangle), \\
            \lVert u + \mathrm{i} v \rVert^2 - \lVert u - \mathrm{i} v \rVert^2 = 4 \operatorname{Re}(\langle u, \mathrm{i} v \rangle) = 4 \operatorname{Re}(-\mathrm{i} \langle u, v \rangle) = 4 \operatorname{Im}(\langle u, v \rangle).
        \end{gather*}
        故等号右侧为 $ \operatorname{Re}(\langle u, v \rangle) + \mathrm{i} \operatorname{Im}(\langle u, v \rangle) = \langle u, v \rangle$. 得证.

        \item
        \begin{gather*}
            \langle T(u + v), u + v\rangle - \langle T(u - v), u - v\rangle = 2\langle Tu, v \rangle + 2\langle Tv, u \rangle,\\
            \langle T(u + \mathrm{i} v), u + \mathrm{i} v\rangle - \langle T(u - \mathrm{i} v), u - \mathrm{i} v\rangle = 2\langle Tu, \mathrm{i} v \rangle + 2\langle \mathrm{i} Tv, u \rangle = 2\mathrm{i} (\langle Tv, u \rangle - \langle Tu, v \rangle),
        \end{gather*}
        代入等式右边可得
        \begin{align*}
            &\quad \frac{1}{4}\left(\left(2\langle Tu, v \rangle + 2\langle Tv, u \rangle\right) + \mathrm{i} \left(2\mathrm{i} (\langle Tv, u \rangle - \langle Tu, v \rangle)\right)\right) \\
            & = \frac{1}{4}\left(2\langle Tu, v \rangle + 2\langle Tv, u \rangle - 2\langle Tv, u \rangle + 2\langle Tu, v \rangle\right) \\
            & = \langle Tu, v \rangle.
        \end{align*}
        得证.
    \end{enumerate}
\end{solve}

\begin{exercise}[补充训练 L19-例-3]\label{prob:L19-例-3}
设 $e_1, \ldots , e_m$ 是 $ V $ 的标准正交组. 设 $ v \in V $. 证明
    \[ \lVert v \rVert^2 = \lvert \langle v, e_1\rangle \rvert^2 + \cdots + \lvert \langle v, e_m\rangle \rvert^2 \]
    当且仅当 $ v \in \operatorname{span}(e_1, \ldots , e_m)$.
\end{exercise}
\noindent 来源：\href{https://github.com/yhwu-is/Linear-Algebra-Left-Undone/blob/new/%E8%AE%B2%E4%B9%89/%E4%B8%93%E9%A2%98/19%20%E5%86%85%E7%A7%AF%E7%A9%BA%E9%97%B4.tex}{L19-例-3}。
\begin{solve}
记 $W = \operatorname{span}(e_1, \ldots , e_m)$.
    \begin{enumerate}
        \item 必要性：构造向量 $u = \displaystyle\sum_{i=1}^{m} \langle v, e_i \rangle e_i$. 显然 $u \in W$.

        令 $w = v - u$，注意到：
        \[
            \langle w, e_j \rangle = \left\langle v - \sum_{i=1}^{m} \langle v, e_i \rangle e_i, e_j \right\rangle = \langle v, e_j \rangle - \sum_{i=1}^{m} \langle v, e_i \rangle \delta_{ij} = \langle v, e_j \rangle - \langle v, e_j \rangle = 0.
        \]
        因此 $w$ 与所有 $e_j$ 正交，由于 $u$ 为 $e_1, \ldots , e_m$ 的线性组合，故 $w$ 与 $u$ 也正交.

        计算 $\lVert w \rVert^2$:
        \begin{align*}
            \|w\|^2 &= \langle w, v-u \rangle = \langle w, v \rangle - \langle w, u \rangle \\
            &= \langle v-u, v \rangle = \langle v, v \rangle - \langle u, v \rangle \\
            &= \|v\|^2 - \sum_{i=1}^{m} \langle v, e_i \rangle \langle e_i, v \rangle \\
            &= \|v\|^2 - \sum_{i=1}^{m} \langle v, e_i \rangle \overline{\langle v, e_i \rangle} \\
            &= \|v\|^2 - \sum_{i=1}^{m} |\langle v, e_i \rangle|^2.
        \end{align*}
        根据假设 $\lVert v \rVert^2 = \displaystyle\sum_{i=1}^{m} \lvert \langle v, e_i \rangle \rvert^2$，故 $\lVert w \rVert^2 = 0$.
        由范数正定性 $w = \vec{0}$，即 $v = u$.

        由于 $u \in W$，所以 $v \in \operatorname{span}(e_1, \ldots , e_m)$. 必要性得证.

        \item 充分性：对子空间 $W$ 使用原讲义“Fourier展开”处 即可. 充分性得证.
    \end{enumerate}
\end{solve}

\begin{exercise}[补充训练 L19-例-4]\label{prob:L19-例-4}
求 $\mathbf{R}[x]_2$ 的一组标准正交基，内积定义为 $\langle p, q \rangle = \displaystyle\int_{-1}^1 p(x)q(x)\,\mathrm{d}x$.
\end{exercise}
\noindent 来源：\href{https://github.com/yhwu-is/Linear-Algebra-Left-Undone/blob/new/%E8%AE%B2%E4%B9%89/%E4%B8%93%E9%A2%98/19%20%E5%86%85%E7%A7%AF%E7%A9%BA%E9%97%B4.tex}{L19-例-4}。
\begin{solve}
取 $\mathbf{R}[x]_3$ 的一组基 $\{v_1, v_2, v_3\} = \{1, x, x^2\}$.

    我们使用 Gram-Schmidt 过程.

    首先计算 $e_1$:
    \[ \lVert v_1 \rVert^2 = \langle 1, 1 \rangle = \int_{-1}^1 1 \,\mathrm{d}x = 2, e_1 = \frac{v_1}{\lVert v_1 \rVert} = \frac{1}{\sqrt{2}}. \]

    正交化：
    \[ u_2 = v_2 - \langle v_2, e_1 \rangle e_1, \]
    计算投影系数：
    \[ \langle v_2, e_1 \rangle = \left\langle x, \frac{1}{\sqrt{2}} \right\rangle = \frac{1}{\sqrt{2}} \int_{-1}^1 x \,\mathrm{d}x = 0, \]
    因此 $u_2 = x - 0 \cdot e_1 = x$.

    单位化：
    \[ \lVert u_2 \rVert^2 = \langle x, x \rangle = \int_{-1}^1 x^2 \,\mathrm{d}x = \left.\left(\frac{1}{3}x^3\right)\right\vert_{-1}^1= \frac{2}{3}, e_2 = \frac{u_2}{\lVert u_2 \rVert} = \frac{x}{\sqrt{2/3}} = \sqrt{\frac{3}{2}} x. \]

    正交化：
    \[ u_3 = v_3 - \langle v_3, e_1 \rangle e_1 - \langle v_3, e_2 \rangle e_2, \]
    计算投影系数：
    \begin{gather*}
        \langle v_3, e_1 \rangle = \left\langle x^2, \frac{1}{\sqrt{2}} \right\rangle = \frac{1}{\sqrt{2}} \int_{-1}^1 x^2 \,\mathrm{d}x = \frac{\sqrt{2}}{3}, \\
        \langle v_3, e_2 \rangle = \left\langle x^2, \sqrt{\frac{3}{2}} x \right\rangle = \sqrt{\frac{3}{2}} \int_{-1}^1 x^3 \,\mathrm{d}x = 0,
    \end{gather*}
    计算 $u_3$:
    \[ u_3 = x^2 - \left(\frac{\sqrt{2}}{3}\right) e_1 - 0 = x^2 - \frac{\sqrt{2}}{3} \cdot \frac{1}{\sqrt{2}} = x^2 - \frac{1}{3}. \]

    单位化：
    \[ \lVert u_3 \rVert^2 = \left\langle x^2 - \frac{1}{3}, x^2 - \frac{1}{3} \right\rangle = \int_{-1}^1 \left(x^4 - \frac{2}{3}x^2 + \frac{1}{9}\right) \,\mathrm{d}x = \frac{8}{45}, \]
    \[ e_3 = \frac{u_3}{\lVert u_3 \rVert} = \frac{x^2 - 1/3}{\sqrt{8/45}} = \sqrt{\frac{45}{8}} \left(x^2 - \frac{1}{3}\right) = \frac{\sqrt{10}}{4} (3x^2 - 1). \]

    综上，$\mathbf{R}[x]_3$ 的一组标准正交基是 $\left\{e_1, e_2, e_3\right\}$，其中
    \[ e_1(x) = \frac{1}{\sqrt{2}}, \quad e_2(x) = \sqrt{\frac{3}{2}} x, \quad e_3(x) = \frac{\sqrt{10}}{4} (3x^2 - 1). \]
\end{solve}

\begin{exercise}[补充训练 L19-例-5]\label{prob:L19-例-5}
求矩阵
    \[
        A = \begin{pmatrix}
            1 & 1 & 0 \\
            1 & 0 & 1 \\
            0 & 1 & 1
        \end{pmatrix}
    \]
    的 QR 分解.
\end{exercise}
\noindent 来源：\href{https://github.com/yhwu-is/Linear-Algebra-Left-Undone/blob/new/%E8%AE%B2%E4%B9%89/%E4%B8%93%E9%A2%98/19%20%E5%86%85%E7%A7%AF%E7%A9%BA%E9%97%B4.tex}{L19-例-5}。
\begin{solve}
记 $A$ 的列向量为
    \[
        \alpha_1 = \begin{pmatrix}
            1 \\
            1 \\
            0
        \end{pmatrix}, \enspace
        \alpha_2 = \begin{pmatrix}
            1 \\
            0 \\
            1
        \end{pmatrix}, \enspace
        \alpha_3 = \begin{pmatrix}
            0 \\
            1 \\
            1
        \end{pmatrix}.
    \]
    对 $\alpha_1, \alpha_2, \alpha_3$ 施行 Gram-Schmidt 过程.

    先有
    \[
        u_1 = \alpha_1, \enspace q_1 = \frac{\alpha_1}{\lVert \alpha_1 \rVert}
        = \frac{1}{\sqrt{2}}
        \begin{pmatrix}
            1 \\
            1 \\
            0
        \end{pmatrix}.
    \]

    再看第二个向量，
    \[
        \langle \alpha_2, q_1 \rangle = \frac{1}{\sqrt{2}}, \enspace
        u_2 = \alpha_2 - \langle \alpha_2, q_1 \rangle q_1
        =
        \begin{pmatrix}
            \frac{1}{2} \\
            -\frac{1}{2} \\
            1
        \end{pmatrix},
    \]
    从而
    \[
        q_2 = \frac{u_2}{\lVert u_2 \rVert}
        = \frac{1}{\sqrt{6}}
        \begin{pmatrix}
            1 \\
            -1 \\
            2
        \end{pmatrix}.
    \]

    对第三个向量，
    \[
        \langle \alpha_3, q_1 \rangle = \frac{1}{\sqrt{2}}, \enspace
        \langle \alpha_3, q_2 \rangle = \frac{1}{\sqrt{6}},
    \]
    因而
    \[
        u_3 = \alpha_3 - \langle \alpha_3, q_1 \rangle q_1 - \langle \alpha_3, q_2 \rangle q_2
        = \frac{2}{3}
        \begin{pmatrix}
            -1 \\
            1 \\
            1
        \end{pmatrix},
    \]
    从而
    \[
        q_3 = \frac{u_3}{\lVert u_3 \rVert}
        = \frac{1}{\sqrt{3}}
        \begin{pmatrix}
            -1 \\
            1 \\
            1
        \end{pmatrix}.
    \]

    因而
    \[
        Q = \begin{pmatrix}
            \dfrac{1}{\sqrt{2}} & \dfrac{1}{\sqrt{6}}  & -\dfrac{1}{\sqrt{3}} \\
            \dfrac{1}{\sqrt{2}} & -\dfrac{1}{\sqrt{6}} & \dfrac{1}{\sqrt{3}}  \\
            0                   & \dfrac{2}{\sqrt{6}}  & \dfrac{1}{\sqrt{3}}
        \end{pmatrix}, \enspace
        R = \begin{pmatrix}
            \sqrt{2} & \dfrac{1}{\sqrt{2}} & \dfrac{1}{\sqrt{2}} \\
            0        & \dfrac{\sqrt{6}}{2} & \dfrac{1}{\sqrt{6}} \\
            0        & 0                   & \dfrac{2}{\sqrt{3}}
        \end{pmatrix}.
    \]
    这里
    \[
        r_{11} = \lVert \alpha_1 \rVert = \sqrt{2}, \enspace
        r_{22} = \lVert u_2 \rVert = \frac{\sqrt{6}}{2}, \enspace
        r_{33} = \lVert u_3 \rVert = \frac{2}{\sqrt{3}}.
    \]
    直接验证可得 $A = QR$，故这就是所求的 QR 分解.
\end{solve}

\begin{exercise}[补充训练 L19-例-7]\label{prob:L19-例-7}
定义在 $ V = \mathbf{R}^3 $ 上的运算
    \[ \langle \vec{x}, \vec{y} \rangle_V = x_1 y_1 + x_2 y_2 + (x_2 + x_3)(y_2 + y_3) \]
    其中 $ \vec{x} = (x_1, x_2, x_3),\enspace \vec{y} = (y_1, y_2, y_3) $.
    \begin{enumerate}
        \item 验证 $ \langle \cdot, \cdot \rangle_V $ 是 $ \mathbf{R}^3 $ 上的一个内积；

        \item 求 $ \mathbf{R}^3 $ 在 $ \langle \cdot, \cdot \rangle_V $ 下的一组标准正交基；

        \item 求 $ \vec{\beta} \in V $ 使得 $ \forall \vec{x} \in V,\enspace x_1 + 2x_2 = \langle \vec{x}, \vec{\beta} \rangle_V $.
    \end{enumerate}
\end{exercise}
\noindent 来源：\href{https://github.com/yhwu-is/Linear-Algebra-Left-Undone/blob/new/%E8%AE%B2%E4%B9%89/%E4%B8%93%E9%A2%98/19%20%E5%86%85%E7%A7%AF%E7%A9%BA%E9%97%B4.tex}{L19-例-7}。
\begin{solve}
\begin{enumerate}
        \item 设 $\vec{x}, \vec{y}, \vec{z} \in \mathbf{R}^3$，$\lambda \in \mathbf{R}$. 先验证线性性与对称性：
        \begin{align*}
            \langle \vec{x} + \vec{y}, \vec{z} \rangle_V
             & = (x_1 + y_1)z_1 + (x_2 + y_2)z_2 + (x_2 + y_2 + x_3 + y_3)(z_2 + z_3) \\
             & = \langle \vec{x}, \vec{z} \rangle_V + \langle \vec{y}, \vec{z} \rangle_V, \\
            \langle \lambda \vec{x}, \vec{y} \rangle_V
             & = \lambda x_1 y_1 + \lambda x_2 y_2 + (\lambda x_2 + \lambda x_3)(y_2 + y_3)
            = \lambda \langle \vec{x}, \vec{y} \rangle_V.
        \end{align*}
        交换 $\vec{x}, \vec{y}$ 后定义式不变，因此 $\langle \vec{x}, \vec{y} \rangle_V = \langle \vec{y}, \vec{x} \rangle_V$.

        再看正定性：$\langle \vec{x}, \vec{x} \rangle_V = x_1^2 + x_2^2 + (x_2 + x_3)^2 \geqslant 0$. 若 $\langle \vec{x}, \vec{x} \rangle_V = 0$，则 $x_1^2 = 0$，$x_2^2 = 0$，$(x_2 + x_3)^2 = 0$，从而 $x_1 = 0$，$x_2 = 0$，$x_3 = 0$. 故 $\vec{x} = \vec{0}$.

        综上，$\langle \cdot, \cdot \rangle_V$ 是 $\mathbf{R}^3$ 上的一个内积.

        \item 取 $\mathbf{R}^3$ 的标准基 $v_1 = (1, 0, 0)^{\mathrm{T}}, \enspace v_2 = (0, 1, 0)^{\mathrm{T}}, \enspace v_3 = (0, 0, 1)^{\mathrm{T}}$.
        对其施行 Gram-Schmidt 过程（具体过程不展开），得到在内积 $\langle \cdot, \cdot \rangle_V$ 下的一组标准正交基可以取为
        \[
            \left(1, 0, 0\right)^{\mathrm{T}}, \enspace \left(0, \frac{1}{\sqrt{2}}, 0\right)^{\mathrm{T}}, \enspace \left(0, -\frac{1}{\sqrt{2}}, \sqrt{2}\right)^{\mathrm{T}}.
        \]

        \item 设 $\varphi \colon V \to \mathbf{R}$ 定义为
        \[
            \varphi(\vec{x}) = x_1 + 2x_2.
        \]
        由原讲义“Riesz”处，存在唯一的 $\vec{\beta} \in V$ 使得对任意 $\vec{x} \in V$ 都有
        \[
            \varphi(\vec{x}) = \langle \vec{x}, \vec{\beta} \rangle_V.
        \]
        利用上一步得到的标准正交基 $e_1, e_2, e_3$，可得
        \[
            \vec{\beta} = \varphi(e_1)e_1 + \varphi(e_2)e_2 + \varphi(e_3)e_3 = (1, 2, -2)^{\mathrm{T}}.
        \]
    \end{enumerate}
\end{solve}

\begin{exercise}[补充训练 S-QR-1]\label{prob:S-QR-1}
对矩阵 $A=\begin{pmatrix}1&0&0\\0&1&0\\0&-4&1\end{pmatrix}$ 的列依次进行施密特正交化，给出正交基、标准正交基与 $A=QR$ 分解，并说明为何每一步均不会得到零向量。
\end{exercise}
\noindent 编者补充：S-QR-1；相关主题见原讲义对应讲，题目数据与解答由本选编另行构造。
\begin{solve}
设列为 $a_1,a_2,a_3$。按
\[u_j=a_j-\sum_{i<j}\frac{u_i^{\mathrm T}a_j}{u_i^{\mathrm T}u_i}u_i\]
依次计算，得到 $U=(u_1,u_2,u_3)=\begin{pmatrix}1&0&0\\0&1&\frac{4}{17}\\0&-4&\frac{1}{17}\end{pmatrix}$，且
\[U^{\mathrm T}U=\operatorname{diag}(1,17,\frac{1}{17}).\]
所以 $Q=\left(u_1/\sqrt{1},u_2/\sqrt{17},u_3/\sqrt{\frac{1}{17}}\right)$ 是标准正交矩阵，而
\[R=Q^{\mathrm T}A=\begin{pmatrix}\frac{1}{\sqrt{1}}&0&0\\0&\frac{17}{\sqrt{17}}&\frac{-4}{\sqrt{17}}\\0&0&\frac{\frac{1}{17}}{\sqrt{\frac{1}{17}}}\end{pmatrix}.\]
由构造可得 $A=QR$。这里投影系数的分母是未归一化向量的长度平方，不能省略。因为 $\det A=1$，原列向量线性无关；若某个 $u_j=0$，则 $a_j$ 属于前面的列的张成空间，矛盾。
\end{solve}

\begin{exercise}[补充训练 S-LS-1]\label{prob:S-LS-1}
对数据点 $(-1,1),(0,2),(2,-1)$，用直线 $y=c_0+c_1t$ 作最小二乘拟合。写出设计矩阵、正规方程、拟合值与残差，并计算最小残差平方和。
\end{exercise}
\noindent 编者补充：S-LS-1；相关主题见原讲义对应讲，题目数据与解答由本选编另行构造。
\begin{solve}
设计矩阵与观测向量为 $A=\begin{pmatrix}1&-1\\1&0\\1&2\end{pmatrix}$、$b=\begin{pmatrix}1\\2\\-1\end{pmatrix}$。目标为 $\min_c\|Ac-b\|_2^2$，正规方程为
\[A^{\mathrm T}Ac=A^{\mathrm T}b,\qquad \begin{pmatrix}3&1\\1&5\end{pmatrix}c=\begin{pmatrix}2\\-3\end{pmatrix}.\]
因为两列线性无关，$A^{\mathrm T}A$ 正定且可逆，得到唯一解 $c=\begin{pmatrix}\frac{13}{14}\\-\frac{11}{14}\end{pmatrix}$。进一步，
\[\widehat b=Ac=\begin{pmatrix}\frac{12}{7}\\\frac{13}{14}\\-\frac{9}{14}\end{pmatrix},\qquad e=b-\widehat b=\begin{pmatrix}-\frac{5}{7}\\\frac{15}{14}\\-\frac{5}{14}\end{pmatrix},\qquad \|e\|_2^2=\frac{25}{14}.\]
直接检验 $A^{\mathrm T}e=0$，说明残差正交于列空间。对任意 $d$，有 $\|b-A(c+d)\|^2=\|e\|^2+\|Ad\|^2$，从而上述驻点确为全局最小点，并由满列秩得到唯一性。
\end{solve}

\section{补充训练：内积空间上的算子}

本节包含 3 道训练题，建议在完成相关专题真题后选择性练习。

\begin{exercise}[补充训练 L20-例-8]\label{prob:L20-例-8}
求矩阵
    \[
        A = \begin{pmatrix}
            1 & 1  \\
            1 & 1  \\
            1 & -1
        \end{pmatrix}
    \]
    的奇异值分解.
\end{exercise}
\noindent 来源：\href{https://github.com/yhwu-is/Linear-Algebra-Left-Undone/blob/new/%E8%AE%B2%E4%B9%89/%E4%B8%93%E9%A2%98/20%20%E5%86%85%E7%A7%AF%E7%A9%BA%E9%97%B4%E4%B8%8A%E7%9A%84%E7%AE%97%E5%AD%90.tex}{L20-例-8}。
\begin{solve}
我们先总结一套求矩阵奇异值分解的通用步骤，然后在本题中逐步执行.

    \begin{enumerate}
        \item 先计算 $A^{\mathrm{T}}A$（或复数域中的 $A^*A$），求出其全部特征值与一组标准正交特征向量.

        \item 将这些特征值开平方，得到奇异值，并按从大到小排序.

        \item 与各个正奇异值 $s_i$ 对应的单位特征向量 $v_i$ 就是右奇异向量，再由公式
              \[
                  u_i = \frac{1}{s_i}Av_i
              \]
              求出左奇异向量.

        \item 若 $U$ 或 $V$ 的列向量还不够构成一组标准正交基，就将其补齐.

        \item 以左奇异向量为列组成 $U$，以右奇异向量为列组成 $V$，再把奇异值按顺序填入 $\Sigma$ 的主对角线，即得
              \[
                  A = U\Sigma V^{\mathrm{T}}.
              \]
    \end{enumerate}

    下面按此步骤计算.

    \begin{enumerate}
        \item 先求 $A^{\mathrm{T}}A$ 的特征值与特征向量，不难得到其特征值为 $4$ 和 $2$，对应的特征向量分别为 $(1, 1)^{\mathrm{T}}$ 和 $(1, -1)^{\mathrm{T}}$. 归一化后得到一组标准正交特征向量
              \[
                  v_1 = \frac{1}{\sqrt{2}}\begin{pmatrix}
                      1 \\
                      1
                  \end{pmatrix},\qquad
                  v_2 = \frac{1}{\sqrt{2}}\begin{pmatrix}
                      1  \\
                      -1
                  \end{pmatrix}.
              \]

        \item 因此 $A$ 的奇异值为 $s_1 = \sqrt{4} = 2, s_2 = \sqrt{2}$.

        \item 由 $u_i = \dfrac{1}{s_i}Av_i$ 求出左奇异向量为
              \[ u_1 = \frac{1}{2}\begin{pmatrix}
                      1 \\
                      1 \\
                      0
                  \end{pmatrix}, \qquad u_2 = \frac{1}{\sqrt{2}}\begin{pmatrix}
                      0 \\
                      0 \\
                      1
                  \end{pmatrix}. \]

        \item 由于 $A$ 是 $3 \times 2$ 矩阵，$U$ 还需要第三个列向量. 我们只需补一个与 $u_1, u_2$ 都正交的单位向量即可，例如取
              \[
                  u_3 = \frac{1}{\sqrt{2}}
                  \begin{pmatrix}
                      1  \\
                      -1 \\
                      0
                  \end{pmatrix}.
              \]
              于是
              \[
                  U =
                  \begin{pmatrix}
                      \frac{1}{\sqrt{2}} & 0 & \frac{1}{\sqrt{2}}  \\
                      \frac{1}{\sqrt{2}} & 0 & -\frac{1}{\sqrt{2}} \\
                      0                  & 1 & 0
                  \end{pmatrix},\qquad
                  V =
                  \frac{1}{\sqrt{2}}
                  \begin{pmatrix}
                      1 & 1  \\
                      1 & -1
                  \end{pmatrix}.
              \]

        \item 最后构造
              \[
                  \Sigma =
                  \begin{pmatrix}
                      2 & 0        \\
                      0 & \sqrt{2} \\
                      0 & 0
                  \end{pmatrix}.
              \]
              因而得到 $A = U\Sigma V^{\mathrm{T}}$，至此，奇异值分解求得.
    \end{enumerate}
\end{solve}

\begin{exercise}[补充训练 S-SPECTRAL-1]\label{prob:S-SPECTRAL-1}
设实对称矩阵 $A=\begin{pmatrix}\frac{23}{25}&-\frac{36}{25}&0\\-\frac{36}{25}&\frac{2}{25}&0\\0&0&4\end{pmatrix}$。求正交对角化，判断其正定性，求二次型 $x^{\mathrm T}Ax$ 在 $x^{\mathrm T}x=1$ 下的最大值、最小值及取到这些值的向量。
\end{exercise}
\noindent 编者补充：S-SPECTRAL-1；相关主题见原讲义对应讲，题目数据与解答由本选编另行构造。
\begin{solve}
取正交矩阵 $Q=\begin{pmatrix}\frac{3}{5}&-\frac{4}{5}&0\\\frac{4}{5}&\frac{3}{5}&0\\0&0&1\end{pmatrix}$，相乘可验证
\[Q^{\mathrm T}AQ=\operatorname{diag}(-1,2,4).\]
由于既有正特征值又有负特征值，矩阵不定，正、负惯性指数分别为 $2,1$。令 $y=Q^{\mathrm T}x$，正交变换保持单位球面，因此
\[x^{\mathrm T}Ax=(-1)y_1^2+2y_2^2+4y_3^2,\qquad y_1^2+y_2^2+y_3^2=1.\]
这是三个特征值的凸组合，最小值为 $-1$，最大值为 $4$。分别在 $x=\pm Qe_1$ 与 $x=\pm Qe_3$ 处取到；因为这些特征值互异，等号条件没有额外方向。
\end{solve}

\begin{exercise}[补充训练 S-SVD-1]\label{prob:S-SVD-1}
对 $A=\begin{pmatrix}\frac{36}{25}&\frac{48}{25}\\\frac{48}{25}&\frac{64}{25}\\-\frac{8}{5}&\frac{6}{5}\end{pmatrix}$ 求一个紧致奇异值分解，给出秩、两个奇异值、谱范数及 Moore--Penrose 广义逆，并验证 $A^+A=I_2$。
\end{exercise}
\noindent 编者补充：S-SVD-1；相关主题见原讲义对应讲，题目数据与解答由本选编另行构造。
\begin{solve}
取
\[U=\begin{pmatrix}\frac{3}{5}&0\\\frac{4}{5}&0\\0&1\end{pmatrix},\quad\Sigma=\operatorname{diag}(4,2),\quad V=\begin{pmatrix}\frac{3}{5}&-\frac{4}{5}\\\frac{4}{5}&\frac{3}{5}\end{pmatrix}.\]
直接核对 $U^{\mathrm T}U=I_2$、$V^{\mathrm T}V=I_2$ 和 $A=U\Sigma V^{\mathrm T}$，所以这就是紧致奇异值分解。两个奇异值均为正，故秩为 $2$；谱范数为较大的奇异值 $4$。广义逆为
\[A^+=V\Sigma^{-1}U^{\mathrm T}=\begin{pmatrix}\frac{9}{100}&\frac{3}{25}&-\frac{2}{5}\\\frac{3}{25}&\frac{4}{25}&\frac{3}{10}\end{pmatrix}.\]
于是 $A^+A=V\Sigma^{-1}\Sigma V^{\mathrm T}=I_2$，而 $AA^+=UU^{\mathrm T}$ 是到 $A$ 列空间的正交投影。后者通常不等于 $I_3$，这与矩阵只有两列相符。
\end{solve}

\section{补充训练：二次型}

本节包含 12 道训练题，建议在完成相关专题真题后选择性练习。

\begin{exercise}[补充训练 L22-例-1]\label{prob:L22-例-1}
回答以下问题：
    \begin{enumerate}
        \item 已知$A$是一个$n$阶矩阵，则$A$为反对称矩阵的充要条件是对任意$n$元列向量$X$都有$X^\mathrm{T}AX=0$；

        \item 若二次型$f(x_1,x_2,\ldots,x_n)=X^\mathrm{T}AX$对任意的$n$元列向量$X$都有$f(x_1,x_2,\ldots,x_n)=0$，证明：$A=O$；

        \item 设二次型$f(x_1,x_2,\ldots,x_n)=X^\mathrm{T}AX,\enspace g(x_1,x_2,\ldots,x_n)=X^\mathrm{T}BX$.\\
              证明：若$f(x_1,x_2,\ldots,x_n)=g(x_1,x_2,\ldots,x_n)$，则$A=B$.
    \end{enumerate}
\end{exercise}
\noindent 来源：\href{https://github.com/yhwu-is/Linear-Algebra-Left-Undone/blob/new/%E8%AE%B2%E4%B9%89/%E4%B8%93%E9%A2%98/22%20%E4%BA%8C%E6%AC%A1%E5%9E%8B.tex}{L22-例-1}。
\begin{solve}
\begin{enumerate}
        \item 先证充分性. 若 $A^\mathrm{T}=-A$，则对任意列向量 $X$ 都有
              \[
                  (X^\mathrm{T}AX)^\mathrm{T}=X^\mathrm{T}A^\mathrm{T}X=-X^\mathrm{T}AX.
              \]
              由于 $X^\mathrm{T}AX$ 是一个 $1$ 阶矩阵，也就是一个数，它等于自身的转置，于是
              \[
                  X^\mathrm{T}AX=-X^\mathrm{T}AX,
              \]
              从而 $X^\mathrm{T}AX=0$.

              再证必要性. 若对任意 $X$ 都有 $X^\mathrm{T}AX=0$，取 $X=e_i$ 得 $a_{ii}=0$. 再取 $X=e_i+e_j\enspace(i\neq j)$，则
              \[
                  0=(e_i+e_j)^\mathrm{T}A(e_i+e_j)=a_{ii}+a_{ij}+a_{ji}+a_{jj}=a_{ij}+a_{ji}.
              \]
              故 $a_{ji}=-a_{ij}$ 对任意 $i,j$ 都成立，即 $A^\mathrm{T}=-A$.

        \item 因为 $f(X)=X^\mathrm{T}AX$ 是二次型，所以这里的 $A$ 是对称矩阵. 已知对任意 $X$ 都有 $X^\mathrm{T}AX=0$，取 $X=e_i$，得到 $a_{ii}=0$. 再取 $X=e_i+e_j\enspace(i\neq j)$，则
              \[
                  0=(e_i+e_j)^\mathrm{T}A(e_i+e_j)=a_{ii}+2a_{ij}+a_{jj}=2a_{ij}.
              \]
              故 $a_{ij}=0$. 因而 $A=O$.

        \item 由 $f(X)=g(X)$ 可知对任意 $X$ 都有
              \[
                  X^\mathrm{T}(A-B)X=0.
              \]
              又 $A,B$ 都是二次型的相伴矩阵，所以它们都是对称矩阵，从而 $A-B$ 也是对称矩阵. 由第 2 问立刻得到 $A-B=O$，即 $A=B$.
    \end{enumerate}
\end{solve}

\begin{exercise}[补充训练 L22-例-2]\label{prob:L22-例-2}
已知二次型
    \[f(X)=(x_1,x_2,x_3,x_4)\begin{pmatrix}
            1 & 2 & 3 & -4 \\ 3 & 2 & 1 & 4 \\ -4 & 3 & -7 & 2 \\ 0 & -6 & 8 & 4
        \end{pmatrix}\begin{pmatrix}
            x_1 \\ x_2 \\ x_3 \\ x_4
        \end{pmatrix}\]
    写出二次型$f(X)$的相伴矩阵.
\end{exercise}
\noindent 来源：\href{https://github.com/yhwu-is/Linear-Algebra-Left-Undone/blob/new/%E8%AE%B2%E4%B9%89/%E4%B8%93%E9%A2%98/22%20%E4%BA%8C%E6%AC%A1%E5%9E%8B.tex}{L22-例-2}。
\begin{solve}
设题中给出的矩阵为
    \[
        M=\begin{pmatrix}
            1 & 2 & 3 & -4 \\
            3 & 2 & 1 & 4 \\
            -4 & 3 & -7 & 2 \\
            0 & -6 & 8 & 4
        \end{pmatrix}.
    \]
    因为
    \[
        X^\mathrm{T}MX=X^\mathrm{T}\left(\frac{M+M^\mathrm{T}}{2}\right)X+X^\mathrm{T}\left(\frac{M-M^\mathrm{T}}{2}\right)X,
    \]
    而 $ \dfrac{M-M^\mathrm{T}}{2} $ 是反对称矩阵，故对任意 $X$ 都有
    \[
        X^\mathrm{T}\left(\frac{M-M^\mathrm{T}}{2}\right)X=0.
    \]
    因此二次型 $f(X)$ 的相伴矩阵就是 $M$ 的对称部分
    \[
        A=\frac{M+M^\mathrm{T}}{2}
        =\begin{pmatrix}
            1 & \dfrac{5}{2} & -\dfrac{1}{2} & -2 \\
            \dfrac{5}{2} & 2 & 2 & -1 \\
            -\dfrac{1}{2} & 2 & -7 & 5 \\
            -2 & -1 & 5 & 4
        \end{pmatrix}.
    \]
\end{solve}

\begin{exercise}[补充训练 L22-例-3]\label{prob:L22-例-3}
利用初等变换法将二次型 $f(x_1,x_2,x_3)=2x_1x_2-2x_1x_3+2x_2x_3$ 化为标准形，并且求出坐标变换矩阵$C$.
\end{exercise}
\noindent 来源：\href{https://github.com/yhwu-is/Linear-Algebra-Left-Undone/blob/new/%E8%AE%B2%E4%B9%89/%E4%B8%93%E9%A2%98/22%20%E4%BA%8C%E6%AC%A1%E5%9E%8B.tex}{L22-例-3}。
\begin{solve}
类似于利用初等变换求矩阵的逆，我们将相伴矩阵和单位矩阵拼接为 $\begin{pmatrix}
        A \\ E
    \end{pmatrix}$，对 $A$ 进行初等行列变换，使其变为对角矩阵，同时对 $E$ 进行对应的列变换，最终得到的 $\begin{pmatrix}
        \Lambda \\ C
    \end{pmatrix}$ 就是我们要求的结果.

    \begin{align*}
        \left(\begin{array}{c}
            A \\ \hdashline E
        \end{array}\right) ={} & \left(\begin{array}{ccc}
            0 & 1 & -1 \\
            1 & 0 & 1  \\
            -1 & 1 & 0 \\
            \hdashline
            1 & 0 & 0 \\
            0 & 1 & 0 \\
            0 & 0 & 1
        \end{array}\right)
        \xrightarrow[r_1+r_2]{c_1+c_2} \left(\begin{array}{ccc}
            2 & 1 & 0 \\
            1 & 0 & 1  \\
            0 & 1 & 0 \\
            \hdashline
            1 & 0 & 0 \\
            1 & 1 & 0 \\
            0 & 0 & 1
        \end{array}\right) & \\
        \xrightarrow[r_2-\dfrac{1}{2}r_1]{c_2-\dfrac{1}{2}c_1} & \left(\begin{array}{ccc}
            2 & 0 & 0 \\
            0 & -\dfrac{1}{2} & 1  \\
            0 & 1 & 0 \\
            \hdashline
            1 & -\dfrac{1}{2} & 0 \\
            1 & \dfrac{1}{2} & 0 \\
            0 & 0 & 1
        \end{array}\right)
        \xrightarrow[r_3+2r_2]{c_3+2c_2} \left(\begin{array}{ccc}
            2 & 0 & 0 \\
            0 & -\dfrac{1}{2} & 0  \\
            0 & 0 & 2 \\
            \hdashline
            1 & -\dfrac{1}{2} & -1 \\
            1 & \dfrac{1}{2} & 1 \\
            0 & 0 & 1
        \end{array}\right) & \\
    \end{align*}
\end{solve}

\begin{exercise}[补充训练 L22-例-4]\label{prob:L22-例-4}
用配方法把三元二次型
    \[f(x_1,x_2,x_3)=2x_1^2+3x_2^2+x_3^2+4x_1x_2-4x_1x_3-8x_2x_3\]
    化为标准形，并求所用的坐标变换$X=CY$即变换矩阵$C$.
\end{exercise}
\noindent 来源：\href{https://github.com/yhwu-is/Linear-Algebra-Left-Undone/blob/new/%E8%AE%B2%E4%B9%89/%E4%B8%93%E9%A2%98/22%20%E4%BA%8C%E6%AC%A1%E5%9E%8B.tex}{L22-例-4}。
\begin{solve}
直接配方可得
    \begin{align*}
        f(x_1,x_2,x_3)
         & = 2x_1^2+4x_1x_2-4x_1x_3+3x_2^2-8x_2x_3+x_3^2 \\
         & = 2(x_1+x_2-x_3)^2+x_2^2-4x_2x_3-3x_3^2         \\
         & = 2(x_1+x_2-x_3)^2+(x_2-2x_3)^2-5x_3^2.
    \end{align*}
    令
    \[
        \begin{cases}
            y_1 = x_1+x_2-x_3, \\
            y_2 = x_2-2x_3,    \\
            y_3 = x_3,
        \end{cases}
    \]
    则
    \[
        f(y_1,y_2,y_3)=2y_1^2+y_2^2-5y_3^2.
    \]
    由上式反解得
    \[
        \begin{cases}
            x_3 = y_3,             \\
            x_2 = y_2+2y_3,        \\
            x_1 = y_1-y_2-y_3,
        \end{cases}
    \]
    因而所求坐标变换 $X=CY$ 为
    \[
        \begin{pmatrix}
            x_1 \\
            x_2 \\
            x_3
        \end{pmatrix}
        =
        \begin{pmatrix}
            1 & -1 & -1 \\
            0 & 1 & 2   \\
            0 & 0 & 1
        \end{pmatrix}
        \begin{pmatrix}
            y_1 \\
            y_2 \\
            y_3
        \end{pmatrix}.
    \]
\end{solve}

\begin{exercise}[补充训练 L22-例-5]\label{prob:L22-例-5}
用配方法把二次型
    \[f(x_1,x_2,x_3)=2x_1x_2-2x_1x_3+2x_2x_3\]
    化为标准形，并求所用的坐标变换$X=CY$即变换矩阵$C$.
\end{exercise}
\noindent 来源：\href{https://github.com/yhwu-is/Linear-Algebra-Left-Undone/blob/new/%E8%AE%B2%E4%B9%89/%E4%B8%93%E9%A2%98/22%20%E4%BA%8C%E6%AC%A1%E5%9E%8B.tex}{L22-例-5}。
\begin{solve}
对于 $f(x_1,x_2,x_3)=2x_1x_2-2x_1x_3+2x_2x_3$，原式中没有平方项，因此不能直接像通常那样配方. 这时先作一个可逆线性变换
    \[
        \begin{cases}
            y_1 = x_1+x_2, \\
            y_2 = x_1-x_2, \\
            y_3 = x_3.
        \end{cases}
    \]
    由此反解得
    \[
        \begin{cases}
            x_1 = \dfrac{y_1+y_2}{2}, \\
            x_2 = \dfrac{y_1-y_2}{2}, \\
            x_3 = y_3.
        \end{cases}
    \]
    代入原二次型，有
    \[ f(y_1, y_2, y_3) = \frac{1}{2}y_1^2-\frac{1}{2}y_2^2-2y_2y_3. \]
    经过这一步，二次型中已经出现了平方项，于是便可以继续配方：
    \[ f(y_1, y_2, y_3) = \frac{1}{2}y_1^2-\frac{1}{2}(y_2+2y_3)^2+2y_3^2. \]
    再令
    \[
        z_1=y_1,\qquad z_2=y_2+2y_3,\qquad z_3=y_3,
    \]
    便得到标准形
    \[
        f(z_1,z_2,z_3)=\frac{1}{2}z_1^2-\frac{1}{2}z_2^2+2z_3^2.
    \]
    如果写成从新坐标到旧坐标的形式，则由
    \[
        y_1=z_1,\qquad y_2=z_2-2z_3,\qquad y_3=z_3
    \]
    以及
    \[
        x_1=\frac{y_1+y_2}{2},\qquad x_2=\frac{y_1-y_2}{2},\qquad x_3=y_3
    \]
    可得
    \[
        \begin{cases}
            x_1 = \dfrac{1}{2}z_1+\dfrac{1}{2}z_2-z_3, \\
            x_2 = \dfrac{1}{2}z_1-\dfrac{1}{2}z_2+z_3, \\
            x_3 = z_3.
        \end{cases}
    \]
    因此坐标变换矩阵可取为
    \[
        X=
        \begin{pmatrix}
            \dfrac{1}{2} & \dfrac{1}{2} & -1 \\
            \dfrac{1}{2} & -\dfrac{1}{2} & 1 \\
            0 & 0 & 1
        \end{pmatrix}
        Z.
    \]
    这说明：当原二次型没有平方项、无法直接配方时，可以先通过一个简单的可逆线性变换人为造出平方项，再继续做配方法.
\end{solve}

\begin{exercise}[补充训练 L22-例-6]\label{prob:L22-例-6}
用正交变换法将二次型 $f(x_1,x_2) = 2x_1^2 + 5x_2^2 + 4x_1x_2$ 化为标准形，并给出变换矩阵 $P$.
\end{exercise}
\noindent 来源：\href{https://github.com/yhwu-is/Linear-Algebra-Left-Undone/blob/new/%E8%AE%B2%E4%B9%89/%E4%B8%93%E9%A2%98/22%20%E4%BA%8C%E6%AC%A1%E5%9E%8B.tex}{L22-例-6}。
\begin{solve}
二次型的相伴矩阵为
    \[
        A=\begin{pmatrix}
            2 & 2 \\
            2 & 5
        \end{pmatrix}.
    \]
    其特征多项式为
    \[
        |\lambda E-A|=\begin{vmatrix}
            \lambda-2 & -2        \\
            -2        & \lambda-5
        \end{vmatrix}
        =\lambda^2-7\lambda+6=(\lambda-1)(\lambda-6),
    \]
    故特征值为 $1,6$.

    对于特征值 $1$，解方程 $(A-E)X=0$ 得特征向量可取为 $(-2,1)^\mathrm{T}$；对于特征值 $6$，解方程 $(A-6E)X=0$ 得特征向量可取为 $(1,2)^\mathrm{T}$. 将它们单位化，得到两个单位特征向量
    \[
        u_1=\frac{1}{\sqrt{5}}\begin{pmatrix}
            -2 \\
            1
        \end{pmatrix},
        \qquad
        u_2=\frac{1}{\sqrt{5}}\begin{pmatrix}
            1 \\
            2
        \end{pmatrix}.
    \]
    令
    \[
        P=\begin{pmatrix}
            -\dfrac{2}{\sqrt{5}} & \dfrac{1}{\sqrt{5}} \\
            \dfrac{1}{\sqrt{5}}  & \dfrac{2}{\sqrt{5}}
        \end{pmatrix},
    \]
    则 $P$ 为正交矩阵，且
    \[
        P^\mathrm{T}AP=\operatorname{diag}(1,6).
    \]
    因此在正交变换 $X=PY$ 下，
    \[
        f(X)=Y^\mathrm{T}(P^\mathrm{T}AP)Y=y_1^2+6y_2^2.
    \]
    这就是所求标准形.
\end{solve}

\begin{exercise}[补充训练 L22-例-7]\label{prob:L22-例-7}
确定二次型$f(x_1,x_2,\ldots,x_{10})=x_1x_2+x_3x_4+x_5x_6+x_7x_8+x_9x_{10}$的秩以及正、负惯性指数.
\end{exercise}
\noindent 来源：\href{https://github.com/yhwu-is/Linear-Algebra-Left-Undone/blob/new/%E8%AE%B2%E4%B9%89/%E4%B8%93%E9%A2%98/22%20%E4%BA%8C%E6%AC%A1%E5%9E%8B.tex}{L22-例-7}。
\begin{solve}
对每一对变量分别配方：
    \[
        x_{2k-1}x_{2k}=\frac{1}{4}\bigl((x_{2k-1}+x_{2k})^2-(x_{2k-1}-x_{2k})^2\bigr), \quad k=1,2,\ldots,5.
    \]
    因而
    \[
        f=\frac{1}{4}\sum_{k=1}^5\bigl((x_{2k-1}+x_{2k})^2-(x_{2k-1}-x_{2k})^2\bigr).
    \]
    令
    \[
        y_{2k-1}=x_{2k-1}+x_{2k}, \qquad y_{2k}=x_{2k-1}-x_{2k}, \qquad k=1,2,\ldots,5,
    \]
    则这是一个可逆线性变换，并且
    \[
        f=\frac{1}{4}(y_1^2-y_2^2+y_3^2-y_4^2+y_5^2-y_6^2+y_7^2-y_8^2+y_9^2-y_{10}^2).
    \]
    因而该二次型的相合规范形中有 $5$ 个正项和 $5$ 个负项，没有零项，所以其秩为 $10$，正惯性指数为 $5$，负惯性指数为 $5$.
\end{solve}

\begin{exercise}[补充训练 L22-例-8]\label{prob:L22-例-8}
求$t$的取值范围，使得下面的二次型是正定的：
    \[f(x_1,x_2,x_3,x_4) = x_1^2+4x_2^2+4x_3^2+3x_4^2+2tx_1x_2-2x_1x_3+4x_1x_2.\]
\end{exercise}
\noindent 来源：\href{https://github.com/yhwu-is/Linear-Algebra-Left-Undone/blob/new/%E8%AE%B2%E4%B9%89/%E4%B8%93%E9%A2%98/22%20%E4%BA%8C%E6%AC%A1%E5%9E%8B.tex}{L22-例-8}。
\begin{solve}
该二次型的相伴矩阵为
    \[
        A=\begin{pmatrix}
            1 & t+2 & -1 & 0 \\
            t+2 & 4 & 0 & 0 \\
            -1 & 0 & 4 & 0 \\
            0 & 0 & 0 & 3
        \end{pmatrix}.
    \]
    由 Sylvester 判别法，只需计算顺序主子式：
    \[
        D_1=1>0,
    \]
    \[
        D_2=\begin{vmatrix}
            1 & t+2 \\
            t+2 & 4
        \end{vmatrix}
        =4-(t+2)^2,
    \]
    \[
        D_3 = \begin{vmatrix}
                 1 & t+2 & -1 \\
                 t+2 & 4 & 0 \\
                 -1 & 0 & 4
             \end{vmatrix} = 4\bigl(3-(t+2)^2\bigr),
    \]
    \[
        D_4=3D_3=12\bigl(3-(t+2)^2\bigr).
    \]
    因而正定的充要条件是 $D_2>0,\qquad D_3>0,\qquad D_4>0$. 其中 $D_3>0$ 已经推出 $D_2>0$ 和 $D_4>0$，故最终只需 $3-(t+2)^2>0$. 即 $-2-\sqrt{3}<t<-2+\sqrt{3}$.
\end{solve}

\begin{exercise}[补充训练 L22-例-9]\label{prob:L22-例-9}
证明：每个秩为$r$的$n\enspace(r<n)$阶实对称矩阵均可表示为$n-r$个秩为$n-1$的实对称矩阵的乘积.
\end{exercise}
\noindent 来源：\href{https://github.com/yhwu-is/Linear-Algebra-Left-Undone/blob/new/%E8%AE%B2%E4%B9%89/%E4%B8%93%E9%A2%98/22%20%E4%BA%8C%E6%AC%A1%E5%9E%8B.tex}{L22-例-9}。
\begin{solve}
设 $A$ 是秩为 $r$ 的 $n$ 阶实对称矩阵，且 $r<n$. 由原讲义“实谱定理”处，存在正交矩阵 $Q$ 使得 $A=Q^\mathrm{T}\Lambda Q$，其中 $\Lambda=\operatorname{diag}(\lambda_1,\ldots,\lambda_r,0,\ldots,0)$ 且 $\lambda_1,\ldots,\lambda_r\neq 0$.

    下面将对角矩阵 $\Lambda$ 分解成 $n-r$ 个秩为 $n-1$ 的对角矩阵的乘积. 定义
    \[
        D_1=\operatorname{diag}(\lambda_1,\ldots,\lambda_r,0,1,\ldots,1),
    \]
    对于 $s=2,\ldots,n-r$，定义
    \[
        D_s=\operatorname{diag}(\underbrace{1,\ldots,1}_{r+s-1},0,1,\ldots,1).
    \]
    每个 $D_s$ 都是实对称矩阵，且恰有一个对角元为 $0$，因此 $r(D_s)=n-1$ 并且显然有 $\Lambda=D_1D_2\cdots D_{n-r}$. 令 $A_s=Q^\mathrm{T}D_sQ$（$s=1,2,\ldots,n-r$），则每个 $A_s$ 都是实对称矩阵，并且由于正交相似保持秩不变，有 $r(A_s)=n-1$. 同时
    \[
        A_1A_2\cdots A_{n-r}
        =Q^\mathrm{T}D_1Q\cdots Q^\mathrm{T}D_{n-r}Q
        =Q^\mathrm{T}(D_1D_2\cdots D_{n-r})Q
        =Q^\mathrm{T}\Lambda Q
        =A.
    \]
    因而 $A$ 可表示为 $n-r$ 个秩为 $n-1$ 的实对称矩阵的乘积.
\end{solve}

\begin{exercise}[补充训练 L22-例-10]\label{prob:L22-例-10}
求二次型
    \[
        f(x_1,x_2)=2x_1^2+5x_2^2+4x_1x_2
    \]
    在约束 $x_1^2+x_2^2=1$ 下的最大值与最小值.
\end{exercise}
\noindent 来源：\href{https://github.com/yhwu-is/Linear-Algebra-Left-Undone/blob/new/%E8%AE%B2%E4%B9%89/%E4%B8%93%E9%A2%98/22%20%E4%BA%8C%E6%AC%A1%E5%9E%8B.tex}{L22-例-10}。
\begin{solve}
该二次型的相伴矩阵为
    \[
        A=\begin{pmatrix}
            2 & 2 \\
            2 & 5
        \end{pmatrix}.
    \]
    在原讲义“正交变换法例1”处已经算出它的特征值为 $1,6$. 由于约束条件 $x_1^2+x_2^2=1$ 等价于 $\|X\|=1$，故 $f(X)=X^\mathrm{T}AX=R_A(X)$. 由原讲义“Rayleigh-Ritz”处定理立刻得到
    \[
        \min_{\|X\|=1}f(X)=1,\qquad \max_{\|X\|=1}f(X)=6.
    \]
\end{solve}

\begin{exercise}[补充训练 S-QUAD-1]\label{prob:S-QUAD-1}
设二次型 $f(x)=x^{\mathrm T}Ax$，$A=\begin{pmatrix}2&2&2\\2&0&4\\2&4&0\end{pmatrix}$。用可逆线性变量替换化为标准形，写明 $x=Py$ 的矩阵 $P$，求秩、正负惯性指数与零指数。
\end{exercise}
\noindent 编者补充：S-QUAD-1；相关主题见原讲义对应讲，题目数据与解答由本选编另行构造。
\begin{solve}
观察配方可取 $y=Cx$，其中
\[C=\begin{pmatrix}1&1&1\\0&1&-1\\0&0&1\end{pmatrix},\qquad P=C^{-1}=\begin{pmatrix}1&-1&-2\\0&1&1\\0&0&1\end{pmatrix}.\]
直接核对 $A=C^{\mathrm T}DC$，其中 $D=\operatorname{diag}(2,-2,0)$。因此 $x=Py$ 时，
\[P^{\mathrm T}AP=D,\qquad f(x)=2y_1^2-2y_2^2+0y_3^2.\]
正、负、零惯性指数依次为 $1,1,1$，秩为 $2$。因为 $\det C=1$，变换可逆，Sylvester 惯性定理保证这些指数与所选坐标无关。这里使用的是合同变换 $P^{\mathrm T}AP$；不能把它直接替换成相似变换 $P^{-1}AP$。
\end{solve}

\begin{exercise}[补充训练 S-CHOLESKY-1]\label{prob:S-CHOLESKY-1}
设 $A=\begin{pmatrix}1&1&1\\1&5&-1\\1&-1&11\end{pmatrix}$。用顺序主子式判定正定性，求对角元为正的 Cholesky 分解，并利用两次三角方程求解 $Ax=\begin{pmatrix}1\\7\\-11\end{pmatrix}$。
\end{exercise}
\noindent 编者补充：S-CHOLESKY-1；相关主题见原讲义对应讲，题目数据与解答由本选编另行构造。
\begin{solve}
三个顺序主子式依次为 $1,4,36$，均为正，故 $A$ 正定。按 Cholesky 递推逐列计算得到
\[A=LL^{\mathrm T},\qquad L=\begin{pmatrix}1&0&0\\1&2&0\\1&-1&3\end{pmatrix}.\]
先求解 $Ly=b$，再求解 $L^{\mathrm T}x=y$。这两步分别是前代和回代，得到
\[y=\begin{pmatrix}1\\3\\-3\end{pmatrix},\qquad x=\begin{pmatrix}1\\1\\-1\end{pmatrix}.\]
因为 $L$ 的三个对角元均为正且非零，两次三角方程都有唯一解；将所得 $x$ 代回原式可核验。另由 $z^{\mathrm T}Az=\|L^{\mathrm T}z\|^2$，对任意非零 $z$ 此值严格为正，这从分解角度再次证明了正定性。
\end{solve}

\part{检索索引与核验说明}

\section{按届次检索真题}\label{sec:sources}

下表记录本书实际收录范围，所有编号均为所引题解的整理序号。初赛题面对应数学A类及早期数学专业卷；决赛共同题只计一次，高低年级来源保存在逐题链接对应的目录中。未列出的题目不在本选编的100道真题之内。

\begin{longtable}{p{0.30\textwidth}p{0.23\textwidth}p{0.35\textwidth}}

\toprule 场次 & 选入来源序号 & 题目页码\\\midrule\endhead

初赛第 1 届（2009 年度） & 2, 3 & \pageref{ex:pre01-math:2}, \pageref{ex:pre01-math:3}\\

初赛第 2 届（2010 年度） & 2, 6 & \pageref{ex:pre02-math:2}, \pageref{ex:pre02-math:6}\\

初赛第 3 届（2011 年度） & 3, 6 & \pageref{ex:pre03-math:3}, \pageref{ex:pre03-math:6}\\

初赛第 4 届（2012 年度） & 4, 7 & \pageref{ex:pre04-math:4}, \pageref{ex:pre04-math:7}\\

初赛第 5 届（2013 年度） & 2, 6 & \pageref{ex:pre05-math:2}, \pageref{ex:pre05-math:6}\\

初赛第 6 届（2014 年度） & 3, 5 & \pageref{ex:pre06-math:3}, \pageref{ex:pre06-math:5}\\

初赛第 7 届（2015 年度） & 2, 3 & \pageref{ex:pre07-math:2}, \pageref{ex:pre07-math:3}\\

初赛第 8 届（2016 年度） & 2, 3 & \pageref{ex:pre08-math:2}, \pageref{ex:pre08-math:3}\\

初赛第 9 届（2017 年度） & 3, 4 & \pageref{ex:pre09-math:3}, \pageref{ex:pre09-math:4}\\

初赛第 10 届（2018 年度） & 2, 3 & \pageref{ex:pre10-math:2}, \pageref{ex:pre10-math:3}\\

初赛第 11 届（2019 年度） & 3, 4 & \pageref{ex:ch10-11th-A:3}, \pageref{ex:ch10-11th-A:4}\\

初赛第 12 届（2020 年度） & 3, 4 & \pageref{ex:ch09-12th-A:3}, \pageref{ex:ch09-12th-A:4}\\

初赛第 13 届（2021 年度） & 3, 4 & \pageref{ex:ch07-13th-A:3}, \pageref{ex:ch07-13th-A:4}\\

初赛第 13 届补赛（2021 年度） & 3, 4 & \pageref{ex:ch08-13th-A-makeup:3}, \pageref{ex:ch08-13th-A-makeup:4}\\

初赛第 14 届（2022 年度） & 3, 4 & \pageref{ex:ch04-14th-A:3}, \pageref{ex:ch04-14th-A:4}\\

初赛第 14 届补赛（2022 年度） & 3, 4 & \pageref{ex:ch05-14th-A-makeup:3}, \pageref{ex:ch05-14th-A-makeup:4}\\

初赛第 14 届第二次补赛（2022 年度） & 3, 4 & \pageref{ex:ch06-14th-A-makeup2:3}, \pageref{ex:ch06-14th-A-makeup2:4}\\

初赛第 15 届（2023 年度） & 3, 4 & \pageref{ex:ch03-15th-A:3}, \pageref{ex:ch03-15th-A:4}\\

初赛第 16 届（2024 年度） & 3, 4 & \pageref{ex:ch02-16th-A:3}, \pageref{ex:ch02-16th-A:4}\\

初赛第 17 届（2025 年度） & 3, 4 & \pageref{ex:ch01-17th-A:3}, \pageref{ex:ch01-17th-A:4}\\

决赛第 1 届（2010 年） & 4, 10, 11 & \pageref{ex:fin01-all:4}, \pageref{ex:fin01-all:10}, \pageref{ex:fin01-all:11}\\

决赛第 2 届（2011 年） & 1, 4, 6 & \pageref{ex:fin02-all:1}, \pageref{ex:fin02-all:4}, \pageref{ex:fin02-all:6}\\

决赛第 3 届（2012 年） & 5, 6 & \pageref{ex:fin03-all:5}, \pageref{ex:fin03-all:6}\\

决赛第 4 届（2013 年） & 3, 5 & \pageref{ex:fin04-all:3}, \pageref{ex:fin04-all:5}\\

决赛第 5 届（2014 年） & 2, 4, 6, 12 & \pageref{ex:fin05-hi:2}, \pageref{ex:fin05-hi:4}, \pageref{ex:fin05-lo:3}, \pageref{ex:fin05-hi:11}\\

决赛第 6 届（2015 年） & 1, 4, 6, 13 & \pageref{ex:fin06-lo:1}, \pageref{ex:fin06-lo:4}, \pageref{ex:fin06-lo:6}, \pageref{ex:fin06-hi:12}\\

决赛第 7 届（2016 年） & 1, 4, 6 & \pageref{ex:fin07-hi:1}, \pageref{ex:fin07-hi:4}, \pageref{ex:fin07-hi:6}\\

决赛第 8 届（2017 年） & 1, 2, 4, 6 & \pageref{ex:fin08-hi:1}, \pageref{ex:fin08-hi:2}, \pageref{ex:fin08-hi:4}, \pageref{ex:fin08-hi:6}\\

决赛第 9 届（2018 年） & 1, 4, 6, 8, 13 & \pageref{ex:fin09-hi:1}, \pageref{ex:fin09-hi:4}, \pageref{ex:fin09-hi:6}, \pageref{ex:fin09-lo:8}, \pageref{ex:fin09-hi:11}\\

决赛第 10 届（2019 年） & 1, 4, 6 & \pageref{ex:fin10-hi:1}, \pageref{ex:fin10-hi:4}, \pageref{ex:fin10-hi:6}\\

决赛第 11 届（2021 年） & 3, 4, 6, 8, 15 & \pageref{ex:fin11-high:3}, \pageref{ex:fin11-high:4}, \pageref{ex:fin11-high:6}, \pageref{ex:fin11-low:8}, \pageref{ex:fin11-high:13}\\

决赛第 12 届（2021 年） & 3, 4, 6 & \pageref{ex:fin12-high:3}, \pageref{ex:fin12-high:4}, \pageref{ex:fin12-high:6}\\

决赛第 13 届（2023 年） & 3, 6, 8, 16 & \pageref{ex:fin13-high:3}, \pageref{ex:fin13-high:6}, \pageref{ex:fin13-low:8}, \pageref{ex:fin13-high:13}\\

决赛第 14 届（2023 年） & 3, 5, 6, 8, 16 & \pageref{ex:fin14-high:3}, \pageref{ex:fin14-high:5}, \pageref{ex:fin14-high:6}, \pageref{ex:fin14-low:8}, \pageref{ex:fin14-high:13}\\

决赛第 15 届（2024 年） & 1, 6, 8 & \pageref{ex:fin15-high:1}, \pageref{ex:fin15-high:6}, \pageref{ex:fin15-low:8}\\

决赛第 16 届（2025 年） & 1, 6, 8 & \pageref{ex:fin16-high:1}, \pageref{ex:fin16-high:6}, \pageref{ex:fin16-low:8}\\

决赛第 17 届（2026 年） & 1, 6, 8, 16 & \pageref{ex:fin17-high:1}, \pageref{ex:fin17-high:6}, \pageref{ex:fin17-low:6}, \pageref{ex:fin17-high:13}\\

\bottomrule\end{longtable}

\section{来源定位与整理范围}

可编辑题解与编号说明见\href{https://github.com/lzxcr/cmc-mathematics/blob/64916fa1fdbc3e444a652eca204d5669058ee321/docs/COVERAGE.md}{原项目收录与编号说明}。原项目共登记37个场次，本书从中按课程内容筛选，不以原项目总题数作为本书题数。题目对应的题解源文件链接已置于题后，固定版本可用于复查本次使用的文字。

高校检索目录见\href{https://gdjpkc.xmu.edu.cn/DocumentList.aspx?cID=22&dID=383}{厦门大学课程资料页}。目录列出数学A、B类与非数学类试卷，也列出其他竞赛，故必须依据条目标题确认赛事、届次和组别。该站部分内容以旧版Flash展示，目录可用于定位，不把“目录存在”等同于“已读取全部原卷”。网络材料整理站的\href{https://www.cmathc.org.cn/shiti/}{竞赛试题栏目}仅用作辅助检索入口，本书不据域名认定其为主办方官方网站。

第13届决赛的来源目录年份为2022年，但\href{https://www.cms.org.cn/Home/news/news_details/id/1048.html}{中国数学会关于第十三届决赛的举办报道}记录活动于2023年3月25至26日在华东师范大学举行。本书将相应题目、解析和索引统一更正为2023年。第12届决赛举办于2021年5月，其记录见\href{https://www.cms.org.cn/Home/comp/comp_details/id/193.html}{中国数学会关于第十二届决赛的报道}。本书不以最初计划时间替代实际举办时间。

第14届决赛来源序号6的有限阶矩阵复合逆题，原题解末句称构造所得矩阵为实矩阵。本书明确：当扰动参数 $\varepsilon$ 取任意非零复数时，得到的是复矩阵；当 $\varepsilon$ 为实数时，构造可以在实数域内进行。这一条件区分不影响复矩阵存在性的证明。

本书以公开题解中的数学表达为起点，重排为10个课程专题，建立100组题目与解析链接，再附方法稿及课程训练。检查内容包括题答配对、去重统计、环境与括号、标签及引用、编译错误和缺字；对涉及数值计算和特殊边界的选定条目另作核验。这些检查不能替代对全部来源解答的独立数学审稿。

所有正文均包含在本TeX文件内，没有插入整本PDF、页面图片或不可编辑附件。编译依赖项目根目录的共享样式和常规宏包，支持从项目根目录、高等代数目录或当前文件夹运行XeLaTeX。目录、页码和交叉引用须至少编译两次。

\end{document}